% Semantic, multi-shade mathematical notation for the Davis--Kahan paper.
% The six core families follow visualizations/dkvis/palette.py (LIGHT).
% Subshades distinguish objects that are RELATED but not IDENTICAL.
% Values were chosen for legibility against white manuscript paper; some
% amber, green, and vermilion anchors are a little darker than the slides.
%
% Color control (independent of hover reminders):
%   \DKMathColorsfalse   (after % Semantic, multi-shade mathematical notation for the Davis--Kahan paper.
% The six core families follow visualizations/dkvis/palette.py (LIGHT).
% Subshades distinguish objects that are RELATED but not IDENTICAL.
% Values were chosen for legibility against white manuscript paper; some
% amber, green, and vermilion anchors are a little darker than the slides.
%
% Color control (independent of hover reminders):
%   \DKMathColorsfalse   (after % Semantic, multi-shade mathematical notation for the Davis--Kahan paper.
% The six core families follow visualizations/dkvis/palette.py (LIGHT).
% Subshades distinguish objects that are RELATED but not IDENTICAL.
% Values were chosen for legibility against white manuscript paper; some
% amber, green, and vermilion anchors are a little darker than the slides.
%
% Color control (independent of hover reminders):
%   \DKMathColorsfalse   (after % Semantic, multi-shade mathematical notation for the Davis--Kahan paper.
% The six core families follow visualizations/dkvis/palette.py (LIGHT).
% Subshades distinguish objects that are RELATED but not IDENTICAL.
% Values were chosen for legibility against white manuscript paper; some
% amber, green, and vermilion anchors are a little darker than the slides.
%
% Color control (independent of hover reminders):
%   \DKMathColorsfalse   (after \input{math_symbol_colors.tex} in paper.tex)
% or pass \def\DKMonochromeMath{1} before \input{paper.tex}, without editing.
% EVERY \sym, \like, and \dkrole obeys this switch, including the glossary.
% No color is required to identify a mathematical object: the glyphs, indices,
% definitions and adjacent prose also distinguish them.
%
% Families (base hue): exact/target=blue, trial/source=amber, exact unwanted=
% cyan, angles and their operators=magenta, residual/displacement=green,
% gaps=violet. Norm gauges=slate; provenance retains brown/steel/vermilion.
% P_ex is the exact-target projection in sine-theta; P is the source
% projection in transport. R denotes a residual and D a direct rotation.
% Semantic keys retain their roles even when the printed notation changes.
\newif\ifDKMathColors
\DKMathColorstrue
\ifdefined\DKMonochromeMath
  \DKMathColorsfalse
\fi

% Family anchors (print-compatible versions of the Manim light palette).
\definecolor{dkfg}{HTML}{1B1F24}
\definecolor{dkmuted}{HTML}{6B7280}
\definecolor{dkwanted}{HTML}{1F6FD1}
\definecolor{dktrial}{HTML}{A66705}
\definecolor{dkunwanted}{HTML}{007C91}
\definecolor{dksine}{HTML}{D12A7B}
\definecolor{dkresid}{HTML}{087D43}
\definecolor{dkgap}{HTML}{6E4BD8}
\definecolor{dkold}{HTML}{705519}
\definecolor{dkcurrent}{HTML}{054669}
\definecolor{dkperturb}{HTML}{BB491B}
\definecolor{dkrefuted}{HTML}{C62828}

% Related yet distinguishable shades. These are NOT new unrelated categories:
% objects inherit their parent family; a different shade signals a different
% mathematical object (e.g., a frame versus its subspace or its projection).
\definecolor{dkwantedframe}{HTML}{1254A9}
\definecolor{dkwantedprojection}{HTML}{08618D}
\definecolor{dkwantedrotation}{HTML}{4248B0}
\definecolor{dkwantedspectrum}{HTML}{345F8F}
\definecolor{dkwantedvector}{HTML}{1468AA}
\definecolor{dktrialframe}{HTML}{905600}
\definecolor{dktrialprojection}{HTML}{79511B}
\definecolor{dktrialcompetitor}{HTML}{A04D10}
\definecolor{dktrialspectrum}{HTML}{79591C}
\definecolor{dktrialvector}{HTML}{9B522A}
\definecolor{dkunwantedframe}{HTML}{006477}
\definecolor{dkunwantedspectrum}{HTML}{00546B}
\definecolor{dkunwantedprojection}{HTML}{145B6D}
\definecolor{dkunwantedvector}{HTML}{256879}
\definecolor{dkangleoperator}{HTML}{923478}
\definecolor{dksineoperator}{HTML}{AA275A}
\definecolor{dksineblock}{HTML}{B43F70}
\definecolor{dksylvester}{HTML}{804074}
\definecolor{dkanglemax}{HTML}{B32185}
\definecolor{dkanglecos}{HTML}{824276}
\definecolor{dkresidvector}{HTML}{086038}
\definecolor{dkdisplacement}{HTML}{226C50}
\definecolor{dkgaporiginal}{HTML}{4E4D89}
\definecolor{dkgapstat}{HTML}{7039A7}
\definecolor{dkgapinterval}{HTML}{574A82}
\definecolor{dkambient}{HTML}{36516A}
\definecolor{dknorm}{HTML}{475875}
\definecolor{dknormaux}{HTML}{66526E}
\definecolor{dkintertwiner}{HTML}{4A6280}
\definecolor{dkturn}{HTML}{47647D}
\definecolor{dkmix}{HTML}{505B65}

% Coloring primitives. "plain" suppresses color without suppressing hovers.
% All descendants can retain a semantic role through \like{key}{expression}.
\expandafter\def\csname dknocolor@plain\endcsname{1}
\newcommand{\dkrole}[2]{%
  \ifDKMathColors
    \ifcsname dknocolor@#1\endcsname #2\else\textcolor{dk#1}{#2}\fi
  \else #2\fi
}
% An optional plain-text reminder supplies definitions and types on hover.
% Required fields: {semantic key}{printed glyph}{color role}{glossary|local}
%                  {short definition}{mathematical type}. The "glossary" flag records whether an
% entry appears in the curated printed notation table; "local" does not.
% This registry is the sole authority for symbol meaning and printed form.
% Use role "plain" for an uncolored symbol that still carries hover text.
% Changing the glyph is a separate editorial decision: this audit preserves
% the current E, F, F_ex, G, and other existing printed choices.
\newcommand{\dkDeclareSymbol}[7][]{%
  \expandafter\def\csname dkgloss@#2\endcsname{#6}%
  \expandafter\def\csname dkglossaryflag@#2\endcsname{#5}%
  \if\relax\detokenize{#1}\relax
    \expandafter\pdfstringdef\csname dkpdfgloss@#2\endcsname{#7: #6}%
  \else
    \expandafter\pdfstringdef\csname dkpdfgloss@#2\endcsname{#7: \detokenize{#1}}%
  \fi
  \expandafter\def\csname dktype@#2\endcsname{#7}%
  \expandafter\def\csname dkglyph@#2\endcsname{#3}%
  \expandafter\def\csname dkrole@#2\endcsname{#4}%
  \expandafter\def\csname dksym@#2\endcsname{\dkrole{#4}{#3}}%
}
% Unstyled printed glyph, for variants constructed through \like.
% Example: \like{E0}{\symglyph{E0}^*} derives E_0^* from E0's
% one canonical glyph; retyping E_0 here would defeat central control.
\DeclareRobustCommand{\symglyph}[1]{\csname dkglyph@#1\endcsname}
\DeclareRobustCommand{\sym}[1]{%
  \ifcsname dksym@#1\endcsname
    \dkSymbolTooltip{#1}{\csname dksym@#1\endcsname}%
  \else
    \PackageError{dk-math-colors}{Unknown symbol key #1}{See math_symbol_colors.tex}%
  \fi
}
\DeclareRobustCommand{\like}[2]{%
  \ifcsname dkrole@#1\endcsname
    \dkSymbolTooltip{#1}{\dkrole{\csname dkrole@#1\endcsname}{#2}}%
  \else
    \PackageError{dk-math-colors}{Unknown role key #1}{See math_symbol_colors.tex}%
  \fi
}

% A short description belongs to the semantic key, not the bare letter.
% The same descriptions are available to the notation appendix.
\DeclareRobustCommand{\dkSymbolDescription}[1]{\csname dkgloss@#1\endcsname}
\DeclareRobustCommand{\dkSymbolGlossaryStatus}[1]{\csname dkglossaryflag@#1\endcsname}
\DeclareRobustCommand{\dkSymbolType}[1]{\csname dktype@#1\endcsname}

% Citable alternative notation is metadata, not a request to rename glyphs.
% The source identifier, the source's exact printed glyph, and the specific
% source location are all mandatory. No alternative is inferred from a name.
% One key may have alternatives from several works or multiple passages.
\newcommand{\dkDeclareNotationSource}[3]{%
  \expandafter\def\csname dksourcetitle@#1\endcsname{#2}%
  \expandafter\def\csname dksourceurl@#1\endcsname{#3}%
}
\newcommand{\dkDavisKahanRows}{}
\makeatletter
\newcommand{\dkDeclareNotationAlternative}[5]{%
  \ifcsname dkglyph@#1\endcsname\else
    \PackageError{dk-math-colors}{Unknown alternative key #1}{Declare the symbol first}%
  \fi
  \ifcsname dksourcetitle@#2\endcsname\else
    \PackageError{dk-math-colors}{Unknown notation source #2}{Declare the source first}%
  \fi
  \expandafter\def\csname dkaltglyph@#2@#1\endcsname{#3}%
  \expandafter\def\csname dkaltlocator@#2@#1\endcsname{#4}%
  \expandafter\def\csname dkaltnote@#2@#1\endcsname{#5}%
  \def\dk@source{#2}%
  \def\dk@original{DK1970}%
  \ifx\dk@source\dk@original
    \g@addto@macro\dkDavisKahanRows{\dkDavisKahanRow{#1}}%
  \fi
}
\makeatother
% A correspondence row is produced from the SAME registry records used by
% \sym and by the source-census script. It is never spelled a second time.
\newcommand{\dkDavisKahanRow}[1]{%
  \csname dkaltnote@DK1970@#1\endcsname &
  $\csname dkaltglyph@DK1970@#1\endcsname$ &
  $\sym{#1}$ &
  \csname dkaltlocator@DK1970@#1\endcsname\\
}

\newif\ifDKSymbolTooltips
\DKSymbolTooltipstrue
\ifdefined\DKNoSymbolTooltips\DKSymbolTooltipsfalse\fi
\newsavebox{\DKTooltipBox}
\newcount\DKTooltipCount
% Register only shipped widgets through hyperref's existing form-field
% machinery. Poppler/Okular need the catalog's field list, not just page
% annotations. The deferred aux write also discards measuring-only copies.
\makeatletter
\newcommand{\dkRegisterTooltipField}{%
  \expandafter\HyField@@AddToFields\expandafter{\number\pdffeedback lastannot}%
}
\makeatother
\AtBeginDocument{\ifDKSymbolTooltips\Form[NeedAppearances=false]\fi}
\AtEndDocument{\ifDKSymbolTooltips\endForm\fi}
% Native LuaLaTeX PDF widgets: no navigation action or JavaScript. Avoid
% pdfcomment's soulpos dependency, which previously broke local Vim builds.
% Measure in the CURRENT math style, then emit the original expression as-is:
% boxing the visible math would change spacing, scripts and line breaking.
\newcommand{\dkSymbolTooltip}[2]{%
  \ifDKSymbolTooltips
    \begingroup
      \DKSymbolTooltipsfalse % suppress nested annotations while measuring
      \ifmmode
        \edef\DKTooltipStyle{\ifcase\mathstyle
          \displaystyle\or\displaystyle\or\textstyle\or\textstyle\or
          \scriptstyle\or\scriptstyle\or\scriptscriptstyle\or\scriptscriptstyle\fi}%
        \sbox{\DKTooltipBox}{$\DKTooltipStyle#2$}%
      \else
        \leavevmode\sbox{\DKTooltipBox}{#2}%
      \fi
      \global\advance\DKTooltipCount by 1
      \pdfextension annot width\wd\DKTooltipBox
        height\ht\DKTooltipBox depth\dp\DKTooltipBox
        {/Subtype /Widget /FT /Btn /Ff 65536 /F 512
         /T (dk-symbol-\the\DKTooltipCount) /TU (\csname dkpdfgloss@#1\endcsname)
         /H /N /BS << /W 0 >> /Border [0 0 0]}%
      \dkRegisterTooltipField
      #2%
    \endgroup
  \else
    #2%
  \fi
}
\pdfstringdefDisableCommands{%
  \def\sym#1{\csname dkglyph@#1\endcsname}%
  \def\like#1#2{#2}%
  \def\symglyph#1{\csname dkglyph@#1\endcsname}%
}

% Declarations: [detailed hover], key, glyph, role, glossary|local, short text, type.
% A local entry has full semantic/hover support, but no notation-table row.
% When two keys have the same glyph, it is an auditable collision, not an
% instruction to change the printed notation automatically.
% This registry includes local proof notation. The notation appendix selects
% the symbols useful for reference; registration does not add a printed row.
% Manim's original symbol keys, now split into semantic subshades.
% Ambient and perturbed operators (provenance is independent of geometry).
\dkDeclareSymbol[Reference self-adjoint operator A on the ambient Hilbert space; the perturbed operator is Ahat = A + H. In a single-operator eigenproblem, A itself is the operator being approximated. An unbounded A acts only on domain(A).]
  {A}{A}{old}{glossary}{Reference self-adjoint operator before perturbation.}{self-adjoint-operator}
\dkDeclareSymbol{A1}{A_1}{old}{local}{Complementary block of the reference operator.}{self-adjoint-operator}
\dkDeclareSymbol{a1}{a_1}{old}{local}{Reference eigenvalue or complementary spectral value.}{scalar}
\dkDeclareSymbol[Reference block A_0: the restriction of A to the range of the reducing source projection P, expressed in source coordinates. In Theorem 8.1 it lies at or below alpha; in Theorem 8.2 its spectrum lies in the enlarged central interval. The perturbed blocks are denoted Lambda_0 and Lambda_1.]
  {ReferenceBlock0}{A_0}{old}{glossary}{Reference operator restricted to the source subspace in Section 8.}{self-adjoint-operator}
\dkDeclareSymbol[Reference block A_1: the restriction of A to the kernel of the reducing source projection P, expressed in complementary coordinates. In Theorem 8.1 it lies at or above alpha + delta in quadratic-form order.]
  {ReferenceBlock1}{A_1}{old}{glossary}{Reference operator restricted to the source complement in Section 8.}{self-adjoint-operator}
\dkDeclareSymbol[Self-adjoint ambient operator Ahat being studied; in the perturbation setting Ahat = A + H. Its action on a trial vector is compared with the proposed eigenvalue times that vector.]
  {Ahat}{\widehat A}{current}{local}{Operator under study; in the perturbation setting, the reference operator plus its perturbation.}{self-adjoint-operator}
\dkDeclareSymbol[Self-adjoint ambient operator Ahat being studied. In the perturbation setting Ahat = A + H. The trial residual is R = Ahat E_0 - E_0 A_tr, and the exact complementary block satisfies Ahat F_1 = F_1 Lambda_1 on its domain.]
  {Acurrent}{\widehat A}{current}{glossary}{Operator under study; in the perturbation setting, the reference operator plus its perturbation.}{self-adjoint-operator}
\dkDeclareSymbol[Additive self-adjoint perturbation H = Ahat - A. A bound on its norm can control the residual of an eigenspace of one operator relative to the other.]
  {H}{H}{perturb}{glossary}{Additive self-adjoint perturbation of the reference operator.}{self-adjoint-operator}
\dkDeclareSymbol{H0}{H_0}{perturb}{local}{Wanted diagonal block of the perturbation.}{self-adjoint-operator}
\dkDeclareSymbol{H1}{H_1}{perturb}{local}{Unwanted diagonal block of the perturbation.}{self-adjoint-operator}
\dkDeclareSymbol{B}{B}{perturb}{local}{Off-diagonal perturbation or coupling block.}{linear-operator}
\dkDeclareSymbol{b}{b}{perturb}{local}{Scalar perturbation or coupling coefficient.}{scalar}
% Trial/source family: subspace, frame, compressed operator, projector, vector.
\dkDeclareSymbol{U}{U}{wanted}{local}{Desired exact subspace in the local geometric notation.}{linear-subspace}% Manim's U remains wanted in local slide notation.
\dkDeclareSymbol{V}{V}{trial}{local}{Trial subspace in the local geometric notation.}{linear-subspace}% Manim's V remains trial in local slide notation.
\dkDeclareSymbol[Source subspace U to be carried onto the target V by an orthogonal or unitary map W. In the transport sections P projects onto U and Q onto V. The matching theorem also names the source subspace K.]
  {Usource}{U}{trial}{glossary}{Source subspace to be transported onto the target subspace.}{linear-subspace}
\dkDeclareSymbol[Trial frame E_0: F -> E, a linear isometry from trial-coordinate space F into ambient space E. Its range is the trial subspace; E_0* returns trial coordinates and E_0* E_0 = I_F. In rank one, E_0(x) = x v.]
  {E0}{E_0}{trialframe}{glossary}{Isometric embedding of the trial subspace; its adjoint returns trial coordinates.}{isometric-embedding}
\dkDeclareSymbol[Self-adjoint trial operator A_tr on trial-coordinate space F. The residual compares Ahat E_0 with E_0 A_tr. Rayleigh--Ritz chooses A_tr = E_0* Ahat E_0 when this compression is defined; the theorem allows other trial operators.]
  {A0}{A_{\mathrm{tr}}}{trialspectrum}{glossary}{Self-adjoint operator on trial coordinates; its eigenvalues are the trial spectral values.}{self-adjoint-operator}
\dkDeclareSymbol[Trial vector: a proposed approximate eigenvector. In the rank-one geometry v has unit length and E_0(x) = x v. Its residual for trial scalar rho is r = Ahat v - rho v.]
  {v}{v}{trialvector}{local}{Trial vector or vector in the trial subspace.}{vector}
\dkDeclareSymbol[Scalar trial operator in rank one, A_tr(x) = rho x. Figure 1 uses the Rayleigh value rho = <v, Ahat v> for unit v; then rho v is the projection of Ahat v onto the trial line and r is perpendicular to v.]
  {rho}{\rho}{trialspectrum}{local}{Trial eigenvalue or scalar trial operator.}{scalar}
\dkDeclareSymbol[Approximate eigenvalue assigned to a trial vector v. For the local operator A, its residual is A v - mu v. For unit v, the Rayleigh--Ritz choice is mu = <A v, v>.]
  {mu}{\mu}{trialspectrum}{local}{Trial spectral value.}{scalar}
\dkDeclareSymbol{a0}{a_0}{trialspectrum}{local}{Wanted trial eigenvalue or trial spectral value.}{scalar}
\dkDeclareSymbol[Ambient orthogonal or unitary transporter W with W(U) = V, equivalently W P W* = Q. Its full displacement is I - W on the whole ambient space. The real matching theorem minimizes its trace norm over every such orthogonal W.]
  {Wcompetitor}{W}{trialcompetitor}{glossary}{Orthogonal or unitary competitor carrying the source subspace onto the target.}{unitary-operator}
\dkDeclareSymbol[Orthogonal projection P onto the source subspace U in the transport problem; P = P* = P squared. A feasible orthogonal or unitary W satisfies W P W* = Q. Restricted displacement (I - W)P tests only source vectors.]
  {Psource}{P}{trialprojection}{glossary}{Orthogonal projection onto the source subspace in the transport problem.}{orthogonal-projection}
\dkDeclareSymbol{ProjSource}{P_U}{trialprojection}{local}{Orthogonal projection onto the source subspace in the transport problem.}{orthogonal-projection}
\dkDeclareSymbol[Vector x in the source subspace U. For principal vectors, unit x_j in U and y_j in V are chosen with <x_i, y_j> = 0 for i different from j and <x_j, y_j> = cos(theta_j).]
  {SourceVector}{x}{trialvector}{local}{Vector in the source subspace.}{vector}
\dkDeclareSymbol[Orthonormal coordinate frame X for the source subspace, so X* X = I. Together with a target frame Y, the singular values of X* Y are the principal cosines. In the four-dimensional counterexample X is the vertical stack of I_2 and 0.]
  {Xframe}{X}{trialframe}{glossary}{Orthonormal coordinate matrix for the source subspace in the counterexample.}{isometric-embedding}
\dkDeclareSymbol[Hilbert space F of trial coordinates. E_0 embeds F isometrically into ambient space E; A_tr and the positive sine operator sin(Theta_0) act on F, while the residual R maps F into E.]
  {Ftrial}{F}{trialframe}{glossary}{Hilbert space of trial coordinates.}{hilbert-space}
\dkDeclareSymbol[Ambient Hilbert space E over the real or complex scalars. It contains the trial subspace range(E_0) and the orthogonal exact subspaces range(F_0) and range(F_1). The source theorem permits infinite-dimensional separable E.]
  {Eambient}{E}{ambient}{glossary}{Ambient Hilbert space containing the trial and exact subspaces.}{hilbert-space}
% Zero and all-ones vectors are distinguished from scalar 0 and 1.
\dkDeclareSymbol[Zero vector of the vector or Hilbert space in context: the additive identity. The zero subspace consists only of this vector.]
  {ZeroVector}{\mathbf{0}}{fg}{glossary}{Zero vector; the additive identity of the ambient vector space.}{vector}
\dkDeclareSymbol[All-ones vector in a specified finite-dimensional coordinate space: every coordinate equals 1. This vector depends on the chosen coordinates and dimension.]
  {OnesVector}{\mathbf{1}}{fg}{local}{All-ones vector in a chosen finite-dimensional coordinate space.}{vector}
% Exact/wanted family: the target subspace, its frame, projection and rotation.
\dkDeclareSymbol[Target subspace V in the transport problem: feasible maps satisfy W(U) = V, or equivalently W P W* = Q. The matching theorem also names this subspace L.]
  {Vtarget}{V}{wanted}{glossary}{Target subspace in the transport problem.}{linear-subspace}
\dkDeclareSymbol[Exact wanted frame F_0: F_ex -> E, a linear isometry onto the desired exact spectral subspace. F_0* returns its coordinates. P_ex = F_0 F_0* is its ambient orthogonal projection; in Figure 1, F_0(h) = (h, 0).]
  {F0}{F_0}{wantedframe}{glossary}{Isometric coordinates for the desired exact spectral subspace; star denotes the adjoint.}{isometric-embedding}
\dkDeclareSymbol[Self-adjoint operator block on coordinates for the desired exact spectral subspace. Its spectrum consists of the wanted exact spectral values; in the corresponding decomposition Ahat F_0 = F_0 Lambda_0 on the operator domain.]
  {Lambda0}{\Lambda_0}{wantedspectrum}{local}{Self-adjoint block on the desired exact spectral subspace.}{self-adjoint-operator}
\dkDeclareSymbol{u}{u}{wantedvector}{local}{Vector in the desired exact subspace.}{vector}
\dkDeclareSymbol{f0}{f_0}{wantedvector}{local}{Component or coordinate vector in the desired exact subspace.}{vector}
\dkDeclareSymbol{lambda0}{\lambda_0}{wantedspectrum}{local}{Wanted exact eigenvalue or wanted spectral value.}{scalar}
\dkDeclareSymbol{lambdau}{\lambda_u}{wantedspectrum}{local}{Eigenvalue associated with a desired exact vector.}{scalar}
\dkDeclareSymbol[Exact wanted orthogonal projection P_ex: E -> E, with P_ex = F_0 F_0*. It retains the desired exact spectral component. Its complementary projection is P_ex_perp = I - P_ex = F_1 F_1*; the projected trial frame is S = P_ex_perp E_0.]
  {Pexact}{P_{\mathrm{ex}}}{wantedprojection}{glossary}{Orthogonal projection onto the desired exact subspace in the sine-theta theorem.}{orthogonal-projection}
\dkDeclareSymbol[Exact complementary orthogonal projection P_ex_perp = I - P_ex = F_1 F_1*: E -> E. It retains the unwanted exact component. In Figure 1, P_ex_perp v = S(1), whereas P_ex_perp r is the unwanted component of the residual.]
  {PexactPerp}{P_{\mathrm{ex}}^\perp}{unwantedprojection}{local}{Orthogonal projection onto the exact complementary subspace.}{orthogonal-projection}
\dkDeclareSymbol{ReflectionSigma}{\Sigma}{wantedprojection}{local}{Reflection associated with the desired exact spectral projection.}{unitary-operator}
\dkDeclareSymbol[Orthogonal projection Q onto the target subspace V in the transport problem. A feasible orthogonal or unitary W satisfies W P W* = Q, where P projects onto the source U.]
  {Qtarget}{Q}{wantedprojection}{glossary}{Orthogonal projection onto the target subspace in the transport problem.}{orthogonal-projection}
\dkDeclareSymbol{ProjTarget}{P_V}{wantedprojection}{local}{Orthogonal projection onto the target subspace in the transport problem.}{orthogonal-projection}
\dkDeclareSymbol[Canonical direct rotation D from the source U to the target V for an acute pair. It is the unitary polar factor of T = QP + (I - Q)(I - P), and acts by the principal angle on each principal plane. R is reserved for the trial residual.]
  {Rrotation}{D}{wantedrotation}{glossary}{Canonical direct rotation carrying the source subspace onto the target.}{unitary-operator}
\dkDeclareSymbol[Direct rotation D between the ranges of orthogonal projections P and Q, including the nonacute case. Relative to range(P) and kernel(P), its diagonal blocks are positive and its off-diagonal blocks are opposite adjoints. It exists when the two crossed defect spaces have equal dimension, and is unique when both defects vanish.]
  {SourceDirectRotation}{D}{wantedrotation}{local}{Direct rotation, including the nonacute case in the source-result table.}{unitary-operator}
\dkDeclareSymbol[Orthonormal coordinate frame Y for the target subspace, so Y* Y = I. In the counterexample Y is the vertical stack of cos(theta) I_2 and sin(theta) J_0; hence X* Y = cos(theta) I_2 and both principal angles equal theta.]
  {Yframe}{Y}{wantedframe}{glossary}{Orthonormal coordinate matrix for the target subspace in the counterexample.}{isometric-embedding}
\dkDeclareSymbol[Vector y in the target subspace V. For principal vectors, unit x_j in U and y_j in V are paired so that <x_j, y_j> = cos(theta_j); distinct pairs have zero cross inner product.]
  {TargetVector}{y}{wantedvector}{local}{Vector in the target subspace.}{vector}
\dkDeclareSymbol[Hilbert space F_ex of exact wanted coordinates. The isometry F_0: F_ex -> E identifies it with the desired exact spectral subspace. This space differs from H, the additive perturbation.]
  {Hexact}{F_{\mathrm{ex}}}{wantedframe}{glossary}{Hilbert space of coordinates for the desired exact spectral subspace.}{hilbert-space}% Coordinate Hilbert space (not perturbation).
% Exact/unwanted family; the trial error operators are MAGENTA, not cyan.
\dkDeclareSymbol{Uperp}{U^{\perp}}{unwanted}{local}{Orthogonal complement of the desired exact subspace in the local notation.}{linear-subspace}
\dkDeclareSymbol[Exact complementary frame F_1: G -> E, a linear isometry onto the orthogonal complement of range(F_0). F_1* returns unwanted coordinates; F_1 F_1* = I - P_ex. In Figure 1, F_1(g) = (0, g).]
  {F1}{F_1}{unwantedframe}{glossary}{Isometric coordinates for the exact complementary subspace; star denotes the adjoint.}{isometric-embedding}
\dkDeclareSymbol[Self-adjoint unwanted block Lambda_1 on G, satisfying Ahat F_1 = F_1 Lambda_1 on its domain. Its spectrum is separated from the trial spectrum of A_tr by the specified positive delta. Figure 1 uses the scalar block 0.7.]
  {Lambda1}{\Lambda_1}{unwantedspectrum}{glossary}{Self-adjoint operator on exact complementary coordinates: the unwanted spectral block.}{self-adjoint-operator}
% Spectral measure and set (distinct from the ambient space E and sine block S).
\dkDeclareSymbol[Projection-valued spectral measure E associated with a self-adjoint operator A: A = integral lambda dE(lambda). For a Borel set S, E(S) is the orthogonal projection onto the corresponding spectral subspace.]
  {SpectralMeasure}{E}{wantedprojection}{glossary}{Projection-valued spectral measure of a self-adjoint operator.}{projection-valued-measure}
\dkDeclareSymbol[Borel subset S of the real spectral axis. The spectral projection E(S) selects the part of a self-adjoint operator with spectral values in S; its range is a reducing subspace.]
  {SpectralSet}{S}{wantedspectrum}{glossary}{Borel spectral set used to select a spectral subspace.}{set}
% Generic Sylvester coefficients and right side, by source/target role.
\dkDeclareSymbol[Left coefficient B of the Sylvester map X -> B X - X C; B acts on the output space of X. For the sine bound, B = Lambda_1 on the unwanted-coordinate space G.]
  {SylvesterLeft}{B}{unwantedspectrum}{glossary}{Left coefficient of the Sylvester equation; corresponds to the unwanted spectral block.}{self-adjoint-operator}
\dkDeclareSymbol[Right coefficient C of the Sylvester map X -> B X - X C; C acts on the input space of X. For the sine bound, C = A_tr on trial space F. The historical equation also denotes this coefficient by A.]
  {SylvesterRight}{C}{trialspectrum}{glossary}{Right coefficient of the Sylvester equation; corresponds to the trial operator.}{self-adjoint-operator}
\dkDeclareSymbol[Right-hand side Y in B X - X C = Y, with the same input and output spaces as X. For the sine bound, Y = F_1* R: F -> G, the residual expressed in unwanted exact coordinates.]
  {SylvesterRHS}{Y}{resid}{glossary}{Right-hand side of the Sylvester equation, obtained by projecting the residual.}{linear-operator}
\dkDeclareSymbol{w}{w}{unwantedvector}{local}{Vector in an unwanted or complementary spectral subspace.}{vector}
\dkDeclareSymbol{f1}{f_1}{unwantedvector}{local}{Component or coordinate vector in the unwanted exact subspace.}{vector}
\dkDeclareSymbol{lambda1}{\lambda_1}{unwantedspectrum}{local}{Unwanted exact eigenvalue or complementary spectral value.}{scalar}
\dkDeclareSymbol{lambdaw}{\lambda_w}{unwantedspectrum}{local}{Eigenvalue associated with an unwanted exact vector.}{scalar}
\dkDeclareSymbol[Hilbert space G of exact complementary coordinates. F_1: G -> E embeds it onto the unwanted exact subspace; Lambda_1 acts on G and X = F_1* E_0 maps trial coordinates into G.]
  {Gunwanted}{G}{unwantedframe}{glossary}{Hilbert space of exact complementary coordinates.}{hilbert-space}
\dkDeclareSymbol{E1}{E_1}{unwantedframe}{local}{Isometric coordinates for the complementary trial subspace.}{isometric-embedding}
\dkDeclareSymbol{Vperp}{V^{\perp}}{unwanted}{local}{Orthogonal complement of the trial subspace in the local notation.}{linear-subspace}
% Angle and error geometry, including distinct rectangular and positive maps.
\dkDeclareSymbol[Individual principal angle, between 0 and pi/2. For orthonormal subspace frames X and Y, the singular values of X* Y are the principal cosines. In Figure 1, theta is the angle between the unit trial vector v and the exact wanted line.]
  {theta}{\theta}{sine}{glossary}{Principal angle between the subspaces under comparison.}{angle}
\dkDeclareSymbol[Full positive angle operator Theta = diag(Theta_0, Theta_1) for the source/target decompositions. Theta_0 compares the source with the target; Theta_1 compares their orthogonal complements. Its operator norm is the largest spectral angle. For positive c, the conditions Theta <= c and Theta < c mean operator norm at most c and strictly below c, respectively.]
  {Theta}{\Theta}{angleoperator}{glossary}{Full angle operator, the diagonal sum of the source and complementary angle blocks.}{angle-operator}
\dkDeclareSymbol[Directed principal-angle operator Theta_0 on source or trial coordinates. For a trial frame E_0 and exact wanted projection P_ex, cos(Theta_0) = sqrt(E_0* P_ex E_0), and sin(Theta_0) = sqrt(I_F - E_0* P_ex E_0). In the two-projection geometry it comes from the positive cosine compression P Q P restricted to range(P).]
  {Theta0}{\Theta_0}{angleoperator}{glossary}{Directed principal-angle operator acting on source or trial coordinates.}{angle-operator}
\dkDeclareSymbol[Complementary directed angle operator Theta_1, obtained from the projections I-P and I-Q by the same cosine-compression construction as Theta_0. It acts on the source complement, or its isometric coordinates. The full angle operator is diag(Theta_0, Theta_1).]
  {Theta1}{\Theta_1}{angleoperator}{glossary}{Directed angle operator comparing the orthogonal complements.}{angle-operator}
\dkDeclareSymbol[Right-angle source defect D_0 = range(P) intersect kernel(Q). These are source vectors orthogonal to the target. The full-angle conclusion in the Section 8 smallness criterion assumes dim(D_0) = dim(D_1); in Lean their identification is CrossedDefectsEquivalent.]
  {CrossedSourceDefect}{D_0}{fg}{glossary}{Source directions perpendicular to the target subspace.}{linear-subspace}
\dkDeclareSymbol[Right-angle target defect D_1 = kernel(P) intersect range(Q). These are target vectors orthogonal to the source. Equality of its Hilbert dimension with that of D_0 is assumed for the full-angle conclusion of Theorem 8.2.]
  {CrossedTargetDefect}{D_1}{fg}{glossary}{Target directions perpendicular to the source subspace.}{linear-subspace}
\dkDeclareSymbol[Positive complementary cosine block C_1 = cos(Theta_1). The reference complement is identified with the perturbed complement by the direct rotation D, so A_1 and Lambda_1 act on common coordinates. Theorem 8.1 compares A_1 - alpha I with C_1 (Lambda_1 - alpha I) C_1 as quadratic forms.]
  {ComplementCosineBlock}{C_1}{anglecos}{local}{Cosine of the complementary angle block in the spectral-repulsion inequality.}{self-adjoint-operator}
\dkDeclareSymbol[Scalar sine of a principal angle. For a unit principal trial vector v and exact wanted projection P_ex, sin(theta) = ||(I - P_ex)v||. In Figure 1 it is the length of the pink leg S(1).]
  {sintheta}{\sin\theta}{sine}{local}{Sine of an individual principal angle: the length of the complementary component of a unit principal trial vector.}{scalar}
\dkDeclareSymbol[Positive sine operator sin(Theta_0) = sqrt(I_F - E_0* P_ex E_0): F -> F. It equals |S| for S = (I - P_ex)E_0: F -> E. Its eigenvalues in finite dimensions are the principal sines; the sine theorem bounds its ideal norm by N(R)/delta.]
  {sinTheta0}{\sin\Theta_0}{sineoperator}{glossary}{Positive sine operator on trial coordinates; the modulus of the rectangular complementary error.}{angle-operator}
\dkDeclareSymbol[Rectangular complementary error S = (I - P_ex) E_0: F -> E. It records how much each trial vector leaks outside the exact wanted subspace. Its modulus |S| = sqrt(S* S) is sin(Theta_0): F -> F. In rank one S(1) = P_ex_perp v.]
  {S}{S}{sineblock}{glossary}{Rectangular complementary error: the trial frame projected onto the unwanted exact subspace.}{linear-operator}
\dkDeclareSymbol[Complementary-coordinate map X = F_1* E_0: F -> G, with S = F_1 X. It satisfies Lambda_1 X - X A_tr = F_1* R on domain(A_tr). Isometric multiplication gives N(X) = N(S); the separate modulus identity abs(S) = sin(Theta_0) gives N(S) = N(sin(Theta_0)), whenever these norms are finite.]
  {X}{X}{sylvester}{glossary}{Solution of the Sylvester equation; for the sine bound, the trial frame in unwanted exact coordinates.}{linear-operator}
\dkDeclareSymbol{alphaAngle}{\alpha}{anglemax}{local}{Largest principal angle.}{angle}
\dkDeclareSymbol[Scalar c = cos(theta) in the counterexample. In the general two-projection decomposition, the same role denotes a positive cosine operator C; the cosine and sine blocks commute and satisfy C squared + S squared = I.]
  {CosAngle}{c}{anglecos}{local}{Cosine of a principal angle or the corresponding positive cosine block.}{scalar-or-operator}
\dkDeclareSymbol[Scalar s = sin(theta) in the counterexample. In the general two-projection decomposition, the same role denotes a positive sine operator S; it commutes with the positive cosine block C and satisfies C squared + S squared = I.]
  {SinAngle}{s}{sineoperator}{local}{Sine of a principal angle or the corresponding positive sine block.}{scalar-or-operator}
\dkDeclareSymbol[Positive small-angle parameter in the local subspace comparison. It controls the allowed angular deviation; it is an angle parameter rather than a perturbation-operator norm.]
  {EpsilonAngle}{\varepsilon}{anglemax}{local}{Angle parameter controlling the maximum subspace deviation.}{angle}
\dkDeclareSymbol{PhiAngle}{\varphi}{angleoperator}{local}{Auxiliary angle or angle operator in the two-subspace comparison.}{angle}
\dkDeclareSymbol{Bblock}{B}{sineblock}{local}{Off-diagonal block in the two-subspace decomposition.}{linear-operator}
\dkDeclareSymbol{Cblock}{C}{sylvester}{local}{Complementary or Sylvester block in the two-subspace decomposition.}{linear-operator}
\dkDeclareSymbol{Ablock}{A}{anglecos}{local}{Positive cosine-related block in the two-subspace decomposition.}{linear-operator}
% Residual/error-displacement family.
\dkDeclareSymbol[Residual vector r = Ahat v - rho v = R(1) for the rank-one frame E_0(x) = x v. Figure 1 draws it from rho v to Ahat v. For unit v and the Rayleigh value rho = <v, Ahat v>, <v, r> = 0; a general trial scalar need not give orthogonality.]
  {r}{r}{residvector}{local}{Residual vector of an approximate eigenpair.}{vector}
\dkDeclareSymbol[Bounded residual R: F -> E, agreeing with Ahat E_0 - E_0 A_tr on domain(A_tr). The domain condition requires E_0(domain(A_tr)) inside domain(Ahat). R = 0 means that E_0 intertwines A_tr with Ahat; its norm measures the defect.]
  {R}{R}{resid}{glossary}{Bounded residual of the trial eigenproblem, agreeing with the operator difference on the trial domain.}{linear-operator}
\dkDeclareSymbol[Full displacement Delta_W = I - W on the ambient space. In the failed Ky Fan comparison W = W_(pi/4); Delta_W has singular values sqrt(2), sqrt(2), 0, 0. Multiplying by a principal-plane projection restricts its input to that plane.]
  {Kdisplacement}{\Delta_W}{displacement}{glossary}{Full displacement of a transporter from the identity.}{operator-valued-function}
\dkDeclareSymbol[Squared full displacement Dcal(W) = (I-W)* (I-W), a positive operator. Proposition 4.3 compares its unitarily invariant norm for a direct rotation D with that for every feasible transporter W.]
  {SquaredDisplacement}{\mathcal D}{displacement}{local}{Squared full displacement as a function of the transporter.}{operator-valued-function}
% Matching proof: the gradient and its spectral level are not a space or gap.
\dkDeclareSymbol[Skew-symmetric spectral gradient G of Phi_p. It commutes with W and represents first variations by d/dt Phi_p(W exp(tB)) at t = 0 equal to -tr(G B) for skew-symmetric B. Feasible stationarity gives J_src G J_src = -G.]
  {SpectralGradient}{G}{displacement}{local}{Skew-symmetric spectral gradient of the displacement power sum in the matching proof.}{linear-operator}
\dkDeclareSymbol[Positive spectral level g of the skew-symmetric gradient G: on this invariant subspace, -G squared = g squared I. Consequently G/g squares to -I and supplies the complex structure used in the real matching proof.]
  {GradientLevel}{g}{displacement}{local}{Positive gradient level; the gradient squared is minus this level squared times the identity.}{scalar}
\dkDeclareSymbol[Complex structure I_g = G/g on the positive gradient-level space V_g. It is a real isometry with square -I. Multiplication by a+ib means a times identity plus b times I_g; source reflection J_src reverses I_g and acts as conjugation.]
  {GradientComplexStructure}{\mathcal I_g}{displacement}{local}{Complex multiplication constructed from a positive gradient level.}{linear-operator}
\dkDeclareSymbol[Real skew-symmetric matrix B, with B transpose = -B, giving the orthogonal variation W exp(tB). If B commutes with the source projection P, this curve preserves the constraint W(U) = V.]
  {SkewDirection}{B}{ambient}{local}{Skew-symmetric direction for a first variation of an orthogonal transporter.}{linear-operator}
% Spectral separation family: gap and original gap are NOT identical.
\dkDeclareSymbol[Positive spectral separation in the stated configuration. In the sine theorem it separates the trial operator from the exact complementary block; Theorem 8.1 uses it for the ordered reference gap, and Theorem 8.2 for the perturbed central-band separation. Figure 1 uses delta = abs(rho - 0.7).]
  {delta}{\delta}{gap}{glossary}{Positive spectral separation in the stated gap configuration.}{scalar}
\dkDeclareSymbol{g}{g}{gaporiginal}{local}{Original spectral gap of the reference operator; distinct from the certified sine-theta separation.}{scalar}
\dkDeclareSymbol{Delta}{\Delta}{gapstat}{local}{Population or statistical spectral gap.}{scalar}
\dkDeclareSymbol[Spectral endpoint alpha. In the interval/exterior configuration, one spectrum lies between beta and alpha, the other at or below beta - delta or at or above alpha + delta. Theorem 8.2 places the selected perturbed block in the central interval. In Theorem 8.1, alpha bounds the lower reference block and alpha + delta bounds the upper block from below.]
  {alphaEndpoint}{\alpha}{gapinterval}{glossary}{Upper endpoint of the central band, or lower-block cut in the ordered gap.}{scalar}
\dkDeclareSymbol[Lower endpoint beta of the finite interval from beta to alpha in the gap hypothesis. One spectrum lies in that interval; the other lies at or below beta - delta or at or above alpha + delta.]
  {betaEndpoint}{\beta}{gapinterval}{glossary}{Lower endpoint of the finite interval in the interval/exterior gap hypothesis.}{scalar}
% Spectral/projector and auxiliary matrix objects, context-scoped.
\dkDeclareSymbol{P}{P}{fg}{local}{Orthogonal projection; its source or exact-target role is specified locally.}{orthogonal-projection}% Neutral historical Manim key; prefer Pexact/Psource.
\dkDeclareSymbol{Q}{Q}{fg}{local}{Orthogonal projection; its role is specified locally.}{orthogonal-projection}% Neutral historical Manim key; prefer Qtarget.
\dkDeclareSymbol[Two-projection intertwiner T = QP + (I - Q)(I - P), satisfying TP = QT and T* T = I - (P - Q) squared. For ||P - Q|| < 1 it is invertible; its unitary polar factor T |T| inverse is the canonical direct rotation.]
  {IntertwinerT}{T}{intertwiner}{glossary}{Two-projection intertwiner; its unitary polar factor is the acute direct rotation.}{linear-operator}
\dkDeclareSymbol[Planar pi/2-rotation matrix J_0 = ((0, -1), (1, 0)), with J_0 squared = -I_2 and J_0 transpose = -J_0. It defines the rotation C = cos(2 theta) I_2 + sin(2 theta) J_0 used to construct the real competitor.]
  {Jturn}{J_0}{turn}{glossary}{Planar rotation through $\pi/2$, with square minus the identity.}{matrix}
\dkDeclareSymbol[Partial isometry V_pol in the polar decomposition S = V_pol |S|, where |S| = sqrt(S* S). It is isometric on the orthogonal complement of kernel(S) and zero on kernel(S). For an invertible operator its polar factor is unitary.]
  {PolarU}{V_{\mathrm{pol}}}{intertwiner}{glossary}{Partial isometry in the polar decomposition; unitary when the operator is invertible.}{partial-isometry}
\dkDeclareSymbol[Planar rotation C = cos(2 theta) I_2 + sin(2 theta) J_0 through twice the principal angle. The competitor W_theta is orthogonally conjugate to C direct-sum I_2, so it rotates one two-plane and fixes the other.]
  {Cturn}{C}{turn}{glossary}{Planar rotation through twice the principal angle in the counterexample.}{matrix}
\dkDeclareSymbol[Orthogonal mixing matrix H_mix = (1/sqrt(2)) times the block matrix ((I_2, I_2), (I_2, -I_2)). It constructs W_theta = H_mix (C direct-sum I_2) H_mix transpose.]
  {Hmix}{H_{\mathrm{mix}}}{mix}{glossary}{Orthogonal mixing matrix in the counterexample.}{matrix}
\dkDeclareSymbol[Source reflection J_src = 2P - I: it acts as +I on the source subspace and -I on its orthogonal complement, so its square is I. In the matching proof stationarity gives J_src G J_src = -G for the skew-symmetric spectral gradient G.]
  {Jreflect}{J_{\mathrm{src}}}{trialprojection}{glossary}{Reflection equal to twice the source projection minus the identity.}{linear-operator}
\dkDeclareSymbol[Skew angle generator J_ang for a direct rotation D = exp(J_ang Theta). On each principal plane it is the oriented rotation through pi/2; it vanishes on the common fixed directions. It commutes with the full angle operator Theta.]
  {AngleGenerator}{J_{\mathrm{ang}}}{turn}{glossary}{Skew angle generator in the direct-rotation result summaries.}{linear-operator}
\dkDeclareSymbol{Dblock}{D}{mix}{local}{Auxiliary diagonal or block matrix in the counterexample.}{matrix}
\dkDeclareSymbol[Orthogonal projection onto the first principal plane span(e_0, e_3) in the four-dimensional counterexample. For Delta_W = I - W_(pi/4), the Ky Fan 2-norm of Delta_W Omega_1 equals 2.]
  {OmegaOne}{\Omega_1}{turn}{glossary}{First principal-plane projection in the failed Ky Fan comparison.}{orthogonal-projection}
\dkDeclareSymbol[Orthogonal projection onto the second principal plane span(e_1, e_2) in the four-dimensional counterexample. For Delta_W = I - W_(pi/4), the Ky Fan 2-norm of Delta_W Omega_2 equals 2, although the trace norm of Delta_W is only 2 sqrt(2).]
  {OmegaTwo}{\Omega_2}{intertwiner}{glossary}{Second principal-plane projection in the failed Ky Fan comparison.}{orthogonal-projection}
\dkDeclareSymbol{Dplanar}{D}{wantedrotation}{local}{Planar direct-rotation block.}{unitary-operator}
\dkDeclareSymbol{Hreflection}{H}{trialcompetitor}{local}{Planar reflection carrying the source line to the target.}{linear-operator}
% Unitarily invariant gauge family; N, M are NOT a norm versus its input.
\dkDeclareSymbol[Unitarily invariant norm N, or a symmetric operator-ideal norm on its domain of finiteness. It satisfies N(U T V) = N(T) for unitaries and N(B T) <= ||B||op N(T). In the sine theorem it is normalized by N(u v*) = ||u|| ||v|| and has Ky Fan dominance where both norms are finite. Both R and sin(Theta_0) must belong to its ideal.]
  {NormN}{N}{norm}{glossary}{Unitarily invariant norm on its domain of finiteness; in the sine theorem, normalized by $N(uv^*)=\|u\|\|v\|$.}{norm}
\dkDeclareSymbol[Sine-theta separation G_delta(B,C): finite interval/exterior separation in either spectral order, or ordered separation B <= a I and C >= (a+delta) I, or the reverse. The ordered alternatives allow both spectra to be unbounded.]
  {SineGap}{\mathcal G_\delta}{gap}{local}{Interval/exterior separation or either ordered separation, as in the sine-theta theorem.}{predicate}
\dkDeclareSymbol[Auxiliary unitarily invariant norm M witnessing the Q-norm property N(T) squared = M(T* T). For the Frobenius norm N, one can take M to be the trace norm.]
  {NormM}{M}{normaux}{local}{Auxiliary unitarily invariant norm used to define a Q-norm through the squared modulus.}{norm}
\dkDeclareSymbol[Q-norm property: a unitarily invariant norm N is a Q-norm when N(T) squared = M(T* T) for some unitarily invariant norm M. The Frobenius norm is an example, with M the trace norm; the repaired full-displacement theorem uses this norm class.]
  {NormQ}{Q}{normaux}{glossary}{Q-norm: its square is a unitarily invariant norm of the squared modulus.}{norm}
\dkDeclareSymbol[Singleton trace cost s_i = min(4 sin(theta_i/2), 2) for angle label i. It transports one real principal plane independently, choosing a rotation of cost 4 sin(theta_i/2) or a reflection of cost 2.]
  {SingletonCost}{s}{norm}{glossary}{Minimum trace displacement for one real principal plane: choose a rotation or a reflection.}{scalar-cost}
\dkDeclareSymbol[Paired trace cost e_ij = 4 sin((theta_i + theta_j)/2) for distinct angle labels i,j. It is attained by coupling their principal planes in a real four-dimensional block. The edge weight s_i + s_j - e_ij is the difference between the cost of two singletons and this paired cost.]
  {PairCost}{e}{norm}{glossary}{Trace displacement cost of coupling two principal planes in a four-dimensional block.}{scalar-cost}
\dkDeclareSymbol[Scalar phase cost f_p(t) = (2 sin(t/2)) to the power p, for 0 <= t <= pi. A real planar rotation through t contributes two copies, so its displacement power cost is 2 f_p(t).]
  {MatchingPhaseCost}{f_p}{norm}{local}{Displacement power cost for one scalar phase.}{scalar-cost}
\dkDeclareSymbol[Singleton power cost a_i(p) = min(2 f_p(theta_i), 2 to the power p). These are the rotation and reflection choices for one principal plane. At p=1 it equals the trace-cost weight s_i.]
  {SingletonPowerCost}{a}{norm}{local}{Minimum displacement power cost for one angle.}{scalar-cost}
\dkDeclareSymbol[Paired power cost b_ij(p) = min over d <= tau <= h of 2(f_p(h+tau) + f_p(h-tau)), where h=(theta_i+theta_j)/2 and d=abs(theta_i-theta_j)/2. This is the positive-determinant pair minimum; unrestricted block optimization also compares two singletons. At p=1 it equals e_ij.]
  {PairPowerCost}{b}{norm}{local}{Positive-determinant displacement power cost for a pair of angles.}{scalar-cost}
\dkDeclareSymbol[Matching M on the finite general graph of positive principal-angle labels. Each label belongs to at most one edge; unmatched labels use singleton transporters. Minimizing paired and singleton costs is equivalent to maximum-weight matching with edge weights g_ij = s_i + s_j - e_ij. The empty matching is allowed.]
  {Matching}{M}{normaux}{glossary}{Matching of positive principal-angle labels: each label is used in at most one pair.}{matching}
\dkDeclareSymbol[Displacement power sum Phi_p(W) = sum_j s_j(I - W)^p. The matching proof first minimizes it for 1 < p < 2, then lets p decrease to 1, where it becomes the trace norm ||I - W||_1.]
  {PowerCost}{\Phi}{norm}{local}{Sum of the p-th powers of displacement singular values.}{scalar-cost}
% Auxiliary matching-proof notation: neutral ink, with the same hover registry.
\dkDeclareSymbol[Off-diagonal gradient block G_01 mapping the source complement K perpendicular into K. Stationarity makes the diagonal blocks of G zero, so G = ((0,G_01),(-G_01 transpose,0)) in the source/complement decomposition.]
  {GradientOffDiagonal}{G_{01}}{fg}{local}{Complement-to-source block of the skew-symmetric gradient.}{linear-operator}
\dkDeclareSymbol[Orthogonal group O(n): real n by n matrices W satisfying W transpose times W = I. Its elements preserve Euclidean inner products and lengths. The notation O(2) refers to the same group in dimension two.]
  {OrthogonalGroup}{O(n)}{fg}{local}{Group of real orthogonal matrices.}{group}
\dkDeclareSymbol[Feasible transporters F = {W in O(n) : W(K) = L}. Equivalently W P_K W transpose = P_L. This nonempty compact set is the optimization domain throughout the matching proof.]
  {FeasibleTransport}{\mathcal F}{fg}{local}{Orthogonal maps carrying the source subspace onto the target.}{set}
\dkDeclareSymbol[Singular value s_j(T), the nonnegative square root of an eigenvalue of T transpose times T, counted with multiplicity. Summing s_j(I-W) gives the trace displacement; summing its p-th powers gives Phi_p(W). These singular values differ from the singleton weights s_i.]
  {TransportSingularValues}{\symglyph{SingularValueBase}_j}{fg}{local}{A displacement singular value in the matching proof.}{scalar}
\dkDeclareSymbol[Scalar singular-value magnitude s between zero and two for the displacement I-W of an orthogonal matrix. Uniform convergence of s to the power p to s on this interval permits the trace-norm limit. This variable differs from the rotation parameter s used earlier in the proof.]
  {DisplacementSize}{s}{fg}{local}{Scalar displacement singular-value magnitude in the uniform limit.}{scalar}
\dkDeclareSymbol[All-singleton cost C_0 = sum_i s_i, before any two angle labels are paired. A matching has cost C_0 minus the sum of its selected edge weights.]
  {SingletonBaseline}{C_0}{fg}{local}{Trace displacement when every angle uses its singleton transporter.}{scalar-cost}
\dkDeclareSymbol[Matching power cost C_p(M): sum of b_ij(p) over selected pairs plus a_i(p) over unmatched labels. Its minimum equals the minimum of Phi_p over feasible transporters. At p=1 it is C(M).]
  {MatchingPowerCost}{C_p}{fg}{local}{Total paired and singleton power cost of a matching.}{scalar-cost}
\dkDeclareSymbol[Invariant orthogonal block V_a preserved by the source projection, target projection, and minimizing transporter. The displacement cost adds over these blocks, and their principal-angle lists combine into the original list.]
  {MatchingBlock}{V_a}{fg}{local}{Simultaneous invariant block in the matching proof.}{linear-subspace}
\dkDeclareSymbol[Gradient-level space V_g: the eigenspace of -G squared for a positive eigenvalue g squared. It is preserved by G, W, and the source reflection J_src. Multiplication by i on it is G/g.]
  {GradientLevelSpace}{V_g}{fg}{local}{Positive gradient-level eigenspace.}{linear-subspace}
\dkDeclareSymbol[Eigenphase varphi of W on V_g viewed as a complex space with multiplication by i = G/g. It lies between zero and pi and satisfies f_p prime(varphi) = g. There are at most two such phases.]
  {GradientEigenphase}{\varphi}{fg}{local}{One of the at most two phases at a positive gradient level.}{angle}
\dkDeclareSymbol[Orthogonal projection E onto the first phase space of W on V_g. Here W = exp(i varphi_1) E + exp(i varphi_2) (I-E). In coordinates where the source reflection is conjugation, E = X + iY.]
  {GradientPhaseProjection}{E}{fg}{local}{Projection onto the first phase space in a gradient-level space.}{orthogonal-projection}
\dkDeclareSymbol[Real part X of the phase projection E = X + iY. It is a real symmetric matrix with spectrum in the closed interval from zero to one, satisfying X squared - Y squared = X and XY + YX = Y. This X differs from the source-coordinate and Sylvester matrices elsewhere in the paper.]
  {ProjectionRealPart}{X}{fg}{local}{Real symmetric part of the phase projection.}{linear-operator}
\dkDeclareSymbol[Imaginary part Y of the phase projection E = X + iY. It is a real skew-symmetric matrix. It maps the lambda eigenspace of X to the 1-lambda eigenspace and constructs the small invariant planes.]
  {ProjectionImaginaryPart}{Y}{fg}{local}{Real skew-symmetric imaginary part of the phase projection.}{linear-operator}
\dkDeclareSymbol[Eigenvalue lambda in the closed interval from zero to one of the real part X of the phase projection. It is local to the block construction and differs from the original operator's spectral parameter.]
  {ProjectionEigenvalue}{\lambda}{fg}{local}{Eigenvalue of the real part of the phase projection.}{scalar}
\dkDeclareSymbol[Plane coupling eta = sqrt(lambda (1-lambda)) for 0 < lambda < 1. If u is a unit lambda eigenvector of X, then v = Yu/eta completes its invariant real plane.]
  {ProjectionCoupling}{\eta}{fg}{local}{Normalization used to construct an invariant real plane.}{scalar}
\dkDeclareSymbol[Pair mean angle h = (theta_i + theta_j)/2. Together with d = (theta_i - theta_j)/2 after ordering theta_i >= theta_j, it parametrizes the two-angle block cost.]
  {PairMeanAngle}{h}{fg}{local}{Mean of the two ordered principal angles.}{angle}
\dkDeclareSymbol[Pair half-difference d = (theta_i - theta_j)/2 after ordering theta_i >= theta_j. It satisfies 0 <= d <= h and h+d <= pi/2.]
  {PairHalfDifference}{d}{fg}{local}{Half the difference of the two ordered principal angles.}{angle}
\dkDeclareSymbol[Scalar pair parameter tau in the closed interval from d to h. The paired power cost is the minimum of 2(f_p(h+tau) + f_p(h-tau)) over this interval. At p=1 a minimum is attained at tau=h.]
  {PairPhaseParameter}{\tau}{fg}{local}{Scalar parameter in the two-angle power-cost minimum.}{scalar}
\dkDeclareSymbol[Attaining rotation angle gamma = arccos(cos(tau)/cos(d)). Choosing U_0 = U_1 = R_gamma realizes phase parameters x=h and y=tau in the four-dimensional block.]
  {PairAttainingAngle}{\gamma}{fg}{local}{Rotation angle realizing a chosen scalar pair parameter.}{angle}
\dkDeclareSymbol[Four-dimensional direct rotation D = ((C,-S),(S,C)) in source and complementary principal coordinates, with C the diagonal principal cosines and S the diagonal principal sines. Every transporter on this block is D diag(U_0,U_1).]
  {PairDirectRotation}{D}{fg}{local}{Direct rotation on the two principal planes.}{unitary-operator}
\dkDeclareSymbol[Diagonal cosine block C = diag(cos(theta_i), cos(theta_j)) in the four-dimensional direct rotation D. This matrix differs from the scalar displacement cost C(M).]
  {PairCosineBlock}{C}{fg}{local}{Diagonal principal-cosine block.}{matrix}
\dkDeclareSymbol[Diagonal sine block S = diag(sin(theta_i), sin(theta_j)) in the four-dimensional direct rotation D. This matrix is local to the matching proof.]
  {PairSineBlock}{S}{fg}{local}{Diagonal principal-sine block.}{matrix}
\dkDeclareSymbol[Free planar orthogonal matrices U_0 and U_1 acting on the two-dimensional source and its two-dimensional complement, respectively. Multiplying D by diag(U_0,U_1) parametrizes every transporter on this four-dimensional block.]
  {PairFreeOrthogonal}{U}{fg}{local}{Free source or complementary planar orthogonal factor.}{unitary-operator}
\dkDeclareSymbol[Phase parameter x in the closed interval from h to pi-h for a positive-determinant four-dimensional transporter. Its eigenvalues are exp(plus or minus i(x+y)) and exp(plus or minus i(x-y)).]
  {PairPhaseX}{x}{fg}{local}{First phase parameter of a four-dimensional transporter.}{angle}
\dkDeclareSymbol[Phase parameter y in the closed interval from d to pi-d for a positive-determinant four-dimensional transporter. Together with x it gives rotational phases x+y and x-y, folded into the closed interval from zero to pi for the displacement cost.]
  {PairPhaseY}{y}{fg}{local}{Second phase parameter of a four-dimensional transporter.}{angle}
\dkDeclareSymbol[Scalar a = cos(x) in the four-dimensional phase calculation. Together with b = cos(y), it expresses the trace as 4ab and determines the reciprocal characteristic polynomial.]
  {PairPhaseCosineX}{a}{fg}{local}{Cosine of the first phase parameter in the trace calculation.}{scalar}
\dkDeclareSymbol[Scalar b = cos(y) in the four-dimensional phase calculation, paired with a = cos(x). These scalars are local to the characteristic-polynomial calculation.]
  {PairPhaseCosineY}{b}{fg}{local}{Cosine of the second phase parameter in the trace calculation.}{scalar}
\dkDeclareSymbol[Rotation parameter r: the two free orthogonal factors use planar rotation angles r+s and r-s. It appears in cos(y) = cos(d) cos(r) for two rotations.]
  {PairRotationR}{r}{fg}{local}{Half the sum of the two free planar rotation angles.}{angle}
\dkDeclareSymbol[Rotation parameter s: the two free orthogonal factors use planar rotation angles r+s and r-s. It appears in cos(x) = cos(h) cos(s) for two rotations. This scalar differs from the singular values and singleton costs.]
  {PairRotationS}{s}{fg}{local}{Half the difference of the two free planar rotation angles.}{angle}
\dkDeclareSymbol[Planar reflection J_2 = diag(1,-1). Every two-dimensional orthogonal matrix of determinant -1 can be written as a planar rotation times J_2.]
  {PairPlanarReflection}{J_2}{fg}{local}{Fixed planar reflection used to parametrize determinant -1 factors.}{unitary-operator}
\dkDeclareSymbol[Real planar rotation R_t = ((cos(t),-sin(t)),(sin(t),cos(t))). Its two displacement singular values are both 2 times abs(sin(t/2)). In the matching proof a rotation through t in the closed interval from zero to pi contributes 2 f_p(t).]
  {MatchingPlanarRotation}{\symglyph{PlanarRotationBase}_t}{fg}{local}{Planar rotation through the indicated angle.}{unitary-operator}
% Statistical application / operator provenance.
\dkDeclareSymbol{Sigma}{\Sigma}{old}{local}{Population covariance matrix.}{matrix}
\dkDeclareSymbol{SigmaHat}{\widehat\Sigma}{current}{local}{Estimated or perturbed covariance matrix.}{matrix}
\dkDeclareSymbol{YWSV}{V}{wanted}{local}{Orthonormal coordinates for the population eigenspace in the statistical application.}{isometric-embedding}
\dkDeclareSymbol{YWSVhat}{\widehat V}{trialframe}{local}{Estimated eigenspace frame in the statistical application.}{isometric-embedding}

% Figure-specific objects retain their own meanings when bare letters overlap.
\dkDeclareSymbol[Signed rotation angle phi of the displayed map on its invariant real two-plane. At theta = pi/4 the direct rotation turns its two principal planes by opposite signed angles of magnitude pi/4; the competitor turns its moving plane by pi/2 and fixes its other plane.]
  {PlaneRotationAngle}{\varphi}{sine}{local}{Signed rotation angle of the orthogonal map on the displayed invariant plane.}{angle}
\dkDeclareSymbol[Unit vector q in the displayed invariant plane. The map T is the direct rotation D or the competitor W in that panel; Tq is the image and (I - T)q = q - Tq is the displacement.]
  {FigureInput}{q}{ambient}{local}{Unit input to the orthogonal map shown in this invariant-plane panel.}{vector}
\dkDeclareSymbol[Image Tq of the panel unit input q under T = D or W. Both q and Tq lie on the unit circle; the green displacement arrow joins Tq to q, so it represents q - Tq.]
  {FigureImage}{Tq}{wantedvector}{local}{Image of the unit input under the orthogonal map shown in this panel.}{vector}
\dkDeclareSymbol[Displacement (I - T)q = q - Tq, drawn from the image Tq to the unit input q. On a plane rotated through phi its length is 2 |sin(phi/2)|; it vanishes on a fixed plane.]
  {DisplacementVector}{(I-T)q}{displacement}{local}{Difference between a unit input and its image under the displayed orthogonal map.}{vector}
\dkDeclareSymbol[First principal plane Pi_1 = span(e_0, e_3) for the source-target pair in the real four-dimensional counterexample. At theta = pi/4 the direct rotation acts on this plane by a quarter of pi, with displacement singular value sqrt(2 - sqrt(2)).]
  {CounterPrincipalOne}{\Pi_1}{fg}{glossary}{First principal plane in the real counterexample: the span of e0 and e3.}{linear-subspace}
\dkDeclareSymbol[Second principal plane Pi_2 = span(e_1, e_2) for the source-target pair in the real four-dimensional counterexample. The direct rotation acts by the same angle magnitude as on Pi_1, with the opposite sign in the displayed plane coordinates.]
  {CounterPrincipalTwo}{\Pi_2}{fg}{glossary}{Second principal plane in the real counterexample: the span of e1 and e2.}{linear-subspace}
\dkDeclareSymbol[Moving plane M = range(W - I) = span(e_0 + e_2, e_1 + e_3) for the competitor. W rotates it by 2 theta; at theta = pi/4 the two displacement singular values contributed by this plane are both sqrt(2).]
  {CounterMovingPlane}{M}{ambient}{glossary}{Moving plane of the competing transporter: the range of W minus the identity.}{linear-subspace}
\dkDeclareSymbol[Fixed plane Z = kernel(W - I) = span(e_0 - e_2, e_1 - e_3) for the competitor. W acts as the identity here, so this plane contributes two zero displacement singular values.]
  {CounterFixedPlane}{Z}{ambient}{glossary}{Fixed plane of the competing transporter: the kernel of W minus the identity.}{linear-subspace}
\dkDeclareSymbol[Direct-rotation displacement singular value a = 2 sin(pi/8) = sqrt(2 - sqrt(2)) at theta = pi/4. Each of its two principal planes contributes two copies, giving trace displacement 4a.]
  {CounterChordA}{a}{displacement}{glossary}{Direct-rotation displacement singular value in the counterexample: sqrt(2 - sqrt(2)).}{scalar}
\dkDeclareSymbol[Competitor displacement singular value b = 2 sin(pi/4) = sqrt(2) at theta = pi/4. Its moving plane contributes two copies; the fixed plane contributes two zeros, giving trace displacement 2b.]
  {CounterChordB}{b}{displacement}{glossary}{Competing-transporter displacement singular value in the counterexample: square root of 2.}{scalar}
\dkDeclareSymbol[Singular value s_j(I - T) of the full ambient displacement, ordered from largest to smallest. Each rotating real two-plane contributes two values 2 |sin(phi/2)|; fixed directions contribute zeros. Their sum is the trace cost.]
  {DisplacementSingular}{\symglyph{SingularValueBase}_j}{displacement}{local}{An individual singular value of the full displacement from the identity.}{scalar}
\dkDeclareSymbol[Edge weight g_ij = s_i + s_j - e_ij, the difference between the cost of two singletons and the paired cost. A matching has total cost C_0 minus the sum of its edge weights, with C_0 = sum_i s_i. Maximizing this sum gives the minimum transport cost; edges of negative weight may be omitted.]
  {MatchingEdgeWeight}{g_{ij}}{current}{glossary}{Edge weight: the cost of two singletons minus their paired cost.}{scalar-weight}
\dkDeclareSymbol[Total trace-norm displacement cost. For a transporter W, C(W) = ||I - W||_1, the sum of its displacement singular values. For a matching M, C(M) is the sum of its paired and unmatched singleton costs.]
  {TraceCost}{C}{norm}{glossary}{Trace-norm displacement cost of the indicated orthogonal transporter or matching.}{scalar-cost}

% Spectral values have their own key, distinct from an operator block.
\dkDeclareSymbol[Spectral parameter lambda of a self-adjoint operator. In a finite-dimensional eigenvalue list lambda_j(A) it denotes an individual eigenvalue; in A = integral lambda dE(lambda), it is the real integration variable.]{SpectralValue}{\lambda}{wantedspectrum}{local}{Real spectral parameter or individual eigenvalue of a self-adjoint operator.}{scalar}



% Conventional mathematical notation. Hover text applies even when ink is plain.
% These are local (non-glossary) unless promoted explicitly by the author.
% IdentityTwo and IdentityFour derive their I from Identity. Singular-value
% registrations derive s_j from SingularValueBase; SingletonCost has its own s.
\dkDeclareSymbol[Identity operator on the vector or Hilbert space in context. Its dimension is inferred from the expression and need not be finite.]{Identity}{I}{plain}{local}{Identity operator on the space in context.}{identity-operator}
\dkDeclareSymbol[Two-by-two identity matrix, acting on a two-dimensional coordinate space.]{IdentityTwo}{\symglyph{Identity}_2}{plain}{local}{Two-by-two identity matrix.}{identity-operator}
\dkDeclareSymbol[Four-by-four identity matrix, acting on a four-dimensional coordinate space.]{IdentityFour}{\symglyph{Identity}_4}{plain}{local}{Four-by-four identity matrix.}{identity-operator}
\dkDeclareSymbol[Two-by-two zero matrix. Unlike the bold zero vector, this is a linear operator/matrix with two rows and two columns.]{ZeroMatrixTwo}{0_2}{plain}{local}{Two-by-two zero matrix.}{matrix}
\dkDeclareSymbol[The constant pi, the ratio of the circumference of a circle to its diameter. Angles in this article use radians.]{PiConstant}{\pi}{plain}{local}{The constant pi (angles are in radians).}{constant}
\dkDeclareSymbol[The field of real numbers.]{RealField}{\mathbb R}{plain}{local}{The real numbers.}{field}
\dkDeclareSymbol[The field of complex numbers.]{ComplexField}{\mathbb C}{plain}{local}{The complex numbers.}{field}
\dkDeclareSymbol[Generic Hilbert space used in the two-projection discussion; distinct from the paper's ambient space symbol E in the sine-theta theorem.]{GeneralHilbertSpace}{\mathcal H}{plain}{local}{Generic Hilbert space in a local discussion.}{hilbert-space}
\dkDeclareSymbol[The spectrum sigma(T) of an operator T: the set of complex scalars for which T minus the scalar identity is not boundedly invertible. For self-adjoint T the spectrum is real.]{SpectrumSymbol}{\sigma}{plain}{local}{Spectrum of an operator.}{operator-notation}
\dkDeclareSymbol[Base letter for singular values, presently s. It is distinct from the base s for matching singleton costs.]{SingularValueBase}{s}{plain}{local}{Base letter for singular values.}{notation-base}
\dkDeclareSymbol[Ordered singular value s_j(T), equal to the j-th eigenvalue of |T| for a finite matrix; distinguish this from the matching singleton cost s_j.]{SingularValueS}{\symglyph{SingularValueBase}_j}{plain}{glossary}{The j-th singular value, under the currently printed s_j convention.}{scalar}
\dkDeclareSymbol[Alternate indexed singular-value notation sigma_j(T). The current manuscript also uses s_j(T); this registered alias is available for a later deliberate notation decision.]{SingularValueSigma}{\sigma_j}{plain}{local}{Indexed singular value under the sigma_j convention.}{scalar}
\dkDeclareSymbol[Operator norm ||T||_op, the supremum of ||Tx|| over unit vectors x.]{OperatorNormNotation}{\|\cdot\|_{\mathrm{op}}}{plain}{glossary}{Operator norm.}{norm}
\dkDeclareSymbol[Frobenius norm ||T||_F, the square root of the sum of squared matrix entries, equal to the finite-dimensional Hilbert--Schmidt norm.]{FrobeniusNormNotation}{\|\cdot\|_F}{plain}{glossary}{Frobenius norm.}{norm}
\dkDeclareSymbol[Trace norm ||T||_1, the sum of singular values when finite.]{TraceNormNotation}{\|\cdot\|_1}{plain}{glossary}{Trace (nuclear) norm.}{norm}
\dkDeclareSymbol[Hilbert--Schmidt norm ||T||_HS, the square root of the sum of squared singular values when finite.]{HilbertSchmidtNormNotation}{\|\cdot\|_{\mathrm{HS}}}{plain}{local}{Hilbert--Schmidt norm.}{norm}
\dkDeclareSymbol[Ky Fan j-norm, the sum of the largest j singular values.]{KyFanNormNotation}{\|\cdot\|_{(j)}}{plain}{glossary}{Ky Fan j-norm.}{norm}
\dkDeclareSymbol[Interval/exterior spectral separation predicate I_delta(B,C), specifying a central closed interval separated by delta from the other operator's spectrum.]{IntervalSeparation}{\mathcal I}{plain}{local}{Interval/exterior separation predicate.}{predicate}
\dkDeclareSymbol[Ordered spectral separation predicate O_delta(B,C), requiring one self-adjoint operator to lie below a spectral threshold and the other above it by at least delta.]{OrderedSeparation}{\mathcal O}{plain}{local}{Ordered spectral separation predicate.}{predicate}
\dkDeclareSymbol[Positive parameter epsilon controlling a small perturbation in the two-dimensional eigenvalue example.]{EigenvalueExampleEpsilon}{\varepsilon}{plain}{local}{Small positive parameter in the eigenvalue-example matrices.}{scalar}
\dkDeclareSymbol[Positive lower bound epsilon for the trial-coordinate map E_0, as in E_0* E_0 >= epsilon squared times the identity.]{FrameLowerBoundEpsilon}{\varepsilon}{plain}{local}{Lower-bound parameter for the trial-coordinate embedding.}{scalar}
\dkDeclareSymbol[Rotation angle vartheta of a planar orthogonal block in the matching-proof first-variation calculation.]{MatchingPlanarAngle}{\vartheta}{plain}{local}{Planar rotation angle in the matching-proof calculation.}{angle}


% Ordinary named mathematical operators already used throughout the manuscript.
% Keeping these here makes their spelling a single editorial choice.
\dkDeclareSymbol[Range of a linear operator: the set of its output vectors.]{RangeOperator}{\operatorname{ran}}{plain}{local}{Range (image) of an operator.}{operator-notation}
\dkDeclareSymbol[Kernel (null space) of a linear operator: inputs mapped to zero.]{KernelOperator}{\operatorname{ker}}{plain}{local}{Kernel (null space) of an operator.}{operator-notation}
\dkDeclareSymbol[Domain of an operator: the vectors on which the operator is defined, important for unbounded operators.]{DomainOperator}{\operatorname{dom}}{plain}{local}{Domain of an operator.}{operator-notation}
\DeclareRobustCommand{\DKRange}{\sym{RangeOperator}}
\DeclareRobustCommand{\DKKernel}{\sym{KernelOperator}}
\DeclareRobustCommand{\DKDom}{\sym{DomainOperator}}

% Norm-formatting macros: the operand changes by occurrence but the notation
% lives here. The outer hover describes the norm and does not add ink color.
\DeclareRobustCommand{\dkOperatorNorm}[1]{\like{OperatorNormNotation}{\|#1\|_{\mathrm{op}}}}
\DeclareRobustCommand{\dkTraceNorm}[1]{\like{TraceNormNotation}{\|#1\|_1}}
\DeclareRobustCommand{\dkFrobeniusNorm}[1]{\like{FrobeniusNormNotation}{\|#1\|_F}}
\DeclareRobustCommand{\dkHilbertSchmidtNorm}[1]{\like{HilbertSchmidtNormNotation}{\|#1\|_{\mathrm{HS}}}}
\DeclareRobustCommand{\dkKyFanNorm}[2]{\like{KyFanNormNotation}{\|#1\|_{(#2)}}}
\DeclareRobustCommand{\dkSchattenNorm}[2]{\like{SchattenNormNotation}{\|#1\|_{S_{#2}}}}

\dkDeclareSymbol[Schatten p-norm ||T||_{S_p}, obtained from the l^p-norm of the singular values and finite for operators in the Schatten p-class.]{SchattenNormNotation}{\|\cdot\|_{S_p}}{plain}{local}{Schatten p-norm.}{norm}


\dkDeclareSymbol[Positive modulus of an operator T: |T| = (T* T)^(1/2).]{OperatorModulusNotation}{|\cdot|}{plain}{glossary}{Positive modulus of an operator.}{operator-notation}
\dkDeclareSymbol[Linear span of the indicated vectors, i.e. the smallest linear subspace that contains them.]{SpanOperator}{\operatorname{span}}{plain}{local}{Linear span of vectors.}{operator-notation}
\dkDeclareSymbol[Diagonal or block-diagonal matrix/operator with the displayed entries or blocks.]{DiagonalOperator}{\operatorname{diag}}{plain}{local}{Diagonal or block-diagonal operator.}{linear-operator}
\dkDeclareSymbol[Trace of a trace-class operator, or ordinary matrix trace in finite dimensions.]{TraceOperator}{\operatorname{tr}}{plain}{local}{Trace of an operator.}{operator-notation}
\dkDeclareSymbol[Dimension of a vector space or subspace.]{DimensionOperator}{\dim}{plain}{local}{Dimension of a vector space.}{operator-notation}
\dkDeclareSymbol[Rank of a linear map or matrix: dimension of its range.]{RankOperator}{\operatorname{rank}}{plain}{local}{Rank of a linear map or matrix.}{operator-notation}
\dkDeclareSymbol[Distance between sets of spectral values: the infimum of pairwise distances.]{SpectralDistanceOperator}{\operatorname{dist}}{plain}{local}{Distance between two spectral sets.}{operator-notation}
\DeclareRobustCommand{\dkModulus}[1]{\like{OperatorModulusNotation}{|#1|}}

% Context-specific glyphs previously hidden in variants of other semantic keys.
% Keeping both the roles and the exact printed forms makes collisions explicit.
\dkDeclareSymbol[Nonzero vector x tested in the Rayleigh quotient, before imposing unit length.]{RayleighTestVector}{x}{trialvector}{local}{Test vector in the Rayleigh quotient.}{vector}
\dkDeclareSymbol[Positive cosine block C in the generic two-projection decomposition of Halmos; C and S commute and C squared plus S squared is the identity.]{HalmosCosineBlock}{C}{anglecos}{local}{Cosine block in the Halmos two-projection representation.}{linear-operator}
\dkDeclareSymbol[Positive sine block S in the generic two-projection decomposition of Halmos; C and S commute and C squared plus S squared is the identity.]{HalmosSineBlock}{S}{sineoperator}{local}{Sine block in the Halmos two-projection representation.}{linear-operator}
\dkDeclareSymbol[Riesz spectral projection P_Gamma defined by a contour Gamma that encloses selected spectral values.]{RieszProjection}{P_{\symglyph{RieszContour}}}{wantedprojection}{local}{Spectral projection associated with a contour Gamma.}{orthogonal-projection}
\dkDeclareSymbol[Historical Sylvester right-hand coefficient matrix A, in B X - X A = Y. Distinct from the current paper's reference operator A.]{HistoricalSylvesterRight}{A}{trialspectrum}{local}{Right-hand coefficient A in the historical Sylvester equation.}{self-adjoint-operator}
\dkDeclareSymbol[Scalar coefficient b in the Sylvester prototype (b-a)x=y, with separated scalar spectral values a and b.]{SylvesterScalarLeft}{b}{unwantedspectrum}{local}{Scalar left coefficient b in a Sylvester equation.}{scalar}
\dkDeclareSymbol[Scalar coefficient a in the Sylvester prototype (b-a)x=y.]{SylvesterScalarRight}{a}{trialspectrum}{local}{Scalar right coefficient a in a Sylvester equation.}{scalar}
\dkDeclareSymbol[Scalar unknown x in the Sylvester prototype (b-a)x=y.]{SylvesterScalarUnknown}{x}{sylvester}{local}{Scalar unknown x in a Sylvester equation.}{scalar}
\dkDeclareSymbol[Scalar right side y in the Sylvester prototype (b-a)x=y.]{SylvesterScalarRHS}{y}{resid}{local}{Scalar right side y in a Sylvester equation.}{scalar}
\dkDeclareSymbol[Source subspace K of a real Euclidean space in the matching proof; an orthogonal transporter sends K onto L.]{MatchingSourceSubspace}{K}{trial}{local}{Source subspace K in the matching proof.}{linear-subspace}
\dkDeclareSymbol[Target subspace L of a real Euclidean space in the matching proof; an orthogonal transporter sends K onto L.]{MatchingTargetSubspace}{L}{wanted}{local}{Target subspace L in the matching proof.}{linear-subspace}
\dkDeclareSymbol[Orthogonal projection P_L onto the matching proof's target subspace L.]{MatchingTargetProjection}{P_L}{wantedprojection}{local}{Orthogonal projection onto the matching target L.}{orthogonal-projection}
\dkDeclareSymbol[Real two-dimensional orthogonal group O(2): planar rotations and reflections.]{PlanarOrthogonalGroup}{O(2)}{fg}{local}{Orthogonal matrices of size two.}{group}
\dkDeclareSymbol[Base letter R used to denote planar rotation operators R_t in the matching proof.]{PlanarRotationBase}{R}{fg}{local}{Base letter for planar rotations.}{notation-base}
\dkDeclareSymbol[Planar rotation R_(r+s) through the sum of scalar phases r and s.]{MatchingPlanarRotationSum}{\symglyph{PlanarRotationBase}_{r+s}}{fg}{local}{Planar rotation through phase r+s.}{unitary-operator}
\dkDeclareSymbol[Planar rotation R_gamma through phase gamma in the block rotation calculation.]{MatchingPlanarRotationGamma}{\symglyph{PlanarRotationBase}_{\symglyph{MatchingPhaseGamma}}}{fg}{local}{Planar rotation through phase gamma.}{unitary-operator}

% Specific symbols in the historical and appendix examples. They retain their
% existing black-ink notation but receive locally accurate hover definitions.
\dkDeclareSymbol[Contour Gamma, a positively oriented closed curve in the resolvent set enclosing a selected isolated spectral cluster.]{RieszContour}{\Gamma}{plain}{local}{Contour used in the Riesz spectral-projection formula.}{contour}
\dkDeclareSymbol[Even integer nu in a tested Ky Fan comparison, with nu/2 principal-plane terms on the right side.]{KyFanEvenIndex}{\nu}{plain}{local}{Even integer in the principal-plane Ky Fan comparison.}{scalar}
\dkDeclareSymbol[Orthogonal projection Omega_k onto the k-th principal plane in the published displacement comparison.]{PrincipalPlaneProjectionBase}{\Omega}{plain}{local}{Projection onto the indicated principal plane.}{notation-base}
\dkDeclareSymbol[Scalar gap threshold gamma in the ordered Sylvester spectral separation B >= (gamma+delta)I and C <= gamma I.]{SylvesterGapThreshold}{\gamma}{plain}{local}{Spectral separation threshold in a Sylvester estimate.}{scalar}
\dkDeclareSymbol[Projection Omega_t from a net or sequence that converges strongly to the identity in the singular-value approximation lemma.]{StrongApproxProjection}{\Omega_t}{plain}{local}{Orthogonal projection converging strongly to the identity.}{orthogonal-projection}
\dkDeclareSymbol[First orthogonal projection Omega used in the two-block pinching inequality for a unitarily invariant norm.]{PinchingProjectionOne}{\Omega}{plain}{local}{First projection in the pinching inequality.}{orthogonal-projection}
\dkDeclareSymbol[Second orthogonal projection Upsilon used in the two-block pinching inequality for a unitarily invariant norm.]{PinchingProjectionTwo}{\Upsilon}{plain}{local}{Second projection in the pinching inequality.}{orthogonal-projection}
\dkDeclareSymbol[Finite-rank projection Gamma used to retain the image of a chosen finite-rank compression K Psi.]{CompressionProjectionImage}{\Gamma}{plain}{local}{Range projection in the finite-rank compression lemma.}{orthogonal-projection}
\dkDeclareSymbol[Finite-rank projection Psi used to select input vectors in a nearly maximizing singular-value compression.]{CompressionProjectionInput}{\Psi}{plain}{local}{Domain projection in the finite-rank compression lemma.}{orthogonal-projection}
\dkDeclareSymbol[Small positive tolerance eta in the finite-rank compression approximation estimate.]{CompressionTolerance}{\eta}{plain}{local}{Tolerance controlling the off-diagonal compression norm.}{scalar}

\dkDeclareSymbol[Rotation phase gamma in a planar block from the matching proof.]{MatchingPhaseGamma}{\gamma}{plain}{local}{Rotation phase gamma in the matching calculation.}{scalar}

% Two common indices with local mathematical meanings distinct from other deltas.
\dkDeclareSymbol[Kronecker delta: delta_ij is one when the indices i and j agree and zero otherwise. This is not the spectral-gap parameter delta.]{KroneckerDelta}{\delta}{plain}{local}{Kronecker delta: equals 1 for equal indices and 0 otherwise.}{scalar}
\dkDeclareSymbol[Sign variable epsilon taking the values -1 and +1 in the factorization over two phase choices.]{PhaseSignChoice}{\epsilon}{plain}{local}{Sign variable taking values minus one and plus one.}{scalar}

% Documented notation correspondence. The source column is a faithful
% transcription of each cited original context, NOT an asserted universal DK
% meaning. The original A_0 is deliberately listed for two distinct roles:
% its reference block and the comparison/trial operator. The source uses the
% same printed symbol for both; this paper is free to distinguish them.
\dkDeclareNotationSource{DK1970}{Davis and Kahan (1970), The Rotation of Eigenvectors by a Perturbation. III}{https://doi.org/10.1137/0707001}
\dkDeclareNotationSource{Seelmann2019}{Seelmann (2019), Semidefinite Perturbations in the Subspace Perturbation Problem}{https://arxiv.org/abs/1708.02463}
\dkDeclareNotationSource{Blechta2020}{Blechta (2020), Stability of Linear GMRES Convergence with Respect to Compact Perturbations}{https://arxiv.org/abs/2005.12960}
% Arguments: {key}{source ID}{source glyph}{section or equation}{role-specific qualification}
\dkDeclareNotationAlternative{A}{DK1970}{A}{\S1, (1.2)}{Reference self-adjoint operator.}
\dkDeclareNotationAlternative{H}{DK1970}{H}{\S1, (1.2)}{Additive perturbation.}
\dkDeclareNotationAlternative{Acurrent}{DK1970}{A+H}{\S1}{Perturbed operator.}
\dkDeclareNotationAlternative{E0}{DK1970}{E_0}{\S1, (1.1)}{Source isometric frame.}
\dkDeclareNotationAlternative{E1}{DK1970}{E_1}{\S1, (1.1)}{Reference complement frame.}
\dkDeclareNotationAlternative{F0}{DK1970}{F_0}{\S1}{Exact-target frame.}
\dkDeclareNotationAlternative{F1}{DK1970}{F_1}{\S1}{Exact-target complement frame.}
\dkDeclareNotationAlternative{ReferenceBlock0}{DK1970}{A_0}{\S1, (1.2)}{Reference source block.}
\dkDeclareNotationAlternative{ReferenceBlock1}{DK1970}{A_1}{\S1, (1.2)}{Reference complement block.}
\dkDeclareNotationAlternative{A0}{DK1970}{A_0}{\S1, residual discussion}{Residual comparison operator (the reference block in DK).}
% We deliberately leave Lambda0 and Lambda1 without DK source-glyph metadata:
% the available source transcription does not reliably preserve the source
% accents/decorations on the perturbed spectral-block symbols.
\dkDeclareNotationAlternative{Psource}{DK1970}{P}{\S1}{Reference/source projection.}
\dkDeclareNotationAlternative{Qtarget}{DK1970}{Q}{\S1}{Target projection.}
\dkDeclareNotationAlternative{Pexact}{DK1970}{Q}{\S1}{Exact-target projection.}
\dkDeclareNotationAlternative{B}{DK1970}{B}{\S1, (1.2)}{Off-diagonal coupling block.}
\dkDeclareNotationAlternative{Theta0}{DK1970}{\Theta_0}{\S1}{Source-side angle operator.}
\dkDeclareNotationAlternative{Theta1}{DK1970}{\Theta_1}{\S1}{Complementary angle operator.}
\dkDeclareNotationAlternative{Theta}{DK1970}{\Theta}{\S1}{Full angle operator.}
\dkDeclareNotationAlternative{R}{DK1970}{R}{\S1, residual discussion}{Trial residual.}
\dkDeclareNotationAlternative{Rrotation}{DK1970}{U}{\S1, (1.18)}{Canonical direct rotation.}
\dkDeclareNotationAlternative{Wcompetitor}{DK1970}{V}{\S4}{Competing unitary transporter.}
\dkDeclareNotationAlternative{delta}{DK1970}{\delta}{\S2}{Spectral separation parameter.}
% Contemporary alternatives are recorded separately and do not enter the DK
% correspondence table. Their citation locations must be individually checked.
\dkDeclareNotationAlternative{Eambient}{Seelmann2019}{\mathcal H}{p.~1}{Ambient Hilbert space.}
\dkDeclareNotationAlternative{SingularValueS}{Blechta2020}{\sigma_j}{Thm.~1.1}{Ordered singular values of a compact operator.}
 in paper.tex)
% or pass \def\DKMonochromeMath{1} before % One switch: false = journal style; true = colored symbols and author--year citations.
\newif\ifDKAuthorStyle

%\DKAuthorStylefalse
\DKAuthorStyletrue

\ifdefined\DKUseAuthorStyle\DKAuthorStyletrue\fi
\ifdefined\DKJournalStyle\DKAuthorStylefalse\fi
\ifDKAuthorStyle
  \PassOptionsToClass{sn-mathphys-ay}{sn-jnl}
\else
  \PassOptionsToClass{sn-mathphys-num}{sn-jnl}
\fi
\documentclass[pdflatex]{sn-jnl} % PDF engine option; compiled with LuaLaTeX.
% PDF version used for the manuscript and its symbol annotations.
\pdfvariable minorversion=7\relax

\usepackage{amsmath,amssymb,amsthm,mathtools}
\usepackage{booktabs,tabularx,array,longtable}
\usepackage[section]{placeins}
\usepackage{float}
\usepackage{microtype}
\usepackage{graphicx}
\usepackage{pgf} % Live vector figures with LaTeX symbol labels and tooltips
\usepackage{xurl}
\ifDKAuthorStyle
  \setcitestyle{authoryear,round,aysep={,}}
  \hypersetup{colorlinks=true,allcolors=blue}
\else
  \setcitestyle{numbers,square}
  \hypersetup{hidelinks}
\fi
\usepackage[nameinlink]{cleveref}
\usepackage[T1]{fontenc}
\usepackage{lmodern} % Scalable LaTeX lettering, shared with the Manim slides
\usepackage{xcolor}
% Draft controls. Set these independently for review, print, or arXiv versions.
\newif\ifDKAnonymous
\DKAnonymousfalse
\ifdefined\DKAnonymousReview\DKAnonymoustrue\fi
\newif\ifDKColoredSymbols
\ifDKAuthorStyle\DKColoredSymbolstrue\else\DKColoredSymbolsfalse\fi
\ifdefined\DKMonochromeMath\DKColoredSymbolsfalse\fi
% The colors retain the Manim glossary's light-theme mathematical roles.
\input{math_symbol_colors.tex}
\ifDKColoredSymbols\else\DKMathColorsfalse\fi
\ifDKAuthorStyle\else\DKSymbolTooltipsfalse\fi




\usepackage[most]{tcolorbox}
\usepackage{listings}
\usepackage{listingsutf8}
\usepackage{fontspec}
\IfFileExists{emoji.sty}{%
  \IfFontExistsTF{Noto Color Emoji}{\usepackage{emoji}}{}%
}{}
\providecommand{\emoji}[1]{\textnormal{[#1]}} % Review notes also work without an emoji font.

\setlength{\emergencystretch}{3em}
\raggedbottom

% ArXiv-friendly Lean code: no shell escape, minted, or external highlighter.
\definecolor{leankeyword}{HTML}{243B6B}
\definecolor{leantactic}{HTML}{70477A}
\definecolor{leantype}{HTML}{1F6473}
\definecolor{leancomment}{HTML}{39734D}
\definecolor{leanstring}{HTML}{8A4B2A}
\definecolor{leannumber}{HTML}{777777}
\definecolor{leanbackground}{HTML}{F7F7F3}
\lstdefinelanguage{lean}{%
  sensitive=true,
  morecomment=[l]{--},
  morecomment=[s]{/-}{-/},
  morestring=[b]",
  morekeywords=[1]{import,open,namespace,section,end,noncomputable,variable,variables,
    universe,universes,set_option,attribute,axiom,opaque,theorem,lemma,def,abbrev,
    example,instance,structure,class,inductive,notation,local,deriving,extends,
    by,where,let,have,show,fun,match,with,if,then,else,do,return,in,forall},
  morekeywords=[2]{simp,simpa,rw,erw,dsimp,unfold,ring,ring_nf,nlinarith,linarith,
    positivity,norm_num,field_simp,omega,exact,refine,apply,assumption,intro,intros,
    ext,funext,cases,rcases,obtain,constructor,left,right,classical,by_contra,
    contrapose,convert,using,suffices,calc},
  morekeywords=[3]{Prop,Type,Nat,Int,Real,Fin,Set,Submodule,Matrix,EuclideanSpace,
    LinearMap,ContinuousLinearMap,LinearPMap,RCLike,IsTrialResidual,
    IsExactSpectralDecomposition,FormBoundedSylvesterGap,
    NormalizedSymmetricOperatorIdealFamily},
}
\lstdefinestyle{leanpaper}{%
  language=lean,
  columns=fullflexible,
  backgroundcolor=\color{leanbackground},
  basicstyle=\ttfamily\scriptsize,
  commentstyle=\color{leancomment}\itshape,
  keywordstyle=[1]\color{leankeyword}\bfseries,
  keywordstyle=[2]\color{leantactic},
  keywordstyle=[3]\color{leantype},
  stringstyle=\color{leanstring},
  numberstyle=\tiny\color{leannumber},
  breakatwhitespace=false,
  breaklines=true,
  captionpos=b,
  keepspaces=true,
  numbers=left,
  numbersep=6pt,
  showspaces=false,
  showstringspaces=false,
  showtabs=false,
  tabsize=2,
  aboveskip=0.8em,
  belowskip=0.8em,
  literate=
    {\\KField}{{$\Bbbk$}}1
    {∃}{{$\exists$}}1
    {∧}{{$\land$}}1
    {∀}{{$\forall$}}1
    {≠}{{$\neq$}}1
    {⊤}{{$\top$}}1
}
\crefname{lstlisting}{formalization}{formalizations}
\Crefname{lstlisting}{Formalization}{Formalizations}
\renewcommand{\lstlistingname}{Formalization}

\theoremstyle{thmstyleone}
\newtheorem{theorem}{Theorem}
\newtheorem{proposition}[theorem]{Proposition}
\newtheorem{lemma}[theorem]{Lemma}
\newtheorem{remark}[theorem]{Remark}
\newtheorem{corollary}[theorem]{Corollary}

\newcommand{\Lean}{\textsc{Lean}}
\newcommand{\DK}{Davis--Kahan}
\newcommand{\TauCeti}{Tau Ceti}
% Use \ie{} and \eg{}; supply surrounding punctuation at the call site.
\newcommand{\ie}{i.e.\@}
\newcommand{\eg}{e.g.\@}
\newcommand{\DKResultCount}{29}
\newcommand{\DKProvedCount}{28}
\newcommand{\DKRefutedCount}{1}

% Usage: \margincomment{Check this formula!}
%\newcommand{\jonsays}[1]{%
%  \marginpar{%
%    \raggedright\color{red}\tiny% Change to \footnotesize if text is too small
%    \emoji{person-raising-hand} \textbf{Comment:}\\#1%
%  }%
%}

% Usage: \margincomment{Your comment that now wraps perfectly!}
\newcommand{\jonsays}[1]{%
  \marginpar{%
    \raggedright% Forces text to wrap beautifully without ugly hyphenation
    \color{red}\scriptsize% Using \scriptsize makes the text tighter
    \emoji{person-raising-hand}: \\#1%
  }%
}
%\newcommand{\jonsays}[1]{{\color{red}\small\textbf{[\emoji{person-raising-hand}: #1]}}}

% Usage: \margincomment{Your comment that now wraps perfectly!}
%\newcommand{\agentsays}[1]{%
%  \marginpar{%
%    \raggedright% Forces text to wrap beautifully without ugly hyphenation
%    \color{red}\scriptsize% Using \scriptsize makes the text tighter
%    \emoji{robot}: \\#1%
%  }%
%}
\newcommand{\agentsays}[1]{{\color{red}\small\textbf{[\emoji{robot}: #1]}}}

\newcommand{\jonrsp}[1]{{\color{red}\small\textbf{[\emoji{person-raising-hand}: #1]}}}

\title[Formalizing Davis--Kahan Perturbation Theory]{Formalizing Davis--Kahan Perturbation Theory in Lean 4}
% Additional authors are expected; keep the confirmed ordering below stable.
\ifDKAnonymous
  \author{Anonymous authors}
\else
  \author*[1]{\fnm{Jonathan} \sur{Crall}}\email{jon.crall@kitware.com}
  \author[2]{\fnm{Edward} \sur{Wang}}\email{ewang43@jhu.edu}
  \author[1]{\fnm{Brian} \sur{Hu}}\email{brian.hu@kitware.com}
  \author[2]{\fnm{Carey E.} \sur{Priebe}}\email{cep@jhu.edu}
  \affil[1]{\orgname{Kitware, Inc.},
    \orgaddress{\city{Clifton Park}, \state{NY}, \country{USA}}}
  \affil[2]{\orgdiv{Department of Applied Mathematics and Statistics},
    \orgname{Johns Hopkins University},
    \orgaddress{\city{Baltimore}, \state{MD}, \country{USA}}}
\fi

\abstract{We formalize the spectral-subspace perturbation theory of Davis and
Kahan's 1970 paper in Lean~4. Of \DKResultCount{} selected results, we prove
\DKProvedCount{} under their stated hypotheses and refute the remaining
printed statement. The development includes the sine and tangent theorems,
direct rotations, principal-angle operators, Sylvester separation, and
operator-ideal norms on real and complex Hilbert spaces, including unbounded
self-adjoint operators. We explain the mathematical lineage connecting these
constructions. The refuted statement concerns minimization of full ambient
displacement by direct rotation. We give real four-dimensional counterexamples
at every prescribed positive acute angle and identify the failed
singular-value comparison. A corrected optimality statement holds for Q-norms,
whose squares are unitarily invariant norms of the squared modulus, including
the Frobenius norm. The formalization builds on Mathlib and adapts Spectra's
projection-valued-measure and Borel-calculus constructions. Portions of the
operator theory developed here were subsequently adapted into Tau Ceti.
These foundations also support an exact reduction of trace-norm optimal real
subspace transport to weighted matching on principal angles.}

\keywords{Davis--Kahan theory, spectral perturbation, formal verification,
Lean 4, direct rotation, principal angles}
\pacs[MSC Classification (2020)]{68V20, 47A55, 15A42}

\begin{document}
\maketitle

\ifDKAnonymous\else
\enlargethispage{2\baselineskip}
\vspace{-18pt}
\noindent\begin{minipage}{\linewidth}\footnotesize
\noindent\textbf{Acknowledgments.}
This material is based upon work supported by the Defense Advanced Research
Projects Agency (DARPA) under Contract No. HR001125CE017. Any opinions, findings
and conclusions or recommendations expressed in this material are those of the
author(s) and do not necessarily reflect the views of the Defense Advanced
Research Projects Agency.
\end{minipage}
\fi
\clearpage

\section{Introduction}
\label{sec:introduction}

An approximate eigenvector gives a candidate direction for the action of an
operator. If a unit vector $\sym{v}$ is paired with an approximate eigenvalue $\sym{mu}$,
the residual $\sym{A}\sym{v}-\sym{mu} \sym{v}$ measures how closely it satisfies the eigenvector
equation. A small residual can identify an accurate approximation to an
invariant subspace even when repeated or nearby eigenvalues make individual
eigenvectors difficult to distinguish. Davis--Kahan theory bounds the angles
between approximate and exact spectral subspaces in terms of a residual or
perturbation and a separation between the relevant spectra
\citep{davis1970rotation}.

In applications, one often works with a perturbed operator
  $\sym{Acurrent}=\sym{A}+\sym{H}$ while the reference operator $\sym{A}$ is
  unknown.
Taking an eigenspace of $\sym{Acurrent}$ as a trial subspace for $\sym{A}$, the norm of its
  residual relative to $\sym{A}$ is bounded by the norm of $\sym{H}$, allowing
  Davis--Kahan theory to bound the angle between the two subspaces without
  requiring $\sym{A}$ to be known or its residuals to be evaluated directly.

Our project began with a variant of this bound needed for
\ifDKAnonymous a spectral-estimation application\else the data kernel
perspective space (DKPS) \citep{helm2026query}\fi.

We extended the formalization to the results established in the original 1970
paper, including real and complex Hilbert spaces, unbounded self-adjoint
operators, and unitarily invariant operator-ideal norms. In particular, the
spectral-separation hypotheses allow half-infinite intervals
(e.g.\ $(-\infty,a]$). These settings require operator domains, spectral
calculus, and norm estimates beyond finite matrices.

We formalized the four headline theorems \citep[Section~2]{davis1970rotation} and every
named theorem, proposition, lemma, and corollary established in the 1970 paper.

This gives \DKResultCount{} targets, enumerated in
Appendix~\ref{app:dk-inventory}. We prove \DKProvedCount{} of them and refute
the printed statement \citep[Proposition~4.4]{davis1970rotation}.
That proposition asserts a full-displacement
optimality property for real direct rotations. We give a family of real
four-dimensional examples exhibiting trace-norm failure at every prescribed
positive acute angle and identify
the invalid singular-value comparison in its proof. The phase construction
behind the counterexample occurs over the complex field in the authors'
example \citep[Example~4.2]{davis1970rotation}. To the best of our knowledge,
its contradiction of the printed real, small-angle claim has not previously
been made explicit.
For the constructed pair, whose two principal angles both equal
$\sym{theta}$, the competing transporter attains the global trace-norm minimum
$4\sin\sym{theta}$ for the constructed subspaces. Section~\ref{sec:prop44}
compares these statements and repairs the optimality claim by restricting
the norm class to Q-norms (N is a Q-norm when $N(T)^2 = M(T^* T)$ for some
unitarily invariant norm $M$). This repair follows from the squared-displacement
comparison \citep[Proposition~4.3]{davis1970rotation};
earlier Q-norm minimization results appear in \citet{bolotnikov2001minimal}.

% TODO (lineage): Consolidate the scattered lineage records through the
% maintained dev/tauceti/dkps-upstream-contribution-census.md and the evidence
% index in AUTHORING_NOTES.md. Audit each attribution against an immutable
% donor commit/path, the accepted upstream commit/PR, and the corresponding
% declaration statements and proofs. Preserve Mathlib prerequisites, Spectra
% Spectra-derived constructions where applicable, and independent work as distinct relationships.
% Courant--Fischer/Weyl checked 2026-10-28: developed here in June 2026;
% August donor df78246f965881705c9301f96feddc270c98ceb1 predates the
% accepted Tau Ceti port 7ebf17ba04725c1a575ecf3d099ce0630cd949a0 (#13599).
% See COURANT_FISCHER_LINEAGE_AUDIT.md for the proof comparison and sources.

The proofs required us to formalize a significant amount of operator
theory, namely: principal angles and direct rotations, approximation numbers
and symmetric operator ideals, Borel spectral calculus, and Sylvester
estimates for unbounded operators. These developments build on Mathlib and
adapt Spectra's projection-valued-measure and Borel-calculus constructions
\citep{mathlib2020,bornemann2026spectra}.
These operator-theory foundations include our
Courant--Fischer and Weyl proofs, rectangular unitarily invariant seminorms,
and projection-valued measures, which were subsequently adapted into Tau
Ceti \citep{tauceti2026}.
Section~\ref{sec:related} explains the mathematical history and formal
provenance; Section~\ref{sec:foundations} describes the foundations used in
our proofs.

The same principal-angle and operator-ideal results support a further
optimization theorem. In Section~\ref{sec:example-transport}, we prove that
minimizing trace-norm displacement over all real orthogonal maps carrying
one subspace onto another reduces to weighted matching on principal angles.

\noindent\textbf{Contributions.}
\begin{enumerate}
\item \textbf{Davis--Kahan formalization.} We prove 28 of the 29 source
results and refute the full-displacement assertion
\citep[Proposition~4.4]{davis1970rotation}. The sine-theta theorem covers real and
complex Hilbert spaces, unbounded self-adjoint operators, the source's
spectral-separation alternatives, and operator-ideal norms where finite
(\cref{sec:coverage,sec:sine-theta}).
\item \textbf{Formalized operator-theoretic foundations.} We formalized the
geometric and analytic results required by these theorems, adapted existing
formal mathematics, and contributed foundations later used in Tau Ceti
(\cref{sec:related-formal,sec:foundations}).
\item \textbf{Source correction.} We identify a false comparison
of Ky Fan sums in the printed claim \citep[Proposition~4.4]{davis1970rotation},
give a real counterexample at every positive acute angle,
and prove the Q-norm consequence of the valid squared-displacement theorem
(\cref{sec:prop44}).
\item \textbf{A reduction to weighted matching.} We reduce unrestricted
real trace-norm transport to a weighted-matching problem on principal angles
(\cref{sec:example-transport}).

\end{enumerate}

\section{Mathematical background}
\label{sec:related}

The Davis--Kahan angle estimates use three constructions. Principal angles
measure the error between subspaces; a residual measures how well a proposed
subspace represents the operator; and a Sylvester estimate bounds the first
quantity by the second under spectral separation. Operator ideals allow the
same argument to express several norms. We describe the development of these
constructions before stating the formalized sine-theta theorem.

Existing historical surveys cover these developments. Paige and Wei trace
the cosine--sine decomposition from Jordan's subspace geometry through the
principal ideas assembled by Davis and Kahan to Stewart's later completion
and explicit formulation \citep{paige1994history}.
Simon's historical survey of Kato's work in non-relativistic quantum mechanics
traces Sz.-Nagy's unitary construction for pairs of projections and Kato's
later algebraic formulation \citep[\S5]{simon2018kato}.
We use this account for the history of the projection transport described in
\cref{sec:related-projections}.
B\"ottcher and Spitkovsky survey pairs of projections and their operator algebras.
Their historical discussion relates Halmos's 1969 block representation to the
earlier formulations of Dixmier (1948) and Davis (1958), and to the related
projection constructions in Davis and Kahan's 1970 paper
\citep[Remark~1.3]{bottcher2010gentle}.
We use these accounts together with the primary
sources. Table~\ref{tab:historical-lineage} summarizes the mathematical
constructions used in our formalization.

\begin{table}[!htbp]
\centering
\small
\setlength{\tabcolsep}{5pt}
\renewcommand{\arraystretch}{1.18}
\begin{tabularx}{\linewidth}{@{}>{\raggedright\arraybackslash}p{0.26\linewidth}>{\raggedright\arraybackslash}X>{\raggedright\arraybackslash}p{0.28\linewidth}@{}}
\toprule
\textbf{Lineage} & \textbf{Mathematical role} & \textbf{Formal dependency}\\
\midrule
Jordan (1875); Davis (1958); Halmos (1969) & Principal angles, invariant subspace geometry,
  and pairs of projections & Angle operators; projection blocks\\
Rayleigh (1877); Ritz (1909); Weyl (1912) & Approximate eigenpairs,
  variational spectra, and eigenvalue stability & Trial coordinates;
  min--max and eigenvalue inequalities\\
Rellich (1937--1942); Sz.-Nagy (1947); Kato (1966) & Moving spectral projections and unitary
  transport between nearby subspaces & Spectral calculus; polar
  direct rotation\\
Sylvester (1884); Rosenblum (1956, 1969) & Spectral separation controls an
  off-diagonal operator & Domain-aware Sylvester estimates\\
Von Neumann (1937); Schatten (1950); Ky Fan (1951) & Comparison of spectral errors
  through singular-value gauges & Approximation numbers;
  symmetric operator ideals\\
Davis (1963, 1965); Davis--Kahan (1969, 1970) & Residual-to-angle bounds, tangent estimates,
  rotation extremality, and spectral branch selection & Source-level
  sine, tangent, and direct-rotation theorems\\
\bottomrule
\end{tabularx}
\caption{Mathematical sources used in the Davis--Kahan formalization.
Dates identify the publications discussed in the text.}
\label{tab:historical-lineage}
\end{table}
\FloatBarrier

\subsection{Principal angles}
\label{sec:related-classical}

Camille Jordan studied the geometry of linear subspaces in $n$ dimensions in
  1875 \citep{jordan1875essai}.

For two lines, a single acute angle describes their relative position.
For subspaces of higher dimension, one needs a sequence of \emph{principal
  angles}.
The successive principal vectors are selected under orthogonality constraints,
  so their angles are independent of the coordinates used to describe either
  subspace.
This separates the geometry of a subspace from the arbitrary choice of its
  orthonormal basis.

For finite-dimensional subspaces $\sym{Usource}$ and $\sym{Vtarget}$ of equal dimension, choose
isometries $\sym{Xframe}$ and $\sym{Yframe}$ with those ranges.
The singular values of the overlap matrix are the principal cosines:
\begin{equation}
  \sym{SingularValue}(\like{Xframe}{\symglyph{Xframe}^*}\sym{Yframe})=\cos\like{theta}{\symglyph{theta}_j},
  \qquad 0\leq\like{theta}{\symglyph{theta}_j}\leq\frac{\sym{PiConstant}}{2}.
  \label{eq:history-principal-cosines}
\end{equation}

The corresponding principal vectors $\like{SourceVector}{\symglyph{SourceVector}_j}\in \sym{Usource}$ and $\like{TargetVector}{\symglyph{TargetVector}_j}\in \sym{Vtarget}$ satisfy
$\langle \like{SourceVector}{\symglyph{SourceVector}_i},\like{TargetVector}{\symglyph{TargetVector}_j}\rangle=\like{KroneckerDelta}{\symglyph{KroneckerDelta}_{ij}}\cos\like{theta}{\symglyph{theta}_j}$ in suitably oriented
coordinates.  For nonzero acute angles, each pair determines a real
two-dimensional plane in which the relative position is an ordinary
rotation.  An entire pair of subspaces is described by these planes,
together with their shared and orthogonal components.

Principal angles also arise in statistics: Hotelling's 1936 canonical
  correlation analysis sought pairs of maximally correlated linear combinations
  of two sets of variables \citep{hotelling1936relations}.
After appropriate whitening, the canonical correlations are singular values of a
  cross-coordinate map.
In numerical linear algebra, Bj\"orck and Golub developed algorithms for
  computing principal angles and vectors in 1973 \citep{bjorck1973numerical}.
Their computational formulation makes the basis-invariant angles available from
  orthonormal matrix factorizations.

A further formulation uses the cosine--sine decomposition of a
partitioned orthogonal or unitary matrix.
The blocks of the decomposition contain compatible diagonal cosine and sine
  factors.
Davis and Kahan describe the geometry using orthogonal projections,
angle operators, and the direct rotation carrying one subspace onto the other
\citep[Section~3]{davis1970rotation}. Their work assembled the principal ideas
of the cosine--sine decomposition; Stewart completed and explicitly formulated
the matrix decomposition in 1977
\citep{stewart1977perturbation,paige1994history}.
Paige and Wei explain how this decomposition and direct rotation represent
the same principal-angle data
\citep{paige1994history}. Our formalization uses projections and functional
calculus for the Hilbert-space results and principal-plane coordinates for
the finite-dimensional counterexample.

\subsection{Eigenvalue approximation}
\label{sec:related-ritz-weyl}

Jordan's angles answered how to compare subspaces once they are given.
An independent question came from mechanics and numerical analysis:
how can one approximate an eigenvalue or eigenmode when the underlying
differential equation cannot be solved exactly?  In the nineteenth
century Rayleigh estimated vibration frequencies using trial shapes
and their energy ratios \citep{rayleigh1877sound}.  Ritz developed a systematic finite-dimensional
variational approximation in the work published in 1909
\citep{ritz1909variations,leissa2005historical}.  The method selects
a trial space and studies the action of the operator through that
space. A trial space is a chosen space of candidate approximate solutions.
For an eigenvalue problem, one searches for an approximate eigenvector among
its vectors. This reduces an infinite-dimensional problem to a smaller
problem whose residual can then be assessed.

For a Hermitian matrix $\sym{A}$, the Rayleigh quotient
$\langle \sym{A}\sym{RayleighTestVector},\sym{RayleighTestVector}\rangle/\langle \sym{RayleighTestVector},\sym{RayleighTestVector}\rangle$ expresses the scalar action
of $\sym{A}$ along a trial vector. A normalized trial vector $\sym{v}$ gives the
Rayleigh--Ritz approximate eigenvalue $\sym{mu}=\langle \sym{A}\sym{v},\sym{v}\rangle$ and the
residual $\sym{A}\sym{v}-\sym{mu} \sym{v}$. The Courant--Fischer min--max
principle organizes the eigenvalues through extremal quotients over
subspaces \citep{bhatia1997matrix}.  Weyl's eigenvalue comparison
inequalities then control the motion of ordered eigenvalues under a
Hermitian perturbation \citep{weyl1912eigenwerte,bhatia1997matrix}.  In modern notation,
if $\sym{A}$ and $\sym{Acurrent}=\sym{A}+\sym{H}$ are Hermitian $n$-by-$n$ matrices and their eigenvalues
are arranged in the same order, then
\begin{equation}
  \max_j |\like{SpectralValue}{\symglyph{SpectralValue}_j}(\sym{A})-\like{SpectralValue}{\symglyph{SpectralValue}_j}(\sym{Acurrent})|\leq\dkOperatorNorm{\sym{H}}.
  \label{eq:history-weyl}
\end{equation}
This is an estimate on \emph{eigenvalues}; it does not tell us which
individual eigenvector of $\sym{A}$ corresponds to which eigenvector of $\sym{Acurrent}$.

The limitation already appears in two dimensions.  For any $\sym{EigenvalueExampleEpsilon}>0$,
put
\begin{equation}
 \like{A}{\symglyph{A}_\varepsilon}=\begin{pmatrix}0&0\\0&\sym{EigenvalueExampleEpsilon}\end{pmatrix},
 \qquad
 \like{Acurrent}{\symglyph{Acurrent}_\varepsilon}=\begin{pmatrix}\sym{EigenvalueExampleEpsilon}&0\\0&0\end{pmatrix}.
 \label{eq:history-swap-example}
\end{equation}
The matrices have identical ordered eigenvalues and
$\dkOperatorNorm{\like{A}{\symglyph{A}_\varepsilon}-\like{Acurrent}{\symglyph{Acurrent}_\varepsilon}}=\sym{EigenvalueExampleEpsilon}$,
yet their lowest-eigenvalue lines are perpendicular. The gap between the
two eigenvalues of each matrix is $\sym{EigenvalueExampleEpsilon}$, equal to the perturbation
norm. An eigenspace estimate therefore needs to compare the perturbation
size with the spectral separation.  For a cluster with repeated eigenvalues, an
entire invariant subspace is likewise more natural than a selected
orthonormal eigenbasis.

Rellich developed a rigorous perturbation theory of self-adjoint operators
in papers published from 1937 to 1942. Kato's 1966 book gives a systematic
account \citep{simon2018kato,kato1966perturbation}. The Davis--Kahan residual
bounds address a related approximation question: given a trial space and its
residual, spectral separation bounds its angle from an exact spectral
subspace. The estimate can therefore assess the accuracy of a
Rayleigh--Ritz approximation.

\subsection{Two orthogonal projections}
\label{sec:related-two-projections}

For a closed subspace of a Hilbert space $\sym{GeneralHilbertSpace}$, the orthogonal
projection remains a bounded operator even when the subspace is
infinite dimensional.  Denote the projections onto two closed
subspaces by $\sym{Psource}$ and $\sym{Qtarget}$. Compressing an operator to $\DKRange(\sym{Psource})$ means
restricting its input to that subspace and projecting its output back with
$\sym{Psource}$. The positive compression of $\sym{Qtarget}$ is therefore $\sym{Psource}\sym{Qtarget}\sym{Psource}$ restricted to
$\DKRange(\sym{Psource})$. It is the operator version of the squared overlap matrix: in finite dimensions its eigenvalues are
$\cos^2\like{theta}{\symglyph{theta}_j}$.  In Hilbert space it need not have a discrete
spectrum.  Functional calculus nevertheless defines its positive
square root and the associated cosine and sine of an angle operator.

Write $\sym{Theta0}$ for this directed angle operator on
$\DKRange(\sym{Psource})$, and $\sym{Theta1}$ for the corresponding
operator for the complementary pair $\sym{Identity}-\sym{Psource},\sym{Identity}-\sym{Qtarget}$.
The full angle operator is
\[
 \sym{Theta}=\sym{DiagonalOperator}(\sym{Theta0},\sym{Theta1})
 \quad\text{on }\DKRange(\sym{Psource})\oplus\DKKernel(\sym{Psource}).
\]
For $c>0$, the comparisons $\sym{Theta}\le c$ and $\sym{Theta}<c$ mean that its
operator norm is at most, or strictly less than, $c$, respectively.
In finite dimensions this norm is the largest eigenvalue of either directed
angle block, comparing the source with the target and their orthogonal
complements.
This includes right angles from unmatched directions: if a source vector is
orthogonal to the target, its cosine compression is zero and its directed
angle is $\sym{PiConstant}/2$. The compression construction applies even when the subspace
dimensions differ. Target directions orthogonal to the source appear as
right angles in the complementary block. Later, isometric frames express these angle blocks on
coordinate Hilbert spaces.

The relative position of two closed subspaces was studied by Dixmier
in 1948 and by Davis in his 1958 paper \emph{Separation of Two Linear
Subspaces} \citep{dixmier1948position,davis1958separation}.  Davis's
closeness and separation operators place the cosine/sine data in an
operator-theoretic setting.  Halmos's 1969 \emph{Two Subspaces} gives a
block representation of the two projections
\citep{halmos1969two}.  One first splits off the four intersection
spaces corresponding to the combinations of membership in the ranges
and kernels of $\sym{Psource}$ and $\sym{Qtarget}$.  On their remaining, generic part, the pair
is unitarily equivalent to
\begin{equation}
 \like{Psource}{\symglyph{Psource}_{\mathrm{gen}}}=\begin{pmatrix}\sym{Identity}&0\\0&0\end{pmatrix},
 \qquad
 \like{Qtarget}{\symglyph{Qtarget}_{\mathrm{gen}}}=\begin{pmatrix}\like{HalmosCosineBlock}{\symglyph{HalmosCosineBlock}^2}&\sym{HalmosCosineBlock}\sym{HalmosSineBlock}\\\sym{HalmosCosineBlock}\sym{HalmosSineBlock}&\like{HalmosSineBlock}{\symglyph{HalmosSineBlock}^2}\end{pmatrix},
 \qquad \like{HalmosCosineBlock}{\symglyph{HalmosCosineBlock}^2}+\like{HalmosSineBlock}{\symglyph{HalmosSineBlock}^2}=\sym{Identity},
 \quad \sym{HalmosCosineBlock}\sym{HalmosSineBlock}=\sym{HalmosSineBlock}\sym{HalmosCosineBlock},
 \label{eq:history-halmos-block}
\end{equation}
where $\sym{HalmosCosineBlock}$ and $\sym{HalmosSineBlock}$ are commuting positive contractions on a common
Hilbert space.  Thus an operator-valued angle replaces a finite list
of principal angles.  B\"ottcher and Spitkovsky survey this
representation and its consequences for the algebra generated by two
projections \citep{bottcher2010gentle}.

Among these four intersections, the two right-angle intersection
subspaces (called defect spaces below) are
\[
 \sym{CrossedSourceDefect}=\DKRange(\sym{Psource})\cap\DKKernel(\sym{Qtarget}),
 \qquad
 \sym{CrossedTargetDefect}=\DKKernel(\sym{Psource})\cap\DKRange(\sym{Qtarget}).
\]
These contain source directions orthogonal to the target, and target directions
orthogonal to the source, respectively. The pair is called \emph{acute} when
both are $\{\sym{ZeroVector}\}$, the zero subspace consisting only of the zero
vector. For a general pair these spaces affect whether a specified
intertwining operator exists. An operator $\sym{IntertwinerT}$ intertwines $\sym{Psource}$ and $\sym{Qtarget}$ when
$\sym{IntertwinerT}\sym{Psource}=\sym{Qtarget}\sym{IntertwinerT}$: applying $\sym{IntertwinerT}$ after the source projection agrees with applying the
target projection after $\sym{IntertwinerT}$. In particular, $\sym{IntertwinerT}$ maps the source range into
the target range and its orthogonal complement into the target complement.
The existence criteria for direct rotations include conditions on the
dimensions of these defect spaces \citep[Section~3]{davis1970rotation}.

A finite-dimensional orthogonal matrix decomposes into invariant real
lines and two-dimensional rotation planes.  In a general Hilbert space an
angle operator can instead have continuous spectrum.  The spectral
calculus of the pair of projections takes the place of a finite list
of angles, and the intersection and defect spaces cannot be discarded.
Our Hilbert-space direct-rotation proofs therefore use projection identities
and functional calculus. In finite dimensions these identities recover the
principal-plane construction.

\subsection{Transport between nearby spectral subspaces}
\label{sec:related-projections}

A perturbation can change the invariant subspace associated with an isolated
part of a self-adjoint operator's spectrum. Let $\sym{Psource}$ and
$\sym{Qtarget}$ be the orthogonal projections onto the corresponding spectral
subspaces before and after perturbation. To compare these subspaces, we seek
a unitary carrying $\DKRange(\sym{Psource})$ onto
$\DKRange(\sym{Qtarget})$ that depends continuously, or analytically, on the
projections.
Sz.-Nagy developed the nearby-projection construction in work
published in 1947 \citep{sznagy1947perturbations,simon2018kato}.  Kato gave a streamlined
algebraic construction, described in his 1955 report and later in
\emph{Perturbation Theory for Linear Operators}
\citep[Chapter~I, \S4.6, equations~(4.33)--(4.42)]{kato1966perturbation};
see also \citet[Theorem~5.1]{simon2018kato}.

For orthogonal projections the useful intertwiner is
\begin{equation}
 \sym{IntertwinerT}=\sym{Qtarget}\sym{Psource}+(\sym{Identity}-\sym{Qtarget})(\sym{Identity}-\sym{Psource}), \qquad \sym{IntertwinerT}\sym{Psource}=\sym{Qtarget}\sym{IntertwinerT}.
 \label{eq:history-kato-T}
\end{equation}
The identity
\begin{equation}
 \like{IntertwinerT}{\symglyph{IntertwinerT}^*}\sym{IntertwinerT}=\sym{Identity}-(\sym{Psource}-\sym{Qtarget})^2
 \label{eq:history-kato-square}
\end{equation}
shows why $\|\sym{Psource}-\sym{Qtarget}\|<1$ is the natural threshold: it makes $\like{IntertwinerT}{\symglyph{IntertwinerT}^*}\sym{IntertwinerT}$
strictly positive and invertible.  Hence
\begin{equation}
 \sym{Rrotation}=\sym{IntertwinerT}\bigl(\sym{Identity}-(\sym{Psource}-\sym{Qtarget})^2\bigr)^{-1/2},
 \qquad \sym{Rrotation}\sym{Psource}\like{Rrotation}{\symglyph{Rrotation}^*}=\sym{Qtarget}
 \label{eq:history-kato-unitary}
\end{equation}
is unitary. This is the polar decomposition $\sym{IntertwinerT}=\sym{Rrotation}|\sym{IntertwinerT}|$, where
$|\sym{IntertwinerT}|=(\like{IntertwinerT}{\symglyph{IntertwinerT}^*}\sym{IntertwinerT})^{1/2}$ is positive. The factor $|\sym{IntertwinerT}|$ describes stretching;
the unitary factor $\sym{Rrotation}$ preserves lengths. Here $\sym{IntertwinerT}$ is invertible, so
$\sym{Rrotation}=\sym{IntertwinerT}|\sym{IntertwinerT}|^{-1}$. In principal-plane coordinates it is the direct rotation through the
principal angle, acting trivially on common directions.

The transport equation
$\sym{Wcompetitor}\sym{Psource}\sym{Wcompetitor}^*=\sym{Qtarget}$
has many unitary solutions.
The polar-factor formula selects one determined by the two projections
and continuously connected to the identity under the norm bound above. That
choice makes the construction valuable in perturbation theory: the
transport identifies moving spectral subspaces while avoiding an
arbitrary choice of orthonormal bases. In infinite dimensions,
$\|\sym{Psource}-\sym{Qtarget}\|<1$ is stronger than acuteness: it also
bounds the operator angles away from $\sym{PiConstant}/2$. For equal-dimensional
finite subspaces the conditions coincide. For an orthogonal or unitary transporter $\sym{Wcompetitor}$, the operator
$\sym{Identity}-\sym{Wcompetitor}$ records the displacement of every vector. Davis and Kahan ask which
norms of this displacement are minimized by the direct rotation
\citep[Section~4]{davis1970rotation}.

Kato used such transformations to identify nearby spectral
subspaces of analytic operator families, thereby reducing the study
of an isolated moving eigenvalue cluster to a fixed coordinate
space.  Simon's historical account explains the earlier Sz.-Nagy
construction, Kato's formula, and subsequent work on pairs of
projections \citep[\S5]{simon2018kato}. Our proofs of direct-rotation
existence, uniqueness, and symmetry use the square-root and
polar-decomposition identities in
\cref{eq:history-kato-square,eq:history-kato-unitary}. The condition
$\|\sym{Psource}-\sym{Qtarget}\|<1$ applies in arbitrary Hilbert dimension.

\subsection{Spectral projections and residuals}
\label{sec:related-spectral}

The spectral theorem represents a self-adjoint operator as
$\sym{A}=\int\sym{SpectralValue}\,d\sym{SpectralMeasure}(\sym{SpectralValue})$ for a projection-valued measure $\sym{SpectralMeasure}$.
For a Borel set $\sym{SpectralSet}$, the operator $\sym{SpectralMeasure}(\sym{SpectralSet})$ is the orthogonal projection onto
the corresponding spectral subspace. This gives a second way of obtaining
a subspace: select a region of the spectrum and take the range of its
spectral projection. The functional-calculus
framework developed from the work of Riesz, von Neumann, and others;
Kato made spectral projections central to the perturbation theory of
closed operators \citep{kato1966perturbation,simon2018kato}.
Unlike the finite-dimensional eigenspace picture, a spectral
subspace need not be spanned by a finite list of eigenvectors.
A self-adjoint operator may be unbounded, while each of its spectral
projections is bounded and everywhere defined. The Hellinger--Toeplitz
theorem states that an everywhere-defined symmetric operator on a Hilbert
space is bounded \citep{hellinger1910grundlagen}. Thus an unbounded
self-adjoint operator has a proper domain. Its action $\sym{A}x$ is defined only
for $x\in\DKDom(\sym{A})$, so residual and intertwining equations must include
the required domain inclusions.

In the complex case, a number $z$ belongs to the resolvent set when
$z\sym{Identity}-\sym{A}$ has a bounded, everywhere-defined inverse. That inverse,
$(z\sym{Identity}-\sym{A})^{-1}$, is the resolvent at $z$. For an isolated bounded spectral
cluster, choose a positively oriented contour $\sym{RieszContour}$ in the resolvent
set enclosing the cluster and no other spectrum. Its projection is
\begin{equation}
 \sym{RieszProjection}=\frac{1}{2\sym{PiConstant} i}\int_{\sym{RieszContour}}(z\sym{Identity}-\sym{A})^{-1}\,dz.
 \label{eq:history-riesz-projection}
\end{equation}
The resolvent is bounded at each point of $\sym{RieszContour}$ even when $\sym{A}$ is unbounded.
This makes the spectral projection a stable object for perturbation analysis,
  while preserving the separate domain requirements whenever $\sym{A}$ itself appears
  in an identity.
More general Borel spectral projections cover finite and half-infinite intervals,
  including spectral configurations for which a single isolated
  finite-eigenvalue cluster is not the right model.
Our formalization uses Borel functional calculus to construct projections for
  these spectral sets, including for unbounded self-adjoint operators.

The Courant--Fischer min--max principle and Weyl-type eigenvalue
perturbation bounds provide complementary information about spectral
\emph{values} \citep{bhatia1997matrix,kato1966perturbation}.  They do
not by themselves identify a stable basis inside a multiple-eigenvalue
cluster, or quantify the movement of its invariant subspace.  The Courant--Fischer and Weyl results that we formalized describe extremal
Rayleigh quotients and eigenvalue stability. The angle estimates add control
of the associated subspaces.

Davis studied the rotation of eigenvectors under perturbation in
1963 and 1965 \citep{davis1963rotation,davis1965rotation}.  The
1969 Davis--Kahan announcement preceded their extended 1970 treatment
\citep{daviskahan1969bounds,davis1970rotation}.  Their problem has
two closely related formulations.  In one, the original self-adjoint
operator is perturbed and the movement of an exact spectral subspace
is estimated.  In the other, one has an approximate trial subspace and
a computable residual: how far can it be from the desired exact
spectral subspace?  Unlike a comparison of individual eigenvectors,
these are meaningful even when an eigenvalue has multiplicity and its
orthonormal eigenbasis is nonunique.

Write $\sym{A}$ for a reference operator and $\sym{Acurrent}=\sym{A}+\sym{H}$ for the perturbed
ambient operator. We use this convention in the sine-theta statement below.
Let $\sym{E0}$ embed a trial-coordinate Hilbert space into the ambient space
isometrically, and let $\sym{TrialOperator}$ be a self-adjoint operator on those coordinates.
The residual is the bounded operator agreeing, where defined, with
\begin{equation}
 \sym{R}=\sym{Acurrent}\sym{E0}-\sym{E0}\sym{TrialOperator}.
 \label{eq:history-residual}
\end{equation}
If $\sym{R}=0$, applying $\sym{Acurrent}$ after the embedding agrees with applying
$\sym{TrialOperator}$ before it: the trial subspace is invariant and $\sym{TrialOperator}$ represents its
operator. For $\sym{R}\ne0$, the residual measures the failure of this identity.  A second isometric coordinate map for the
desired exact spectral space defines its orthogonal projection $\sym{Pexact}$.
The positive operator
\begin{equation}
 \sym{sinTheta0}=\bigl(\sym{Identity}-\like{E0}{\symglyph{E0}^*}\sym{Pexact}\sym{E0}\bigr)^{1/2}
 \label{eq:history-angle-sine}
\end{equation}
measures the component of a trial vector lying outside the exact
subspace.  The critical distinction for the formal proof is that this
positive sine acts on the trial-coordinate space, whereas
$(\sym{Identity}-\sym{Pexact})\sym{E0}$ maps into the ambient space.  For an operator $B$, its positive modulus $|B|=(B^*B)^{1/2}$ acts on its
input space and records its stretching factors. The modulus of $(\sym{Identity}-\sym{Pexact})\sym{E0}$
is $\sym{sinTheta0}$, so the two maps have equal unitarily invariant ideal
norms whenever those norms are finite.

\subsection{Sylvester estimates}
\label{sec:related-sylvester}

The matrix equation $\sym{SylvesterLeft}\sym{X}-\sym{X}\sym{HistoricalSylvesterRight}=\sym{SylvesterRHS}$ is associated with Sylvester's 1884
work \citep{sylvester1884equation}.  Rosenblum studied separated
operator spectra in 1956 and the equation with self-adjoint coefficients
in 1969 \citep{rosenblum1956equation,rosenblum1969selfadjoint}.

The scalar prototype is $(\sym{SylvesterScalarLeft}-\sym{SylvesterScalarRight})\sym{SylvesterScalarUnknown}=\sym{SylvesterScalarRHS}$: if $|\sym{SylvesterScalarLeft}-\sym{SylvesterScalarRight}|\geq\sym{delta}>0$,
then $\sym{delta}|\sym{SylvesterScalarUnknown}|\leq|\sym{SylvesterScalarRHS}|$.  In the operator problem the scalar
multipliers are replaced by the spectral actions of two operators.
The operator estimate has the form $\sym{delta} \sym{NormN}(\sym{X})\leq \sym{NormN}(\sym{SylvesterLeft}\sym{X}-\sym{X}\sym{HistoricalSylvesterRight})$.
Its hypotheses specify the spectral positions of $\sym{HistoricalSylvesterRight}$ and $\sym{SylvesterLeft}$, their domains
when unbounded, and the unitarily invariant norm $\sym{NormN}$.

Let $\sym{F1}$ give isometric coordinates for the exact complementary
spectral subspace, with self-adjoint block $\sym{Lambda1}$.  Writing
$\sym{X}=\like{F1}{\symglyph{F1}^*}\sym{E0}$ converts the residual identity into a Sylvester equation:
\begin{equation}
 \sym{Lambda1}\sym{X}-\sym{X}\sym{TrialOperator}=\like{F1}{\symglyph{F1}^*}\sym{R}
 \label{eq:history-sylvester}
\end{equation}
on the appropriate domain. Thus $\sym{X}$ represents the trial frame's
``leakage'' (i.e., its orthogonal projection onto the exact complementary
spectral subspace, expressed in $\sym{F1}$ coordinates), while the right
side is controlled by the computable residual.  Under the interval/exterior
or ordered spectral-gap hypotheses of the sine-theta theorem,
separated Sylvester estimates give
\begin{equation}
 \sym{delta}\,\sym{NormN}(\sym{X})\leq \sym{NormN}(\like{F1}{\symglyph{F1}^*}\sym{R})\leq \sym{NormN}(\sym{R}).
 \label{eq:history-separation-bound}
\end{equation}
This is the bridge from a spectral-gap assumption to the bound on the
sine of a subspace angle.

The coefficient of $\sym{NormN}(\sym{R})/\sym{delta}$ in this bound is one.
This coefficient uses the ordered or interval/exterior spectral
positions specified by Davis and Kahan. Arbitrarily interlaced spectra can
require a different constant. Our theorem includes both of the source's
spectral configurations, including ordered cases where both operators are
unbounded.  A formal implementation
therefore needs the domains of unbounded self-adjoint operators,
their spectral projections, strong spectral cutoffs, and the
ideal-norm inequalities used to pass through bounded approximations.

\subsection{Operator ideals and displacement norms}
\label{sec:related-norms}

The geometry of a pair of subspaces is fixed, but there are many ways
to assign a cost to a unitary transporting one onto the other.
Von Neumann's 1937 study of unitarily invariant matrix norms identified
the role of symmetric gauge functions of singular values
\citep{vonneumann1937matrix}. Schatten's 1950 monograph developed normed
operator ideals; Ky Fan's 1951 eigenvalue and singular-value inequalities
provided the partial-sum comparisons underlying majorization
arguments \citep{schatten1950crossspaces,fan1951maximum,bhatia1997matrix}.
The singular values $\sym{SingularValue}(Z)$ are the nonnegative square roots of the
eigenvalues of $Z^*Z$, ordered decreasingly. They record the stretching
factors of the matrix $Z$. For $1\leq p<\infty$, the Schatten norms include
\begin{equation}
 \|Z\|_{S_p}=\left(\sum_j \sym{SingularValue}(Z)^p\right)^{1/p},
 \qquad
 \|Z\|_{S_1}=\sum_j \sym{SingularValue}(Z),
 \quad
 \|Z\|_{S_2}=\left(\sum_j \sym{SingularValue}(Z)^2\right)^{1/2}.
 \label{eq:history-schatten}
\end{equation}
An operator ideal is a linear class of operators closed under multiplication
on either side by bounded operators. A normed ideal also controls the norm
of these products. For compact operators in infinite dimensions, the
Schatten $p$ ideal consists of those whose singular values have a finite
$p$-power sum. Thus a trace-norm inequality requires trace-class membership;
this is an instance of the source's qualification that its norms be defined.

The largest singular value gives the operator norm, the Euclidean
norm of the singular-value sequence gives the Hilbert--Schmidt norm
(the Frobenius norm $\dkFrobeniusNorm{Z}=\|Z\|_{S_2}$ for matrices),
and its sum gives the trace norm. These norms can rank transporters
differently. In Section~\ref{sec:prop44}, the two displacement singular-value
lists are $(2\sin(\sym{theta}/2),\ldots,2\sin(\sym{theta}/2))$ with four entries and
$(2\sin\sym{theta},2\sin\sym{theta},0,0)$. The second has a smaller sum and a larger
maximum. Ky Fan norms compare partial sums of the ordered singular values;
these comparisons determine which norm inequalities follow.

Davis established Hilbert--Schmidt and finite-dimensional Schatten
extremal properties for exponents at least two in his 1958 subspace
paper \citep[Theorems~7.3--7.4]{davis1958separation}.  Davis and
Kahan compare three objectives \citep[Section~4]{davis1970rotation} for a transporter $\sym{Wcompetitor}$ carrying
$\sym{Usource}$ onto $\sym{Vtarget}$: its restricted displacement $(\sym{Identity}-\sym{Wcompetitor})\like{Psource}{\symglyph{Psource}_U}$, its squared full
displacement $(\sym{Identity}-\sym{Wcompetitor})^*(\sym{Identity}-\sym{Wcompetitor})$, and its full displacement $\sym{Identity}-\sym{Wcompetitor}$. The last acts
on every ambient vector, including vectors outside $\sym{Usource}$. They gave a
complex-phase modification that decreases its trace norm \citep[Example~4.2]{davis1970rotation}.
The full-displacement claim \citep[Proposition~4.4]{davis1970rotation}
asserts that direct rotation minimizes
full displacement for real subspaces if their principal angles are no
larger than $\sym{PiConstant}/3$.  The four-dimensional example in
\cref{sec:prop44} gives a smaller trace norm at every positive acute
angle.  It exposes the failed step: adding singular-value sums of
restrictions to separate principal planes is not a valid global lower
bound for an arbitrary ambient transporter.

The subsequent minimal-distortion theory of Bolotnikov, Li, and Rodman
studied the same class of subspace-transport constraints, including
Frobenius and Q-norm optimality and explicit subquadratic
counterexamples \citep{bolotnikov2001minimal}.  Later work on
Grassmannian metrics and canonical-angle inequalities also separates
intrinsic paths between subspaces from minimizing the extrinsic
operator displacement \citep{qiu2006grassmann,zhang2010canonicalangles}.
These results give the norm-theoretic context for the source
correction and for the trace-norm matching problem in
\cref{sec:example-transport}.

\subsection{The Davis--Kahan angle estimates}
\label{sec:related-synthesis}

Davis's papers of 1963 and 1965 and the Davis--Kahan announcement of 1969
preceded the extended 1970 treatment
\citep{davis1963rotation,davis1965rotation,daviskahan1969bounds,davis1970rotation}.
For the sine-theta estimate, let $\sym{R}=\sym{Acurrent} \sym{E0}-\sym{E0}\sym{TrialOperator}$ be the residual
and $\sym{Pexact}$ the desired exact spectral projection. The operator $(\sym{Identity}-\sym{Pexact})\sym{E0}$
extracts the part of a trial vector outside that spectral subspace. In
isometric coordinates $\sym{F1}$ for the complementary space, this component is
represented by $\sym{X}=\like{F1}{\symglyph{F1}^*}\sym{E0}$. The residual identity gives
\cref{eq:history-sylvester}, and a separated Sylvester estimate bounds $\sym{NormN}(\sym{X})$.
Since $\sym{S}=\sym{F1}\sym{X}$ and $\sym{F1}$ is an isometry,
$\sym{NormN}(\sym{X})=\sym{NormN}(\sym{S})$.
The modulus identity $|\sym{S}|=\sym{sinTheta0}$ gives
$\sym{NormN}(\sym{S})=\sym{NormN}(\sym{sinTheta0})$.

The resulting chain of equalities and inequalities is
\begin{equation}
 \underbrace{\sym{delta} \sym{NormN}(\sym{sinTheta0})}_{\text{subspace error}}
 =\sym{delta} \sym{NormN}(\sym{X})
 \leq \sym{NormN}(\like{F1}{\symglyph{F1}^*}\sym{R})
 \leq\underbrace{\sym{NormN}(\sym{R})}_{\text{computable residual}}.
 \label{eq:history-unified-chain}
\end{equation}
The first equality follows from the positive-modulus identity for the
complementary projection. The first inequality applies the Sylvester
estimate to $\sym{X}$; the second uses $\sym{NormN}(BC)\leq\dkOperatorNorm{B}\sym{NormN}(C)$ and
$\dkOperatorNorm{\like{F1}{\symglyph{F1}^*}}\leq1$. For unbounded operators the Sylvester
identity holds on the trial operator's domain, with a bounded residual.
These are the geometric, spectral, and norm estimates formalized in our
proof of the theorem.

The tangent and double-angle estimates use related projection blocks.
For example, an estimate for $\sin(2\sym{Theta})$ allows angles near zero and
near $\sym{PiConstant}/2$. The spectral-selection results \citep[Section~8]{davis1970rotation}
identify which spectral subspaces give the small-angle alternative.
The direct-rotation analysis \citep[Section~4]{davis1970rotation}
studies the orthogonal or unitary map carrying one subspace onto another
and asks which displacement norms the direct rotation minimizes.

\subsection{Numerical and statistical applications}
\label{sec:related-numerical}

The 1970 theory influenced more than the analysis of exact
self-adjoint eigenvectors.  Wedin's 1972 singular-subspace
perturbation bounds address the corresponding problem for singular
vectors of general rectangular matrices
\citep{wedin1972perturbation}.  Bj\"orck and Golub's angle algorithms
made the subspace quantities computable
\citep{bjorck1973numerical}.  Stewart's work on pseudoinverses,
projections, and least-squares sensitivity helped connect these
geometric notions with matrix computations
\citep{stewart1977perturbation,stewart1990matrix}.
Paige and Wei's 1994 history shows how the cosine--sine decomposition
links the principal-angle geometry of Davis and Kahan with numerical
matrix factorizations \citep{paige1994history}.

The same subspace-versus-eigenvector distinction remains central in
statistical spectral estimation.  For example, Yu, Wang, and
Samworth give a statistical variant whose eigenspace error bound uses
separation in the population spectrum \citep{yu2015davis}, while Chen and
collaborators survey spectral methods across statistical learning
problems \citep{chen2021spectral}.  In that setting, an empirical covariance matrix approximates a population
operator and the perturbation norm controls the error of estimated
eigenspaces. Our general theorem supplies the deterministic angle estimate;
using it in a statistical result also requires the corresponding probabilistic
control of the perturbation.

\subsection{Foundations used in our formalization}
\label{sec:related-formal}

Each historical component above corresponds to a specific requirement
in the formalization.  Jordan's angle geometry and Halmos's
projection blocks lead to positive principal-angle operators rather
than chosen eigenbases.  The Sz.-Nagy--Kato intertwiner requires
polar decomposition, square roots, and real/complex transfer.
Spectral projections require Borel functional calculus, including for
unbounded self-adjoint operators.  The Sylvester proof requires a
representation carrying operator domains and a controlled passage
through spectral cutoffs.  Schatten and Ky Fan theory requires
approximation numbers and symmetric operator ideals to express
inequalities with the source's norm generality.

We formalized principal-angle geometry, symmetric operator ideals, and
unbounded Sylvester estimates using Lean~4 and Mathlib
\citep{demoura2021lean4,mathlib2020}. Mathlib supplies the algebraic,
topological, and Hilbert-space framework. Its continuous functional calculus
is described in the preprint by Anatole Dedecker and Jireh Loreaux
(\href{https://arxiv.org/abs/2501.15639}{arXiv:2501.15639}).
Our projection-valued-measure and Borel-calculus
work also adapts Adam Bornemann's Spectra formalization
\citep{bornemann2026spectra}.

Our Courant--Fischer and Weyl development, rectangular unitarily invariant
seminorms, projection-valued measures, and supporting polar-decomposition
results were subsequently adapted into Tau Ceti \citep{tauceti2026}.
% Editorial note (not manuscript prose): Do not add a sentence emphasizing
% the current Tau Ceti canonical polar-decomposition interface here.
% This paragraph records work developed in DKPS and subsequently
% contributed upstream; current library interfaces are discussed elsewhere.
% Evidence: dev/tauceti/dkps-upstream-contribution-census.md, PRs 13599,
% 14241, 17527 (adaptations), and 16127/16450 (independent polar proofs).
% AUTHORING_NOTES.md indexes accepted-source and campaign evidence.
% Mathematical origin and formal proof provenance are separate attributions.

This history also separates the two mathematical questions that meet
in the project.  The sine-theta theorem estimates \emph{how far} an
approximate or perturbed spectral subspace lies from an exact one.
The direct-rotation problem asks \emph{which} ambient transformation
carries one given subspace onto another at minimum displacement cost.
Our proofs of the optimization results use the principal-angle,
projection, and operator-ideal foundations developed for the formalization.
The matching theorem has its own proof; the shared foundations explain the
relationship between the two parts of the work. The resulting counterexample
to the full-displacement claim \citep[Proposition~4.4]{davis1970rotation}
and the follow-on matching theorem are examined in
\cref{sec:prop44,sec:example-transport}.

\section{Davis--Kahan formalization and the sine-theta theorem}
\label{sec:coverage}

Our targets are the four headline theorems \citep[Section~2]{davis1970rotation} and every
named theorem, proposition, lemma, and corollary established by Davis and
Kahan in the 1970 paper. These \DKResultCount{} results include the sine and
tangent inequalities, the geometry and extremal properties of direct
rotations, bounded and unbounded Sylvester estimates, the generalized
perturbation bounds \citep[Section~6]{davis1970rotation}, and the spectral-selection
theorems \citep[Section~8]{davis1970rotation}. We prove \DKProvedCount{} results and refute the printed
full-displacement claim \citep[Proposition~4.4]{davis1970rotation}.
Appendix~\ref{app:dk-inventory} lists these targets.

We develop the sine-theta theorem and its proof below; the remaining source
results are summarized in Appendix~\ref{app:dk-inventory}, with hypotheses
and conclusions in Section~\ref{app:dk-result-summaries}. The main text also
explains the spectral-selection results needed to interpret the angle bounds.
Appendix~\ref{app:notation} collects the notation used here and in the
direct-rotation argument.

\subsection{The sine-theta theorem for self-adjoint operators}
\label{sec:sine-theta}

\subsubsection{Trial and exact subspaces}

We use $\sym{A}$ for the reference operator and
$\sym{Acurrent}=\sym{A}+\sym{H}$ for the perturbed ambient operator.
The symbol $\sym{Acurrent}$ corresponds to $\sym{A}+\sym{H}$ in the source.

For a unit trial vector $\sym{v}$ and an approximate eigenvalue $\sym{mu}$, the residual
is $\sym{Acurrent}\sym{v}-\sym{mu} \sym{v}$. An isometric coordinate map $\sym{E0}$ generalizes
this construction to a trial subspace: its range consists of the candidate
approximate vectors, and a self-adjoint operator $\sym{TrialOperator}$ represents their
approximate dynamics. For a finite-dimensional Rayleigh--Ritz choice,
$\sym{TrialOperator}=\like{E0}{\symglyph{E0}^*}\sym{Acurrent}\sym{E0}$. The theorem allows other self-adjoint trial
operators as well. We write $\sym{TrialOperator}$ for the general trial operator,
reserving $\sym{ReferenceBlock0}$ for the reference spectral block that
Davis and Kahan also call $A_0$.

Let $\sym{AmbientHilbertSpace}$ be the ambient Hilbert space.  Let $\sym{TrialOperator}$ be a self-adjoint operator on a
trial-coordinate space $\sym{TrialCoordinateSpace}$, and let
\[
  \sym{E0}:\sym{TrialCoordinateSpace}\to \sym{AmbientHilbertSpace}
\]
be an isometric coordinate map for the trial subspace.  The residual compares
applying the ambient operator after embedding with applying the trial operator
before embedding:
\begin{equation}
  \sym{R} = \sym{Acurrent}\sym{E0}-\sym{E0}\sym{TrialOperator}.
  \label{eq:residual}
\end{equation}
When the operators are unbounded, \cref{eq:residual} is asserted on
$\DKDom(\sym{TrialOperator})$ and $\sym{R}$ is the bounded residual that agrees there.
The Lean predicate \texttt{IsTrialResidual} packages the isometry, the required
domain inclusion, and this residual identity.

Let $\sym{F0}:\sym{ExactCoordinateSpace}\to \sym{AmbientHilbertSpace}$ be isometric coordinates for the desired exact spectral
subspace and let $\sym{F1}:\sym{ComplementaryCoordinateSpace}\to \sym{AmbientHilbertSpace}$ be isometric coordinates for its orthogonal
complement.  A self-adjoint operator $\sym{Lambda1}$ on $\sym{ComplementaryCoordinateSpace}$ describes the
complementary spectral block through
\[
  \sym{Acurrent}\sym{F1} = \sym{F1}\sym{Lambda1}
\]
on its domain. This intertwining identity says that $\sym{Lambda1}$ gives
the action of $\sym{Acurrent}$ in complementary coordinates. The predicate
\texttt{IsExactSpectralDecomposition} also requires the domain inclusion
$\sym{F1}(\DKDom(\sym{Lambda1}))\subseteq\DKDom(\sym{Acurrent})$
and states orthogonality and completeness through
\[
 \like{F0}{\symglyph{F0}^*}\sym{F1}=0,\qquad
 \sym{F0}\like{F0}{\symglyph{F0}^*}+\sym{F1}\like{F1}{\symglyph{F1}^*}=\sym{Identity}.
\]

\subsubsection{The sine operator and its rectangular representative}

Here $\sym{Theta0}$ acts on $\sym{TrialCoordinateSpace}$; under the isometry
$\sym{E0}$ it represents the directed angle operator on the trial subspace
introduced in \cref{sec:related-two-projections}. If
$\sym{Pexact}=\sym{F0}\like{F0}{\symglyph{F0}^*}$ is the orthogonal projection onto the desired exact subspace, then
\[
  \cos\sym{Theta0} = (\like{E0}{\symglyph{E0}^*}\sym{Pexact}\sym{E0})^{1/2},
  \qquad
  \sym{sinTheta0} = (\sym{Identity}-\like{E0}{\symglyph{E0}^*}\sym{Pexact}\sym{E0})^{1/2}.
\]
The analytic proof uses the rectangular complementary block
\begin{equation}
  \sym{S}=(\sym{Identity}-\sym{Pexact})\sym{E0}.
  \label{eq:rectangular-sine}
\end{equation}
The theorem uses the positive sine on the trial-coordinate space.  The rectangular block has codomain $\sym{AmbientHilbertSpace}$, while the
positive sine has codomain $\sym{TrialCoordinateSpace}$.  Their relationship is
\[
  \sym{S}^*\sym{S}
  = \like{E0}{\symglyph{E0}^*}(\sym{Identity}-\sym{Pexact})\sym{E0}
  = \sym{Identity}-\like{E0}{\symglyph{E0}^*}\sym{Pexact}\sym{E0},
\]
so
\begin{equation}
  |\sym{S}|=(\sym{S}^*\sym{S})^{1/2}=\sym{sinTheta0}.
  \label{eq:modulus-sine}
\end{equation}
% It might be worth noting that our formalization actually defines \sym{Theta0}
% and proves eq:modulus-sine whereas most treatments of the subject seam to
% just start from |\sym{S}|. Maybe it's worth enough to put in the paper, but
% it at least deserves a latex comment!
The polar decomposition $\sym{S}=\sym{PolarU}|\sym{S}|$ separates the positive
  modulus from a partial isometry.
Unitarily invariant ideal norms assign the same value to an operator and its
  modulus, so \[ \sym{NormN}(\sym{S})=\sym{NormN}(\sym{sinTheta0}) \] whenever
  the norm is finite.
This equality transfers an estimate for the complementary block $\sym{S}$ to the
  positive sine operator.

\subsubsection{Gap geometry and the inequality}
\label{sec:gap-geometry}

Fix $\sym{delta}>0$. There are two geometric forms of separation.
\emph{Interval/exterior separation:} for some real
$\sym{betaEndpoint}\leq\sym{alphaEndpoint}$, one of the two spectra
$\sym{SpectrumSymbol}(\sym{TrialOperator})$ and $\sym{SpectrumSymbol}(\sym{Lambda1})$ lies in $[\sym{betaEndpoint},\sym{alphaEndpoint}]$ and the other
lies in
\[
 (-\infty,\sym{betaEndpoint}-\sym{delta}]\;\cup\;[\sym{alphaEndpoint}+\sym{delta},\infty).
\]
\emph{Ordered separation:} for some real $b$, the operator forms satisfy either
\[
 \sym{TrialOperator}\geq (b+\sym{delta})\sym{Identity},\quad \sym{Lambda1}\leq bI,
 \qquad\text{or}\qquad
 \sym{TrialOperator}\leq bI,\quad \sym{Lambda1}\geq (b+\sym{delta})\sym{Identity}.
\]
Each form inequality is asserted on the domain of the corresponding operator;
for example, $\sym{TrialOperator}\geq aI$ means
$\langle \sym{TrialOperator}x,x\rangle\geq a\|x\|^2$ for every
$x\in\DKDom(\sym{TrialOperator})$. For self-adjoint operators, the source's ordered
hypotheses on half-infinite spectral intervals, such as $(-\infty,b]$ and
$[b+\sym{delta},\infty)$, imply these form bounds. These cases are included
in the source's unbounded sine theorem
\citep[Section~6, discussion of unbounded operators]{davis1970rotation}.
Write $\sym{SineGap}(\sym{TrialOperator},\sym{Lambda1})$ for their union.
The Lean predicate \texttt{FormBoundedSylvesterGap} has three cases:
the finite-interval alternative, which includes either spectral order,
and the two ordered form alternatives.

Both blocks may be unbounded in an ordered configuration. For instance,
on $\ell^2(\mathbb N_{\geq1})$, take diagonal operators
$\sym{TrialOperator}e_n=(n+\sym{delta})e_n$ and $\sym{Lambda1}e_n=-ne_n$, on the domains where their
images are square summable. They satisfy the first ordered configuration
with $b=0$, while neither spectrum lies in a finite interval. This example
explains why retaining the half-infinite alternatives changes the theorem's
applicability.

For the unbounded trial configuration, the source assumes that
$\sym{Acurrent}\sym{E0}$ and $\sym{E0}\sym{TrialOperator}$ have a common dense domain.
Since $\sym{E0}$ is everywhere defined, this gives
\[
 \DKDom(\sym{Acurrent}\sym{E0})=\DKDom(\sym{TrialOperator}),
 \qquad
 \sym{E0}(\DKDom(\sym{TrialOperator}))\subseteq\DKDom(\sym{Acurrent}).
\]
The residual agrees with the difference on this domain and extends to a
bounded operator. The formal predicate \texttt{IsTrialResidual} includes the
domain inclusion and the residual identity.

The norm $\sym{NormN}$ is a normalized symmetric operator-ideal norm.
Normalization means that the rank-one map $uv^*:x\mapsto u\langle v,x\rangle$
has norm
\[
 \sym{NormN}(uv^*)=\|u\|\,\|v\|;
\]
in particular, its norm is one when $u,v$ are unit vectors.
Its hypotheses also include Ky Fan dominance whenever both ideal norms
are finite: domination of all finite singular-value sums implies domination
of the ideal norms. This comparison does not assert ideal membership.
The norm is invariant under isometries and positive modulus, and satisfies
$\sym{NormN}(BT)\leq\dkOperatorNorm{B}\sym{NormN}(T)$. Trace-class and
Hilbert--Schmidt operators with their respective norms are examples; the
operator norm on bounded operators is another. For an ideal norm, both the
residual and the positive sine must belong to its domain.

In two dimensions, take a one-dimensional trial frame $\sym{E0}(x)=x\sym{v}$ with
$\|\sym{v}\|=1$. A scalar trial operator $\sym{TrialOperator}=\sym{rho}$ gives
$\sym{r}=\sym{R}(1)=\sym{Acurrent}\sym{v}-\sym{rho}\sym{v}$.
Choosing $\sym{rho}=\langle\sym{v},\sym{Acurrent}\sym{v}\rangle$
makes $\sym{rho}\sym{v}$ the orthogonal projection of
$\sym{Acurrent}\sym{v}$ onto the trial line. The residual is therefore
perpendicular to that line. Figure~\ref{fig:sine-theta-residual} pairs this
residual triangle with the exact complementary projection that measures
the subspace error.

\begin{figure}[tb]
\centering
\resizebox{119mm}{!}{\input{figures/Fig1.pgf}}
\caption{Rayleigh residual geometry for a one-dimensional trial space in
$\sym{RealField}^2$ (the rank-one case). With $\sym{E0}(x)=x\sym{v}$, the residual
operator satisfies $\sym{R}(x)=x\sym{r}$, where
$\sym{r}=\sym{Acurrent}\sym{v}-\sym{rho}\sym{v}$.
For a higher-dimensional trial space,
$\sym{R}=\sym{Acurrent}\sym{E0}-\sym{E0}\sym{TrialOperator}$ assigns a residual vector
to each trial-coordinate vector. In a finite-dimensional eigenbasis of
$\sym{TrialOperator}$, its columns are the individual eigenpair residuals.
The exact frames in the drawing are
$\sym{F0}(h)=(h,0)$ and $\sym{F1}(g)=(0,g)$, with
$\sym{Pexact}=\sym{F0}\like{F0}{\symglyph{F0}^*}$ and $\sym{PexactPerp}=\sym{Identity}-\sym{Pexact}$.
(a) The pink leg $\sym{S}(1)=\sym{PexactPerp}\sym{v}$ has length $\sin\sym{theta}$.
(b) The green residual joins $\sym{rho}\sym{v}$ to $\sym{Acurrent}\sym{v}$;
for the Rayleigh choice it is perpendicular to the trial line. The dashed
legs are its exact wanted and unwanted components. Their lengths give
$\sym{delta}\sin\sym{theta}=\|\sym{PexactPerp}\sym{r}\|_2
\leq\dkOperatorNorm{\sym{r}\|_2=\|\sym{R}}$.
The drawing uses $\sym{Acurrent}=\sym{DiagonalOperator}(1.9,0.7)$,
$\sym{theta}=40^\circ$, $\sym{rho}=\langle\sym{v},\sym{Acurrent}\sym{v}\rangle$,
and $\sym{delta}=\sym{rho}-0.7$}
\label{fig:sine-theta-residual}
\end{figure}
% Place the residual figure before the theorem and its Lean signature.
\FloatBarrier

Appendix~\ref{app:lean-guide} explains the general Lean syntax in the
following theorem signature; the geometric and spectral hypotheses have been
defined above.

\begin{theorem}[Sine-theta estimate]
\label{thm:sine-theta}
Let the scalar field be $\sym{RealField}$ or $\sym{ComplexField}$, let $\sym{AmbientHilbertSpace}$ be separable,
and let $\sym{Acurrent}$, $\sym{TrialOperator}$, and $\sym{Lambda1}$ be self-adjoint operators on the
Hilbert spaces above. Suppose the isometries, domain inclusions, residual
identity, exact complementary decomposition, and the separation condition
$\sym{SineGap}(\sym{TrialOperator},\sym{Lambda1})$ of
Section~\ref{sec:gap-geometry} hold. Let $\sym{NormN}$ be a normalized symmetric
operator-ideal norm. Whenever $\sym{sinTheta0}$ and $\sym{R}$ belong to its domain,
\begin{equation}
 \sym{delta}\,\sym{NormN}(\sym{sinTheta0})\leq \sym{NormN}(\sym{R}).
 \label{eq:dk-sine}
\end{equation}
\end{theorem}

% Keep the Lean signature and its caption together across page breaks.
\noindent\begin{minipage}{\linewidth}
\begin{lstlisting}[style=leanpaper,
  caption={Lean signature for the sine-theta estimate. The surrounding
  scalar-field and Hilbert-space context is omitted.
  This is a readable ASCII transcription of the Lean declaration,
  rather than a verbatim copy of its Unicode binders and operator arrows.
  The paper-to-Lean names are registered centrally and displayed in
  Table~\ref{tab:paper-lean-notation}. In particular, the Lean binder
  displayed here as \texttt{A0} names our $\sym{TrialOperator}$, not the
  distinct reference block $\sym{ReferenceBlock0}$.
  The positive trial-coordinate sine is
  \texttt{sourceDirectedSinThetaOperator E0 F0};
  \texttt{FormBoundedSylvesterGap} encodes the interval/exterior and ordered
  half-infinite separation hypotheses described above.},
  label={lst:dk-sintheta-headline}]
theorem sinTheta_unbounded_formGap_whereDefinedUIN_rclike
    [TopologicalSpace.SeparableSpace E]
    (N : NormalizedSymmetricOperatorIdealFamily K)
    (A : LinearPMap K E E)
    (A0 : LinearPMap K F F)
    (Lambda1 : LinearPMap K G G)
    (E0 : ContinuousLinearMap K F E)
    (F0 : ContinuousLinearMap K H E)
    (F1 : ContinuousLinearMap K G E)
    (R : ContinuousLinearMap K F E)
    (hA : IsSelfAdjoint A)
    (hA0 : IsSelfAdjoint A0)
    (hLambda1 : IsSelfAdjoint Lambda1)
    (htrial : IsTrialResidual A A0 E0 R)
    (hexact : IsExactSpectralDecomposition A Lambda1 F0 F1)
    (delta : Real) (hdelta : 0 < delta)
    (hgap : FormBoundedSylvesterGap A0 Lambda1 delta) :
    N.Mem (sourceDirectedSinThetaOperator E0 F0) ->
    N.Mem R ->
      delta * N.gaugeReal (sourceDirectedSinThetaOperator E0 F0) <=
        N.gaugeReal R
\end{lstlisting}
\end{minipage}\par

The two ideal-membership hypotheses ensure that the norm values in the
conclusion are finite. The gap predicate includes the finite
interval/exterior case and both ordered cases. In the ordered cases,
$\sym{TrialOperator}$ and $\sym{Lambda1}$ may both be unbounded.

\subsubsection{Why the Sylvester estimate gives the angle bound}

Set $\sym{X}=\like{F1}{\symglyph{F1}^*}\sym{E0}$. The exact complementary decomposition gives
$\sym{S}=\sym{F1}\sym{X}$, and multiplication by the isometry $\sym{F1}$ preserves the
singular-value gauges. Self-adjointness, the domain hypotheses, and the
residual identity yield the Sylvester relation
\[
 \sym{Lambda1} \sym{X}-\sym{X}\sym{TrialOperator}=\like{F1}{\symglyph{F1}^*}\sym{R}
\]
on $\DKDom(\sym{TrialOperator})$, with the left side represented by the bounded
right side. The separated Sylvester estimate gives
\[
 \sym{delta} \sym{NormN}(\sym{X})\leq \sym{NormN}(\like{F1}{\symglyph{F1}^*}\sym{R})\leq \sym{NormN}(\sym{R}).
\]
Combining this with $\sym{NormN}(\sym{X})=\sym{NormN}(\sym{S})=\sym{NormN}(\sym{sinTheta0})$ proves the conclusion.
For unbounded operators, deriving the domain compatibility for $\sym{X}$ is part of
the argument. The displayed relation holds on $\DKDom(\sym{TrialOperator})$ and
has a bounded right side. The specified interval/exterior or ordered
spectral positions give the factor $1/\sym{delta}$, with no additional
multiplicative constant.

For a one-dimensional trial space, take
$\sym{Acurrent}=\sym{DiagonalOperator}(0,\sym{delta})$, the trial scalar
$\sym{TrialOperator}=0$, and the unit vector
$\like{v}{\symglyph{v}_{\symglyph{theta}}}=(\cos\sym{theta},\sin\sym{theta})$.
Choose the desired exact subspace $\sym{SpanOperator}\{(1,0)\}$;
its orthogonal complement is $\sym{SpanOperator}\{(0,1)\}$, on which
$\sym{Lambda1}=\sym{delta}$.
The complementary component is $(0,\sin\sym{theta})$ and the
residual is $(0,\sym{delta}\sin\sym{theta})$. Thus equality holds
in the sine-theta estimate for this finite-dimensional configuration.

\section{Formalized operator-theoretic foundations}
\label{sec:foundations}

We formalized the geometry of two subspaces, symmetric operator ideals and
majorization, and Sylvester estimates for unbounded self-adjoint operators.
These results supply the constructions and inequalities used in the
Davis--Kahan proofs.

\subsection{Direct rotations and two-subspace geometry}

Let $\sym{Psource}$ and $\sym{Qtarget}$ project onto closed subspaces
$\sym{Usource}$ and $\sym{Vtarget}$. For an acute pair,
$\|\sym{Psource}-\sym{Qtarget}\|<1$. The compression
$\sym{Psource}\sym{Qtarget}\sym{Psource}$ on $\sym{Usource}$ determines
the positive cosine operator. Its square root and the complementary
compression give the cosine and sine factors in the direct-rotation
construction. Polar decomposition gives the isometries between the
corresponding parts of the two subspaces.

The resulting orthogonal or unitary map $\sym{Rrotation}$ carries
$\sym{Usource}$ onto $\sym{Vtarget}$. In finite-dimensional principal-plane
coordinates it rotates each plane through its principal angle and fixes
common directions. The projection formula from
\cref{eq:history-kato-unitary} defines it in Hilbert space without requiring
a discrete list of principal vectors.

Our formalization proves acute existence and uniqueness, the
principal-square-root characterization, commutation with the angle
operators, and reversal symmetry: interchanging the two subspaces replaces
the rotation by its adjoint. For general pairs, the existence criteria
include the right-angle defect spaces described in
\cref{sec:related-two-projections}. When the angle operator is compact,
its spectral decomposition recovers a sequence of principal angles
\citep[Section~3]{davis1970rotation}.

The double-angle theorems for $\sin(2\sym{Theta})$ and $\tan(2\sym{Theta})$
  require selecting the branch on which $\sym{Theta}\le\sym{PiConstant}/4$.
Here $\sym{Theta}$ includes both the source and complementary angle blocks;
$\sym{Theta0}$ alone is the directed angle on source coordinates.
Because the function $\sin(2\sym{theta})$ is small both when $\sym{theta}$ is
  near $0$ and when $\sym{theta}$ is near $\sym{PiConstant}/2$, a bound on
  $\sin(2\sym{Theta})$ alone cannot determine whether an approximate subspace is
  close to or nearly orthogonal to the target.
The branch-selection theorem \citep[Theorem~8.1]{davis1970rotation}
uses an ordered reference gap. Let $\sym{Psource}$ reduce a self-adjoint
reference operator $\sym{A}$, with
blocks $\sym{ReferenceBlock0}$ on $\DKRange(\sym{Psource})$ and
$\sym{ReferenceBlock1}$ on $\DKKernel(\sym{Psource})$, satisfying
\[
 \sym{ReferenceBlock0}\le\sym{alphaEndpoint}\sym{Identity},
 \qquad
 \sym{ReferenceBlock1}\ge(\sym{alphaEndpoint}+\sym{delta})\sym{Identity},
 \qquad \sym{delta}>0,
\]
in the quadratic-form sense. Suppose $\sym{H}$ is bounded and self-adjoint,
with both diagonal blocks zero:
$\sym{Psource}\sym{H}\sym{Psource}
 =(\sym{Identity}-\sym{Psource})\sym{H}(\sym{Identity}-\sym{Psource})=0$.
These hypotheses imply that $\sym{Acurrent}=\sym{A}+\sym{H}$ has no spectrum
in $(\sym{alphaEndpoint},\sym{alphaEndpoint}+\sym{delta})$.

Fix a reducing subspace $M$ of $\sym{Acurrent}$, and write
$\sym{Lambda0},\sym{Lambda1}$ for the blocks of $\sym{Acurrent}$ on $M,M^\perp$.
The full angle $\sym{Theta}$ comparing $\DKRange(\sym{Psource})$ with $M$
satisfies $\sym{Theta}\le\sym{PiConstant}/4$ if and only if
$\sym{Lambda0}\le\sym{alphaEndpoint}\sym{Identity}$ and
$\sym{Lambda1}\ge(\sym{alphaEndpoint}+\sym{delta})\sym{Identity}$.
Such a subspace exists: take the spectral subspace of $\sym{Acurrent}$
for $(-\infty,\sym{alphaEndpoint}]$. Our formalization also proves that
it is the unique reducing subspace satisfying this angle bound.

To compare the reference and perturbed blocks on common coordinate spaces,
choose isometric frames for $\DKRange(\sym{Psource})$ and
$\DKKernel(\sym{Psource})$. Apply the direct rotation $\sym{Rrotation}$ to
these frames to obtain frames for $M$ and $M^\perp$. Express
$\sym{ReferenceBlock0},\sym{ReferenceBlock1}$ in the original frames and
$\sym{Lambda0},\sym{Lambda1}$ in the transported frames. In these coordinates
the branch satisfies the stronger quadratic-form inequality
\[
 \sym{ReferenceBlock1}-\sym{alphaEndpoint}\sym{Identity}
 \le\sym{ComplementCosineBlock}
       (\sym{Lambda1}-\sym{alphaEndpoint}\sym{Identity})\sym{ComplementCosineBlock},
 \qquad \sym{ComplementCosineBlock}=\cos\sym{Theta1},
\]
with the lower-block analogue. These inequalities quantify spectral
repulsion in terms of the angles and yield eigenvalue and symmetric-gauge
comparisons in finite dimensions.

The smallness criterion \citep[Theorem~8.2]{davis1970rotation} uses a central
spectral interval instead of an ordered gap. Let $\sym{Psource}$ reduce
the self-adjoint operator $\sym{A}$ and $\sym{Qtarget}$ reduce the
self-adjoint operator $\sym{Acurrent}$. The perturbed block on
$\DKRange(\sym{Qtarget})$ has spectrum in
$[\sym{betaEndpoint},\sym{alphaEndpoint}]$, while its complementary spectrum
lies in
\[
 (-\infty,\sym{betaEndpoint}-\sym{delta}]
 \ \cup\ [\sym{alphaEndpoint}+\sym{delta},\infty),
 \qquad \sym{betaEndpoint}\le\sym{alphaEndpoint},\quad\sym{delta}>0.
\]
The reference block $\sym{ReferenceBlock0}$, representing
$\sym{A}|_{\DKRange(\sym{Psource})}$ in source coordinates, must have spectrum
in $[\sym{betaEndpoint}-\sym{delta}/2,\sym{alphaEndpoint}+\sym{delta}/2]$.
For an isometric frame $\sym{E0}$ of the reference subspace, its residual
is $\sym{R}=\sym{Acurrent}\sym{E0}-\sym{E0}\sym{ReferenceBlock0}$
on the operator domain, as in \cref{eq:history-residual}.
Assume
$\sym{DimensionOperator}\sym{CrossedSourceDefect}=\sym{DimensionOperator}\sym{CrossedTargetDefect}$: the space of source
vectors orthogonal to the target and the space of target vectors orthogonal
to the source have equal Hilbert dimension.

If $\dkOperatorNorm{\sym{H}}<\sym{delta}/2$ for a
bounded self-adjoint perturbation, or $\dkOperatorNorm{\sym{R}}<\sym{delta}/2$
for the bounded residual on the reference subspace, then
$\sym{Theta}<\sym{PiConstant}/4$. For each ideal norm whose displayed arguments belong
to its domain, the corresponding double-angle estimate also holds:
\[
 \sym{delta}\sym{NormN}(\sin(2\sym{Theta}))\le2\sym{NormN}(\sym{H}),
 \qquad\text{or}\qquad
 \sym{delta}\sym{NormN}(\sin(2\sym{Theta0}))\le2\sym{NormN}(\sym{R}).
\]
% Note (not for print): the factor 2 in the residual form is correct and
% matches the 1970 Section 2 statement. The 1970 Section 7 proof, after (7.5),
% rewrites delta||sin 2Theta|| <= ||H - XHX|| with right side ||[[0,B*],[B,0]]||
% and concludes delta||sin 2Theta_0|| <= ||B|| <= ||R||; since H - XHX has
% off-diagonal block 2B, the right side should be 2||B|| <= 2||R||. The
% Section 7 example (A = diag(0,d), H = [[0,1],[1,-d]], delta = 2,
% theta = pi/4, ||R|| = 1) attains delta sin 2theta = 2 = 2||R||. Proof-line
% slip only; the stated theorem is unaffected.
The perturbation estimate uses the full angle, and the residual estimate
the directed angle on source coordinates. In the unbounded residual form,
$\sym{A}$ and $\sym{Acurrent}$ are self-adjoint on a common domain; only
the residual is required to be bounded.

We formalized both theorems in Hilbert space using spectral calculus.

\subsection{Approximation numbers and operator-ideal norms}

In finite dimensions, a unitarily invariant norm is a symmetric gauge
function of the singular values. Approximation numbers extend the
singular-value comparisons needed for bounded operators; for compact
operators they give the usual singular-value sequence. We formalized
symmetric norming functions and Ky Fan partial sums to pass from
majorization to norm inequalities.

For example, let $\sym{Rrotation}$ be the direct rotation and
$\sym{Wcompetitor}$ any unitary carrying $\sym{Usource}$ onto
$\sym{Vtarget}$. The restricted displacement of the direct rotation has
singular-value minimality, which implies minimality for unitarily invariant
norms. The squared-displacement result compares
$(\sym{Identity}-\sym{Rrotation})^*(\sym{Identity}-\sym{Rrotation})$ and
$(\sym{Identity}-\sym{Wcompetitor})^*(\sym{Identity}-\sym{Wcompetitor})$
\citep[Proposition~4.3]{davis1970rotation}. Their singular values are the
squares of those of the unsquared displacements. Section~\ref{sec:prop44}
shows why the claimed trace-norm comparison fails and derives the valid
Q-norm conclusion.
% Additional formalized quantitative results, omitted from this article:
% DavisKahan/FiniteDimensional/DirectRotation/SharpDisplacement.lean and
% SharpDisplacementLimit.lean. AUTHORING_NOTES.md indexes the evidence.

The same norm inequalities enter the sine-theta proof. The assertion
$\sym{delta}\sym{NormN}(\sym{S})\leq\sym{NormN}(\sym{R})$ applies when
$\sym{S}$ and $\sym{R}$ have finite ideal norm. Our ideal-membership
hypotheses state this condition explicitly in arbitrary Hilbert dimension.

\subsection{Sylvester estimates and unbounded operators}

For operators $\sym{SylvesterLeft}$ and $\sym{SylvesterRight}$ on the target and source spaces of $\sym{X}$, the
Sylvester map is $\sym{X}\mapsto \sym{SylvesterLeft}\sym{X}-\sym{X}\sym{SylvesterRight}$. Under the interval/exterior or ordered
spectral-separation hypotheses, the estimate has the form
\[
 \sym{delta} \sym{NormN}(\sym{X})\leq \sym{NormN}(\sym{SylvesterLeft}\sym{X}-\sym{X}\sym{SylvesterRight}).
\]
The sine-theta proof uses $\sym{SylvesterLeft}=\sym{Lambda1}$, $\sym{SylvesterRight}=\sym{TrialOperator}$, and
$\sym{X}=\like{F1}{\symglyph{F1}^*}\sym{E0}$. Other perturbation estimates choose related
projection blocks for $\sym{X}$.

For unbounded $\sym{SylvesterLeft}$ and $\sym{SylvesterRight}$, the identity is asserted for
$x\in\DKDom(\sym{SylvesterRight})$, with $\sym{X}x\in\DKDom(\sym{SylvesterLeft})$, and $\sym{SylvesterLeft}\sym{X}-\sym{X}\sym{SylvesterRight}$ agrees there with a
bounded residual. We formalized these domain conditions using linear
partial maps. Spectral cutoffs give bounded approximations; convergence
of approximation numbers passes the norm estimates to the unbounded
operators \citep[Section~5 and Appendix to Section~6]{davis1970rotation}.

The generalized sine and tangent theorems also compare block operators.
Reflection pinching replaces a block matrix $T$ by
$(T+\sym{Jreflect}T\sym{Jreflect})/2$, where $\sym{Jreflect}$ acts as $+\sym{Identity}$ on one coordinate subspace and $-\sym{Identity}$
on the other. This removes the off-diagonal blocks and contracts a unitarily
invariant norm. Direct-sum comparisons relate the norms of the component
blocks to the whole operator. Together with Sylvester separation these give
the generalized angle estimates
\citep[Lemmas~6.1--6.2 and Theorems~6.1--6.3]{davis1970rotation}.

\subsection{Real and complex Hilbert spaces}

The spectral constructions use complex functional calculus. We formalized
the scalar-transfer results needed to carry these arguments to real Hilbert
spaces through complexification. This permits the source theorems to be
stated over either field, including their domain and ideal-norm hypotheses.
The underlying algebra, topology, and continuous functional calculus come
from Mathlib; the adapted Spectra constructions and later Tau Ceti
contributions are described in Section~\ref{sec:related-formal}.

\input{counterexample.tex}

\input{formalization_enabled_extension.tex}

\section{Discussion}
\label{sec:discussion}

The formalization covers the original perturbation theory on real and
complex Hilbert spaces, including unbounded self-adjoint operators and
operator-ideal norms. The unbounded statements require explicit domain
conditions and bounded residuals. The operator-ideal statements require
finite norm values. These hypotheses determine which approximation problems
the estimates address.

The correction of the full-displacement claim
\citep[Proposition~4.4]{davis1970rotation} separates two extremal statements.
The valid comparison \citep[Proposition~4.3]{davis1970rotation}
gives squared-displacement minimality; its Q-norm consequence
is valid. Our real four-dimensional family exhibits full trace-norm failure
at every prescribed positive acute angle. In each example, the competitor has
two nonzero displacement singular values while the direct rotation has four.
The sums and squared sums of these values account for their different norm
comparisons. The earlier complex example supplied the construction; our
calculation makes its contradiction of the printed real claim explicit.

The matching theorem addresses the resulting optimization problem for any
equal-dimensional pair of finite-dimensional real subspaces. Its proof uses
principal-angle geometry, real orthogonal block decompositions, and
singular-value norms. These are also among the foundations required for the
Davis--Kahan formalization. The proof of the matching theorem uses additional
optimization arguments, including a passage to the trace-norm endpoint.

\section{Conclusion}
\label{sec:conclusion}

We formalized 28 of the 29 selected Davis--Kahan results and refuted the
remaining printed statement \citep[Proposition~4.4]{davis1970rotation}.
The sine-theta theorem includes
unbounded self-adjoint operators and symmetric operator-ideal norms with
their domain and finiteness assumptions. The direct-rotation calculation
identifies the failed singular-value comparison and gives an explicit real
counterexample and a valid Q-norm consequence.

We also formalized the operator theory used in these results: principal
angles and direct rotations, Borel spectral calculus, approximation numbers
and symmetric ideals, and Sylvester separation. Parts of this development
were subsequently adapted into Tau Ceti. Using these foundations, we proved
an exact weighted-matching formula for the real trace-norm transport
minimum.

\backmatter
\section*{Statements and Declarations}

\subsection*{Funding}
\ifDKAnonymous
Funding information is omitted for anonymous review.
\else
This work was supported by the Defense Advanced Research Projects Agency
under Contract No. HR001125CE017.
\fi

\subsection*{Competing interests}
The authors declare no competing interests.

\subsection*{Data and code availability}
\ifDKAnonymous
Artifact information is omitted for anonymous review.
\else
The Lean proofs, result-to-Lean mapping, provenance records, figure-generation
sources, and manuscript sources are available at
\url{https://github.com/AIQ-Kitware/aiq-dkps-formalization}.
\fi

\subsection*{AI assistance}
ChatGPT and Claude assisted with the formalization and manuscript preparation.
Formalizations were carried out by large language models with human oversight.
The matching result was derived using GPT-6.1 Sol and reviewed by the human
  authors.

\bibliography{references}

\clearpage
\appendix

\input{appendices.tex}

\end{document}
, without editing.
% EVERY \sym, \like, and \dkrole obeys this switch, including the glossary.
% No color is required to identify a mathematical object: the glyphs, indices,
% definitions and adjacent prose also distinguish them.
%
% Families (base hue): exact/target=blue, trial/source=amber, exact unwanted=
% cyan, angles and their operators=magenta, residual/displacement=green,
% gaps=violet. Norm gauges=slate; provenance retains brown/steel/vermilion.
% P_ex is the exact-target projection in sine-theta; P is the source
% projection in transport. R denotes a residual and D a direct rotation.
% Semantic keys retain their roles even when the printed notation changes.
\newif\ifDKMathColors
\DKMathColorstrue
\ifdefined\DKMonochromeMath
  \DKMathColorsfalse
\fi

% Family anchors (print-compatible versions of the Manim light palette).
\definecolor{dkfg}{HTML}{1B1F24}
\definecolor{dkmuted}{HTML}{6B7280}
\definecolor{dkwanted}{HTML}{1F6FD1}
\definecolor{dktrial}{HTML}{A66705}
\definecolor{dkunwanted}{HTML}{007C91}
\definecolor{dksine}{HTML}{D12A7B}
\definecolor{dkresid}{HTML}{087D43}
\definecolor{dkgap}{HTML}{6E4BD8}
\definecolor{dkold}{HTML}{705519}
\definecolor{dkcurrent}{HTML}{054669}
\definecolor{dkperturb}{HTML}{BB491B}
\definecolor{dkrefuted}{HTML}{C62828}

% Related yet distinguishable shades. These are NOT new unrelated categories:
% objects inherit their parent family; a different shade signals a different
% mathematical object (e.g., a frame versus its subspace or its projection).
\definecolor{dkwantedframe}{HTML}{1254A9}
\definecolor{dkwantedprojection}{HTML}{08618D}
\definecolor{dkwantedrotation}{HTML}{4248B0}
\definecolor{dkwantedspectrum}{HTML}{345F8F}
\definecolor{dkwantedvector}{HTML}{1468AA}
\definecolor{dktrialframe}{HTML}{905600}
\definecolor{dktrialprojection}{HTML}{79511B}
\definecolor{dktrialcompetitor}{HTML}{A04D10}
\definecolor{dktrialspectrum}{HTML}{79591C}
\definecolor{dktrialvector}{HTML}{9B522A}
\definecolor{dkunwantedframe}{HTML}{006477}
\definecolor{dkunwantedspectrum}{HTML}{00546B}
\definecolor{dkunwantedprojection}{HTML}{145B6D}
\definecolor{dkunwantedvector}{HTML}{256879}
\definecolor{dkangleoperator}{HTML}{923478}
\definecolor{dksineoperator}{HTML}{AA275A}
\definecolor{dksineblock}{HTML}{B43F70}
\definecolor{dksylvester}{HTML}{804074}
\definecolor{dkanglemax}{HTML}{B32185}
\definecolor{dkanglecos}{HTML}{824276}
\definecolor{dkresidvector}{HTML}{086038}
\definecolor{dkdisplacement}{HTML}{226C50}
\definecolor{dkgaporiginal}{HTML}{4E4D89}
\definecolor{dkgapstat}{HTML}{7039A7}
\definecolor{dkgapinterval}{HTML}{574A82}
\definecolor{dkambient}{HTML}{36516A}
\definecolor{dknorm}{HTML}{475875}
\definecolor{dknormaux}{HTML}{66526E}
\definecolor{dkintertwiner}{HTML}{4A6280}
\definecolor{dkturn}{HTML}{47647D}
\definecolor{dkmix}{HTML}{505B65}

% Coloring primitives. "plain" suppresses color without suppressing hovers.
% All descendants can retain a semantic role through \like{key}{expression}.
\expandafter\def\csname dknocolor@plain\endcsname{1}
\newcommand{\dkrole}[2]{%
  \ifDKMathColors
    \ifcsname dknocolor@#1\endcsname #2\else\textcolor{dk#1}{#2}\fi
  \else #2\fi
}
% An optional plain-text reminder supplies definitions and types on hover.
% Required fields: {semantic key}{printed glyph}{color role}{glossary|local}
%                  {short definition}{mathematical type}. The "glossary" flag records whether an
% entry appears in the curated printed notation table; "local" does not.
% This registry is the sole authority for symbol meaning and printed form.
% Use role "plain" for an uncolored symbol that still carries hover text.
% Changing the glyph is a separate editorial decision: this audit preserves
% the current E, F, F_ex, G, and other existing printed choices.
\newcommand{\dkDeclareSymbol}[7][]{%
  \expandafter\def\csname dkgloss@#2\endcsname{#6}%
  \expandafter\def\csname dkglossaryflag@#2\endcsname{#5}%
  \if\relax\detokenize{#1}\relax
    \expandafter\pdfstringdef\csname dkpdfgloss@#2\endcsname{#7: #6}%
  \else
    \expandafter\pdfstringdef\csname dkpdfgloss@#2\endcsname{#7: \detokenize{#1}}%
  \fi
  \expandafter\def\csname dktype@#2\endcsname{#7}%
  \expandafter\def\csname dkglyph@#2\endcsname{#3}%
  \expandafter\def\csname dkrole@#2\endcsname{#4}%
  \expandafter\def\csname dksym@#2\endcsname{\dkrole{#4}{#3}}%
}
% Unstyled printed glyph, for variants constructed through \like.
% Example: \like{E0}{\symglyph{E0}^*} derives E_0^* from E0's
% one canonical glyph; retyping E_0 here would defeat central control.
\DeclareRobustCommand{\symglyph}[1]{\csname dkglyph@#1\endcsname}
\DeclareRobustCommand{\sym}[1]{%
  \ifcsname dksym@#1\endcsname
    \dkSymbolTooltip{#1}{\csname dksym@#1\endcsname}%
  \else
    \PackageError{dk-math-colors}{Unknown symbol key #1}{See math_symbol_colors.tex}%
  \fi
}
\DeclareRobustCommand{\like}[2]{%
  \ifcsname dkrole@#1\endcsname
    \dkSymbolTooltip{#1}{\dkrole{\csname dkrole@#1\endcsname}{#2}}%
  \else
    \PackageError{dk-math-colors}{Unknown role key #1}{See math_symbol_colors.tex}%
  \fi
}

% A short description belongs to the semantic key, not the bare letter.
% The same descriptions are available to the notation appendix.
\DeclareRobustCommand{\dkSymbolDescription}[1]{\csname dkgloss@#1\endcsname}
\DeclareRobustCommand{\dkSymbolGlossaryStatus}[1]{\csname dkglossaryflag@#1\endcsname}
\DeclareRobustCommand{\dkSymbolType}[1]{\csname dktype@#1\endcsname}

% Citable alternative notation is metadata, not a request to rename glyphs.
% The source identifier, the source's exact printed glyph, and the specific
% source location are all mandatory. No alternative is inferred from a name.
% One key may have alternatives from several works or multiple passages.
\newcommand{\dkDeclareNotationSource}[3]{%
  \expandafter\def\csname dksourcetitle@#1\endcsname{#2}%
  \expandafter\def\csname dksourceurl@#1\endcsname{#3}%
}
\newcommand{\dkDavisKahanRows}{}
\makeatletter
\newcommand{\dkDeclareNotationAlternative}[5]{%
  \ifcsname dkglyph@#1\endcsname\else
    \PackageError{dk-math-colors}{Unknown alternative key #1}{Declare the symbol first}%
  \fi
  \ifcsname dksourcetitle@#2\endcsname\else
    \PackageError{dk-math-colors}{Unknown notation source #2}{Declare the source first}%
  \fi
  \expandafter\def\csname dkaltglyph@#2@#1\endcsname{#3}%
  \expandafter\def\csname dkaltlocator@#2@#1\endcsname{#4}%
  \expandafter\def\csname dkaltnote@#2@#1\endcsname{#5}%
  \def\dk@source{#2}%
  \def\dk@original{DK1970}%
  \ifx\dk@source\dk@original
    \g@addto@macro\dkDavisKahanRows{\dkDavisKahanRow{#1}}%
  \fi
}
\makeatother
% A correspondence row is produced from the SAME registry records used by
% \sym and by the source-census script. It is never spelled a second time.
\newcommand{\dkDavisKahanRow}[1]{%
  \csname dkaltnote@DK1970@#1\endcsname &
  $\csname dkaltglyph@DK1970@#1\endcsname$ &
  $\sym{#1}$ &
  \csname dkaltlocator@DK1970@#1\endcsname\\
}

\newif\ifDKSymbolTooltips
\DKSymbolTooltipstrue
\ifdefined\DKNoSymbolTooltips\DKSymbolTooltipsfalse\fi
\newsavebox{\DKTooltipBox}
\newcount\DKTooltipCount
% Register only shipped widgets through hyperref's existing form-field
% machinery. Poppler/Okular need the catalog's field list, not just page
% annotations. The deferred aux write also discards measuring-only copies.
\makeatletter
\newcommand{\dkRegisterTooltipField}{%
  \expandafter\HyField@@AddToFields\expandafter{\number\pdffeedback lastannot}%
}
\makeatother
\AtBeginDocument{\ifDKSymbolTooltips\Form[NeedAppearances=false]\fi}
\AtEndDocument{\ifDKSymbolTooltips\endForm\fi}
% Native LuaLaTeX PDF widgets: no navigation action or JavaScript. Avoid
% pdfcomment's soulpos dependency, which previously broke local Vim builds.
% Measure in the CURRENT math style, then emit the original expression as-is:
% boxing the visible math would change spacing, scripts and line breaking.
\newcommand{\dkSymbolTooltip}[2]{%
  \ifDKSymbolTooltips
    \begingroup
      \DKSymbolTooltipsfalse % suppress nested annotations while measuring
      \ifmmode
        \edef\DKTooltipStyle{\ifcase\mathstyle
          \displaystyle\or\displaystyle\or\textstyle\or\textstyle\or
          \scriptstyle\or\scriptstyle\or\scriptscriptstyle\or\scriptscriptstyle\fi}%
        \sbox{\DKTooltipBox}{$\DKTooltipStyle#2$}%
      \else
        \leavevmode\sbox{\DKTooltipBox}{#2}%
      \fi
      \global\advance\DKTooltipCount by 1
      \pdfextension annot width\wd\DKTooltipBox
        height\ht\DKTooltipBox depth\dp\DKTooltipBox
        {/Subtype /Widget /FT /Btn /Ff 65536 /F 512
         /T (dk-symbol-\the\DKTooltipCount) /TU (\csname dkpdfgloss@#1\endcsname)
         /H /N /BS << /W 0 >> /Border [0 0 0]}%
      \dkRegisterTooltipField
      #2%
    \endgroup
  \else
    #2%
  \fi
}
\pdfstringdefDisableCommands{%
  \def\sym#1{\csname dkglyph@#1\endcsname}%
  \def\like#1#2{#2}%
  \def\symglyph#1{\csname dkglyph@#1\endcsname}%
}

% Declarations: [detailed hover], key, glyph, role, glossary|local, short text, type.
% A local entry has full semantic/hover support, but no notation-table row.
% When two keys have the same glyph, it is an auditable collision, not an
% instruction to change the printed notation automatically.
% This registry includes local proof notation. The notation appendix selects
% the symbols useful for reference; registration does not add a printed row.
% Manim's original symbol keys, now split into semantic subshades.
% Ambient and perturbed operators (provenance is independent of geometry).
\dkDeclareSymbol[Reference self-adjoint operator A on the ambient Hilbert space; the perturbed operator is Ahat = A + H. In a single-operator eigenproblem, A itself is the operator being approximated. An unbounded A acts only on domain(A).]
  {A}{A}{old}{glossary}{Reference self-adjoint operator before perturbation.}{self-adjoint-operator}
\dkDeclareSymbol{A1}{A_1}{old}{local}{Complementary block of the reference operator.}{self-adjoint-operator}
\dkDeclareSymbol{a1}{a_1}{old}{local}{Reference eigenvalue or complementary spectral value.}{scalar}
\dkDeclareSymbol[Reference block A_0: the restriction of A to the range of the reducing source projection P, expressed in source coordinates. In Theorem 8.1 it lies at or below alpha; in Theorem 8.2 its spectrum lies in the enlarged central interval. The perturbed blocks are denoted Lambda_0 and Lambda_1.]
  {ReferenceBlock0}{A_0}{old}{glossary}{Reference operator restricted to the source subspace in Section 8.}{self-adjoint-operator}
\dkDeclareSymbol[Reference block A_1: the restriction of A to the kernel of the reducing source projection P, expressed in complementary coordinates. In Theorem 8.1 it lies at or above alpha + delta in quadratic-form order.]
  {ReferenceBlock1}{A_1}{old}{glossary}{Reference operator restricted to the source complement in Section 8.}{self-adjoint-operator}
\dkDeclareSymbol[Self-adjoint ambient operator Ahat being studied; in the perturbation setting Ahat = A + H. Its action on a trial vector is compared with the proposed eigenvalue times that vector.]
  {Ahat}{\widehat A}{current}{local}{Operator under study; in the perturbation setting, the reference operator plus its perturbation.}{self-adjoint-operator}
\dkDeclareSymbol[Self-adjoint ambient operator Ahat being studied. In the perturbation setting Ahat = A + H. The trial residual is R = Ahat E_0 - E_0 A_tr, and the exact complementary block satisfies Ahat F_1 = F_1 Lambda_1 on its domain.]
  {Acurrent}{\widehat A}{current}{glossary}{Operator under study; in the perturbation setting, the reference operator plus its perturbation.}{self-adjoint-operator}
\dkDeclareSymbol[Additive self-adjoint perturbation H = Ahat - A. A bound on its norm can control the residual of an eigenspace of one operator relative to the other.]
  {H}{H}{perturb}{glossary}{Additive self-adjoint perturbation of the reference operator.}{self-adjoint-operator}
\dkDeclareSymbol{H0}{H_0}{perturb}{local}{Wanted diagonal block of the perturbation.}{self-adjoint-operator}
\dkDeclareSymbol{H1}{H_1}{perturb}{local}{Unwanted diagonal block of the perturbation.}{self-adjoint-operator}
\dkDeclareSymbol{B}{B}{perturb}{local}{Off-diagonal perturbation or coupling block.}{linear-operator}
\dkDeclareSymbol{b}{b}{perturb}{local}{Scalar perturbation or coupling coefficient.}{scalar}
% Trial/source family: subspace, frame, compressed operator, projector, vector.
\dkDeclareSymbol{U}{U}{wanted}{local}{Desired exact subspace in the local geometric notation.}{linear-subspace}% Manim's U remains wanted in local slide notation.
\dkDeclareSymbol{V}{V}{trial}{local}{Trial subspace in the local geometric notation.}{linear-subspace}% Manim's V remains trial in local slide notation.
\dkDeclareSymbol[Source subspace U to be carried onto the target V by an orthogonal or unitary map W. In the transport sections P projects onto U and Q onto V. The matching theorem also names the source subspace K.]
  {Usource}{U}{trial}{glossary}{Source subspace to be transported onto the target subspace.}{linear-subspace}
\dkDeclareSymbol[Trial frame E_0: F -> E, a linear isometry from trial-coordinate space F into ambient space E. Its range is the trial subspace; E_0* returns trial coordinates and E_0* E_0 = I_F. In rank one, E_0(x) = x v.]
  {E0}{E_0}{trialframe}{glossary}{Isometric embedding of the trial subspace; its adjoint returns trial coordinates.}{isometric-embedding}
\dkDeclareSymbol[Self-adjoint trial operator A_tr on trial-coordinate space F. The residual compares Ahat E_0 with E_0 A_tr. Rayleigh--Ritz chooses A_tr = E_0* Ahat E_0 when this compression is defined; the theorem allows other trial operators.]
  {A0}{A_{\mathrm{tr}}}{trialspectrum}{glossary}{Self-adjoint operator on trial coordinates; its eigenvalues are the trial spectral values.}{self-adjoint-operator}
\dkDeclareSymbol[Trial vector: a proposed approximate eigenvector. In the rank-one geometry v has unit length and E_0(x) = x v. Its residual for trial scalar rho is r = Ahat v - rho v.]
  {v}{v}{trialvector}{local}{Trial vector or vector in the trial subspace.}{vector}
\dkDeclareSymbol[Scalar trial operator in rank one, A_tr(x) = rho x. Figure 1 uses the Rayleigh value rho = <v, Ahat v> for unit v; then rho v is the projection of Ahat v onto the trial line and r is perpendicular to v.]
  {rho}{\rho}{trialspectrum}{local}{Trial eigenvalue or scalar trial operator.}{scalar}
\dkDeclareSymbol[Approximate eigenvalue assigned to a trial vector v. For the local operator A, its residual is A v - mu v. For unit v, the Rayleigh--Ritz choice is mu = <A v, v>.]
  {mu}{\mu}{trialspectrum}{local}{Trial spectral value.}{scalar}
\dkDeclareSymbol{a0}{a_0}{trialspectrum}{local}{Wanted trial eigenvalue or trial spectral value.}{scalar}
\dkDeclareSymbol[Ambient orthogonal or unitary transporter W with W(U) = V, equivalently W P W* = Q. Its full displacement is I - W on the whole ambient space. The real matching theorem minimizes its trace norm over every such orthogonal W.]
  {Wcompetitor}{W}{trialcompetitor}{glossary}{Orthogonal or unitary competitor carrying the source subspace onto the target.}{unitary-operator}
\dkDeclareSymbol[Orthogonal projection P onto the source subspace U in the transport problem; P = P* = P squared. A feasible orthogonal or unitary W satisfies W P W* = Q. Restricted displacement (I - W)P tests only source vectors.]
  {Psource}{P}{trialprojection}{glossary}{Orthogonal projection onto the source subspace in the transport problem.}{orthogonal-projection}
\dkDeclareSymbol{ProjSource}{P_U}{trialprojection}{local}{Orthogonal projection onto the source subspace in the transport problem.}{orthogonal-projection}
\dkDeclareSymbol[Vector x in the source subspace U. For principal vectors, unit x_j in U and y_j in V are chosen with <x_i, y_j> = 0 for i different from j and <x_j, y_j> = cos(theta_j).]
  {SourceVector}{x}{trialvector}{local}{Vector in the source subspace.}{vector}
\dkDeclareSymbol[Orthonormal coordinate frame X for the source subspace, so X* X = I. Together with a target frame Y, the singular values of X* Y are the principal cosines. In the four-dimensional counterexample X is the vertical stack of I_2 and 0.]
  {Xframe}{X}{trialframe}{glossary}{Orthonormal coordinate matrix for the source subspace in the counterexample.}{isometric-embedding}
\dkDeclareSymbol[Hilbert space F of trial coordinates. E_0 embeds F isometrically into ambient space E; A_tr and the positive sine operator sin(Theta_0) act on F, while the residual R maps F into E.]
  {Ftrial}{F}{trialframe}{glossary}{Hilbert space of trial coordinates.}{hilbert-space}
\dkDeclareSymbol[Ambient Hilbert space E over the real or complex scalars. It contains the trial subspace range(E_0) and the orthogonal exact subspaces range(F_0) and range(F_1). The source theorem permits infinite-dimensional separable E.]
  {Eambient}{E}{ambient}{glossary}{Ambient Hilbert space containing the trial and exact subspaces.}{hilbert-space}
% Zero and all-ones vectors are distinguished from scalar 0 and 1.
\dkDeclareSymbol[Zero vector of the vector or Hilbert space in context: the additive identity. The zero subspace consists only of this vector.]
  {ZeroVector}{\mathbf{0}}{fg}{glossary}{Zero vector; the additive identity of the ambient vector space.}{vector}
\dkDeclareSymbol[All-ones vector in a specified finite-dimensional coordinate space: every coordinate equals 1. This vector depends on the chosen coordinates and dimension.]
  {OnesVector}{\mathbf{1}}{fg}{local}{All-ones vector in a chosen finite-dimensional coordinate space.}{vector}
% Exact/wanted family: the target subspace, its frame, projection and rotation.
\dkDeclareSymbol[Target subspace V in the transport problem: feasible maps satisfy W(U) = V, or equivalently W P W* = Q. The matching theorem also names this subspace L.]
  {Vtarget}{V}{wanted}{glossary}{Target subspace in the transport problem.}{linear-subspace}
\dkDeclareSymbol[Exact wanted frame F_0: F_ex -> E, a linear isometry onto the desired exact spectral subspace. F_0* returns its coordinates. P_ex = F_0 F_0* is its ambient orthogonal projection; in Figure 1, F_0(h) = (h, 0).]
  {F0}{F_0}{wantedframe}{glossary}{Isometric coordinates for the desired exact spectral subspace; star denotes the adjoint.}{isometric-embedding}
\dkDeclareSymbol[Self-adjoint operator block on coordinates for the desired exact spectral subspace. Its spectrum consists of the wanted exact spectral values; in the corresponding decomposition Ahat F_0 = F_0 Lambda_0 on the operator domain.]
  {Lambda0}{\Lambda_0}{wantedspectrum}{local}{Self-adjoint block on the desired exact spectral subspace.}{self-adjoint-operator}
\dkDeclareSymbol{u}{u}{wantedvector}{local}{Vector in the desired exact subspace.}{vector}
\dkDeclareSymbol{f0}{f_0}{wantedvector}{local}{Component or coordinate vector in the desired exact subspace.}{vector}
\dkDeclareSymbol{lambda0}{\lambda_0}{wantedspectrum}{local}{Wanted exact eigenvalue or wanted spectral value.}{scalar}
\dkDeclareSymbol{lambdau}{\lambda_u}{wantedspectrum}{local}{Eigenvalue associated with a desired exact vector.}{scalar}
\dkDeclareSymbol[Exact wanted orthogonal projection P_ex: E -> E, with P_ex = F_0 F_0*. It retains the desired exact spectral component. Its complementary projection is P_ex_perp = I - P_ex = F_1 F_1*; the projected trial frame is S = P_ex_perp E_0.]
  {Pexact}{P_{\mathrm{ex}}}{wantedprojection}{glossary}{Orthogonal projection onto the desired exact subspace in the sine-theta theorem.}{orthogonal-projection}
\dkDeclareSymbol[Exact complementary orthogonal projection P_ex_perp = I - P_ex = F_1 F_1*: E -> E. It retains the unwanted exact component. In Figure 1, P_ex_perp v = S(1), whereas P_ex_perp r is the unwanted component of the residual.]
  {PexactPerp}{P_{\mathrm{ex}}^\perp}{unwantedprojection}{local}{Orthogonal projection onto the exact complementary subspace.}{orthogonal-projection}
\dkDeclareSymbol{ReflectionSigma}{\Sigma}{wantedprojection}{local}{Reflection associated with the desired exact spectral projection.}{unitary-operator}
\dkDeclareSymbol[Orthogonal projection Q onto the target subspace V in the transport problem. A feasible orthogonal or unitary W satisfies W P W* = Q, where P projects onto the source U.]
  {Qtarget}{Q}{wantedprojection}{glossary}{Orthogonal projection onto the target subspace in the transport problem.}{orthogonal-projection}
\dkDeclareSymbol{ProjTarget}{P_V}{wantedprojection}{local}{Orthogonal projection onto the target subspace in the transport problem.}{orthogonal-projection}
\dkDeclareSymbol[Canonical direct rotation D from the source U to the target V for an acute pair. It is the unitary polar factor of T = QP + (I - Q)(I - P), and acts by the principal angle on each principal plane. R is reserved for the trial residual.]
  {Rrotation}{D}{wantedrotation}{glossary}{Canonical direct rotation carrying the source subspace onto the target.}{unitary-operator}
\dkDeclareSymbol[Direct rotation D between the ranges of orthogonal projections P and Q, including the nonacute case. Relative to range(P) and kernel(P), its diagonal blocks are positive and its off-diagonal blocks are opposite adjoints. It exists when the two crossed defect spaces have equal dimension, and is unique when both defects vanish.]
  {SourceDirectRotation}{D}{wantedrotation}{local}{Direct rotation, including the nonacute case in the source-result table.}{unitary-operator}
\dkDeclareSymbol[Orthonormal coordinate frame Y for the target subspace, so Y* Y = I. In the counterexample Y is the vertical stack of cos(theta) I_2 and sin(theta) J_0; hence X* Y = cos(theta) I_2 and both principal angles equal theta.]
  {Yframe}{Y}{wantedframe}{glossary}{Orthonormal coordinate matrix for the target subspace in the counterexample.}{isometric-embedding}
\dkDeclareSymbol[Vector y in the target subspace V. For principal vectors, unit x_j in U and y_j in V are paired so that <x_j, y_j> = cos(theta_j); distinct pairs have zero cross inner product.]
  {TargetVector}{y}{wantedvector}{local}{Vector in the target subspace.}{vector}
\dkDeclareSymbol[Hilbert space F_ex of exact wanted coordinates. The isometry F_0: F_ex -> E identifies it with the desired exact spectral subspace. This space differs from H, the additive perturbation.]
  {Hexact}{F_{\mathrm{ex}}}{wantedframe}{glossary}{Hilbert space of coordinates for the desired exact spectral subspace.}{hilbert-space}% Coordinate Hilbert space (not perturbation).
% Exact/unwanted family; the trial error operators are MAGENTA, not cyan.
\dkDeclareSymbol{Uperp}{U^{\perp}}{unwanted}{local}{Orthogonal complement of the desired exact subspace in the local notation.}{linear-subspace}
\dkDeclareSymbol[Exact complementary frame F_1: G -> E, a linear isometry onto the orthogonal complement of range(F_0). F_1* returns unwanted coordinates; F_1 F_1* = I - P_ex. In Figure 1, F_1(g) = (0, g).]
  {F1}{F_1}{unwantedframe}{glossary}{Isometric coordinates for the exact complementary subspace; star denotes the adjoint.}{isometric-embedding}
\dkDeclareSymbol[Self-adjoint unwanted block Lambda_1 on G, satisfying Ahat F_1 = F_1 Lambda_1 on its domain. Its spectrum is separated from the trial spectrum of A_tr by the specified positive delta. Figure 1 uses the scalar block 0.7.]
  {Lambda1}{\Lambda_1}{unwantedspectrum}{glossary}{Self-adjoint operator on exact complementary coordinates: the unwanted spectral block.}{self-adjoint-operator}
% Spectral measure and set (distinct from the ambient space E and sine block S).
\dkDeclareSymbol[Projection-valued spectral measure E associated with a self-adjoint operator A: A = integral lambda dE(lambda). For a Borel set S, E(S) is the orthogonal projection onto the corresponding spectral subspace.]
  {SpectralMeasure}{E}{wantedprojection}{glossary}{Projection-valued spectral measure of a self-adjoint operator.}{projection-valued-measure}
\dkDeclareSymbol[Borel subset S of the real spectral axis. The spectral projection E(S) selects the part of a self-adjoint operator with spectral values in S; its range is a reducing subspace.]
  {SpectralSet}{S}{wantedspectrum}{glossary}{Borel spectral set used to select a spectral subspace.}{set}
% Generic Sylvester coefficients and right side, by source/target role.
\dkDeclareSymbol[Left coefficient B of the Sylvester map X -> B X - X C; B acts on the output space of X. For the sine bound, B = Lambda_1 on the unwanted-coordinate space G.]
  {SylvesterLeft}{B}{unwantedspectrum}{glossary}{Left coefficient of the Sylvester equation; corresponds to the unwanted spectral block.}{self-adjoint-operator}
\dkDeclareSymbol[Right coefficient C of the Sylvester map X -> B X - X C; C acts on the input space of X. For the sine bound, C = A_tr on trial space F. The historical equation also denotes this coefficient by A.]
  {SylvesterRight}{C}{trialspectrum}{glossary}{Right coefficient of the Sylvester equation; corresponds to the trial operator.}{self-adjoint-operator}
\dkDeclareSymbol[Right-hand side Y in B X - X C = Y, with the same input and output spaces as X. For the sine bound, Y = F_1* R: F -> G, the residual expressed in unwanted exact coordinates.]
  {SylvesterRHS}{Y}{resid}{glossary}{Right-hand side of the Sylvester equation, obtained by projecting the residual.}{linear-operator}
\dkDeclareSymbol{w}{w}{unwantedvector}{local}{Vector in an unwanted or complementary spectral subspace.}{vector}
\dkDeclareSymbol{f1}{f_1}{unwantedvector}{local}{Component or coordinate vector in the unwanted exact subspace.}{vector}
\dkDeclareSymbol{lambda1}{\lambda_1}{unwantedspectrum}{local}{Unwanted exact eigenvalue or complementary spectral value.}{scalar}
\dkDeclareSymbol{lambdaw}{\lambda_w}{unwantedspectrum}{local}{Eigenvalue associated with an unwanted exact vector.}{scalar}
\dkDeclareSymbol[Hilbert space G of exact complementary coordinates. F_1: G -> E embeds it onto the unwanted exact subspace; Lambda_1 acts on G and X = F_1* E_0 maps trial coordinates into G.]
  {Gunwanted}{G}{unwantedframe}{glossary}{Hilbert space of exact complementary coordinates.}{hilbert-space}
\dkDeclareSymbol{E1}{E_1}{unwantedframe}{local}{Isometric coordinates for the complementary trial subspace.}{isometric-embedding}
\dkDeclareSymbol{Vperp}{V^{\perp}}{unwanted}{local}{Orthogonal complement of the trial subspace in the local notation.}{linear-subspace}
% Angle and error geometry, including distinct rectangular and positive maps.
\dkDeclareSymbol[Individual principal angle, between 0 and pi/2. For orthonormal subspace frames X and Y, the singular values of X* Y are the principal cosines. In Figure 1, theta is the angle between the unit trial vector v and the exact wanted line.]
  {theta}{\theta}{sine}{glossary}{Principal angle between the subspaces under comparison.}{angle}
\dkDeclareSymbol[Full positive angle operator Theta = diag(Theta_0, Theta_1) for the source/target decompositions. Theta_0 compares the source with the target; Theta_1 compares their orthogonal complements. Its operator norm is the largest spectral angle. For positive c, the conditions Theta <= c and Theta < c mean operator norm at most c and strictly below c, respectively.]
  {Theta}{\Theta}{angleoperator}{glossary}{Full angle operator, the diagonal sum of the source and complementary angle blocks.}{angle-operator}
\dkDeclareSymbol[Directed principal-angle operator Theta_0 on source or trial coordinates. For a trial frame E_0 and exact wanted projection P_ex, cos(Theta_0) = sqrt(E_0* P_ex E_0), and sin(Theta_0) = sqrt(I_F - E_0* P_ex E_0). In the two-projection geometry it comes from the positive cosine compression P Q P restricted to range(P).]
  {Theta0}{\Theta_0}{angleoperator}{glossary}{Directed principal-angle operator acting on source or trial coordinates.}{angle-operator}
\dkDeclareSymbol[Complementary directed angle operator Theta_1, obtained from the projections I-P and I-Q by the same cosine-compression construction as Theta_0. It acts on the source complement, or its isometric coordinates. The full angle operator is diag(Theta_0, Theta_1).]
  {Theta1}{\Theta_1}{angleoperator}{glossary}{Directed angle operator comparing the orthogonal complements.}{angle-operator}
\dkDeclareSymbol[Right-angle source defect D_0 = range(P) intersect kernel(Q). These are source vectors orthogonal to the target. The full-angle conclusion in the Section 8 smallness criterion assumes dim(D_0) = dim(D_1); in Lean their identification is CrossedDefectsEquivalent.]
  {CrossedSourceDefect}{D_0}{fg}{glossary}{Source directions perpendicular to the target subspace.}{linear-subspace}
\dkDeclareSymbol[Right-angle target defect D_1 = kernel(P) intersect range(Q). These are target vectors orthogonal to the source. Equality of its Hilbert dimension with that of D_0 is assumed for the full-angle conclusion of Theorem 8.2.]
  {CrossedTargetDefect}{D_1}{fg}{glossary}{Target directions perpendicular to the source subspace.}{linear-subspace}
\dkDeclareSymbol[Positive complementary cosine block C_1 = cos(Theta_1). The reference complement is identified with the perturbed complement by the direct rotation D, so A_1 and Lambda_1 act on common coordinates. Theorem 8.1 compares A_1 - alpha I with C_1 (Lambda_1 - alpha I) C_1 as quadratic forms.]
  {ComplementCosineBlock}{C_1}{anglecos}{local}{Cosine of the complementary angle block in the spectral-repulsion inequality.}{self-adjoint-operator}
\dkDeclareSymbol[Scalar sine of a principal angle. For a unit principal trial vector v and exact wanted projection P_ex, sin(theta) = ||(I - P_ex)v||. In Figure 1 it is the length of the pink leg S(1).]
  {sintheta}{\sin\theta}{sine}{local}{Sine of an individual principal angle: the length of the complementary component of a unit principal trial vector.}{scalar}
\dkDeclareSymbol[Positive sine operator sin(Theta_0) = sqrt(I_F - E_0* P_ex E_0): F -> F. It equals |S| for S = (I - P_ex)E_0: F -> E. Its eigenvalues in finite dimensions are the principal sines; the sine theorem bounds its ideal norm by N(R)/delta.]
  {sinTheta0}{\sin\Theta_0}{sineoperator}{glossary}{Positive sine operator on trial coordinates; the modulus of the rectangular complementary error.}{angle-operator}
\dkDeclareSymbol[Rectangular complementary error S = (I - P_ex) E_0: F -> E. It records how much each trial vector leaks outside the exact wanted subspace. Its modulus |S| = sqrt(S* S) is sin(Theta_0): F -> F. In rank one S(1) = P_ex_perp v.]
  {S}{S}{sineblock}{glossary}{Rectangular complementary error: the trial frame projected onto the unwanted exact subspace.}{linear-operator}
\dkDeclareSymbol[Complementary-coordinate map X = F_1* E_0: F -> G, with S = F_1 X. It satisfies Lambda_1 X - X A_tr = F_1* R on domain(A_tr). Isometric multiplication gives N(X) = N(S); the separate modulus identity abs(S) = sin(Theta_0) gives N(S) = N(sin(Theta_0)), whenever these norms are finite.]
  {X}{X}{sylvester}{glossary}{Solution of the Sylvester equation; for the sine bound, the trial frame in unwanted exact coordinates.}{linear-operator}
\dkDeclareSymbol{alphaAngle}{\alpha}{anglemax}{local}{Largest principal angle.}{angle}
\dkDeclareSymbol[Scalar c = cos(theta) in the counterexample. In the general two-projection decomposition, the same role denotes a positive cosine operator C; the cosine and sine blocks commute and satisfy C squared + S squared = I.]
  {CosAngle}{c}{anglecos}{local}{Cosine of a principal angle or the corresponding positive cosine block.}{scalar-or-operator}
\dkDeclareSymbol[Scalar s = sin(theta) in the counterexample. In the general two-projection decomposition, the same role denotes a positive sine operator S; it commutes with the positive cosine block C and satisfies C squared + S squared = I.]
  {SinAngle}{s}{sineoperator}{local}{Sine of a principal angle or the corresponding positive sine block.}{scalar-or-operator}
\dkDeclareSymbol[Positive small-angle parameter in the local subspace comparison. It controls the allowed angular deviation; it is an angle parameter rather than a perturbation-operator norm.]
  {EpsilonAngle}{\varepsilon}{anglemax}{local}{Angle parameter controlling the maximum subspace deviation.}{angle}
\dkDeclareSymbol{PhiAngle}{\varphi}{angleoperator}{local}{Auxiliary angle or angle operator in the two-subspace comparison.}{angle}
\dkDeclareSymbol{Bblock}{B}{sineblock}{local}{Off-diagonal block in the two-subspace decomposition.}{linear-operator}
\dkDeclareSymbol{Cblock}{C}{sylvester}{local}{Complementary or Sylvester block in the two-subspace decomposition.}{linear-operator}
\dkDeclareSymbol{Ablock}{A}{anglecos}{local}{Positive cosine-related block in the two-subspace decomposition.}{linear-operator}
% Residual/error-displacement family.
\dkDeclareSymbol[Residual vector r = Ahat v - rho v = R(1) for the rank-one frame E_0(x) = x v. Figure 1 draws it from rho v to Ahat v. For unit v and the Rayleigh value rho = <v, Ahat v>, <v, r> = 0; a general trial scalar need not give orthogonality.]
  {r}{r}{residvector}{local}{Residual vector of an approximate eigenpair.}{vector}
\dkDeclareSymbol[Bounded residual R: F -> E, agreeing with Ahat E_0 - E_0 A_tr on domain(A_tr). The domain condition requires E_0(domain(A_tr)) inside domain(Ahat). R = 0 means that E_0 intertwines A_tr with Ahat; its norm measures the defect.]
  {R}{R}{resid}{glossary}{Bounded residual of the trial eigenproblem, agreeing with the operator difference on the trial domain.}{linear-operator}
\dkDeclareSymbol[Full displacement Delta_W = I - W on the ambient space. In the failed Ky Fan comparison W = W_(pi/4); Delta_W has singular values sqrt(2), sqrt(2), 0, 0. Multiplying by a principal-plane projection restricts its input to that plane.]
  {Kdisplacement}{\Delta_W}{displacement}{glossary}{Full displacement of a transporter from the identity.}{operator-valued-function}
\dkDeclareSymbol[Squared full displacement Dcal(W) = (I-W)* (I-W), a positive operator. Proposition 4.3 compares its unitarily invariant norm for a direct rotation D with that for every feasible transporter W.]
  {SquaredDisplacement}{\mathcal D}{displacement}{local}{Squared full displacement as a function of the transporter.}{operator-valued-function}
% Matching proof: the gradient and its spectral level are not a space or gap.
\dkDeclareSymbol[Skew-symmetric spectral gradient G of Phi_p. It commutes with W and represents first variations by d/dt Phi_p(W exp(tB)) at t = 0 equal to -tr(G B) for skew-symmetric B. Feasible stationarity gives J_src G J_src = -G.]
  {SpectralGradient}{G}{displacement}{local}{Skew-symmetric spectral gradient of the displacement power sum in the matching proof.}{linear-operator}
\dkDeclareSymbol[Positive spectral level g of the skew-symmetric gradient G: on this invariant subspace, -G squared = g squared I. Consequently G/g squares to -I and supplies the complex structure used in the real matching proof.]
  {GradientLevel}{g}{displacement}{local}{Positive gradient level; the gradient squared is minus this level squared times the identity.}{scalar}
\dkDeclareSymbol[Complex structure I_g = G/g on the positive gradient-level space V_g. It is a real isometry with square -I. Multiplication by a+ib means a times identity plus b times I_g; source reflection J_src reverses I_g and acts as conjugation.]
  {GradientComplexStructure}{\mathcal I_g}{displacement}{local}{Complex multiplication constructed from a positive gradient level.}{linear-operator}
\dkDeclareSymbol[Real skew-symmetric matrix B, with B transpose = -B, giving the orthogonal variation W exp(tB). If B commutes with the source projection P, this curve preserves the constraint W(U) = V.]
  {SkewDirection}{B}{ambient}{local}{Skew-symmetric direction for a first variation of an orthogonal transporter.}{linear-operator}
% Spectral separation family: gap and original gap are NOT identical.
\dkDeclareSymbol[Positive spectral separation in the stated configuration. In the sine theorem it separates the trial operator from the exact complementary block; Theorem 8.1 uses it for the ordered reference gap, and Theorem 8.2 for the perturbed central-band separation. Figure 1 uses delta = abs(rho - 0.7).]
  {delta}{\delta}{gap}{glossary}{Positive spectral separation in the stated gap configuration.}{scalar}
\dkDeclareSymbol{g}{g}{gaporiginal}{local}{Original spectral gap of the reference operator; distinct from the certified sine-theta separation.}{scalar}
\dkDeclareSymbol{Delta}{\Delta}{gapstat}{local}{Population or statistical spectral gap.}{scalar}
\dkDeclareSymbol[Spectral endpoint alpha. In the interval/exterior configuration, one spectrum lies between beta and alpha, the other at or below beta - delta or at or above alpha + delta. Theorem 8.2 places the selected perturbed block in the central interval. In Theorem 8.1, alpha bounds the lower reference block and alpha + delta bounds the upper block from below.]
  {alphaEndpoint}{\alpha}{gapinterval}{glossary}{Upper endpoint of the central band, or lower-block cut in the ordered gap.}{scalar}
\dkDeclareSymbol[Lower endpoint beta of the finite interval from beta to alpha in the gap hypothesis. One spectrum lies in that interval; the other lies at or below beta - delta or at or above alpha + delta.]
  {betaEndpoint}{\beta}{gapinterval}{glossary}{Lower endpoint of the finite interval in the interval/exterior gap hypothesis.}{scalar}
% Spectral/projector and auxiliary matrix objects, context-scoped.
\dkDeclareSymbol{P}{P}{fg}{local}{Orthogonal projection; its source or exact-target role is specified locally.}{orthogonal-projection}% Neutral historical Manim key; prefer Pexact/Psource.
\dkDeclareSymbol{Q}{Q}{fg}{local}{Orthogonal projection; its role is specified locally.}{orthogonal-projection}% Neutral historical Manim key; prefer Qtarget.
\dkDeclareSymbol[Two-projection intertwiner T = QP + (I - Q)(I - P), satisfying TP = QT and T* T = I - (P - Q) squared. For ||P - Q|| < 1 it is invertible; its unitary polar factor T |T| inverse is the canonical direct rotation.]
  {IntertwinerT}{T}{intertwiner}{glossary}{Two-projection intertwiner; its unitary polar factor is the acute direct rotation.}{linear-operator}
\dkDeclareSymbol[Planar pi/2-rotation matrix J_0 = ((0, -1), (1, 0)), with J_0 squared = -I_2 and J_0 transpose = -J_0. It defines the rotation C = cos(2 theta) I_2 + sin(2 theta) J_0 used to construct the real competitor.]
  {Jturn}{J_0}{turn}{glossary}{Planar rotation through $\pi/2$, with square minus the identity.}{matrix}
\dkDeclareSymbol[Partial isometry V_pol in the polar decomposition S = V_pol |S|, where |S| = sqrt(S* S). It is isometric on the orthogonal complement of kernel(S) and zero on kernel(S). For an invertible operator its polar factor is unitary.]
  {PolarU}{V_{\mathrm{pol}}}{intertwiner}{glossary}{Partial isometry in the polar decomposition; unitary when the operator is invertible.}{partial-isometry}
\dkDeclareSymbol[Planar rotation C = cos(2 theta) I_2 + sin(2 theta) J_0 through twice the principal angle. The competitor W_theta is orthogonally conjugate to C direct-sum I_2, so it rotates one two-plane and fixes the other.]
  {Cturn}{C}{turn}{glossary}{Planar rotation through twice the principal angle in the counterexample.}{matrix}
\dkDeclareSymbol[Orthogonal mixing matrix H_mix = (1/sqrt(2)) times the block matrix ((I_2, I_2), (I_2, -I_2)). It constructs W_theta = H_mix (C direct-sum I_2) H_mix transpose.]
  {Hmix}{H_{\mathrm{mix}}}{mix}{glossary}{Orthogonal mixing matrix in the counterexample.}{matrix}
\dkDeclareSymbol[Source reflection J_src = 2P - I: it acts as +I on the source subspace and -I on its orthogonal complement, so its square is I. In the matching proof stationarity gives J_src G J_src = -G for the skew-symmetric spectral gradient G.]
  {Jreflect}{J_{\mathrm{src}}}{trialprojection}{glossary}{Reflection equal to twice the source projection minus the identity.}{linear-operator}
\dkDeclareSymbol[Skew angle generator J_ang for a direct rotation D = exp(J_ang Theta). On each principal plane it is the oriented rotation through pi/2; it vanishes on the common fixed directions. It commutes with the full angle operator Theta.]
  {AngleGenerator}{J_{\mathrm{ang}}}{turn}{glossary}{Skew angle generator in the direct-rotation result summaries.}{linear-operator}
\dkDeclareSymbol{Dblock}{D}{mix}{local}{Auxiliary diagonal or block matrix in the counterexample.}{matrix}
\dkDeclareSymbol[Orthogonal projection onto the first principal plane span(e_0, e_3) in the four-dimensional counterexample. For Delta_W = I - W_(pi/4), the Ky Fan 2-norm of Delta_W Omega_1 equals 2.]
  {OmegaOne}{\Omega_1}{turn}{glossary}{First principal-plane projection in the failed Ky Fan comparison.}{orthogonal-projection}
\dkDeclareSymbol[Orthogonal projection onto the second principal plane span(e_1, e_2) in the four-dimensional counterexample. For Delta_W = I - W_(pi/4), the Ky Fan 2-norm of Delta_W Omega_2 equals 2, although the trace norm of Delta_W is only 2 sqrt(2).]
  {OmegaTwo}{\Omega_2}{intertwiner}{glossary}{Second principal-plane projection in the failed Ky Fan comparison.}{orthogonal-projection}
\dkDeclareSymbol{Dplanar}{D}{wantedrotation}{local}{Planar direct-rotation block.}{unitary-operator}
\dkDeclareSymbol{Hreflection}{H}{trialcompetitor}{local}{Planar reflection carrying the source line to the target.}{linear-operator}
% Unitarily invariant gauge family; N, M are NOT a norm versus its input.
\dkDeclareSymbol[Unitarily invariant norm N, or a symmetric operator-ideal norm on its domain of finiteness. It satisfies N(U T V) = N(T) for unitaries and N(B T) <= ||B||op N(T). In the sine theorem it is normalized by N(u v*) = ||u|| ||v|| and has Ky Fan dominance where both norms are finite. Both R and sin(Theta_0) must belong to its ideal.]
  {NormN}{N}{norm}{glossary}{Unitarily invariant norm on its domain of finiteness; in the sine theorem, normalized by $N(uv^*)=\|u\|\|v\|$.}{norm}
\dkDeclareSymbol[Sine-theta separation G_delta(B,C): finite interval/exterior separation in either spectral order, or ordered separation B <= a I and C >= (a+delta) I, or the reverse. The ordered alternatives allow both spectra to be unbounded.]
  {SineGap}{\mathcal G_\delta}{gap}{local}{Interval/exterior separation or either ordered separation, as in the sine-theta theorem.}{predicate}
\dkDeclareSymbol[Auxiliary unitarily invariant norm M witnessing the Q-norm property N(T) squared = M(T* T). For the Frobenius norm N, one can take M to be the trace norm.]
  {NormM}{M}{normaux}{local}{Auxiliary unitarily invariant norm used to define a Q-norm through the squared modulus.}{norm}
\dkDeclareSymbol[Q-norm property: a unitarily invariant norm N is a Q-norm when N(T) squared = M(T* T) for some unitarily invariant norm M. The Frobenius norm is an example, with M the trace norm; the repaired full-displacement theorem uses this norm class.]
  {NormQ}{Q}{normaux}{glossary}{Q-norm: its square is a unitarily invariant norm of the squared modulus.}{norm}
\dkDeclareSymbol[Singleton trace cost s_i = min(4 sin(theta_i/2), 2) for angle label i. It transports one real principal plane independently, choosing a rotation of cost 4 sin(theta_i/2) or a reflection of cost 2.]
  {SingletonCost}{s}{norm}{glossary}{Minimum trace displacement for one real principal plane: choose a rotation or a reflection.}{scalar-cost}
\dkDeclareSymbol[Paired trace cost e_ij = 4 sin((theta_i + theta_j)/2) for distinct angle labels i,j. It is attained by coupling their principal planes in a real four-dimensional block. The edge weight s_i + s_j - e_ij is the difference between the cost of two singletons and this paired cost.]
  {PairCost}{e}{norm}{glossary}{Trace displacement cost of coupling two principal planes in a four-dimensional block.}{scalar-cost}
\dkDeclareSymbol[Scalar phase cost f_p(t) = (2 sin(t/2)) to the power p, for 0 <= t <= pi. A real planar rotation through t contributes two copies, so its displacement power cost is 2 f_p(t).]
  {MatchingPhaseCost}{f_p}{norm}{local}{Displacement power cost for one scalar phase.}{scalar-cost}
\dkDeclareSymbol[Singleton power cost a_i(p) = min(2 f_p(theta_i), 2 to the power p). These are the rotation and reflection choices for one principal plane. At p=1 it equals the trace-cost weight s_i.]
  {SingletonPowerCost}{a}{norm}{local}{Minimum displacement power cost for one angle.}{scalar-cost}
\dkDeclareSymbol[Paired power cost b_ij(p) = min over d <= tau <= h of 2(f_p(h+tau) + f_p(h-tau)), where h=(theta_i+theta_j)/2 and d=abs(theta_i-theta_j)/2. This is the positive-determinant pair minimum; unrestricted block optimization also compares two singletons. At p=1 it equals e_ij.]
  {PairPowerCost}{b}{norm}{local}{Positive-determinant displacement power cost for a pair of angles.}{scalar-cost}
\dkDeclareSymbol[Matching M on the finite general graph of positive principal-angle labels. Each label belongs to at most one edge; unmatched labels use singleton transporters. Minimizing paired and singleton costs is equivalent to maximum-weight matching with edge weights g_ij = s_i + s_j - e_ij. The empty matching is allowed.]
  {Matching}{M}{normaux}{glossary}{Matching of positive principal-angle labels: each label is used in at most one pair.}{matching}
\dkDeclareSymbol[Displacement power sum Phi_p(W) = sum_j s_j(I - W)^p. The matching proof first minimizes it for 1 < p < 2, then lets p decrease to 1, where it becomes the trace norm ||I - W||_1.]
  {PowerCost}{\Phi}{norm}{local}{Sum of the p-th powers of displacement singular values.}{scalar-cost}
% Auxiliary matching-proof notation: neutral ink, with the same hover registry.
\dkDeclareSymbol[Off-diagonal gradient block G_01 mapping the source complement K perpendicular into K. Stationarity makes the diagonal blocks of G zero, so G = ((0,G_01),(-G_01 transpose,0)) in the source/complement decomposition.]
  {GradientOffDiagonal}{G_{01}}{fg}{local}{Complement-to-source block of the skew-symmetric gradient.}{linear-operator}
\dkDeclareSymbol[Orthogonal group O(n): real n by n matrices W satisfying W transpose times W = I. Its elements preserve Euclidean inner products and lengths. The notation O(2) refers to the same group in dimension two.]
  {OrthogonalGroup}{O(n)}{fg}{local}{Group of real orthogonal matrices.}{group}
\dkDeclareSymbol[Feasible transporters F = {W in O(n) : W(K) = L}. Equivalently W P_K W transpose = P_L. This nonempty compact set is the optimization domain throughout the matching proof.]
  {FeasibleTransport}{\mathcal F}{fg}{local}{Orthogonal maps carrying the source subspace onto the target.}{set}
\dkDeclareSymbol[Singular value s_j(T), the nonnegative square root of an eigenvalue of T transpose times T, counted with multiplicity. Summing s_j(I-W) gives the trace displacement; summing its p-th powers gives Phi_p(W). These singular values differ from the singleton weights s_i.]
  {TransportSingularValues}{\symglyph{SingularValueBase}_j}{fg}{local}{A displacement singular value in the matching proof.}{scalar}
\dkDeclareSymbol[Scalar singular-value magnitude s between zero and two for the displacement I-W of an orthogonal matrix. Uniform convergence of s to the power p to s on this interval permits the trace-norm limit. This variable differs from the rotation parameter s used earlier in the proof.]
  {DisplacementSize}{s}{fg}{local}{Scalar displacement singular-value magnitude in the uniform limit.}{scalar}
\dkDeclareSymbol[All-singleton cost C_0 = sum_i s_i, before any two angle labels are paired. A matching has cost C_0 minus the sum of its selected edge weights.]
  {SingletonBaseline}{C_0}{fg}{local}{Trace displacement when every angle uses its singleton transporter.}{scalar-cost}
\dkDeclareSymbol[Matching power cost C_p(M): sum of b_ij(p) over selected pairs plus a_i(p) over unmatched labels. Its minimum equals the minimum of Phi_p over feasible transporters. At p=1 it is C(M).]
  {MatchingPowerCost}{C_p}{fg}{local}{Total paired and singleton power cost of a matching.}{scalar-cost}
\dkDeclareSymbol[Invariant orthogonal block V_a preserved by the source projection, target projection, and minimizing transporter. The displacement cost adds over these blocks, and their principal-angle lists combine into the original list.]
  {MatchingBlock}{V_a}{fg}{local}{Simultaneous invariant block in the matching proof.}{linear-subspace}
\dkDeclareSymbol[Gradient-level space V_g: the eigenspace of -G squared for a positive eigenvalue g squared. It is preserved by G, W, and the source reflection J_src. Multiplication by i on it is G/g.]
  {GradientLevelSpace}{V_g}{fg}{local}{Positive gradient-level eigenspace.}{linear-subspace}
\dkDeclareSymbol[Eigenphase varphi of W on V_g viewed as a complex space with multiplication by i = G/g. It lies between zero and pi and satisfies f_p prime(varphi) = g. There are at most two such phases.]
  {GradientEigenphase}{\varphi}{fg}{local}{One of the at most two phases at a positive gradient level.}{angle}
\dkDeclareSymbol[Orthogonal projection E onto the first phase space of W on V_g. Here W = exp(i varphi_1) E + exp(i varphi_2) (I-E). In coordinates where the source reflection is conjugation, E = X + iY.]
  {GradientPhaseProjection}{E}{fg}{local}{Projection onto the first phase space in a gradient-level space.}{orthogonal-projection}
\dkDeclareSymbol[Real part X of the phase projection E = X + iY. It is a real symmetric matrix with spectrum in the closed interval from zero to one, satisfying X squared - Y squared = X and XY + YX = Y. This X differs from the source-coordinate and Sylvester matrices elsewhere in the paper.]
  {ProjectionRealPart}{X}{fg}{local}{Real symmetric part of the phase projection.}{linear-operator}
\dkDeclareSymbol[Imaginary part Y of the phase projection E = X + iY. It is a real skew-symmetric matrix. It maps the lambda eigenspace of X to the 1-lambda eigenspace and constructs the small invariant planes.]
  {ProjectionImaginaryPart}{Y}{fg}{local}{Real skew-symmetric imaginary part of the phase projection.}{linear-operator}
\dkDeclareSymbol[Eigenvalue lambda in the closed interval from zero to one of the real part X of the phase projection. It is local to the block construction and differs from the original operator's spectral parameter.]
  {ProjectionEigenvalue}{\lambda}{fg}{local}{Eigenvalue of the real part of the phase projection.}{scalar}
\dkDeclareSymbol[Plane coupling eta = sqrt(lambda (1-lambda)) for 0 < lambda < 1. If u is a unit lambda eigenvector of X, then v = Yu/eta completes its invariant real plane.]
  {ProjectionCoupling}{\eta}{fg}{local}{Normalization used to construct an invariant real plane.}{scalar}
\dkDeclareSymbol[Pair mean angle h = (theta_i + theta_j)/2. Together with d = (theta_i - theta_j)/2 after ordering theta_i >= theta_j, it parametrizes the two-angle block cost.]
  {PairMeanAngle}{h}{fg}{local}{Mean of the two ordered principal angles.}{angle}
\dkDeclareSymbol[Pair half-difference d = (theta_i - theta_j)/2 after ordering theta_i >= theta_j. It satisfies 0 <= d <= h and h+d <= pi/2.]
  {PairHalfDifference}{d}{fg}{local}{Half the difference of the two ordered principal angles.}{angle}
\dkDeclareSymbol[Scalar pair parameter tau in the closed interval from d to h. The paired power cost is the minimum of 2(f_p(h+tau) + f_p(h-tau)) over this interval. At p=1 a minimum is attained at tau=h.]
  {PairPhaseParameter}{\tau}{fg}{local}{Scalar parameter in the two-angle power-cost minimum.}{scalar}
\dkDeclareSymbol[Attaining rotation angle gamma = arccos(cos(tau)/cos(d)). Choosing U_0 = U_1 = R_gamma realizes phase parameters x=h and y=tau in the four-dimensional block.]
  {PairAttainingAngle}{\gamma}{fg}{local}{Rotation angle realizing a chosen scalar pair parameter.}{angle}
\dkDeclareSymbol[Four-dimensional direct rotation D = ((C,-S),(S,C)) in source and complementary principal coordinates, with C the diagonal principal cosines and S the diagonal principal sines. Every transporter on this block is D diag(U_0,U_1).]
  {PairDirectRotation}{D}{fg}{local}{Direct rotation on the two principal planes.}{unitary-operator}
\dkDeclareSymbol[Diagonal cosine block C = diag(cos(theta_i), cos(theta_j)) in the four-dimensional direct rotation D. This matrix differs from the scalar displacement cost C(M).]
  {PairCosineBlock}{C}{fg}{local}{Diagonal principal-cosine block.}{matrix}
\dkDeclareSymbol[Diagonal sine block S = diag(sin(theta_i), sin(theta_j)) in the four-dimensional direct rotation D. This matrix is local to the matching proof.]
  {PairSineBlock}{S}{fg}{local}{Diagonal principal-sine block.}{matrix}
\dkDeclareSymbol[Free planar orthogonal matrices U_0 and U_1 acting on the two-dimensional source and its two-dimensional complement, respectively. Multiplying D by diag(U_0,U_1) parametrizes every transporter on this four-dimensional block.]
  {PairFreeOrthogonal}{U}{fg}{local}{Free source or complementary planar orthogonal factor.}{unitary-operator}
\dkDeclareSymbol[Phase parameter x in the closed interval from h to pi-h for a positive-determinant four-dimensional transporter. Its eigenvalues are exp(plus or minus i(x+y)) and exp(plus or minus i(x-y)).]
  {PairPhaseX}{x}{fg}{local}{First phase parameter of a four-dimensional transporter.}{angle}
\dkDeclareSymbol[Phase parameter y in the closed interval from d to pi-d for a positive-determinant four-dimensional transporter. Together with x it gives rotational phases x+y and x-y, folded into the closed interval from zero to pi for the displacement cost.]
  {PairPhaseY}{y}{fg}{local}{Second phase parameter of a four-dimensional transporter.}{angle}
\dkDeclareSymbol[Scalar a = cos(x) in the four-dimensional phase calculation. Together with b = cos(y), it expresses the trace as 4ab and determines the reciprocal characteristic polynomial.]
  {PairPhaseCosineX}{a}{fg}{local}{Cosine of the first phase parameter in the trace calculation.}{scalar}
\dkDeclareSymbol[Scalar b = cos(y) in the four-dimensional phase calculation, paired with a = cos(x). These scalars are local to the characteristic-polynomial calculation.]
  {PairPhaseCosineY}{b}{fg}{local}{Cosine of the second phase parameter in the trace calculation.}{scalar}
\dkDeclareSymbol[Rotation parameter r: the two free orthogonal factors use planar rotation angles r+s and r-s. It appears in cos(y) = cos(d) cos(r) for two rotations.]
  {PairRotationR}{r}{fg}{local}{Half the sum of the two free planar rotation angles.}{angle}
\dkDeclareSymbol[Rotation parameter s: the two free orthogonal factors use planar rotation angles r+s and r-s. It appears in cos(x) = cos(h) cos(s) for two rotations. This scalar differs from the singular values and singleton costs.]
  {PairRotationS}{s}{fg}{local}{Half the difference of the two free planar rotation angles.}{angle}
\dkDeclareSymbol[Planar reflection J_2 = diag(1,-1). Every two-dimensional orthogonal matrix of determinant -1 can be written as a planar rotation times J_2.]
  {PairPlanarReflection}{J_2}{fg}{local}{Fixed planar reflection used to parametrize determinant -1 factors.}{unitary-operator}
\dkDeclareSymbol[Real planar rotation R_t = ((cos(t),-sin(t)),(sin(t),cos(t))). Its two displacement singular values are both 2 times abs(sin(t/2)). In the matching proof a rotation through t in the closed interval from zero to pi contributes 2 f_p(t).]
  {MatchingPlanarRotation}{\symglyph{PlanarRotationBase}_t}{fg}{local}{Planar rotation through the indicated angle.}{unitary-operator}
% Statistical application / operator provenance.
\dkDeclareSymbol{Sigma}{\Sigma}{old}{local}{Population covariance matrix.}{matrix}
\dkDeclareSymbol{SigmaHat}{\widehat\Sigma}{current}{local}{Estimated or perturbed covariance matrix.}{matrix}
\dkDeclareSymbol{YWSV}{V}{wanted}{local}{Orthonormal coordinates for the population eigenspace in the statistical application.}{isometric-embedding}
\dkDeclareSymbol{YWSVhat}{\widehat V}{trialframe}{local}{Estimated eigenspace frame in the statistical application.}{isometric-embedding}

% Figure-specific objects retain their own meanings when bare letters overlap.
\dkDeclareSymbol[Signed rotation angle phi of the displayed map on its invariant real two-plane. At theta = pi/4 the direct rotation turns its two principal planes by opposite signed angles of magnitude pi/4; the competitor turns its moving plane by pi/2 and fixes its other plane.]
  {PlaneRotationAngle}{\varphi}{sine}{local}{Signed rotation angle of the orthogonal map on the displayed invariant plane.}{angle}
\dkDeclareSymbol[Unit vector q in the displayed invariant plane. The map T is the direct rotation D or the competitor W in that panel; Tq is the image and (I - T)q = q - Tq is the displacement.]
  {FigureInput}{q}{ambient}{local}{Unit input to the orthogonal map shown in this invariant-plane panel.}{vector}
\dkDeclareSymbol[Image Tq of the panel unit input q under T = D or W. Both q and Tq lie on the unit circle; the green displacement arrow joins Tq to q, so it represents q - Tq.]
  {FigureImage}{Tq}{wantedvector}{local}{Image of the unit input under the orthogonal map shown in this panel.}{vector}
\dkDeclareSymbol[Displacement (I - T)q = q - Tq, drawn from the image Tq to the unit input q. On a plane rotated through phi its length is 2 |sin(phi/2)|; it vanishes on a fixed plane.]
  {DisplacementVector}{(I-T)q}{displacement}{local}{Difference between a unit input and its image under the displayed orthogonal map.}{vector}
\dkDeclareSymbol[First principal plane Pi_1 = span(e_0, e_3) for the source-target pair in the real four-dimensional counterexample. At theta = pi/4 the direct rotation acts on this plane by a quarter of pi, with displacement singular value sqrt(2 - sqrt(2)).]
  {CounterPrincipalOne}{\Pi_1}{fg}{glossary}{First principal plane in the real counterexample: the span of e0 and e3.}{linear-subspace}
\dkDeclareSymbol[Second principal plane Pi_2 = span(e_1, e_2) for the source-target pair in the real four-dimensional counterexample. The direct rotation acts by the same angle magnitude as on Pi_1, with the opposite sign in the displayed plane coordinates.]
  {CounterPrincipalTwo}{\Pi_2}{fg}{glossary}{Second principal plane in the real counterexample: the span of e1 and e2.}{linear-subspace}
\dkDeclareSymbol[Moving plane M = range(W - I) = span(e_0 + e_2, e_1 + e_3) for the competitor. W rotates it by 2 theta; at theta = pi/4 the two displacement singular values contributed by this plane are both sqrt(2).]
  {CounterMovingPlane}{M}{ambient}{glossary}{Moving plane of the competing transporter: the range of W minus the identity.}{linear-subspace}
\dkDeclareSymbol[Fixed plane Z = kernel(W - I) = span(e_0 - e_2, e_1 - e_3) for the competitor. W acts as the identity here, so this plane contributes two zero displacement singular values.]
  {CounterFixedPlane}{Z}{ambient}{glossary}{Fixed plane of the competing transporter: the kernel of W minus the identity.}{linear-subspace}
\dkDeclareSymbol[Direct-rotation displacement singular value a = 2 sin(pi/8) = sqrt(2 - sqrt(2)) at theta = pi/4. Each of its two principal planes contributes two copies, giving trace displacement 4a.]
  {CounterChordA}{a}{displacement}{glossary}{Direct-rotation displacement singular value in the counterexample: sqrt(2 - sqrt(2)).}{scalar}
\dkDeclareSymbol[Competitor displacement singular value b = 2 sin(pi/4) = sqrt(2) at theta = pi/4. Its moving plane contributes two copies; the fixed plane contributes two zeros, giving trace displacement 2b.]
  {CounterChordB}{b}{displacement}{glossary}{Competing-transporter displacement singular value in the counterexample: square root of 2.}{scalar}
\dkDeclareSymbol[Singular value s_j(I - T) of the full ambient displacement, ordered from largest to smallest. Each rotating real two-plane contributes two values 2 |sin(phi/2)|; fixed directions contribute zeros. Their sum is the trace cost.]
  {DisplacementSingular}{\symglyph{SingularValueBase}_j}{displacement}{local}{An individual singular value of the full displacement from the identity.}{scalar}
\dkDeclareSymbol[Edge weight g_ij = s_i + s_j - e_ij, the difference between the cost of two singletons and the paired cost. A matching has total cost C_0 minus the sum of its edge weights, with C_0 = sum_i s_i. Maximizing this sum gives the minimum transport cost; edges of negative weight may be omitted.]
  {MatchingEdgeWeight}{g_{ij}}{current}{glossary}{Edge weight: the cost of two singletons minus their paired cost.}{scalar-weight}
\dkDeclareSymbol[Total trace-norm displacement cost. For a transporter W, C(W) = ||I - W||_1, the sum of its displacement singular values. For a matching M, C(M) is the sum of its paired and unmatched singleton costs.]
  {TraceCost}{C}{norm}{glossary}{Trace-norm displacement cost of the indicated orthogonal transporter or matching.}{scalar-cost}

% Spectral values have their own key, distinct from an operator block.
\dkDeclareSymbol[Spectral parameter lambda of a self-adjoint operator. In a finite-dimensional eigenvalue list lambda_j(A) it denotes an individual eigenvalue; in A = integral lambda dE(lambda), it is the real integration variable.]{SpectralValue}{\lambda}{wantedspectrum}{local}{Real spectral parameter or individual eigenvalue of a self-adjoint operator.}{scalar}



% Conventional mathematical notation. Hover text applies even when ink is plain.
% These are local (non-glossary) unless promoted explicitly by the author.
% IdentityTwo and IdentityFour derive their I from Identity. Singular-value
% registrations derive s_j from SingularValueBase; SingletonCost has its own s.
\dkDeclareSymbol[Identity operator on the vector or Hilbert space in context. Its dimension is inferred from the expression and need not be finite.]{Identity}{I}{plain}{local}{Identity operator on the space in context.}{identity-operator}
\dkDeclareSymbol[Two-by-two identity matrix, acting on a two-dimensional coordinate space.]{IdentityTwo}{\symglyph{Identity}_2}{plain}{local}{Two-by-two identity matrix.}{identity-operator}
\dkDeclareSymbol[Four-by-four identity matrix, acting on a four-dimensional coordinate space.]{IdentityFour}{\symglyph{Identity}_4}{plain}{local}{Four-by-four identity matrix.}{identity-operator}
\dkDeclareSymbol[Two-by-two zero matrix. Unlike the bold zero vector, this is a linear operator/matrix with two rows and two columns.]{ZeroMatrixTwo}{0_2}{plain}{local}{Two-by-two zero matrix.}{matrix}
\dkDeclareSymbol[The constant pi, the ratio of the circumference of a circle to its diameter. Angles in this article use radians.]{PiConstant}{\pi}{plain}{local}{The constant pi (angles are in radians).}{constant}
\dkDeclareSymbol[The field of real numbers.]{RealField}{\mathbb R}{plain}{local}{The real numbers.}{field}
\dkDeclareSymbol[The field of complex numbers.]{ComplexField}{\mathbb C}{plain}{local}{The complex numbers.}{field}
\dkDeclareSymbol[Generic Hilbert space used in the two-projection discussion; distinct from the paper's ambient space symbol E in the sine-theta theorem.]{GeneralHilbertSpace}{\mathcal H}{plain}{local}{Generic Hilbert space in a local discussion.}{hilbert-space}
\dkDeclareSymbol[The spectrum sigma(T) of an operator T: the set of complex scalars for which T minus the scalar identity is not boundedly invertible. For self-adjoint T the spectrum is real.]{SpectrumSymbol}{\sigma}{plain}{local}{Spectrum of an operator.}{operator-notation}
\dkDeclareSymbol[Base letter for singular values, presently s. It is distinct from the base s for matching singleton costs.]{SingularValueBase}{s}{plain}{local}{Base letter for singular values.}{notation-base}
\dkDeclareSymbol[Ordered singular value s_j(T), equal to the j-th eigenvalue of |T| for a finite matrix; distinguish this from the matching singleton cost s_j.]{SingularValueS}{\symglyph{SingularValueBase}_j}{plain}{glossary}{The j-th singular value, under the currently printed s_j convention.}{scalar}
\dkDeclareSymbol[Alternate indexed singular-value notation sigma_j(T). The current manuscript also uses s_j(T); this registered alias is available for a later deliberate notation decision.]{SingularValueSigma}{\sigma_j}{plain}{local}{Indexed singular value under the sigma_j convention.}{scalar}
\dkDeclareSymbol[Operator norm ||T||_op, the supremum of ||Tx|| over unit vectors x.]{OperatorNormNotation}{\|\cdot\|_{\mathrm{op}}}{plain}{glossary}{Operator norm.}{norm}
\dkDeclareSymbol[Frobenius norm ||T||_F, the square root of the sum of squared matrix entries, equal to the finite-dimensional Hilbert--Schmidt norm.]{FrobeniusNormNotation}{\|\cdot\|_F}{plain}{glossary}{Frobenius norm.}{norm}
\dkDeclareSymbol[Trace norm ||T||_1, the sum of singular values when finite.]{TraceNormNotation}{\|\cdot\|_1}{plain}{glossary}{Trace (nuclear) norm.}{norm}
\dkDeclareSymbol[Hilbert--Schmidt norm ||T||_HS, the square root of the sum of squared singular values when finite.]{HilbertSchmidtNormNotation}{\|\cdot\|_{\mathrm{HS}}}{plain}{local}{Hilbert--Schmidt norm.}{norm}
\dkDeclareSymbol[Ky Fan j-norm, the sum of the largest j singular values.]{KyFanNormNotation}{\|\cdot\|_{(j)}}{plain}{glossary}{Ky Fan j-norm.}{norm}
\dkDeclareSymbol[Interval/exterior spectral separation predicate I_delta(B,C), specifying a central closed interval separated by delta from the other operator's spectrum.]{IntervalSeparation}{\mathcal I}{plain}{local}{Interval/exterior separation predicate.}{predicate}
\dkDeclareSymbol[Ordered spectral separation predicate O_delta(B,C), requiring one self-adjoint operator to lie below a spectral threshold and the other above it by at least delta.]{OrderedSeparation}{\mathcal O}{plain}{local}{Ordered spectral separation predicate.}{predicate}
\dkDeclareSymbol[Positive parameter epsilon controlling a small perturbation in the two-dimensional eigenvalue example.]{EigenvalueExampleEpsilon}{\varepsilon}{plain}{local}{Small positive parameter in the eigenvalue-example matrices.}{scalar}
\dkDeclareSymbol[Positive lower bound epsilon for the trial-coordinate map E_0, as in E_0* E_0 >= epsilon squared times the identity.]{FrameLowerBoundEpsilon}{\varepsilon}{plain}{local}{Lower-bound parameter for the trial-coordinate embedding.}{scalar}
\dkDeclareSymbol[Rotation angle vartheta of a planar orthogonal block in the matching-proof first-variation calculation.]{MatchingPlanarAngle}{\vartheta}{plain}{local}{Planar rotation angle in the matching-proof calculation.}{angle}


% Ordinary named mathematical operators already used throughout the manuscript.
% Keeping these here makes their spelling a single editorial choice.
\dkDeclareSymbol[Range of a linear operator: the set of its output vectors.]{RangeOperator}{\operatorname{ran}}{plain}{local}{Range (image) of an operator.}{operator-notation}
\dkDeclareSymbol[Kernel (null space) of a linear operator: inputs mapped to zero.]{KernelOperator}{\operatorname{ker}}{plain}{local}{Kernel (null space) of an operator.}{operator-notation}
\dkDeclareSymbol[Domain of an operator: the vectors on which the operator is defined, important for unbounded operators.]{DomainOperator}{\operatorname{dom}}{plain}{local}{Domain of an operator.}{operator-notation}
\DeclareRobustCommand{\DKRange}{\sym{RangeOperator}}
\DeclareRobustCommand{\DKKernel}{\sym{KernelOperator}}
\DeclareRobustCommand{\DKDom}{\sym{DomainOperator}}

% Norm-formatting macros: the operand changes by occurrence but the notation
% lives here. The outer hover describes the norm and does not add ink color.
\DeclareRobustCommand{\dkOperatorNorm}[1]{\like{OperatorNormNotation}{\|#1\|_{\mathrm{op}}}}
\DeclareRobustCommand{\dkTraceNorm}[1]{\like{TraceNormNotation}{\|#1\|_1}}
\DeclareRobustCommand{\dkFrobeniusNorm}[1]{\like{FrobeniusNormNotation}{\|#1\|_F}}
\DeclareRobustCommand{\dkHilbertSchmidtNorm}[1]{\like{HilbertSchmidtNormNotation}{\|#1\|_{\mathrm{HS}}}}
\DeclareRobustCommand{\dkKyFanNorm}[2]{\like{KyFanNormNotation}{\|#1\|_{(#2)}}}
\DeclareRobustCommand{\dkSchattenNorm}[2]{\like{SchattenNormNotation}{\|#1\|_{S_{#2}}}}

\dkDeclareSymbol[Schatten p-norm ||T||_{S_p}, obtained from the l^p-norm of the singular values and finite for operators in the Schatten p-class.]{SchattenNormNotation}{\|\cdot\|_{S_p}}{plain}{local}{Schatten p-norm.}{norm}


\dkDeclareSymbol[Positive modulus of an operator T: |T| = (T* T)^(1/2).]{OperatorModulusNotation}{|\cdot|}{plain}{glossary}{Positive modulus of an operator.}{operator-notation}
\dkDeclareSymbol[Linear span of the indicated vectors, i.e. the smallest linear subspace that contains them.]{SpanOperator}{\operatorname{span}}{plain}{local}{Linear span of vectors.}{operator-notation}
\dkDeclareSymbol[Diagonal or block-diagonal matrix/operator with the displayed entries or blocks.]{DiagonalOperator}{\operatorname{diag}}{plain}{local}{Diagonal or block-diagonal operator.}{linear-operator}
\dkDeclareSymbol[Trace of a trace-class operator, or ordinary matrix trace in finite dimensions.]{TraceOperator}{\operatorname{tr}}{plain}{local}{Trace of an operator.}{operator-notation}
\dkDeclareSymbol[Dimension of a vector space or subspace.]{DimensionOperator}{\dim}{plain}{local}{Dimension of a vector space.}{operator-notation}
\dkDeclareSymbol[Rank of a linear map or matrix: dimension of its range.]{RankOperator}{\operatorname{rank}}{plain}{local}{Rank of a linear map or matrix.}{operator-notation}
\dkDeclareSymbol[Distance between sets of spectral values: the infimum of pairwise distances.]{SpectralDistanceOperator}{\operatorname{dist}}{plain}{local}{Distance between two spectral sets.}{operator-notation}
\DeclareRobustCommand{\dkModulus}[1]{\like{OperatorModulusNotation}{|#1|}}

% Context-specific glyphs previously hidden in variants of other semantic keys.
% Keeping both the roles and the exact printed forms makes collisions explicit.
\dkDeclareSymbol[Nonzero vector x tested in the Rayleigh quotient, before imposing unit length.]{RayleighTestVector}{x}{trialvector}{local}{Test vector in the Rayleigh quotient.}{vector}
\dkDeclareSymbol[Positive cosine block C in the generic two-projection decomposition of Halmos; C and S commute and C squared plus S squared is the identity.]{HalmosCosineBlock}{C}{anglecos}{local}{Cosine block in the Halmos two-projection representation.}{linear-operator}
\dkDeclareSymbol[Positive sine block S in the generic two-projection decomposition of Halmos; C and S commute and C squared plus S squared is the identity.]{HalmosSineBlock}{S}{sineoperator}{local}{Sine block in the Halmos two-projection representation.}{linear-operator}
\dkDeclareSymbol[Riesz spectral projection P_Gamma defined by a contour Gamma that encloses selected spectral values.]{RieszProjection}{P_{\symglyph{RieszContour}}}{wantedprojection}{local}{Spectral projection associated with a contour Gamma.}{orthogonal-projection}
\dkDeclareSymbol[Historical Sylvester right-hand coefficient matrix A, in B X - X A = Y. Distinct from the current paper's reference operator A.]{HistoricalSylvesterRight}{A}{trialspectrum}{local}{Right-hand coefficient A in the historical Sylvester equation.}{self-adjoint-operator}
\dkDeclareSymbol[Scalar coefficient b in the Sylvester prototype (b-a)x=y, with separated scalar spectral values a and b.]{SylvesterScalarLeft}{b}{unwantedspectrum}{local}{Scalar left coefficient b in a Sylvester equation.}{scalar}
\dkDeclareSymbol[Scalar coefficient a in the Sylvester prototype (b-a)x=y.]{SylvesterScalarRight}{a}{trialspectrum}{local}{Scalar right coefficient a in a Sylvester equation.}{scalar}
\dkDeclareSymbol[Scalar unknown x in the Sylvester prototype (b-a)x=y.]{SylvesterScalarUnknown}{x}{sylvester}{local}{Scalar unknown x in a Sylvester equation.}{scalar}
\dkDeclareSymbol[Scalar right side y in the Sylvester prototype (b-a)x=y.]{SylvesterScalarRHS}{y}{resid}{local}{Scalar right side y in a Sylvester equation.}{scalar}
\dkDeclareSymbol[Source subspace K of a real Euclidean space in the matching proof; an orthogonal transporter sends K onto L.]{MatchingSourceSubspace}{K}{trial}{local}{Source subspace K in the matching proof.}{linear-subspace}
\dkDeclareSymbol[Target subspace L of a real Euclidean space in the matching proof; an orthogonal transporter sends K onto L.]{MatchingTargetSubspace}{L}{wanted}{local}{Target subspace L in the matching proof.}{linear-subspace}
\dkDeclareSymbol[Orthogonal projection P_L onto the matching proof's target subspace L.]{MatchingTargetProjection}{P_L}{wantedprojection}{local}{Orthogonal projection onto the matching target L.}{orthogonal-projection}
\dkDeclareSymbol[Real two-dimensional orthogonal group O(2): planar rotations and reflections.]{PlanarOrthogonalGroup}{O(2)}{fg}{local}{Orthogonal matrices of size two.}{group}
\dkDeclareSymbol[Base letter R used to denote planar rotation operators R_t in the matching proof.]{PlanarRotationBase}{R}{fg}{local}{Base letter for planar rotations.}{notation-base}
\dkDeclareSymbol[Planar rotation R_(r+s) through the sum of scalar phases r and s.]{MatchingPlanarRotationSum}{\symglyph{PlanarRotationBase}_{r+s}}{fg}{local}{Planar rotation through phase r+s.}{unitary-operator}
\dkDeclareSymbol[Planar rotation R_gamma through phase gamma in the block rotation calculation.]{MatchingPlanarRotationGamma}{\symglyph{PlanarRotationBase}_{\symglyph{MatchingPhaseGamma}}}{fg}{local}{Planar rotation through phase gamma.}{unitary-operator}

% Specific symbols in the historical and appendix examples. They retain their
% existing black-ink notation but receive locally accurate hover definitions.
\dkDeclareSymbol[Contour Gamma, a positively oriented closed curve in the resolvent set enclosing a selected isolated spectral cluster.]{RieszContour}{\Gamma}{plain}{local}{Contour used in the Riesz spectral-projection formula.}{contour}
\dkDeclareSymbol[Even integer nu in a tested Ky Fan comparison, with nu/2 principal-plane terms on the right side.]{KyFanEvenIndex}{\nu}{plain}{local}{Even integer in the principal-plane Ky Fan comparison.}{scalar}
\dkDeclareSymbol[Orthogonal projection Omega_k onto the k-th principal plane in the published displacement comparison.]{PrincipalPlaneProjectionBase}{\Omega}{plain}{local}{Projection onto the indicated principal plane.}{notation-base}
\dkDeclareSymbol[Scalar gap threshold gamma in the ordered Sylvester spectral separation B >= (gamma+delta)I and C <= gamma I.]{SylvesterGapThreshold}{\gamma}{plain}{local}{Spectral separation threshold in a Sylvester estimate.}{scalar}
\dkDeclareSymbol[Projection Omega_t from a net or sequence that converges strongly to the identity in the singular-value approximation lemma.]{StrongApproxProjection}{\Omega_t}{plain}{local}{Orthogonal projection converging strongly to the identity.}{orthogonal-projection}
\dkDeclareSymbol[First orthogonal projection Omega used in the two-block pinching inequality for a unitarily invariant norm.]{PinchingProjectionOne}{\Omega}{plain}{local}{First projection in the pinching inequality.}{orthogonal-projection}
\dkDeclareSymbol[Second orthogonal projection Upsilon used in the two-block pinching inequality for a unitarily invariant norm.]{PinchingProjectionTwo}{\Upsilon}{plain}{local}{Second projection in the pinching inequality.}{orthogonal-projection}
\dkDeclareSymbol[Finite-rank projection Gamma used to retain the image of a chosen finite-rank compression K Psi.]{CompressionProjectionImage}{\Gamma}{plain}{local}{Range projection in the finite-rank compression lemma.}{orthogonal-projection}
\dkDeclareSymbol[Finite-rank projection Psi used to select input vectors in a nearly maximizing singular-value compression.]{CompressionProjectionInput}{\Psi}{plain}{local}{Domain projection in the finite-rank compression lemma.}{orthogonal-projection}
\dkDeclareSymbol[Small positive tolerance eta in the finite-rank compression approximation estimate.]{CompressionTolerance}{\eta}{plain}{local}{Tolerance controlling the off-diagonal compression norm.}{scalar}

\dkDeclareSymbol[Rotation phase gamma in a planar block from the matching proof.]{MatchingPhaseGamma}{\gamma}{plain}{local}{Rotation phase gamma in the matching calculation.}{scalar}

% Two common indices with local mathematical meanings distinct from other deltas.
\dkDeclareSymbol[Kronecker delta: delta_ij is one when the indices i and j agree and zero otherwise. This is not the spectral-gap parameter delta.]{KroneckerDelta}{\delta}{plain}{local}{Kronecker delta: equals 1 for equal indices and 0 otherwise.}{scalar}
\dkDeclareSymbol[Sign variable epsilon taking the values -1 and +1 in the factorization over two phase choices.]{PhaseSignChoice}{\epsilon}{plain}{local}{Sign variable taking values minus one and plus one.}{scalar}

% Documented notation correspondence. The source column is a faithful
% transcription of each cited original context, NOT an asserted universal DK
% meaning. The original A_0 is deliberately listed for two distinct roles:
% its reference block and the comparison/trial operator. The source uses the
% same printed symbol for both; this paper is free to distinguish them.
\dkDeclareNotationSource{DK1970}{Davis and Kahan (1970), The Rotation of Eigenvectors by a Perturbation. III}{https://doi.org/10.1137/0707001}
\dkDeclareNotationSource{Seelmann2019}{Seelmann (2019), Semidefinite Perturbations in the Subspace Perturbation Problem}{https://arxiv.org/abs/1708.02463}
\dkDeclareNotationSource{Blechta2020}{Blechta (2020), Stability of Linear GMRES Convergence with Respect to Compact Perturbations}{https://arxiv.org/abs/2005.12960}
% Arguments: {key}{source ID}{source glyph}{section or equation}{role-specific qualification}
\dkDeclareNotationAlternative{A}{DK1970}{A}{\S1, (1.2)}{Reference self-adjoint operator.}
\dkDeclareNotationAlternative{H}{DK1970}{H}{\S1, (1.2)}{Additive perturbation.}
\dkDeclareNotationAlternative{Acurrent}{DK1970}{A+H}{\S1}{Perturbed operator.}
\dkDeclareNotationAlternative{E0}{DK1970}{E_0}{\S1, (1.1)}{Source isometric frame.}
\dkDeclareNotationAlternative{E1}{DK1970}{E_1}{\S1, (1.1)}{Reference complement frame.}
\dkDeclareNotationAlternative{F0}{DK1970}{F_0}{\S1}{Exact-target frame.}
\dkDeclareNotationAlternative{F1}{DK1970}{F_1}{\S1}{Exact-target complement frame.}
\dkDeclareNotationAlternative{ReferenceBlock0}{DK1970}{A_0}{\S1, (1.2)}{Reference source block.}
\dkDeclareNotationAlternative{ReferenceBlock1}{DK1970}{A_1}{\S1, (1.2)}{Reference complement block.}
\dkDeclareNotationAlternative{A0}{DK1970}{A_0}{\S1, residual discussion}{Residual comparison operator (the reference block in DK).}
% We deliberately leave Lambda0 and Lambda1 without DK source-glyph metadata:
% the available source transcription does not reliably preserve the source
% accents/decorations on the perturbed spectral-block symbols.
\dkDeclareNotationAlternative{Psource}{DK1970}{P}{\S1}{Reference/source projection.}
\dkDeclareNotationAlternative{Qtarget}{DK1970}{Q}{\S1}{Target projection.}
\dkDeclareNotationAlternative{Pexact}{DK1970}{Q}{\S1}{Exact-target projection.}
\dkDeclareNotationAlternative{B}{DK1970}{B}{\S1, (1.2)}{Off-diagonal coupling block.}
\dkDeclareNotationAlternative{Theta0}{DK1970}{\Theta_0}{\S1}{Source-side angle operator.}
\dkDeclareNotationAlternative{Theta1}{DK1970}{\Theta_1}{\S1}{Complementary angle operator.}
\dkDeclareNotationAlternative{Theta}{DK1970}{\Theta}{\S1}{Full angle operator.}
\dkDeclareNotationAlternative{R}{DK1970}{R}{\S1, residual discussion}{Trial residual.}
\dkDeclareNotationAlternative{Rrotation}{DK1970}{U}{\S1, (1.18)}{Canonical direct rotation.}
\dkDeclareNotationAlternative{Wcompetitor}{DK1970}{V}{\S4}{Competing unitary transporter.}
\dkDeclareNotationAlternative{delta}{DK1970}{\delta}{\S2}{Spectral separation parameter.}
% Contemporary alternatives are recorded separately and do not enter the DK
% correspondence table. Their citation locations must be individually checked.
\dkDeclareNotationAlternative{Eambient}{Seelmann2019}{\mathcal H}{p.~1}{Ambient Hilbert space.}
\dkDeclareNotationAlternative{SingularValueS}{Blechta2020}{\sigma_j}{Thm.~1.1}{Ordered singular values of a compact operator.}
 in paper.tex)
% or pass \def\DKMonochromeMath{1} before % One switch: false = journal style; true = colored symbols and author--year citations.
\newif\ifDKAuthorStyle

%\DKAuthorStylefalse
\DKAuthorStyletrue

\ifdefined\DKUseAuthorStyle\DKAuthorStyletrue\fi
\ifdefined\DKJournalStyle\DKAuthorStylefalse\fi
\ifDKAuthorStyle
  \PassOptionsToClass{sn-mathphys-ay}{sn-jnl}
\else
  \PassOptionsToClass{sn-mathphys-num}{sn-jnl}
\fi
\documentclass[pdflatex]{sn-jnl} % PDF engine option; compiled with LuaLaTeX.
% PDF version used for the manuscript and its symbol annotations.
\pdfvariable minorversion=7\relax

\usepackage{amsmath,amssymb,amsthm,mathtools}
\usepackage{booktabs,tabularx,array,longtable}
\usepackage[section]{placeins}
\usepackage{float}
\usepackage{microtype}
\usepackage{graphicx}
\usepackage{pgf} % Live vector figures with LaTeX symbol labels and tooltips
\usepackage{xurl}
\ifDKAuthorStyle
  \setcitestyle{authoryear,round,aysep={,}}
  \hypersetup{colorlinks=true,allcolors=blue}
\else
  \setcitestyle{numbers,square}
  \hypersetup{hidelinks}
\fi
\usepackage[nameinlink]{cleveref}
\usepackage[T1]{fontenc}
\usepackage{lmodern} % Scalable LaTeX lettering, shared with the Manim slides
\usepackage{xcolor}
% Draft controls. Set these independently for review, print, or arXiv versions.
\newif\ifDKAnonymous
\DKAnonymousfalse
\ifdefined\DKAnonymousReview\DKAnonymoustrue\fi
\newif\ifDKColoredSymbols
\ifDKAuthorStyle\DKColoredSymbolstrue\else\DKColoredSymbolsfalse\fi
\ifdefined\DKMonochromeMath\DKColoredSymbolsfalse\fi
% The colors retain the Manim glossary's light-theme mathematical roles.
% Semantic, multi-shade mathematical notation for the Davis--Kahan paper.
% The six core families follow visualizations/dkvis/palette.py (LIGHT).
% Subshades distinguish objects that are RELATED but not IDENTICAL.
% Values were chosen for legibility against white manuscript paper; some
% amber, green, and vermilion anchors are a little darker than the slides.
%
% Color control (independent of hover reminders):
%   \DKMathColorsfalse   (after \input{math_symbol_colors.tex} in paper.tex)
% or pass \def\DKMonochromeMath{1} before \input{paper.tex}, without editing.
% EVERY \sym, \like, and \dkrole obeys this switch, including the glossary.
% No color is required to identify a mathematical object: the glyphs, indices,
% definitions and adjacent prose also distinguish them.
%
% Families (base hue): exact/target=blue, trial/source=amber, exact unwanted=
% cyan, angles and their operators=magenta, residual/displacement=green,
% gaps=violet. Norm gauges=slate; provenance retains brown/steel/vermilion.
% P_ex is the exact-target projection in sine-theta; P is the source
% projection in transport. R denotes a residual and D a direct rotation.
% Semantic keys retain their roles even when the printed notation changes.
\newif\ifDKMathColors
\DKMathColorstrue
\ifdefined\DKMonochromeMath
  \DKMathColorsfalse
\fi

% Family anchors (print-compatible versions of the Manim light palette).
\definecolor{dkfg}{HTML}{1B1F24}
\definecolor{dkmuted}{HTML}{6B7280}
\definecolor{dkwanted}{HTML}{1F6FD1}
\definecolor{dktrial}{HTML}{A66705}
\definecolor{dkunwanted}{HTML}{007C91}
\definecolor{dksine}{HTML}{D12A7B}
\definecolor{dkresid}{HTML}{087D43}
\definecolor{dkgap}{HTML}{6E4BD8}
\definecolor{dkold}{HTML}{705519}
\definecolor{dkcurrent}{HTML}{054669}
\definecolor{dkperturb}{HTML}{BB491B}
\definecolor{dkrefuted}{HTML}{C62828}

% Related yet distinguishable shades. These are NOT new unrelated categories:
% objects inherit their parent family; a different shade signals a different
% mathematical object (e.g., a frame versus its subspace or its projection).
\definecolor{dkwantedframe}{HTML}{1254A9}
\definecolor{dkwantedprojection}{HTML}{08618D}
\definecolor{dkwantedrotation}{HTML}{4248B0}
\definecolor{dkwantedspectrum}{HTML}{345F8F}
\definecolor{dkwantedvector}{HTML}{1468AA}
\definecolor{dktrialframe}{HTML}{905600}
\definecolor{dktrialprojection}{HTML}{79511B}
\definecolor{dktrialcompetitor}{HTML}{A04D10}
\definecolor{dktrialspectrum}{HTML}{79591C}
\definecolor{dktrialvector}{HTML}{9B522A}
\definecolor{dkunwantedframe}{HTML}{006477}
\definecolor{dkunwantedspectrum}{HTML}{00546B}
\definecolor{dkunwantedprojection}{HTML}{145B6D}
\definecolor{dkunwantedvector}{HTML}{256879}
\definecolor{dkangleoperator}{HTML}{923478}
\definecolor{dksineoperator}{HTML}{AA275A}
\definecolor{dksineblock}{HTML}{B43F70}
\definecolor{dksylvester}{HTML}{804074}
\definecolor{dkanglemax}{HTML}{B32185}
\definecolor{dkanglecos}{HTML}{824276}
\definecolor{dkresidvector}{HTML}{086038}
\definecolor{dkdisplacement}{HTML}{226C50}
\definecolor{dkgaporiginal}{HTML}{4E4D89}
\definecolor{dkgapstat}{HTML}{7039A7}
\definecolor{dkgapinterval}{HTML}{574A82}
\definecolor{dkambient}{HTML}{36516A}
\definecolor{dknorm}{HTML}{475875}
\definecolor{dknormaux}{HTML}{66526E}
\definecolor{dkintertwiner}{HTML}{4A6280}
\definecolor{dkturn}{HTML}{47647D}
\definecolor{dkmix}{HTML}{505B65}

% Coloring primitives. "plain" suppresses color without suppressing hovers.
% All descendants can retain a semantic role through \like{key}{expression}.
\expandafter\def\csname dknocolor@plain\endcsname{1}
\newcommand{\dkrole}[2]{%
  \ifDKMathColors
    \ifcsname dknocolor@#1\endcsname #2\else\textcolor{dk#1}{#2}\fi
  \else #2\fi
}
% An optional plain-text reminder supplies definitions and types on hover.
% Required fields: {semantic key}{printed glyph}{color role}{glossary|local}
%                  {short definition}{mathematical type}. The "glossary" flag records whether an
% entry appears in the curated printed notation table; "local" does not.
% This registry is the sole authority for symbol meaning and printed form.
% Use role "plain" for an uncolored symbol that still carries hover text.
% Changing the glyph is a separate editorial decision: this audit preserves
% the current E, F, F_ex, G, and other existing printed choices.
\newcommand{\dkDeclareSymbol}[7][]{%
  \expandafter\def\csname dkgloss@#2\endcsname{#6}%
  \expandafter\def\csname dkglossaryflag@#2\endcsname{#5}%
  \if\relax\detokenize{#1}\relax
    \expandafter\pdfstringdef\csname dkpdfgloss@#2\endcsname{#7: #6}%
  \else
    \expandafter\pdfstringdef\csname dkpdfgloss@#2\endcsname{#7: \detokenize{#1}}%
  \fi
  \expandafter\def\csname dktype@#2\endcsname{#7}%
  \expandafter\def\csname dkglyph@#2\endcsname{#3}%
  \expandafter\def\csname dkrole@#2\endcsname{#4}%
  \expandafter\def\csname dksym@#2\endcsname{\dkrole{#4}{#3}}%
}
% Unstyled printed glyph, for variants constructed through \like.
% Example: \like{E0}{\symglyph{E0}^*} derives E_0^* from E0's
% one canonical glyph; retyping E_0 here would defeat central control.
\DeclareRobustCommand{\symglyph}[1]{\csname dkglyph@#1\endcsname}
\DeclareRobustCommand{\sym}[1]{%
  \ifcsname dksym@#1\endcsname
    \dkSymbolTooltip{#1}{\csname dksym@#1\endcsname}%
  \else
    \PackageError{dk-math-colors}{Unknown symbol key #1}{See math_symbol_colors.tex}%
  \fi
}
\DeclareRobustCommand{\like}[2]{%
  \ifcsname dkrole@#1\endcsname
    \dkSymbolTooltip{#1}{\dkrole{\csname dkrole@#1\endcsname}{#2}}%
  \else
    \PackageError{dk-math-colors}{Unknown role key #1}{See math_symbol_colors.tex}%
  \fi
}

% A short description belongs to the semantic key, not the bare letter.
% The same descriptions are available to the notation appendix.
\DeclareRobustCommand{\dkSymbolDescription}[1]{\csname dkgloss@#1\endcsname}
\DeclareRobustCommand{\dkSymbolGlossaryStatus}[1]{\csname dkglossaryflag@#1\endcsname}
\DeclareRobustCommand{\dkSymbolType}[1]{\csname dktype@#1\endcsname}

% Citable alternative notation is metadata, not a request to rename glyphs.
% The source identifier, the source's exact printed glyph, and the specific
% source location are all mandatory. No alternative is inferred from a name.
% One key may have alternatives from several works or multiple passages.
\newcommand{\dkDeclareNotationSource}[3]{%
  \expandafter\def\csname dksourcetitle@#1\endcsname{#2}%
  \expandafter\def\csname dksourceurl@#1\endcsname{#3}%
}
\newcommand{\dkDavisKahanRows}{}
\makeatletter
\newcommand{\dkDeclareNotationAlternative}[5]{%
  \ifcsname dkglyph@#1\endcsname\else
    \PackageError{dk-math-colors}{Unknown alternative key #1}{Declare the symbol first}%
  \fi
  \ifcsname dksourcetitle@#2\endcsname\else
    \PackageError{dk-math-colors}{Unknown notation source #2}{Declare the source first}%
  \fi
  \expandafter\def\csname dkaltglyph@#2@#1\endcsname{#3}%
  \expandafter\def\csname dkaltlocator@#2@#1\endcsname{#4}%
  \expandafter\def\csname dkaltnote@#2@#1\endcsname{#5}%
  \def\dk@source{#2}%
  \def\dk@original{DK1970}%
  \ifx\dk@source\dk@original
    \g@addto@macro\dkDavisKahanRows{\dkDavisKahanRow{#1}}%
  \fi
}
\makeatother
% A correspondence row is produced from the SAME registry records used by
% \sym and by the source-census script. It is never spelled a second time.
\newcommand{\dkDavisKahanRow}[1]{%
  \csname dkaltnote@DK1970@#1\endcsname &
  $\csname dkaltglyph@DK1970@#1\endcsname$ &
  $\sym{#1}$ &
  \csname dkaltlocator@DK1970@#1\endcsname\\
}

\newif\ifDKSymbolTooltips
\DKSymbolTooltipstrue
\ifdefined\DKNoSymbolTooltips\DKSymbolTooltipsfalse\fi
\newsavebox{\DKTooltipBox}
\newcount\DKTooltipCount
% Register only shipped widgets through hyperref's existing form-field
% machinery. Poppler/Okular need the catalog's field list, not just page
% annotations. The deferred aux write also discards measuring-only copies.
\makeatletter
\newcommand{\dkRegisterTooltipField}{%
  \expandafter\HyField@@AddToFields\expandafter{\number\pdffeedback lastannot}%
}
\makeatother
\AtBeginDocument{\ifDKSymbolTooltips\Form[NeedAppearances=false]\fi}
\AtEndDocument{\ifDKSymbolTooltips\endForm\fi}
% Native LuaLaTeX PDF widgets: no navigation action or JavaScript. Avoid
% pdfcomment's soulpos dependency, which previously broke local Vim builds.
% Measure in the CURRENT math style, then emit the original expression as-is:
% boxing the visible math would change spacing, scripts and line breaking.
\newcommand{\dkSymbolTooltip}[2]{%
  \ifDKSymbolTooltips
    \begingroup
      \DKSymbolTooltipsfalse % suppress nested annotations while measuring
      \ifmmode
        \edef\DKTooltipStyle{\ifcase\mathstyle
          \displaystyle\or\displaystyle\or\textstyle\or\textstyle\or
          \scriptstyle\or\scriptstyle\or\scriptscriptstyle\or\scriptscriptstyle\fi}%
        \sbox{\DKTooltipBox}{$\DKTooltipStyle#2$}%
      \else
        \leavevmode\sbox{\DKTooltipBox}{#2}%
      \fi
      \global\advance\DKTooltipCount by 1
      \pdfextension annot width\wd\DKTooltipBox
        height\ht\DKTooltipBox depth\dp\DKTooltipBox
        {/Subtype /Widget /FT /Btn /Ff 65536 /F 512
         /T (dk-symbol-\the\DKTooltipCount) /TU (\csname dkpdfgloss@#1\endcsname)
         /H /N /BS << /W 0 >> /Border [0 0 0]}%
      \dkRegisterTooltipField
      #2%
    \endgroup
  \else
    #2%
  \fi
}
\pdfstringdefDisableCommands{%
  \def\sym#1{\csname dkglyph@#1\endcsname}%
  \def\like#1#2{#2}%
  \def\symglyph#1{\csname dkglyph@#1\endcsname}%
}

% Declarations: [detailed hover], key, glyph, role, glossary|local, short text, type.
% A local entry has full semantic/hover support, but no notation-table row.
% When two keys have the same glyph, it is an auditable collision, not an
% instruction to change the printed notation automatically.
% This registry includes local proof notation. The notation appendix selects
% the symbols useful for reference; registration does not add a printed row.
% Manim's original symbol keys, now split into semantic subshades.
% Ambient and perturbed operators (provenance is independent of geometry).
\dkDeclareSymbol[Reference self-adjoint operator A on the ambient Hilbert space; the perturbed operator is Ahat = A + H. In a single-operator eigenproblem, A itself is the operator being approximated. An unbounded A acts only on domain(A).]
  {A}{A}{old}{glossary}{Reference self-adjoint operator before perturbation.}{self-adjoint-operator}
\dkDeclareSymbol{A1}{A_1}{old}{local}{Complementary block of the reference operator.}{self-adjoint-operator}
\dkDeclareSymbol{a1}{a_1}{old}{local}{Reference eigenvalue or complementary spectral value.}{scalar}
\dkDeclareSymbol[Reference block A_0: the restriction of A to the range of the reducing source projection P, expressed in source coordinates. In Theorem 8.1 it lies at or below alpha; in Theorem 8.2 its spectrum lies in the enlarged central interval. The perturbed blocks are denoted Lambda_0 and Lambda_1.]
  {ReferenceBlock0}{A_0}{old}{glossary}{Reference operator restricted to the source subspace in Section 8.}{self-adjoint-operator}
\dkDeclareSymbol[Reference block A_1: the restriction of A to the kernel of the reducing source projection P, expressed in complementary coordinates. In Theorem 8.1 it lies at or above alpha + delta in quadratic-form order.]
  {ReferenceBlock1}{A_1}{old}{glossary}{Reference operator restricted to the source complement in Section 8.}{self-adjoint-operator}
\dkDeclareSymbol[Self-adjoint ambient operator Ahat being studied; in the perturbation setting Ahat = A + H. Its action on a trial vector is compared with the proposed eigenvalue times that vector.]
  {Ahat}{\widehat A}{current}{local}{Operator under study; in the perturbation setting, the reference operator plus its perturbation.}{self-adjoint-operator}
\dkDeclareSymbol[Self-adjoint ambient operator Ahat being studied. In the perturbation setting Ahat = A + H. The trial residual is R = Ahat E_0 - E_0 A_tr, and the exact complementary block satisfies Ahat F_1 = F_1 Lambda_1 on its domain.]
  {Acurrent}{\widehat A}{current}{glossary}{Operator under study; in the perturbation setting, the reference operator plus its perturbation.}{self-adjoint-operator}
\dkDeclareSymbol[Additive self-adjoint perturbation H = Ahat - A. A bound on its norm can control the residual of an eigenspace of one operator relative to the other.]
  {H}{H}{perturb}{glossary}{Additive self-adjoint perturbation of the reference operator.}{self-adjoint-operator}
\dkDeclareSymbol{H0}{H_0}{perturb}{local}{Wanted diagonal block of the perturbation.}{self-adjoint-operator}
\dkDeclareSymbol{H1}{H_1}{perturb}{local}{Unwanted diagonal block of the perturbation.}{self-adjoint-operator}
\dkDeclareSymbol{B}{B}{perturb}{local}{Off-diagonal perturbation or coupling block.}{linear-operator}
\dkDeclareSymbol{b}{b}{perturb}{local}{Scalar perturbation or coupling coefficient.}{scalar}
% Trial/source family: subspace, frame, compressed operator, projector, vector.
\dkDeclareSymbol{U}{U}{wanted}{local}{Desired exact subspace in the local geometric notation.}{linear-subspace}% Manim's U remains wanted in local slide notation.
\dkDeclareSymbol{V}{V}{trial}{local}{Trial subspace in the local geometric notation.}{linear-subspace}% Manim's V remains trial in local slide notation.
\dkDeclareSymbol[Source subspace U to be carried onto the target V by an orthogonal or unitary map W. In the transport sections P projects onto U and Q onto V. The matching theorem also names the source subspace K.]
  {Usource}{U}{trial}{glossary}{Source subspace to be transported onto the target subspace.}{linear-subspace}
\dkDeclareSymbol[Trial frame E_0: F -> E, a linear isometry from trial-coordinate space F into ambient space E. Its range is the trial subspace; E_0* returns trial coordinates and E_0* E_0 = I_F. In rank one, E_0(x) = x v.]
  {E0}{E_0}{trialframe}{glossary}{Isometric embedding of the trial subspace; its adjoint returns trial coordinates.}{isometric-embedding}
\dkDeclareSymbol[Self-adjoint trial operator A_tr on trial-coordinate space F. The residual compares Ahat E_0 with E_0 A_tr. Rayleigh--Ritz chooses A_tr = E_0* Ahat E_0 when this compression is defined; the theorem allows other trial operators.]
  {A0}{A_{\mathrm{tr}}}{trialspectrum}{glossary}{Self-adjoint operator on trial coordinates; its eigenvalues are the trial spectral values.}{self-adjoint-operator}
\dkDeclareSymbol[Trial vector: a proposed approximate eigenvector. In the rank-one geometry v has unit length and E_0(x) = x v. Its residual for trial scalar rho is r = Ahat v - rho v.]
  {v}{v}{trialvector}{local}{Trial vector or vector in the trial subspace.}{vector}
\dkDeclareSymbol[Scalar trial operator in rank one, A_tr(x) = rho x. Figure 1 uses the Rayleigh value rho = <v, Ahat v> for unit v; then rho v is the projection of Ahat v onto the trial line and r is perpendicular to v.]
  {rho}{\rho}{trialspectrum}{local}{Trial eigenvalue or scalar trial operator.}{scalar}
\dkDeclareSymbol[Approximate eigenvalue assigned to a trial vector v. For the local operator A, its residual is A v - mu v. For unit v, the Rayleigh--Ritz choice is mu = <A v, v>.]
  {mu}{\mu}{trialspectrum}{local}{Trial spectral value.}{scalar}
\dkDeclareSymbol{a0}{a_0}{trialspectrum}{local}{Wanted trial eigenvalue or trial spectral value.}{scalar}
\dkDeclareSymbol[Ambient orthogonal or unitary transporter W with W(U) = V, equivalently W P W* = Q. Its full displacement is I - W on the whole ambient space. The real matching theorem minimizes its trace norm over every such orthogonal W.]
  {Wcompetitor}{W}{trialcompetitor}{glossary}{Orthogonal or unitary competitor carrying the source subspace onto the target.}{unitary-operator}
\dkDeclareSymbol[Orthogonal projection P onto the source subspace U in the transport problem; P = P* = P squared. A feasible orthogonal or unitary W satisfies W P W* = Q. Restricted displacement (I - W)P tests only source vectors.]
  {Psource}{P}{trialprojection}{glossary}{Orthogonal projection onto the source subspace in the transport problem.}{orthogonal-projection}
\dkDeclareSymbol{ProjSource}{P_U}{trialprojection}{local}{Orthogonal projection onto the source subspace in the transport problem.}{orthogonal-projection}
\dkDeclareSymbol[Vector x in the source subspace U. For principal vectors, unit x_j in U and y_j in V are chosen with <x_i, y_j> = 0 for i different from j and <x_j, y_j> = cos(theta_j).]
  {SourceVector}{x}{trialvector}{local}{Vector in the source subspace.}{vector}
\dkDeclareSymbol[Orthonormal coordinate frame X for the source subspace, so X* X = I. Together with a target frame Y, the singular values of X* Y are the principal cosines. In the four-dimensional counterexample X is the vertical stack of I_2 and 0.]
  {Xframe}{X}{trialframe}{glossary}{Orthonormal coordinate matrix for the source subspace in the counterexample.}{isometric-embedding}
\dkDeclareSymbol[Hilbert space F of trial coordinates. E_0 embeds F isometrically into ambient space E; A_tr and the positive sine operator sin(Theta_0) act on F, while the residual R maps F into E.]
  {Ftrial}{F}{trialframe}{glossary}{Hilbert space of trial coordinates.}{hilbert-space}
\dkDeclareSymbol[Ambient Hilbert space E over the real or complex scalars. It contains the trial subspace range(E_0) and the orthogonal exact subspaces range(F_0) and range(F_1). The source theorem permits infinite-dimensional separable E.]
  {Eambient}{E}{ambient}{glossary}{Ambient Hilbert space containing the trial and exact subspaces.}{hilbert-space}
% Zero and all-ones vectors are distinguished from scalar 0 and 1.
\dkDeclareSymbol[Zero vector of the vector or Hilbert space in context: the additive identity. The zero subspace consists only of this vector.]
  {ZeroVector}{\mathbf{0}}{fg}{glossary}{Zero vector; the additive identity of the ambient vector space.}{vector}
\dkDeclareSymbol[All-ones vector in a specified finite-dimensional coordinate space: every coordinate equals 1. This vector depends on the chosen coordinates and dimension.]
  {OnesVector}{\mathbf{1}}{fg}{local}{All-ones vector in a chosen finite-dimensional coordinate space.}{vector}
% Exact/wanted family: the target subspace, its frame, projection and rotation.
\dkDeclareSymbol[Target subspace V in the transport problem: feasible maps satisfy W(U) = V, or equivalently W P W* = Q. The matching theorem also names this subspace L.]
  {Vtarget}{V}{wanted}{glossary}{Target subspace in the transport problem.}{linear-subspace}
\dkDeclareSymbol[Exact wanted frame F_0: F_ex -> E, a linear isometry onto the desired exact spectral subspace. F_0* returns its coordinates. P_ex = F_0 F_0* is its ambient orthogonal projection; in Figure 1, F_0(h) = (h, 0).]
  {F0}{F_0}{wantedframe}{glossary}{Isometric coordinates for the desired exact spectral subspace; star denotes the adjoint.}{isometric-embedding}
\dkDeclareSymbol[Self-adjoint operator block on coordinates for the desired exact spectral subspace. Its spectrum consists of the wanted exact spectral values; in the corresponding decomposition Ahat F_0 = F_0 Lambda_0 on the operator domain.]
  {Lambda0}{\Lambda_0}{wantedspectrum}{local}{Self-adjoint block on the desired exact spectral subspace.}{self-adjoint-operator}
\dkDeclareSymbol{u}{u}{wantedvector}{local}{Vector in the desired exact subspace.}{vector}
\dkDeclareSymbol{f0}{f_0}{wantedvector}{local}{Component or coordinate vector in the desired exact subspace.}{vector}
\dkDeclareSymbol{lambda0}{\lambda_0}{wantedspectrum}{local}{Wanted exact eigenvalue or wanted spectral value.}{scalar}
\dkDeclareSymbol{lambdau}{\lambda_u}{wantedspectrum}{local}{Eigenvalue associated with a desired exact vector.}{scalar}
\dkDeclareSymbol[Exact wanted orthogonal projection P_ex: E -> E, with P_ex = F_0 F_0*. It retains the desired exact spectral component. Its complementary projection is P_ex_perp = I - P_ex = F_1 F_1*; the projected trial frame is S = P_ex_perp E_0.]
  {Pexact}{P_{\mathrm{ex}}}{wantedprojection}{glossary}{Orthogonal projection onto the desired exact subspace in the sine-theta theorem.}{orthogonal-projection}
\dkDeclareSymbol[Exact complementary orthogonal projection P_ex_perp = I - P_ex = F_1 F_1*: E -> E. It retains the unwanted exact component. In Figure 1, P_ex_perp v = S(1), whereas P_ex_perp r is the unwanted component of the residual.]
  {PexactPerp}{P_{\mathrm{ex}}^\perp}{unwantedprojection}{local}{Orthogonal projection onto the exact complementary subspace.}{orthogonal-projection}
\dkDeclareSymbol{ReflectionSigma}{\Sigma}{wantedprojection}{local}{Reflection associated with the desired exact spectral projection.}{unitary-operator}
\dkDeclareSymbol[Orthogonal projection Q onto the target subspace V in the transport problem. A feasible orthogonal or unitary W satisfies W P W* = Q, where P projects onto the source U.]
  {Qtarget}{Q}{wantedprojection}{glossary}{Orthogonal projection onto the target subspace in the transport problem.}{orthogonal-projection}
\dkDeclareSymbol{ProjTarget}{P_V}{wantedprojection}{local}{Orthogonal projection onto the target subspace in the transport problem.}{orthogonal-projection}
\dkDeclareSymbol[Canonical direct rotation D from the source U to the target V for an acute pair. It is the unitary polar factor of T = QP + (I - Q)(I - P), and acts by the principal angle on each principal plane. R is reserved for the trial residual.]
  {Rrotation}{D}{wantedrotation}{glossary}{Canonical direct rotation carrying the source subspace onto the target.}{unitary-operator}
\dkDeclareSymbol[Direct rotation D between the ranges of orthogonal projections P and Q, including the nonacute case. Relative to range(P) and kernel(P), its diagonal blocks are positive and its off-diagonal blocks are opposite adjoints. It exists when the two crossed defect spaces have equal dimension, and is unique when both defects vanish.]
  {SourceDirectRotation}{D}{wantedrotation}{local}{Direct rotation, including the nonacute case in the source-result table.}{unitary-operator}
\dkDeclareSymbol[Orthonormal coordinate frame Y for the target subspace, so Y* Y = I. In the counterexample Y is the vertical stack of cos(theta) I_2 and sin(theta) J_0; hence X* Y = cos(theta) I_2 and both principal angles equal theta.]
  {Yframe}{Y}{wantedframe}{glossary}{Orthonormal coordinate matrix for the target subspace in the counterexample.}{isometric-embedding}
\dkDeclareSymbol[Vector y in the target subspace V. For principal vectors, unit x_j in U and y_j in V are paired so that <x_j, y_j> = cos(theta_j); distinct pairs have zero cross inner product.]
  {TargetVector}{y}{wantedvector}{local}{Vector in the target subspace.}{vector}
\dkDeclareSymbol[Hilbert space F_ex of exact wanted coordinates. The isometry F_0: F_ex -> E identifies it with the desired exact spectral subspace. This space differs from H, the additive perturbation.]
  {Hexact}{F_{\mathrm{ex}}}{wantedframe}{glossary}{Hilbert space of coordinates for the desired exact spectral subspace.}{hilbert-space}% Coordinate Hilbert space (not perturbation).
% Exact/unwanted family; the trial error operators are MAGENTA, not cyan.
\dkDeclareSymbol{Uperp}{U^{\perp}}{unwanted}{local}{Orthogonal complement of the desired exact subspace in the local notation.}{linear-subspace}
\dkDeclareSymbol[Exact complementary frame F_1: G -> E, a linear isometry onto the orthogonal complement of range(F_0). F_1* returns unwanted coordinates; F_1 F_1* = I - P_ex. In Figure 1, F_1(g) = (0, g).]
  {F1}{F_1}{unwantedframe}{glossary}{Isometric coordinates for the exact complementary subspace; star denotes the adjoint.}{isometric-embedding}
\dkDeclareSymbol[Self-adjoint unwanted block Lambda_1 on G, satisfying Ahat F_1 = F_1 Lambda_1 on its domain. Its spectrum is separated from the trial spectrum of A_tr by the specified positive delta. Figure 1 uses the scalar block 0.7.]
  {Lambda1}{\Lambda_1}{unwantedspectrum}{glossary}{Self-adjoint operator on exact complementary coordinates: the unwanted spectral block.}{self-adjoint-operator}
% Spectral measure and set (distinct from the ambient space E and sine block S).
\dkDeclareSymbol[Projection-valued spectral measure E associated with a self-adjoint operator A: A = integral lambda dE(lambda). For a Borel set S, E(S) is the orthogonal projection onto the corresponding spectral subspace.]
  {SpectralMeasure}{E}{wantedprojection}{glossary}{Projection-valued spectral measure of a self-adjoint operator.}{projection-valued-measure}
\dkDeclareSymbol[Borel subset S of the real spectral axis. The spectral projection E(S) selects the part of a self-adjoint operator with spectral values in S; its range is a reducing subspace.]
  {SpectralSet}{S}{wantedspectrum}{glossary}{Borel spectral set used to select a spectral subspace.}{set}
% Generic Sylvester coefficients and right side, by source/target role.
\dkDeclareSymbol[Left coefficient B of the Sylvester map X -> B X - X C; B acts on the output space of X. For the sine bound, B = Lambda_1 on the unwanted-coordinate space G.]
  {SylvesterLeft}{B}{unwantedspectrum}{glossary}{Left coefficient of the Sylvester equation; corresponds to the unwanted spectral block.}{self-adjoint-operator}
\dkDeclareSymbol[Right coefficient C of the Sylvester map X -> B X - X C; C acts on the input space of X. For the sine bound, C = A_tr on trial space F. The historical equation also denotes this coefficient by A.]
  {SylvesterRight}{C}{trialspectrum}{glossary}{Right coefficient of the Sylvester equation; corresponds to the trial operator.}{self-adjoint-operator}
\dkDeclareSymbol[Right-hand side Y in B X - X C = Y, with the same input and output spaces as X. For the sine bound, Y = F_1* R: F -> G, the residual expressed in unwanted exact coordinates.]
  {SylvesterRHS}{Y}{resid}{glossary}{Right-hand side of the Sylvester equation, obtained by projecting the residual.}{linear-operator}
\dkDeclareSymbol{w}{w}{unwantedvector}{local}{Vector in an unwanted or complementary spectral subspace.}{vector}
\dkDeclareSymbol{f1}{f_1}{unwantedvector}{local}{Component or coordinate vector in the unwanted exact subspace.}{vector}
\dkDeclareSymbol{lambda1}{\lambda_1}{unwantedspectrum}{local}{Unwanted exact eigenvalue or complementary spectral value.}{scalar}
\dkDeclareSymbol{lambdaw}{\lambda_w}{unwantedspectrum}{local}{Eigenvalue associated with an unwanted exact vector.}{scalar}
\dkDeclareSymbol[Hilbert space G of exact complementary coordinates. F_1: G -> E embeds it onto the unwanted exact subspace; Lambda_1 acts on G and X = F_1* E_0 maps trial coordinates into G.]
  {Gunwanted}{G}{unwantedframe}{glossary}{Hilbert space of exact complementary coordinates.}{hilbert-space}
\dkDeclareSymbol{E1}{E_1}{unwantedframe}{local}{Isometric coordinates for the complementary trial subspace.}{isometric-embedding}
\dkDeclareSymbol{Vperp}{V^{\perp}}{unwanted}{local}{Orthogonal complement of the trial subspace in the local notation.}{linear-subspace}
% Angle and error geometry, including distinct rectangular and positive maps.
\dkDeclareSymbol[Individual principal angle, between 0 and pi/2. For orthonormal subspace frames X and Y, the singular values of X* Y are the principal cosines. In Figure 1, theta is the angle between the unit trial vector v and the exact wanted line.]
  {theta}{\theta}{sine}{glossary}{Principal angle between the subspaces under comparison.}{angle}
\dkDeclareSymbol[Full positive angle operator Theta = diag(Theta_0, Theta_1) for the source/target decompositions. Theta_0 compares the source with the target; Theta_1 compares their orthogonal complements. Its operator norm is the largest spectral angle. For positive c, the conditions Theta <= c and Theta < c mean operator norm at most c and strictly below c, respectively.]
  {Theta}{\Theta}{angleoperator}{glossary}{Full angle operator, the diagonal sum of the source and complementary angle blocks.}{angle-operator}
\dkDeclareSymbol[Directed principal-angle operator Theta_0 on source or trial coordinates. For a trial frame E_0 and exact wanted projection P_ex, cos(Theta_0) = sqrt(E_0* P_ex E_0), and sin(Theta_0) = sqrt(I_F - E_0* P_ex E_0). In the two-projection geometry it comes from the positive cosine compression P Q P restricted to range(P).]
  {Theta0}{\Theta_0}{angleoperator}{glossary}{Directed principal-angle operator acting on source or trial coordinates.}{angle-operator}
\dkDeclareSymbol[Complementary directed angle operator Theta_1, obtained from the projections I-P and I-Q by the same cosine-compression construction as Theta_0. It acts on the source complement, or its isometric coordinates. The full angle operator is diag(Theta_0, Theta_1).]
  {Theta1}{\Theta_1}{angleoperator}{glossary}{Directed angle operator comparing the orthogonal complements.}{angle-operator}
\dkDeclareSymbol[Right-angle source defect D_0 = range(P) intersect kernel(Q). These are source vectors orthogonal to the target. The full-angle conclusion in the Section 8 smallness criterion assumes dim(D_0) = dim(D_1); in Lean their identification is CrossedDefectsEquivalent.]
  {CrossedSourceDefect}{D_0}{fg}{glossary}{Source directions perpendicular to the target subspace.}{linear-subspace}
\dkDeclareSymbol[Right-angle target defect D_1 = kernel(P) intersect range(Q). These are target vectors orthogonal to the source. Equality of its Hilbert dimension with that of D_0 is assumed for the full-angle conclusion of Theorem 8.2.]
  {CrossedTargetDefect}{D_1}{fg}{glossary}{Target directions perpendicular to the source subspace.}{linear-subspace}
\dkDeclareSymbol[Positive complementary cosine block C_1 = cos(Theta_1). The reference complement is identified with the perturbed complement by the direct rotation D, so A_1 and Lambda_1 act on common coordinates. Theorem 8.1 compares A_1 - alpha I with C_1 (Lambda_1 - alpha I) C_1 as quadratic forms.]
  {ComplementCosineBlock}{C_1}{anglecos}{local}{Cosine of the complementary angle block in the spectral-repulsion inequality.}{self-adjoint-operator}
\dkDeclareSymbol[Scalar sine of a principal angle. For a unit principal trial vector v and exact wanted projection P_ex, sin(theta) = ||(I - P_ex)v||. In Figure 1 it is the length of the pink leg S(1).]
  {sintheta}{\sin\theta}{sine}{local}{Sine of an individual principal angle: the length of the complementary component of a unit principal trial vector.}{scalar}
\dkDeclareSymbol[Positive sine operator sin(Theta_0) = sqrt(I_F - E_0* P_ex E_0): F -> F. It equals |S| for S = (I - P_ex)E_0: F -> E. Its eigenvalues in finite dimensions are the principal sines; the sine theorem bounds its ideal norm by N(R)/delta.]
  {sinTheta0}{\sin\Theta_0}{sineoperator}{glossary}{Positive sine operator on trial coordinates; the modulus of the rectangular complementary error.}{angle-operator}
\dkDeclareSymbol[Rectangular complementary error S = (I - P_ex) E_0: F -> E. It records how much each trial vector leaks outside the exact wanted subspace. Its modulus |S| = sqrt(S* S) is sin(Theta_0): F -> F. In rank one S(1) = P_ex_perp v.]
  {S}{S}{sineblock}{glossary}{Rectangular complementary error: the trial frame projected onto the unwanted exact subspace.}{linear-operator}
\dkDeclareSymbol[Complementary-coordinate map X = F_1* E_0: F -> G, with S = F_1 X. It satisfies Lambda_1 X - X A_tr = F_1* R on domain(A_tr). Isometric multiplication gives N(X) = N(S); the separate modulus identity abs(S) = sin(Theta_0) gives N(S) = N(sin(Theta_0)), whenever these norms are finite.]
  {X}{X}{sylvester}{glossary}{Solution of the Sylvester equation; for the sine bound, the trial frame in unwanted exact coordinates.}{linear-operator}
\dkDeclareSymbol{alphaAngle}{\alpha}{anglemax}{local}{Largest principal angle.}{angle}
\dkDeclareSymbol[Scalar c = cos(theta) in the counterexample. In the general two-projection decomposition, the same role denotes a positive cosine operator C; the cosine and sine blocks commute and satisfy C squared + S squared = I.]
  {CosAngle}{c}{anglecos}{local}{Cosine of a principal angle or the corresponding positive cosine block.}{scalar-or-operator}
\dkDeclareSymbol[Scalar s = sin(theta) in the counterexample. In the general two-projection decomposition, the same role denotes a positive sine operator S; it commutes with the positive cosine block C and satisfies C squared + S squared = I.]
  {SinAngle}{s}{sineoperator}{local}{Sine of a principal angle or the corresponding positive sine block.}{scalar-or-operator}
\dkDeclareSymbol[Positive small-angle parameter in the local subspace comparison. It controls the allowed angular deviation; it is an angle parameter rather than a perturbation-operator norm.]
  {EpsilonAngle}{\varepsilon}{anglemax}{local}{Angle parameter controlling the maximum subspace deviation.}{angle}
\dkDeclareSymbol{PhiAngle}{\varphi}{angleoperator}{local}{Auxiliary angle or angle operator in the two-subspace comparison.}{angle}
\dkDeclareSymbol{Bblock}{B}{sineblock}{local}{Off-diagonal block in the two-subspace decomposition.}{linear-operator}
\dkDeclareSymbol{Cblock}{C}{sylvester}{local}{Complementary or Sylvester block in the two-subspace decomposition.}{linear-operator}
\dkDeclareSymbol{Ablock}{A}{anglecos}{local}{Positive cosine-related block in the two-subspace decomposition.}{linear-operator}
% Residual/error-displacement family.
\dkDeclareSymbol[Residual vector r = Ahat v - rho v = R(1) for the rank-one frame E_0(x) = x v. Figure 1 draws it from rho v to Ahat v. For unit v and the Rayleigh value rho = <v, Ahat v>, <v, r> = 0; a general trial scalar need not give orthogonality.]
  {r}{r}{residvector}{local}{Residual vector of an approximate eigenpair.}{vector}
\dkDeclareSymbol[Bounded residual R: F -> E, agreeing with Ahat E_0 - E_0 A_tr on domain(A_tr). The domain condition requires E_0(domain(A_tr)) inside domain(Ahat). R = 0 means that E_0 intertwines A_tr with Ahat; its norm measures the defect.]
  {R}{R}{resid}{glossary}{Bounded residual of the trial eigenproblem, agreeing with the operator difference on the trial domain.}{linear-operator}
\dkDeclareSymbol[Full displacement Delta_W = I - W on the ambient space. In the failed Ky Fan comparison W = W_(pi/4); Delta_W has singular values sqrt(2), sqrt(2), 0, 0. Multiplying by a principal-plane projection restricts its input to that plane.]
  {Kdisplacement}{\Delta_W}{displacement}{glossary}{Full displacement of a transporter from the identity.}{operator-valued-function}
\dkDeclareSymbol[Squared full displacement Dcal(W) = (I-W)* (I-W), a positive operator. Proposition 4.3 compares its unitarily invariant norm for a direct rotation D with that for every feasible transporter W.]
  {SquaredDisplacement}{\mathcal D}{displacement}{local}{Squared full displacement as a function of the transporter.}{operator-valued-function}
% Matching proof: the gradient and its spectral level are not a space or gap.
\dkDeclareSymbol[Skew-symmetric spectral gradient G of Phi_p. It commutes with W and represents first variations by d/dt Phi_p(W exp(tB)) at t = 0 equal to -tr(G B) for skew-symmetric B. Feasible stationarity gives J_src G J_src = -G.]
  {SpectralGradient}{G}{displacement}{local}{Skew-symmetric spectral gradient of the displacement power sum in the matching proof.}{linear-operator}
\dkDeclareSymbol[Positive spectral level g of the skew-symmetric gradient G: on this invariant subspace, -G squared = g squared I. Consequently G/g squares to -I and supplies the complex structure used in the real matching proof.]
  {GradientLevel}{g}{displacement}{local}{Positive gradient level; the gradient squared is minus this level squared times the identity.}{scalar}
\dkDeclareSymbol[Complex structure I_g = G/g on the positive gradient-level space V_g. It is a real isometry with square -I. Multiplication by a+ib means a times identity plus b times I_g; source reflection J_src reverses I_g and acts as conjugation.]
  {GradientComplexStructure}{\mathcal I_g}{displacement}{local}{Complex multiplication constructed from a positive gradient level.}{linear-operator}
\dkDeclareSymbol[Real skew-symmetric matrix B, with B transpose = -B, giving the orthogonal variation W exp(tB). If B commutes with the source projection P, this curve preserves the constraint W(U) = V.]
  {SkewDirection}{B}{ambient}{local}{Skew-symmetric direction for a first variation of an orthogonal transporter.}{linear-operator}
% Spectral separation family: gap and original gap are NOT identical.
\dkDeclareSymbol[Positive spectral separation in the stated configuration. In the sine theorem it separates the trial operator from the exact complementary block; Theorem 8.1 uses it for the ordered reference gap, and Theorem 8.2 for the perturbed central-band separation. Figure 1 uses delta = abs(rho - 0.7).]
  {delta}{\delta}{gap}{glossary}{Positive spectral separation in the stated gap configuration.}{scalar}
\dkDeclareSymbol{g}{g}{gaporiginal}{local}{Original spectral gap of the reference operator; distinct from the certified sine-theta separation.}{scalar}
\dkDeclareSymbol{Delta}{\Delta}{gapstat}{local}{Population or statistical spectral gap.}{scalar}
\dkDeclareSymbol[Spectral endpoint alpha. In the interval/exterior configuration, one spectrum lies between beta and alpha, the other at or below beta - delta or at or above alpha + delta. Theorem 8.2 places the selected perturbed block in the central interval. In Theorem 8.1, alpha bounds the lower reference block and alpha + delta bounds the upper block from below.]
  {alphaEndpoint}{\alpha}{gapinterval}{glossary}{Upper endpoint of the central band, or lower-block cut in the ordered gap.}{scalar}
\dkDeclareSymbol[Lower endpoint beta of the finite interval from beta to alpha in the gap hypothesis. One spectrum lies in that interval; the other lies at or below beta - delta or at or above alpha + delta.]
  {betaEndpoint}{\beta}{gapinterval}{glossary}{Lower endpoint of the finite interval in the interval/exterior gap hypothesis.}{scalar}
% Spectral/projector and auxiliary matrix objects, context-scoped.
\dkDeclareSymbol{P}{P}{fg}{local}{Orthogonal projection; its source or exact-target role is specified locally.}{orthogonal-projection}% Neutral historical Manim key; prefer Pexact/Psource.
\dkDeclareSymbol{Q}{Q}{fg}{local}{Orthogonal projection; its role is specified locally.}{orthogonal-projection}% Neutral historical Manim key; prefer Qtarget.
\dkDeclareSymbol[Two-projection intertwiner T = QP + (I - Q)(I - P), satisfying TP = QT and T* T = I - (P - Q) squared. For ||P - Q|| < 1 it is invertible; its unitary polar factor T |T| inverse is the canonical direct rotation.]
  {IntertwinerT}{T}{intertwiner}{glossary}{Two-projection intertwiner; its unitary polar factor is the acute direct rotation.}{linear-operator}
\dkDeclareSymbol[Planar pi/2-rotation matrix J_0 = ((0, -1), (1, 0)), with J_0 squared = -I_2 and J_0 transpose = -J_0. It defines the rotation C = cos(2 theta) I_2 + sin(2 theta) J_0 used to construct the real competitor.]
  {Jturn}{J_0}{turn}{glossary}{Planar rotation through $\pi/2$, with square minus the identity.}{matrix}
\dkDeclareSymbol[Partial isometry V_pol in the polar decomposition S = V_pol |S|, where |S| = sqrt(S* S). It is isometric on the orthogonal complement of kernel(S) and zero on kernel(S). For an invertible operator its polar factor is unitary.]
  {PolarU}{V_{\mathrm{pol}}}{intertwiner}{glossary}{Partial isometry in the polar decomposition; unitary when the operator is invertible.}{partial-isometry}
\dkDeclareSymbol[Planar rotation C = cos(2 theta) I_2 + sin(2 theta) J_0 through twice the principal angle. The competitor W_theta is orthogonally conjugate to C direct-sum I_2, so it rotates one two-plane and fixes the other.]
  {Cturn}{C}{turn}{glossary}{Planar rotation through twice the principal angle in the counterexample.}{matrix}
\dkDeclareSymbol[Orthogonal mixing matrix H_mix = (1/sqrt(2)) times the block matrix ((I_2, I_2), (I_2, -I_2)). It constructs W_theta = H_mix (C direct-sum I_2) H_mix transpose.]
  {Hmix}{H_{\mathrm{mix}}}{mix}{glossary}{Orthogonal mixing matrix in the counterexample.}{matrix}
\dkDeclareSymbol[Source reflection J_src = 2P - I: it acts as +I on the source subspace and -I on its orthogonal complement, so its square is I. In the matching proof stationarity gives J_src G J_src = -G for the skew-symmetric spectral gradient G.]
  {Jreflect}{J_{\mathrm{src}}}{trialprojection}{glossary}{Reflection equal to twice the source projection minus the identity.}{linear-operator}
\dkDeclareSymbol[Skew angle generator J_ang for a direct rotation D = exp(J_ang Theta). On each principal plane it is the oriented rotation through pi/2; it vanishes on the common fixed directions. It commutes with the full angle operator Theta.]
  {AngleGenerator}{J_{\mathrm{ang}}}{turn}{glossary}{Skew angle generator in the direct-rotation result summaries.}{linear-operator}
\dkDeclareSymbol{Dblock}{D}{mix}{local}{Auxiliary diagonal or block matrix in the counterexample.}{matrix}
\dkDeclareSymbol[Orthogonal projection onto the first principal plane span(e_0, e_3) in the four-dimensional counterexample. For Delta_W = I - W_(pi/4), the Ky Fan 2-norm of Delta_W Omega_1 equals 2.]
  {OmegaOne}{\Omega_1}{turn}{glossary}{First principal-plane projection in the failed Ky Fan comparison.}{orthogonal-projection}
\dkDeclareSymbol[Orthogonal projection onto the second principal plane span(e_1, e_2) in the four-dimensional counterexample. For Delta_W = I - W_(pi/4), the Ky Fan 2-norm of Delta_W Omega_2 equals 2, although the trace norm of Delta_W is only 2 sqrt(2).]
  {OmegaTwo}{\Omega_2}{intertwiner}{glossary}{Second principal-plane projection in the failed Ky Fan comparison.}{orthogonal-projection}
\dkDeclareSymbol{Dplanar}{D}{wantedrotation}{local}{Planar direct-rotation block.}{unitary-operator}
\dkDeclareSymbol{Hreflection}{H}{trialcompetitor}{local}{Planar reflection carrying the source line to the target.}{linear-operator}
% Unitarily invariant gauge family; N, M are NOT a norm versus its input.
\dkDeclareSymbol[Unitarily invariant norm N, or a symmetric operator-ideal norm on its domain of finiteness. It satisfies N(U T V) = N(T) for unitaries and N(B T) <= ||B||op N(T). In the sine theorem it is normalized by N(u v*) = ||u|| ||v|| and has Ky Fan dominance where both norms are finite. Both R and sin(Theta_0) must belong to its ideal.]
  {NormN}{N}{norm}{glossary}{Unitarily invariant norm on its domain of finiteness; in the sine theorem, normalized by $N(uv^*)=\|u\|\|v\|$.}{norm}
\dkDeclareSymbol[Sine-theta separation G_delta(B,C): finite interval/exterior separation in either spectral order, or ordered separation B <= a I and C >= (a+delta) I, or the reverse. The ordered alternatives allow both spectra to be unbounded.]
  {SineGap}{\mathcal G_\delta}{gap}{local}{Interval/exterior separation or either ordered separation, as in the sine-theta theorem.}{predicate}
\dkDeclareSymbol[Auxiliary unitarily invariant norm M witnessing the Q-norm property N(T) squared = M(T* T). For the Frobenius norm N, one can take M to be the trace norm.]
  {NormM}{M}{normaux}{local}{Auxiliary unitarily invariant norm used to define a Q-norm through the squared modulus.}{norm}
\dkDeclareSymbol[Q-norm property: a unitarily invariant norm N is a Q-norm when N(T) squared = M(T* T) for some unitarily invariant norm M. The Frobenius norm is an example, with M the trace norm; the repaired full-displacement theorem uses this norm class.]
  {NormQ}{Q}{normaux}{glossary}{Q-norm: its square is a unitarily invariant norm of the squared modulus.}{norm}
\dkDeclareSymbol[Singleton trace cost s_i = min(4 sin(theta_i/2), 2) for angle label i. It transports one real principal plane independently, choosing a rotation of cost 4 sin(theta_i/2) or a reflection of cost 2.]
  {SingletonCost}{s}{norm}{glossary}{Minimum trace displacement for one real principal plane: choose a rotation or a reflection.}{scalar-cost}
\dkDeclareSymbol[Paired trace cost e_ij = 4 sin((theta_i + theta_j)/2) for distinct angle labels i,j. It is attained by coupling their principal planes in a real four-dimensional block. The edge weight s_i + s_j - e_ij is the difference between the cost of two singletons and this paired cost.]
  {PairCost}{e}{norm}{glossary}{Trace displacement cost of coupling two principal planes in a four-dimensional block.}{scalar-cost}
\dkDeclareSymbol[Scalar phase cost f_p(t) = (2 sin(t/2)) to the power p, for 0 <= t <= pi. A real planar rotation through t contributes two copies, so its displacement power cost is 2 f_p(t).]
  {MatchingPhaseCost}{f_p}{norm}{local}{Displacement power cost for one scalar phase.}{scalar-cost}
\dkDeclareSymbol[Singleton power cost a_i(p) = min(2 f_p(theta_i), 2 to the power p). These are the rotation and reflection choices for one principal plane. At p=1 it equals the trace-cost weight s_i.]
  {SingletonPowerCost}{a}{norm}{local}{Minimum displacement power cost for one angle.}{scalar-cost}
\dkDeclareSymbol[Paired power cost b_ij(p) = min over d <= tau <= h of 2(f_p(h+tau) + f_p(h-tau)), where h=(theta_i+theta_j)/2 and d=abs(theta_i-theta_j)/2. This is the positive-determinant pair minimum; unrestricted block optimization also compares two singletons. At p=1 it equals e_ij.]
  {PairPowerCost}{b}{norm}{local}{Positive-determinant displacement power cost for a pair of angles.}{scalar-cost}
\dkDeclareSymbol[Matching M on the finite general graph of positive principal-angle labels. Each label belongs to at most one edge; unmatched labels use singleton transporters. Minimizing paired and singleton costs is equivalent to maximum-weight matching with edge weights g_ij = s_i + s_j - e_ij. The empty matching is allowed.]
  {Matching}{M}{normaux}{glossary}{Matching of positive principal-angle labels: each label is used in at most one pair.}{matching}
\dkDeclareSymbol[Displacement power sum Phi_p(W) = sum_j s_j(I - W)^p. The matching proof first minimizes it for 1 < p < 2, then lets p decrease to 1, where it becomes the trace norm ||I - W||_1.]
  {PowerCost}{\Phi}{norm}{local}{Sum of the p-th powers of displacement singular values.}{scalar-cost}
% Auxiliary matching-proof notation: neutral ink, with the same hover registry.
\dkDeclareSymbol[Off-diagonal gradient block G_01 mapping the source complement K perpendicular into K. Stationarity makes the diagonal blocks of G zero, so G = ((0,G_01),(-G_01 transpose,0)) in the source/complement decomposition.]
  {GradientOffDiagonal}{G_{01}}{fg}{local}{Complement-to-source block of the skew-symmetric gradient.}{linear-operator}
\dkDeclareSymbol[Orthogonal group O(n): real n by n matrices W satisfying W transpose times W = I. Its elements preserve Euclidean inner products and lengths. The notation O(2) refers to the same group in dimension two.]
  {OrthogonalGroup}{O(n)}{fg}{local}{Group of real orthogonal matrices.}{group}
\dkDeclareSymbol[Feasible transporters F = {W in O(n) : W(K) = L}. Equivalently W P_K W transpose = P_L. This nonempty compact set is the optimization domain throughout the matching proof.]
  {FeasibleTransport}{\mathcal F}{fg}{local}{Orthogonal maps carrying the source subspace onto the target.}{set}
\dkDeclareSymbol[Singular value s_j(T), the nonnegative square root of an eigenvalue of T transpose times T, counted with multiplicity. Summing s_j(I-W) gives the trace displacement; summing its p-th powers gives Phi_p(W). These singular values differ from the singleton weights s_i.]
  {TransportSingularValues}{\symglyph{SingularValueBase}_j}{fg}{local}{A displacement singular value in the matching proof.}{scalar}
\dkDeclareSymbol[Scalar singular-value magnitude s between zero and two for the displacement I-W of an orthogonal matrix. Uniform convergence of s to the power p to s on this interval permits the trace-norm limit. This variable differs from the rotation parameter s used earlier in the proof.]
  {DisplacementSize}{s}{fg}{local}{Scalar displacement singular-value magnitude in the uniform limit.}{scalar}
\dkDeclareSymbol[All-singleton cost C_0 = sum_i s_i, before any two angle labels are paired. A matching has cost C_0 minus the sum of its selected edge weights.]
  {SingletonBaseline}{C_0}{fg}{local}{Trace displacement when every angle uses its singleton transporter.}{scalar-cost}
\dkDeclareSymbol[Matching power cost C_p(M): sum of b_ij(p) over selected pairs plus a_i(p) over unmatched labels. Its minimum equals the minimum of Phi_p over feasible transporters. At p=1 it is C(M).]
  {MatchingPowerCost}{C_p}{fg}{local}{Total paired and singleton power cost of a matching.}{scalar-cost}
\dkDeclareSymbol[Invariant orthogonal block V_a preserved by the source projection, target projection, and minimizing transporter. The displacement cost adds over these blocks, and their principal-angle lists combine into the original list.]
  {MatchingBlock}{V_a}{fg}{local}{Simultaneous invariant block in the matching proof.}{linear-subspace}
\dkDeclareSymbol[Gradient-level space V_g: the eigenspace of -G squared for a positive eigenvalue g squared. It is preserved by G, W, and the source reflection J_src. Multiplication by i on it is G/g.]
  {GradientLevelSpace}{V_g}{fg}{local}{Positive gradient-level eigenspace.}{linear-subspace}
\dkDeclareSymbol[Eigenphase varphi of W on V_g viewed as a complex space with multiplication by i = G/g. It lies between zero and pi and satisfies f_p prime(varphi) = g. There are at most two such phases.]
  {GradientEigenphase}{\varphi}{fg}{local}{One of the at most two phases at a positive gradient level.}{angle}
\dkDeclareSymbol[Orthogonal projection E onto the first phase space of W on V_g. Here W = exp(i varphi_1) E + exp(i varphi_2) (I-E). In coordinates where the source reflection is conjugation, E = X + iY.]
  {GradientPhaseProjection}{E}{fg}{local}{Projection onto the first phase space in a gradient-level space.}{orthogonal-projection}
\dkDeclareSymbol[Real part X of the phase projection E = X + iY. It is a real symmetric matrix with spectrum in the closed interval from zero to one, satisfying X squared - Y squared = X and XY + YX = Y. This X differs from the source-coordinate and Sylvester matrices elsewhere in the paper.]
  {ProjectionRealPart}{X}{fg}{local}{Real symmetric part of the phase projection.}{linear-operator}
\dkDeclareSymbol[Imaginary part Y of the phase projection E = X + iY. It is a real skew-symmetric matrix. It maps the lambda eigenspace of X to the 1-lambda eigenspace and constructs the small invariant planes.]
  {ProjectionImaginaryPart}{Y}{fg}{local}{Real skew-symmetric imaginary part of the phase projection.}{linear-operator}
\dkDeclareSymbol[Eigenvalue lambda in the closed interval from zero to one of the real part X of the phase projection. It is local to the block construction and differs from the original operator's spectral parameter.]
  {ProjectionEigenvalue}{\lambda}{fg}{local}{Eigenvalue of the real part of the phase projection.}{scalar}
\dkDeclareSymbol[Plane coupling eta = sqrt(lambda (1-lambda)) for 0 < lambda < 1. If u is a unit lambda eigenvector of X, then v = Yu/eta completes its invariant real plane.]
  {ProjectionCoupling}{\eta}{fg}{local}{Normalization used to construct an invariant real plane.}{scalar}
\dkDeclareSymbol[Pair mean angle h = (theta_i + theta_j)/2. Together with d = (theta_i - theta_j)/2 after ordering theta_i >= theta_j, it parametrizes the two-angle block cost.]
  {PairMeanAngle}{h}{fg}{local}{Mean of the two ordered principal angles.}{angle}
\dkDeclareSymbol[Pair half-difference d = (theta_i - theta_j)/2 after ordering theta_i >= theta_j. It satisfies 0 <= d <= h and h+d <= pi/2.]
  {PairHalfDifference}{d}{fg}{local}{Half the difference of the two ordered principal angles.}{angle}
\dkDeclareSymbol[Scalar pair parameter tau in the closed interval from d to h. The paired power cost is the minimum of 2(f_p(h+tau) + f_p(h-tau)) over this interval. At p=1 a minimum is attained at tau=h.]
  {PairPhaseParameter}{\tau}{fg}{local}{Scalar parameter in the two-angle power-cost minimum.}{scalar}
\dkDeclareSymbol[Attaining rotation angle gamma = arccos(cos(tau)/cos(d)). Choosing U_0 = U_1 = R_gamma realizes phase parameters x=h and y=tau in the four-dimensional block.]
  {PairAttainingAngle}{\gamma}{fg}{local}{Rotation angle realizing a chosen scalar pair parameter.}{angle}
\dkDeclareSymbol[Four-dimensional direct rotation D = ((C,-S),(S,C)) in source and complementary principal coordinates, with C the diagonal principal cosines and S the diagonal principal sines. Every transporter on this block is D diag(U_0,U_1).]
  {PairDirectRotation}{D}{fg}{local}{Direct rotation on the two principal planes.}{unitary-operator}
\dkDeclareSymbol[Diagonal cosine block C = diag(cos(theta_i), cos(theta_j)) in the four-dimensional direct rotation D. This matrix differs from the scalar displacement cost C(M).]
  {PairCosineBlock}{C}{fg}{local}{Diagonal principal-cosine block.}{matrix}
\dkDeclareSymbol[Diagonal sine block S = diag(sin(theta_i), sin(theta_j)) in the four-dimensional direct rotation D. This matrix is local to the matching proof.]
  {PairSineBlock}{S}{fg}{local}{Diagonal principal-sine block.}{matrix}
\dkDeclareSymbol[Free planar orthogonal matrices U_0 and U_1 acting on the two-dimensional source and its two-dimensional complement, respectively. Multiplying D by diag(U_0,U_1) parametrizes every transporter on this four-dimensional block.]
  {PairFreeOrthogonal}{U}{fg}{local}{Free source or complementary planar orthogonal factor.}{unitary-operator}
\dkDeclareSymbol[Phase parameter x in the closed interval from h to pi-h for a positive-determinant four-dimensional transporter. Its eigenvalues are exp(plus or minus i(x+y)) and exp(plus or minus i(x-y)).]
  {PairPhaseX}{x}{fg}{local}{First phase parameter of a four-dimensional transporter.}{angle}
\dkDeclareSymbol[Phase parameter y in the closed interval from d to pi-d for a positive-determinant four-dimensional transporter. Together with x it gives rotational phases x+y and x-y, folded into the closed interval from zero to pi for the displacement cost.]
  {PairPhaseY}{y}{fg}{local}{Second phase parameter of a four-dimensional transporter.}{angle}
\dkDeclareSymbol[Scalar a = cos(x) in the four-dimensional phase calculation. Together with b = cos(y), it expresses the trace as 4ab and determines the reciprocal characteristic polynomial.]
  {PairPhaseCosineX}{a}{fg}{local}{Cosine of the first phase parameter in the trace calculation.}{scalar}
\dkDeclareSymbol[Scalar b = cos(y) in the four-dimensional phase calculation, paired with a = cos(x). These scalars are local to the characteristic-polynomial calculation.]
  {PairPhaseCosineY}{b}{fg}{local}{Cosine of the second phase parameter in the trace calculation.}{scalar}
\dkDeclareSymbol[Rotation parameter r: the two free orthogonal factors use planar rotation angles r+s and r-s. It appears in cos(y) = cos(d) cos(r) for two rotations.]
  {PairRotationR}{r}{fg}{local}{Half the sum of the two free planar rotation angles.}{angle}
\dkDeclareSymbol[Rotation parameter s: the two free orthogonal factors use planar rotation angles r+s and r-s. It appears in cos(x) = cos(h) cos(s) for two rotations. This scalar differs from the singular values and singleton costs.]
  {PairRotationS}{s}{fg}{local}{Half the difference of the two free planar rotation angles.}{angle}
\dkDeclareSymbol[Planar reflection J_2 = diag(1,-1). Every two-dimensional orthogonal matrix of determinant -1 can be written as a planar rotation times J_2.]
  {PairPlanarReflection}{J_2}{fg}{local}{Fixed planar reflection used to parametrize determinant -1 factors.}{unitary-operator}
\dkDeclareSymbol[Real planar rotation R_t = ((cos(t),-sin(t)),(sin(t),cos(t))). Its two displacement singular values are both 2 times abs(sin(t/2)). In the matching proof a rotation through t in the closed interval from zero to pi contributes 2 f_p(t).]
  {MatchingPlanarRotation}{\symglyph{PlanarRotationBase}_t}{fg}{local}{Planar rotation through the indicated angle.}{unitary-operator}
% Statistical application / operator provenance.
\dkDeclareSymbol{Sigma}{\Sigma}{old}{local}{Population covariance matrix.}{matrix}
\dkDeclareSymbol{SigmaHat}{\widehat\Sigma}{current}{local}{Estimated or perturbed covariance matrix.}{matrix}
\dkDeclareSymbol{YWSV}{V}{wanted}{local}{Orthonormal coordinates for the population eigenspace in the statistical application.}{isometric-embedding}
\dkDeclareSymbol{YWSVhat}{\widehat V}{trialframe}{local}{Estimated eigenspace frame in the statistical application.}{isometric-embedding}

% Figure-specific objects retain their own meanings when bare letters overlap.
\dkDeclareSymbol[Signed rotation angle phi of the displayed map on its invariant real two-plane. At theta = pi/4 the direct rotation turns its two principal planes by opposite signed angles of magnitude pi/4; the competitor turns its moving plane by pi/2 and fixes its other plane.]
  {PlaneRotationAngle}{\varphi}{sine}{local}{Signed rotation angle of the orthogonal map on the displayed invariant plane.}{angle}
\dkDeclareSymbol[Unit vector q in the displayed invariant plane. The map T is the direct rotation D or the competitor W in that panel; Tq is the image and (I - T)q = q - Tq is the displacement.]
  {FigureInput}{q}{ambient}{local}{Unit input to the orthogonal map shown in this invariant-plane panel.}{vector}
\dkDeclareSymbol[Image Tq of the panel unit input q under T = D or W. Both q and Tq lie on the unit circle; the green displacement arrow joins Tq to q, so it represents q - Tq.]
  {FigureImage}{Tq}{wantedvector}{local}{Image of the unit input under the orthogonal map shown in this panel.}{vector}
\dkDeclareSymbol[Displacement (I - T)q = q - Tq, drawn from the image Tq to the unit input q. On a plane rotated through phi its length is 2 |sin(phi/2)|; it vanishes on a fixed plane.]
  {DisplacementVector}{(I-T)q}{displacement}{local}{Difference between a unit input and its image under the displayed orthogonal map.}{vector}
\dkDeclareSymbol[First principal plane Pi_1 = span(e_0, e_3) for the source-target pair in the real four-dimensional counterexample. At theta = pi/4 the direct rotation acts on this plane by a quarter of pi, with displacement singular value sqrt(2 - sqrt(2)).]
  {CounterPrincipalOne}{\Pi_1}{fg}{glossary}{First principal plane in the real counterexample: the span of e0 and e3.}{linear-subspace}
\dkDeclareSymbol[Second principal plane Pi_2 = span(e_1, e_2) for the source-target pair in the real four-dimensional counterexample. The direct rotation acts by the same angle magnitude as on Pi_1, with the opposite sign in the displayed plane coordinates.]
  {CounterPrincipalTwo}{\Pi_2}{fg}{glossary}{Second principal plane in the real counterexample: the span of e1 and e2.}{linear-subspace}
\dkDeclareSymbol[Moving plane M = range(W - I) = span(e_0 + e_2, e_1 + e_3) for the competitor. W rotates it by 2 theta; at theta = pi/4 the two displacement singular values contributed by this plane are both sqrt(2).]
  {CounterMovingPlane}{M}{ambient}{glossary}{Moving plane of the competing transporter: the range of W minus the identity.}{linear-subspace}
\dkDeclareSymbol[Fixed plane Z = kernel(W - I) = span(e_0 - e_2, e_1 - e_3) for the competitor. W acts as the identity here, so this plane contributes two zero displacement singular values.]
  {CounterFixedPlane}{Z}{ambient}{glossary}{Fixed plane of the competing transporter: the kernel of W minus the identity.}{linear-subspace}
\dkDeclareSymbol[Direct-rotation displacement singular value a = 2 sin(pi/8) = sqrt(2 - sqrt(2)) at theta = pi/4. Each of its two principal planes contributes two copies, giving trace displacement 4a.]
  {CounterChordA}{a}{displacement}{glossary}{Direct-rotation displacement singular value in the counterexample: sqrt(2 - sqrt(2)).}{scalar}
\dkDeclareSymbol[Competitor displacement singular value b = 2 sin(pi/4) = sqrt(2) at theta = pi/4. Its moving plane contributes two copies; the fixed plane contributes two zeros, giving trace displacement 2b.]
  {CounterChordB}{b}{displacement}{glossary}{Competing-transporter displacement singular value in the counterexample: square root of 2.}{scalar}
\dkDeclareSymbol[Singular value s_j(I - T) of the full ambient displacement, ordered from largest to smallest. Each rotating real two-plane contributes two values 2 |sin(phi/2)|; fixed directions contribute zeros. Their sum is the trace cost.]
  {DisplacementSingular}{\symglyph{SingularValueBase}_j}{displacement}{local}{An individual singular value of the full displacement from the identity.}{scalar}
\dkDeclareSymbol[Edge weight g_ij = s_i + s_j - e_ij, the difference between the cost of two singletons and the paired cost. A matching has total cost C_0 minus the sum of its edge weights, with C_0 = sum_i s_i. Maximizing this sum gives the minimum transport cost; edges of negative weight may be omitted.]
  {MatchingEdgeWeight}{g_{ij}}{current}{glossary}{Edge weight: the cost of two singletons minus their paired cost.}{scalar-weight}
\dkDeclareSymbol[Total trace-norm displacement cost. For a transporter W, C(W) = ||I - W||_1, the sum of its displacement singular values. For a matching M, C(M) is the sum of its paired and unmatched singleton costs.]
  {TraceCost}{C}{norm}{glossary}{Trace-norm displacement cost of the indicated orthogonal transporter or matching.}{scalar-cost}

% Spectral values have their own key, distinct from an operator block.
\dkDeclareSymbol[Spectral parameter lambda of a self-adjoint operator. In a finite-dimensional eigenvalue list lambda_j(A) it denotes an individual eigenvalue; in A = integral lambda dE(lambda), it is the real integration variable.]{SpectralValue}{\lambda}{wantedspectrum}{local}{Real spectral parameter or individual eigenvalue of a self-adjoint operator.}{scalar}



% Conventional mathematical notation. Hover text applies even when ink is plain.
% These are local (non-glossary) unless promoted explicitly by the author.
% IdentityTwo and IdentityFour derive their I from Identity. Singular-value
% registrations derive s_j from SingularValueBase; SingletonCost has its own s.
\dkDeclareSymbol[Identity operator on the vector or Hilbert space in context. Its dimension is inferred from the expression and need not be finite.]{Identity}{I}{plain}{local}{Identity operator on the space in context.}{identity-operator}
\dkDeclareSymbol[Two-by-two identity matrix, acting on a two-dimensional coordinate space.]{IdentityTwo}{\symglyph{Identity}_2}{plain}{local}{Two-by-two identity matrix.}{identity-operator}
\dkDeclareSymbol[Four-by-four identity matrix, acting on a four-dimensional coordinate space.]{IdentityFour}{\symglyph{Identity}_4}{plain}{local}{Four-by-four identity matrix.}{identity-operator}
\dkDeclareSymbol[Two-by-two zero matrix. Unlike the bold zero vector, this is a linear operator/matrix with two rows and two columns.]{ZeroMatrixTwo}{0_2}{plain}{local}{Two-by-two zero matrix.}{matrix}
\dkDeclareSymbol[The constant pi, the ratio of the circumference of a circle to its diameter. Angles in this article use radians.]{PiConstant}{\pi}{plain}{local}{The constant pi (angles are in radians).}{constant}
\dkDeclareSymbol[The field of real numbers.]{RealField}{\mathbb R}{plain}{local}{The real numbers.}{field}
\dkDeclareSymbol[The field of complex numbers.]{ComplexField}{\mathbb C}{plain}{local}{The complex numbers.}{field}
\dkDeclareSymbol[Generic Hilbert space used in the two-projection discussion; distinct from the paper's ambient space symbol E in the sine-theta theorem.]{GeneralHilbertSpace}{\mathcal H}{plain}{local}{Generic Hilbert space in a local discussion.}{hilbert-space}
\dkDeclareSymbol[The spectrum sigma(T) of an operator T: the set of complex scalars for which T minus the scalar identity is not boundedly invertible. For self-adjoint T the spectrum is real.]{SpectrumSymbol}{\sigma}{plain}{local}{Spectrum of an operator.}{operator-notation}
\dkDeclareSymbol[Base letter for singular values, presently s. It is distinct from the base s for matching singleton costs.]{SingularValueBase}{s}{plain}{local}{Base letter for singular values.}{notation-base}
\dkDeclareSymbol[Ordered singular value s_j(T), equal to the j-th eigenvalue of |T| for a finite matrix; distinguish this from the matching singleton cost s_j.]{SingularValueS}{\symglyph{SingularValueBase}_j}{plain}{glossary}{The j-th singular value, under the currently printed s_j convention.}{scalar}
\dkDeclareSymbol[Alternate indexed singular-value notation sigma_j(T). The current manuscript also uses s_j(T); this registered alias is available for a later deliberate notation decision.]{SingularValueSigma}{\sigma_j}{plain}{local}{Indexed singular value under the sigma_j convention.}{scalar}
\dkDeclareSymbol[Operator norm ||T||_op, the supremum of ||Tx|| over unit vectors x.]{OperatorNormNotation}{\|\cdot\|_{\mathrm{op}}}{plain}{glossary}{Operator norm.}{norm}
\dkDeclareSymbol[Frobenius norm ||T||_F, the square root of the sum of squared matrix entries, equal to the finite-dimensional Hilbert--Schmidt norm.]{FrobeniusNormNotation}{\|\cdot\|_F}{plain}{glossary}{Frobenius norm.}{norm}
\dkDeclareSymbol[Trace norm ||T||_1, the sum of singular values when finite.]{TraceNormNotation}{\|\cdot\|_1}{plain}{glossary}{Trace (nuclear) norm.}{norm}
\dkDeclareSymbol[Hilbert--Schmidt norm ||T||_HS, the square root of the sum of squared singular values when finite.]{HilbertSchmidtNormNotation}{\|\cdot\|_{\mathrm{HS}}}{plain}{local}{Hilbert--Schmidt norm.}{norm}
\dkDeclareSymbol[Ky Fan j-norm, the sum of the largest j singular values.]{KyFanNormNotation}{\|\cdot\|_{(j)}}{plain}{glossary}{Ky Fan j-norm.}{norm}
\dkDeclareSymbol[Interval/exterior spectral separation predicate I_delta(B,C), specifying a central closed interval separated by delta from the other operator's spectrum.]{IntervalSeparation}{\mathcal I}{plain}{local}{Interval/exterior separation predicate.}{predicate}
\dkDeclareSymbol[Ordered spectral separation predicate O_delta(B,C), requiring one self-adjoint operator to lie below a spectral threshold and the other above it by at least delta.]{OrderedSeparation}{\mathcal O}{plain}{local}{Ordered spectral separation predicate.}{predicate}
\dkDeclareSymbol[Positive parameter epsilon controlling a small perturbation in the two-dimensional eigenvalue example.]{EigenvalueExampleEpsilon}{\varepsilon}{plain}{local}{Small positive parameter in the eigenvalue-example matrices.}{scalar}
\dkDeclareSymbol[Positive lower bound epsilon for the trial-coordinate map E_0, as in E_0* E_0 >= epsilon squared times the identity.]{FrameLowerBoundEpsilon}{\varepsilon}{plain}{local}{Lower-bound parameter for the trial-coordinate embedding.}{scalar}
\dkDeclareSymbol[Rotation angle vartheta of a planar orthogonal block in the matching-proof first-variation calculation.]{MatchingPlanarAngle}{\vartheta}{plain}{local}{Planar rotation angle in the matching-proof calculation.}{angle}


% Ordinary named mathematical operators already used throughout the manuscript.
% Keeping these here makes their spelling a single editorial choice.
\dkDeclareSymbol[Range of a linear operator: the set of its output vectors.]{RangeOperator}{\operatorname{ran}}{plain}{local}{Range (image) of an operator.}{operator-notation}
\dkDeclareSymbol[Kernel (null space) of a linear operator: inputs mapped to zero.]{KernelOperator}{\operatorname{ker}}{plain}{local}{Kernel (null space) of an operator.}{operator-notation}
\dkDeclareSymbol[Domain of an operator: the vectors on which the operator is defined, important for unbounded operators.]{DomainOperator}{\operatorname{dom}}{plain}{local}{Domain of an operator.}{operator-notation}
\DeclareRobustCommand{\DKRange}{\sym{RangeOperator}}
\DeclareRobustCommand{\DKKernel}{\sym{KernelOperator}}
\DeclareRobustCommand{\DKDom}{\sym{DomainOperator}}

% Norm-formatting macros: the operand changes by occurrence but the notation
% lives here. The outer hover describes the norm and does not add ink color.
\DeclareRobustCommand{\dkOperatorNorm}[1]{\like{OperatorNormNotation}{\|#1\|_{\mathrm{op}}}}
\DeclareRobustCommand{\dkTraceNorm}[1]{\like{TraceNormNotation}{\|#1\|_1}}
\DeclareRobustCommand{\dkFrobeniusNorm}[1]{\like{FrobeniusNormNotation}{\|#1\|_F}}
\DeclareRobustCommand{\dkHilbertSchmidtNorm}[1]{\like{HilbertSchmidtNormNotation}{\|#1\|_{\mathrm{HS}}}}
\DeclareRobustCommand{\dkKyFanNorm}[2]{\like{KyFanNormNotation}{\|#1\|_{(#2)}}}
\DeclareRobustCommand{\dkSchattenNorm}[2]{\like{SchattenNormNotation}{\|#1\|_{S_{#2}}}}

\dkDeclareSymbol[Schatten p-norm ||T||_{S_p}, obtained from the l^p-norm of the singular values and finite for operators in the Schatten p-class.]{SchattenNormNotation}{\|\cdot\|_{S_p}}{plain}{local}{Schatten p-norm.}{norm}


\dkDeclareSymbol[Positive modulus of an operator T: |T| = (T* T)^(1/2).]{OperatorModulusNotation}{|\cdot|}{plain}{glossary}{Positive modulus of an operator.}{operator-notation}
\dkDeclareSymbol[Linear span of the indicated vectors, i.e. the smallest linear subspace that contains them.]{SpanOperator}{\operatorname{span}}{plain}{local}{Linear span of vectors.}{operator-notation}
\dkDeclareSymbol[Diagonal or block-diagonal matrix/operator with the displayed entries or blocks.]{DiagonalOperator}{\operatorname{diag}}{plain}{local}{Diagonal or block-diagonal operator.}{linear-operator}
\dkDeclareSymbol[Trace of a trace-class operator, or ordinary matrix trace in finite dimensions.]{TraceOperator}{\operatorname{tr}}{plain}{local}{Trace of an operator.}{operator-notation}
\dkDeclareSymbol[Dimension of a vector space or subspace.]{DimensionOperator}{\dim}{plain}{local}{Dimension of a vector space.}{operator-notation}
\dkDeclareSymbol[Rank of a linear map or matrix: dimension of its range.]{RankOperator}{\operatorname{rank}}{plain}{local}{Rank of a linear map or matrix.}{operator-notation}
\dkDeclareSymbol[Distance between sets of spectral values: the infimum of pairwise distances.]{SpectralDistanceOperator}{\operatorname{dist}}{plain}{local}{Distance between two spectral sets.}{operator-notation}
\DeclareRobustCommand{\dkModulus}[1]{\like{OperatorModulusNotation}{|#1|}}

% Context-specific glyphs previously hidden in variants of other semantic keys.
% Keeping both the roles and the exact printed forms makes collisions explicit.
\dkDeclareSymbol[Nonzero vector x tested in the Rayleigh quotient, before imposing unit length.]{RayleighTestVector}{x}{trialvector}{local}{Test vector in the Rayleigh quotient.}{vector}
\dkDeclareSymbol[Positive cosine block C in the generic two-projection decomposition of Halmos; C and S commute and C squared plus S squared is the identity.]{HalmosCosineBlock}{C}{anglecos}{local}{Cosine block in the Halmos two-projection representation.}{linear-operator}
\dkDeclareSymbol[Positive sine block S in the generic two-projection decomposition of Halmos; C and S commute and C squared plus S squared is the identity.]{HalmosSineBlock}{S}{sineoperator}{local}{Sine block in the Halmos two-projection representation.}{linear-operator}
\dkDeclareSymbol[Riesz spectral projection P_Gamma defined by a contour Gamma that encloses selected spectral values.]{RieszProjection}{P_{\symglyph{RieszContour}}}{wantedprojection}{local}{Spectral projection associated with a contour Gamma.}{orthogonal-projection}
\dkDeclareSymbol[Historical Sylvester right-hand coefficient matrix A, in B X - X A = Y. Distinct from the current paper's reference operator A.]{HistoricalSylvesterRight}{A}{trialspectrum}{local}{Right-hand coefficient A in the historical Sylvester equation.}{self-adjoint-operator}
\dkDeclareSymbol[Scalar coefficient b in the Sylvester prototype (b-a)x=y, with separated scalar spectral values a and b.]{SylvesterScalarLeft}{b}{unwantedspectrum}{local}{Scalar left coefficient b in a Sylvester equation.}{scalar}
\dkDeclareSymbol[Scalar coefficient a in the Sylvester prototype (b-a)x=y.]{SylvesterScalarRight}{a}{trialspectrum}{local}{Scalar right coefficient a in a Sylvester equation.}{scalar}
\dkDeclareSymbol[Scalar unknown x in the Sylvester prototype (b-a)x=y.]{SylvesterScalarUnknown}{x}{sylvester}{local}{Scalar unknown x in a Sylvester equation.}{scalar}
\dkDeclareSymbol[Scalar right side y in the Sylvester prototype (b-a)x=y.]{SylvesterScalarRHS}{y}{resid}{local}{Scalar right side y in a Sylvester equation.}{scalar}
\dkDeclareSymbol[Source subspace K of a real Euclidean space in the matching proof; an orthogonal transporter sends K onto L.]{MatchingSourceSubspace}{K}{trial}{local}{Source subspace K in the matching proof.}{linear-subspace}
\dkDeclareSymbol[Target subspace L of a real Euclidean space in the matching proof; an orthogonal transporter sends K onto L.]{MatchingTargetSubspace}{L}{wanted}{local}{Target subspace L in the matching proof.}{linear-subspace}
\dkDeclareSymbol[Orthogonal projection P_L onto the matching proof's target subspace L.]{MatchingTargetProjection}{P_L}{wantedprojection}{local}{Orthogonal projection onto the matching target L.}{orthogonal-projection}
\dkDeclareSymbol[Real two-dimensional orthogonal group O(2): planar rotations and reflections.]{PlanarOrthogonalGroup}{O(2)}{fg}{local}{Orthogonal matrices of size two.}{group}
\dkDeclareSymbol[Base letter R used to denote planar rotation operators R_t in the matching proof.]{PlanarRotationBase}{R}{fg}{local}{Base letter for planar rotations.}{notation-base}
\dkDeclareSymbol[Planar rotation R_(r+s) through the sum of scalar phases r and s.]{MatchingPlanarRotationSum}{\symglyph{PlanarRotationBase}_{r+s}}{fg}{local}{Planar rotation through phase r+s.}{unitary-operator}
\dkDeclareSymbol[Planar rotation R_gamma through phase gamma in the block rotation calculation.]{MatchingPlanarRotationGamma}{\symglyph{PlanarRotationBase}_{\symglyph{MatchingPhaseGamma}}}{fg}{local}{Planar rotation through phase gamma.}{unitary-operator}

% Specific symbols in the historical and appendix examples. They retain their
% existing black-ink notation but receive locally accurate hover definitions.
\dkDeclareSymbol[Contour Gamma, a positively oriented closed curve in the resolvent set enclosing a selected isolated spectral cluster.]{RieszContour}{\Gamma}{plain}{local}{Contour used in the Riesz spectral-projection formula.}{contour}
\dkDeclareSymbol[Even integer nu in a tested Ky Fan comparison, with nu/2 principal-plane terms on the right side.]{KyFanEvenIndex}{\nu}{plain}{local}{Even integer in the principal-plane Ky Fan comparison.}{scalar}
\dkDeclareSymbol[Orthogonal projection Omega_k onto the k-th principal plane in the published displacement comparison.]{PrincipalPlaneProjectionBase}{\Omega}{plain}{local}{Projection onto the indicated principal plane.}{notation-base}
\dkDeclareSymbol[Scalar gap threshold gamma in the ordered Sylvester spectral separation B >= (gamma+delta)I and C <= gamma I.]{SylvesterGapThreshold}{\gamma}{plain}{local}{Spectral separation threshold in a Sylvester estimate.}{scalar}
\dkDeclareSymbol[Projection Omega_t from a net or sequence that converges strongly to the identity in the singular-value approximation lemma.]{StrongApproxProjection}{\Omega_t}{plain}{local}{Orthogonal projection converging strongly to the identity.}{orthogonal-projection}
\dkDeclareSymbol[First orthogonal projection Omega used in the two-block pinching inequality for a unitarily invariant norm.]{PinchingProjectionOne}{\Omega}{plain}{local}{First projection in the pinching inequality.}{orthogonal-projection}
\dkDeclareSymbol[Second orthogonal projection Upsilon used in the two-block pinching inequality for a unitarily invariant norm.]{PinchingProjectionTwo}{\Upsilon}{plain}{local}{Second projection in the pinching inequality.}{orthogonal-projection}
\dkDeclareSymbol[Finite-rank projection Gamma used to retain the image of a chosen finite-rank compression K Psi.]{CompressionProjectionImage}{\Gamma}{plain}{local}{Range projection in the finite-rank compression lemma.}{orthogonal-projection}
\dkDeclareSymbol[Finite-rank projection Psi used to select input vectors in a nearly maximizing singular-value compression.]{CompressionProjectionInput}{\Psi}{plain}{local}{Domain projection in the finite-rank compression lemma.}{orthogonal-projection}
\dkDeclareSymbol[Small positive tolerance eta in the finite-rank compression approximation estimate.]{CompressionTolerance}{\eta}{plain}{local}{Tolerance controlling the off-diagonal compression norm.}{scalar}

\dkDeclareSymbol[Rotation phase gamma in a planar block from the matching proof.]{MatchingPhaseGamma}{\gamma}{plain}{local}{Rotation phase gamma in the matching calculation.}{scalar}

% Two common indices with local mathematical meanings distinct from other deltas.
\dkDeclareSymbol[Kronecker delta: delta_ij is one when the indices i and j agree and zero otherwise. This is not the spectral-gap parameter delta.]{KroneckerDelta}{\delta}{plain}{local}{Kronecker delta: equals 1 for equal indices and 0 otherwise.}{scalar}
\dkDeclareSymbol[Sign variable epsilon taking the values -1 and +1 in the factorization over two phase choices.]{PhaseSignChoice}{\epsilon}{plain}{local}{Sign variable taking values minus one and plus one.}{scalar}

% Documented notation correspondence. The source column is a faithful
% transcription of each cited original context, NOT an asserted universal DK
% meaning. The original A_0 is deliberately listed for two distinct roles:
% its reference block and the comparison/trial operator. The source uses the
% same printed symbol for both; this paper is free to distinguish them.
\dkDeclareNotationSource{DK1970}{Davis and Kahan (1970), The Rotation of Eigenvectors by a Perturbation. III}{https://doi.org/10.1137/0707001}
\dkDeclareNotationSource{Seelmann2019}{Seelmann (2019), Semidefinite Perturbations in the Subspace Perturbation Problem}{https://arxiv.org/abs/1708.02463}
\dkDeclareNotationSource{Blechta2020}{Blechta (2020), Stability of Linear GMRES Convergence with Respect to Compact Perturbations}{https://arxiv.org/abs/2005.12960}
% Arguments: {key}{source ID}{source glyph}{section or equation}{role-specific qualification}
\dkDeclareNotationAlternative{A}{DK1970}{A}{\S1, (1.2)}{Reference self-adjoint operator.}
\dkDeclareNotationAlternative{H}{DK1970}{H}{\S1, (1.2)}{Additive perturbation.}
\dkDeclareNotationAlternative{Acurrent}{DK1970}{A+H}{\S1}{Perturbed operator.}
\dkDeclareNotationAlternative{E0}{DK1970}{E_0}{\S1, (1.1)}{Source isometric frame.}
\dkDeclareNotationAlternative{E1}{DK1970}{E_1}{\S1, (1.1)}{Reference complement frame.}
\dkDeclareNotationAlternative{F0}{DK1970}{F_0}{\S1}{Exact-target frame.}
\dkDeclareNotationAlternative{F1}{DK1970}{F_1}{\S1}{Exact-target complement frame.}
\dkDeclareNotationAlternative{ReferenceBlock0}{DK1970}{A_0}{\S1, (1.2)}{Reference source block.}
\dkDeclareNotationAlternative{ReferenceBlock1}{DK1970}{A_1}{\S1, (1.2)}{Reference complement block.}
\dkDeclareNotationAlternative{A0}{DK1970}{A_0}{\S1, residual discussion}{Residual comparison operator (the reference block in DK).}
% We deliberately leave Lambda0 and Lambda1 without DK source-glyph metadata:
% the available source transcription does not reliably preserve the source
% accents/decorations on the perturbed spectral-block symbols.
\dkDeclareNotationAlternative{Psource}{DK1970}{P}{\S1}{Reference/source projection.}
\dkDeclareNotationAlternative{Qtarget}{DK1970}{Q}{\S1}{Target projection.}
\dkDeclareNotationAlternative{Pexact}{DK1970}{Q}{\S1}{Exact-target projection.}
\dkDeclareNotationAlternative{B}{DK1970}{B}{\S1, (1.2)}{Off-diagonal coupling block.}
\dkDeclareNotationAlternative{Theta0}{DK1970}{\Theta_0}{\S1}{Source-side angle operator.}
\dkDeclareNotationAlternative{Theta1}{DK1970}{\Theta_1}{\S1}{Complementary angle operator.}
\dkDeclareNotationAlternative{Theta}{DK1970}{\Theta}{\S1}{Full angle operator.}
\dkDeclareNotationAlternative{R}{DK1970}{R}{\S1, residual discussion}{Trial residual.}
\dkDeclareNotationAlternative{Rrotation}{DK1970}{U}{\S1, (1.18)}{Canonical direct rotation.}
\dkDeclareNotationAlternative{Wcompetitor}{DK1970}{V}{\S4}{Competing unitary transporter.}
\dkDeclareNotationAlternative{delta}{DK1970}{\delta}{\S2}{Spectral separation parameter.}
% Contemporary alternatives are recorded separately and do not enter the DK
% correspondence table. Their citation locations must be individually checked.
\dkDeclareNotationAlternative{Eambient}{Seelmann2019}{\mathcal H}{p.~1}{Ambient Hilbert space.}
\dkDeclareNotationAlternative{SingularValueS}{Blechta2020}{\sigma_j}{Thm.~1.1}{Ordered singular values of a compact operator.}

\ifDKColoredSymbols\else\DKMathColorsfalse\fi
\ifDKAuthorStyle\else\DKSymbolTooltipsfalse\fi




\usepackage[most]{tcolorbox}
\usepackage{listings}
\usepackage{listingsutf8}
\usepackage{fontspec}
\IfFileExists{emoji.sty}{%
  \IfFontExistsTF{Noto Color Emoji}{\usepackage{emoji}}{}%
}{}
\providecommand{\emoji}[1]{\textnormal{[#1]}} % Review notes also work without an emoji font.

\setlength{\emergencystretch}{3em}
\raggedbottom

% ArXiv-friendly Lean code: no shell escape, minted, or external highlighter.
\definecolor{leankeyword}{HTML}{243B6B}
\definecolor{leantactic}{HTML}{70477A}
\definecolor{leantype}{HTML}{1F6473}
\definecolor{leancomment}{HTML}{39734D}
\definecolor{leanstring}{HTML}{8A4B2A}
\definecolor{leannumber}{HTML}{777777}
\definecolor{leanbackground}{HTML}{F7F7F3}
\lstdefinelanguage{lean}{%
  sensitive=true,
  morecomment=[l]{--},
  morecomment=[s]{/-}{-/},
  morestring=[b]",
  morekeywords=[1]{import,open,namespace,section,end,noncomputable,variable,variables,
    universe,universes,set_option,attribute,axiom,opaque,theorem,lemma,def,abbrev,
    example,instance,structure,class,inductive,notation,local,deriving,extends,
    by,where,let,have,show,fun,match,with,if,then,else,do,return,in,forall},
  morekeywords=[2]{simp,simpa,rw,erw,dsimp,unfold,ring,ring_nf,nlinarith,linarith,
    positivity,norm_num,field_simp,omega,exact,refine,apply,assumption,intro,intros,
    ext,funext,cases,rcases,obtain,constructor,left,right,classical,by_contra,
    contrapose,convert,using,suffices,calc},
  morekeywords=[3]{Prop,Type,Nat,Int,Real,Fin,Set,Submodule,Matrix,EuclideanSpace,
    LinearMap,ContinuousLinearMap,LinearPMap,RCLike,IsTrialResidual,
    IsExactSpectralDecomposition,FormBoundedSylvesterGap,
    NormalizedSymmetricOperatorIdealFamily},
}
\lstdefinestyle{leanpaper}{%
  language=lean,
  columns=fullflexible,
  backgroundcolor=\color{leanbackground},
  basicstyle=\ttfamily\scriptsize,
  commentstyle=\color{leancomment}\itshape,
  keywordstyle=[1]\color{leankeyword}\bfseries,
  keywordstyle=[2]\color{leantactic},
  keywordstyle=[3]\color{leantype},
  stringstyle=\color{leanstring},
  numberstyle=\tiny\color{leannumber},
  breakatwhitespace=false,
  breaklines=true,
  captionpos=b,
  keepspaces=true,
  numbers=left,
  numbersep=6pt,
  showspaces=false,
  showstringspaces=false,
  showtabs=false,
  tabsize=2,
  aboveskip=0.8em,
  belowskip=0.8em,
  literate=
    {\\KField}{{$\Bbbk$}}1
    {∃}{{$\exists$}}1
    {∧}{{$\land$}}1
    {∀}{{$\forall$}}1
    {≠}{{$\neq$}}1
    {⊤}{{$\top$}}1
}
\crefname{lstlisting}{formalization}{formalizations}
\Crefname{lstlisting}{Formalization}{Formalizations}
\renewcommand{\lstlistingname}{Formalization}

\theoremstyle{thmstyleone}
\newtheorem{theorem}{Theorem}
\newtheorem{proposition}[theorem]{Proposition}
\newtheorem{lemma}[theorem]{Lemma}
\newtheorem{remark}[theorem]{Remark}
\newtheorem{corollary}[theorem]{Corollary}

\newcommand{\Lean}{\textsc{Lean}}
\newcommand{\DK}{Davis--Kahan}
\newcommand{\TauCeti}{Tau Ceti}
% Use \ie{} and \eg{}; supply surrounding punctuation at the call site.
\newcommand{\ie}{i.e.\@}
\newcommand{\eg}{e.g.\@}
\newcommand{\DKResultCount}{29}
\newcommand{\DKProvedCount}{28}
\newcommand{\DKRefutedCount}{1}

% Usage: \margincomment{Check this formula!}
%\newcommand{\jonsays}[1]{%
%  \marginpar{%
%    \raggedright\color{red}\tiny% Change to \footnotesize if text is too small
%    \emoji{person-raising-hand} \textbf{Comment:}\\#1%
%  }%
%}

% Usage: \margincomment{Your comment that now wraps perfectly!}
\newcommand{\jonsays}[1]{%
  \marginpar{%
    \raggedright% Forces text to wrap beautifully without ugly hyphenation
    \color{red}\scriptsize% Using \scriptsize makes the text tighter
    \emoji{person-raising-hand}: \\#1%
  }%
}
%\newcommand{\jonsays}[1]{{\color{red}\small\textbf{[\emoji{person-raising-hand}: #1]}}}

% Usage: \margincomment{Your comment that now wraps perfectly!}
%\newcommand{\agentsays}[1]{%
%  \marginpar{%
%    \raggedright% Forces text to wrap beautifully without ugly hyphenation
%    \color{red}\scriptsize% Using \scriptsize makes the text tighter
%    \emoji{robot}: \\#1%
%  }%
%}
\newcommand{\agentsays}[1]{{\color{red}\small\textbf{[\emoji{robot}: #1]}}}

\newcommand{\jonrsp}[1]{{\color{red}\small\textbf{[\emoji{person-raising-hand}: #1]}}}

\title[Formalizing Davis--Kahan Perturbation Theory]{Formalizing Davis--Kahan Perturbation Theory in Lean 4}
% Additional authors are expected; keep the confirmed ordering below stable.
\ifDKAnonymous
  \author{Anonymous authors}
\else
  \author*[1]{\fnm{Jonathan} \sur{Crall}}\email{jon.crall@kitware.com}
  \author[2]{\fnm{Edward} \sur{Wang}}\email{ewang43@jhu.edu}
  \author[1]{\fnm{Brian} \sur{Hu}}\email{brian.hu@kitware.com}
  \author[2]{\fnm{Carey E.} \sur{Priebe}}\email{cep@jhu.edu}
  \affil[1]{\orgname{Kitware, Inc.},
    \orgaddress{\city{Clifton Park}, \state{NY}, \country{USA}}}
  \affil[2]{\orgdiv{Department of Applied Mathematics and Statistics},
    \orgname{Johns Hopkins University},
    \orgaddress{\city{Baltimore}, \state{MD}, \country{USA}}}
\fi

\abstract{We formalize the spectral-subspace perturbation theory of Davis and
Kahan's 1970 paper in Lean~4. Of \DKResultCount{} selected results, we prove
\DKProvedCount{} under their stated hypotheses and refute the remaining
printed statement. The development includes the sine and tangent theorems,
direct rotations, principal-angle operators, Sylvester separation, and
operator-ideal norms on real and complex Hilbert spaces, including unbounded
self-adjoint operators. We explain the mathematical lineage connecting these
constructions. The refuted statement concerns minimization of full ambient
displacement by direct rotation. We give real four-dimensional counterexamples
at every prescribed positive acute angle and identify the failed
singular-value comparison. A corrected optimality statement holds for Q-norms,
whose squares are unitarily invariant norms of the squared modulus, including
the Frobenius norm. The formalization builds on Mathlib and adapts Spectra's
projection-valued-measure and Borel-calculus constructions. Portions of the
operator theory developed here were subsequently adapted into Tau Ceti.
These foundations also support an exact reduction of trace-norm optimal real
subspace transport to weighted matching on principal angles.}

\keywords{Davis--Kahan theory, spectral perturbation, formal verification,
Lean 4, direct rotation, principal angles}
\pacs[MSC Classification (2020)]{68V20, 47A55, 15A42}

\begin{document}
\maketitle

\ifDKAnonymous\else
\enlargethispage{2\baselineskip}
\vspace{-18pt}
\noindent\begin{minipage}{\linewidth}\footnotesize
\noindent\textbf{Acknowledgments.}
This material is based upon work supported by the Defense Advanced Research
Projects Agency (DARPA) under Contract No. HR001125CE017. Any opinions, findings
and conclusions or recommendations expressed in this material are those of the
author(s) and do not necessarily reflect the views of the Defense Advanced
Research Projects Agency.
\end{minipage}
\fi
\clearpage

\section{Introduction}
\label{sec:introduction}

An approximate eigenvector gives a candidate direction for the action of an
operator. If a unit vector $\sym{v}$ is paired with an approximate eigenvalue $\sym{mu}$,
the residual $\sym{A}\sym{v}-\sym{mu} \sym{v}$ measures how closely it satisfies the eigenvector
equation. A small residual can identify an accurate approximation to an
invariant subspace even when repeated or nearby eigenvalues make individual
eigenvectors difficult to distinguish. Davis--Kahan theory bounds the angles
between approximate and exact spectral subspaces in terms of a residual or
perturbation and a separation between the relevant spectra
\citep{davis1970rotation}.

In applications, one often works with a perturbed operator
  $\sym{Acurrent}=\sym{A}+\sym{H}$ while the reference operator $\sym{A}$ is
  unknown.
Taking an eigenspace of $\sym{Acurrent}$ as a trial subspace for $\sym{A}$, the norm of its
  residual relative to $\sym{A}$ is bounded by the norm of $\sym{H}$, allowing
  Davis--Kahan theory to bound the angle between the two subspaces without
  requiring $\sym{A}$ to be known or its residuals to be evaluated directly.

Our project began with a variant of this bound needed for
\ifDKAnonymous a spectral-estimation application\else the data kernel
perspective space (DKPS) \citep{helm2026query}\fi.

We extended the formalization to the results established in the original 1970
paper, including real and complex Hilbert spaces, unbounded self-adjoint
operators, and unitarily invariant operator-ideal norms. In particular, the
spectral-separation hypotheses allow half-infinite intervals
(e.g.\ $(-\infty,a]$). These settings require operator domains, spectral
calculus, and norm estimates beyond finite matrices.

We formalized the four headline theorems \citep[Section~2]{davis1970rotation} and every
named theorem, proposition, lemma, and corollary established in the 1970 paper.

This gives \DKResultCount{} targets, enumerated in
Appendix~\ref{app:dk-inventory}. We prove \DKProvedCount{} of them and refute
the printed statement \citep[Proposition~4.4]{davis1970rotation}.
That proposition asserts a full-displacement
optimality property for real direct rotations. We give a family of real
four-dimensional examples exhibiting trace-norm failure at every prescribed
positive acute angle and identify
the invalid singular-value comparison in its proof. The phase construction
behind the counterexample occurs over the complex field in the authors'
example \citep[Example~4.2]{davis1970rotation}. To the best of our knowledge,
its contradiction of the printed real, small-angle claim has not previously
been made explicit.
For the constructed pair, whose two principal angles both equal
$\sym{theta}$, the competing transporter attains the global trace-norm minimum
$4\sin\sym{theta}$ for the constructed subspaces. Section~\ref{sec:prop44}
compares these statements and repairs the optimality claim by restricting
the norm class to Q-norms (N is a Q-norm when $N(T)^2 = M(T^* T)$ for some
unitarily invariant norm $M$). This repair follows from the squared-displacement
comparison \citep[Proposition~4.3]{davis1970rotation};
earlier Q-norm minimization results appear in \citet{bolotnikov2001minimal}.

% TODO (lineage): Consolidate the scattered lineage records through the
% maintained dev/tauceti/dkps-upstream-contribution-census.md and the evidence
% index in AUTHORING_NOTES.md. Audit each attribution against an immutable
% donor commit/path, the accepted upstream commit/PR, and the corresponding
% declaration statements and proofs. Preserve Mathlib prerequisites, Spectra
% Spectra-derived constructions where applicable, and independent work as distinct relationships.
% Courant--Fischer/Weyl checked 2026-10-28: developed here in June 2026;
% August donor df78246f965881705c9301f96feddc270c98ceb1 predates the
% accepted Tau Ceti port 7ebf17ba04725c1a575ecf3d099ce0630cd949a0 (#13599).
% See COURANT_FISCHER_LINEAGE_AUDIT.md for the proof comparison and sources.

The proofs required us to formalize a significant amount of operator
theory, namely: principal angles and direct rotations, approximation numbers
and symmetric operator ideals, Borel spectral calculus, and Sylvester
estimates for unbounded operators. These developments build on Mathlib and
adapt Spectra's projection-valued-measure and Borel-calculus constructions
\citep{mathlib2020,bornemann2026spectra}.
These operator-theory foundations include our
Courant--Fischer and Weyl proofs, rectangular unitarily invariant seminorms,
and projection-valued measures, which were subsequently adapted into Tau
Ceti \citep{tauceti2026}.
Section~\ref{sec:related} explains the mathematical history and formal
provenance; Section~\ref{sec:foundations} describes the foundations used in
our proofs.

The same principal-angle and operator-ideal results support a further
optimization theorem. In Section~\ref{sec:example-transport}, we prove that
minimizing trace-norm displacement over all real orthogonal maps carrying
one subspace onto another reduces to weighted matching on principal angles.

\noindent\textbf{Contributions.}
\begin{enumerate}
\item \textbf{Davis--Kahan formalization.} We prove 28 of the 29 source
results and refute the full-displacement assertion
\citep[Proposition~4.4]{davis1970rotation}. The sine-theta theorem covers real and
complex Hilbert spaces, unbounded self-adjoint operators, the source's
spectral-separation alternatives, and operator-ideal norms where finite
(\cref{sec:coverage,sec:sine-theta}).
\item \textbf{Formalized operator-theoretic foundations.} We formalized the
geometric and analytic results required by these theorems, adapted existing
formal mathematics, and contributed foundations later used in Tau Ceti
(\cref{sec:related-formal,sec:foundations}).
\item \textbf{Source correction.} We identify a false comparison
of Ky Fan sums in the printed claim \citep[Proposition~4.4]{davis1970rotation},
give a real counterexample at every positive acute angle,
and prove the Q-norm consequence of the valid squared-displacement theorem
(\cref{sec:prop44}).
\item \textbf{A reduction to weighted matching.} We reduce unrestricted
real trace-norm transport to a weighted-matching problem on principal angles
(\cref{sec:example-transport}).

\end{enumerate}

\section{Mathematical background}
\label{sec:related}

The Davis--Kahan angle estimates use three constructions. Principal angles
measure the error between subspaces; a residual measures how well a proposed
subspace represents the operator; and a Sylvester estimate bounds the first
quantity by the second under spectral separation. Operator ideals allow the
same argument to express several norms. We describe the development of these
constructions before stating the formalized sine-theta theorem.

Existing historical surveys cover these developments. Paige and Wei trace
the cosine--sine decomposition from Jordan's subspace geometry through the
principal ideas assembled by Davis and Kahan to Stewart's later completion
and explicit formulation \citep{paige1994history}.
Simon's historical survey of Kato's work in non-relativistic quantum mechanics
traces Sz.-Nagy's unitary construction for pairs of projections and Kato's
later algebraic formulation \citep[\S5]{simon2018kato}.
We use this account for the history of the projection transport described in
\cref{sec:related-projections}.
B\"ottcher and Spitkovsky survey pairs of projections and their operator algebras.
Their historical discussion relates Halmos's 1969 block representation to the
earlier formulations of Dixmier (1948) and Davis (1958), and to the related
projection constructions in Davis and Kahan's 1970 paper
\citep[Remark~1.3]{bottcher2010gentle}.
We use these accounts together with the primary
sources. Table~\ref{tab:historical-lineage} summarizes the mathematical
constructions used in our formalization.

\begin{table}[!htbp]
\centering
\small
\setlength{\tabcolsep}{5pt}
\renewcommand{\arraystretch}{1.18}
\begin{tabularx}{\linewidth}{@{}>{\raggedright\arraybackslash}p{0.26\linewidth}>{\raggedright\arraybackslash}X>{\raggedright\arraybackslash}p{0.28\linewidth}@{}}
\toprule
\textbf{Lineage} & \textbf{Mathematical role} & \textbf{Formal dependency}\\
\midrule
Jordan (1875); Davis (1958); Halmos (1969) & Principal angles, invariant subspace geometry,
  and pairs of projections & Angle operators; projection blocks\\
Rayleigh (1877); Ritz (1909); Weyl (1912) & Approximate eigenpairs,
  variational spectra, and eigenvalue stability & Trial coordinates;
  min--max and eigenvalue inequalities\\
Rellich (1937--1942); Sz.-Nagy (1947); Kato (1966) & Moving spectral projections and unitary
  transport between nearby subspaces & Spectral calculus; polar
  direct rotation\\
Sylvester (1884); Rosenblum (1956, 1969) & Spectral separation controls an
  off-diagonal operator & Domain-aware Sylvester estimates\\
Von Neumann (1937); Schatten (1950); Ky Fan (1951) & Comparison of spectral errors
  through singular-value gauges & Approximation numbers;
  symmetric operator ideals\\
Davis (1963, 1965); Davis--Kahan (1969, 1970) & Residual-to-angle bounds, tangent estimates,
  rotation extremality, and spectral branch selection & Source-level
  sine, tangent, and direct-rotation theorems\\
\bottomrule
\end{tabularx}
\caption{Mathematical sources used in the Davis--Kahan formalization.
Dates identify the publications discussed in the text.}
\label{tab:historical-lineage}
\end{table}
\FloatBarrier

\subsection{Principal angles}
\label{sec:related-classical}

Camille Jordan studied the geometry of linear subspaces in $n$ dimensions in
  1875 \citep{jordan1875essai}.

For two lines, a single acute angle describes their relative position.
For subspaces of higher dimension, one needs a sequence of \emph{principal
  angles}.
The successive principal vectors are selected under orthogonality constraints,
  so their angles are independent of the coordinates used to describe either
  subspace.
This separates the geometry of a subspace from the arbitrary choice of its
  orthonormal basis.

For finite-dimensional subspaces $\sym{Usource}$ and $\sym{Vtarget}$ of equal dimension, choose
isometries $\sym{Xframe}$ and $\sym{Yframe}$ with those ranges.
The singular values of the overlap matrix are the principal cosines:
\begin{equation}
  \sym{SingularValue}(\like{Xframe}{\symglyph{Xframe}^*}\sym{Yframe})=\cos\like{theta}{\symglyph{theta}_j},
  \qquad 0\leq\like{theta}{\symglyph{theta}_j}\leq\frac{\sym{PiConstant}}{2}.
  \label{eq:history-principal-cosines}
\end{equation}

The corresponding principal vectors $\like{SourceVector}{\symglyph{SourceVector}_j}\in \sym{Usource}$ and $\like{TargetVector}{\symglyph{TargetVector}_j}\in \sym{Vtarget}$ satisfy
$\langle \like{SourceVector}{\symglyph{SourceVector}_i},\like{TargetVector}{\symglyph{TargetVector}_j}\rangle=\like{KroneckerDelta}{\symglyph{KroneckerDelta}_{ij}}\cos\like{theta}{\symglyph{theta}_j}$ in suitably oriented
coordinates.  For nonzero acute angles, each pair determines a real
two-dimensional plane in which the relative position is an ordinary
rotation.  An entire pair of subspaces is described by these planes,
together with their shared and orthogonal components.

Principal angles also arise in statistics: Hotelling's 1936 canonical
  correlation analysis sought pairs of maximally correlated linear combinations
  of two sets of variables \citep{hotelling1936relations}.
After appropriate whitening, the canonical correlations are singular values of a
  cross-coordinate map.
In numerical linear algebra, Bj\"orck and Golub developed algorithms for
  computing principal angles and vectors in 1973 \citep{bjorck1973numerical}.
Their computational formulation makes the basis-invariant angles available from
  orthonormal matrix factorizations.

A further formulation uses the cosine--sine decomposition of a
partitioned orthogonal or unitary matrix.
The blocks of the decomposition contain compatible diagonal cosine and sine
  factors.
Davis and Kahan describe the geometry using orthogonal projections,
angle operators, and the direct rotation carrying one subspace onto the other
\citep[Section~3]{davis1970rotation}. Their work assembled the principal ideas
of the cosine--sine decomposition; Stewart completed and explicitly formulated
the matrix decomposition in 1977
\citep{stewart1977perturbation,paige1994history}.
Paige and Wei explain how this decomposition and direct rotation represent
the same principal-angle data
\citep{paige1994history}. Our formalization uses projections and functional
calculus for the Hilbert-space results and principal-plane coordinates for
the finite-dimensional counterexample.

\subsection{Eigenvalue approximation}
\label{sec:related-ritz-weyl}

Jordan's angles answered how to compare subspaces once they are given.
An independent question came from mechanics and numerical analysis:
how can one approximate an eigenvalue or eigenmode when the underlying
differential equation cannot be solved exactly?  In the nineteenth
century Rayleigh estimated vibration frequencies using trial shapes
and their energy ratios \citep{rayleigh1877sound}.  Ritz developed a systematic finite-dimensional
variational approximation in the work published in 1909
\citep{ritz1909variations,leissa2005historical}.  The method selects
a trial space and studies the action of the operator through that
space. A trial space is a chosen space of candidate approximate solutions.
For an eigenvalue problem, one searches for an approximate eigenvector among
its vectors. This reduces an infinite-dimensional problem to a smaller
problem whose residual can then be assessed.

For a Hermitian matrix $\sym{A}$, the Rayleigh quotient
$\langle \sym{A}\sym{RayleighTestVector},\sym{RayleighTestVector}\rangle/\langle \sym{RayleighTestVector},\sym{RayleighTestVector}\rangle$ expresses the scalar action
of $\sym{A}$ along a trial vector. A normalized trial vector $\sym{v}$ gives the
Rayleigh--Ritz approximate eigenvalue $\sym{mu}=\langle \sym{A}\sym{v},\sym{v}\rangle$ and the
residual $\sym{A}\sym{v}-\sym{mu} \sym{v}$. The Courant--Fischer min--max
principle organizes the eigenvalues through extremal quotients over
subspaces \citep{bhatia1997matrix}.  Weyl's eigenvalue comparison
inequalities then control the motion of ordered eigenvalues under a
Hermitian perturbation \citep{weyl1912eigenwerte,bhatia1997matrix}.  In modern notation,
if $\sym{A}$ and $\sym{Acurrent}=\sym{A}+\sym{H}$ are Hermitian $n$-by-$n$ matrices and their eigenvalues
are arranged in the same order, then
\begin{equation}
  \max_j |\like{SpectralValue}{\symglyph{SpectralValue}_j}(\sym{A})-\like{SpectralValue}{\symglyph{SpectralValue}_j}(\sym{Acurrent})|\leq\dkOperatorNorm{\sym{H}}.
  \label{eq:history-weyl}
\end{equation}
This is an estimate on \emph{eigenvalues}; it does not tell us which
individual eigenvector of $\sym{A}$ corresponds to which eigenvector of $\sym{Acurrent}$.

The limitation already appears in two dimensions.  For any $\sym{EigenvalueExampleEpsilon}>0$,
put
\begin{equation}
 \like{A}{\symglyph{A}_\varepsilon}=\begin{pmatrix}0&0\\0&\sym{EigenvalueExampleEpsilon}\end{pmatrix},
 \qquad
 \like{Acurrent}{\symglyph{Acurrent}_\varepsilon}=\begin{pmatrix}\sym{EigenvalueExampleEpsilon}&0\\0&0\end{pmatrix}.
 \label{eq:history-swap-example}
\end{equation}
The matrices have identical ordered eigenvalues and
$\dkOperatorNorm{\like{A}{\symglyph{A}_\varepsilon}-\like{Acurrent}{\symglyph{Acurrent}_\varepsilon}}=\sym{EigenvalueExampleEpsilon}$,
yet their lowest-eigenvalue lines are perpendicular. The gap between the
two eigenvalues of each matrix is $\sym{EigenvalueExampleEpsilon}$, equal to the perturbation
norm. An eigenspace estimate therefore needs to compare the perturbation
size with the spectral separation.  For a cluster with repeated eigenvalues, an
entire invariant subspace is likewise more natural than a selected
orthonormal eigenbasis.

Rellich developed a rigorous perturbation theory of self-adjoint operators
in papers published from 1937 to 1942. Kato's 1966 book gives a systematic
account \citep{simon2018kato,kato1966perturbation}. The Davis--Kahan residual
bounds address a related approximation question: given a trial space and its
residual, spectral separation bounds its angle from an exact spectral
subspace. The estimate can therefore assess the accuracy of a
Rayleigh--Ritz approximation.

\subsection{Two orthogonal projections}
\label{sec:related-two-projections}

For a closed subspace of a Hilbert space $\sym{GeneralHilbertSpace}$, the orthogonal
projection remains a bounded operator even when the subspace is
infinite dimensional.  Denote the projections onto two closed
subspaces by $\sym{Psource}$ and $\sym{Qtarget}$. Compressing an operator to $\DKRange(\sym{Psource})$ means
restricting its input to that subspace and projecting its output back with
$\sym{Psource}$. The positive compression of $\sym{Qtarget}$ is therefore $\sym{Psource}\sym{Qtarget}\sym{Psource}$ restricted to
$\DKRange(\sym{Psource})$. It is the operator version of the squared overlap matrix: in finite dimensions its eigenvalues are
$\cos^2\like{theta}{\symglyph{theta}_j}$.  In Hilbert space it need not have a discrete
spectrum.  Functional calculus nevertheless defines its positive
square root and the associated cosine and sine of an angle operator.

Write $\sym{Theta0}$ for this directed angle operator on
$\DKRange(\sym{Psource})$, and $\sym{Theta1}$ for the corresponding
operator for the complementary pair $\sym{Identity}-\sym{Psource},\sym{Identity}-\sym{Qtarget}$.
The full angle operator is
\[
 \sym{Theta}=\sym{DiagonalOperator}(\sym{Theta0},\sym{Theta1})
 \quad\text{on }\DKRange(\sym{Psource})\oplus\DKKernel(\sym{Psource}).
\]
For $c>0$, the comparisons $\sym{Theta}\le c$ and $\sym{Theta}<c$ mean that its
operator norm is at most, or strictly less than, $c$, respectively.
In finite dimensions this norm is the largest eigenvalue of either directed
angle block, comparing the source with the target and their orthogonal
complements.
This includes right angles from unmatched directions: if a source vector is
orthogonal to the target, its cosine compression is zero and its directed
angle is $\sym{PiConstant}/2$. The compression construction applies even when the subspace
dimensions differ. Target directions orthogonal to the source appear as
right angles in the complementary block. Later, isometric frames express these angle blocks on
coordinate Hilbert spaces.

The relative position of two closed subspaces was studied by Dixmier
in 1948 and by Davis in his 1958 paper \emph{Separation of Two Linear
Subspaces} \citep{dixmier1948position,davis1958separation}.  Davis's
closeness and separation operators place the cosine/sine data in an
operator-theoretic setting.  Halmos's 1969 \emph{Two Subspaces} gives a
block representation of the two projections
\citep{halmos1969two}.  One first splits off the four intersection
spaces corresponding to the combinations of membership in the ranges
and kernels of $\sym{Psource}$ and $\sym{Qtarget}$.  On their remaining, generic part, the pair
is unitarily equivalent to
\begin{equation}
 \like{Psource}{\symglyph{Psource}_{\mathrm{gen}}}=\begin{pmatrix}\sym{Identity}&0\\0&0\end{pmatrix},
 \qquad
 \like{Qtarget}{\symglyph{Qtarget}_{\mathrm{gen}}}=\begin{pmatrix}\like{HalmosCosineBlock}{\symglyph{HalmosCosineBlock}^2}&\sym{HalmosCosineBlock}\sym{HalmosSineBlock}\\\sym{HalmosCosineBlock}\sym{HalmosSineBlock}&\like{HalmosSineBlock}{\symglyph{HalmosSineBlock}^2}\end{pmatrix},
 \qquad \like{HalmosCosineBlock}{\symglyph{HalmosCosineBlock}^2}+\like{HalmosSineBlock}{\symglyph{HalmosSineBlock}^2}=\sym{Identity},
 \quad \sym{HalmosCosineBlock}\sym{HalmosSineBlock}=\sym{HalmosSineBlock}\sym{HalmosCosineBlock},
 \label{eq:history-halmos-block}
\end{equation}
where $\sym{HalmosCosineBlock}$ and $\sym{HalmosSineBlock}$ are commuting positive contractions on a common
Hilbert space.  Thus an operator-valued angle replaces a finite list
of principal angles.  B\"ottcher and Spitkovsky survey this
representation and its consequences for the algebra generated by two
projections \citep{bottcher2010gentle}.

Among these four intersections, the two right-angle intersection
subspaces (called defect spaces below) are
\[
 \sym{CrossedSourceDefect}=\DKRange(\sym{Psource})\cap\DKKernel(\sym{Qtarget}),
 \qquad
 \sym{CrossedTargetDefect}=\DKKernel(\sym{Psource})\cap\DKRange(\sym{Qtarget}).
\]
These contain source directions orthogonal to the target, and target directions
orthogonal to the source, respectively. The pair is called \emph{acute} when
both are $\{\sym{ZeroVector}\}$, the zero subspace consisting only of the zero
vector. For a general pair these spaces affect whether a specified
intertwining operator exists. An operator $\sym{IntertwinerT}$ intertwines $\sym{Psource}$ and $\sym{Qtarget}$ when
$\sym{IntertwinerT}\sym{Psource}=\sym{Qtarget}\sym{IntertwinerT}$: applying $\sym{IntertwinerT}$ after the source projection agrees with applying the
target projection after $\sym{IntertwinerT}$. In particular, $\sym{IntertwinerT}$ maps the source range into
the target range and its orthogonal complement into the target complement.
The existence criteria for direct rotations include conditions on the
dimensions of these defect spaces \citep[Section~3]{davis1970rotation}.

A finite-dimensional orthogonal matrix decomposes into invariant real
lines and two-dimensional rotation planes.  In a general Hilbert space an
angle operator can instead have continuous spectrum.  The spectral
calculus of the pair of projections takes the place of a finite list
of angles, and the intersection and defect spaces cannot be discarded.
Our Hilbert-space direct-rotation proofs therefore use projection identities
and functional calculus. In finite dimensions these identities recover the
principal-plane construction.

\subsection{Transport between nearby spectral subspaces}
\label{sec:related-projections}

A perturbation can change the invariant subspace associated with an isolated
part of a self-adjoint operator's spectrum. Let $\sym{Psource}$ and
$\sym{Qtarget}$ be the orthogonal projections onto the corresponding spectral
subspaces before and after perturbation. To compare these subspaces, we seek
a unitary carrying $\DKRange(\sym{Psource})$ onto
$\DKRange(\sym{Qtarget})$ that depends continuously, or analytically, on the
projections.
Sz.-Nagy developed the nearby-projection construction in work
published in 1947 \citep{sznagy1947perturbations,simon2018kato}.  Kato gave a streamlined
algebraic construction, described in his 1955 report and later in
\emph{Perturbation Theory for Linear Operators}
\citep[Chapter~I, \S4.6, equations~(4.33)--(4.42)]{kato1966perturbation};
see also \citet[Theorem~5.1]{simon2018kato}.

For orthogonal projections the useful intertwiner is
\begin{equation}
 \sym{IntertwinerT}=\sym{Qtarget}\sym{Psource}+(\sym{Identity}-\sym{Qtarget})(\sym{Identity}-\sym{Psource}), \qquad \sym{IntertwinerT}\sym{Psource}=\sym{Qtarget}\sym{IntertwinerT}.
 \label{eq:history-kato-T}
\end{equation}
The identity
\begin{equation}
 \like{IntertwinerT}{\symglyph{IntertwinerT}^*}\sym{IntertwinerT}=\sym{Identity}-(\sym{Psource}-\sym{Qtarget})^2
 \label{eq:history-kato-square}
\end{equation}
shows why $\|\sym{Psource}-\sym{Qtarget}\|<1$ is the natural threshold: it makes $\like{IntertwinerT}{\symglyph{IntertwinerT}^*}\sym{IntertwinerT}$
strictly positive and invertible.  Hence
\begin{equation}
 \sym{Rrotation}=\sym{IntertwinerT}\bigl(\sym{Identity}-(\sym{Psource}-\sym{Qtarget})^2\bigr)^{-1/2},
 \qquad \sym{Rrotation}\sym{Psource}\like{Rrotation}{\symglyph{Rrotation}^*}=\sym{Qtarget}
 \label{eq:history-kato-unitary}
\end{equation}
is unitary. This is the polar decomposition $\sym{IntertwinerT}=\sym{Rrotation}|\sym{IntertwinerT}|$, where
$|\sym{IntertwinerT}|=(\like{IntertwinerT}{\symglyph{IntertwinerT}^*}\sym{IntertwinerT})^{1/2}$ is positive. The factor $|\sym{IntertwinerT}|$ describes stretching;
the unitary factor $\sym{Rrotation}$ preserves lengths. Here $\sym{IntertwinerT}$ is invertible, so
$\sym{Rrotation}=\sym{IntertwinerT}|\sym{IntertwinerT}|^{-1}$. In principal-plane coordinates it is the direct rotation through the
principal angle, acting trivially on common directions.

The transport equation
$\sym{Wcompetitor}\sym{Psource}\sym{Wcompetitor}^*=\sym{Qtarget}$
has many unitary solutions.
The polar-factor formula selects one determined by the two projections
and continuously connected to the identity under the norm bound above. That
choice makes the construction valuable in perturbation theory: the
transport identifies moving spectral subspaces while avoiding an
arbitrary choice of orthonormal bases. In infinite dimensions,
$\|\sym{Psource}-\sym{Qtarget}\|<1$ is stronger than acuteness: it also
bounds the operator angles away from $\sym{PiConstant}/2$. For equal-dimensional
finite subspaces the conditions coincide. For an orthogonal or unitary transporter $\sym{Wcompetitor}$, the operator
$\sym{Identity}-\sym{Wcompetitor}$ records the displacement of every vector. Davis and Kahan ask which
norms of this displacement are minimized by the direct rotation
\citep[Section~4]{davis1970rotation}.

Kato used such transformations to identify nearby spectral
subspaces of analytic operator families, thereby reducing the study
of an isolated moving eigenvalue cluster to a fixed coordinate
space.  Simon's historical account explains the earlier Sz.-Nagy
construction, Kato's formula, and subsequent work on pairs of
projections \citep[\S5]{simon2018kato}. Our proofs of direct-rotation
existence, uniqueness, and symmetry use the square-root and
polar-decomposition identities in
\cref{eq:history-kato-square,eq:history-kato-unitary}. The condition
$\|\sym{Psource}-\sym{Qtarget}\|<1$ applies in arbitrary Hilbert dimension.

\subsection{Spectral projections and residuals}
\label{sec:related-spectral}

The spectral theorem represents a self-adjoint operator as
$\sym{A}=\int\sym{SpectralValue}\,d\sym{SpectralMeasure}(\sym{SpectralValue})$ for a projection-valued measure $\sym{SpectralMeasure}$.
For a Borel set $\sym{SpectralSet}$, the operator $\sym{SpectralMeasure}(\sym{SpectralSet})$ is the orthogonal projection onto
the corresponding spectral subspace. This gives a second way of obtaining
a subspace: select a region of the spectrum and take the range of its
spectral projection. The functional-calculus
framework developed from the work of Riesz, von Neumann, and others;
Kato made spectral projections central to the perturbation theory of
closed operators \citep{kato1966perturbation,simon2018kato}.
Unlike the finite-dimensional eigenspace picture, a spectral
subspace need not be spanned by a finite list of eigenvectors.
A self-adjoint operator may be unbounded, while each of its spectral
projections is bounded and everywhere defined. The Hellinger--Toeplitz
theorem states that an everywhere-defined symmetric operator on a Hilbert
space is bounded \citep{hellinger1910grundlagen}. Thus an unbounded
self-adjoint operator has a proper domain. Its action $\sym{A}x$ is defined only
for $x\in\DKDom(\sym{A})$, so residual and intertwining equations must include
the required domain inclusions.

In the complex case, a number $z$ belongs to the resolvent set when
$z\sym{Identity}-\sym{A}$ has a bounded, everywhere-defined inverse. That inverse,
$(z\sym{Identity}-\sym{A})^{-1}$, is the resolvent at $z$. For an isolated bounded spectral
cluster, choose a positively oriented contour $\sym{RieszContour}$ in the resolvent
set enclosing the cluster and no other spectrum. Its projection is
\begin{equation}
 \sym{RieszProjection}=\frac{1}{2\sym{PiConstant} i}\int_{\sym{RieszContour}}(z\sym{Identity}-\sym{A})^{-1}\,dz.
 \label{eq:history-riesz-projection}
\end{equation}
The resolvent is bounded at each point of $\sym{RieszContour}$ even when $\sym{A}$ is unbounded.
This makes the spectral projection a stable object for perturbation analysis,
  while preserving the separate domain requirements whenever $\sym{A}$ itself appears
  in an identity.
More general Borel spectral projections cover finite and half-infinite intervals,
  including spectral configurations for which a single isolated
  finite-eigenvalue cluster is not the right model.
Our formalization uses Borel functional calculus to construct projections for
  these spectral sets, including for unbounded self-adjoint operators.

The Courant--Fischer min--max principle and Weyl-type eigenvalue
perturbation bounds provide complementary information about spectral
\emph{values} \citep{bhatia1997matrix,kato1966perturbation}.  They do
not by themselves identify a stable basis inside a multiple-eigenvalue
cluster, or quantify the movement of its invariant subspace.  The Courant--Fischer and Weyl results that we formalized describe extremal
Rayleigh quotients and eigenvalue stability. The angle estimates add control
of the associated subspaces.

Davis studied the rotation of eigenvectors under perturbation in
1963 and 1965 \citep{davis1963rotation,davis1965rotation}.  The
1969 Davis--Kahan announcement preceded their extended 1970 treatment
\citep{daviskahan1969bounds,davis1970rotation}.  Their problem has
two closely related formulations.  In one, the original self-adjoint
operator is perturbed and the movement of an exact spectral subspace
is estimated.  In the other, one has an approximate trial subspace and
a computable residual: how far can it be from the desired exact
spectral subspace?  Unlike a comparison of individual eigenvectors,
these are meaningful even when an eigenvalue has multiplicity and its
orthonormal eigenbasis is nonunique.

Write $\sym{A}$ for a reference operator and $\sym{Acurrent}=\sym{A}+\sym{H}$ for the perturbed
ambient operator. We use this convention in the sine-theta statement below.
Let $\sym{E0}$ embed a trial-coordinate Hilbert space into the ambient space
isometrically, and let $\sym{TrialOperator}$ be a self-adjoint operator on those coordinates.
The residual is the bounded operator agreeing, where defined, with
\begin{equation}
 \sym{R}=\sym{Acurrent}\sym{E0}-\sym{E0}\sym{TrialOperator}.
 \label{eq:history-residual}
\end{equation}
If $\sym{R}=0$, applying $\sym{Acurrent}$ after the embedding agrees with applying
$\sym{TrialOperator}$ before it: the trial subspace is invariant and $\sym{TrialOperator}$ represents its
operator. For $\sym{R}\ne0$, the residual measures the failure of this identity.  A second isometric coordinate map for the
desired exact spectral space defines its orthogonal projection $\sym{Pexact}$.
The positive operator
\begin{equation}
 \sym{sinTheta0}=\bigl(\sym{Identity}-\like{E0}{\symglyph{E0}^*}\sym{Pexact}\sym{E0}\bigr)^{1/2}
 \label{eq:history-angle-sine}
\end{equation}
measures the component of a trial vector lying outside the exact
subspace.  The critical distinction for the formal proof is that this
positive sine acts on the trial-coordinate space, whereas
$(\sym{Identity}-\sym{Pexact})\sym{E0}$ maps into the ambient space.  For an operator $B$, its positive modulus $|B|=(B^*B)^{1/2}$ acts on its
input space and records its stretching factors. The modulus of $(\sym{Identity}-\sym{Pexact})\sym{E0}$
is $\sym{sinTheta0}$, so the two maps have equal unitarily invariant ideal
norms whenever those norms are finite.

\subsection{Sylvester estimates}
\label{sec:related-sylvester}

The matrix equation $\sym{SylvesterLeft}\sym{X}-\sym{X}\sym{HistoricalSylvesterRight}=\sym{SylvesterRHS}$ is associated with Sylvester's 1884
work \citep{sylvester1884equation}.  Rosenblum studied separated
operator spectra in 1956 and the equation with self-adjoint coefficients
in 1969 \citep{rosenblum1956equation,rosenblum1969selfadjoint}.

The scalar prototype is $(\sym{SylvesterScalarLeft}-\sym{SylvesterScalarRight})\sym{SylvesterScalarUnknown}=\sym{SylvesterScalarRHS}$: if $|\sym{SylvesterScalarLeft}-\sym{SylvesterScalarRight}|\geq\sym{delta}>0$,
then $\sym{delta}|\sym{SylvesterScalarUnknown}|\leq|\sym{SylvesterScalarRHS}|$.  In the operator problem the scalar
multipliers are replaced by the spectral actions of two operators.
The operator estimate has the form $\sym{delta} \sym{NormN}(\sym{X})\leq \sym{NormN}(\sym{SylvesterLeft}\sym{X}-\sym{X}\sym{HistoricalSylvesterRight})$.
Its hypotheses specify the spectral positions of $\sym{HistoricalSylvesterRight}$ and $\sym{SylvesterLeft}$, their domains
when unbounded, and the unitarily invariant norm $\sym{NormN}$.

Let $\sym{F1}$ give isometric coordinates for the exact complementary
spectral subspace, with self-adjoint block $\sym{Lambda1}$.  Writing
$\sym{X}=\like{F1}{\symglyph{F1}^*}\sym{E0}$ converts the residual identity into a Sylvester equation:
\begin{equation}
 \sym{Lambda1}\sym{X}-\sym{X}\sym{TrialOperator}=\like{F1}{\symglyph{F1}^*}\sym{R}
 \label{eq:history-sylvester}
\end{equation}
on the appropriate domain. Thus $\sym{X}$ represents the trial frame's
``leakage'' (i.e., its orthogonal projection onto the exact complementary
spectral subspace, expressed in $\sym{F1}$ coordinates), while the right
side is controlled by the computable residual.  Under the interval/exterior
or ordered spectral-gap hypotheses of the sine-theta theorem,
separated Sylvester estimates give
\begin{equation}
 \sym{delta}\,\sym{NormN}(\sym{X})\leq \sym{NormN}(\like{F1}{\symglyph{F1}^*}\sym{R})\leq \sym{NormN}(\sym{R}).
 \label{eq:history-separation-bound}
\end{equation}
This is the bridge from a spectral-gap assumption to the bound on the
sine of a subspace angle.

The coefficient of $\sym{NormN}(\sym{R})/\sym{delta}$ in this bound is one.
This coefficient uses the ordered or interval/exterior spectral
positions specified by Davis and Kahan. Arbitrarily interlaced spectra can
require a different constant. Our theorem includes both of the source's
spectral configurations, including ordered cases where both operators are
unbounded.  A formal implementation
therefore needs the domains of unbounded self-adjoint operators,
their spectral projections, strong spectral cutoffs, and the
ideal-norm inequalities used to pass through bounded approximations.

\subsection{Operator ideals and displacement norms}
\label{sec:related-norms}

The geometry of a pair of subspaces is fixed, but there are many ways
to assign a cost to a unitary transporting one onto the other.
Von Neumann's 1937 study of unitarily invariant matrix norms identified
the role of symmetric gauge functions of singular values
\citep{vonneumann1937matrix}. Schatten's 1950 monograph developed normed
operator ideals; Ky Fan's 1951 eigenvalue and singular-value inequalities
provided the partial-sum comparisons underlying majorization
arguments \citep{schatten1950crossspaces,fan1951maximum,bhatia1997matrix}.
The singular values $\sym{SingularValue}(Z)$ are the nonnegative square roots of the
eigenvalues of $Z^*Z$, ordered decreasingly. They record the stretching
factors of the matrix $Z$. For $1\leq p<\infty$, the Schatten norms include
\begin{equation}
 \|Z\|_{S_p}=\left(\sum_j \sym{SingularValue}(Z)^p\right)^{1/p},
 \qquad
 \|Z\|_{S_1}=\sum_j \sym{SingularValue}(Z),
 \quad
 \|Z\|_{S_2}=\left(\sum_j \sym{SingularValue}(Z)^2\right)^{1/2}.
 \label{eq:history-schatten}
\end{equation}
An operator ideal is a linear class of operators closed under multiplication
on either side by bounded operators. A normed ideal also controls the norm
of these products. For compact operators in infinite dimensions, the
Schatten $p$ ideal consists of those whose singular values have a finite
$p$-power sum. Thus a trace-norm inequality requires trace-class membership;
this is an instance of the source's qualification that its norms be defined.

The largest singular value gives the operator norm, the Euclidean
norm of the singular-value sequence gives the Hilbert--Schmidt norm
(the Frobenius norm $\dkFrobeniusNorm{Z}=\|Z\|_{S_2}$ for matrices),
and its sum gives the trace norm. These norms can rank transporters
differently. In Section~\ref{sec:prop44}, the two displacement singular-value
lists are $(2\sin(\sym{theta}/2),\ldots,2\sin(\sym{theta}/2))$ with four entries and
$(2\sin\sym{theta},2\sin\sym{theta},0,0)$. The second has a smaller sum and a larger
maximum. Ky Fan norms compare partial sums of the ordered singular values;
these comparisons determine which norm inequalities follow.

Davis established Hilbert--Schmidt and finite-dimensional Schatten
extremal properties for exponents at least two in his 1958 subspace
paper \citep[Theorems~7.3--7.4]{davis1958separation}.  Davis and
Kahan compare three objectives \citep[Section~4]{davis1970rotation} for a transporter $\sym{Wcompetitor}$ carrying
$\sym{Usource}$ onto $\sym{Vtarget}$: its restricted displacement $(\sym{Identity}-\sym{Wcompetitor})\like{Psource}{\symglyph{Psource}_U}$, its squared full
displacement $(\sym{Identity}-\sym{Wcompetitor})^*(\sym{Identity}-\sym{Wcompetitor})$, and its full displacement $\sym{Identity}-\sym{Wcompetitor}$. The last acts
on every ambient vector, including vectors outside $\sym{Usource}$. They gave a
complex-phase modification that decreases its trace norm \citep[Example~4.2]{davis1970rotation}.
The full-displacement claim \citep[Proposition~4.4]{davis1970rotation}
asserts that direct rotation minimizes
full displacement for real subspaces if their principal angles are no
larger than $\sym{PiConstant}/3$.  The four-dimensional example in
\cref{sec:prop44} gives a smaller trace norm at every positive acute
angle.  It exposes the failed step: adding singular-value sums of
restrictions to separate principal planes is not a valid global lower
bound for an arbitrary ambient transporter.

The subsequent minimal-distortion theory of Bolotnikov, Li, and Rodman
studied the same class of subspace-transport constraints, including
Frobenius and Q-norm optimality and explicit subquadratic
counterexamples \citep{bolotnikov2001minimal}.  Later work on
Grassmannian metrics and canonical-angle inequalities also separates
intrinsic paths between subspaces from minimizing the extrinsic
operator displacement \citep{qiu2006grassmann,zhang2010canonicalangles}.
These results give the norm-theoretic context for the source
correction and for the trace-norm matching problem in
\cref{sec:example-transport}.

\subsection{The Davis--Kahan angle estimates}
\label{sec:related-synthesis}

Davis's papers of 1963 and 1965 and the Davis--Kahan announcement of 1969
preceded the extended 1970 treatment
\citep{davis1963rotation,davis1965rotation,daviskahan1969bounds,davis1970rotation}.
For the sine-theta estimate, let $\sym{R}=\sym{Acurrent} \sym{E0}-\sym{E0}\sym{TrialOperator}$ be the residual
and $\sym{Pexact}$ the desired exact spectral projection. The operator $(\sym{Identity}-\sym{Pexact})\sym{E0}$
extracts the part of a trial vector outside that spectral subspace. In
isometric coordinates $\sym{F1}$ for the complementary space, this component is
represented by $\sym{X}=\like{F1}{\symglyph{F1}^*}\sym{E0}$. The residual identity gives
\cref{eq:history-sylvester}, and a separated Sylvester estimate bounds $\sym{NormN}(\sym{X})$.
Since $\sym{S}=\sym{F1}\sym{X}$ and $\sym{F1}$ is an isometry,
$\sym{NormN}(\sym{X})=\sym{NormN}(\sym{S})$.
The modulus identity $|\sym{S}|=\sym{sinTheta0}$ gives
$\sym{NormN}(\sym{S})=\sym{NormN}(\sym{sinTheta0})$.

The resulting chain of equalities and inequalities is
\begin{equation}
 \underbrace{\sym{delta} \sym{NormN}(\sym{sinTheta0})}_{\text{subspace error}}
 =\sym{delta} \sym{NormN}(\sym{X})
 \leq \sym{NormN}(\like{F1}{\symglyph{F1}^*}\sym{R})
 \leq\underbrace{\sym{NormN}(\sym{R})}_{\text{computable residual}}.
 \label{eq:history-unified-chain}
\end{equation}
The first equality follows from the positive-modulus identity for the
complementary projection. The first inequality applies the Sylvester
estimate to $\sym{X}$; the second uses $\sym{NormN}(BC)\leq\dkOperatorNorm{B}\sym{NormN}(C)$ and
$\dkOperatorNorm{\like{F1}{\symglyph{F1}^*}}\leq1$. For unbounded operators the Sylvester
identity holds on the trial operator's domain, with a bounded residual.
These are the geometric, spectral, and norm estimates formalized in our
proof of the theorem.

The tangent and double-angle estimates use related projection blocks.
For example, an estimate for $\sin(2\sym{Theta})$ allows angles near zero and
near $\sym{PiConstant}/2$. The spectral-selection results \citep[Section~8]{davis1970rotation}
identify which spectral subspaces give the small-angle alternative.
The direct-rotation analysis \citep[Section~4]{davis1970rotation}
studies the orthogonal or unitary map carrying one subspace onto another
and asks which displacement norms the direct rotation minimizes.

\subsection{Numerical and statistical applications}
\label{sec:related-numerical}

The 1970 theory influenced more than the analysis of exact
self-adjoint eigenvectors.  Wedin's 1972 singular-subspace
perturbation bounds address the corresponding problem for singular
vectors of general rectangular matrices
\citep{wedin1972perturbation}.  Bj\"orck and Golub's angle algorithms
made the subspace quantities computable
\citep{bjorck1973numerical}.  Stewart's work on pseudoinverses,
projections, and least-squares sensitivity helped connect these
geometric notions with matrix computations
\citep{stewart1977perturbation,stewart1990matrix}.
Paige and Wei's 1994 history shows how the cosine--sine decomposition
links the principal-angle geometry of Davis and Kahan with numerical
matrix factorizations \citep{paige1994history}.

The same subspace-versus-eigenvector distinction remains central in
statistical spectral estimation.  For example, Yu, Wang, and
Samworth give a statistical variant whose eigenspace error bound uses
separation in the population spectrum \citep{yu2015davis}, while Chen and
collaborators survey spectral methods across statistical learning
problems \citep{chen2021spectral}.  In that setting, an empirical covariance matrix approximates a population
operator and the perturbation norm controls the error of estimated
eigenspaces. Our general theorem supplies the deterministic angle estimate;
using it in a statistical result also requires the corresponding probabilistic
control of the perturbation.

\subsection{Foundations used in our formalization}
\label{sec:related-formal}

Each historical component above corresponds to a specific requirement
in the formalization.  Jordan's angle geometry and Halmos's
projection blocks lead to positive principal-angle operators rather
than chosen eigenbases.  The Sz.-Nagy--Kato intertwiner requires
polar decomposition, square roots, and real/complex transfer.
Spectral projections require Borel functional calculus, including for
unbounded self-adjoint operators.  The Sylvester proof requires a
representation carrying operator domains and a controlled passage
through spectral cutoffs.  Schatten and Ky Fan theory requires
approximation numbers and symmetric operator ideals to express
inequalities with the source's norm generality.

We formalized principal-angle geometry, symmetric operator ideals, and
unbounded Sylvester estimates using Lean~4 and Mathlib
\citep{demoura2021lean4,mathlib2020}. Mathlib supplies the algebraic,
topological, and Hilbert-space framework. Its continuous functional calculus
is described in the preprint by Anatole Dedecker and Jireh Loreaux
(\href{https://arxiv.org/abs/2501.15639}{arXiv:2501.15639}).
Our projection-valued-measure and Borel-calculus
work also adapts Adam Bornemann's Spectra formalization
\citep{bornemann2026spectra}.

Our Courant--Fischer and Weyl development, rectangular unitarily invariant
seminorms, projection-valued measures, and supporting polar-decomposition
results were subsequently adapted into Tau Ceti \citep{tauceti2026}.
% Editorial note (not manuscript prose): Do not add a sentence emphasizing
% the current Tau Ceti canonical polar-decomposition interface here.
% This paragraph records work developed in DKPS and subsequently
% contributed upstream; current library interfaces are discussed elsewhere.
% Evidence: dev/tauceti/dkps-upstream-contribution-census.md, PRs 13599,
% 14241, 17527 (adaptations), and 16127/16450 (independent polar proofs).
% AUTHORING_NOTES.md indexes accepted-source and campaign evidence.
% Mathematical origin and formal proof provenance are separate attributions.

This history also separates the two mathematical questions that meet
in the project.  The sine-theta theorem estimates \emph{how far} an
approximate or perturbed spectral subspace lies from an exact one.
The direct-rotation problem asks \emph{which} ambient transformation
carries one given subspace onto another at minimum displacement cost.
Our proofs of the optimization results use the principal-angle,
projection, and operator-ideal foundations developed for the formalization.
The matching theorem has its own proof; the shared foundations explain the
relationship between the two parts of the work. The resulting counterexample
to the full-displacement claim \citep[Proposition~4.4]{davis1970rotation}
and the follow-on matching theorem are examined in
\cref{sec:prop44,sec:example-transport}.

\section{Davis--Kahan formalization and the sine-theta theorem}
\label{sec:coverage}

Our targets are the four headline theorems \citep[Section~2]{davis1970rotation} and every
named theorem, proposition, lemma, and corollary established by Davis and
Kahan in the 1970 paper. These \DKResultCount{} results include the sine and
tangent inequalities, the geometry and extremal properties of direct
rotations, bounded and unbounded Sylvester estimates, the generalized
perturbation bounds \citep[Section~6]{davis1970rotation}, and the spectral-selection
theorems \citep[Section~8]{davis1970rotation}. We prove \DKProvedCount{} results and refute the printed
full-displacement claim \citep[Proposition~4.4]{davis1970rotation}.
Appendix~\ref{app:dk-inventory} lists these targets.

We develop the sine-theta theorem and its proof below; the remaining source
results are summarized in Appendix~\ref{app:dk-inventory}, with hypotheses
and conclusions in Section~\ref{app:dk-result-summaries}. The main text also
explains the spectral-selection results needed to interpret the angle bounds.
Appendix~\ref{app:notation} collects the notation used here and in the
direct-rotation argument.

\subsection{The sine-theta theorem for self-adjoint operators}
\label{sec:sine-theta}

\subsubsection{Trial and exact subspaces}

We use $\sym{A}$ for the reference operator and
$\sym{Acurrent}=\sym{A}+\sym{H}$ for the perturbed ambient operator.
The symbol $\sym{Acurrent}$ corresponds to $\sym{A}+\sym{H}$ in the source.

For a unit trial vector $\sym{v}$ and an approximate eigenvalue $\sym{mu}$, the residual
is $\sym{Acurrent}\sym{v}-\sym{mu} \sym{v}$. An isometric coordinate map $\sym{E0}$ generalizes
this construction to a trial subspace: its range consists of the candidate
approximate vectors, and a self-adjoint operator $\sym{TrialOperator}$ represents their
approximate dynamics. For a finite-dimensional Rayleigh--Ritz choice,
$\sym{TrialOperator}=\like{E0}{\symglyph{E0}^*}\sym{Acurrent}\sym{E0}$. The theorem allows other self-adjoint trial
operators as well. We write $\sym{TrialOperator}$ for the general trial operator,
reserving $\sym{ReferenceBlock0}$ for the reference spectral block that
Davis and Kahan also call $A_0$.

Let $\sym{AmbientHilbertSpace}$ be the ambient Hilbert space.  Let $\sym{TrialOperator}$ be a self-adjoint operator on a
trial-coordinate space $\sym{TrialCoordinateSpace}$, and let
\[
  \sym{E0}:\sym{TrialCoordinateSpace}\to \sym{AmbientHilbertSpace}
\]
be an isometric coordinate map for the trial subspace.  The residual compares
applying the ambient operator after embedding with applying the trial operator
before embedding:
\begin{equation}
  \sym{R} = \sym{Acurrent}\sym{E0}-\sym{E0}\sym{TrialOperator}.
  \label{eq:residual}
\end{equation}
When the operators are unbounded, \cref{eq:residual} is asserted on
$\DKDom(\sym{TrialOperator})$ and $\sym{R}$ is the bounded residual that agrees there.
The Lean predicate \texttt{IsTrialResidual} packages the isometry, the required
domain inclusion, and this residual identity.

Let $\sym{F0}:\sym{ExactCoordinateSpace}\to \sym{AmbientHilbertSpace}$ be isometric coordinates for the desired exact spectral
subspace and let $\sym{F1}:\sym{ComplementaryCoordinateSpace}\to \sym{AmbientHilbertSpace}$ be isometric coordinates for its orthogonal
complement.  A self-adjoint operator $\sym{Lambda1}$ on $\sym{ComplementaryCoordinateSpace}$ describes the
complementary spectral block through
\[
  \sym{Acurrent}\sym{F1} = \sym{F1}\sym{Lambda1}
\]
on its domain. This intertwining identity says that $\sym{Lambda1}$ gives
the action of $\sym{Acurrent}$ in complementary coordinates. The predicate
\texttt{IsExactSpectralDecomposition} also requires the domain inclusion
$\sym{F1}(\DKDom(\sym{Lambda1}))\subseteq\DKDom(\sym{Acurrent})$
and states orthogonality and completeness through
\[
 \like{F0}{\symglyph{F0}^*}\sym{F1}=0,\qquad
 \sym{F0}\like{F0}{\symglyph{F0}^*}+\sym{F1}\like{F1}{\symglyph{F1}^*}=\sym{Identity}.
\]

\subsubsection{The sine operator and its rectangular representative}

Here $\sym{Theta0}$ acts on $\sym{TrialCoordinateSpace}$; under the isometry
$\sym{E0}$ it represents the directed angle operator on the trial subspace
introduced in \cref{sec:related-two-projections}. If
$\sym{Pexact}=\sym{F0}\like{F0}{\symglyph{F0}^*}$ is the orthogonal projection onto the desired exact subspace, then
\[
  \cos\sym{Theta0} = (\like{E0}{\symglyph{E0}^*}\sym{Pexact}\sym{E0})^{1/2},
  \qquad
  \sym{sinTheta0} = (\sym{Identity}-\like{E0}{\symglyph{E0}^*}\sym{Pexact}\sym{E0})^{1/2}.
\]
The analytic proof uses the rectangular complementary block
\begin{equation}
  \sym{S}=(\sym{Identity}-\sym{Pexact})\sym{E0}.
  \label{eq:rectangular-sine}
\end{equation}
The theorem uses the positive sine on the trial-coordinate space.  The rectangular block has codomain $\sym{AmbientHilbertSpace}$, while the
positive sine has codomain $\sym{TrialCoordinateSpace}$.  Their relationship is
\[
  \sym{S}^*\sym{S}
  = \like{E0}{\symglyph{E0}^*}(\sym{Identity}-\sym{Pexact})\sym{E0}
  = \sym{Identity}-\like{E0}{\symglyph{E0}^*}\sym{Pexact}\sym{E0},
\]
so
\begin{equation}
  |\sym{S}|=(\sym{S}^*\sym{S})^{1/2}=\sym{sinTheta0}.
  \label{eq:modulus-sine}
\end{equation}
% It might be worth noting that our formalization actually defines \sym{Theta0}
% and proves eq:modulus-sine whereas most treatments of the subject seam to
% just start from |\sym{S}|. Maybe it's worth enough to put in the paper, but
% it at least deserves a latex comment!
The polar decomposition $\sym{S}=\sym{PolarU}|\sym{S}|$ separates the positive
  modulus from a partial isometry.
Unitarily invariant ideal norms assign the same value to an operator and its
  modulus, so \[ \sym{NormN}(\sym{S})=\sym{NormN}(\sym{sinTheta0}) \] whenever
  the norm is finite.
This equality transfers an estimate for the complementary block $\sym{S}$ to the
  positive sine operator.

\subsubsection{Gap geometry and the inequality}
\label{sec:gap-geometry}

Fix $\sym{delta}>0$. There are two geometric forms of separation.
\emph{Interval/exterior separation:} for some real
$\sym{betaEndpoint}\leq\sym{alphaEndpoint}$, one of the two spectra
$\sym{SpectrumSymbol}(\sym{TrialOperator})$ and $\sym{SpectrumSymbol}(\sym{Lambda1})$ lies in $[\sym{betaEndpoint},\sym{alphaEndpoint}]$ and the other
lies in
\[
 (-\infty,\sym{betaEndpoint}-\sym{delta}]\;\cup\;[\sym{alphaEndpoint}+\sym{delta},\infty).
\]
\emph{Ordered separation:} for some real $b$, the operator forms satisfy either
\[
 \sym{TrialOperator}\geq (b+\sym{delta})\sym{Identity},\quad \sym{Lambda1}\leq bI,
 \qquad\text{or}\qquad
 \sym{TrialOperator}\leq bI,\quad \sym{Lambda1}\geq (b+\sym{delta})\sym{Identity}.
\]
Each form inequality is asserted on the domain of the corresponding operator;
for example, $\sym{TrialOperator}\geq aI$ means
$\langle \sym{TrialOperator}x,x\rangle\geq a\|x\|^2$ for every
$x\in\DKDom(\sym{TrialOperator})$. For self-adjoint operators, the source's ordered
hypotheses on half-infinite spectral intervals, such as $(-\infty,b]$ and
$[b+\sym{delta},\infty)$, imply these form bounds. These cases are included
in the source's unbounded sine theorem
\citep[Section~6, discussion of unbounded operators]{davis1970rotation}.
Write $\sym{SineGap}(\sym{TrialOperator},\sym{Lambda1})$ for their union.
The Lean predicate \texttt{FormBoundedSylvesterGap} has three cases:
the finite-interval alternative, which includes either spectral order,
and the two ordered form alternatives.

Both blocks may be unbounded in an ordered configuration. For instance,
on $\ell^2(\mathbb N_{\geq1})$, take diagonal operators
$\sym{TrialOperator}e_n=(n+\sym{delta})e_n$ and $\sym{Lambda1}e_n=-ne_n$, on the domains where their
images are square summable. They satisfy the first ordered configuration
with $b=0$, while neither spectrum lies in a finite interval. This example
explains why retaining the half-infinite alternatives changes the theorem's
applicability.

For the unbounded trial configuration, the source assumes that
$\sym{Acurrent}\sym{E0}$ and $\sym{E0}\sym{TrialOperator}$ have a common dense domain.
Since $\sym{E0}$ is everywhere defined, this gives
\[
 \DKDom(\sym{Acurrent}\sym{E0})=\DKDom(\sym{TrialOperator}),
 \qquad
 \sym{E0}(\DKDom(\sym{TrialOperator}))\subseteq\DKDom(\sym{Acurrent}).
\]
The residual agrees with the difference on this domain and extends to a
bounded operator. The formal predicate \texttt{IsTrialResidual} includes the
domain inclusion and the residual identity.

The norm $\sym{NormN}$ is a normalized symmetric operator-ideal norm.
Normalization means that the rank-one map $uv^*:x\mapsto u\langle v,x\rangle$
has norm
\[
 \sym{NormN}(uv^*)=\|u\|\,\|v\|;
\]
in particular, its norm is one when $u,v$ are unit vectors.
Its hypotheses also include Ky Fan dominance whenever both ideal norms
are finite: domination of all finite singular-value sums implies domination
of the ideal norms. This comparison does not assert ideal membership.
The norm is invariant under isometries and positive modulus, and satisfies
$\sym{NormN}(BT)\leq\dkOperatorNorm{B}\sym{NormN}(T)$. Trace-class and
Hilbert--Schmidt operators with their respective norms are examples; the
operator norm on bounded operators is another. For an ideal norm, both the
residual and the positive sine must belong to its domain.

In two dimensions, take a one-dimensional trial frame $\sym{E0}(x)=x\sym{v}$ with
$\|\sym{v}\|=1$. A scalar trial operator $\sym{TrialOperator}=\sym{rho}$ gives
$\sym{r}=\sym{R}(1)=\sym{Acurrent}\sym{v}-\sym{rho}\sym{v}$.
Choosing $\sym{rho}=\langle\sym{v},\sym{Acurrent}\sym{v}\rangle$
makes $\sym{rho}\sym{v}$ the orthogonal projection of
$\sym{Acurrent}\sym{v}$ onto the trial line. The residual is therefore
perpendicular to that line. Figure~\ref{fig:sine-theta-residual} pairs this
residual triangle with the exact complementary projection that measures
the subspace error.

\begin{figure}[tb]
\centering
\resizebox{119mm}{!}{\input{figures/Fig1.pgf}}
\caption{Rayleigh residual geometry for a one-dimensional trial space in
$\sym{RealField}^2$ (the rank-one case). With $\sym{E0}(x)=x\sym{v}$, the residual
operator satisfies $\sym{R}(x)=x\sym{r}$, where
$\sym{r}=\sym{Acurrent}\sym{v}-\sym{rho}\sym{v}$.
For a higher-dimensional trial space,
$\sym{R}=\sym{Acurrent}\sym{E0}-\sym{E0}\sym{TrialOperator}$ assigns a residual vector
to each trial-coordinate vector. In a finite-dimensional eigenbasis of
$\sym{TrialOperator}$, its columns are the individual eigenpair residuals.
The exact frames in the drawing are
$\sym{F0}(h)=(h,0)$ and $\sym{F1}(g)=(0,g)$, with
$\sym{Pexact}=\sym{F0}\like{F0}{\symglyph{F0}^*}$ and $\sym{PexactPerp}=\sym{Identity}-\sym{Pexact}$.
(a) The pink leg $\sym{S}(1)=\sym{PexactPerp}\sym{v}$ has length $\sin\sym{theta}$.
(b) The green residual joins $\sym{rho}\sym{v}$ to $\sym{Acurrent}\sym{v}$;
for the Rayleigh choice it is perpendicular to the trial line. The dashed
legs are its exact wanted and unwanted components. Their lengths give
$\sym{delta}\sin\sym{theta}=\|\sym{PexactPerp}\sym{r}\|_2
\leq\dkOperatorNorm{\sym{r}\|_2=\|\sym{R}}$.
The drawing uses $\sym{Acurrent}=\sym{DiagonalOperator}(1.9,0.7)$,
$\sym{theta}=40^\circ$, $\sym{rho}=\langle\sym{v},\sym{Acurrent}\sym{v}\rangle$,
and $\sym{delta}=\sym{rho}-0.7$}
\label{fig:sine-theta-residual}
\end{figure}
% Place the residual figure before the theorem and its Lean signature.
\FloatBarrier

Appendix~\ref{app:lean-guide} explains the general Lean syntax in the
following theorem signature; the geometric and spectral hypotheses have been
defined above.

\begin{theorem}[Sine-theta estimate]
\label{thm:sine-theta}
Let the scalar field be $\sym{RealField}$ or $\sym{ComplexField}$, let $\sym{AmbientHilbertSpace}$ be separable,
and let $\sym{Acurrent}$, $\sym{TrialOperator}$, and $\sym{Lambda1}$ be self-adjoint operators on the
Hilbert spaces above. Suppose the isometries, domain inclusions, residual
identity, exact complementary decomposition, and the separation condition
$\sym{SineGap}(\sym{TrialOperator},\sym{Lambda1})$ of
Section~\ref{sec:gap-geometry} hold. Let $\sym{NormN}$ be a normalized symmetric
operator-ideal norm. Whenever $\sym{sinTheta0}$ and $\sym{R}$ belong to its domain,
\begin{equation}
 \sym{delta}\,\sym{NormN}(\sym{sinTheta0})\leq \sym{NormN}(\sym{R}).
 \label{eq:dk-sine}
\end{equation}
\end{theorem}

% Keep the Lean signature and its caption together across page breaks.
\noindent\begin{minipage}{\linewidth}
\begin{lstlisting}[style=leanpaper,
  caption={Lean signature for the sine-theta estimate. The surrounding
  scalar-field and Hilbert-space context is omitted.
  This is a readable ASCII transcription of the Lean declaration,
  rather than a verbatim copy of its Unicode binders and operator arrows.
  The paper-to-Lean names are registered centrally and displayed in
  Table~\ref{tab:paper-lean-notation}. In particular, the Lean binder
  displayed here as \texttt{A0} names our $\sym{TrialOperator}$, not the
  distinct reference block $\sym{ReferenceBlock0}$.
  The positive trial-coordinate sine is
  \texttt{sourceDirectedSinThetaOperator E0 F0};
  \texttt{FormBoundedSylvesterGap} encodes the interval/exterior and ordered
  half-infinite separation hypotheses described above.},
  label={lst:dk-sintheta-headline}]
theorem sinTheta_unbounded_formGap_whereDefinedUIN_rclike
    [TopologicalSpace.SeparableSpace E]
    (N : NormalizedSymmetricOperatorIdealFamily K)
    (A : LinearPMap K E E)
    (A0 : LinearPMap K F F)
    (Lambda1 : LinearPMap K G G)
    (E0 : ContinuousLinearMap K F E)
    (F0 : ContinuousLinearMap K H E)
    (F1 : ContinuousLinearMap K G E)
    (R : ContinuousLinearMap K F E)
    (hA : IsSelfAdjoint A)
    (hA0 : IsSelfAdjoint A0)
    (hLambda1 : IsSelfAdjoint Lambda1)
    (htrial : IsTrialResidual A A0 E0 R)
    (hexact : IsExactSpectralDecomposition A Lambda1 F0 F1)
    (delta : Real) (hdelta : 0 < delta)
    (hgap : FormBoundedSylvesterGap A0 Lambda1 delta) :
    N.Mem (sourceDirectedSinThetaOperator E0 F0) ->
    N.Mem R ->
      delta * N.gaugeReal (sourceDirectedSinThetaOperator E0 F0) <=
        N.gaugeReal R
\end{lstlisting}
\end{minipage}\par

The two ideal-membership hypotheses ensure that the norm values in the
conclusion are finite. The gap predicate includes the finite
interval/exterior case and both ordered cases. In the ordered cases,
$\sym{TrialOperator}$ and $\sym{Lambda1}$ may both be unbounded.

\subsubsection{Why the Sylvester estimate gives the angle bound}

Set $\sym{X}=\like{F1}{\symglyph{F1}^*}\sym{E0}$. The exact complementary decomposition gives
$\sym{S}=\sym{F1}\sym{X}$, and multiplication by the isometry $\sym{F1}$ preserves the
singular-value gauges. Self-adjointness, the domain hypotheses, and the
residual identity yield the Sylvester relation
\[
 \sym{Lambda1} \sym{X}-\sym{X}\sym{TrialOperator}=\like{F1}{\symglyph{F1}^*}\sym{R}
\]
on $\DKDom(\sym{TrialOperator})$, with the left side represented by the bounded
right side. The separated Sylvester estimate gives
\[
 \sym{delta} \sym{NormN}(\sym{X})\leq \sym{NormN}(\like{F1}{\symglyph{F1}^*}\sym{R})\leq \sym{NormN}(\sym{R}).
\]
Combining this with $\sym{NormN}(\sym{X})=\sym{NormN}(\sym{S})=\sym{NormN}(\sym{sinTheta0})$ proves the conclusion.
For unbounded operators, deriving the domain compatibility for $\sym{X}$ is part of
the argument. The displayed relation holds on $\DKDom(\sym{TrialOperator})$ and
has a bounded right side. The specified interval/exterior or ordered
spectral positions give the factor $1/\sym{delta}$, with no additional
multiplicative constant.

For a one-dimensional trial space, take
$\sym{Acurrent}=\sym{DiagonalOperator}(0,\sym{delta})$, the trial scalar
$\sym{TrialOperator}=0$, and the unit vector
$\like{v}{\symglyph{v}_{\symglyph{theta}}}=(\cos\sym{theta},\sin\sym{theta})$.
Choose the desired exact subspace $\sym{SpanOperator}\{(1,0)\}$;
its orthogonal complement is $\sym{SpanOperator}\{(0,1)\}$, on which
$\sym{Lambda1}=\sym{delta}$.
The complementary component is $(0,\sin\sym{theta})$ and the
residual is $(0,\sym{delta}\sin\sym{theta})$. Thus equality holds
in the sine-theta estimate for this finite-dimensional configuration.

\section{Formalized operator-theoretic foundations}
\label{sec:foundations}

We formalized the geometry of two subspaces, symmetric operator ideals and
majorization, and Sylvester estimates for unbounded self-adjoint operators.
These results supply the constructions and inequalities used in the
Davis--Kahan proofs.

\subsection{Direct rotations and two-subspace geometry}

Let $\sym{Psource}$ and $\sym{Qtarget}$ project onto closed subspaces
$\sym{Usource}$ and $\sym{Vtarget}$. For an acute pair,
$\|\sym{Psource}-\sym{Qtarget}\|<1$. The compression
$\sym{Psource}\sym{Qtarget}\sym{Psource}$ on $\sym{Usource}$ determines
the positive cosine operator. Its square root and the complementary
compression give the cosine and sine factors in the direct-rotation
construction. Polar decomposition gives the isometries between the
corresponding parts of the two subspaces.

The resulting orthogonal or unitary map $\sym{Rrotation}$ carries
$\sym{Usource}$ onto $\sym{Vtarget}$. In finite-dimensional principal-plane
coordinates it rotates each plane through its principal angle and fixes
common directions. The projection formula from
\cref{eq:history-kato-unitary} defines it in Hilbert space without requiring
a discrete list of principal vectors.

Our formalization proves acute existence and uniqueness, the
principal-square-root characterization, commutation with the angle
operators, and reversal symmetry: interchanging the two subspaces replaces
the rotation by its adjoint. For general pairs, the existence criteria
include the right-angle defect spaces described in
\cref{sec:related-two-projections}. When the angle operator is compact,
its spectral decomposition recovers a sequence of principal angles
\citep[Section~3]{davis1970rotation}.

The double-angle theorems for $\sin(2\sym{Theta})$ and $\tan(2\sym{Theta})$
  require selecting the branch on which $\sym{Theta}\le\sym{PiConstant}/4$.
Here $\sym{Theta}$ includes both the source and complementary angle blocks;
$\sym{Theta0}$ alone is the directed angle on source coordinates.
Because the function $\sin(2\sym{theta})$ is small both when $\sym{theta}$ is
  near $0$ and when $\sym{theta}$ is near $\sym{PiConstant}/2$, a bound on
  $\sin(2\sym{Theta})$ alone cannot determine whether an approximate subspace is
  close to or nearly orthogonal to the target.
The branch-selection theorem \citep[Theorem~8.1]{davis1970rotation}
uses an ordered reference gap. Let $\sym{Psource}$ reduce a self-adjoint
reference operator $\sym{A}$, with
blocks $\sym{ReferenceBlock0}$ on $\DKRange(\sym{Psource})$ and
$\sym{ReferenceBlock1}$ on $\DKKernel(\sym{Psource})$, satisfying
\[
 \sym{ReferenceBlock0}\le\sym{alphaEndpoint}\sym{Identity},
 \qquad
 \sym{ReferenceBlock1}\ge(\sym{alphaEndpoint}+\sym{delta})\sym{Identity},
 \qquad \sym{delta}>0,
\]
in the quadratic-form sense. Suppose $\sym{H}$ is bounded and self-adjoint,
with both diagonal blocks zero:
$\sym{Psource}\sym{H}\sym{Psource}
 =(\sym{Identity}-\sym{Psource})\sym{H}(\sym{Identity}-\sym{Psource})=0$.
These hypotheses imply that $\sym{Acurrent}=\sym{A}+\sym{H}$ has no spectrum
in $(\sym{alphaEndpoint},\sym{alphaEndpoint}+\sym{delta})$.

Fix a reducing subspace $M$ of $\sym{Acurrent}$, and write
$\sym{Lambda0},\sym{Lambda1}$ for the blocks of $\sym{Acurrent}$ on $M,M^\perp$.
The full angle $\sym{Theta}$ comparing $\DKRange(\sym{Psource})$ with $M$
satisfies $\sym{Theta}\le\sym{PiConstant}/4$ if and only if
$\sym{Lambda0}\le\sym{alphaEndpoint}\sym{Identity}$ and
$\sym{Lambda1}\ge(\sym{alphaEndpoint}+\sym{delta})\sym{Identity}$.
Such a subspace exists: take the spectral subspace of $\sym{Acurrent}$
for $(-\infty,\sym{alphaEndpoint}]$. Our formalization also proves that
it is the unique reducing subspace satisfying this angle bound.

To compare the reference and perturbed blocks on common coordinate spaces,
choose isometric frames for $\DKRange(\sym{Psource})$ and
$\DKKernel(\sym{Psource})$. Apply the direct rotation $\sym{Rrotation}$ to
these frames to obtain frames for $M$ and $M^\perp$. Express
$\sym{ReferenceBlock0},\sym{ReferenceBlock1}$ in the original frames and
$\sym{Lambda0},\sym{Lambda1}$ in the transported frames. In these coordinates
the branch satisfies the stronger quadratic-form inequality
\[
 \sym{ReferenceBlock1}-\sym{alphaEndpoint}\sym{Identity}
 \le\sym{ComplementCosineBlock}
       (\sym{Lambda1}-\sym{alphaEndpoint}\sym{Identity})\sym{ComplementCosineBlock},
 \qquad \sym{ComplementCosineBlock}=\cos\sym{Theta1},
\]
with the lower-block analogue. These inequalities quantify spectral
repulsion in terms of the angles and yield eigenvalue and symmetric-gauge
comparisons in finite dimensions.

The smallness criterion \citep[Theorem~8.2]{davis1970rotation} uses a central
spectral interval instead of an ordered gap. Let $\sym{Psource}$ reduce
the self-adjoint operator $\sym{A}$ and $\sym{Qtarget}$ reduce the
self-adjoint operator $\sym{Acurrent}$. The perturbed block on
$\DKRange(\sym{Qtarget})$ has spectrum in
$[\sym{betaEndpoint},\sym{alphaEndpoint}]$, while its complementary spectrum
lies in
\[
 (-\infty,\sym{betaEndpoint}-\sym{delta}]
 \ \cup\ [\sym{alphaEndpoint}+\sym{delta},\infty),
 \qquad \sym{betaEndpoint}\le\sym{alphaEndpoint},\quad\sym{delta}>0.
\]
The reference block $\sym{ReferenceBlock0}$, representing
$\sym{A}|_{\DKRange(\sym{Psource})}$ in source coordinates, must have spectrum
in $[\sym{betaEndpoint}-\sym{delta}/2,\sym{alphaEndpoint}+\sym{delta}/2]$.
For an isometric frame $\sym{E0}$ of the reference subspace, its residual
is $\sym{R}=\sym{Acurrent}\sym{E0}-\sym{E0}\sym{ReferenceBlock0}$
on the operator domain, as in \cref{eq:history-residual}.
Assume
$\sym{DimensionOperator}\sym{CrossedSourceDefect}=\sym{DimensionOperator}\sym{CrossedTargetDefect}$: the space of source
vectors orthogonal to the target and the space of target vectors orthogonal
to the source have equal Hilbert dimension.

If $\dkOperatorNorm{\sym{H}}<\sym{delta}/2$ for a
bounded self-adjoint perturbation, or $\dkOperatorNorm{\sym{R}}<\sym{delta}/2$
for the bounded residual on the reference subspace, then
$\sym{Theta}<\sym{PiConstant}/4$. For each ideal norm whose displayed arguments belong
to its domain, the corresponding double-angle estimate also holds:
\[
 \sym{delta}\sym{NormN}(\sin(2\sym{Theta}))\le2\sym{NormN}(\sym{H}),
 \qquad\text{or}\qquad
 \sym{delta}\sym{NormN}(\sin(2\sym{Theta0}))\le2\sym{NormN}(\sym{R}).
\]
% Note (not for print): the factor 2 in the residual form is correct and
% matches the 1970 Section 2 statement. The 1970 Section 7 proof, after (7.5),
% rewrites delta||sin 2Theta|| <= ||H - XHX|| with right side ||[[0,B*],[B,0]]||
% and concludes delta||sin 2Theta_0|| <= ||B|| <= ||R||; since H - XHX has
% off-diagonal block 2B, the right side should be 2||B|| <= 2||R||. The
% Section 7 example (A = diag(0,d), H = [[0,1],[1,-d]], delta = 2,
% theta = pi/4, ||R|| = 1) attains delta sin 2theta = 2 = 2||R||. Proof-line
% slip only; the stated theorem is unaffected.
The perturbation estimate uses the full angle, and the residual estimate
the directed angle on source coordinates. In the unbounded residual form,
$\sym{A}$ and $\sym{Acurrent}$ are self-adjoint on a common domain; only
the residual is required to be bounded.

We formalized both theorems in Hilbert space using spectral calculus.

\subsection{Approximation numbers and operator-ideal norms}

In finite dimensions, a unitarily invariant norm is a symmetric gauge
function of the singular values. Approximation numbers extend the
singular-value comparisons needed for bounded operators; for compact
operators they give the usual singular-value sequence. We formalized
symmetric norming functions and Ky Fan partial sums to pass from
majorization to norm inequalities.

For example, let $\sym{Rrotation}$ be the direct rotation and
$\sym{Wcompetitor}$ any unitary carrying $\sym{Usource}$ onto
$\sym{Vtarget}$. The restricted displacement of the direct rotation has
singular-value minimality, which implies minimality for unitarily invariant
norms. The squared-displacement result compares
$(\sym{Identity}-\sym{Rrotation})^*(\sym{Identity}-\sym{Rrotation})$ and
$(\sym{Identity}-\sym{Wcompetitor})^*(\sym{Identity}-\sym{Wcompetitor})$
\citep[Proposition~4.3]{davis1970rotation}. Their singular values are the
squares of those of the unsquared displacements. Section~\ref{sec:prop44}
shows why the claimed trace-norm comparison fails and derives the valid
Q-norm conclusion.
% Additional formalized quantitative results, omitted from this article:
% DavisKahan/FiniteDimensional/DirectRotation/SharpDisplacement.lean and
% SharpDisplacementLimit.lean. AUTHORING_NOTES.md indexes the evidence.

The same norm inequalities enter the sine-theta proof. The assertion
$\sym{delta}\sym{NormN}(\sym{S})\leq\sym{NormN}(\sym{R})$ applies when
$\sym{S}$ and $\sym{R}$ have finite ideal norm. Our ideal-membership
hypotheses state this condition explicitly in arbitrary Hilbert dimension.

\subsection{Sylvester estimates and unbounded operators}

For operators $\sym{SylvesterLeft}$ and $\sym{SylvesterRight}$ on the target and source spaces of $\sym{X}$, the
Sylvester map is $\sym{X}\mapsto \sym{SylvesterLeft}\sym{X}-\sym{X}\sym{SylvesterRight}$. Under the interval/exterior or ordered
spectral-separation hypotheses, the estimate has the form
\[
 \sym{delta} \sym{NormN}(\sym{X})\leq \sym{NormN}(\sym{SylvesterLeft}\sym{X}-\sym{X}\sym{SylvesterRight}).
\]
The sine-theta proof uses $\sym{SylvesterLeft}=\sym{Lambda1}$, $\sym{SylvesterRight}=\sym{TrialOperator}$, and
$\sym{X}=\like{F1}{\symglyph{F1}^*}\sym{E0}$. Other perturbation estimates choose related
projection blocks for $\sym{X}$.

For unbounded $\sym{SylvesterLeft}$ and $\sym{SylvesterRight}$, the identity is asserted for
$x\in\DKDom(\sym{SylvesterRight})$, with $\sym{X}x\in\DKDom(\sym{SylvesterLeft})$, and $\sym{SylvesterLeft}\sym{X}-\sym{X}\sym{SylvesterRight}$ agrees there with a
bounded residual. We formalized these domain conditions using linear
partial maps. Spectral cutoffs give bounded approximations; convergence
of approximation numbers passes the norm estimates to the unbounded
operators \citep[Section~5 and Appendix to Section~6]{davis1970rotation}.

The generalized sine and tangent theorems also compare block operators.
Reflection pinching replaces a block matrix $T$ by
$(T+\sym{Jreflect}T\sym{Jreflect})/2$, where $\sym{Jreflect}$ acts as $+\sym{Identity}$ on one coordinate subspace and $-\sym{Identity}$
on the other. This removes the off-diagonal blocks and contracts a unitarily
invariant norm. Direct-sum comparisons relate the norms of the component
blocks to the whole operator. Together with Sylvester separation these give
the generalized angle estimates
\citep[Lemmas~6.1--6.2 and Theorems~6.1--6.3]{davis1970rotation}.

\subsection{Real and complex Hilbert spaces}

The spectral constructions use complex functional calculus. We formalized
the scalar-transfer results needed to carry these arguments to real Hilbert
spaces through complexification. This permits the source theorems to be
stated over either field, including their domain and ideal-norm hypotheses.
The underlying algebra, topology, and continuous functional calculus come
from Mathlib; the adapted Spectra constructions and later Tau Ceti
contributions are described in Section~\ref{sec:related-formal}.

\section{A counterexample to full-displacement optimality}
\label{sec:prop44}

This section gives an explicit real counterexample to the small-angle
optimality assertion \citep[Proposition~4.4]{davis1970rotation} and
identifies the invalid inequality in its proof. To the best of our knowledge,
this contradiction of the printed real-space claim has not previously been
published. Its complex-phase construction is already in the authors'
example \citep[Example~4.2]{davis1970rotation}; the distinction between trace-norm failure and Q-norm optimality
also has earlier treatments \citep{bolotnikov2001minimal}. We give the
counterexample, the failed comparison, and a Q-norm repair below.

Write $\sym{Psource}$ and $\sym{Qtarget}$ for the orthogonal projections
onto subspaces $\sym{Usource}$ and $\sym{Vtarget}$. For an orthogonal or
unitary map $\sym{Wcompetitor}$ with
$\sym{Wcompetitor}(\sym{Usource})=\sym{Vtarget}$, the vector
$(\sym{Identity}-\sym{Wcompetitor})x$ is the displacement of $x$ up to sign. Full
displacement measures this operator on the whole ambient space; restricted
displacement measures $(\sym{Identity}-\sym{Wcompetitor})\sym{Psource}$, which tests only
vectors in the source subspace. We use the trace norm
$\dkTraceNorm{Z}=\sum_j \sym{SingularValue}(Z)$.

For a finite-dimensional acute pair, the direct rotation $\sym{Rrotation}$
is the orthogonal or unitary polar factor of
\[
 \sym{IntertwinerT}=\sym{Qtarget}\sym{Psource}
 +(\sym{Identity}-\sym{Qtarget})(\sym{Identity}-\sym{Psource}).
\]
The printed claim \citep[Proposition~4.4]{davis1970rotation} asserts that,
for real subspaces with largest principal angle
at most $\sym{PiConstant}/3$, it minimizes $\sym{NormN}(\sym{Identity}-\sym{Wcompetitor})$ for every
unitarily invariant norm $\sym{NormN}$.

\subsection{The real four-dimensional counterexample}

The competitor below rotates one real two-plane and fixes its orthogonal
complement. The direct rotation moves both principal planes. The competitor
has a smaller trace norm but a larger operator norm of its displacement.

\begin{theorem}[Direct rotation can fail to minimize the trace norm at any acute angle]
\label{thm:small-angle}
For every $0<\sym{theta}<\sym{PiConstant}/2$, there are two-dimensional subspaces
$\sym{Usource},\like{Vtarget}{\symglyph{Vtarget}_{\symglyph{theta}}}\subset\sym{RealField}^4$ and an orthogonal map $\like{Wcompetitor}{\symglyph{Wcompetitor}_{\symglyph{theta}}}$ with
$\like{Wcompetitor}{\symglyph{Wcompetitor}_{\symglyph{theta}}} \sym{Usource}=\like{Vtarget}{\symglyph{Vtarget}_{\symglyph{theta}}}$, such that both principal angles equal $\sym{theta}$ and
the canonical direct rotation $\like{Rrotation}{\symglyph{Rrotation}_{\symglyph{theta}}}$ satisfies
\begin{align}
 \|\sym{Identity}-\like{Wcompetitor}{\symglyph{Wcompetitor}_{\symglyph{theta}}}\|_1&=4\sin\sym{theta},\label{eq:competitor-trace}\\
 \|\sym{Identity}-\like{Rrotation}{\symglyph{Rrotation}_{\symglyph{theta}}}\|_1&=8\sin(\sym{theta}/2)>4\sin\sym{theta}.
 \label{eq:direct-trace}
\end{align}
\end{theorem}

\begin{proof}
Put $\sym{CosAngle}=\cos\sym{theta}$, $\sym{SinAngle}=\sin\sym{theta}$, and let
\[
 \sym{Jturn}=\begin{pmatrix}0&-1\\1&0\end{pmatrix},\quad
 \sym{Cturn}=\cos(2\sym{theta})\sym{IdentityTwo}+\sin(2\sym{theta})\sym{Jturn},\quad
 \sym{Hmix}=\frac{1}{\sqrt2}\begin{pmatrix}\sym{IdentityTwo}&\sym{IdentityTwo}\\I_2&-\sym{IdentityTwo}\end{pmatrix}.
\]
Define
\begin{equation}
 \like{Wcompetitor}{\symglyph{Wcompetitor}_{\symglyph{theta}}}=\sym{Hmix}\begin{pmatrix}\sym{Cturn}&0\\0&\sym{IdentityTwo}\end{pmatrix}\sym{Hmix}^T
 =\begin{pmatrix}\like{CosAngle}{\symglyph{CosAngle}^2}\sym{IdentityTwo}+\sym{SinAngle}\sym{CosAngle}\sym{Jturn}&-\like{SinAngle}{\symglyph{SinAngle}^2}\sym{IdentityTwo}+\sym{SinAngle}\sym{CosAngle}\sym{Jturn}\\
                  -\like{SinAngle}{\symglyph{SinAngle}^2}\sym{IdentityTwo}+\sym{SinAngle}\sym{CosAngle}\sym{Jturn}&\like{CosAngle}{\symglyph{CosAngle}^2}\sym{IdentityTwo}+\sym{SinAngle}\sym{CosAngle}\sym{Jturn}\end{pmatrix}.
 \label{eq:wtheta}
\end{equation}
Both $\sym{Hmix}$ and $\sym{Cturn}$ are orthogonal, so $\like{Wcompetitor}{\symglyph{Wcompetitor}_{\symglyph{theta}}}$ is orthogonal. Let
$\sym{Usource}=\sym{RealField}^2\oplus\{\sym{ZeroVector}\}$ and $\like{Vtarget}{\symglyph{Vtarget}_{\symglyph{theta}}}=\like{Wcompetitor}{\symglyph{Wcompetitor}_{\symglyph{theta}}} \sym{Usource}$.
The first two columns of $\like{Wcompetitor}{\symglyph{Wcompetitor}_{\symglyph{theta}}}$ factor as
\[
 \begin{pmatrix}\like{CosAngle}{\symglyph{CosAngle}^2}\sym{IdentityTwo}+\sym{SinAngle}\sym{CosAngle}\sym{Jturn}\\-\like{SinAngle}{\symglyph{SinAngle}^2}\sym{IdentityTwo}+\sym{SinAngle}\sym{CosAngle}\sym{Jturn}\end{pmatrix}
 =\begin{pmatrix}\sym{CosAngle}\sym{IdentityTwo}\\\sym{SinAngle}\sym{Jturn}\end{pmatrix}(\sym{CosAngle}\sym{IdentityTwo}+\sym{SinAngle}\sym{Jturn}).
\]
Since $\sym{CosAngle}\sym{IdentityTwo}+\sym{SinAngle}\sym{Jturn}$ is orthogonal, the columns of
$\sym{Yframe}=\binom{\sym{CosAngle}\sym{IdentityTwo}}{\sym{SinAngle}\sym{Jturn}}$ form an orthonormal basis of $\like{Vtarget}{\symglyph{Vtarget}_{\symglyph{theta}}}$.
With $\sym{Xframe}=\binom{\sym{IdentityTwo}}{0}$, one has $\sym{Xframe}^T\sym{Yframe}=\sym{CosAngle}\sym{IdentityTwo}$.
Thus both principal cosines are $\sym{CosAngle}>0$, and both principal angles are
$\sym{theta}$.

The corresponding projections and canonical intertwiner are
\[
 \sym{Psource}=\begin{pmatrix}\sym{IdentityTwo}&0\\0&0\end{pmatrix},\qquad
 \sym{Qtarget}=\sym{Yframe}\sym{Yframe}^T=\begin{pmatrix}\like{CosAngle}{\symglyph{CosAngle}^2}\sym{IdentityTwo}&-\sym{SinAngle}\sym{CosAngle}\sym{Jturn}\\\sym{SinAngle}\sym{CosAngle}\sym{Jturn}&\like{SinAngle}{\symglyph{SinAngle}^2}\sym{IdentityTwo}\end{pmatrix},
\]
\begin{equation}
 \sym{IntertwinerT}=\sym{Qtarget}\sym{Psource}+(\sym{Identity}-\sym{Qtarget})(\sym{Identity}-\sym{Psource})
   =\sym{CosAngle}\begin{pmatrix}\sym{CosAngle}\sym{IdentityTwo}&\sym{SinAngle}\sym{Jturn}\\\sym{SinAngle}\sym{Jturn}&\sym{CosAngle}\sym{IdentityTwo}\end{pmatrix}.
 \label{eq:intertwiner-theta}
\end{equation}
The last block matrix is orthogonal and $\sym{CosAngle}>0$, so $|\sym{IntertwinerT}|=\sym{CosAngle}\sym{IdentityFour}$ and
its polar factor is exactly
\[
 \like{Rrotation}{\symglyph{Rrotation}_{\symglyph{theta}}}=\begin{pmatrix}\sym{CosAngle}\sym{IdentityTwo}&\sym{SinAngle}\sym{Jturn}\\\sym{SinAngle}\sym{Jturn}&\sym{CosAngle}\sym{IdentityTwo}\end{pmatrix}.
\]
By the polar-factor characterization in \cref{eq:history-kato-unitary},
this is the canonical direct rotation for the constructed pair.
Since $\like{Rrotation}{\symglyph{Rrotation}_{\symglyph{theta}}}+\like{Rrotation}{\symglyph{Rrotation}_{\symglyph{theta}}}^T=2\sym{CosAngle}\sym{IdentityFour}$,
\[
 (\sym{Identity}-\like{Rrotation}{\symglyph{Rrotation}_{\symglyph{theta}}})^T(\sym{Identity}-\like{Rrotation}{\symglyph{Rrotation}_{\symglyph{theta}}})=(2-2\sym{CosAngle})\sym{IdentityFour}.
\]
Consequently all four displacement singular values equal
$\sqrt{2-2\sym{CosAngle}}=2\sin(\sym{theta}/2)$.

On the other hand, \cref{eq:wtheta} shows that the singular values of
$\sym{Identity}-\like{Wcompetitor}{\symglyph{Wcompetitor}_{\symglyph{theta}}}$ are those of $(\sym{IdentityTwo}-\sym{Cturn})\oplus0_2$.
The identity
$(\sym{IdentityTwo}-\sym{Cturn})^T(\sym{IdentityTwo}-\sym{Cturn})=4\like{SinAngle}{\symglyph{SinAngle}^2}\sym{IdentityTwo}$ gives the list
$(2\sym{SinAngle},2\sym{SinAngle},0,0)$. The two trace formulas follow. Finally,
\[
 4\sin\sym{theta}
 =8\sin(\sym{theta}/2)\cos(\sym{theta}/2)
 <8\sin(\sym{theta}/2),
\]
because $0<\sym{theta}/2<\sym{PiConstant}/4$.
\end{proof}

\begin{corollary}[No positive small-angle threshold]
\label{cor:no-threshold}
For every $\sym{EpsilonAngle}>0$, there are real two-dimensional subspaces of
$\sym{RealField}^4$ whose largest principal angle is strictly between zero and
$\sym{EpsilonAngle}$ and is at most $\sym{PiConstant}/3$, together with an orthogonal competitor
having strictly smaller trace-norm displacement than the direct rotation.
\end{corollary}

\begin{proof}
Apply \cref{thm:small-angle} with
$\sym{theta}=\min\{\sym{EpsilonAngle}/2,\sym{PiConstant}/4\}$.
\end{proof}

\noindent\begin{minipage}{\linewidth}
Take both principal angles to be $\sym{theta}=30^\circ$. The orthogonal map
in \cref{eq:wtheta} has two nonzero displacement singular values, both
equal to $1$. The direct rotation has four displacement singular values,
each equal to $2\sin 15^\circ$. Hence
\begin{equation}\label{eq:prop44-thirty-degrees}
 \|\sym{Identity}-\like{Wcompetitor}{\symglyph{Wcompetitor}_{\symglyph{theta}}}\|_1
 =4\sin 30^\circ=2
 <8\sin 15^\circ
 =\|\sym{Identity}-\like{Rrotation}{\symglyph{Rrotation}_{\symglyph{theta}}}\|_1 .
\end{equation}
The direct rotation costs approximately $2.071$, and the largest angle
$30^\circ$ is below the printed bound of $60^\circ$.
\Cref{lst:prop44-refutation} verifies the transport equation, the principal
angles, and the strict trace-norm comparison for every positive acute angle.

% Source: DavisKahan/FiniteDimensional/DirectRotation/SmallAngleCounterexample.lean,
% TauCeti.DavisKahan.FiniteDimensional.SmallAngleCounterexample.explicit_traceNorm_counterexample.
% U4 is defined in ShortRotationCounterexample.lean.
\begin{lstlisting}[style=leanpaper,
  caption={The explicit counterexample to Proposition~4.4.
  In $\sym{RealField}^4$, \texttt{U4} is the first coordinate plane,
  \texttt{Wequiv theta} is the orthogonal map in~\eqref{eq:wtheta},
  and \texttt{V4 theta} is its image of \texttt{U4}.
  Principal angles are indexed from zero in decreasing order.
  \texttt{traceDisplacement W} is $\|\sym{Identity}-W\|_1$.},
  label={lst:prop44-refutation}]
theorem explicit_traceNorm_counterexample
    (theta : Real) (hpos : 0 < theta)
    (hangle : theta < Real.pi / 2) :
    let V := V4 theta
    let W := Wequiv theta
    let R := directRotation U4 V
      (acute theta (cos_pos theta hpos hangle))
    U4.map W.toLinearMap = V (*@$\land$@*)
      U4.principalAngles V 0 = theta (*@$\land$@*)
      U4.principalAngles V 1 = theta (*@$\land$@*)
      traceDisplacement W < traceDisplacement R := by (*@\laaettr@*)
\end{lstlisting}
\end{minipage}\par

The competitor rotates a diagonal real two-plane through $2\sym{theta}$ and
fixes its orthogonal complement; its displacement has rank two, compared
with rank four for $\like{Rrotation}{\symglyph{Rrotation}_{\symglyph{theta}}}$. The operator norm favors the canonical
rotation instead:
$\dkOperatorNorm{\sym{Identity}-\like{Rrotation}{\symglyph{Rrotation}_{\symglyph{theta}}}}=2\sin(\sym{theta}/2)
 <2\sin\sym{theta}=\dkOperatorNorm{\sym{Identity}-\like{Wcompetitor}{\symglyph{Wcompetitor}_{\symglyph{theta}}}}$.
The competitor attains the global trace-norm minimum $4\sin\sym{theta}$.
Indeed, the projection difference has trace norm
$\|\sym{Psource}-\sym{Qtarget}\|_1=4\sin\sym{theta}$, and every orthogonal
transporter $\sym{Wcompetitor}$ satisfies
$\|\sym{Psource}-\sym{Qtarget}\|_1\leq\|\sym{Identity}-\sym{Wcompetitor}\|_1$.
Equation~\eqref{eq:competitor-trace} attains this lower bound.
% Formal proof: DavisKahan/FiniteDimensional/DirectRotation/TraceOptimality.lean,
% SmallAngleCounterexample.traceDisplacement_lower_bound and Wequiv_traceMinimizer.
% Dimension minimality is now proved in DirectRotation/DimensionMinimality.lean:
% nuclear_directRotation_le_of_finrank_le_three compares the actual direct
% rotation with every ambient real orthogonal transporter in dimension <= 3,
% under the acute-pair and largest-principal-angle <= pi/3 hypotheses.
% isLeast_real_shortAngle_nuclear_counterexample_dimension proves the least
% dimension is four; the same file retains arbitrarily small positive witnesses.
% Proposed corollary for the next prose review: Four is the least real ambient
% dimension for trace-norm failure under the printed small-angle hypotheses.

At $\sym{theta}=\sym{PiConstant}/4$, \cref{fig:prop44-displacement-planes}
shows the different displacement patterns of the canonical direct rotation
and the competing transporter. Both carry the same source subspace onto the
same target, with both principal angles equal to $45^\circ$. The competitor
fixes half of the ambient space; the direct rotation moves both principal planes.

\begin{figure}[tbp]
\centering
\resizebox{119mm}{!}{\input{figures/Fig2.pgf}}
\caption{A counterexample to the full-displacement claim
\citep[Proposition~4.4]{davis1970rotation} at
$\sym{theta}=\sym{PiConstant}/4$. Both orthogonal maps carry $\sym{Usource}$ onto
$\sym{Vtarget}$. Each circle shows the action on one invariant real
two-plane: (a,b) the principal planes $\Pi_1,\Pi_2$ for $\sym{Rrotation}$,
and (c,d) the moving and fixed planes $M,Z$ for $\sym{Wcompetitor}$.
The unit inputs are $u_1=e_0$, $u_2=e_1$,
$x=(e_0+e_2)/\sqrt2$, and $z=(e_0-e_2)/\sqrt2$.
For $T=\sym{Rrotation}$ or $\sym{Wcompetitor}$, the green arrow from
the image $Tq$ to the input $q$ is the displacement $(\sym{Identity}-T)q$.
The bars list the four singular values of $\sym{Identity}-T$, grouped by contributing
plane, on the same vertical scale. Their sums give
$\|\sym{Identity}-\sym{Wcompetitor}\|_1=2\sqrt2<4\sqrt{2-\sqrt2}=\|\sym{Identity}-\sym{Rrotation}\|_1$.
The two maps have different invariant-plane decompositions}
\label{fig:prop44-displacement-planes}
\end{figure}

\subsection{Failure of the published proof}

At $\sym{theta}=\sym{PiConstant}/4$, the family gives
\[
 \like{Wcompetitor}{\symglyph{Wcompetitor}_{\symglyph{PiConstant}/4}}=\frac12
 \begin{pmatrix}
 1&-1&-1&-1\\1&1&1&-1\\-1&-1&1&-1\\1&-1&1&1
 \end{pmatrix}.
\]
The two trace norms are
\[
 \|\sym{Identity}-\like{Wcompetitor}{\symglyph{Wcompetitor}_{\symglyph{PiConstant}/4}}\|_1=2\sqrt2
 <4\sqrt{2-\sqrt2}=\|\sym{Identity}-\like{Rrotation}{\symglyph{Rrotation}_{\symglyph{PiConstant}/4}}\|_1.
\]

For $\sym{Kdisplacement}=\sym{Identity}-\like{Wcompetitor}{\symglyph{Wcompetitor}_{\symglyph{PiConstant}/4}}$, the even-index part
of the published comparison \citep[equation~(4.3)]{davis1970rotation}
reads, in our notation,
\[
 \|\sym{Kdisplacement}\|_{(\sym{KyFanEvenIndex})}\geq
 \sum_{k=1}^{\sym{KyFanEvenIndex}/2}\|\sym{Kdisplacement}\like{PrincipalPlaneProjectionBase}{\symglyph{PrincipalPlaneProjectionBase}_k}\|_{(2)},
 \qquad \sym{KyFanEvenIndex}\ \text{even},
\]
where $\sym{KyFanNormNotation}$ is the Ky Fan $j$-norm and $\like{PrincipalPlaneProjectionBase}{\symglyph{PrincipalPlaneProjectionBase}_k}$ projects onto
the $k$th principal plane. Choose $\sym{KyFanEvenIndex}=4$ and let
$\sym{OmegaOne}$ and $\sym{OmegaTwo}$ project onto
$\sym{SpanOperator}\{e_0,e_3\}$ and
$\sym{SpanOperator}\{e_1,e_2\}$. These are principal planes for the pair.
The comparison then requires
\begin{equation}
 \|\sym{Kdisplacement}\|_{(4)}\geq
 \|\sym{Kdisplacement}\sym{OmegaOne}\|_{(2)}+\|\sym{Kdisplacement}\sym{OmegaTwo}\|_{(2)},
 \label{eq:dk43-false}
\end{equation}
Direct calculation gives
\[
 \|\sym{Kdisplacement}\|_{(4)}=2\sqrt2,\qquad
 \|\sym{Kdisplacement}\sym{OmegaOne}\|_{(2)}=\|\sym{Kdisplacement}\sym{OmegaTwo}\|_{(2)}=2.
\]
Thus \cref{eq:dk43-false} would imply $2\sqrt2\geq4$.
The displayed matrix also gives
\[
 \DKRange(\sym{Kdisplacement}\sym{OmegaOne})
 =\DKRange(\sym{Kdisplacement}\sym{OmegaTwo})
 =\sym{CounterMovingPlane}
 =\sym{SpanOperator}\{e_0+e_2,e_1+e_3\}.
\]
Both restrictions map onto the same two-dimensional output plane, even
though their input planes are orthogonal. Each has singular values $1,1$,
whereas the full displacement has just two nonzero singular values,
$\sqrt2,\sqrt2$. The comparison in \cref{eq:dk43-false} would follow if the
singular-value list of the full displacement were the concatenation of
the two lists for the restrictions. This argument would be valid for an
orthogonal direct sum, but it does not apply here because the output
planes coincide.

\subsection{Repair of Davis and Kahan's Proposition~4.4: direct-rotation optimality for Q-norms}

Call a unitarily invariant norm $\sym{NormN}$ a $\sym{NormQ}$-norm when there is a unitarily
invariant norm $\sym{NormM}$ satisfying
\begin{equation}
 \sym{NormN}(Z)^2=\sym{NormM}(Z^*Z)\qquad\text{for every }Z.
 \label{eq:qnorm-def}
\end{equation}

For example, the Frobenius norm is a $\sym{NormQ}$-norm because
$\dkFrobeniusNorm{Z}^2=\dkTraceNorm{Z^*Z}$: its squared norm is the sum of the eigenvalues of
the positive matrix $Z^*Z$. This identity uses the trace norm for
$\sym{NormM}$. The squared-displacement comparison
\citep[Proposition~4.3]{davis1970rotation} gives the following Q-norm consequence.
Earlier Q-norm minimization results appear in
\citet[Theorem~3.1]{bolotnikov2001minimal}.

\begin{proposition}[$\sym{NormQ}$-norm consequence]
\label{prop:qnorm-repair}
Let $\sym{Usource},\sym{Vtarget}$ be an acute pair of subspaces of a finite-dimensional real or
complex Hilbert space, and let $\sym{Rrotation}$ be their canonical direct rotation.
For every $\sym{NormQ}$-norm $\sym{NormN}$ and every orthogonal or unitary map $\sym{Wcompetitor}$ with $\sym{Wcompetitor}\sym{Usource}=\sym{Vtarget}$,
\begin{equation}
 \sym{NormN}(\sym{Identity}-\sym{Rrotation})\leq \sym{NormN}(\sym{Identity}-\sym{Wcompetitor}).
 \label{eq:qnorm-repair}
\end{equation}
No $\sym{PiConstant}/3$ restriction is required.
\end{proposition}

\begin{proof}
The valid squared-displacement theorem gives
\[
 \sym{NormM}((\sym{Identity}-\sym{Rrotation})^*(\sym{Identity}-\sym{Rrotation}))\leq \sym{NormM}((\sym{Identity}-\sym{Wcompetitor})^*(\sym{Identity}-\sym{Wcompetitor})).
\]
By \cref{eq:qnorm-def}, these quantities are $\sym{NormN}(\sym{Identity}-\sym{Rrotation})^2$ and $\sym{NormN}(\sym{Identity}-\sym{Wcompetitor})^2$.
Both norm values are nonnegative, so taking square roots proves the claim.
\end{proof}

The trace norm on the counterexample space does not satisfy the
$\sym{NormQ}$-norm condition. Its smaller competing displacement in
\cref{thm:small-angle} contradicts the minimality that this condition
would imply.

\subsection{Historical context}

Davis and Kahan already exhibited a complex two-dimensional phase
modification that decreases the trace norm of full displacement
\citep[Example~4.2]{davis1970rotation}. Our real family is its
realification. Indeed, under the identification $\sym{ComplexField}^2\simeq
\sym{RealField}^2\oplus\sym{RealField}^2$ in which multiplication by $i$ is
represented by $\sym{Jturn}$ from \cref{eq:wtheta}, the matrices
\[
 \like{Rrotation}{\symglyph{Rrotation}_{\symglyph{theta}}}^{\sym{ComplexField}}=
 \begin{pmatrix}\sym{CosAngle}&i\sym{SinAngle}\\i\sym{SinAngle}&\sym{CosAngle}\end{pmatrix},\qquad
 \like{Wcompetitor}{\symglyph{Wcompetitor}_{\symglyph{theta}}}^{\sym{ComplexField}}=e^{i\sym{theta}}\like{Rrotation}{\symglyph{Rrotation}_{\symglyph{theta}}}^{\sym{ComplexField}}
\]
become precisely $\like{Rrotation}{\symglyph{Rrotation}_{\symglyph{theta}}}$ and $\like{Wcompetitor}{\symglyph{Wcompetitor}_{\symglyph{theta}}}$ above. Realification doubles
the complex singular-value multiplicities. Thus the real small-angle
restriction in the full-displacement claim
\citep[Proposition~4.4]{davis1970rotation} does not eliminate the mechanism of the
authors' preceding example.

\citet{bolotnikov2001minimal} studied related unitary completion and
subspace-mapping problems. Their results give a $\sym{NormQ}$-norm
minimization theorem and a counterexample for Schatten $p$-norms with
$1\leq p<2$ \citep[Theorem~3.1 and Lemma~3.3]{bolotnikov2001minimal}.
Those results establish the earlier norm-class
distinction. \citet[Section~5]{qiu2006grassmann} subsequently repeated
the $\sym{PiConstant}/3$ assertion, whereas
\citet[Section~4, p.~1615]{zhang2010canonicalangles} explicitly noted
general trace-norm nonminimality without addressing that particular
real-space restriction. The present refutation identifies the invalid
inequality in the printed full-displacement claim
\citep[Proposition~4.4]{davis1970rotation}. The Q-norm conclusion above follows
from the earlier squared-displacement extremality statement.


\section{A follow-on theorem enabled by the formal development}
\label{sec:example-transport}

We prove an exact formula for the smallest trace-norm displacement of a real
orthogonal map carrying one given subspace onto another in finite dimension.
The formula reduces this optimization to maximum-weight matching on a finite
general graph whose vertices label the positive principal angles.
Its proof uses the principal-angle geometry, real and complex operator
  constructions, and singular-value norms developed for the formalization,
  together with additional optimization arguments.

Let $\sym{MatchingSourceSubspace},\sym{MatchingTargetSubspace}\subseteq\sym{RealField}^n$ have the same dimension and let
$\like{theta}{\symglyph{theta}_1},\ldots,\like{theta}{\symglyph{theta}_m}\in(0,\sym{PiConstant}/2]$ be their positive principal angles,
counted with multiplicity. A matching $\sym{Matching}$ chooses unordered
pairs of labels from $\{1,\ldots,m\}$, using each label at most once.
Labels that are not used in a pair are unmatched.
For three labels, $\sym{Matching}=\{\{1,2\}\}$ pairs labels $1,2$ and
leaves $3$ unmatched. The choice $\{\{1,2\},\{2,3\}\}$ is disallowed:
label $2$ would occur in two pairs.

Each positive principal angle describes a source line and a target line
in a two-dimensional principal plane. Transporting that plane independently
allows a rotation through $\like{theta}{\symglyph{theta}_i}$, with trace displacement
$4\sin(\like{theta}{\symglyph{theta}_i}/2)$, or a reflection carrying the source line
to the target, with trace displacement $2$. Pairing two labels $\{i,j\}$
allows a single orthogonal map to act on both principal planes together,
on their four-dimensional sum.
For each possible pair of distinct labels $i,j$, define its paired cost
$\like{PairCost}{\symglyph{PairCost}_{ij}}$; for each unmatched label $i$, define its singleton
cost $\like{SingletonCost}{\symglyph{SingletonCost}_i}$:
\begin{equation}\label{eq:nuclear-matching-costs}
 \like{PairCost}{\symglyph{PairCost}_{ij}}=4\sin\bigl((\like{theta}{\symglyph{theta}_i}+\like{theta}{\symglyph{theta}_j})/2\bigr),\qquad
 \like{SingletonCost}{\symglyph{SingletonCost}_i}=\min\{4\sin(\like{theta}{\symglyph{theta}_i}/2),2\}.
\end{equation}
The paired cost is attained by an orthogonal transporter on the two-plane sum.
It fixes two real directions and rotates the other two through
$\like{theta}{\symglyph{theta}_i}+\like{theta}{\symglyph{theta}_j}$, as the competitor does for
equal angles in \cref{fig:prop44-displacement-planes}. Step~6 below derives
this construction and its cost from the block minimum.
Write $\sym{OrthogonalGroup}$ for the orthogonal group of real
$n\times n$ matrices, defined by $\sym{Wcompetitor}^T\sym{Wcompetitor}=\sym{Identity}$.

% Note:
% this corresponds to exists_minimizing_real_nuclear_transport_matching
% in ../../DavisKahan/FiniteDimensional/DirectRotation/RealNuclearMatching.lean
%  It states that an orthogonal transporter attains the matching cost and
%  minimizes trace displacement over every orthogonal competitor.
%  The matching costs are defined in RealMatchingExponent.lean realNuclearMatchingCost
%  (DavisKahan/FiniteDimensional/DirectRotation/RealMatchingExponent.lean:136).
%  The endpoint proof uses the 1<p<2 theorem in RealMatchingAttainment.lean:296
%  (DavisKahan/FiniteDimensional/DirectRotation/RealMatchingAttainment.lean:296)
%  exists_minimizing_real_schatten_transport_matching
\begin{theorem}[Real trace-norm matching reduction]%
\label{thm:real-matching-example}
For equal-dimensional real subspaces $\sym{MatchingSourceSubspace},\sym{MatchingTargetSubspace}$ as above,
\begin{equation}\label{eq:real-nuclear-match}
\min_{\substack{\sym{Wcompetitor}\in \sym{OrthogonalGroup}\\\sym{Wcompetitor}(\sym{MatchingSourceSubspace})=\sym{MatchingTargetSubspace}}}\|\sym{Identity}-\sym{Wcompetitor}\|_1
=
\min_{\sym{Matching}\text{ a matching}}
\left(
\sum_{\{i,j\}\in \sym{Matching}}\like{PairCost}{\symglyph{PairCost}_{ij}}
+
\sum_{i\text{ unmatched}}\like{SingletonCost}{\symglyph{SingletonCost}_i}
\right).
\end{equation}
Both minima are attained.  The equality includes repeated angles,
right angles, and common ambient directions; the left side has no
principal-plane or determinant constraint.
\end{theorem}

We formalized this theorem in Lean~4, together with the corresponding
matching reduction for the Schatten power cost
$\|\sym{Identity}-\sym{Wcompetitor}\|_p^p$ for $1\le p<2$.
The exponent-dependent matching costs are developed in the proof below.
\Cref{lst:real-transport-matching} states both results.

% Source declarations (namespace TauCeti.DavisKahan.FiniteDimensional.RealTransport):
% DavisKahan/FiniteDimensional/DirectRotation/RealNuclearMatching.lean:
%   positivePrincipalAngles, realTransportCosts,
%   real_nuclear_transport_matching, real_schatten_transport_matching.
% The endpoint identity is RealMatchingExponent.lean:
%   realPositiveSchattenMatchingMinimum_one.
\noindent\begin{minipage}{\linewidth}
\begin{lstlisting}[style=leanpaper,
  caption={Lean presentation theorems for the trace-norm and Schatten
  matching reductions. The omitted context supplies a finite-dimensional
  real inner-product space \texttt{E}. The set \texttt{realTransportCosts p K L}
  contains the power costs $\|\sym{Identity}-T\|_p^p$ of orthogonal transporters
  $T$ carrying $K$ onto $L$. The function \texttt{positivePrincipalAngles}
  retains the labels of the strictly positive principal angles.},
  label={lst:real-transport-matching}]
theorem real_nuclear_transport_matching
    (K L : Submodule Real E)
    (hKL : Module.finrank Real K = Module.finrank Real L) :
    IsLeast (realTransportCosts 1 K L)
      (realNuclearMatchingMinimum
        (positivePrincipalAngles K L)) := by
  (*@\laaettr@*)

theorem real_schatten_transport_matching
    (p : Real) (hp1 : 1 <= p) (hp2 : p < 2)
    (K L : Submodule Real E)
    (hKL : Module.finrank Real K = Module.finrank Real L) :
    IsLeast (realTransportCosts p K L)
      (realPositiveSchattenMatchingMinimum p
        (positivePrincipalAngles K L)) := by
  (*@\laaettr@*)
\end{lstlisting}
\end{minipage}\par

In Lean, \texttt{IsLeast S x} states that $x$ belongs to $S$ and is less
than or equal to every element of $S$. Thus both signatures assert
attainment and global optimality. At $p=1$ the two matching minima coincide.

To express the formula as maximum-weight matching, start from the
all-singleton cost $\sym{SingletonBaseline}=\sum_i \like{SingletonCost}{\symglyph{SingletonCost}_i}$ and assign each edge $\{i,j\}$ the weight
\[
 \sym{MatchingEdgeWeight}=\like{SingletonCost}{\symglyph{SingletonCost}_i}+\like{SingletonCost}{\symglyph{SingletonCost}_j}-\like{PairCost}{\symglyph{PairCost}_{ij}}.
\]
This weight is the difference between the cost of two singletons and the
paired cost. Thus a matching $\sym{Matching}$ has cost
\[
 \like{TraceCost}{\symglyph{TraceCost}(M)}=\sym{SingletonBaseline}-
       \sum_{\{i,j\}\in\sym{Matching}}\sym{MatchingEdgeWeight}.
\]
Minimizing this cost is equivalent to maximizing the sum of edge weights.
The empty matching is allowed, and edges of negative weight can be omitted.
Figure~\ref{fig:real-nuclear-matching} shows a three-angle example.

\begin{figure}[tb]
\centering
\resizebox{119mm}{!}{\input{figures/Fig3.pgf}}
\caption{
Graph of principal-angle pairing choices for two three-dimensional subspaces
of $\sym{RealField}^6$ whose intersection is $\{\sym{ZeroVector}\}$ and whose principal angles are
$30^\circ$, $45^\circ$, and $75^\circ$.
(a) Vertex positions are arbitrary; each vertex represents an angle label and
each edge a possible pairing. Vertices carry singleton costs $\like{SingletonCost}{\symglyph{SingletonCost}_i}$;
edges carry weights $\sym{MatchingEdgeWeight}=\like{SingletonCost}{\symglyph{SingletonCost}_i}+\like{SingletonCost}{\symglyph{SingletonCost}_j}-\like{PairCost}{\symglyph{PairCost}_{ij}}$. The solid edge
pairs labels $1,2$ and leaves label $3$ unmatched. (b) The table lists all four
matchings; costs and edge weights are rounded to four decimals. Geometrically,
the selected pair is realized by a four-dimensional transporter and the
unmatched angle by a two-dimensional reflection. Their orthogonal sum
attains the minimum on $\sym{RealField}^6$, with trace displacement
$4\sin(37.5^\circ)+2\approx4.4350$}
\label{fig:real-nuclear-matching}
\end{figure}

\begin{proof}[Proof and explanation]
The trace norm, corresponding to $p=1$, need not be differentiable.
We first replace it by a differentiable displacement cost with $1<p<2$.
At a global minimizer, derivatives along constraint-preserving variations
vanish. We use these equations to decompose the minimizer into orthogonal
blocks carrying one or two principal angles. The blocks determine a matching;
conversely, every matching constructs a transporter with the same cost.
This proves equality of the two minimum values for $1<p<2$.
We then let $p$ approach $1$ from above to recover the trace norm.

\emph{Step 1: choose a differentiable cost and a global minimizer.}
For $1<p<2$, set
\[
 \sym{FeasibleTransport}=\{\sym{Wcompetitor}\in \sym{OrthogonalGroup}:
                 \sym{Wcompetitor}(\sym{MatchingSourceSubspace})=\sym{MatchingTargetSubspace}\},
 \qquad
 \like{PowerCost}{\symglyph{PowerCost}_p}(\sym{Wcompetitor})
       =\sum_{j=1}^n \sym{TransportSingularValues}(\sym{Identity}-\sym{Wcompetitor})^p.
\]
Here $\sym{TransportSingularValues}(T)$ denotes a singular value of $T$. For $p>1$ this cost is
differentiable, including at zero displacement singular values.

\emph{(a) The cost attains a global minimum on $\sym{FeasibleTransport}$.}
Write $\like{Psource}{\symglyph{Psource}_K}$ and $\sym{MatchingTargetProjection}$ for the orthogonal
projections onto the source and target.
\emph{Justification.} An orthogonal map carrying an
orthonormal basis of $\sym{MatchingSourceSubspace}$ to one of $\sym{MatchingTargetSubspace}$,
extended to their orthogonal complements, belongs to $\sym{FeasibleTransport}$.
Thus $\sym{FeasibleTransport}$ is nonempty. The group $\sym{OrthogonalGroup}$ is closed and bounded
in $\sym{RealField}^{n\times n}$, hence compact by Heine--Borel.
The equivalent transport equation
$\sym{Wcompetitor}\like{Psource}{\symglyph{Psource}_K}\sym{Wcompetitor}^T
 =\sym{MatchingTargetProjection}$ is a closed constraint.
Since $\like{PowerCost}{\symglyph{PowerCost}_p}$ is continuous, the extreme value theorem
gives a global minimizer $\sym{Wcompetitor}\in\sym{FeasibleTransport}$.

\emph{(b) A simultaneous invariant decomposition splits both the cost
and the principal-angle list.}
The decomposition we seek is an orthogonal sum
$\sym{RealField}^n=\bigoplus_a\sym{MatchingBlock}$ in which each $\sym{MatchingBlock}$ is preserved by
$\like{Psource}{\symglyph{Psource}_K}$, $\sym{MatchingTargetProjection}$, and $\sym{Wcompetitor}$.
\emph{Justification.} In a basis adapted to this sum all three matrices are block diagonal.
Consequently, the source and target split as the sums of their
intersections with the $\sym{MatchingBlock}$, and
\begin{equation}\label{eq:matching-block-additivity}
 \like{PowerCost}{\symglyph{PowerCost}_p}(\sym{Wcompetitor})
 =\sum_a\like{PowerCost}{\symglyph{PowerCost}_p}(\sym{Wcompetitor}|_{\sym{MatchingBlock}}).
\end{equation}
Indeed, the displacement is block diagonal too, and its singular-value
list is the union of the block singular-value lists. The principal-angle
lists also combine in this way: the compression
$\like{Psource}{\symglyph{Psource}_K}\sym{MatchingTargetProjection}\like{Psource}{\symglyph{Psource}_K}$
is block diagonal, and its restriction to $\sym{MatchingSourceSubspace}$ has the
squared principal cosines as eigenvalues.
These two additivity facts will turn the block decomposition into a
matching of the original angle labels.

\emph{Step 2: use derivatives along maps that still transport the source.}
Consider smooth curves $\sym{Wcompetitor}(t)\in\sym{FeasibleTransport}$, for real $t$
near zero, with $\sym{Wcompetitor}(0)=\sym{Wcompetitor}$.
The cost derivative at zero vanishes along every such curve because
$\sym{Wcompetitor}$ is a minimizer. This is the stationarity condition.
We use it as a necessary condition on the global minimizer already
obtained in Step~1.

\emph{(a) Source-preserving rotations give feasible curves.}
Let $\sym{SkewDirection}$ be skew-symmetric (\ie{},
$\sym{SkewDirection}^T=-\sym{SkewDirection}$).
\emph{Justification.} For real $t$, its exponential is orthogonal, since
\[
 \exp(t\sym{SkewDirection})^T
 =\exp(-t\sym{SkewDirection})
 =\exp(t\sym{SkewDirection})^{-1}.
\]
If $\sym{SkewDirection}$ commutes with $\like{Psource}{\symglyph{Psource}_K}$, so does its
exponential (every power of $\sym{SkewDirection}$ commutes with
$\like{Psource}{\symglyph{Psource}_K}$, hence so does the exponential power series).
Its inverse $\exp(-t\sym{SkewDirection})$ also commutes with
$\like{Psource}{\symglyph{Psource}_K}$, so the exponential maps the source onto itself.
Consequently,
$\sym{Wcompetitor}\exp(t\sym{SkewDirection})$ remains in $\sym{FeasibleTransport}$.
This gives explicit curves on which to test stationarity.

\emph{(b) The first variation is a trace pairing with a spectral gradient.}
Define the scalar phase cost
\[
 \sym{MatchingPhaseCost}(t)=(2\sin(t/2))^p,\qquad 0\le t\le\sym{PiConstant}.
\]
A planar rotation $\sym{MatchingPlanarRotation}$ through $t$ satisfies
\[
 (\sym{Identity}-\sym{MatchingPlanarRotation})^T(\sym{Identity}-\sym{MatchingPlanarRotation})
   =2\sym{Identity}-\sym{MatchingPlanarRotation}-\sym{MatchingPlanarRotation}^T
   =4\sin^2(t/2)\sym{Identity}.
\]
Its two displacement singular values are therefore both $2\sin(t/2)$,
and its contribution is $2\sym{MatchingPhaseCost}(t)$.
On a plane of rotation angle $\sym{MatchingPlanarAngle}$, use the planar rotation matrix through $\sym{PiConstant}/2$
$\sym{Jturn}=\bigl(\begin{smallmatrix}0&-1\\1&0\end{smallmatrix}\bigr)$
and define the skew-symmetric spectral gradient by
\[
 \sym{Wcompetitor}|_{\mathrm{plane}}
   =\cos\sym{MatchingPlanarAngle}\,\sym{IdentityTwo}+\sin\sym{MatchingPlanarAngle}\,\sym{Jturn},
 \qquad
 \sym{SpectralGradient}|_{\mathrm{plane}}
   =\sym{MatchingPhaseCost}'(\sym{MatchingPlanarAngle})\sym{Jturn}.
\]
On the $+1$ and $-1$ eigenspaces of the transporter, set the spectral
gradient $\sym{SpectralGradient}$ to zero.
Every real orthogonal matrix splits into these scalar blocks and planar
rotations: diagonalize over $\sym{ComplexField}$ and pair conjugate eigenvectors
to obtain the real planes. This defines $\sym{SpectralGradient}$ on the
whole space. On each plane the displayed blocks are combinations of
$\sym{IdentityTwo}$ and the same $\sym{Jturn}$, so they commute; on the scalar blocks
the gradient is zero. Thus $\sym{SpectralGradient}\sym{Wcompetitor}
=\sym{Wcompetitor}\sym{SpectralGradient}$ on the whole space.
These rotation planes may mix source and complementary directions;
stationarity will relate them to $\like{Psource}{\symglyph{Psource}_K}$.
Then for every skew-symmetric $\sym{SkewDirection}$,
\begin{equation}\label{eq:matching-first-variation}
 \left.\frac{d}{dt}
   \like{PowerCost}{\symglyph{PowerCost}_p}(\sym{Wcompetitor}\exp(t\sym{SkewDirection}))
 \right|_{t=0}
 =-\sym{TraceOperator}(\sym{SpectralGradient}\sym{SkewDirection}).
\end{equation}
\emph{Justification.} For $p>1$ the cost is differentiable, including at
zero displacement singular values, and its derivative is linear in
$\sym{SkewDirection}$. Change the sign on any whole rotation plane or
scalar line. This orthogonal sign matrix commutes with
$\sym{Wcompetitor}$, so conjugating $\sym{SkewDirection}$ by it conjugates
the whole curve and leaves its cost unchanged. Average over these sign
choices. Linearity of the derivative leaves its value unchanged, while
every off-diagonal block of $\sym{SkewDirection}$ cancels. This averaging
computes the derivative for arbitrary skew directions; feasibility is
imposed afterward, when we apply stationarity.

On a rotation plane of angle $\sym{MatchingPlanarAngle}$, the diagonal skew block is
$b\sym{Jturn}$. Since
$\exp(tb\sym{Jturn})$ is a rotation through $tb$, the phase becomes
$\sym{MatchingPlanarAngle}+tb$ and the cost derivative is
$2b\sym{MatchingPhaseCost}'(\sym{MatchingPlanarAngle})$. The trace pairing gives exactly
the same value:
\[
 -\sym{TraceOperator}\bigl(
   (\sym{MatchingPhaseCost}'(\sym{MatchingPlanarAngle})\sym{Jturn})(b\sym{Jturn})
  \bigr)
 =-b\sym{MatchingPhaseCost}'(\sym{MatchingPlanarAngle})\sym{TraceOperator}(-\sym{IdentityTwo})
 =2b\sym{MatchingPhaseCost}'(\sym{MatchingPlanarAngle}).
\]
On a scalar line the diagonal skew-symmetric block is zero.
Summing proves \cref{eq:matching-first-variation}, including directions
that mix the rotation planes.

\emph{(c) Stationarity forces the gradient to couple only source and
complementary directions.}
\emph{Justification.} Take the particular direction
$\sym{SkewDirection}=\like{Psource}{\symglyph{Psource}_K}\sym{SpectralGradient}\like{Psource}{\symglyph{Psource}_K}$.
It is skew-symmetric and commutes with $\like{Psource}{\symglyph{Psource}_K}$, so stationarity
and cyclicity of trace give
\[
 0=-\sym{TraceOperator}\bigl(
      (\like{Psource}{\symglyph{Psource}_K}\sym{SpectralGradient}\like{Psource}{\symglyph{Psource}_K})^2
     \bigr)
   =\|\like{Psource}{\symglyph{Psource}_K}\sym{SpectralGradient}\like{Psource}{\symglyph{Psource}_K}\|_F^2.
\]
The source diagonal block of the gradient is therefore zero.
Using $\sym{Identity}-\like{Psource}{\symglyph{Psource}_K}$ instead gives the same conclusion for the
complementary diagonal block. Relative to
$\sym{MatchingSourceSubspace}\oplus\sym{MatchingSourceSubspace}^{\perp}$, the gradient and source
reflection thus have the forms
\[
 \sym{SpectralGradient}
  =\begin{pmatrix}0&\sym{GradientOffDiagonal}\\-\like{GradientOffDiagonal}{\symglyph{GradientOffDiagonal}^T}&0\end{pmatrix},
 \qquad
 \sym{Jreflect}=2\like{Psource}{\symglyph{Psource}_K}-\sym{Identity}
  =\begin{pmatrix}\sym{Identity}&0\\0&-\sym{Identity}\end{pmatrix}.
\]
Multiplication gives
$\sym{Jreflect}\sym{SpectralGradient}\sym{Jreflect}
 =-\sym{SpectralGradient}$.
This identity is what makes the source projection compatible with the
decomposition constructed next.

\emph{Step 3: obtain blocks with one or two source directions.}
The rotation planes used to define the gradient need not preserve the
source. We now use the relation between the gradient and the source
reflection to construct blocks that preserve the source as well as the
minimizer. Their size will explain why the final optimization needs only
single labels and pairs.
\emph{(a) Each positive gradient level (an eigenspace of
$-\sym{SpectralGradient}^2$ with positive eigenvalue) admits a complex structure.}
\emph{Justification.}
The matrix $-\like{SpectralGradient}{\symglyph{SpectralGradient}^2}$ is symmetric and nonnegative,
since
$\langle-\like{SpectralGradient}{\symglyph{SpectralGradient}^2}x,x\rangle
  =\|\sym{SpectralGradient}x\|^2$.
Decompose the space into its eigenspaces. The gradient, transporter,
and source reflection all preserve them. On an eigenspace $\sym{GradientLevelSpace}$
with positive eigenvalue $\like{GradientLevel}{\symglyph{GradientLevel}^2}$, put
$\sym{GradientComplexStructure}=\sym{SpectralGradient}/\sym{GradientLevel}$.
It satisfies $\sym{GradientComplexStructure}^2=-\sym{Identity}$ and preserves lengths.
We can therefore treat $\sym{GradientLevelSpace}$ as a complex space by defining
\[
 (a+ib)x=ax+b\sym{GradientComplexStructure}x.
\]
A real plane spanned by $x,\sym{GradientComplexStructure}x$ becomes one
complex line. Commutation with the gradient makes the transporter
complex-linear and unitary in these coordinates.
The source reflection reverses the complex structure:
\[
 \sym{Jreflect}\sym{GradientComplexStructure}
   =-\sym{GradientComplexStructure}\sym{Jreflect}.
\]
Consequently
$\sym{Jreflect}((a+ib)x)=(a-ib)\sym{Jreflect}x$.
Together with $\sym{Jreflect}^2=\sym{Identity}$, this says that it acts as conjugation.

\emph{(b) The transporter has at most two phases on each positive level.}
\emph{Justification.} The definition of the gradient says that an eigenphase $t$ of the
transporter on this complex space satisfies
$\sym{MatchingPhaseCost}'(t)=\sym{GradientLevel}$, with $0<t<\sym{PiConstant}$.
There are at most two such phases. To check this directly,
\[
 \sym{MatchingPhaseCost}'(t)
   =p(2\sin(t/2))^{p-1}\cos(t/2),
 \qquad
 \sym{MatchingPhaseCost}''(t)
   =p(2\sin(t/2))^{p-2}
       \bigl((p-1)\cos^2(t/2)-\sin^2(t/2)\bigr).
\]
The second derivative changes sign once, where
$\tan^2(t/2)=p-1$. The first derivative therefore increases and then
decreases, so a positive horizontal level has at most two intersections.
Writing the two phases as $\like{GradientEigenphase}{\symglyph{GradientEigenphase}_1},\like{GradientEigenphase}{\symglyph{GradientEigenphase}_2}$, with equality allowed,
we obtain
\[
 \sym{Wcompetitor}|_{\sym{GradientLevelSpace}}
   =e^{i\like{GradientEigenphase}{\symglyph{GradientEigenphase}_1}}\sym{GradientPhaseProjection}+e^{i\like{GradientEigenphase}{\symglyph{GradientEigenphase}_2}}(\sym{Identity}-\sym{GradientPhaseProjection}),
\]
where $\sym{GradientPhaseProjection}$ is the orthogonal projection onto the first phase space.
A subspace preserved by $\sym{GradientPhaseProjection}$ is preserved by this transporter.
To preserve the source projection as well, we require invariance under
the conjugation $\sym{Jreflect}$.

\emph{(c) The phase projection and conjugation preserve common blocks of
real dimension at most four.}
\emph{Justification.} Choose complex coordinates in which $\sym{Jreflect}$ is coordinate
conjugation and write $\sym{GradientPhaseProjection}=\sym{ProjectionRealPart}+i\sym{ProjectionImaginaryPart}$ with real matrices $\sym{ProjectionRealPart},\sym{ProjectionImaginaryPart}$.
The equations $\sym{GradientPhaseProjection}^*=\sym{GradientPhaseProjection}$ and $\sym{GradientPhaseProjection}^2=\sym{GradientPhaseProjection}$ give
\[
 \sym{ProjectionRealPart}^T=\sym{ProjectionRealPart},\qquad \sym{ProjectionImaginaryPart}^T=-\sym{ProjectionImaginaryPart},\qquad
 \sym{ProjectionRealPart}^2-\sym{ProjectionImaginaryPart}^2=\sym{ProjectionRealPart},\qquad \sym{ProjectionRealPart}\sym{ProjectionImaginaryPart}+\sym{ProjectionImaginaryPart}\sym{ProjectionRealPart}=\sym{ProjectionImaginaryPart}.
\]
For real coordinate vectors $x$,
$\langle \sym{ProjectionRealPart}x,x\rangle=\operatorname{Re}\langle \sym{GradientPhaseProjection}x,x\rangle$
lies between zero and $\|x\|^2$, so the eigenvalues of $\sym{ProjectionRealPart}$ lie in $[0,1]$.
If $\sym{ProjectionRealPart}u=\sym{ProjectionEigenvalue} u$, the last two identities imply
\[
 \sym{ProjectionRealPart}(\sym{ProjectionImaginaryPart}u)=(1-\sym{ProjectionEigenvalue})\sym{ProjectionImaginaryPart}u,\qquad
 \sym{ProjectionImaginaryPart}^2u=-\sym{ProjectionEigenvalue}(1-\sym{ProjectionEigenvalue})u.
\]
For a unit $u$ with $0<\sym{ProjectionEigenvalue}<1$, set
$\sym{ProjectionCoupling}=\sqrt{\sym{ProjectionEigenvalue}(1-\sym{ProjectionEigenvalue})}$ and $v=\sym{ProjectionImaginaryPart}u/\sym{ProjectionCoupling}$.
Skew-symmetry gives
\[
 \langle u,\sym{ProjectionImaginaryPart}u\rangle=0,\qquad
 \|\sym{ProjectionImaginaryPart}u\|^2
   =-\langle u,\sym{ProjectionImaginaryPart}^2u\rangle
   =\sym{ProjectionEigenvalue}(1-\sym{ProjectionEigenvalue})
   =\sym{ProjectionCoupling}^2.
\]
Thus $u,v$ are orthonormal. The preceding identities now become
\[
 \sym{ProjectionRealPart}u=\sym{ProjectionEigenvalue} u,\quad \sym{ProjectionRealPart}v=(1-\sym{ProjectionEigenvalue})v,\quad
 \sym{ProjectionImaginaryPart}u=\sym{ProjectionCoupling} v,\quad \sym{ProjectionImaginaryPart}v=-\sym{ProjectionCoupling} u.
\]
Thus their real span is preserved by both $\sym{ProjectionRealPart}$ and $\sym{ProjectionImaginaryPart}$.
This includes $\sym{ProjectionEigenvalue}=1/2$. At $\sym{ProjectionEigenvalue}=0$ or $1$, $\sym{ProjectionImaginaryPart}u=0$ and a
single eigenvector gives an invariant real line.
By symmetry of $\sym{ProjectionRealPart}$ and skew-symmetry of $\sym{ProjectionImaginaryPart}$, the orthogonal complement
of each chosen line or plane is again preserved by $\sym{ProjectionRealPart},\sym{ProjectionImaginaryPart}$.
The construction can be repeated until the real
coordinate space is exhausted.

The complex spans of these lines and planes are invariant under
$\sym{GradientPhaseProjection},\sym{Jreflect}$, and $\sym{Wcompetitor}$. Returning to the original
real space gives blocks of dimension two or four. Each is also
preserved by
$\like{Psource}{\symglyph{Psource}_K}=(\sym{Jreflect}+\sym{Identity})/2$, hence by
$\sym{MatchingTargetProjection}
 =\sym{Wcompetitor}\like{Psource}{\symglyph{Psource}_K}\sym{Wcompetitor}^T$.
For example, the complex span of $u,v$ is the real span of
$u,v,\sym{GradientComplexStructure}u,\sym{GradientComplexStructure}v$.
The first two vectors are fixed by conjugation and hence belong to the
source; conjugation negates the last two, which belong to its complement.
This gives a four-dimensional block with two source directions.
An invariant real coordinate line similarly gives a two-dimensional block
with one source direction. Their principal-angle lists therefore have
two and one entries, respectively.

\emph{(d) The zero-gradient space contributes only fixed directions and
one-angle planes.}
\emph{Justification.} On $\ker\sym{SpectralGradient}$ the transporter has eigenvalues $+1,-1$,
because these are exactly the phases where the scalar derivative vanishes.
It is therefore a reflection there. Apply the two-projection decomposition
to $\like{Psource}{\symglyph{Psource}_K}$ and the projection onto its $-1$ eigenspace:
this gives real lines and planes preserved by both projections and the
transporter. On an invariant line the source and target are either both
zero or both the whole line. Replacing a one-dimensional $-1$ block by
$+1$ preserves them and decreases the cost by $2^p$. Such a block cannot
occur at a global minimum. The remaining blocks are fixed directions
or planes with one source direction.
This completes the decomposition needed for
\cref{eq:matching-block-additivity}.

\emph{Step 4: calculate the costs on a plane and on two planes.}
\emph{(a) The optimal map on a one-angle block is a rotation or a reflection.}
\emph{Justification.} For a principal angle $\like{theta}{\symglyph{theta}_i}$, the least rotation cost
is $2\sym{MatchingPhaseCost}(\like{theta}{\symglyph{theta}_i})$.
The alternative rotation sends a source unit vector to the opposite
target unit vector and turns through $\sym{PiConstant}-\like{theta}{\symglyph{theta}_i}$,
which is at least $\like{theta}{\symglyph{theta}_i}$.
A planar reflection has eigenvalues $+1,-1$, hence displacement singular values $0,2$,
so its cost is $2^p$. The one-angle minimum is
\[
 \like{SingletonPowerCost}{\symglyph{SingletonPowerCost}_i(p)}
   =\min\{2\sym{MatchingPhaseCost}(\like{theta}{\symglyph{theta}_i}),\,2^p\}.
\]

\emph{(b) A positive-determinant two-angle block reduces to a scalar minimum.}
\emph{Justification.} For two angles, relabel within the pair so that
$\like{theta}{\symglyph{theta}_i}\ge\like{theta}{\symglyph{theta}_j}$, and set
\[
 \sym{PairMeanAngle}=(\like{theta}{\symglyph{theta}_i}+\like{theta}{\symglyph{theta}_j})/2,\qquad
 \sym{PairHalfDifference}=(\like{theta}{\symglyph{theta}_i}-\like{theta}{\symglyph{theta}_j})/2.
\]
Then $0\le \sym{PairHalfDifference}\le \sym{PairMeanAngle}$ and $\sym{PairMeanAngle}+\sym{PairHalfDifference}\le\sym{PiConstant}/2$.
In principal coordinates the direct rotation is
\[
 \sym{PairDirectRotation}=\begin{pmatrix}\sym{PairCosineBlock}&-\sym{PairSineBlock}\\\sym{PairSineBlock}&\sym{PairCosineBlock}\end{pmatrix},\qquad
 \sym{PairCosineBlock}=\sym{DiagonalOperator}(\cos\like{theta}{\symglyph{theta}_i},
                      \cos\like{theta}{\symglyph{theta}_j}),\quad
 \sym{PairSineBlock}=\sym{DiagonalOperator}(\sin\like{theta}{\symglyph{theta}_i},
                      \sin\like{theta}{\symglyph{theta}_j}).
\]
Every orthogonal transporter on this four-dimensional space is
$\sym{PairDirectRotation}\sym{DiagonalOperator}(\like{PairFreeOrthogonal}{\symglyph{PairFreeOrthogonal}_0},\like{PairFreeOrthogonal}{\symglyph{PairFreeOrthogonal}_1})$ with $\like{PairFreeOrthogonal}{\symglyph{PairFreeOrthogonal}_0},\like{PairFreeOrthogonal}{\symglyph{PairFreeOrthogonal}_1}\in \sym{PlanarOrthogonalGroup}$:
multiplying by $\sym{PairDirectRotation}^T$ gives an orthogonal map preserving the source
and its complement. This parametrization allows mixing between the
principal planes.

For determinant $+1$, the free factors are either both rotations or both
reflections. Write their rotation angles as
$\sym{PairRotationR}+\sym{PairRotationS}$ and
$\sym{PairRotationR}-\sym{PairRotationS}$: the parameters
$\sym{PairRotationR},\sym{PairRotationS}$ are the half-sum and half-difference
of those angles. Define $\sym{PairPhaseX},\sym{PairPhaseY}\in[0,\sym{PiConstant}]$ by
\[
 \begin{array}{c|cc}
  &\cos\sym{PairPhaseX}&\cos\sym{PairPhaseY}\\ \hline
  \text{two rotations}
    &\cos\sym{PairMeanAngle}\cos\sym{PairRotationS}
    &\cos\sym{PairHalfDifference}\cos\sym{PairRotationR}\\
  \text{two reflections}
    &-\sin\sym{PairHalfDifference}\cos\sym{PairRotationR}
    &\sin\sym{PairMeanAngle}\cos\sym{PairRotationS}
 \end{array}
\]
Here a reflection factor is a planar rotation followed by
$\sym{PairPlanarReflection}=\sym{DiagonalOperator}(1,-1)$; for example,
$\like{PairFreeOrthogonal}{\symglyph{PairFreeOrthogonal}_0}
 =\sym{MatchingPlanarRotationSum}\sym{PairPlanarReflection}$.
The cosine bounds imply
\[
 \sym{PairMeanAngle}\le \sym{PairPhaseX}\le\sym{PiConstant}-\sym{PairMeanAngle},\qquad
 \sym{PairHalfDifference}\le \sym{PairPhaseY}\le\sym{PiConstant}-\sym{PairHalfDifference}.
\]
For reflections, the initial bounds are
$\sym{PiConstant}/2-\sym{PairHalfDifference}\le\sym{PairPhaseX}\le\sym{PiConstant}/2+\sym{PairHalfDifference}$
and $\sym{PiConstant}/2-\sym{PairMeanAngle}\le\sym{PairPhaseY}\le\sym{PiConstant}/2+\sym{PairMeanAngle}$;
these lie in the displayed ranges because
$\sym{PairMeanAngle}+\sym{PairHalfDifference}\le\sym{PiConstant}/2$.

To check that these parameters give the actual eigenphases, put
$\sym{PairPhaseCosineX}=\cos\sym{PairPhaseX}$ and
$\sym{PairPhaseCosineY}=\cos\sym{PairPhaseY}$. Block multiplication gives
\begin{align*}
 \sym{TraceOperator}\sym{Wcompetitor}
   &=\sym{TraceOperator}(\sym{PairCosineBlock}\like{PairFreeOrthogonal}{\symglyph{PairFreeOrthogonal}_0})
     +\sym{TraceOperator}(\sym{PairCosineBlock}\like{PairFreeOrthogonal}{\symglyph{PairFreeOrthogonal}_1})
     =4\sym{PairPhaseCosineX}\sym{PairPhaseCosineY},\\
 \sym{TraceOperator}(\sym{Wcompetitor}^2)
   &=\sym{TraceOperator}((\sym{PairCosineBlock}\like{PairFreeOrthogonal}{\symglyph{PairFreeOrthogonal}_0})^2)
     +\sym{TraceOperator}((\sym{PairCosineBlock}\like{PairFreeOrthogonal}{\symglyph{PairFreeOrthogonal}_1})^2)
     -2\sym{TraceOperator}(\sym{PairSineBlock}\like{PairFreeOrthogonal}{\symglyph{PairFreeOrthogonal}_1}
                        \sym{PairSineBlock}\like{PairFreeOrthogonal}{\symglyph{PairFreeOrthogonal}_0})\\
   &=16\sym{PairPhaseCosineX}^2\sym{PairPhaseCosineY}^2
     -8(\sym{PairPhaseCosineX}^2+\sym{PairPhaseCosineY}^2)+4.
\end{align*}
The coefficient of $z^2$ in the characteristic polynomial is
$((\sym{TraceOperator}\sym{Wcompetitor})^2
   -\sym{TraceOperator}(\sym{Wcompetitor}^2))/2$.
Orthogonality and determinant one make this polynomial reciprocal, so
\begin{align*}
 \det(zI-\sym{Wcompetitor})
   &=z^4-4\sym{PairPhaseCosineX}\sym{PairPhaseCosineY}z^3
     +\bigl(4(\sym{PairPhaseCosineX}^2+\sym{PairPhaseCosineY}^2)-2\bigr)z^2
     -4\sym{PairPhaseCosineX}\sym{PairPhaseCosineY}z+1\\
   &=\prod_{\sym{PhaseSignChoice}=\pm1}
       \bigl(z^2-2\cos(\sym{PairPhaseX}+\sym{PhaseSignChoice}\sym{PairPhaseY})z+1\bigr).
\end{align*}
Its roots are therefore
$e^{\pm i(\sym{PairPhaseX}+\sym{PairPhaseY})},
 e^{\pm i(\sym{PairPhaseX}-\sym{PairPhaseY})}$.
Folding their angles into $[0,\sym{PiConstant}]$ gives the displacement cost
\[
 2\bigl(\sym{MatchingPhaseCost}(\min\{\sym{PairPhaseX}+\sym{PairPhaseY},
                            2\sym{PiConstant}-\sym{PairPhaseX}-\sym{PairPhaseY}\})
       +\sym{MatchingPhaseCost}(|\sym{PairPhaseX}-\sym{PairPhaseY}|)\bigr).
\]

We can replace $(\sym{PairPhaseX},\sym{PairPhaseY})$ by $(\sym{PiConstant}-\sym{PairPhaseX},\sym{PiConstant}-\sym{PairPhaseY})$ when $\sym{PairPhaseX}+\sym{PairPhaseY}>\sym{PiConstant}$ without
changing this cost. Suppose therefore that $\sym{PairPhaseX}+\sym{PairPhaseY}\le\sym{PiConstant}$.
To replace the first parameter by $\sym{PairMeanAngle}$ without increasing the smaller
phase, we need $\sym{PairMeanAngle}-\sym{PairPhaseParameter}\le|\sym{PairPhaseX}-\sym{PairPhaseY}|$, as well as $\sym{PairPhaseParameter}\ge \sym{PairHalfDifference}$.
Choose the smallest such parameter,
$\sym{PairPhaseParameter}=\max\{\sym{PairHalfDifference},\sym{PairMeanAngle}-|\sym{PairPhaseX}-\sym{PairPhaseY}|\}$.
The identities $\sym{PairPhaseX}+\sym{PairPhaseY}+|\sym{PairPhaseX}-\sym{PairPhaseY}|=2\max\{\sym{PairPhaseX},\sym{PairPhaseY}\}\ge2\sym{PairMeanAngle}$ and $\sym{PairPhaseX}+\sym{PairPhaseY}\ge \sym{PairMeanAngle}+\sym{PairHalfDifference}$
also give $\sym{PairMeanAngle}+\sym{PairPhaseParameter}\le \sym{PairPhaseX}+\sym{PairPhaseY}$. Thus
\[
 \sym{PairHalfDifference}\le\sym{PairPhaseParameter}\le \sym{PairMeanAngle},\qquad
 \sym{PairMeanAngle}+\sym{PairPhaseParameter}\le \sym{PairPhaseX}+\sym{PairPhaseY},\qquad \sym{PairMeanAngle}-\sym{PairPhaseParameter}\le|\sym{PairPhaseX}-\sym{PairPhaseY}|.
\]
Since $\sym{MatchingPhaseCost}$ is increasing on $[0,\sym{PiConstant}]$, this
transporter costs at least
$2\bigl(\sym{MatchingPhaseCost}(\sym{PairMeanAngle}+\sym{PairPhaseParameter})+
        \sym{MatchingPhaseCost}(\sym{PairMeanAngle}-\sym{PairPhaseParameter})\bigr)$.
Every $\sym{PairPhaseParameter}\in[\sym{PairHalfDifference},\sym{PairMeanAngle}]$ is attained: take
$\like{PairFreeOrthogonal}{\symglyph{PairFreeOrthogonal}_0}=\like{PairFreeOrthogonal}{\symglyph{PairFreeOrthogonal}_1}=\sym{MatchingPlanarRotationGamma}$, where
$\sym{PairAttainingAngle}=\arccos(\cos\sym{PairPhaseParameter}/\cos \sym{PairHalfDifference})$.
The denominator is positive since $\sym{PairHalfDifference}\le\sym{PiConstant}/4$.
The preceding phase equations give $\sym{PairPhaseX}=\sym{PairMeanAngle},\sym{PairPhaseY}=\sym{PairPhaseParameter}$.
Thus the positive-determinant pair minimum is
\begin{equation}\label{eq:matching-pair-power-cost}
 \like{PairPowerCost}{\symglyph{PairPowerCost}_{ij}(p)}
 =\min_{\sym{PairPhaseParameter}\in[\sym{PairHalfDifference},\sym{PairMeanAngle}]}
     2\bigl(\sym{MatchingPhaseCost}(\sym{PairMeanAngle}+\sym{PairPhaseParameter})+
             \sym{MatchingPhaseCost}(\sym{PairMeanAngle}-\sym{PairPhaseParameter})\bigr).
\end{equation}

\emph{(c) The unrestricted two-angle minimum also compares two singletons.}
\emph{Justification.} For determinant $-1$, the eigenvalues are $+1,-1$ and a conjugate
rotational pair $e^{\pm it}$, with $0\le t\le\sym{PiConstant}$.
One free planar factor is a rotation and the other a reflection, so
the absolute trace is at most
$|\cos\like{theta}{\symglyph{theta}_i}+\cos\like{theta}{\symglyph{theta}_j}|
 +|\cos\like{theta}{\symglyph{theta}_i}-\cos\like{theta}{\symglyph{theta}_j}|
 =2\cos\like{theta}{\symglyph{theta}_j}$; thus
$|\sym{TraceOperator}\sym{Wcompetitor}|\le2\cos\like{theta}{\symglyph{theta}_j}$.
Since this trace is $2\cos t$, the rotational phase satisfies
$t\ge\like{theta}{\symglyph{theta}_j}$.
Its least cost is therefore
$2^p+2\sym{MatchingPhaseCost}(\like{theta}{\symglyph{theta}_j})$.
It is attained by reflecting on the larger-angle plane and rotating
on the smaller-angle plane.
Taking $\sym{PairPhaseParameter}=\sym{PairHalfDifference}$ bounds the positive-determinant minimum by the cost
of two rotations. Taking $\sym{PairPhaseParameter}=\sym{PairMeanAngle}$ bounds it by
$2\sym{MatchingPhaseCost}(\like{theta}{\symglyph{theta}_i}+\like{theta}{\symglyph{theta}_j})
 \le2\cdot2^p$, the cost of two reflections.
The less expensive mixed rotation/reflection choice is the
negative-determinant minimum just computed.
Comparing these four independent choices shows that the unrestricted
four-dimensional minimum is
\[
 \min\{\like{PairPowerCost}{\symglyph{PairPowerCost}_{ij}(p)},\,
         \like{SingletonPowerCost}{\symglyph{SingletonPowerCost}_i(p)}
          +\like{SingletonPowerCost}{\symglyph{SingletonPowerCost}_j(p)}\}.
\]
A four-dimensional block therefore offers either a paired choice or
two independent singleton choices.

\emph{Step 5: prove equality of the minimum costs and recover a transporter.}
We construct a transporter from every matching and extract a matching
from a minimizing transporter, preserving the cost in both constructions.
Thus solving the matching problem gives the minimum cost and an attaining
transporter, although individual transporters and matchings need not
correspond one-to-one.
For a matching $\sym{Matching}$ define
\[
 \sym{MatchingPowerCost}(\sym{Matching})
   =\sum_{\{i,j\}\in\sym{Matching}}\like{PairPowerCost}{\symglyph{PairPowerCost}_{ij}(p)}
     +\sum_{i\text{ unmatched}}\like{SingletonPowerCost}{\symglyph{SingletonPowerCost}_i(p)}.
\]
\emph{(a) Every matching gives a feasible transporter of the same cost.}
\emph{Justification.} The principal-angle decomposition gives mutually orthogonal planes
for the positive angle labels. Choose an attaining two-angle transporter for each selected pair,
an attaining planar rotation or reflection for each unmatched label,
and identity on the remaining common directions.
Their orthogonal sum is an element of $\sym{FeasibleTransport}$.
By \cref{eq:matching-block-additivity}, its cost is
$\sym{MatchingPowerCost}(\sym{Matching})$. Hence
$\min_{\sym{FeasibleTransport}}\like{PowerCost}{\symglyph{PowerCost}_p}
 \le\min_{\sym{Matching}}\sym{MatchingPowerCost}(\sym{Matching})$.

\emph{(b) A minimizing transporter gives a matching of the same cost.}
\emph{Justification.} Use the global minimizer from Step~1 and its
decomposition from Step~3. Its restriction to each block must minimize
the corresponding block problem: otherwise replacing that restriction
by a cheaper transporter would lower the global cost while retaining
the same source and target. By Step~4, a plane supplies a singleton,
and a four-dimensional block supplies a pair or two singletons.
Each occurrence of an angle belongs to one block, so no label can be used twice.
This produces a matching with the same total cost as the minimizer.

Zero angles do not supply cheaper extra partners. Their singleton cost
is zero; pairing an angle $\sym{theta}$ with zero forces $\sym{PairHalfDifference}=\sym{PairMeanAngle}=\sym{theta}/2$
in \cref{eq:matching-pair-power-cost}, giving $2\sym{MatchingPhaseCost}(\sym{theta})$,
which is at least its singleton cost.
Thus we can leave zero-angle directions fixed and use only the original
positive labels. Removing such pairs cannot strictly lower the minimizing
cost, by the construction in (a), so it preserves that cost. We have proved
\begin{equation}\label{eq:matching-power-equality}
 \min_{\sym{FeasibleTransport}}\like{PowerCost}{\symglyph{PowerCost}_p}
   =\min_{\sym{Matching}}\sym{MatchingPowerCost}(\sym{Matching}),
 \qquad 1<p<2.
\end{equation}
For the matching in \cref{fig:real-nuclear-matching}, this construction
uses one four-dimensional block for labels $1,2$ and one plane for label~$3$.

\emph{Step 6: pass to the trace norm and evaluate the weights.}
\emph{(a) The minimum values converge as $p$ approaches $1$ from above.}
\emph{Justification.} Every displacement singular value lies in $[0,2]$.
Since $\sym{DisplacementSize}^p$ converges uniformly to $\sym{DisplacementSize}$ on this interval as $p$ approaches $1$ from above,
\[
 \sup_{\sym{Wcompetitor}\in\sym{FeasibleTransport}}
   \left|\like{PowerCost}{\symglyph{PowerCost}_p}(\sym{Wcompetitor})
                  -\|\sym{Identity}-\sym{Wcompetitor}\|_1\right|
 \le n\sup_{0\le\sym{DisplacementSize}\le2}|\sym{DisplacementSize}^p-\sym{DisplacementSize}|\ \longrightarrow\ 0.
\]
The minimum values on the left of \cref{eq:matching-power-equality}
therefore converge to the minimum trace norm.
The scalar block minima also converge: the one-angle minimum is a
minimum of two continuous quantities, and the pair minimum is over
the fixed compact interval $[\sym{PairHalfDifference},\sym{PairMeanAngle}]$ with a uniformly convergent cost.
There are finitely many matchings, so their minimum values converge too.

\emph{(b) The limiting block weights are the stated trace-norm costs.}
\emph{Justification.} At $p=1$ the singleton formula gives
$\like{SingletonPowerCost}{\symglyph{SingletonPowerCost}_i(1)}
 =\min\{4\sin(\like{theta}{\symglyph{theta}_i}/2),2\}
 =\like{SingletonCost}{\symglyph{SingletonCost}_i}$.
For the paired cost, the scalar expression becomes
\[
 2\bigl(\sym{MatchingPhaseCost}(\sym{PairMeanAngle}+\sym{PairPhaseParameter})
          +\sym{MatchingPhaseCost}(\sym{PairMeanAngle}-\sym{PairPhaseParameter})\bigr)\big|_{p=1}
 =8\sin(\sym{PairMeanAngle}/2)\cos(\sym{PairPhaseParameter}/2).
\]
It is nonincreasing as $\sym{PairPhaseParameter}$ runs from $\sym{PairHalfDifference}$ to $\sym{PairMeanAngle}$, so its minimum is at $\sym{PairPhaseParameter}=\sym{PairMeanAngle}$,
with value $4\sin \sym{PairMeanAngle}=\like{PairCost}{\symglyph{PairCost}_{ij}}$.
At this endpoint one rotational phase is $2\sym{PairMeanAngle}$ and the other is zero:
the paired transporter fixes two real directions and rotates the other
two through $\like{theta}{\symglyph{theta}_i}+\like{theta}{\symglyph{theta}_j}$.
Taking the limit in \cref{eq:matching-power-equality} now gives
\cref{eq:real-nuclear-match}. Compactness supplies an attaining transporter
for the trace norm as well.
\end{proof}
% Formal proof locations: RealNuclearMatching.lean (endpoint continuation),
% SchattenRealDecomposition.lean / SchattenRealRestrictions.lean (block
% reduction), SchattenRealMatchingWeights.lean / RealMatchingAttainment.lean
% (matching and attainment), SchattenRealAngleCosts.lean (angle costs),
% RealFourCoordinates.lean / RealFourPositive.lean / RealFourOptimization.lean
% (all two-angle competitors and the scalar edge reduction).
% All paths are under DavisKahan/FiniteDimensional/DirectRotation/.
% See AUTHORING_NOTES.md for the mathematical and provenance evidence.

For $\like{theta}{\symglyph{theta}_1}=\like{theta}{\symglyph{theta}_2}=\sym{PiConstant}/6$, the direct rotation costs
$8\sin(\sym{PiConstant}/12)>2$, whereas the paired competitor costs
$4\sin(\sym{PiConstant}/6)=2$. The theorem includes comparison with every ambient
orthogonal transporter, so this paired cost is globally minimal in the
example.

Davis and Kahan's complex-phase example already shows how a phase change
can decrease full trace displacement \citep[Example~4.2]{davis1970rotation}.
Bolotnikov, Li, and Rodman formulate the ambient unitary minimal-distortion
problem and prove Q-norm and Frobenius optimality results
\citep[Problem~1.1 and Theorems~3.1, 4.2]{bolotnikov2001minimal}. Our formula gives the exact minimum trace-norm displacement over all real
orthogonal maps carrying one subspace onto the other, in terms of singleton
and paired angle costs. We did not locate this exact matching formula in
the sources compared here. The real field permits pairing between distinct
principal planes, which is why the answer is an optimization over matchings.


\section{Discussion}
\label{sec:discussion}

The formalization covers the original perturbation theory on real and
complex Hilbert spaces, including unbounded self-adjoint operators and
operator-ideal norms. The unbounded statements require explicit domain
conditions and bounded residuals. The operator-ideal statements require
finite norm values. These hypotheses determine which approximation problems
the estimates address.

The correction of the full-displacement claim
\citep[Proposition~4.4]{davis1970rotation} separates two extremal statements.
The valid comparison \citep[Proposition~4.3]{davis1970rotation}
gives squared-displacement minimality; its Q-norm consequence
is valid. Our real four-dimensional family exhibits full trace-norm failure
at every prescribed positive acute angle. In each example, the competitor has
two nonzero displacement singular values while the direct rotation has four.
The sums and squared sums of these values account for their different norm
comparisons. The earlier complex example supplied the construction; our
calculation makes its contradiction of the printed real claim explicit.

The matching theorem addresses the resulting optimization problem for any
equal-dimensional pair of finite-dimensional real subspaces. Its proof uses
principal-angle geometry, real orthogonal block decompositions, and
singular-value norms. These are also among the foundations required for the
Davis--Kahan formalization. The proof of the matching theorem uses additional
optimization arguments, including a passage to the trace-norm endpoint.

\section{Conclusion}
\label{sec:conclusion}

We formalized 28 of the 29 selected Davis--Kahan results and refuted the
remaining printed statement \citep[Proposition~4.4]{davis1970rotation}.
The sine-theta theorem includes
unbounded self-adjoint operators and symmetric operator-ideal norms with
their domain and finiteness assumptions. The direct-rotation calculation
identifies the failed singular-value comparison and gives an explicit real
counterexample and a valid Q-norm consequence.

We also formalized the operator theory used in these results: principal
angles and direct rotations, Borel spectral calculus, approximation numbers
and symmetric ideals, and Sylvester separation. Parts of this development
were subsequently adapted into Tau Ceti. Using these foundations, we proved
an exact weighted-matching formula for the real trace-norm transport
minimum.

\backmatter
\section*{Statements and Declarations}

\subsection*{Funding}
\ifDKAnonymous
Funding information is omitted for anonymous review.
\else
This work was supported by the Defense Advanced Research Projects Agency
under Contract No. HR001125CE017.
\fi

\subsection*{Competing interests}
The authors declare no competing interests.

\subsection*{Data and code availability}
\ifDKAnonymous
Artifact information is omitted for anonymous review.
\else
The Lean proofs, result-to-Lean mapping, provenance records, figure-generation
sources, and manuscript sources are available at
\url{https://github.com/AIQ-Kitware/aiq-dkps-formalization}.
\fi

\subsection*{AI assistance}
ChatGPT and Claude assisted with the formalization and manuscript preparation.
Formalizations were carried out by large language models with human oversight.
The matching result was derived using GPT-6.1 Sol and reviewed by the human
  authors.

\bibliography{references}

\clearpage
\appendix

\section{Symbols, conventions, and terminology}
\label{app:notation}

We write $\sym{A}$ for a reference operator and
$\sym{Acurrent}=\sym{A}+\sym{H}$ for the perturbed operator. The
sine-theta theorem uses coordinate maps for trial and exact subspaces.
The transport results name the subspaces $\sym{Usource},\sym{Vtarget}$
and their projections $\sym{Psource},\sym{Qtarget}$ directly.
The sine-theta projection is $\sym{Pexact}$, the residual is $\sym{R}$,
and direct rotations are denoted by $\sym{Rrotation}$ throughout.

In the colored version, blue identifies exact or target data, amber trial
or source data, cyan the exact complementary block, magenta angles, green
residuals, and violet spectral separation. The residual $\sym{R}$ and
direct rotation $\sym{Rrotation}$ have distinct roles and colors. The
notation and definitions also apply in monochrome.

\begingroup
\small
\setlength{\tabcolsep}{4pt}
\renewcommand{\arraystretch}{1.15}
\begin{longtable}{@{}>{\raggedright\arraybackslash}p{0.29\textwidth}>{\raggedright\arraybackslash}p{0.65\textwidth}@{}}
\toprule
\textbf{Symbol} & \textbf{Definition or interpretation}\\
\midrule
\endfirsthead
\multicolumn{2}{@{}l}{\textit{Symbols and conventions (continued)}}\\[2pt]
\toprule
\textbf{Symbol} & \textbf{Definition or interpretation}\\
\midrule
\endhead
\bottomrule
\endfoot
$\sym{A},\sym{H},\sym{Acurrent}$ & Reference operator, additive perturbation, and perturbed operator $\sym{Acurrent}=\sym{A}+\sym{H}$.\\
$\sym{AmbientHilbertSpace},\ \sym{TrialCoordinateSpace},\ \sym{ExactCoordinateSpace},\ \sym{ComplementaryCoordinateSpace}$ & Ambient, trial-coordinate, exact-coordinate, and complementary-coordinate Hilbert spaces.\\
$\sym{E0}:\sym{TrialCoordinateSpace}\to\sym{AmbientHilbertSpace}$ & \dkSymbolDescription{E0}\\
$\sym{TrialOperator}$ & Self-adjoint trial operator. A Rayleigh--Ritz choice is $\sym{TrialOperator}=\like{E0}{\symglyph{E0}^*}\sym{Acurrent}\sym{E0}$.\\
$\sym{ReferenceBlock0},\sym{ReferenceBlock1}$ & Reference blocks on $\DKRange(\sym{Psource})$ and $\DKKernel(\sym{Psource})$ in Section~\ref{sec:foundations}.\\
$\sym{F0}:\sym{ExactCoordinateSpace}\to\sym{AmbientHilbertSpace}$ & \dkSymbolDescription{F0}\\
$\sym{F1}:\sym{ComplementaryCoordinateSpace}\to\sym{AmbientHilbertSpace}$ & \dkSymbolDescription{F1}\\
$\sym{Lambda1}$ & Complementary self-adjoint block, with $\sym{Acurrent}\sym{F1}=\sym{F1}\sym{Lambda1}$ on its domain.\\
$\sym{SpectralMeasure},\sym{SpectralSet}$ & Projection-valued measure and Borel spectral set; $\sym{SpectralMeasure}(\sym{SpectralSet})$ is the corresponding spectral projection. These symbols differ from the ambient space $\sym{AmbientHilbertSpace}$ and complementary block $\sym{S}$.\\
$\sym{Pexact}=\sym{F0}\like{F0}{\symglyph{F0}^*}$ & Projection onto the desired exact subspace in the sine-theta theorem.\\
$\sym{Theta0},\sym{sinTheta0}$ & Trial-coordinate angle operator and its positive sine, with $\sym{sinTheta0}=(\sym{Identity}-\like{E0}{\symglyph{E0}^*}\sym{Pexact}\sym{E0})^{1/2}$.\\
$\sym{Theta}=\sym{DiagonalOperator}(\sym{Theta0},\sym{Theta1})$ & Full angle operator; $\sym{Theta1}$ compares the orthogonal complements. Its operator norm is the largest spectral angle.\\
$\sym{CrossedSourceDefect},\sym{CrossedTargetDefect}$ & Source directions perpendicular to the target, and target directions perpendicular to the source. Section~\ref{sec:foundations} uses the condition $\sym{DimensionOperator}\sym{CrossedSourceDefect}=\sym{DimensionOperator}\sym{CrossedTargetDefect}$.\\
$\{\sym{ZeroVector}\}$ & Zero subspace: the set containing only the zero vector.\\
$\sym{S}=(\sym{Identity}-\sym{Pexact})\sym{E0}$ & Complementary block, with modulus $\dkModulus{\sym{S}}=\sym{sinTheta0}$.\\
$\sym{X}=\like{F1}{\symglyph{F1}^*}\sym{E0}$ & Coordinates of the trial component in the unwanted space; $\sym{S}=\sym{F1}\sym{X}$.\\
$\sym{SylvesterLeft},\sym{SylvesterRight},\sym{SylvesterRHS}$ & Generic left and right Sylvester coefficients and right side: $\sym{SylvesterLeft}\sym{X}-\sym{X}\sym{SylvesterRight}=\sym{SylvesterRHS}$. The historical equation writes the right coefficient as $\sym{HistoricalSylvesterRight}$.\\
$\sym{R}=\sym{Acurrent}\sym{E0}-\sym{E0}\sym{TrialOperator}$ & Bounded residual agreeing with this difference on $\DKDom(\sym{TrialOperator})$.\\
$\sym{delta}>0$ & \dkSymbolDescription{delta}\\
$[\sym{betaEndpoint},\sym{alphaEndpoint}]$ & Interval in the interval/exterior gap hypothesis.\\
$\DKRange(T),\DKKernel(T),\DKDom(T)$ & Range, kernel, and domain of an operator.\\
$\sym{Usource},\sym{Vtarget}$ & Source and target subspaces in the displacement problem.\\
$\sym{Psource},\sym{Qtarget}$ & Orthogonal projections onto $\sym{Usource},\sym{Vtarget}$ in that problem.\\
$\sym{Rrotation},\sym{Wcompetitor}$ & Direct rotation and an arbitrary orthogonal or unitary transporter carrying $\sym{Usource}$ onto $\sym{Vtarget}$.\\
$\sym{PolarU}$ & Partial isometry in a polar decomposition, as in $\sym{S}=\sym{PolarU}|\sym{S}|$.\\
$\sym{IntertwinerT}$ & The intertwiner $\sym{Qtarget}\sym{Psource}+(\sym{Identity}-\sym{Qtarget})(\sym{Identity}-\sym{Psource})$, whose unitary polar factor is the acute direct rotation.\\
$\sym{theta},\like{theta}{\symglyph{theta}_j}$ & Individual principal angles, in $[0,\sym{PiConstant}/2]$.\\
$\sym{Xframe},\sym{Yframe}$ & Orthonormal coordinate matrices for the source and target in the real counterexample.\\
$\sym{Jturn},\sym{Cturn},\sym{Hmix}$ & Planar rotation through $\sym{PiConstant}/2$, rotation through $2\sym{theta}$, and orthogonal mixing matrix used in that counterexample.\\
$\sym{Jreflect};\ \sym{AngleGenerator}$ & Source reflection $2\sym{Psource}-\sym{Identity}$, with square $\sym{Identity}$; skew angle generator of a direct rotation. The $\sym{PiConstant}/2$-rotation matrix $\sym{Jturn}$ has square $-\sym{IdentityTwo}$.\\
$\sym{CounterPrincipalOne},\sym{CounterPrincipalTwo};\ \sym{CounterMovingPlane},\sym{CounterFixedPlane}$ & Principal planes of the direct rotation and moving/fixed planes of the competitor in \cref{fig:prop44-displacement-planes}.\\
$\sym{CounterChordA},\sym{CounterChordB}$ & Displacement singular values $\sqrt{2-\sqrt2}$ and $\sqrt2$ in that counterexample.\\
$\sym{Kdisplacement}=\sym{Identity}-\sym{Wcompetitor}$ & Full displacement operator in the failed Ky Fan comparison.\\
$\sym{OmegaOne},\sym{OmegaTwo}$ & Principal-plane projections used to test that comparison.\\
$\dkModulus{T}=(T^*T)^{1/2}$ & Positive modulus of an operator.\\
$\sym{SingularValue}(T)$ & Ordered singular values; approximation numbers provide the corresponding comparisons for bounded operators in Hilbert space.\\
$\sym{NormN}(T)$ & \dkSymbolDescription{NormN}\\
$\dkOperatorNorm{T}$ & Operator norm; in finite dimension, the largest singular value.\\
$\dkFrobeniusNorm{T}$ & Frobenius norm, $\bigl(\sum_j \sym{SingularValue}(T)^2\bigr)^{1/2}$.\\
$\dkTraceNorm{T}$ & Trace norm, $\sum_j \sym{SingularValue}(T)$.\\
$\dkKyFanNorm{T}{j}$ & Ky Fan norm, the sum of the $j$ largest singular values.\\
$\sym{NormQ}\text{-norm}$ & A norm satisfying $\sym{NormN}(T)^2=\sym{NormM}(T^*T)$ for a unitarily invariant norm $\sym{NormM}$.\\
$\sym{Matching},\like{SingletonCost}{\symglyph{SingletonCost}_i},\like{PairCost}{\symglyph{PairCost}_{ij}}$ & Matching of positive principal-angle labels, with singleton and paired costs in \cref{eq:nuclear-matching-costs}.\\
$\sym{MatchingEdgeWeight};\ \like{TraceCost}{\symglyph{TraceCost}(M)}$ & Edge weight $\like{SingletonCost}{\symglyph{SingletonCost}_i}+\like{SingletonCost}{\symglyph{SingletonCost}_j}-\like{PairCost}{\symglyph{PairCost}_{ij}}$; total trace-norm displacement cost of the matching.\\
\end{longtable}
\endgroup

\subsection{Terminology}
\label{app:terminology}

The following terms describe the roles of the mathematical objects used in
\cref{sec:sine-theta,sec:example-transport}.

\begin{description}
\item[Ambient space and operator.] The whole Hilbert space
$\sym{AmbientHilbertSpace}$ in which the subspaces lie, and the operator acting
on that space. In the sine-theta statement the ambient operator is
$\sym{Acurrent}$.

\item[Reference and perturbed operators.] The reference operator is
$\sym{A}$; an additive perturbation $\sym{H}$ gives the perturbed operator
$\sym{Acurrent}=\sym{A}+\sym{H}$. The choice of which operator supplies the
desired spectral subspace depends on the approximation problem.

\item[Exact subspace.] The desired subspace specified by the spectral
decomposition of the operator of interest. In the sine-theta statement,
it is $\DKRange(\sym{F0})$ for $\sym{Acurrent}$, with orthogonal complement
$\DKRange(\sym{F1})$. In the perturbation application in the introduction,
the desired subspace instead belongs to the reference operator $\sym{A}$.
An exact formula, as in the transport theorem, gives an equality for the
quantity being determined.

\item[Trial subspace and trial operator.] The trial subspace is a candidate
approximation to the desired subspace, represented by $\DKRange(\sym{E0})$.
The self-adjoint trial operator $\sym{TrialOperator}$ represents an
approximation to the action of the ambient operator in trial coordinates.

\item[Residual.] The difference
$\sym{Acurrent}\sym{E0}-\sym{E0}\sym{TrialOperator}$ between the ambient
action and the embedded trial action. For unbounded operators this identity
holds on the appropriate domain, with a bounded residual $\sym{R}$.

\item[Spectral separation.] A lower bound $\sym{delta}>0$ on the distance
between the relevant spectra. In the sine-theta theorem these are the
spectra of the trial operator and the complementary block $\sym{Lambda1}$;
the interval and ordered separation hypotheses are given in
\cref{sec:gap-geometry}.

\item[Transporter and competitor.] A transporter is an ambient orthogonal
or unitary map carrying the source subspace onto the target subspace.
A competitor is another map satisfying the same constraint.

\item[Full displacement.] The operator $\sym{Identity}-\sym{Wcompetitor}$
on the whole ambient space for a transporter $\sym{Wcompetitor}$. Its
restriction to source directions is represented by
$(\sym{Identity}-\sym{Wcompetitor})\sym{Psource}$.
\end{description}

\section{Reading the Lean signatures}
\label{app:lean-guide}
\label{app:reading-lean}

A Lean theorem declaration has the form
\[
  \texttt{theorem name (parameters) (hypotheses) : conclusion := proof}.
\]
Parameters introduce the mathematical objects under discussion.  A hypothesis
such as \texttt{(hA : IsSelfAdjoint A)} supplies the corresponding fact;
bundled predicates such as \texttt{IsTrialResidual} and
\texttt{IsExactSpectralDecomposition} collect related assumptions.  Braces mark
arguments normally inferred by Lean, while square brackets request typeclass
instances supplying ambient algebraic, analytic, or topological structure.  The proposition after the final colon is the theorem's conclusion
\citep{demoura2021lean4}.

In \Cref{lst:dk-sintheta-headline}, \texttt{RCLike} is the Mathlib abstraction
covering $\sym{RealField}$ and $\sym{ComplexField}$.  A type
\texttt{ContinuousLinearMap K F E} is a continuous linear map.  A type
\texttt{LinearPMap K E E} is a linear partial map whose
domain is a submodule; this is the representation used for possibly unbounded
operators.  \texttt{N.Mem T} states that the ideal gauge associated with $\sym{NormN}$ is
finite on $T$, and \texttt{N.gaugeReal T} is the corresponding real value.  The
suffix \texttt{Real} refers to the codomain of the gauge and not to the scalar
field.

\subsection{Paper notation and Lean binder names}
\label{app:paper-lean-notation}
The following correspondence is for the particular sine-theta theorem
in \Cref{lst:dk-sintheta-headline}, not a global naming rule for the Lean
library. The Lean source is
\texttt{DavisKahan/Sources/DavisKahan1970/SineTheta/Presentation.lean},
the declaration \texttt{sinTheta\_unbounded\_formGap\_whereDefinedUIN\_rclike}.
Lean's exact source uses Unicode subscripts and Greek letters (for example,
its trial-operator binder is typeset here as $\mathtt{A_0}$). The listing
in the main text uses ASCII transliterations and readable type names.
In particular, the Lean binder $\mathtt{A_0}$ denotes the general
$\sym{TrialOperator}$, while $\sym{ReferenceBlock0}$ is the separate
reference spectral block and does not occur as a parameter of this theorem.
The variable $\mathtt{H}$ is a coordinate-space carrier, not the
perturbation $\sym{H}$. Lean's scalar-field type variable is rendered as \texttt{K} in the
schematic listing (the source uses a blackboard-bold k).

\begingroup
\small
\setlength{\tabcolsep}{4pt}
\renewcommand{\arraystretch}{1.11}
\begin{longtable}{@{}>{\centering\arraybackslash}p{0.22\linewidth}>{\centering\arraybackslash}p{0.16\linewidth}>{\raggedright\arraybackslash}p{0.55\linewidth}@{}}
\caption{Paper symbols and the corresponding Lean binders in the sine-theta declaration.}
\label{tab:paper-lean-notation}\\
\toprule
This paper & Lean binder & Mathematical role \\
\midrule
\endfirsthead
\toprule
This paper & Lean binder & Mathematical role \\
\midrule
\endhead
\bottomrule
\endfoot
\dkLeanSineThetaRows
\end{longtable}
\endgroup

\section{Formalized Davis--Kahan results}
\label{app:dk-inventory}

Table~\ref{tab:dk-full-inventory} lists the \DKResultCount{} source targets:
the four headline theorems \citep[Section~2]{davis1970rotation} and every named theorem,
proposition, lemma, and corollary established in the paper. The table keeps
compact IDs for reference. The artifact maps each ID to the corresponding
Lean declarations.
\ifDKAnonymous\else
The \href{https://github.com/AIQ-Kitware/aiq-dkps-formalization/blob/main/dev/davis-kahan-1970-formalization-result-inventory.json}{result-to-Lean mapping}
contains that correspondence.
\fi

\begin{longtable}{>{\raggedright\arraybackslash}p{0.24\linewidth}>{\raggedright\arraybackslash}p{0.68\linewidth}}
\caption{The \DKResultCount{} Davis--Kahan results selected for formalization.}
\label{tab:dk-full-inventory}\\
\toprule
ID & Result \\
\midrule
\endfirsthead
\toprule
ID & Result \\
\midrule
\endhead
\path{S2-sin-theta} & Single-angle sine theorem \\
\path{S2-tan-theta} & Single-angle tangent theorem \\
\path{S2-sin-two-theta} & Double-angle sine theorem \\
\path{S2-tan-two-theta} & Double-angle tangent theorem \\
\path{DK-3.1-prop} & Acute direct rotation existence and uniqueness \\
\path{DK-3.2-prop} & Nonacute existence criterion \\
\path{DK-3.3-prop} & Principal square-root characterization \\
\path{DK-3.4-prop} & Square as a direct rotation \\
\path{DK-3.1-thm} & Classification of pairs of subspaces \\
\path{DK-3.1-cor} & Compact classification by angle eigenvalues \\
\path{DK-3.5-prop} & Angle commutation and eigenspace geometry \\
\path{DK-3.2-cor} & Reversal symmetry \\
\path{DK-4.1-prop} & Pointwise and singular-value extremality of the direct rotation \\
\path{DK-4.1-cor} & Unitarily invariant norm minimality of restricted displacement \\
\path{DK-4.2-prop} & Basis-angle square-sum extremality \\
\path{DK-4.3-prop} & Squared-displacement unitarily invariant norm minimality \\
\path{DK-4.4-prop} & Full-displacement optimality claim \\
\path{DK-5.1-thm} & Banach-space Sylvester lower bound \\
\path{DK-5.2-thm} & Semibounded self-adjoint Sylvester theorem \\
\path{DK-5.1-lem} & Strong-cutoff convergence of singular values \\
\path{DK-6.1-lem} & Direct-sum norm comparison and converse \\
\path{DK-6.2-lem} & Reflection-pinch contraction \\
\path{DK-6.1-prop} & Sine proof, ambient limitation, and symmetric sine theorem \\
\path{DK-6.1-thm} & Generalized sine theorem \\
\path{DK-6.2-thm} & Pairwise-gap square-norm sine theorem \\
\path{DK-6.3-thm} & Generalized tangent theorem and supporting machinery \\
\path{DK-6.3-lem} & Finite-rank near-maximizer off-diagonal block estimate \\
\path{DK-8.1-thm} & Branch selection and spectral repulsion \\
\path{DK-8.2-thm} & Smallness selects the acute branch \\
\bottomrule
\end{longtable}

A single source result may use several Lean declarations. The IDs identify
source statements independently of the organization of the formal library.
The printed claim \citep[Proposition~4.4]{davis1970rotation}
is refuted by the counterexample in
Section~\ref{sec:prop44}; its Q-norm consequence is stated separately there.

% Detailed table following the compact summary in app:dk-inventory.
% Summaries checked against the registered Lean statements and their supporting
% declarations; see TABLE4_LEAN_AUDIT.md for the row-by-row comparison.
\subsection{Hypotheses and conclusions}
\label{app:dk-result-summaries}

Table~\ref{tab:dk-result-summaries} summarizes the hypotheses and main
conclusions of the \DKResultCount{} results listed in
Table~\ref{tab:dk-full-inventory}, using the same IDs
\citep[Sections~2--8]{davis1970rotation}.
Hilbert spaces are separable and real or complex unless specified otherwise.
$\sym{NormN}$ denotes an arbitrary unitarily invariant norm on its domain of
finiteness; each displayed norm is evaluated on an operator in that domain.
$\dkOperatorNorm{\cdot}$ and $\dkHilbertSchmidtNorm{\cdot}$ denote the operator
and Hilbert--Schmidt norms. Self-adjoint operators may be unbounded:
operator identities hold on their domains, and residuals or right-hand sides
in norm estimates have bounded extensions.
For the ambient tangent bound, the tangent is defined as a bounded operator,
equivalently $\dkOperatorNorm{\sym{Psource}-\sym{Qtarget}}<1$.

For the geometric results, $\sym{Psource},\sym{Qtarget}$ are orthogonal
projections, $\sym{SourceDirectRotation}$ is a direct rotation, and
$\sym{Wcompetitor}$ is any orthogonal or unitary transporter satisfying
$\sym{Wcompetitor}\sym{Psource}\sym{Wcompetitor}^*=\sym{Qtarget}$.
Write
\[
 \sym{CrossedSourceDefect}=\DKRange(\sym{Psource})\cap\DKKernel(\sym{Qtarget}),\qquad
 \sym{CrossedTargetDefect}=\DKKernel(\sym{Psource})\cap\DKRange(\sym{Qtarget}).
\]
The pair is acute when $\sym{CrossedSourceDefect}=\sym{CrossedTargetDefect}=\{\sym{ZeroVector}\}$.
For geometric and displacement rows from \path{DK-3.3-prop} onward,
the ranges and kernels of
$\sym{Psource},\sym{Qtarget}$ have matching dimensions and
$\sym{DimensionOperator} \sym{CrossedSourceDefect}=\sym{DimensionOperator} \sym{CrossedTargetDefect}$; the displacement results also assume
the directed sine block
$(\sym{Identity}-\sym{Qtarget})|_{\DKRange(\sym{Psource})}$ compact.
$\sym{Theta0},\sym{Theta1}$ are the source and complementary angle operators,
$\sym{Theta}$ their ambient direct sum, and $\like{theta}{\symglyph{theta}_j}$ the ordered principal
angles. In the source decomposition the direct rotation has blocks
$\bigl(\begin{smallmatrix}C_0&-S_0^*\\S_0&C_1\end{smallmatrix}\bigr)$
with $C_j\ge0$; $\sym{AngleGenerator}$ is its skew angle generator, so
$\sym{SourceDirectRotation}=e^{\sym{AngleGenerator}\sym{Theta}}$.
The symbol $\sym{SourceDirectRotation}$ includes nonacute direct rotations;
for an acute pair it is the unique rotation denoted by the same letter
in the main text.
For real unitaries, spectral half-plane conditions refer to their
complexifications.

For perturbation results, $\sym{TrialOperator}$ is a self-adjoint trial operator.
In a reducing decomposition of the reference operator $\sym{A}$, the
source and complementary reference blocks are $\sym{ReferenceBlock0}$
and $\sym{ReferenceBlock1}$. The trial operator may be chosen to equal the
source reference block when the trial coordinates identify that subspace.
The operators $\sym{Lambda0},\sym{Lambda1}$ are the exact blocks of
$\sym{Acurrent}=\sym{A}+\sym{H}$.
The maps $\sym{E0},\sym{F0}$ give source and target coordinates,
$\sym{H0},\sym{H1}$ are the diagonal perturbation blocks, and
$\sym{R}=\sym{Acurrent}\sym{E0}-\sym{E0}\sym{TrialOperator}$.
Coordinate maps are isometries except in the two lower-frame rows.
The exact frames have orthogonal ranges whose sum is the ambient space;
reducing projections preserve operator domains.
We abbreviate the two spectral hypotheses, for $\sym{delta}>0$, by
\begin{align*}
 \sym{IntervalSeparation}_{\sym{delta}}(B,C):\quad&
 \sym{SpectrumSymbol}(B)\subset[\sym{betaEndpoint},\sym{alphaEndpoint}],\quad
 \sym{SpectrumSymbol}(C)\cap(\sym{betaEndpoint}-\sym{delta},\sym{alphaEndpoint}+\sym{delta})=\varnothing,
 \quad\text{or the roles interchanged};\\
 \sym{OrderedSeparation}_{\sym{delta}}(B,C):\quad&
 B\le\sym{alphaEndpoint} \sym{Identity},\quad C\ge(\sym{alphaEndpoint}+\sym{delta})\sym{Identity}.
\end{align*}
The endpoints are real, satisfy $\sym{betaEndpoint}\le\sym{alphaEndpoint}$, and are chosen for each row.
The ordered condition $\sym{OrderedSeparation}_{\sym{delta}}(B,C)$ places the spectra in
the half-infinite intervals $(-\infty,\sym{alphaEndpoint}]$ and
$[\sym{alphaEndpoint}+\sym{delta},\infty)$; it allows both operators to be unbounded.
To include both orientations, set
\[
 \sym{SineGap}(B,C)\quad\Longleftrightarrow\quad
 \sym{IntervalSeparation}_{\sym{delta}}(B,C)\quad\text{or}\quad
 \sym{OrderedSeparation}_{\sym{delta}}(B,C)\quad\text{or}\quad
 \sym{OrderedSeparation}_{\sym{delta}}(C,B).
\]
This is the separation condition of Section~\ref{sec:gap-geometry}.

\begingroup
\small
\setlength{\tabcolsep}{4pt}
\renewcommand{\arraystretch}{1.15}
\begin{longtable}{@{}>{\raggedright\arraybackslash}p{0.20\linewidth}>{\raggedright\arraybackslash}p{0.34\linewidth}>{\raggedright\arraybackslash}p{\dimexpr0.46\linewidth-4\tabcolsep\relax}@{}}
\caption{Hypotheses and main conclusions of the \DKResultCount{} Davis--Kahan results.}
\label{tab:dk-result-summaries}\\
\toprule
\textbf{ID} & \textbf{Hypotheses} & \textbf{Conclusion}\\
\midrule
\endfirsthead
\multicolumn{3}{@{}l}{\textit{Hypotheses and conclusions (continued)}}\\[2pt]
\toprule
\textbf{ID} & \textbf{Hypotheses} & \textbf{Conclusion}\\
\midrule
\endhead
\midrule
\multicolumn{3}{r@{}}{\textit{Continued on the next page}}\\
\endfoot
\bottomrule
\endlastfoot
\multicolumn{3}{@{}l}{\textbf{Headline perturbation bounds}}\\*\addlinespace[3pt]
\path{S2-sin-theta} &
 $\sym{SineGap}(\sym{TrialOperator},\sym{Lambda1})$. &
 $\sym{delta}\sym{NormN}(\sin\sym{Theta0})\le\sym{NormN}(\sym{R})$.\\\addlinespace
\path{S2-tan-theta} &
 $\sym{OrderedSeparation}_{\sym{delta}}(\sym{TrialOperator},\sym{Lambda1})$;
 Rayleigh--Ritz condition $\sym{H0}=0$. &
 $\sym{delta}\sym{NormN}(\tan\sym{Theta0})\le\sym{NormN}(\sym{R})$;
 $\sym{delta}\sym{NormN}(\tan\sym{Theta})\le\sym{NormN}(\sym{H})$.\\\addlinespace
\path{S2-sin-two-theta} &
 $\sym{A},\sym{Acurrent}$ have a common domain;
 $\sym{SpectrumSymbol}(\sym{Lambda0})\subset[\sym{betaEndpoint},\sym{alphaEndpoint}]$;
 $\sym{SpectrumSymbol}(\sym{Lambda1})$ outside
 $(\sym{betaEndpoint}-\sym{delta},\sym{alphaEndpoint}+\sym{delta})$. &
 $\sym{delta}\sym{NormN}(\sin(2\sym{Theta0}))\le2\sym{NormN}(\sym{R})$;
 $\sym{delta}\sym{NormN}(\sin(2\sym{Theta}))\le2\sym{NormN}(\sym{H})$.\\\addlinespace
\path{S2-tan-two-theta} &
 $\sym{OrderedSeparation}_{\sym{delta}}(\sym{TrialOperator},\sym{A1})$;
 bounded $\sym{H}$, $\sym{H0}=\sym{H1}=0$. &
 $\sym{delta}\sym{NormN}(\tan(2\sym{Theta0}))\le2\sym{NormN}(\sym{R})$;
 $\sym{delta}\sym{NormN}(\tan(2\sym{Theta}))\le2\sym{NormN}(\sym{H})$.\\\addlinespace[5pt]
\multicolumn{3}{@{}l}{\textbf{Two subspaces and direct rotations}}\\*\addlinespace[3pt]
\path{DK-3.1-prop} &
 $\sym{CrossedSourceDefect}=\sym{CrossedTargetDefect}=\{\sym{ZeroVector}\}$. &
 A unique direct rotation exists. Positive diagonal blocks characterize it
 among transporters.\\\addlinespace
\path{DK-3.2-prop} &
 General pair of projections. &
 A direct rotation exists iff $\sym{DimensionOperator} \sym{CrossedSourceDefect}=\sym{DimensionOperator} \sym{CrossedTargetDefect}$.
 On $\sym{CrossedSourceDefect}\oplus \sym{CrossedTargetDefect}$, $\sym{SourceDirectRotation}^2=-\sym{Identity}$;
 uniqueness can fail.\\\addlinespace
\path{DK-3.3-prop} &
 A direct rotation $\sym{SourceDirectRotation}$; conversely a unitary
 square root with spectrum in the closed right half-plane that maps
 $\sym{CrossedSourceDefect}$ onto $\sym{CrossedTargetDefect}$. &
 $\sym{SourceDirectRotation}^2=(2\sym{Qtarget}-\sym{Identity})(2\sym{Psource}-\sym{Identity})$,
 and its spectrum lies in the closed right half-plane.
 Conversely, such a square root is a direct rotation.\\\addlinespace
\path{DK-3.4-prop} &
 $C_0^2\ge\tfrac12 \sym{Identity}$;
 $Q_-=(2\sym{Psource}-\sym{Identity})\sym{Qtarget}(2\sym{Psource}-\sym{Identity})$. &
 $\sym{SourceDirectRotation}^2$ is a direct rotation from
 $\DKRange(Q_-)$ to $\DKRange(\sym{Qtarget})$.\\\addlinespace
\path{DK-3.1-thm} &
 Matching dimensions as above; angle operators
 $0\le\sym{Theta0},\sym{Theta1}\le\sym{PiConstant}/2$. &
 The four intersection dimensions and the generic angle's spectral
 multiplicity completely classify pairs up to isometric equivalence.
 Any angle data with matching nonzero multiplicities
 and the prescribed total dimension are realizable.\\\addlinespace
\path{DK-3.1-cor} &
 Classification hypotheses;
 $\sym{Psource}(\sym{Identity}-\sym{Qtarget})\sym{Psource}$ compact. &
 Classification reduces to angle eigenvalues with multiplicity:
 $\like{theta}{\symglyph{theta}_j}\downarrow0$. The two angle operators may have different
 zero multiplicities.\\\addlinespace
\path{DK-3.5-prop} &
 A direct rotation and its angle operator $\sym{Theta}$; the maximal
 eigenspace characterization assumes the acute case. &
 $\sym{Theta}$ commutes with $\sym{Psource},\sym{Qtarget},\sym{AngleGenerator},
 \sym{SourceDirectRotation}$. If $x\ne0$ and
 $\sym{Theta}x=\sym{theta} x$, then
 $\angle(x,\sym{SourceDirectRotation}x)=\sym{theta}$.
 In the acute case these eigenspaces are the unique maximal reducing
 subspaces of the prescribed angle.\\\addlinespace
\path{DK-3.2-cor} &
 Interchange $\sym{Psource}$ and $\sym{Qtarget}$. &
 $\sym{Theta}$ is unchanged, $\sym{AngleGenerator}$ changes to $-\sym{AngleGenerator}$, and the corresponding
 direct rotation becomes $\sym{SourceDirectRotation}^*$.\\\addlinespace[5pt]
\multicolumn{3}{@{}l}{\textbf{Displacement extremality}}\\*\addlinespace[3pt]
\path{DK-4.1-prop} &
 Compact angle data; arbitrary transporter $\sym{Wcompetitor}$. &
 For the positive principal angles, there are orthonormal
 $v_j\in\DKRange(\sym{Psource})$ with
 $\angle(v_j,\sym{Wcompetitor}v_j)\ge\like{theta}{\symglyph{theta}_j}$.
 Each singular value of restricted displacement is minimized by
 $\sym{SourceDirectRotation}$, with value $2\sin(\like{theta}{\symglyph{theta}_j}/2)$.\\\addlinespace
\path{DK-4.1-cor} &
 Same displacement setting. &
 $\sym{NormN}((\sym{Identity}-\sym{SourceDirectRotation})\sym{Psource})
 \le\sym{NormN}((\sym{Identity}-\sym{Wcompetitor})\sym{Psource})$.\\\addlinespace
\path{DK-4.2-prop} &
 Any orthonormal basis $(v_j)$ of $\DKRange(\sym{Psource})$;
 arbitrary transporter $\sym{Wcompetitor}$. &
 $\sum_j\sin^2\angle(v_j,\sym{Wcompetitor}v_j)
 \ge\sum_j\sin^2\like{theta}{\symglyph{theta}_j}$, including an infinite right-hand side.\\\addlinespace
\path{DK-4.3-prop} &
 Same displacement setting; put
 $\sym{SquaredDisplacement}(W)=(\sym{Identity}-W)^*(\sym{Identity}-W)$. &
 $\sym{NormN}(\sym{SquaredDisplacement}(\sym{SourceDirectRotation}))
 \le\sym{NormN}(\sym{SquaredDisplacement}(\sym{Wcompetitor}))$:
 the direct rotation minimizes squared full displacement.\\\addlinespace
\path{DK-4.4-prop} &
 Real space; largest principal angle at most $\sym{PiConstant}/3$. &
 The printed assertion that $\sym{SourceDirectRotation}$ minimizes
 every $\sym{NormN}(\sym{Identity}-\sym{Wcompetitor})$ is false.
 A real four-dimensional trace-norm counterexample is given in
 \cref{sec:prop44}; the Q-norm repair is stated there separately.\\\addlinespace[5pt]
\multicolumn{3}{@{}l}{\textbf{Sylvester estimates and spectral cutoffs}}\\*\addlinespace[3pt]
\path{DK-5.1-thm} &
 Banach spaces; bounded $C,X,Y$; $BX-XC=Y$;
 $\dkOperatorNorm{C}\le a$,
 $\dkOperatorNorm{B^{-1}}\le(a+\sym{delta})^{-1}$, $a\ge0$.
 If $B$ is unbounded, it is closed and densely defined with bounded inverse.
 The map norm is compatible with these operator norms. &
 $\sym{delta}\|X\|\le\|Y\|$.
 The coefficient hypotheses may be interchanged.\\\addlinespace
\path{DK-5.2-thm} &
 Self-adjoint $B\ge(\sym{SylvesterGapThreshold}+\sym{delta})\sym{Identity}$,
 $C\le\sym{SylvesterGapThreshold} \sym{Identity}$; bounded $X,Y$;
 $X\DKDom(C)\subset\DKDom(B)$ and $BX-XC=Y$ there. &
 $\sym{delta}\sym{NormN}(X)\le\sym{NormN}(Y)$,
 including unbounded $B,C$.\\\addlinespace
\path{DK-5.1-lem} &
 Orthogonal projections $\sym{StrongApproxProjection}\to \sym{Identity}$ strongly; bounded $K$. &
 $\sym{SingularValue}(K\sym{StrongApproxProjection})\to \sym{SingularValue}(K)$ for each singular-value index $j$.\\\addlinespace[5pt]
\multicolumn{3}{@{}l}{\textbf{Generalized single-angle bounds}}\\*\addlinespace[3pt]
\path{DK-6.1-lem} &
 Each diagonal block of $K$ is dominated by the corresponding block of
 $L$ in every $\sym{NormN}$. &
 The same domination holds for their block-diagonal sums.
 The converse holds if the two $K$ blocks have identical singular values
 (counting multiplicities), as do the two $L$ blocks.\\\addlinespace
\path{DK-6.2-lem} &
 Orthogonal projections $\sym{PinchingProjectionOne},\sym{PinchingProjectionTwo}$; bounded $K$. &
 $\sym{NormN}(\sym{PinchingProjectionOne} K\sym{PinchingProjectionTwo}+
 (\sym{Identity}-\sym{PinchingProjectionOne})K(\sym{Identity}-\sym{PinchingProjectionTwo}))\le\sym{NormN}(K)$.\\\addlinespace
\path{DK-6.1-prop} &
 $\sym{A},\sym{Acurrent}$ have a common domain and bounded difference
 $\sym{H}$;
 both $\sym{SineGap}(\sym{TrialOperator},\sym{Lambda1})$ and
 $\sym{SineGap}(\sym{A1},\sym{Lambda0})$. &
 $\sym{delta}\sym{NormN}(\sin\sym{Theta})\le\sym{NormN}(\sym{H})$.\\\addlinespace
\path{DK-6.1-thm} &
 $\sym{E0}^*\sym{E0}\ge\sym{FrameLowerBoundEpsilon}^2\sym{Identity}$, $\sym{FrameLowerBoundEpsilon}>0$;
 $\sym{SineGap}(\sym{TrialOperator},\sym{Lambda1})$.
 Source and target dimensions may differ. &
 $\sym{delta}\sym{FrameLowerBoundEpsilon}\sym{NormN}(\sin\sym{Theta0})
 \le\sym{NormN}(\sym{R})$.
 The directed sine has the singular values of
 $\sym{Psource}(\sym{Identity}-\sym{Qtarget})$.\\\addlinespace
\path{DK-6.2-thm} &
 Same lower-frame condition; only
 $\sym{SpectralDistanceOperator}(\sym{SpectrumSymbol}(\sym{TrialOperator}),\sym{SpectrumSymbol}(\sym{Lambda1}))
 \ge\sym{delta}$. &
 $\sym{delta}\sym{FrameLowerBoundEpsilon}\|\sin\sym{Theta0}\|_{\mathrm{HS}}
 \le\|\sym{R}\|_{\mathrm{HS}}$.
 If $\sym{R}$ has finite rank, also
 $\sym{delta}\sym{FrameLowerBoundEpsilon}\dkOperatorNorm{\sin\sym{Theta0}}
 \le\dkOperatorNorm{\sym{R}}\sqrt{\sym{RankOperator}\sym{R}}$.\\\addlinespace
\path{DK-6.3-thm} &
 Isometric Rayleigh--Ritz trial data;
 $\sym{OrderedSeparation}_{\sym{delta}}(\sym{TrialOperator},\sym{Lambda1})$.
 The source dimension may be smaller than the target dimension. &
 $\sym{delta}\sym{NormN}(\tan\sym{Theta0})
 \le\sym{NormN}(\sym{R})$.\\\addlinespace
\path{DK-6.3-lem} &
 Bounded $K$; projections $\sym{CompressionProjectionImage},\sym{CompressionProjectionInput}$ of rank at most $n\ge1$;
 $K\sym{CompressionProjectionInput}=\sym{CompressionProjectionImage} K\sym{CompressionProjectionInput}$;
 $\sum_{j\le n}\sym{SingularValue}(K\sym{CompressionProjectionInput})^2$\newline
 ${}>\sum_{j\le n}\sym{SingularValue}(K)^2-\sym{CompressionTolerance}^2$, $\sym{CompressionTolerance}>0$. &
 $\dkOperatorNorm{\sym{CompressionProjectionImage} K(\sym{Identity}-\sym{CompressionProjectionInput})}<\sym{CompressionTolerance}$:
 a nearly maximizing finite-rank compression has small ``leakage''
 (i.e., the norm of this off-diagonal block is below $\sym{CompressionTolerance}$).\\\addlinespace[5pt]
\multicolumn{3}{@{}l}{\textbf{Spectral selection and repulsion}}\\*\addlinespace[3pt]
\path{DK-8.1-thm} &
 Ordered reference gap
 $\sym{OrderedSeparation}_{\sym{delta}}(\sym{ReferenceBlock0},\sym{ReferenceBlock1})$ with endpoint $\sym{alphaEndpoint}$;
 bounded $\sym{H}$, $\sym{H0}=\sym{H1}=0$.
 For the equivalence, $\sym{DimensionOperator}\sym{CrossedSourceDefect}=\sym{DimensionOperator}\sym{CrossedTargetDefect}$. &
 For a chosen reducing target,
 $\sym{Theta}\le\sym{PiConstant}/4$ iff
 $\sym{Lambda0}\le\sym{alphaEndpoint} \sym{Identity}$ and
 $\sym{Lambda1}\ge(\sym{alphaEndpoint}+\sym{delta})\sym{Identity}$; such a reducing target exists.
 With exact frames obtained by applying $\sym{SourceDirectRotation}$ to the reference frames,
 $\sym{ReferenceBlock1}-\sym{alphaEndpoint} \sym{Identity}\le\sym{ComplementCosineBlock}(\sym{Lambda1}-\sym{alphaEndpoint} \sym{Identity})\sym{ComplementCosineBlock}$,
 as quadratic forms, with the lower-block analogue.
 In finite dimensions there are eigenvalue and symmetric-gauge
 repulsion inequalities.\\\addlinespace
\path{DK-8.2-thm} &
 Double-angle sine hypotheses; $\sym{DimensionOperator}\sym{CrossedSourceDefect}=\sym{DimensionOperator}\sym{CrossedTargetDefect}$;
 either bounded $\sym{H}$ with
 $\dkOperatorNorm{\sym{H}}<\sym{delta}/2$, or bounded
 $\sym{R}$ with $\dkOperatorNorm{\sym{R}}<\sym{delta}/2$;
 $\sym{SpectrumSymbol}(\sym{ReferenceBlock0})\subset
 [\sym{betaEndpoint}-\sym{delta}/2,\sym{alphaEndpoint}+\sym{delta}/2]$. &
 The corresponding ambient or directed double-angle sine bound holds
 on the branch with angles strictly below $\sym{PiConstant}/4$.\\
\end{longtable}
\endgroup

\clearpage
\section{Notation correspondence with Davis--Kahan (1970)}
\label{app:dk-notation-correspondence}
This table compares our notation with the corresponding symbols of
Davis and Kahan~\citep[especially \S1]{davis1970rotation}. The correspondences
are by mathematical role, not just by letter. In the original, the reference
block $A_0$ also serves as the residual's comparison operator. Here
$\sym{TrialOperator}$ is a general trial operator, which need not equal that block.
Our additional matching constructions have no direct counterpart in the source.

\begingroup
\small
\setlength{\tabcolsep}{3pt}
\renewcommand{\arraystretch}{1.12}
\begin{longtable}{@{}>{\raggedright\arraybackslash}p{0.43\linewidth}>{\centering\arraybackslash}p{0.12\linewidth}>{\centering\arraybackslash}p{0.16\linewidth}>{\raggedright\arraybackslash}p{0.23\linewidth}@{}}
\caption{Documented correspondence with the source symbols of Davis--Kahan (1970).}
\label{tab:dk-source-notation}\\
\toprule
Mathematical object & DK (1970) & This paper & DK location \\
\midrule
\endfirsthead
\toprule
Mathematical object & DK (1970) & This paper & DK location \\
\midrule
\endhead
\bottomrule
\endfoot
\dkDavisKahanRows
\end{longtable}
\endgroup


\end{document}
, without editing.
% EVERY \sym, \like, and \dkrole obeys this switch, including the glossary.
% No color is required to identify a mathematical object: the glyphs, indices,
% definitions and adjacent prose also distinguish them.
%
% Families (base hue): exact/target=blue, trial/source=amber, exact unwanted=
% cyan, angles and their operators=magenta, residual/displacement=green,
% gaps=violet. Norm gauges=slate; provenance retains brown/steel/vermilion.
% P_ex is the exact-target projection in sine-theta; P is the source
% projection in transport. R denotes a residual and D a direct rotation.
% Semantic keys retain their roles even when the printed notation changes.
\newif\ifDKMathColors
\DKMathColorstrue
\ifdefined\DKMonochromeMath
  \DKMathColorsfalse
\fi

% Family anchors (print-compatible versions of the Manim light palette).
\definecolor{dkfg}{HTML}{1B1F24}
\definecolor{dkmuted}{HTML}{6B7280}
\definecolor{dkwanted}{HTML}{1F6FD1}
\definecolor{dktrial}{HTML}{A66705}
\definecolor{dkunwanted}{HTML}{007C91}
\definecolor{dksine}{HTML}{D12A7B}
\definecolor{dkresid}{HTML}{087D43}
\definecolor{dkgap}{HTML}{6E4BD8}
\definecolor{dkold}{HTML}{705519}
\definecolor{dkcurrent}{HTML}{054669}
\definecolor{dkperturb}{HTML}{BB491B}
\definecolor{dkrefuted}{HTML}{C62828}

% Related yet distinguishable shades. These are NOT new unrelated categories:
% objects inherit their parent family; a different shade signals a different
% mathematical object (e.g., a frame versus its subspace or its projection).
\definecolor{dkwantedframe}{HTML}{1254A9}
\definecolor{dkwantedprojection}{HTML}{08618D}
\definecolor{dkwantedrotation}{HTML}{4248B0}
\definecolor{dkwantedspectrum}{HTML}{345F8F}
\definecolor{dkwantedvector}{HTML}{1468AA}
\definecolor{dktrialframe}{HTML}{905600}
\definecolor{dktrialprojection}{HTML}{79511B}
\definecolor{dktrialcompetitor}{HTML}{A04D10}
\definecolor{dktrialspectrum}{HTML}{79591C}
\definecolor{dktrialvector}{HTML}{9B522A}
\definecolor{dkunwantedframe}{HTML}{006477}
\definecolor{dkunwantedspectrum}{HTML}{00546B}
\definecolor{dkunwantedprojection}{HTML}{145B6D}
\definecolor{dkunwantedvector}{HTML}{256879}
\definecolor{dkangleoperator}{HTML}{923478}
\definecolor{dksineoperator}{HTML}{AA275A}
\definecolor{dksineblock}{HTML}{B43F70}
\definecolor{dksylvester}{HTML}{804074}
\definecolor{dkanglemax}{HTML}{B32185}
\definecolor{dkanglecos}{HTML}{824276}
\definecolor{dkresidvector}{HTML}{086038}
\definecolor{dkdisplacement}{HTML}{226C50}
\definecolor{dkgaporiginal}{HTML}{4E4D89}
\definecolor{dkgapstat}{HTML}{7039A7}
\definecolor{dkgapinterval}{HTML}{574A82}
\definecolor{dkambient}{HTML}{36516A}
\definecolor{dknorm}{HTML}{475875}
\definecolor{dknormaux}{HTML}{66526E}
\definecolor{dkintertwiner}{HTML}{4A6280}
\definecolor{dkturn}{HTML}{47647D}
\definecolor{dkmix}{HTML}{505B65}

% Coloring primitives. "plain" suppresses color without suppressing hovers.
% All descendants can retain a semantic role through \like{key}{expression}.
\expandafter\def\csname dknocolor@plain\endcsname{1}
\newcommand{\dkrole}[2]{%
  \ifDKMathColors
    \ifcsname dknocolor@#1\endcsname #2\else\textcolor{dk#1}{#2}\fi
  \else #2\fi
}
% An optional plain-text reminder supplies definitions and types on hover.
% Required fields: {semantic key}{printed glyph}{color role}{glossary|local}
%                  {short definition}{mathematical type}. The "glossary" flag records whether an
% entry appears in the curated printed notation table; "local" does not.
% This registry is the sole authority for symbol meaning and printed form.
% Use role "plain" for an uncolored symbol that still carries hover text.
% Changing the glyph is a separate editorial decision: this audit preserves
% the current E, F, F_ex, G, and other existing printed choices.
\newcommand{\dkDeclareSymbol}[7][]{%
  \expandafter\def\csname dkgloss@#2\endcsname{#6}%
  \expandafter\def\csname dkglossaryflag@#2\endcsname{#5}%
  \if\relax\detokenize{#1}\relax
    \expandafter\pdfstringdef\csname dkpdfgloss@#2\endcsname{#7: #6}%
  \else
    \expandafter\pdfstringdef\csname dkpdfgloss@#2\endcsname{#7: \detokenize{#1}}%
  \fi
  \expandafter\def\csname dktype@#2\endcsname{#7}%
  \expandafter\def\csname dkglyph@#2\endcsname{#3}%
  \expandafter\def\csname dkrole@#2\endcsname{#4}%
  \expandafter\def\csname dksym@#2\endcsname{\dkrole{#4}{#3}}%
}
% Unstyled printed glyph, for variants constructed through \like.
% Example: \like{E0}{\symglyph{E0}^*} derives E_0^* from E0's
% one canonical glyph; retyping E_0 here would defeat central control.
\DeclareRobustCommand{\symglyph}[1]{\csname dkglyph@#1\endcsname}
\DeclareRobustCommand{\sym}[1]{%
  \ifcsname dksym@#1\endcsname
    \dkSymbolTooltip{#1}{\csname dksym@#1\endcsname}%
  \else
    \PackageError{dk-math-colors}{Unknown symbol key #1}{See math_symbol_colors.tex}%
  \fi
}
\DeclareRobustCommand{\like}[2]{%
  \ifcsname dkrole@#1\endcsname
    \dkSymbolTooltip{#1}{\dkrole{\csname dkrole@#1\endcsname}{#2}}%
  \else
    \PackageError{dk-math-colors}{Unknown role key #1}{See math_symbol_colors.tex}%
  \fi
}

% A short description belongs to the semantic key, not the bare letter.
% The same descriptions are available to the notation appendix.
\DeclareRobustCommand{\dkSymbolDescription}[1]{\csname dkgloss@#1\endcsname}
\DeclareRobustCommand{\dkSymbolGlossaryStatus}[1]{\csname dkglossaryflag@#1\endcsname}
\DeclareRobustCommand{\dkSymbolType}[1]{\csname dktype@#1\endcsname}

% Citable alternative notation is metadata, not a request to rename glyphs.
% The source identifier, the source's exact printed glyph, and the specific
% source location are all mandatory. No alternative is inferred from a name.
% One key may have alternatives from several works or multiple passages.
\newcommand{\dkDeclareNotationSource}[3]{%
  \expandafter\def\csname dksourcetitle@#1\endcsname{#2}%
  \expandafter\def\csname dksourceurl@#1\endcsname{#3}%
}
\newcommand{\dkDavisKahanRows}{}
\makeatletter
\newcommand{\dkDeclareNotationAlternative}[5]{%
  \ifcsname dkglyph@#1\endcsname\else
    \PackageError{dk-math-colors}{Unknown alternative key #1}{Declare the symbol first}%
  \fi
  \ifcsname dksourcetitle@#2\endcsname\else
    \PackageError{dk-math-colors}{Unknown notation source #2}{Declare the source first}%
  \fi
  \expandafter\def\csname dkaltglyph@#2@#1\endcsname{#3}%
  \expandafter\def\csname dkaltlocator@#2@#1\endcsname{#4}%
  \expandafter\def\csname dkaltnote@#2@#1\endcsname{#5}%
  \def\dk@source{#2}%
  \def\dk@original{DK1970}%
  \ifx\dk@source\dk@original
    \g@addto@macro\dkDavisKahanRows{\dkDavisKahanRow{#1}}%
  \fi
}
\makeatother
% A correspondence row is produced from the SAME registry records used by
% \sym and by the source-census script. It is never spelled a second time.
\newcommand{\dkDavisKahanRow}[1]{%
  \csname dkaltnote@DK1970@#1\endcsname &
  $\csname dkaltglyph@DK1970@#1\endcsname$ &
  $\sym{#1}$ &
  \csname dkaltlocator@DK1970@#1\endcsname\\
}

\newif\ifDKSymbolTooltips
\DKSymbolTooltipstrue
\ifdefined\DKNoSymbolTooltips\DKSymbolTooltipsfalse\fi
\newsavebox{\DKTooltipBox}
\newcount\DKTooltipCount
% Register only shipped widgets through hyperref's existing form-field
% machinery. Poppler/Okular need the catalog's field list, not just page
% annotations. The deferred aux write also discards measuring-only copies.
\makeatletter
\newcommand{\dkRegisterTooltipField}{%
  \expandafter\HyField@@AddToFields\expandafter{\number\pdffeedback lastannot}%
}
\makeatother
\AtBeginDocument{\ifDKSymbolTooltips\Form[NeedAppearances=false]\fi}
\AtEndDocument{\ifDKSymbolTooltips\endForm\fi}
% Native LuaLaTeX PDF widgets: no navigation action or JavaScript. Avoid
% pdfcomment's soulpos dependency, which previously broke local Vim builds.
% Measure in the CURRENT math style, then emit the original expression as-is:
% boxing the visible math would change spacing, scripts and line breaking.
\newcommand{\dkSymbolTooltip}[2]{%
  \ifDKSymbolTooltips
    \begingroup
      \DKSymbolTooltipsfalse % suppress nested annotations while measuring
      \ifmmode
        \edef\DKTooltipStyle{\ifcase\mathstyle
          \displaystyle\or\displaystyle\or\textstyle\or\textstyle\or
          \scriptstyle\or\scriptstyle\or\scriptscriptstyle\or\scriptscriptstyle\fi}%
        \sbox{\DKTooltipBox}{$\DKTooltipStyle#2$}%
      \else
        \leavevmode\sbox{\DKTooltipBox}{#2}%
      \fi
      \global\advance\DKTooltipCount by 1
      \pdfextension annot width\wd\DKTooltipBox
        height\ht\DKTooltipBox depth\dp\DKTooltipBox
        {/Subtype /Widget /FT /Btn /Ff 65536 /F 512
         /T (dk-symbol-\the\DKTooltipCount) /TU (\csname dkpdfgloss@#1\endcsname)
         /H /N /BS << /W 0 >> /Border [0 0 0]}%
      \dkRegisterTooltipField
      #2%
    \endgroup
  \else
    #2%
  \fi
}
\pdfstringdefDisableCommands{%
  \def\sym#1{\csname dkglyph@#1\endcsname}%
  \def\like#1#2{#2}%
  \def\symglyph#1{\csname dkglyph@#1\endcsname}%
}

% Declarations: [detailed hover], key, glyph, role, glossary|local, short text, type.
% A local entry has full semantic/hover support, but no notation-table row.
% When two keys have the same glyph, it is an auditable collision, not an
% instruction to change the printed notation automatically.
% This registry includes local proof notation. The notation appendix selects
% the symbols useful for reference; registration does not add a printed row.
% Manim's original symbol keys, now split into semantic subshades.
% Ambient and perturbed operators (provenance is independent of geometry).
\dkDeclareSymbol[Reference self-adjoint operator A on the ambient Hilbert space; the perturbed operator is Ahat = A + H. In a single-operator eigenproblem, A itself is the operator being approximated. An unbounded A acts only on domain(A).]
  {A}{A}{old}{glossary}{Reference self-adjoint operator before perturbation.}{self-adjoint-operator}
\dkDeclareSymbol{A1}{A_1}{old}{local}{Complementary block of the reference operator.}{self-adjoint-operator}
\dkDeclareSymbol{a1}{a_1}{old}{local}{Reference eigenvalue or complementary spectral value.}{scalar}
\dkDeclareSymbol[Reference block A_0: the restriction of A to the range of the reducing source projection P, expressed in source coordinates. In Theorem 8.1 it lies at or below alpha; in Theorem 8.2 its spectrum lies in the enlarged central interval. The perturbed blocks are denoted Lambda_0 and Lambda_1.]
  {ReferenceBlock0}{A_0}{old}{glossary}{Reference operator restricted to the source subspace in Section 8.}{self-adjoint-operator}
\dkDeclareSymbol[Reference block A_1: the restriction of A to the kernel of the reducing source projection P, expressed in complementary coordinates. In Theorem 8.1 it lies at or above alpha + delta in quadratic-form order.]
  {ReferenceBlock1}{A_1}{old}{glossary}{Reference operator restricted to the source complement in Section 8.}{self-adjoint-operator}
\dkDeclareSymbol[Self-adjoint ambient operator Ahat being studied; in the perturbation setting Ahat = A + H. Its action on a trial vector is compared with the proposed eigenvalue times that vector.]
  {Ahat}{\widehat A}{current}{local}{Operator under study; in the perturbation setting, the reference operator plus its perturbation.}{self-adjoint-operator}
\dkDeclareSymbol[Self-adjoint ambient operator Ahat being studied. In the perturbation setting Ahat = A + H. The trial residual is R = Ahat E_0 - E_0 A_tr, and the exact complementary block satisfies Ahat F_1 = F_1 Lambda_1 on its domain.]
  {Acurrent}{\widehat A}{current}{glossary}{Operator under study; in the perturbation setting, the reference operator plus its perturbation.}{self-adjoint-operator}
\dkDeclareSymbol[Additive self-adjoint perturbation H = Ahat - A. A bound on its norm can control the residual of an eigenspace of one operator relative to the other.]
  {H}{H}{perturb}{glossary}{Additive self-adjoint perturbation of the reference operator.}{self-adjoint-operator}
\dkDeclareSymbol{H0}{H_0}{perturb}{local}{Wanted diagonal block of the perturbation.}{self-adjoint-operator}
\dkDeclareSymbol{H1}{H_1}{perturb}{local}{Unwanted diagonal block of the perturbation.}{self-adjoint-operator}
\dkDeclareSymbol{B}{B}{perturb}{local}{Off-diagonal perturbation or coupling block.}{linear-operator}
\dkDeclareSymbol{b}{b}{perturb}{local}{Scalar perturbation or coupling coefficient.}{scalar}
% Trial/source family: subspace, frame, compressed operator, projector, vector.
\dkDeclareSymbol{U}{U}{wanted}{local}{Desired exact subspace in the local geometric notation.}{linear-subspace}% Manim's U remains wanted in local slide notation.
\dkDeclareSymbol{V}{V}{trial}{local}{Trial subspace in the local geometric notation.}{linear-subspace}% Manim's V remains trial in local slide notation.
\dkDeclareSymbol[Source subspace U to be carried onto the target V by an orthogonal or unitary map W. In the transport sections P projects onto U and Q onto V. The matching theorem also names the source subspace K.]
  {Usource}{U}{trial}{glossary}{Source subspace to be transported onto the target subspace.}{linear-subspace}
\dkDeclareSymbol[Trial frame E_0: F -> E, a linear isometry from trial-coordinate space F into ambient space E. Its range is the trial subspace; E_0* returns trial coordinates and E_0* E_0 = I_F. In rank one, E_0(x) = x v.]
  {E0}{E_0}{trialframe}{glossary}{Isometric embedding of the trial subspace; its adjoint returns trial coordinates.}{isometric-embedding}
\dkDeclareSymbol[Self-adjoint trial operator A_tr on trial-coordinate space F. The residual compares Ahat E_0 with E_0 A_tr. Rayleigh--Ritz chooses A_tr = E_0* Ahat E_0 when this compression is defined; the theorem allows other trial operators.]
  {A0}{A_{\mathrm{tr}}}{trialspectrum}{glossary}{Self-adjoint operator on trial coordinates; its eigenvalues are the trial spectral values.}{self-adjoint-operator}
\dkDeclareSymbol[Trial vector: a proposed approximate eigenvector. In the rank-one geometry v has unit length and E_0(x) = x v. Its residual for trial scalar rho is r = Ahat v - rho v.]
  {v}{v}{trialvector}{local}{Trial vector or vector in the trial subspace.}{vector}
\dkDeclareSymbol[Scalar trial operator in rank one, A_tr(x) = rho x. Figure 1 uses the Rayleigh value rho = <v, Ahat v> for unit v; then rho v is the projection of Ahat v onto the trial line and r is perpendicular to v.]
  {rho}{\rho}{trialspectrum}{local}{Trial eigenvalue or scalar trial operator.}{scalar}
\dkDeclareSymbol[Approximate eigenvalue assigned to a trial vector v. For the local operator A, its residual is A v - mu v. For unit v, the Rayleigh--Ritz choice is mu = <A v, v>.]
  {mu}{\mu}{trialspectrum}{local}{Trial spectral value.}{scalar}
\dkDeclareSymbol{a0}{a_0}{trialspectrum}{local}{Wanted trial eigenvalue or trial spectral value.}{scalar}
\dkDeclareSymbol[Ambient orthogonal or unitary transporter W with W(U) = V, equivalently W P W* = Q. Its full displacement is I - W on the whole ambient space. The real matching theorem minimizes its trace norm over every such orthogonal W.]
  {Wcompetitor}{W}{trialcompetitor}{glossary}{Orthogonal or unitary competitor carrying the source subspace onto the target.}{unitary-operator}
\dkDeclareSymbol[Orthogonal projection P onto the source subspace U in the transport problem; P = P* = P squared. A feasible orthogonal or unitary W satisfies W P W* = Q. Restricted displacement (I - W)P tests only source vectors.]
  {Psource}{P}{trialprojection}{glossary}{Orthogonal projection onto the source subspace in the transport problem.}{orthogonal-projection}
\dkDeclareSymbol{ProjSource}{P_U}{trialprojection}{local}{Orthogonal projection onto the source subspace in the transport problem.}{orthogonal-projection}
\dkDeclareSymbol[Vector x in the source subspace U. For principal vectors, unit x_j in U and y_j in V are chosen with <x_i, y_j> = 0 for i different from j and <x_j, y_j> = cos(theta_j).]
  {SourceVector}{x}{trialvector}{local}{Vector in the source subspace.}{vector}
\dkDeclareSymbol[Orthonormal coordinate frame X for the source subspace, so X* X = I. Together with a target frame Y, the singular values of X* Y are the principal cosines. In the four-dimensional counterexample X is the vertical stack of I_2 and 0.]
  {Xframe}{X}{trialframe}{glossary}{Orthonormal coordinate matrix for the source subspace in the counterexample.}{isometric-embedding}
\dkDeclareSymbol[Hilbert space F of trial coordinates. E_0 embeds F isometrically into ambient space E; A_tr and the positive sine operator sin(Theta_0) act on F, while the residual R maps F into E.]
  {Ftrial}{F}{trialframe}{glossary}{Hilbert space of trial coordinates.}{hilbert-space}
\dkDeclareSymbol[Ambient Hilbert space E over the real or complex scalars. It contains the trial subspace range(E_0) and the orthogonal exact subspaces range(F_0) and range(F_1). The source theorem permits infinite-dimensional separable E.]
  {Eambient}{E}{ambient}{glossary}{Ambient Hilbert space containing the trial and exact subspaces.}{hilbert-space}
% Zero and all-ones vectors are distinguished from scalar 0 and 1.
\dkDeclareSymbol[Zero vector of the vector or Hilbert space in context: the additive identity. The zero subspace consists only of this vector.]
  {ZeroVector}{\mathbf{0}}{fg}{glossary}{Zero vector; the additive identity of the ambient vector space.}{vector}
\dkDeclareSymbol[All-ones vector in a specified finite-dimensional coordinate space: every coordinate equals 1. This vector depends on the chosen coordinates and dimension.]
  {OnesVector}{\mathbf{1}}{fg}{local}{All-ones vector in a chosen finite-dimensional coordinate space.}{vector}
% Exact/wanted family: the target subspace, its frame, projection and rotation.
\dkDeclareSymbol[Target subspace V in the transport problem: feasible maps satisfy W(U) = V, or equivalently W P W* = Q. The matching theorem also names this subspace L.]
  {Vtarget}{V}{wanted}{glossary}{Target subspace in the transport problem.}{linear-subspace}
\dkDeclareSymbol[Exact wanted frame F_0: F_ex -> E, a linear isometry onto the desired exact spectral subspace. F_0* returns its coordinates. P_ex = F_0 F_0* is its ambient orthogonal projection; in Figure 1, F_0(h) = (h, 0).]
  {F0}{F_0}{wantedframe}{glossary}{Isometric coordinates for the desired exact spectral subspace; star denotes the adjoint.}{isometric-embedding}
\dkDeclareSymbol[Self-adjoint operator block on coordinates for the desired exact spectral subspace. Its spectrum consists of the wanted exact spectral values; in the corresponding decomposition Ahat F_0 = F_0 Lambda_0 on the operator domain.]
  {Lambda0}{\Lambda_0}{wantedspectrum}{local}{Self-adjoint block on the desired exact spectral subspace.}{self-adjoint-operator}
\dkDeclareSymbol{u}{u}{wantedvector}{local}{Vector in the desired exact subspace.}{vector}
\dkDeclareSymbol{f0}{f_0}{wantedvector}{local}{Component or coordinate vector in the desired exact subspace.}{vector}
\dkDeclareSymbol{lambda0}{\lambda_0}{wantedspectrum}{local}{Wanted exact eigenvalue or wanted spectral value.}{scalar}
\dkDeclareSymbol{lambdau}{\lambda_u}{wantedspectrum}{local}{Eigenvalue associated with a desired exact vector.}{scalar}
\dkDeclareSymbol[Exact wanted orthogonal projection P_ex: E -> E, with P_ex = F_0 F_0*. It retains the desired exact spectral component. Its complementary projection is P_ex_perp = I - P_ex = F_1 F_1*; the projected trial frame is S = P_ex_perp E_0.]
  {Pexact}{P_{\mathrm{ex}}}{wantedprojection}{glossary}{Orthogonal projection onto the desired exact subspace in the sine-theta theorem.}{orthogonal-projection}
\dkDeclareSymbol[Exact complementary orthogonal projection P_ex_perp = I - P_ex = F_1 F_1*: E -> E. It retains the unwanted exact component. In Figure 1, P_ex_perp v = S(1), whereas P_ex_perp r is the unwanted component of the residual.]
  {PexactPerp}{P_{\mathrm{ex}}^\perp}{unwantedprojection}{local}{Orthogonal projection onto the exact complementary subspace.}{orthogonal-projection}
\dkDeclareSymbol{ReflectionSigma}{\Sigma}{wantedprojection}{local}{Reflection associated with the desired exact spectral projection.}{unitary-operator}
\dkDeclareSymbol[Orthogonal projection Q onto the target subspace V in the transport problem. A feasible orthogonal or unitary W satisfies W P W* = Q, where P projects onto the source U.]
  {Qtarget}{Q}{wantedprojection}{glossary}{Orthogonal projection onto the target subspace in the transport problem.}{orthogonal-projection}
\dkDeclareSymbol{ProjTarget}{P_V}{wantedprojection}{local}{Orthogonal projection onto the target subspace in the transport problem.}{orthogonal-projection}
\dkDeclareSymbol[Canonical direct rotation D from the source U to the target V for an acute pair. It is the unitary polar factor of T = QP + (I - Q)(I - P), and acts by the principal angle on each principal plane. R is reserved for the trial residual.]
  {Rrotation}{D}{wantedrotation}{glossary}{Canonical direct rotation carrying the source subspace onto the target.}{unitary-operator}
\dkDeclareSymbol[Direct rotation D between the ranges of orthogonal projections P and Q, including the nonacute case. Relative to range(P) and kernel(P), its diagonal blocks are positive and its off-diagonal blocks are opposite adjoints. It exists when the two crossed defect spaces have equal dimension, and is unique when both defects vanish.]
  {SourceDirectRotation}{D}{wantedrotation}{local}{Direct rotation, including the nonacute case in the source-result table.}{unitary-operator}
\dkDeclareSymbol[Orthonormal coordinate frame Y for the target subspace, so Y* Y = I. In the counterexample Y is the vertical stack of cos(theta) I_2 and sin(theta) J_0; hence X* Y = cos(theta) I_2 and both principal angles equal theta.]
  {Yframe}{Y}{wantedframe}{glossary}{Orthonormal coordinate matrix for the target subspace in the counterexample.}{isometric-embedding}
\dkDeclareSymbol[Vector y in the target subspace V. For principal vectors, unit x_j in U and y_j in V are paired so that <x_j, y_j> = cos(theta_j); distinct pairs have zero cross inner product.]
  {TargetVector}{y}{wantedvector}{local}{Vector in the target subspace.}{vector}
\dkDeclareSymbol[Hilbert space F_ex of exact wanted coordinates. The isometry F_0: F_ex -> E identifies it with the desired exact spectral subspace. This space differs from H, the additive perturbation.]
  {Hexact}{F_{\mathrm{ex}}}{wantedframe}{glossary}{Hilbert space of coordinates for the desired exact spectral subspace.}{hilbert-space}% Coordinate Hilbert space (not perturbation).
% Exact/unwanted family; the trial error operators are MAGENTA, not cyan.
\dkDeclareSymbol{Uperp}{U^{\perp}}{unwanted}{local}{Orthogonal complement of the desired exact subspace in the local notation.}{linear-subspace}
\dkDeclareSymbol[Exact complementary frame F_1: G -> E, a linear isometry onto the orthogonal complement of range(F_0). F_1* returns unwanted coordinates; F_1 F_1* = I - P_ex. In Figure 1, F_1(g) = (0, g).]
  {F1}{F_1}{unwantedframe}{glossary}{Isometric coordinates for the exact complementary subspace; star denotes the adjoint.}{isometric-embedding}
\dkDeclareSymbol[Self-adjoint unwanted block Lambda_1 on G, satisfying Ahat F_1 = F_1 Lambda_1 on its domain. Its spectrum is separated from the trial spectrum of A_tr by the specified positive delta. Figure 1 uses the scalar block 0.7.]
  {Lambda1}{\Lambda_1}{unwantedspectrum}{glossary}{Self-adjoint operator on exact complementary coordinates: the unwanted spectral block.}{self-adjoint-operator}
% Spectral measure and set (distinct from the ambient space E and sine block S).
\dkDeclareSymbol[Projection-valued spectral measure E associated with a self-adjoint operator A: A = integral lambda dE(lambda). For a Borel set S, E(S) is the orthogonal projection onto the corresponding spectral subspace.]
  {SpectralMeasure}{E}{wantedprojection}{glossary}{Projection-valued spectral measure of a self-adjoint operator.}{projection-valued-measure}
\dkDeclareSymbol[Borel subset S of the real spectral axis. The spectral projection E(S) selects the part of a self-adjoint operator with spectral values in S; its range is a reducing subspace.]
  {SpectralSet}{S}{wantedspectrum}{glossary}{Borel spectral set used to select a spectral subspace.}{set}
% Generic Sylvester coefficients and right side, by source/target role.
\dkDeclareSymbol[Left coefficient B of the Sylvester map X -> B X - X C; B acts on the output space of X. For the sine bound, B = Lambda_1 on the unwanted-coordinate space G.]
  {SylvesterLeft}{B}{unwantedspectrum}{glossary}{Left coefficient of the Sylvester equation; corresponds to the unwanted spectral block.}{self-adjoint-operator}
\dkDeclareSymbol[Right coefficient C of the Sylvester map X -> B X - X C; C acts on the input space of X. For the sine bound, C = A_tr on trial space F. The historical equation also denotes this coefficient by A.]
  {SylvesterRight}{C}{trialspectrum}{glossary}{Right coefficient of the Sylvester equation; corresponds to the trial operator.}{self-adjoint-operator}
\dkDeclareSymbol[Right-hand side Y in B X - X C = Y, with the same input and output spaces as X. For the sine bound, Y = F_1* R: F -> G, the residual expressed in unwanted exact coordinates.]
  {SylvesterRHS}{Y}{resid}{glossary}{Right-hand side of the Sylvester equation, obtained by projecting the residual.}{linear-operator}
\dkDeclareSymbol{w}{w}{unwantedvector}{local}{Vector in an unwanted or complementary spectral subspace.}{vector}
\dkDeclareSymbol{f1}{f_1}{unwantedvector}{local}{Component or coordinate vector in the unwanted exact subspace.}{vector}
\dkDeclareSymbol{lambda1}{\lambda_1}{unwantedspectrum}{local}{Unwanted exact eigenvalue or complementary spectral value.}{scalar}
\dkDeclareSymbol{lambdaw}{\lambda_w}{unwantedspectrum}{local}{Eigenvalue associated with an unwanted exact vector.}{scalar}
\dkDeclareSymbol[Hilbert space G of exact complementary coordinates. F_1: G -> E embeds it onto the unwanted exact subspace; Lambda_1 acts on G and X = F_1* E_0 maps trial coordinates into G.]
  {Gunwanted}{G}{unwantedframe}{glossary}{Hilbert space of exact complementary coordinates.}{hilbert-space}
\dkDeclareSymbol{E1}{E_1}{unwantedframe}{local}{Isometric coordinates for the complementary trial subspace.}{isometric-embedding}
\dkDeclareSymbol{Vperp}{V^{\perp}}{unwanted}{local}{Orthogonal complement of the trial subspace in the local notation.}{linear-subspace}
% Angle and error geometry, including distinct rectangular and positive maps.
\dkDeclareSymbol[Individual principal angle, between 0 and pi/2. For orthonormal subspace frames X and Y, the singular values of X* Y are the principal cosines. In Figure 1, theta is the angle between the unit trial vector v and the exact wanted line.]
  {theta}{\theta}{sine}{glossary}{Principal angle between the subspaces under comparison.}{angle}
\dkDeclareSymbol[Full positive angle operator Theta = diag(Theta_0, Theta_1) for the source/target decompositions. Theta_0 compares the source with the target; Theta_1 compares their orthogonal complements. Its operator norm is the largest spectral angle. For positive c, the conditions Theta <= c and Theta < c mean operator norm at most c and strictly below c, respectively.]
  {Theta}{\Theta}{angleoperator}{glossary}{Full angle operator, the diagonal sum of the source and complementary angle blocks.}{angle-operator}
\dkDeclareSymbol[Directed principal-angle operator Theta_0 on source or trial coordinates. For a trial frame E_0 and exact wanted projection P_ex, cos(Theta_0) = sqrt(E_0* P_ex E_0), and sin(Theta_0) = sqrt(I_F - E_0* P_ex E_0). In the two-projection geometry it comes from the positive cosine compression P Q P restricted to range(P).]
  {Theta0}{\Theta_0}{angleoperator}{glossary}{Directed principal-angle operator acting on source or trial coordinates.}{angle-operator}
\dkDeclareSymbol[Complementary directed angle operator Theta_1, obtained from the projections I-P and I-Q by the same cosine-compression construction as Theta_0. It acts on the source complement, or its isometric coordinates. The full angle operator is diag(Theta_0, Theta_1).]
  {Theta1}{\Theta_1}{angleoperator}{glossary}{Directed angle operator comparing the orthogonal complements.}{angle-operator}
\dkDeclareSymbol[Right-angle source defect D_0 = range(P) intersect kernel(Q). These are source vectors orthogonal to the target. The full-angle conclusion in the Section 8 smallness criterion assumes dim(D_0) = dim(D_1); in Lean their identification is CrossedDefectsEquivalent.]
  {CrossedSourceDefect}{D_0}{fg}{glossary}{Source directions perpendicular to the target subspace.}{linear-subspace}
\dkDeclareSymbol[Right-angle target defect D_1 = kernel(P) intersect range(Q). These are target vectors orthogonal to the source. Equality of its Hilbert dimension with that of D_0 is assumed for the full-angle conclusion of Theorem 8.2.]
  {CrossedTargetDefect}{D_1}{fg}{glossary}{Target directions perpendicular to the source subspace.}{linear-subspace}
\dkDeclareSymbol[Positive complementary cosine block C_1 = cos(Theta_1). The reference complement is identified with the perturbed complement by the direct rotation D, so A_1 and Lambda_1 act on common coordinates. Theorem 8.1 compares A_1 - alpha I with C_1 (Lambda_1 - alpha I) C_1 as quadratic forms.]
  {ComplementCosineBlock}{C_1}{anglecos}{local}{Cosine of the complementary angle block in the spectral-repulsion inequality.}{self-adjoint-operator}
\dkDeclareSymbol[Scalar sine of a principal angle. For a unit principal trial vector v and exact wanted projection P_ex, sin(theta) = ||(I - P_ex)v||. In Figure 1 it is the length of the pink leg S(1).]
  {sintheta}{\sin\theta}{sine}{local}{Sine of an individual principal angle: the length of the complementary component of a unit principal trial vector.}{scalar}
\dkDeclareSymbol[Positive sine operator sin(Theta_0) = sqrt(I_F - E_0* P_ex E_0): F -> F. It equals |S| for S = (I - P_ex)E_0: F -> E. Its eigenvalues in finite dimensions are the principal sines; the sine theorem bounds its ideal norm by N(R)/delta.]
  {sinTheta0}{\sin\Theta_0}{sineoperator}{glossary}{Positive sine operator on trial coordinates; the modulus of the rectangular complementary error.}{angle-operator}
\dkDeclareSymbol[Rectangular complementary error S = (I - P_ex) E_0: F -> E. It records how much each trial vector leaks outside the exact wanted subspace. Its modulus |S| = sqrt(S* S) is sin(Theta_0): F -> F. In rank one S(1) = P_ex_perp v.]
  {S}{S}{sineblock}{glossary}{Rectangular complementary error: the trial frame projected onto the unwanted exact subspace.}{linear-operator}
\dkDeclareSymbol[Complementary-coordinate map X = F_1* E_0: F -> G, with S = F_1 X. It satisfies Lambda_1 X - X A_tr = F_1* R on domain(A_tr). Isometric multiplication gives N(X) = N(S); the separate modulus identity abs(S) = sin(Theta_0) gives N(S) = N(sin(Theta_0)), whenever these norms are finite.]
  {X}{X}{sylvester}{glossary}{Solution of the Sylvester equation; for the sine bound, the trial frame in unwanted exact coordinates.}{linear-operator}
\dkDeclareSymbol{alphaAngle}{\alpha}{anglemax}{local}{Largest principal angle.}{angle}
\dkDeclareSymbol[Scalar c = cos(theta) in the counterexample. In the general two-projection decomposition, the same role denotes a positive cosine operator C; the cosine and sine blocks commute and satisfy C squared + S squared = I.]
  {CosAngle}{c}{anglecos}{local}{Cosine of a principal angle or the corresponding positive cosine block.}{scalar-or-operator}
\dkDeclareSymbol[Scalar s = sin(theta) in the counterexample. In the general two-projection decomposition, the same role denotes a positive sine operator S; it commutes with the positive cosine block C and satisfies C squared + S squared = I.]
  {SinAngle}{s}{sineoperator}{local}{Sine of a principal angle or the corresponding positive sine block.}{scalar-or-operator}
\dkDeclareSymbol[Positive small-angle parameter in the local subspace comparison. It controls the allowed angular deviation; it is an angle parameter rather than a perturbation-operator norm.]
  {EpsilonAngle}{\varepsilon}{anglemax}{local}{Angle parameter controlling the maximum subspace deviation.}{angle}
\dkDeclareSymbol{PhiAngle}{\varphi}{angleoperator}{local}{Auxiliary angle or angle operator in the two-subspace comparison.}{angle}
\dkDeclareSymbol{Bblock}{B}{sineblock}{local}{Off-diagonal block in the two-subspace decomposition.}{linear-operator}
\dkDeclareSymbol{Cblock}{C}{sylvester}{local}{Complementary or Sylvester block in the two-subspace decomposition.}{linear-operator}
\dkDeclareSymbol{Ablock}{A}{anglecos}{local}{Positive cosine-related block in the two-subspace decomposition.}{linear-operator}
% Residual/error-displacement family.
\dkDeclareSymbol[Residual vector r = Ahat v - rho v = R(1) for the rank-one frame E_0(x) = x v. Figure 1 draws it from rho v to Ahat v. For unit v and the Rayleigh value rho = <v, Ahat v>, <v, r> = 0; a general trial scalar need not give orthogonality.]
  {r}{r}{residvector}{local}{Residual vector of an approximate eigenpair.}{vector}
\dkDeclareSymbol[Bounded residual R: F -> E, agreeing with Ahat E_0 - E_0 A_tr on domain(A_tr). The domain condition requires E_0(domain(A_tr)) inside domain(Ahat). R = 0 means that E_0 intertwines A_tr with Ahat; its norm measures the defect.]
  {R}{R}{resid}{glossary}{Bounded residual of the trial eigenproblem, agreeing with the operator difference on the trial domain.}{linear-operator}
\dkDeclareSymbol[Full displacement Delta_W = I - W on the ambient space. In the failed Ky Fan comparison W = W_(pi/4); Delta_W has singular values sqrt(2), sqrt(2), 0, 0. Multiplying by a principal-plane projection restricts its input to that plane.]
  {Kdisplacement}{\Delta_W}{displacement}{glossary}{Full displacement of a transporter from the identity.}{operator-valued-function}
\dkDeclareSymbol[Squared full displacement Dcal(W) = (I-W)* (I-W), a positive operator. Proposition 4.3 compares its unitarily invariant norm for a direct rotation D with that for every feasible transporter W.]
  {SquaredDisplacement}{\mathcal D}{displacement}{local}{Squared full displacement as a function of the transporter.}{operator-valued-function}
% Matching proof: the gradient and its spectral level are not a space or gap.
\dkDeclareSymbol[Skew-symmetric spectral gradient G of Phi_p. It commutes with W and represents first variations by d/dt Phi_p(W exp(tB)) at t = 0 equal to -tr(G B) for skew-symmetric B. Feasible stationarity gives J_src G J_src = -G.]
  {SpectralGradient}{G}{displacement}{local}{Skew-symmetric spectral gradient of the displacement power sum in the matching proof.}{linear-operator}
\dkDeclareSymbol[Positive spectral level g of the skew-symmetric gradient G: on this invariant subspace, -G squared = g squared I. Consequently G/g squares to -I and supplies the complex structure used in the real matching proof.]
  {GradientLevel}{g}{displacement}{local}{Positive gradient level; the gradient squared is minus this level squared times the identity.}{scalar}
\dkDeclareSymbol[Complex structure I_g = G/g on the positive gradient-level space V_g. It is a real isometry with square -I. Multiplication by a+ib means a times identity plus b times I_g; source reflection J_src reverses I_g and acts as conjugation.]
  {GradientComplexStructure}{\mathcal I_g}{displacement}{local}{Complex multiplication constructed from a positive gradient level.}{linear-operator}
\dkDeclareSymbol[Real skew-symmetric matrix B, with B transpose = -B, giving the orthogonal variation W exp(tB). If B commutes with the source projection P, this curve preserves the constraint W(U) = V.]
  {SkewDirection}{B}{ambient}{local}{Skew-symmetric direction for a first variation of an orthogonal transporter.}{linear-operator}
% Spectral separation family: gap and original gap are NOT identical.
\dkDeclareSymbol[Positive spectral separation in the stated configuration. In the sine theorem it separates the trial operator from the exact complementary block; Theorem 8.1 uses it for the ordered reference gap, and Theorem 8.2 for the perturbed central-band separation. Figure 1 uses delta = abs(rho - 0.7).]
  {delta}{\delta}{gap}{glossary}{Positive spectral separation in the stated gap configuration.}{scalar}
\dkDeclareSymbol{g}{g}{gaporiginal}{local}{Original spectral gap of the reference operator; distinct from the certified sine-theta separation.}{scalar}
\dkDeclareSymbol{Delta}{\Delta}{gapstat}{local}{Population or statistical spectral gap.}{scalar}
\dkDeclareSymbol[Spectral endpoint alpha. In the interval/exterior configuration, one spectrum lies between beta and alpha, the other at or below beta - delta or at or above alpha + delta. Theorem 8.2 places the selected perturbed block in the central interval. In Theorem 8.1, alpha bounds the lower reference block and alpha + delta bounds the upper block from below.]
  {alphaEndpoint}{\alpha}{gapinterval}{glossary}{Upper endpoint of the central band, or lower-block cut in the ordered gap.}{scalar}
\dkDeclareSymbol[Lower endpoint beta of the finite interval from beta to alpha in the gap hypothesis. One spectrum lies in that interval; the other lies at or below beta - delta or at or above alpha + delta.]
  {betaEndpoint}{\beta}{gapinterval}{glossary}{Lower endpoint of the finite interval in the interval/exterior gap hypothesis.}{scalar}
% Spectral/projector and auxiliary matrix objects, context-scoped.
\dkDeclareSymbol{P}{P}{fg}{local}{Orthogonal projection; its source or exact-target role is specified locally.}{orthogonal-projection}% Neutral historical Manim key; prefer Pexact/Psource.
\dkDeclareSymbol{Q}{Q}{fg}{local}{Orthogonal projection; its role is specified locally.}{orthogonal-projection}% Neutral historical Manim key; prefer Qtarget.
\dkDeclareSymbol[Two-projection intertwiner T = QP + (I - Q)(I - P), satisfying TP = QT and T* T = I - (P - Q) squared. For ||P - Q|| < 1 it is invertible; its unitary polar factor T |T| inverse is the canonical direct rotation.]
  {IntertwinerT}{T}{intertwiner}{glossary}{Two-projection intertwiner; its unitary polar factor is the acute direct rotation.}{linear-operator}
\dkDeclareSymbol[Planar pi/2-rotation matrix J_0 = ((0, -1), (1, 0)), with J_0 squared = -I_2 and J_0 transpose = -J_0. It defines the rotation C = cos(2 theta) I_2 + sin(2 theta) J_0 used to construct the real competitor.]
  {Jturn}{J_0}{turn}{glossary}{Planar rotation through $\pi/2$, with square minus the identity.}{matrix}
\dkDeclareSymbol[Partial isometry V_pol in the polar decomposition S = V_pol |S|, where |S| = sqrt(S* S). It is isometric on the orthogonal complement of kernel(S) and zero on kernel(S). For an invertible operator its polar factor is unitary.]
  {PolarU}{V_{\mathrm{pol}}}{intertwiner}{glossary}{Partial isometry in the polar decomposition; unitary when the operator is invertible.}{partial-isometry}
\dkDeclareSymbol[Planar rotation C = cos(2 theta) I_2 + sin(2 theta) J_0 through twice the principal angle. The competitor W_theta is orthogonally conjugate to C direct-sum I_2, so it rotates one two-plane and fixes the other.]
  {Cturn}{C}{turn}{glossary}{Planar rotation through twice the principal angle in the counterexample.}{matrix}
\dkDeclareSymbol[Orthogonal mixing matrix H_mix = (1/sqrt(2)) times the block matrix ((I_2, I_2), (I_2, -I_2)). It constructs W_theta = H_mix (C direct-sum I_2) H_mix transpose.]
  {Hmix}{H_{\mathrm{mix}}}{mix}{glossary}{Orthogonal mixing matrix in the counterexample.}{matrix}
\dkDeclareSymbol[Source reflection J_src = 2P - I: it acts as +I on the source subspace and -I on its orthogonal complement, so its square is I. In the matching proof stationarity gives J_src G J_src = -G for the skew-symmetric spectral gradient G.]
  {Jreflect}{J_{\mathrm{src}}}{trialprojection}{glossary}{Reflection equal to twice the source projection minus the identity.}{linear-operator}
\dkDeclareSymbol[Skew angle generator J_ang for a direct rotation D = exp(J_ang Theta). On each principal plane it is the oriented rotation through pi/2; it vanishes on the common fixed directions. It commutes with the full angle operator Theta.]
  {AngleGenerator}{J_{\mathrm{ang}}}{turn}{glossary}{Skew angle generator in the direct-rotation result summaries.}{linear-operator}
\dkDeclareSymbol{Dblock}{D}{mix}{local}{Auxiliary diagonal or block matrix in the counterexample.}{matrix}
\dkDeclareSymbol[Orthogonal projection onto the first principal plane span(e_0, e_3) in the four-dimensional counterexample. For Delta_W = I - W_(pi/4), the Ky Fan 2-norm of Delta_W Omega_1 equals 2.]
  {OmegaOne}{\Omega_1}{turn}{glossary}{First principal-plane projection in the failed Ky Fan comparison.}{orthogonal-projection}
\dkDeclareSymbol[Orthogonal projection onto the second principal plane span(e_1, e_2) in the four-dimensional counterexample. For Delta_W = I - W_(pi/4), the Ky Fan 2-norm of Delta_W Omega_2 equals 2, although the trace norm of Delta_W is only 2 sqrt(2).]
  {OmegaTwo}{\Omega_2}{intertwiner}{glossary}{Second principal-plane projection in the failed Ky Fan comparison.}{orthogonal-projection}
\dkDeclareSymbol{Dplanar}{D}{wantedrotation}{local}{Planar direct-rotation block.}{unitary-operator}
\dkDeclareSymbol{Hreflection}{H}{trialcompetitor}{local}{Planar reflection carrying the source line to the target.}{linear-operator}
% Unitarily invariant gauge family; N, M are NOT a norm versus its input.
\dkDeclareSymbol[Unitarily invariant norm N, or a symmetric operator-ideal norm on its domain of finiteness. It satisfies N(U T V) = N(T) for unitaries and N(B T) <= ||B||op N(T). In the sine theorem it is normalized by N(u v*) = ||u|| ||v|| and has Ky Fan dominance where both norms are finite. Both R and sin(Theta_0) must belong to its ideal.]
  {NormN}{N}{norm}{glossary}{Unitarily invariant norm on its domain of finiteness; in the sine theorem, normalized by $N(uv^*)=\|u\|\|v\|$.}{norm}
\dkDeclareSymbol[Sine-theta separation G_delta(B,C): finite interval/exterior separation in either spectral order, or ordered separation B <= a I and C >= (a+delta) I, or the reverse. The ordered alternatives allow both spectra to be unbounded.]
  {SineGap}{\mathcal G_\delta}{gap}{local}{Interval/exterior separation or either ordered separation, as in the sine-theta theorem.}{predicate}
\dkDeclareSymbol[Auxiliary unitarily invariant norm M witnessing the Q-norm property N(T) squared = M(T* T). For the Frobenius norm N, one can take M to be the trace norm.]
  {NormM}{M}{normaux}{local}{Auxiliary unitarily invariant norm used to define a Q-norm through the squared modulus.}{norm}
\dkDeclareSymbol[Q-norm property: a unitarily invariant norm N is a Q-norm when N(T) squared = M(T* T) for some unitarily invariant norm M. The Frobenius norm is an example, with M the trace norm; the repaired full-displacement theorem uses this norm class.]
  {NormQ}{Q}{normaux}{glossary}{Q-norm: its square is a unitarily invariant norm of the squared modulus.}{norm}
\dkDeclareSymbol[Singleton trace cost s_i = min(4 sin(theta_i/2), 2) for angle label i. It transports one real principal plane independently, choosing a rotation of cost 4 sin(theta_i/2) or a reflection of cost 2.]
  {SingletonCost}{s}{norm}{glossary}{Minimum trace displacement for one real principal plane: choose a rotation or a reflection.}{scalar-cost}
\dkDeclareSymbol[Paired trace cost e_ij = 4 sin((theta_i + theta_j)/2) for distinct angle labels i,j. It is attained by coupling their principal planes in a real four-dimensional block. The edge weight s_i + s_j - e_ij is the difference between the cost of two singletons and this paired cost.]
  {PairCost}{e}{norm}{glossary}{Trace displacement cost of coupling two principal planes in a four-dimensional block.}{scalar-cost}
\dkDeclareSymbol[Scalar phase cost f_p(t) = (2 sin(t/2)) to the power p, for 0 <= t <= pi. A real planar rotation through t contributes two copies, so its displacement power cost is 2 f_p(t).]
  {MatchingPhaseCost}{f_p}{norm}{local}{Displacement power cost for one scalar phase.}{scalar-cost}
\dkDeclareSymbol[Singleton power cost a_i(p) = min(2 f_p(theta_i), 2 to the power p). These are the rotation and reflection choices for one principal plane. At p=1 it equals the trace-cost weight s_i.]
  {SingletonPowerCost}{a}{norm}{local}{Minimum displacement power cost for one angle.}{scalar-cost}
\dkDeclareSymbol[Paired power cost b_ij(p) = min over d <= tau <= h of 2(f_p(h+tau) + f_p(h-tau)), where h=(theta_i+theta_j)/2 and d=abs(theta_i-theta_j)/2. This is the positive-determinant pair minimum; unrestricted block optimization also compares two singletons. At p=1 it equals e_ij.]
  {PairPowerCost}{b}{norm}{local}{Positive-determinant displacement power cost for a pair of angles.}{scalar-cost}
\dkDeclareSymbol[Matching M on the finite general graph of positive principal-angle labels. Each label belongs to at most one edge; unmatched labels use singleton transporters. Minimizing paired and singleton costs is equivalent to maximum-weight matching with edge weights g_ij = s_i + s_j - e_ij. The empty matching is allowed.]
  {Matching}{M}{normaux}{glossary}{Matching of positive principal-angle labels: each label is used in at most one pair.}{matching}
\dkDeclareSymbol[Displacement power sum Phi_p(W) = sum_j s_j(I - W)^p. The matching proof first minimizes it for 1 < p < 2, then lets p decrease to 1, where it becomes the trace norm ||I - W||_1.]
  {PowerCost}{\Phi}{norm}{local}{Sum of the p-th powers of displacement singular values.}{scalar-cost}
% Auxiliary matching-proof notation: neutral ink, with the same hover registry.
\dkDeclareSymbol[Off-diagonal gradient block G_01 mapping the source complement K perpendicular into K. Stationarity makes the diagonal blocks of G zero, so G = ((0,G_01),(-G_01 transpose,0)) in the source/complement decomposition.]
  {GradientOffDiagonal}{G_{01}}{fg}{local}{Complement-to-source block of the skew-symmetric gradient.}{linear-operator}
\dkDeclareSymbol[Orthogonal group O(n): real n by n matrices W satisfying W transpose times W = I. Its elements preserve Euclidean inner products and lengths. The notation O(2) refers to the same group in dimension two.]
  {OrthogonalGroup}{O(n)}{fg}{local}{Group of real orthogonal matrices.}{group}
\dkDeclareSymbol[Feasible transporters F = {W in O(n) : W(K) = L}. Equivalently W P_K W transpose = P_L. This nonempty compact set is the optimization domain throughout the matching proof.]
  {FeasibleTransport}{\mathcal F}{fg}{local}{Orthogonal maps carrying the source subspace onto the target.}{set}
\dkDeclareSymbol[Singular value s_j(T), the nonnegative square root of an eigenvalue of T transpose times T, counted with multiplicity. Summing s_j(I-W) gives the trace displacement; summing its p-th powers gives Phi_p(W). These singular values differ from the singleton weights s_i.]
  {TransportSingularValues}{\symglyph{SingularValueBase}_j}{fg}{local}{A displacement singular value in the matching proof.}{scalar}
\dkDeclareSymbol[Scalar singular-value magnitude s between zero and two for the displacement I-W of an orthogonal matrix. Uniform convergence of s to the power p to s on this interval permits the trace-norm limit. This variable differs from the rotation parameter s used earlier in the proof.]
  {DisplacementSize}{s}{fg}{local}{Scalar displacement singular-value magnitude in the uniform limit.}{scalar}
\dkDeclareSymbol[All-singleton cost C_0 = sum_i s_i, before any two angle labels are paired. A matching has cost C_0 minus the sum of its selected edge weights.]
  {SingletonBaseline}{C_0}{fg}{local}{Trace displacement when every angle uses its singleton transporter.}{scalar-cost}
\dkDeclareSymbol[Matching power cost C_p(M): sum of b_ij(p) over selected pairs plus a_i(p) over unmatched labels. Its minimum equals the minimum of Phi_p over feasible transporters. At p=1 it is C(M).]
  {MatchingPowerCost}{C_p}{fg}{local}{Total paired and singleton power cost of a matching.}{scalar-cost}
\dkDeclareSymbol[Invariant orthogonal block V_a preserved by the source projection, target projection, and minimizing transporter. The displacement cost adds over these blocks, and their principal-angle lists combine into the original list.]
  {MatchingBlock}{V_a}{fg}{local}{Simultaneous invariant block in the matching proof.}{linear-subspace}
\dkDeclareSymbol[Gradient-level space V_g: the eigenspace of -G squared for a positive eigenvalue g squared. It is preserved by G, W, and the source reflection J_src. Multiplication by i on it is G/g.]
  {GradientLevelSpace}{V_g}{fg}{local}{Positive gradient-level eigenspace.}{linear-subspace}
\dkDeclareSymbol[Eigenphase varphi of W on V_g viewed as a complex space with multiplication by i = G/g. It lies between zero and pi and satisfies f_p prime(varphi) = g. There are at most two such phases.]
  {GradientEigenphase}{\varphi}{fg}{local}{One of the at most two phases at a positive gradient level.}{angle}
\dkDeclareSymbol[Orthogonal projection E onto the first phase space of W on V_g. Here W = exp(i varphi_1) E + exp(i varphi_2) (I-E). In coordinates where the source reflection is conjugation, E = X + iY.]
  {GradientPhaseProjection}{E}{fg}{local}{Projection onto the first phase space in a gradient-level space.}{orthogonal-projection}
\dkDeclareSymbol[Real part X of the phase projection E = X + iY. It is a real symmetric matrix with spectrum in the closed interval from zero to one, satisfying X squared - Y squared = X and XY + YX = Y. This X differs from the source-coordinate and Sylvester matrices elsewhere in the paper.]
  {ProjectionRealPart}{X}{fg}{local}{Real symmetric part of the phase projection.}{linear-operator}
\dkDeclareSymbol[Imaginary part Y of the phase projection E = X + iY. It is a real skew-symmetric matrix. It maps the lambda eigenspace of X to the 1-lambda eigenspace and constructs the small invariant planes.]
  {ProjectionImaginaryPart}{Y}{fg}{local}{Real skew-symmetric imaginary part of the phase projection.}{linear-operator}
\dkDeclareSymbol[Eigenvalue lambda in the closed interval from zero to one of the real part X of the phase projection. It is local to the block construction and differs from the original operator's spectral parameter.]
  {ProjectionEigenvalue}{\lambda}{fg}{local}{Eigenvalue of the real part of the phase projection.}{scalar}
\dkDeclareSymbol[Plane coupling eta = sqrt(lambda (1-lambda)) for 0 < lambda < 1. If u is a unit lambda eigenvector of X, then v = Yu/eta completes its invariant real plane.]
  {ProjectionCoupling}{\eta}{fg}{local}{Normalization used to construct an invariant real plane.}{scalar}
\dkDeclareSymbol[Pair mean angle h = (theta_i + theta_j)/2. Together with d = (theta_i - theta_j)/2 after ordering theta_i >= theta_j, it parametrizes the two-angle block cost.]
  {PairMeanAngle}{h}{fg}{local}{Mean of the two ordered principal angles.}{angle}
\dkDeclareSymbol[Pair half-difference d = (theta_i - theta_j)/2 after ordering theta_i >= theta_j. It satisfies 0 <= d <= h and h+d <= pi/2.]
  {PairHalfDifference}{d}{fg}{local}{Half the difference of the two ordered principal angles.}{angle}
\dkDeclareSymbol[Scalar pair parameter tau in the closed interval from d to h. The paired power cost is the minimum of 2(f_p(h+tau) + f_p(h-tau)) over this interval. At p=1 a minimum is attained at tau=h.]
  {PairPhaseParameter}{\tau}{fg}{local}{Scalar parameter in the two-angle power-cost minimum.}{scalar}
\dkDeclareSymbol[Attaining rotation angle gamma = arccos(cos(tau)/cos(d)). Choosing U_0 = U_1 = R_gamma realizes phase parameters x=h and y=tau in the four-dimensional block.]
  {PairAttainingAngle}{\gamma}{fg}{local}{Rotation angle realizing a chosen scalar pair parameter.}{angle}
\dkDeclareSymbol[Four-dimensional direct rotation D = ((C,-S),(S,C)) in source and complementary principal coordinates, with C the diagonal principal cosines and S the diagonal principal sines. Every transporter on this block is D diag(U_0,U_1).]
  {PairDirectRotation}{D}{fg}{local}{Direct rotation on the two principal planes.}{unitary-operator}
\dkDeclareSymbol[Diagonal cosine block C = diag(cos(theta_i), cos(theta_j)) in the four-dimensional direct rotation D. This matrix differs from the scalar displacement cost C(M).]
  {PairCosineBlock}{C}{fg}{local}{Diagonal principal-cosine block.}{matrix}
\dkDeclareSymbol[Diagonal sine block S = diag(sin(theta_i), sin(theta_j)) in the four-dimensional direct rotation D. This matrix is local to the matching proof.]
  {PairSineBlock}{S}{fg}{local}{Diagonal principal-sine block.}{matrix}
\dkDeclareSymbol[Free planar orthogonal matrices U_0 and U_1 acting on the two-dimensional source and its two-dimensional complement, respectively. Multiplying D by diag(U_0,U_1) parametrizes every transporter on this four-dimensional block.]
  {PairFreeOrthogonal}{U}{fg}{local}{Free source or complementary planar orthogonal factor.}{unitary-operator}
\dkDeclareSymbol[Phase parameter x in the closed interval from h to pi-h for a positive-determinant four-dimensional transporter. Its eigenvalues are exp(plus or minus i(x+y)) and exp(plus or minus i(x-y)).]
  {PairPhaseX}{x}{fg}{local}{First phase parameter of a four-dimensional transporter.}{angle}
\dkDeclareSymbol[Phase parameter y in the closed interval from d to pi-d for a positive-determinant four-dimensional transporter. Together with x it gives rotational phases x+y and x-y, folded into the closed interval from zero to pi for the displacement cost.]
  {PairPhaseY}{y}{fg}{local}{Second phase parameter of a four-dimensional transporter.}{angle}
\dkDeclareSymbol[Scalar a = cos(x) in the four-dimensional phase calculation. Together with b = cos(y), it expresses the trace as 4ab and determines the reciprocal characteristic polynomial.]
  {PairPhaseCosineX}{a}{fg}{local}{Cosine of the first phase parameter in the trace calculation.}{scalar}
\dkDeclareSymbol[Scalar b = cos(y) in the four-dimensional phase calculation, paired with a = cos(x). These scalars are local to the characteristic-polynomial calculation.]
  {PairPhaseCosineY}{b}{fg}{local}{Cosine of the second phase parameter in the trace calculation.}{scalar}
\dkDeclareSymbol[Rotation parameter r: the two free orthogonal factors use planar rotation angles r+s and r-s. It appears in cos(y) = cos(d) cos(r) for two rotations.]
  {PairRotationR}{r}{fg}{local}{Half the sum of the two free planar rotation angles.}{angle}
\dkDeclareSymbol[Rotation parameter s: the two free orthogonal factors use planar rotation angles r+s and r-s. It appears in cos(x) = cos(h) cos(s) for two rotations. This scalar differs from the singular values and singleton costs.]
  {PairRotationS}{s}{fg}{local}{Half the difference of the two free planar rotation angles.}{angle}
\dkDeclareSymbol[Planar reflection J_2 = diag(1,-1). Every two-dimensional orthogonal matrix of determinant -1 can be written as a planar rotation times J_2.]
  {PairPlanarReflection}{J_2}{fg}{local}{Fixed planar reflection used to parametrize determinant -1 factors.}{unitary-operator}
\dkDeclareSymbol[Real planar rotation R_t = ((cos(t),-sin(t)),(sin(t),cos(t))). Its two displacement singular values are both 2 times abs(sin(t/2)). In the matching proof a rotation through t in the closed interval from zero to pi contributes 2 f_p(t).]
  {MatchingPlanarRotation}{\symglyph{PlanarRotationBase}_t}{fg}{local}{Planar rotation through the indicated angle.}{unitary-operator}
% Statistical application / operator provenance.
\dkDeclareSymbol{Sigma}{\Sigma}{old}{local}{Population covariance matrix.}{matrix}
\dkDeclareSymbol{SigmaHat}{\widehat\Sigma}{current}{local}{Estimated or perturbed covariance matrix.}{matrix}
\dkDeclareSymbol{YWSV}{V}{wanted}{local}{Orthonormal coordinates for the population eigenspace in the statistical application.}{isometric-embedding}
\dkDeclareSymbol{YWSVhat}{\widehat V}{trialframe}{local}{Estimated eigenspace frame in the statistical application.}{isometric-embedding}

% Figure-specific objects retain their own meanings when bare letters overlap.
\dkDeclareSymbol[Signed rotation angle phi of the displayed map on its invariant real two-plane. At theta = pi/4 the direct rotation turns its two principal planes by opposite signed angles of magnitude pi/4; the competitor turns its moving plane by pi/2 and fixes its other plane.]
  {PlaneRotationAngle}{\varphi}{sine}{local}{Signed rotation angle of the orthogonal map on the displayed invariant plane.}{angle}
\dkDeclareSymbol[Unit vector q in the displayed invariant plane. The map T is the direct rotation D or the competitor W in that panel; Tq is the image and (I - T)q = q - Tq is the displacement.]
  {FigureInput}{q}{ambient}{local}{Unit input to the orthogonal map shown in this invariant-plane panel.}{vector}
\dkDeclareSymbol[Image Tq of the panel unit input q under T = D or W. Both q and Tq lie on the unit circle; the green displacement arrow joins Tq to q, so it represents q - Tq.]
  {FigureImage}{Tq}{wantedvector}{local}{Image of the unit input under the orthogonal map shown in this panel.}{vector}
\dkDeclareSymbol[Displacement (I - T)q = q - Tq, drawn from the image Tq to the unit input q. On a plane rotated through phi its length is 2 |sin(phi/2)|; it vanishes on a fixed plane.]
  {DisplacementVector}{(I-T)q}{displacement}{local}{Difference between a unit input and its image under the displayed orthogonal map.}{vector}
\dkDeclareSymbol[First principal plane Pi_1 = span(e_0, e_3) for the source-target pair in the real four-dimensional counterexample. At theta = pi/4 the direct rotation acts on this plane by a quarter of pi, with displacement singular value sqrt(2 - sqrt(2)).]
  {CounterPrincipalOne}{\Pi_1}{fg}{glossary}{First principal plane in the real counterexample: the span of e0 and e3.}{linear-subspace}
\dkDeclareSymbol[Second principal plane Pi_2 = span(e_1, e_2) for the source-target pair in the real four-dimensional counterexample. The direct rotation acts by the same angle magnitude as on Pi_1, with the opposite sign in the displayed plane coordinates.]
  {CounterPrincipalTwo}{\Pi_2}{fg}{glossary}{Second principal plane in the real counterexample: the span of e1 and e2.}{linear-subspace}
\dkDeclareSymbol[Moving plane M = range(W - I) = span(e_0 + e_2, e_1 + e_3) for the competitor. W rotates it by 2 theta; at theta = pi/4 the two displacement singular values contributed by this plane are both sqrt(2).]
  {CounterMovingPlane}{M}{ambient}{glossary}{Moving plane of the competing transporter: the range of W minus the identity.}{linear-subspace}
\dkDeclareSymbol[Fixed plane Z = kernel(W - I) = span(e_0 - e_2, e_1 - e_3) for the competitor. W acts as the identity here, so this plane contributes two zero displacement singular values.]
  {CounterFixedPlane}{Z}{ambient}{glossary}{Fixed plane of the competing transporter: the kernel of W minus the identity.}{linear-subspace}
\dkDeclareSymbol[Direct-rotation displacement singular value a = 2 sin(pi/8) = sqrt(2 - sqrt(2)) at theta = pi/4. Each of its two principal planes contributes two copies, giving trace displacement 4a.]
  {CounterChordA}{a}{displacement}{glossary}{Direct-rotation displacement singular value in the counterexample: sqrt(2 - sqrt(2)).}{scalar}
\dkDeclareSymbol[Competitor displacement singular value b = 2 sin(pi/4) = sqrt(2) at theta = pi/4. Its moving plane contributes two copies; the fixed plane contributes two zeros, giving trace displacement 2b.]
  {CounterChordB}{b}{displacement}{glossary}{Competing-transporter displacement singular value in the counterexample: square root of 2.}{scalar}
\dkDeclareSymbol[Singular value s_j(I - T) of the full ambient displacement, ordered from largest to smallest. Each rotating real two-plane contributes two values 2 |sin(phi/2)|; fixed directions contribute zeros. Their sum is the trace cost.]
  {DisplacementSingular}{\symglyph{SingularValueBase}_j}{displacement}{local}{An individual singular value of the full displacement from the identity.}{scalar}
\dkDeclareSymbol[Edge weight g_ij = s_i + s_j - e_ij, the difference between the cost of two singletons and the paired cost. A matching has total cost C_0 minus the sum of its edge weights, with C_0 = sum_i s_i. Maximizing this sum gives the minimum transport cost; edges of negative weight may be omitted.]
  {MatchingEdgeWeight}{g_{ij}}{current}{glossary}{Edge weight: the cost of two singletons minus their paired cost.}{scalar-weight}
\dkDeclareSymbol[Total trace-norm displacement cost. For a transporter W, C(W) = ||I - W||_1, the sum of its displacement singular values. For a matching M, C(M) is the sum of its paired and unmatched singleton costs.]
  {TraceCost}{C}{norm}{glossary}{Trace-norm displacement cost of the indicated orthogonal transporter or matching.}{scalar-cost}

% Spectral values have their own key, distinct from an operator block.
\dkDeclareSymbol[Spectral parameter lambda of a self-adjoint operator. In a finite-dimensional eigenvalue list lambda_j(A) it denotes an individual eigenvalue; in A = integral lambda dE(lambda), it is the real integration variable.]{SpectralValue}{\lambda}{wantedspectrum}{local}{Real spectral parameter or individual eigenvalue of a self-adjoint operator.}{scalar}



% Conventional mathematical notation. Hover text applies even when ink is plain.
% These are local (non-glossary) unless promoted explicitly by the author.
% IdentityTwo and IdentityFour derive their I from Identity. Singular-value
% registrations derive s_j from SingularValueBase; SingletonCost has its own s.
\dkDeclareSymbol[Identity operator on the vector or Hilbert space in context. Its dimension is inferred from the expression and need not be finite.]{Identity}{I}{plain}{local}{Identity operator on the space in context.}{identity-operator}
\dkDeclareSymbol[Two-by-two identity matrix, acting on a two-dimensional coordinate space.]{IdentityTwo}{\symglyph{Identity}_2}{plain}{local}{Two-by-two identity matrix.}{identity-operator}
\dkDeclareSymbol[Four-by-four identity matrix, acting on a four-dimensional coordinate space.]{IdentityFour}{\symglyph{Identity}_4}{plain}{local}{Four-by-four identity matrix.}{identity-operator}
\dkDeclareSymbol[Two-by-two zero matrix. Unlike the bold zero vector, this is a linear operator/matrix with two rows and two columns.]{ZeroMatrixTwo}{0_2}{plain}{local}{Two-by-two zero matrix.}{matrix}
\dkDeclareSymbol[The constant pi, the ratio of the circumference of a circle to its diameter. Angles in this article use radians.]{PiConstant}{\pi}{plain}{local}{The constant pi (angles are in radians).}{constant}
\dkDeclareSymbol[The field of real numbers.]{RealField}{\mathbb R}{plain}{local}{The real numbers.}{field}
\dkDeclareSymbol[The field of complex numbers.]{ComplexField}{\mathbb C}{plain}{local}{The complex numbers.}{field}
\dkDeclareSymbol[Generic Hilbert space used in the two-projection discussion; distinct from the paper's ambient space symbol E in the sine-theta theorem.]{GeneralHilbertSpace}{\mathcal H}{plain}{local}{Generic Hilbert space in a local discussion.}{hilbert-space}
\dkDeclareSymbol[The spectrum sigma(T) of an operator T: the set of complex scalars for which T minus the scalar identity is not boundedly invertible. For self-adjoint T the spectrum is real.]{SpectrumSymbol}{\sigma}{plain}{local}{Spectrum of an operator.}{operator-notation}
\dkDeclareSymbol[Base letter for singular values, presently s. It is distinct from the base s for matching singleton costs.]{SingularValueBase}{s}{plain}{local}{Base letter for singular values.}{notation-base}
\dkDeclareSymbol[Ordered singular value s_j(T), equal to the j-th eigenvalue of |T| for a finite matrix; distinguish this from the matching singleton cost s_j.]{SingularValueS}{\symglyph{SingularValueBase}_j}{plain}{glossary}{The j-th singular value, under the currently printed s_j convention.}{scalar}
\dkDeclareSymbol[Alternate indexed singular-value notation sigma_j(T). The current manuscript also uses s_j(T); this registered alias is available for a later deliberate notation decision.]{SingularValueSigma}{\sigma_j}{plain}{local}{Indexed singular value under the sigma_j convention.}{scalar}
\dkDeclareSymbol[Operator norm ||T||_op, the supremum of ||Tx|| over unit vectors x.]{OperatorNormNotation}{\|\cdot\|_{\mathrm{op}}}{plain}{glossary}{Operator norm.}{norm}
\dkDeclareSymbol[Frobenius norm ||T||_F, the square root of the sum of squared matrix entries, equal to the finite-dimensional Hilbert--Schmidt norm.]{FrobeniusNormNotation}{\|\cdot\|_F}{plain}{glossary}{Frobenius norm.}{norm}
\dkDeclareSymbol[Trace norm ||T||_1, the sum of singular values when finite.]{TraceNormNotation}{\|\cdot\|_1}{plain}{glossary}{Trace (nuclear) norm.}{norm}
\dkDeclareSymbol[Hilbert--Schmidt norm ||T||_HS, the square root of the sum of squared singular values when finite.]{HilbertSchmidtNormNotation}{\|\cdot\|_{\mathrm{HS}}}{plain}{local}{Hilbert--Schmidt norm.}{norm}
\dkDeclareSymbol[Ky Fan j-norm, the sum of the largest j singular values.]{KyFanNormNotation}{\|\cdot\|_{(j)}}{plain}{glossary}{Ky Fan j-norm.}{norm}
\dkDeclareSymbol[Interval/exterior spectral separation predicate I_delta(B,C), specifying a central closed interval separated by delta from the other operator's spectrum.]{IntervalSeparation}{\mathcal I}{plain}{local}{Interval/exterior separation predicate.}{predicate}
\dkDeclareSymbol[Ordered spectral separation predicate O_delta(B,C), requiring one self-adjoint operator to lie below a spectral threshold and the other above it by at least delta.]{OrderedSeparation}{\mathcal O}{plain}{local}{Ordered spectral separation predicate.}{predicate}
\dkDeclareSymbol[Positive parameter epsilon controlling a small perturbation in the two-dimensional eigenvalue example.]{EigenvalueExampleEpsilon}{\varepsilon}{plain}{local}{Small positive parameter in the eigenvalue-example matrices.}{scalar}
\dkDeclareSymbol[Positive lower bound epsilon for the trial-coordinate map E_0, as in E_0* E_0 >= epsilon squared times the identity.]{FrameLowerBoundEpsilon}{\varepsilon}{plain}{local}{Lower-bound parameter for the trial-coordinate embedding.}{scalar}
\dkDeclareSymbol[Rotation angle vartheta of a planar orthogonal block in the matching-proof first-variation calculation.]{MatchingPlanarAngle}{\vartheta}{plain}{local}{Planar rotation angle in the matching-proof calculation.}{angle}


% Ordinary named mathematical operators already used throughout the manuscript.
% Keeping these here makes their spelling a single editorial choice.
\dkDeclareSymbol[Range of a linear operator: the set of its output vectors.]{RangeOperator}{\operatorname{ran}}{plain}{local}{Range (image) of an operator.}{operator-notation}
\dkDeclareSymbol[Kernel (null space) of a linear operator: inputs mapped to zero.]{KernelOperator}{\operatorname{ker}}{plain}{local}{Kernel (null space) of an operator.}{operator-notation}
\dkDeclareSymbol[Domain of an operator: the vectors on which the operator is defined, important for unbounded operators.]{DomainOperator}{\operatorname{dom}}{plain}{local}{Domain of an operator.}{operator-notation}
\DeclareRobustCommand{\DKRange}{\sym{RangeOperator}}
\DeclareRobustCommand{\DKKernel}{\sym{KernelOperator}}
\DeclareRobustCommand{\DKDom}{\sym{DomainOperator}}

% Norm-formatting macros: the operand changes by occurrence but the notation
% lives here. The outer hover describes the norm and does not add ink color.
\DeclareRobustCommand{\dkOperatorNorm}[1]{\like{OperatorNormNotation}{\|#1\|_{\mathrm{op}}}}
\DeclareRobustCommand{\dkTraceNorm}[1]{\like{TraceNormNotation}{\|#1\|_1}}
\DeclareRobustCommand{\dkFrobeniusNorm}[1]{\like{FrobeniusNormNotation}{\|#1\|_F}}
\DeclareRobustCommand{\dkHilbertSchmidtNorm}[1]{\like{HilbertSchmidtNormNotation}{\|#1\|_{\mathrm{HS}}}}
\DeclareRobustCommand{\dkKyFanNorm}[2]{\like{KyFanNormNotation}{\|#1\|_{(#2)}}}
\DeclareRobustCommand{\dkSchattenNorm}[2]{\like{SchattenNormNotation}{\|#1\|_{S_{#2}}}}

\dkDeclareSymbol[Schatten p-norm ||T||_{S_p}, obtained from the l^p-norm of the singular values and finite for operators in the Schatten p-class.]{SchattenNormNotation}{\|\cdot\|_{S_p}}{plain}{local}{Schatten p-norm.}{norm}


\dkDeclareSymbol[Positive modulus of an operator T: |T| = (T* T)^(1/2).]{OperatorModulusNotation}{|\cdot|}{plain}{glossary}{Positive modulus of an operator.}{operator-notation}
\dkDeclareSymbol[Linear span of the indicated vectors, i.e. the smallest linear subspace that contains them.]{SpanOperator}{\operatorname{span}}{plain}{local}{Linear span of vectors.}{operator-notation}
\dkDeclareSymbol[Diagonal or block-diagonal matrix/operator with the displayed entries or blocks.]{DiagonalOperator}{\operatorname{diag}}{plain}{local}{Diagonal or block-diagonal operator.}{linear-operator}
\dkDeclareSymbol[Trace of a trace-class operator, or ordinary matrix trace in finite dimensions.]{TraceOperator}{\operatorname{tr}}{plain}{local}{Trace of an operator.}{operator-notation}
\dkDeclareSymbol[Dimension of a vector space or subspace.]{DimensionOperator}{\dim}{plain}{local}{Dimension of a vector space.}{operator-notation}
\dkDeclareSymbol[Rank of a linear map or matrix: dimension of its range.]{RankOperator}{\operatorname{rank}}{plain}{local}{Rank of a linear map or matrix.}{operator-notation}
\dkDeclareSymbol[Distance between sets of spectral values: the infimum of pairwise distances.]{SpectralDistanceOperator}{\operatorname{dist}}{plain}{local}{Distance between two spectral sets.}{operator-notation}
\DeclareRobustCommand{\dkModulus}[1]{\like{OperatorModulusNotation}{|#1|}}

% Context-specific glyphs previously hidden in variants of other semantic keys.
% Keeping both the roles and the exact printed forms makes collisions explicit.
\dkDeclareSymbol[Nonzero vector x tested in the Rayleigh quotient, before imposing unit length.]{RayleighTestVector}{x}{trialvector}{local}{Test vector in the Rayleigh quotient.}{vector}
\dkDeclareSymbol[Positive cosine block C in the generic two-projection decomposition of Halmos; C and S commute and C squared plus S squared is the identity.]{HalmosCosineBlock}{C}{anglecos}{local}{Cosine block in the Halmos two-projection representation.}{linear-operator}
\dkDeclareSymbol[Positive sine block S in the generic two-projection decomposition of Halmos; C and S commute and C squared plus S squared is the identity.]{HalmosSineBlock}{S}{sineoperator}{local}{Sine block in the Halmos two-projection representation.}{linear-operator}
\dkDeclareSymbol[Riesz spectral projection P_Gamma defined by a contour Gamma that encloses selected spectral values.]{RieszProjection}{P_{\symglyph{RieszContour}}}{wantedprojection}{local}{Spectral projection associated with a contour Gamma.}{orthogonal-projection}
\dkDeclareSymbol[Historical Sylvester right-hand coefficient matrix A, in B X - X A = Y. Distinct from the current paper's reference operator A.]{HistoricalSylvesterRight}{A}{trialspectrum}{local}{Right-hand coefficient A in the historical Sylvester equation.}{self-adjoint-operator}
\dkDeclareSymbol[Scalar coefficient b in the Sylvester prototype (b-a)x=y, with separated scalar spectral values a and b.]{SylvesterScalarLeft}{b}{unwantedspectrum}{local}{Scalar left coefficient b in a Sylvester equation.}{scalar}
\dkDeclareSymbol[Scalar coefficient a in the Sylvester prototype (b-a)x=y.]{SylvesterScalarRight}{a}{trialspectrum}{local}{Scalar right coefficient a in a Sylvester equation.}{scalar}
\dkDeclareSymbol[Scalar unknown x in the Sylvester prototype (b-a)x=y.]{SylvesterScalarUnknown}{x}{sylvester}{local}{Scalar unknown x in a Sylvester equation.}{scalar}
\dkDeclareSymbol[Scalar right side y in the Sylvester prototype (b-a)x=y.]{SylvesterScalarRHS}{y}{resid}{local}{Scalar right side y in a Sylvester equation.}{scalar}
\dkDeclareSymbol[Source subspace K of a real Euclidean space in the matching proof; an orthogonal transporter sends K onto L.]{MatchingSourceSubspace}{K}{trial}{local}{Source subspace K in the matching proof.}{linear-subspace}
\dkDeclareSymbol[Target subspace L of a real Euclidean space in the matching proof; an orthogonal transporter sends K onto L.]{MatchingTargetSubspace}{L}{wanted}{local}{Target subspace L in the matching proof.}{linear-subspace}
\dkDeclareSymbol[Orthogonal projection P_L onto the matching proof's target subspace L.]{MatchingTargetProjection}{P_L}{wantedprojection}{local}{Orthogonal projection onto the matching target L.}{orthogonal-projection}
\dkDeclareSymbol[Real two-dimensional orthogonal group O(2): planar rotations and reflections.]{PlanarOrthogonalGroup}{O(2)}{fg}{local}{Orthogonal matrices of size two.}{group}
\dkDeclareSymbol[Base letter R used to denote planar rotation operators R_t in the matching proof.]{PlanarRotationBase}{R}{fg}{local}{Base letter for planar rotations.}{notation-base}
\dkDeclareSymbol[Planar rotation R_(r+s) through the sum of scalar phases r and s.]{MatchingPlanarRotationSum}{\symglyph{PlanarRotationBase}_{r+s}}{fg}{local}{Planar rotation through phase r+s.}{unitary-operator}
\dkDeclareSymbol[Planar rotation R_gamma through phase gamma in the block rotation calculation.]{MatchingPlanarRotationGamma}{\symglyph{PlanarRotationBase}_{\symglyph{MatchingPhaseGamma}}}{fg}{local}{Planar rotation through phase gamma.}{unitary-operator}

% Specific symbols in the historical and appendix examples. They retain their
% existing black-ink notation but receive locally accurate hover definitions.
\dkDeclareSymbol[Contour Gamma, a positively oriented closed curve in the resolvent set enclosing a selected isolated spectral cluster.]{RieszContour}{\Gamma}{plain}{local}{Contour used in the Riesz spectral-projection formula.}{contour}
\dkDeclareSymbol[Even integer nu in a tested Ky Fan comparison, with nu/2 principal-plane terms on the right side.]{KyFanEvenIndex}{\nu}{plain}{local}{Even integer in the principal-plane Ky Fan comparison.}{scalar}
\dkDeclareSymbol[Orthogonal projection Omega_k onto the k-th principal plane in the published displacement comparison.]{PrincipalPlaneProjectionBase}{\Omega}{plain}{local}{Projection onto the indicated principal plane.}{notation-base}
\dkDeclareSymbol[Scalar gap threshold gamma in the ordered Sylvester spectral separation B >= (gamma+delta)I and C <= gamma I.]{SylvesterGapThreshold}{\gamma}{plain}{local}{Spectral separation threshold in a Sylvester estimate.}{scalar}
\dkDeclareSymbol[Projection Omega_t from a net or sequence that converges strongly to the identity in the singular-value approximation lemma.]{StrongApproxProjection}{\Omega_t}{plain}{local}{Orthogonal projection converging strongly to the identity.}{orthogonal-projection}
\dkDeclareSymbol[First orthogonal projection Omega used in the two-block pinching inequality for a unitarily invariant norm.]{PinchingProjectionOne}{\Omega}{plain}{local}{First projection in the pinching inequality.}{orthogonal-projection}
\dkDeclareSymbol[Second orthogonal projection Upsilon used in the two-block pinching inequality for a unitarily invariant norm.]{PinchingProjectionTwo}{\Upsilon}{plain}{local}{Second projection in the pinching inequality.}{orthogonal-projection}
\dkDeclareSymbol[Finite-rank projection Gamma used to retain the image of a chosen finite-rank compression K Psi.]{CompressionProjectionImage}{\Gamma}{plain}{local}{Range projection in the finite-rank compression lemma.}{orthogonal-projection}
\dkDeclareSymbol[Finite-rank projection Psi used to select input vectors in a nearly maximizing singular-value compression.]{CompressionProjectionInput}{\Psi}{plain}{local}{Domain projection in the finite-rank compression lemma.}{orthogonal-projection}
\dkDeclareSymbol[Small positive tolerance eta in the finite-rank compression approximation estimate.]{CompressionTolerance}{\eta}{plain}{local}{Tolerance controlling the off-diagonal compression norm.}{scalar}

\dkDeclareSymbol[Rotation phase gamma in a planar block from the matching proof.]{MatchingPhaseGamma}{\gamma}{plain}{local}{Rotation phase gamma in the matching calculation.}{scalar}

% Two common indices with local mathematical meanings distinct from other deltas.
\dkDeclareSymbol[Kronecker delta: delta_ij is one when the indices i and j agree and zero otherwise. This is not the spectral-gap parameter delta.]{KroneckerDelta}{\delta}{plain}{local}{Kronecker delta: equals 1 for equal indices and 0 otherwise.}{scalar}
\dkDeclareSymbol[Sign variable epsilon taking the values -1 and +1 in the factorization over two phase choices.]{PhaseSignChoice}{\epsilon}{plain}{local}{Sign variable taking values minus one and plus one.}{scalar}

% Documented notation correspondence. The source column is a faithful
% transcription of each cited original context, NOT an asserted universal DK
% meaning. The original A_0 is deliberately listed for two distinct roles:
% its reference block and the comparison/trial operator. The source uses the
% same printed symbol for both; this paper is free to distinguish them.
\dkDeclareNotationSource{DK1970}{Davis and Kahan (1970), The Rotation of Eigenvectors by a Perturbation. III}{https://doi.org/10.1137/0707001}
\dkDeclareNotationSource{Seelmann2019}{Seelmann (2019), Semidefinite Perturbations in the Subspace Perturbation Problem}{https://arxiv.org/abs/1708.02463}
\dkDeclareNotationSource{Blechta2020}{Blechta (2020), Stability of Linear GMRES Convergence with Respect to Compact Perturbations}{https://arxiv.org/abs/2005.12960}
% Arguments: {key}{source ID}{source glyph}{section or equation}{role-specific qualification}
\dkDeclareNotationAlternative{A}{DK1970}{A}{\S1, (1.2)}{Reference self-adjoint operator.}
\dkDeclareNotationAlternative{H}{DK1970}{H}{\S1, (1.2)}{Additive perturbation.}
\dkDeclareNotationAlternative{Acurrent}{DK1970}{A+H}{\S1}{Perturbed operator.}
\dkDeclareNotationAlternative{E0}{DK1970}{E_0}{\S1, (1.1)}{Source isometric frame.}
\dkDeclareNotationAlternative{E1}{DK1970}{E_1}{\S1, (1.1)}{Reference complement frame.}
\dkDeclareNotationAlternative{F0}{DK1970}{F_0}{\S1}{Exact-target frame.}
\dkDeclareNotationAlternative{F1}{DK1970}{F_1}{\S1}{Exact-target complement frame.}
\dkDeclareNotationAlternative{ReferenceBlock0}{DK1970}{A_0}{\S1, (1.2)}{Reference source block.}
\dkDeclareNotationAlternative{ReferenceBlock1}{DK1970}{A_1}{\S1, (1.2)}{Reference complement block.}
\dkDeclareNotationAlternative{A0}{DK1970}{A_0}{\S1, residual discussion}{Residual comparison operator (the reference block in DK).}
% We deliberately leave Lambda0 and Lambda1 without DK source-glyph metadata:
% the available source transcription does not reliably preserve the source
% accents/decorations on the perturbed spectral-block symbols.
\dkDeclareNotationAlternative{Psource}{DK1970}{P}{\S1}{Reference/source projection.}
\dkDeclareNotationAlternative{Qtarget}{DK1970}{Q}{\S1}{Target projection.}
\dkDeclareNotationAlternative{Pexact}{DK1970}{Q}{\S1}{Exact-target projection.}
\dkDeclareNotationAlternative{B}{DK1970}{B}{\S1, (1.2)}{Off-diagonal coupling block.}
\dkDeclareNotationAlternative{Theta0}{DK1970}{\Theta_0}{\S1}{Source-side angle operator.}
\dkDeclareNotationAlternative{Theta1}{DK1970}{\Theta_1}{\S1}{Complementary angle operator.}
\dkDeclareNotationAlternative{Theta}{DK1970}{\Theta}{\S1}{Full angle operator.}
\dkDeclareNotationAlternative{R}{DK1970}{R}{\S1, residual discussion}{Trial residual.}
\dkDeclareNotationAlternative{Rrotation}{DK1970}{U}{\S1, (1.18)}{Canonical direct rotation.}
\dkDeclareNotationAlternative{Wcompetitor}{DK1970}{V}{\S4}{Competing unitary transporter.}
\dkDeclareNotationAlternative{delta}{DK1970}{\delta}{\S2}{Spectral separation parameter.}
% Contemporary alternatives are recorded separately and do not enter the DK
% correspondence table. Their citation locations must be individually checked.
\dkDeclareNotationAlternative{Eambient}{Seelmann2019}{\mathcal H}{p.~1}{Ambient Hilbert space.}
\dkDeclareNotationAlternative{SingularValueS}{Blechta2020}{\sigma_j}{Thm.~1.1}{Ordered singular values of a compact operator.}
 in paper.tex)
% or pass \def\DKMonochromeMath{1} before % One switch: false = journal style; true = colored symbols and author--year citations.
\newif\ifDKAuthorStyle

%\DKAuthorStylefalse
\DKAuthorStyletrue

\ifdefined\DKUseAuthorStyle\DKAuthorStyletrue\fi
\ifdefined\DKJournalStyle\DKAuthorStylefalse\fi
\ifDKAuthorStyle
  \PassOptionsToClass{sn-mathphys-ay}{sn-jnl}
\else
  \PassOptionsToClass{sn-mathphys-num}{sn-jnl}
\fi
\documentclass[pdflatex]{sn-jnl} % PDF engine option; compiled with LuaLaTeX.
% PDF version used for the manuscript and its symbol annotations.
\pdfvariable minorversion=7\relax

\usepackage{amsmath,amssymb,amsthm,mathtools}
\usepackage{booktabs,tabularx,array,longtable}
\usepackage[section]{placeins}
\usepackage{float}
\usepackage{microtype}
\usepackage{graphicx}
\usepackage{pgf} % Live vector figures with LaTeX symbol labels and tooltips
\usepackage{xurl}
\ifDKAuthorStyle
  \setcitestyle{authoryear,round,aysep={,}}
  \hypersetup{colorlinks=true,allcolors=blue}
\else
  \setcitestyle{numbers,square}
  \hypersetup{hidelinks}
\fi
\usepackage[nameinlink]{cleveref}
\usepackage[T1]{fontenc}
\usepackage{lmodern} % Scalable LaTeX lettering, shared with the Manim slides
\usepackage{xcolor}
% Draft controls. Set these independently for review, print, or arXiv versions.
\newif\ifDKAnonymous
\DKAnonymousfalse
\ifdefined\DKAnonymousReview\DKAnonymoustrue\fi
\newif\ifDKColoredSymbols
\ifDKAuthorStyle\DKColoredSymbolstrue\else\DKColoredSymbolsfalse\fi
\ifdefined\DKMonochromeMath\DKColoredSymbolsfalse\fi
% The colors retain the Manim glossary's light-theme mathematical roles.
% Semantic, multi-shade mathematical notation for the Davis--Kahan paper.
% The six core families follow visualizations/dkvis/palette.py (LIGHT).
% Subshades distinguish objects that are RELATED but not IDENTICAL.
% Values were chosen for legibility against white manuscript paper; some
% amber, green, and vermilion anchors are a little darker than the slides.
%
% Color control (independent of hover reminders):
%   \DKMathColorsfalse   (after % Semantic, multi-shade mathematical notation for the Davis--Kahan paper.
% The six core families follow visualizations/dkvis/palette.py (LIGHT).
% Subshades distinguish objects that are RELATED but not IDENTICAL.
% Values were chosen for legibility against white manuscript paper; some
% amber, green, and vermilion anchors are a little darker than the slides.
%
% Color control (independent of hover reminders):
%   \DKMathColorsfalse   (after \input{math_symbol_colors.tex} in paper.tex)
% or pass \def\DKMonochromeMath{1} before \input{paper.tex}, without editing.
% EVERY \sym, \like, and \dkrole obeys this switch, including the glossary.
% No color is required to identify a mathematical object: the glyphs, indices,
% definitions and adjacent prose also distinguish them.
%
% Families (base hue): exact/target=blue, trial/source=amber, exact unwanted=
% cyan, angles and their operators=magenta, residual/displacement=green,
% gaps=violet. Norm gauges=slate; provenance retains brown/steel/vermilion.
% P_ex is the exact-target projection in sine-theta; P is the source
% projection in transport. R denotes a residual and D a direct rotation.
% Semantic keys retain their roles even when the printed notation changes.
\newif\ifDKMathColors
\DKMathColorstrue
\ifdefined\DKMonochromeMath
  \DKMathColorsfalse
\fi

% Family anchors (print-compatible versions of the Manim light palette).
\definecolor{dkfg}{HTML}{1B1F24}
\definecolor{dkmuted}{HTML}{6B7280}
\definecolor{dkwanted}{HTML}{1F6FD1}
\definecolor{dktrial}{HTML}{A66705}
\definecolor{dkunwanted}{HTML}{007C91}
\definecolor{dksine}{HTML}{D12A7B}
\definecolor{dkresid}{HTML}{087D43}
\definecolor{dkgap}{HTML}{6E4BD8}
\definecolor{dkold}{HTML}{705519}
\definecolor{dkcurrent}{HTML}{054669}
\definecolor{dkperturb}{HTML}{BB491B}
\definecolor{dkrefuted}{HTML}{C62828}

% Related yet distinguishable shades. These are NOT new unrelated categories:
% objects inherit their parent family; a different shade signals a different
% mathematical object (e.g., a frame versus its subspace or its projection).
\definecolor{dkwantedframe}{HTML}{1254A9}
\definecolor{dkwantedprojection}{HTML}{08618D}
\definecolor{dkwantedrotation}{HTML}{4248B0}
\definecolor{dkwantedspectrum}{HTML}{345F8F}
\definecolor{dkwantedvector}{HTML}{1468AA}
\definecolor{dktrialframe}{HTML}{905600}
\definecolor{dktrialprojection}{HTML}{79511B}
\definecolor{dktrialcompetitor}{HTML}{A04D10}
\definecolor{dktrialspectrum}{HTML}{79591C}
\definecolor{dktrialvector}{HTML}{9B522A}
\definecolor{dkunwantedframe}{HTML}{006477}
\definecolor{dkunwantedspectrum}{HTML}{00546B}
\definecolor{dkunwantedprojection}{HTML}{145B6D}
\definecolor{dkunwantedvector}{HTML}{256879}
\definecolor{dkangleoperator}{HTML}{923478}
\definecolor{dksineoperator}{HTML}{AA275A}
\definecolor{dksineblock}{HTML}{B43F70}
\definecolor{dksylvester}{HTML}{804074}
\definecolor{dkanglemax}{HTML}{B32185}
\definecolor{dkanglecos}{HTML}{824276}
\definecolor{dkresidvector}{HTML}{086038}
\definecolor{dkdisplacement}{HTML}{226C50}
\definecolor{dkgaporiginal}{HTML}{4E4D89}
\definecolor{dkgapstat}{HTML}{7039A7}
\definecolor{dkgapinterval}{HTML}{574A82}
\definecolor{dkambient}{HTML}{36516A}
\definecolor{dknorm}{HTML}{475875}
\definecolor{dknormaux}{HTML}{66526E}
\definecolor{dkintertwiner}{HTML}{4A6280}
\definecolor{dkturn}{HTML}{47647D}
\definecolor{dkmix}{HTML}{505B65}

% Coloring primitives. "plain" suppresses color without suppressing hovers.
% All descendants can retain a semantic role through \like{key}{expression}.
\expandafter\def\csname dknocolor@plain\endcsname{1}
\newcommand{\dkrole}[2]{%
  \ifDKMathColors
    \ifcsname dknocolor@#1\endcsname #2\else\textcolor{dk#1}{#2}\fi
  \else #2\fi
}
% An optional plain-text reminder supplies definitions and types on hover.
% Required fields: {semantic key}{printed glyph}{color role}{glossary|local}
%                  {short definition}{mathematical type}. The "glossary" flag records whether an
% entry appears in the curated printed notation table; "local" does not.
% This registry is the sole authority for symbol meaning and printed form.
% Use role "plain" for an uncolored symbol that still carries hover text.
% Changing the glyph is a separate editorial decision: this audit preserves
% the current E, F, F_ex, G, and other existing printed choices.
\newcommand{\dkDeclareSymbol}[7][]{%
  \expandafter\def\csname dkgloss@#2\endcsname{#6}%
  \expandafter\def\csname dkglossaryflag@#2\endcsname{#5}%
  \if\relax\detokenize{#1}\relax
    \expandafter\pdfstringdef\csname dkpdfgloss@#2\endcsname{#7: #6}%
  \else
    \expandafter\pdfstringdef\csname dkpdfgloss@#2\endcsname{#7: \detokenize{#1}}%
  \fi
  \expandafter\def\csname dktype@#2\endcsname{#7}%
  \expandafter\def\csname dkglyph@#2\endcsname{#3}%
  \expandafter\def\csname dkrole@#2\endcsname{#4}%
  \expandafter\def\csname dksym@#2\endcsname{\dkrole{#4}{#3}}%
}
% Unstyled printed glyph, for variants constructed through \like.
% Example: \like{E0}{\symglyph{E0}^*} derives E_0^* from E0's
% one canonical glyph; retyping E_0 here would defeat central control.
\DeclareRobustCommand{\symglyph}[1]{\csname dkglyph@#1\endcsname}
\DeclareRobustCommand{\sym}[1]{%
  \ifcsname dksym@#1\endcsname
    \dkSymbolTooltip{#1}{\csname dksym@#1\endcsname}%
  \else
    \PackageError{dk-math-colors}{Unknown symbol key #1}{See math_symbol_colors.tex}%
  \fi
}
\DeclareRobustCommand{\like}[2]{%
  \ifcsname dkrole@#1\endcsname
    \dkSymbolTooltip{#1}{\dkrole{\csname dkrole@#1\endcsname}{#2}}%
  \else
    \PackageError{dk-math-colors}{Unknown role key #1}{See math_symbol_colors.tex}%
  \fi
}

% A short description belongs to the semantic key, not the bare letter.
% The same descriptions are available to the notation appendix.
\DeclareRobustCommand{\dkSymbolDescription}[1]{\csname dkgloss@#1\endcsname}
\DeclareRobustCommand{\dkSymbolGlossaryStatus}[1]{\csname dkglossaryflag@#1\endcsname}
\DeclareRobustCommand{\dkSymbolType}[1]{\csname dktype@#1\endcsname}

% Citable alternative notation is metadata, not a request to rename glyphs.
% The source identifier, the source's exact printed glyph, and the specific
% source location are all mandatory. No alternative is inferred from a name.
% One key may have alternatives from several works or multiple passages.
\newcommand{\dkDeclareNotationSource}[3]{%
  \expandafter\def\csname dksourcetitle@#1\endcsname{#2}%
  \expandafter\def\csname dksourceurl@#1\endcsname{#3}%
}
\newcommand{\dkDavisKahanRows}{}
\makeatletter
\newcommand{\dkDeclareNotationAlternative}[5]{%
  \ifcsname dkglyph@#1\endcsname\else
    \PackageError{dk-math-colors}{Unknown alternative key #1}{Declare the symbol first}%
  \fi
  \ifcsname dksourcetitle@#2\endcsname\else
    \PackageError{dk-math-colors}{Unknown notation source #2}{Declare the source first}%
  \fi
  \expandafter\def\csname dkaltglyph@#2@#1\endcsname{#3}%
  \expandafter\def\csname dkaltlocator@#2@#1\endcsname{#4}%
  \expandafter\def\csname dkaltnote@#2@#1\endcsname{#5}%
  \def\dk@source{#2}%
  \def\dk@original{DK1970}%
  \ifx\dk@source\dk@original
    \g@addto@macro\dkDavisKahanRows{\dkDavisKahanRow{#1}}%
  \fi
}
\makeatother
% A correspondence row is produced from the SAME registry records used by
% \sym and by the source-census script. It is never spelled a second time.
\newcommand{\dkDavisKahanRow}[1]{%
  \csname dkaltnote@DK1970@#1\endcsname &
  $\csname dkaltglyph@DK1970@#1\endcsname$ &
  $\sym{#1}$ &
  \csname dkaltlocator@DK1970@#1\endcsname\\
}

\newif\ifDKSymbolTooltips
\DKSymbolTooltipstrue
\ifdefined\DKNoSymbolTooltips\DKSymbolTooltipsfalse\fi
\newsavebox{\DKTooltipBox}
\newcount\DKTooltipCount
% Register only shipped widgets through hyperref's existing form-field
% machinery. Poppler/Okular need the catalog's field list, not just page
% annotations. The deferred aux write also discards measuring-only copies.
\makeatletter
\newcommand{\dkRegisterTooltipField}{%
  \expandafter\HyField@@AddToFields\expandafter{\number\pdffeedback lastannot}%
}
\makeatother
\AtBeginDocument{\ifDKSymbolTooltips\Form[NeedAppearances=false]\fi}
\AtEndDocument{\ifDKSymbolTooltips\endForm\fi}
% Native LuaLaTeX PDF widgets: no navigation action or JavaScript. Avoid
% pdfcomment's soulpos dependency, which previously broke local Vim builds.
% Measure in the CURRENT math style, then emit the original expression as-is:
% boxing the visible math would change spacing, scripts and line breaking.
\newcommand{\dkSymbolTooltip}[2]{%
  \ifDKSymbolTooltips
    \begingroup
      \DKSymbolTooltipsfalse % suppress nested annotations while measuring
      \ifmmode
        \edef\DKTooltipStyle{\ifcase\mathstyle
          \displaystyle\or\displaystyle\or\textstyle\or\textstyle\or
          \scriptstyle\or\scriptstyle\or\scriptscriptstyle\or\scriptscriptstyle\fi}%
        \sbox{\DKTooltipBox}{$\DKTooltipStyle#2$}%
      \else
        \leavevmode\sbox{\DKTooltipBox}{#2}%
      \fi
      \global\advance\DKTooltipCount by 1
      \pdfextension annot width\wd\DKTooltipBox
        height\ht\DKTooltipBox depth\dp\DKTooltipBox
        {/Subtype /Widget /FT /Btn /Ff 65536 /F 512
         /T (dk-symbol-\the\DKTooltipCount) /TU (\csname dkpdfgloss@#1\endcsname)
         /H /N /BS << /W 0 >> /Border [0 0 0]}%
      \dkRegisterTooltipField
      #2%
    \endgroup
  \else
    #2%
  \fi
}
\pdfstringdefDisableCommands{%
  \def\sym#1{\csname dkglyph@#1\endcsname}%
  \def\like#1#2{#2}%
  \def\symglyph#1{\csname dkglyph@#1\endcsname}%
}

% Declarations: [detailed hover], key, glyph, role, glossary|local, short text, type.
% A local entry has full semantic/hover support, but no notation-table row.
% When two keys have the same glyph, it is an auditable collision, not an
% instruction to change the printed notation automatically.
% This registry includes local proof notation. The notation appendix selects
% the symbols useful for reference; registration does not add a printed row.
% Manim's original symbol keys, now split into semantic subshades.
% Ambient and perturbed operators (provenance is independent of geometry).
\dkDeclareSymbol[Reference self-adjoint operator A on the ambient Hilbert space; the perturbed operator is Ahat = A + H. In a single-operator eigenproblem, A itself is the operator being approximated. An unbounded A acts only on domain(A).]
  {A}{A}{old}{glossary}{Reference self-adjoint operator before perturbation.}{self-adjoint-operator}
\dkDeclareSymbol{A1}{A_1}{old}{local}{Complementary block of the reference operator.}{self-adjoint-operator}
\dkDeclareSymbol{a1}{a_1}{old}{local}{Reference eigenvalue or complementary spectral value.}{scalar}
\dkDeclareSymbol[Reference block A_0: the restriction of A to the range of the reducing source projection P, expressed in source coordinates. In Theorem 8.1 it lies at or below alpha; in Theorem 8.2 its spectrum lies in the enlarged central interval. The perturbed blocks are denoted Lambda_0 and Lambda_1.]
  {ReferenceBlock0}{A_0}{old}{glossary}{Reference operator restricted to the source subspace in Section 8.}{self-adjoint-operator}
\dkDeclareSymbol[Reference block A_1: the restriction of A to the kernel of the reducing source projection P, expressed in complementary coordinates. In Theorem 8.1 it lies at or above alpha + delta in quadratic-form order.]
  {ReferenceBlock1}{A_1}{old}{glossary}{Reference operator restricted to the source complement in Section 8.}{self-adjoint-operator}
\dkDeclareSymbol[Self-adjoint ambient operator Ahat being studied; in the perturbation setting Ahat = A + H. Its action on a trial vector is compared with the proposed eigenvalue times that vector.]
  {Ahat}{\widehat A}{current}{local}{Operator under study; in the perturbation setting, the reference operator plus its perturbation.}{self-adjoint-operator}
\dkDeclareSymbol[Self-adjoint ambient operator Ahat being studied. In the perturbation setting Ahat = A + H. The trial residual is R = Ahat E_0 - E_0 A_tr, and the exact complementary block satisfies Ahat F_1 = F_1 Lambda_1 on its domain.]
  {Acurrent}{\widehat A}{current}{glossary}{Operator under study; in the perturbation setting, the reference operator plus its perturbation.}{self-adjoint-operator}
\dkDeclareSymbol[Additive self-adjoint perturbation H = Ahat - A. A bound on its norm can control the residual of an eigenspace of one operator relative to the other.]
  {H}{H}{perturb}{glossary}{Additive self-adjoint perturbation of the reference operator.}{self-adjoint-operator}
\dkDeclareSymbol{H0}{H_0}{perturb}{local}{Wanted diagonal block of the perturbation.}{self-adjoint-operator}
\dkDeclareSymbol{H1}{H_1}{perturb}{local}{Unwanted diagonal block of the perturbation.}{self-adjoint-operator}
\dkDeclareSymbol{B}{B}{perturb}{local}{Off-diagonal perturbation or coupling block.}{linear-operator}
\dkDeclareSymbol{b}{b}{perturb}{local}{Scalar perturbation or coupling coefficient.}{scalar}
% Trial/source family: subspace, frame, compressed operator, projector, vector.
\dkDeclareSymbol{U}{U}{wanted}{local}{Desired exact subspace in the local geometric notation.}{linear-subspace}% Manim's U remains wanted in local slide notation.
\dkDeclareSymbol{V}{V}{trial}{local}{Trial subspace in the local geometric notation.}{linear-subspace}% Manim's V remains trial in local slide notation.
\dkDeclareSymbol[Source subspace U to be carried onto the target V by an orthogonal or unitary map W. In the transport sections P projects onto U and Q onto V. The matching theorem also names the source subspace K.]
  {Usource}{U}{trial}{glossary}{Source subspace to be transported onto the target subspace.}{linear-subspace}
\dkDeclareSymbol[Trial frame E_0: F -> E, a linear isometry from trial-coordinate space F into ambient space E. Its range is the trial subspace; E_0* returns trial coordinates and E_0* E_0 = I_F. In rank one, E_0(x) = x v.]
  {E0}{E_0}{trialframe}{glossary}{Isometric embedding of the trial subspace; its adjoint returns trial coordinates.}{isometric-embedding}
\dkDeclareSymbol[Self-adjoint trial operator A_tr on trial-coordinate space F. The residual compares Ahat E_0 with E_0 A_tr. Rayleigh--Ritz chooses A_tr = E_0* Ahat E_0 when this compression is defined; the theorem allows other trial operators.]
  {A0}{A_{\mathrm{tr}}}{trialspectrum}{glossary}{Self-adjoint operator on trial coordinates; its eigenvalues are the trial spectral values.}{self-adjoint-operator}
\dkDeclareSymbol[Trial vector: a proposed approximate eigenvector. In the rank-one geometry v has unit length and E_0(x) = x v. Its residual for trial scalar rho is r = Ahat v - rho v.]
  {v}{v}{trialvector}{local}{Trial vector or vector in the trial subspace.}{vector}
\dkDeclareSymbol[Scalar trial operator in rank one, A_tr(x) = rho x. Figure 1 uses the Rayleigh value rho = <v, Ahat v> for unit v; then rho v is the projection of Ahat v onto the trial line and r is perpendicular to v.]
  {rho}{\rho}{trialspectrum}{local}{Trial eigenvalue or scalar trial operator.}{scalar}
\dkDeclareSymbol[Approximate eigenvalue assigned to a trial vector v. For the local operator A, its residual is A v - mu v. For unit v, the Rayleigh--Ritz choice is mu = <A v, v>.]
  {mu}{\mu}{trialspectrum}{local}{Trial spectral value.}{scalar}
\dkDeclareSymbol{a0}{a_0}{trialspectrum}{local}{Wanted trial eigenvalue or trial spectral value.}{scalar}
\dkDeclareSymbol[Ambient orthogonal or unitary transporter W with W(U) = V, equivalently W P W* = Q. Its full displacement is I - W on the whole ambient space. The real matching theorem minimizes its trace norm over every such orthogonal W.]
  {Wcompetitor}{W}{trialcompetitor}{glossary}{Orthogonal or unitary competitor carrying the source subspace onto the target.}{unitary-operator}
\dkDeclareSymbol[Orthogonal projection P onto the source subspace U in the transport problem; P = P* = P squared. A feasible orthogonal or unitary W satisfies W P W* = Q. Restricted displacement (I - W)P tests only source vectors.]
  {Psource}{P}{trialprojection}{glossary}{Orthogonal projection onto the source subspace in the transport problem.}{orthogonal-projection}
\dkDeclareSymbol{ProjSource}{P_U}{trialprojection}{local}{Orthogonal projection onto the source subspace in the transport problem.}{orthogonal-projection}
\dkDeclareSymbol[Vector x in the source subspace U. For principal vectors, unit x_j in U and y_j in V are chosen with <x_i, y_j> = 0 for i different from j and <x_j, y_j> = cos(theta_j).]
  {SourceVector}{x}{trialvector}{local}{Vector in the source subspace.}{vector}
\dkDeclareSymbol[Orthonormal coordinate frame X for the source subspace, so X* X = I. Together with a target frame Y, the singular values of X* Y are the principal cosines. In the four-dimensional counterexample X is the vertical stack of I_2 and 0.]
  {Xframe}{X}{trialframe}{glossary}{Orthonormal coordinate matrix for the source subspace in the counterexample.}{isometric-embedding}
\dkDeclareSymbol[Hilbert space F of trial coordinates. E_0 embeds F isometrically into ambient space E; A_tr and the positive sine operator sin(Theta_0) act on F, while the residual R maps F into E.]
  {Ftrial}{F}{trialframe}{glossary}{Hilbert space of trial coordinates.}{hilbert-space}
\dkDeclareSymbol[Ambient Hilbert space E over the real or complex scalars. It contains the trial subspace range(E_0) and the orthogonal exact subspaces range(F_0) and range(F_1). The source theorem permits infinite-dimensional separable E.]
  {Eambient}{E}{ambient}{glossary}{Ambient Hilbert space containing the trial and exact subspaces.}{hilbert-space}
% Zero and all-ones vectors are distinguished from scalar 0 and 1.
\dkDeclareSymbol[Zero vector of the vector or Hilbert space in context: the additive identity. The zero subspace consists only of this vector.]
  {ZeroVector}{\mathbf{0}}{fg}{glossary}{Zero vector; the additive identity of the ambient vector space.}{vector}
\dkDeclareSymbol[All-ones vector in a specified finite-dimensional coordinate space: every coordinate equals 1. This vector depends on the chosen coordinates and dimension.]
  {OnesVector}{\mathbf{1}}{fg}{local}{All-ones vector in a chosen finite-dimensional coordinate space.}{vector}
% Exact/wanted family: the target subspace, its frame, projection and rotation.
\dkDeclareSymbol[Target subspace V in the transport problem: feasible maps satisfy W(U) = V, or equivalently W P W* = Q. The matching theorem also names this subspace L.]
  {Vtarget}{V}{wanted}{glossary}{Target subspace in the transport problem.}{linear-subspace}
\dkDeclareSymbol[Exact wanted frame F_0: F_ex -> E, a linear isometry onto the desired exact spectral subspace. F_0* returns its coordinates. P_ex = F_0 F_0* is its ambient orthogonal projection; in Figure 1, F_0(h) = (h, 0).]
  {F0}{F_0}{wantedframe}{glossary}{Isometric coordinates for the desired exact spectral subspace; star denotes the adjoint.}{isometric-embedding}
\dkDeclareSymbol[Self-adjoint operator block on coordinates for the desired exact spectral subspace. Its spectrum consists of the wanted exact spectral values; in the corresponding decomposition Ahat F_0 = F_0 Lambda_0 on the operator domain.]
  {Lambda0}{\Lambda_0}{wantedspectrum}{local}{Self-adjoint block on the desired exact spectral subspace.}{self-adjoint-operator}
\dkDeclareSymbol{u}{u}{wantedvector}{local}{Vector in the desired exact subspace.}{vector}
\dkDeclareSymbol{f0}{f_0}{wantedvector}{local}{Component or coordinate vector in the desired exact subspace.}{vector}
\dkDeclareSymbol{lambda0}{\lambda_0}{wantedspectrum}{local}{Wanted exact eigenvalue or wanted spectral value.}{scalar}
\dkDeclareSymbol{lambdau}{\lambda_u}{wantedspectrum}{local}{Eigenvalue associated with a desired exact vector.}{scalar}
\dkDeclareSymbol[Exact wanted orthogonal projection P_ex: E -> E, with P_ex = F_0 F_0*. It retains the desired exact spectral component. Its complementary projection is P_ex_perp = I - P_ex = F_1 F_1*; the projected trial frame is S = P_ex_perp E_0.]
  {Pexact}{P_{\mathrm{ex}}}{wantedprojection}{glossary}{Orthogonal projection onto the desired exact subspace in the sine-theta theorem.}{orthogonal-projection}
\dkDeclareSymbol[Exact complementary orthogonal projection P_ex_perp = I - P_ex = F_1 F_1*: E -> E. It retains the unwanted exact component. In Figure 1, P_ex_perp v = S(1), whereas P_ex_perp r is the unwanted component of the residual.]
  {PexactPerp}{P_{\mathrm{ex}}^\perp}{unwantedprojection}{local}{Orthogonal projection onto the exact complementary subspace.}{orthogonal-projection}
\dkDeclareSymbol{ReflectionSigma}{\Sigma}{wantedprojection}{local}{Reflection associated with the desired exact spectral projection.}{unitary-operator}
\dkDeclareSymbol[Orthogonal projection Q onto the target subspace V in the transport problem. A feasible orthogonal or unitary W satisfies W P W* = Q, where P projects onto the source U.]
  {Qtarget}{Q}{wantedprojection}{glossary}{Orthogonal projection onto the target subspace in the transport problem.}{orthogonal-projection}
\dkDeclareSymbol{ProjTarget}{P_V}{wantedprojection}{local}{Orthogonal projection onto the target subspace in the transport problem.}{orthogonal-projection}
\dkDeclareSymbol[Canonical direct rotation D from the source U to the target V for an acute pair. It is the unitary polar factor of T = QP + (I - Q)(I - P), and acts by the principal angle on each principal plane. R is reserved for the trial residual.]
  {Rrotation}{D}{wantedrotation}{glossary}{Canonical direct rotation carrying the source subspace onto the target.}{unitary-operator}
\dkDeclareSymbol[Direct rotation D between the ranges of orthogonal projections P and Q, including the nonacute case. Relative to range(P) and kernel(P), its diagonal blocks are positive and its off-diagonal blocks are opposite adjoints. It exists when the two crossed defect spaces have equal dimension, and is unique when both defects vanish.]
  {SourceDirectRotation}{D}{wantedrotation}{local}{Direct rotation, including the nonacute case in the source-result table.}{unitary-operator}
\dkDeclareSymbol[Orthonormal coordinate frame Y for the target subspace, so Y* Y = I. In the counterexample Y is the vertical stack of cos(theta) I_2 and sin(theta) J_0; hence X* Y = cos(theta) I_2 and both principal angles equal theta.]
  {Yframe}{Y}{wantedframe}{glossary}{Orthonormal coordinate matrix for the target subspace in the counterexample.}{isometric-embedding}
\dkDeclareSymbol[Vector y in the target subspace V. For principal vectors, unit x_j in U and y_j in V are paired so that <x_j, y_j> = cos(theta_j); distinct pairs have zero cross inner product.]
  {TargetVector}{y}{wantedvector}{local}{Vector in the target subspace.}{vector}
\dkDeclareSymbol[Hilbert space F_ex of exact wanted coordinates. The isometry F_0: F_ex -> E identifies it with the desired exact spectral subspace. This space differs from H, the additive perturbation.]
  {Hexact}{F_{\mathrm{ex}}}{wantedframe}{glossary}{Hilbert space of coordinates for the desired exact spectral subspace.}{hilbert-space}% Coordinate Hilbert space (not perturbation).
% Exact/unwanted family; the trial error operators are MAGENTA, not cyan.
\dkDeclareSymbol{Uperp}{U^{\perp}}{unwanted}{local}{Orthogonal complement of the desired exact subspace in the local notation.}{linear-subspace}
\dkDeclareSymbol[Exact complementary frame F_1: G -> E, a linear isometry onto the orthogonal complement of range(F_0). F_1* returns unwanted coordinates; F_1 F_1* = I - P_ex. In Figure 1, F_1(g) = (0, g).]
  {F1}{F_1}{unwantedframe}{glossary}{Isometric coordinates for the exact complementary subspace; star denotes the adjoint.}{isometric-embedding}
\dkDeclareSymbol[Self-adjoint unwanted block Lambda_1 on G, satisfying Ahat F_1 = F_1 Lambda_1 on its domain. Its spectrum is separated from the trial spectrum of A_tr by the specified positive delta. Figure 1 uses the scalar block 0.7.]
  {Lambda1}{\Lambda_1}{unwantedspectrum}{glossary}{Self-adjoint operator on exact complementary coordinates: the unwanted spectral block.}{self-adjoint-operator}
% Spectral measure and set (distinct from the ambient space E and sine block S).
\dkDeclareSymbol[Projection-valued spectral measure E associated with a self-adjoint operator A: A = integral lambda dE(lambda). For a Borel set S, E(S) is the orthogonal projection onto the corresponding spectral subspace.]
  {SpectralMeasure}{E}{wantedprojection}{glossary}{Projection-valued spectral measure of a self-adjoint operator.}{projection-valued-measure}
\dkDeclareSymbol[Borel subset S of the real spectral axis. The spectral projection E(S) selects the part of a self-adjoint operator with spectral values in S; its range is a reducing subspace.]
  {SpectralSet}{S}{wantedspectrum}{glossary}{Borel spectral set used to select a spectral subspace.}{set}
% Generic Sylvester coefficients and right side, by source/target role.
\dkDeclareSymbol[Left coefficient B of the Sylvester map X -> B X - X C; B acts on the output space of X. For the sine bound, B = Lambda_1 on the unwanted-coordinate space G.]
  {SylvesterLeft}{B}{unwantedspectrum}{glossary}{Left coefficient of the Sylvester equation; corresponds to the unwanted spectral block.}{self-adjoint-operator}
\dkDeclareSymbol[Right coefficient C of the Sylvester map X -> B X - X C; C acts on the input space of X. For the sine bound, C = A_tr on trial space F. The historical equation also denotes this coefficient by A.]
  {SylvesterRight}{C}{trialspectrum}{glossary}{Right coefficient of the Sylvester equation; corresponds to the trial operator.}{self-adjoint-operator}
\dkDeclareSymbol[Right-hand side Y in B X - X C = Y, with the same input and output spaces as X. For the sine bound, Y = F_1* R: F -> G, the residual expressed in unwanted exact coordinates.]
  {SylvesterRHS}{Y}{resid}{glossary}{Right-hand side of the Sylvester equation, obtained by projecting the residual.}{linear-operator}
\dkDeclareSymbol{w}{w}{unwantedvector}{local}{Vector in an unwanted or complementary spectral subspace.}{vector}
\dkDeclareSymbol{f1}{f_1}{unwantedvector}{local}{Component or coordinate vector in the unwanted exact subspace.}{vector}
\dkDeclareSymbol{lambda1}{\lambda_1}{unwantedspectrum}{local}{Unwanted exact eigenvalue or complementary spectral value.}{scalar}
\dkDeclareSymbol{lambdaw}{\lambda_w}{unwantedspectrum}{local}{Eigenvalue associated with an unwanted exact vector.}{scalar}
\dkDeclareSymbol[Hilbert space G of exact complementary coordinates. F_1: G -> E embeds it onto the unwanted exact subspace; Lambda_1 acts on G and X = F_1* E_0 maps trial coordinates into G.]
  {Gunwanted}{G}{unwantedframe}{glossary}{Hilbert space of exact complementary coordinates.}{hilbert-space}
\dkDeclareSymbol{E1}{E_1}{unwantedframe}{local}{Isometric coordinates for the complementary trial subspace.}{isometric-embedding}
\dkDeclareSymbol{Vperp}{V^{\perp}}{unwanted}{local}{Orthogonal complement of the trial subspace in the local notation.}{linear-subspace}
% Angle and error geometry, including distinct rectangular and positive maps.
\dkDeclareSymbol[Individual principal angle, between 0 and pi/2. For orthonormal subspace frames X and Y, the singular values of X* Y are the principal cosines. In Figure 1, theta is the angle between the unit trial vector v and the exact wanted line.]
  {theta}{\theta}{sine}{glossary}{Principal angle between the subspaces under comparison.}{angle}
\dkDeclareSymbol[Full positive angle operator Theta = diag(Theta_0, Theta_1) for the source/target decompositions. Theta_0 compares the source with the target; Theta_1 compares their orthogonal complements. Its operator norm is the largest spectral angle. For positive c, the conditions Theta <= c and Theta < c mean operator norm at most c and strictly below c, respectively.]
  {Theta}{\Theta}{angleoperator}{glossary}{Full angle operator, the diagonal sum of the source and complementary angle blocks.}{angle-operator}
\dkDeclareSymbol[Directed principal-angle operator Theta_0 on source or trial coordinates. For a trial frame E_0 and exact wanted projection P_ex, cos(Theta_0) = sqrt(E_0* P_ex E_0), and sin(Theta_0) = sqrt(I_F - E_0* P_ex E_0). In the two-projection geometry it comes from the positive cosine compression P Q P restricted to range(P).]
  {Theta0}{\Theta_0}{angleoperator}{glossary}{Directed principal-angle operator acting on source or trial coordinates.}{angle-operator}
\dkDeclareSymbol[Complementary directed angle operator Theta_1, obtained from the projections I-P and I-Q by the same cosine-compression construction as Theta_0. It acts on the source complement, or its isometric coordinates. The full angle operator is diag(Theta_0, Theta_1).]
  {Theta1}{\Theta_1}{angleoperator}{glossary}{Directed angle operator comparing the orthogonal complements.}{angle-operator}
\dkDeclareSymbol[Right-angle source defect D_0 = range(P) intersect kernel(Q). These are source vectors orthogonal to the target. The full-angle conclusion in the Section 8 smallness criterion assumes dim(D_0) = dim(D_1); in Lean their identification is CrossedDefectsEquivalent.]
  {CrossedSourceDefect}{D_0}{fg}{glossary}{Source directions perpendicular to the target subspace.}{linear-subspace}
\dkDeclareSymbol[Right-angle target defect D_1 = kernel(P) intersect range(Q). These are target vectors orthogonal to the source. Equality of its Hilbert dimension with that of D_0 is assumed for the full-angle conclusion of Theorem 8.2.]
  {CrossedTargetDefect}{D_1}{fg}{glossary}{Target directions perpendicular to the source subspace.}{linear-subspace}
\dkDeclareSymbol[Positive complementary cosine block C_1 = cos(Theta_1). The reference complement is identified with the perturbed complement by the direct rotation D, so A_1 and Lambda_1 act on common coordinates. Theorem 8.1 compares A_1 - alpha I with C_1 (Lambda_1 - alpha I) C_1 as quadratic forms.]
  {ComplementCosineBlock}{C_1}{anglecos}{local}{Cosine of the complementary angle block in the spectral-repulsion inequality.}{self-adjoint-operator}
\dkDeclareSymbol[Scalar sine of a principal angle. For a unit principal trial vector v and exact wanted projection P_ex, sin(theta) = ||(I - P_ex)v||. In Figure 1 it is the length of the pink leg S(1).]
  {sintheta}{\sin\theta}{sine}{local}{Sine of an individual principal angle: the length of the complementary component of a unit principal trial vector.}{scalar}
\dkDeclareSymbol[Positive sine operator sin(Theta_0) = sqrt(I_F - E_0* P_ex E_0): F -> F. It equals |S| for S = (I - P_ex)E_0: F -> E. Its eigenvalues in finite dimensions are the principal sines; the sine theorem bounds its ideal norm by N(R)/delta.]
  {sinTheta0}{\sin\Theta_0}{sineoperator}{glossary}{Positive sine operator on trial coordinates; the modulus of the rectangular complementary error.}{angle-operator}
\dkDeclareSymbol[Rectangular complementary error S = (I - P_ex) E_0: F -> E. It records how much each trial vector leaks outside the exact wanted subspace. Its modulus |S| = sqrt(S* S) is sin(Theta_0): F -> F. In rank one S(1) = P_ex_perp v.]
  {S}{S}{sineblock}{glossary}{Rectangular complementary error: the trial frame projected onto the unwanted exact subspace.}{linear-operator}
\dkDeclareSymbol[Complementary-coordinate map X = F_1* E_0: F -> G, with S = F_1 X. It satisfies Lambda_1 X - X A_tr = F_1* R on domain(A_tr). Isometric multiplication gives N(X) = N(S); the separate modulus identity abs(S) = sin(Theta_0) gives N(S) = N(sin(Theta_0)), whenever these norms are finite.]
  {X}{X}{sylvester}{glossary}{Solution of the Sylvester equation; for the sine bound, the trial frame in unwanted exact coordinates.}{linear-operator}
\dkDeclareSymbol{alphaAngle}{\alpha}{anglemax}{local}{Largest principal angle.}{angle}
\dkDeclareSymbol[Scalar c = cos(theta) in the counterexample. In the general two-projection decomposition, the same role denotes a positive cosine operator C; the cosine and sine blocks commute and satisfy C squared + S squared = I.]
  {CosAngle}{c}{anglecos}{local}{Cosine of a principal angle or the corresponding positive cosine block.}{scalar-or-operator}
\dkDeclareSymbol[Scalar s = sin(theta) in the counterexample. In the general two-projection decomposition, the same role denotes a positive sine operator S; it commutes with the positive cosine block C and satisfies C squared + S squared = I.]
  {SinAngle}{s}{sineoperator}{local}{Sine of a principal angle or the corresponding positive sine block.}{scalar-or-operator}
\dkDeclareSymbol[Positive small-angle parameter in the local subspace comparison. It controls the allowed angular deviation; it is an angle parameter rather than a perturbation-operator norm.]
  {EpsilonAngle}{\varepsilon}{anglemax}{local}{Angle parameter controlling the maximum subspace deviation.}{angle}
\dkDeclareSymbol{PhiAngle}{\varphi}{angleoperator}{local}{Auxiliary angle or angle operator in the two-subspace comparison.}{angle}
\dkDeclareSymbol{Bblock}{B}{sineblock}{local}{Off-diagonal block in the two-subspace decomposition.}{linear-operator}
\dkDeclareSymbol{Cblock}{C}{sylvester}{local}{Complementary or Sylvester block in the two-subspace decomposition.}{linear-operator}
\dkDeclareSymbol{Ablock}{A}{anglecos}{local}{Positive cosine-related block in the two-subspace decomposition.}{linear-operator}
% Residual/error-displacement family.
\dkDeclareSymbol[Residual vector r = Ahat v - rho v = R(1) for the rank-one frame E_0(x) = x v. Figure 1 draws it from rho v to Ahat v. For unit v and the Rayleigh value rho = <v, Ahat v>, <v, r> = 0; a general trial scalar need not give orthogonality.]
  {r}{r}{residvector}{local}{Residual vector of an approximate eigenpair.}{vector}
\dkDeclareSymbol[Bounded residual R: F -> E, agreeing with Ahat E_0 - E_0 A_tr on domain(A_tr). The domain condition requires E_0(domain(A_tr)) inside domain(Ahat). R = 0 means that E_0 intertwines A_tr with Ahat; its norm measures the defect.]
  {R}{R}{resid}{glossary}{Bounded residual of the trial eigenproblem, agreeing with the operator difference on the trial domain.}{linear-operator}
\dkDeclareSymbol[Full displacement Delta_W = I - W on the ambient space. In the failed Ky Fan comparison W = W_(pi/4); Delta_W has singular values sqrt(2), sqrt(2), 0, 0. Multiplying by a principal-plane projection restricts its input to that plane.]
  {Kdisplacement}{\Delta_W}{displacement}{glossary}{Full displacement of a transporter from the identity.}{operator-valued-function}
\dkDeclareSymbol[Squared full displacement Dcal(W) = (I-W)* (I-W), a positive operator. Proposition 4.3 compares its unitarily invariant norm for a direct rotation D with that for every feasible transporter W.]
  {SquaredDisplacement}{\mathcal D}{displacement}{local}{Squared full displacement as a function of the transporter.}{operator-valued-function}
% Matching proof: the gradient and its spectral level are not a space or gap.
\dkDeclareSymbol[Skew-symmetric spectral gradient G of Phi_p. It commutes with W and represents first variations by d/dt Phi_p(W exp(tB)) at t = 0 equal to -tr(G B) for skew-symmetric B. Feasible stationarity gives J_src G J_src = -G.]
  {SpectralGradient}{G}{displacement}{local}{Skew-symmetric spectral gradient of the displacement power sum in the matching proof.}{linear-operator}
\dkDeclareSymbol[Positive spectral level g of the skew-symmetric gradient G: on this invariant subspace, -G squared = g squared I. Consequently G/g squares to -I and supplies the complex structure used in the real matching proof.]
  {GradientLevel}{g}{displacement}{local}{Positive gradient level; the gradient squared is minus this level squared times the identity.}{scalar}
\dkDeclareSymbol[Complex structure I_g = G/g on the positive gradient-level space V_g. It is a real isometry with square -I. Multiplication by a+ib means a times identity plus b times I_g; source reflection J_src reverses I_g and acts as conjugation.]
  {GradientComplexStructure}{\mathcal I_g}{displacement}{local}{Complex multiplication constructed from a positive gradient level.}{linear-operator}
\dkDeclareSymbol[Real skew-symmetric matrix B, with B transpose = -B, giving the orthogonal variation W exp(tB). If B commutes with the source projection P, this curve preserves the constraint W(U) = V.]
  {SkewDirection}{B}{ambient}{local}{Skew-symmetric direction for a first variation of an orthogonal transporter.}{linear-operator}
% Spectral separation family: gap and original gap are NOT identical.
\dkDeclareSymbol[Positive spectral separation in the stated configuration. In the sine theorem it separates the trial operator from the exact complementary block; Theorem 8.1 uses it for the ordered reference gap, and Theorem 8.2 for the perturbed central-band separation. Figure 1 uses delta = abs(rho - 0.7).]
  {delta}{\delta}{gap}{glossary}{Positive spectral separation in the stated gap configuration.}{scalar}
\dkDeclareSymbol{g}{g}{gaporiginal}{local}{Original spectral gap of the reference operator; distinct from the certified sine-theta separation.}{scalar}
\dkDeclareSymbol{Delta}{\Delta}{gapstat}{local}{Population or statistical spectral gap.}{scalar}
\dkDeclareSymbol[Spectral endpoint alpha. In the interval/exterior configuration, one spectrum lies between beta and alpha, the other at or below beta - delta or at or above alpha + delta. Theorem 8.2 places the selected perturbed block in the central interval. In Theorem 8.1, alpha bounds the lower reference block and alpha + delta bounds the upper block from below.]
  {alphaEndpoint}{\alpha}{gapinterval}{glossary}{Upper endpoint of the central band, or lower-block cut in the ordered gap.}{scalar}
\dkDeclareSymbol[Lower endpoint beta of the finite interval from beta to alpha in the gap hypothesis. One spectrum lies in that interval; the other lies at or below beta - delta or at or above alpha + delta.]
  {betaEndpoint}{\beta}{gapinterval}{glossary}{Lower endpoint of the finite interval in the interval/exterior gap hypothesis.}{scalar}
% Spectral/projector and auxiliary matrix objects, context-scoped.
\dkDeclareSymbol{P}{P}{fg}{local}{Orthogonal projection; its source or exact-target role is specified locally.}{orthogonal-projection}% Neutral historical Manim key; prefer Pexact/Psource.
\dkDeclareSymbol{Q}{Q}{fg}{local}{Orthogonal projection; its role is specified locally.}{orthogonal-projection}% Neutral historical Manim key; prefer Qtarget.
\dkDeclareSymbol[Two-projection intertwiner T = QP + (I - Q)(I - P), satisfying TP = QT and T* T = I - (P - Q) squared. For ||P - Q|| < 1 it is invertible; its unitary polar factor T |T| inverse is the canonical direct rotation.]
  {IntertwinerT}{T}{intertwiner}{glossary}{Two-projection intertwiner; its unitary polar factor is the acute direct rotation.}{linear-operator}
\dkDeclareSymbol[Planar pi/2-rotation matrix J_0 = ((0, -1), (1, 0)), with J_0 squared = -I_2 and J_0 transpose = -J_0. It defines the rotation C = cos(2 theta) I_2 + sin(2 theta) J_0 used to construct the real competitor.]
  {Jturn}{J_0}{turn}{glossary}{Planar rotation through $\pi/2$, with square minus the identity.}{matrix}
\dkDeclareSymbol[Partial isometry V_pol in the polar decomposition S = V_pol |S|, where |S| = sqrt(S* S). It is isometric on the orthogonal complement of kernel(S) and zero on kernel(S). For an invertible operator its polar factor is unitary.]
  {PolarU}{V_{\mathrm{pol}}}{intertwiner}{glossary}{Partial isometry in the polar decomposition; unitary when the operator is invertible.}{partial-isometry}
\dkDeclareSymbol[Planar rotation C = cos(2 theta) I_2 + sin(2 theta) J_0 through twice the principal angle. The competitor W_theta is orthogonally conjugate to C direct-sum I_2, so it rotates one two-plane and fixes the other.]
  {Cturn}{C}{turn}{glossary}{Planar rotation through twice the principal angle in the counterexample.}{matrix}
\dkDeclareSymbol[Orthogonal mixing matrix H_mix = (1/sqrt(2)) times the block matrix ((I_2, I_2), (I_2, -I_2)). It constructs W_theta = H_mix (C direct-sum I_2) H_mix transpose.]
  {Hmix}{H_{\mathrm{mix}}}{mix}{glossary}{Orthogonal mixing matrix in the counterexample.}{matrix}
\dkDeclareSymbol[Source reflection J_src = 2P - I: it acts as +I on the source subspace and -I on its orthogonal complement, so its square is I. In the matching proof stationarity gives J_src G J_src = -G for the skew-symmetric spectral gradient G.]
  {Jreflect}{J_{\mathrm{src}}}{trialprojection}{glossary}{Reflection equal to twice the source projection minus the identity.}{linear-operator}
\dkDeclareSymbol[Skew angle generator J_ang for a direct rotation D = exp(J_ang Theta). On each principal plane it is the oriented rotation through pi/2; it vanishes on the common fixed directions. It commutes with the full angle operator Theta.]
  {AngleGenerator}{J_{\mathrm{ang}}}{turn}{glossary}{Skew angle generator in the direct-rotation result summaries.}{linear-operator}
\dkDeclareSymbol{Dblock}{D}{mix}{local}{Auxiliary diagonal or block matrix in the counterexample.}{matrix}
\dkDeclareSymbol[Orthogonal projection onto the first principal plane span(e_0, e_3) in the four-dimensional counterexample. For Delta_W = I - W_(pi/4), the Ky Fan 2-norm of Delta_W Omega_1 equals 2.]
  {OmegaOne}{\Omega_1}{turn}{glossary}{First principal-plane projection in the failed Ky Fan comparison.}{orthogonal-projection}
\dkDeclareSymbol[Orthogonal projection onto the second principal plane span(e_1, e_2) in the four-dimensional counterexample. For Delta_W = I - W_(pi/4), the Ky Fan 2-norm of Delta_W Omega_2 equals 2, although the trace norm of Delta_W is only 2 sqrt(2).]
  {OmegaTwo}{\Omega_2}{intertwiner}{glossary}{Second principal-plane projection in the failed Ky Fan comparison.}{orthogonal-projection}
\dkDeclareSymbol{Dplanar}{D}{wantedrotation}{local}{Planar direct-rotation block.}{unitary-operator}
\dkDeclareSymbol{Hreflection}{H}{trialcompetitor}{local}{Planar reflection carrying the source line to the target.}{linear-operator}
% Unitarily invariant gauge family; N, M are NOT a norm versus its input.
\dkDeclareSymbol[Unitarily invariant norm N, or a symmetric operator-ideal norm on its domain of finiteness. It satisfies N(U T V) = N(T) for unitaries and N(B T) <= ||B||op N(T). In the sine theorem it is normalized by N(u v*) = ||u|| ||v|| and has Ky Fan dominance where both norms are finite. Both R and sin(Theta_0) must belong to its ideal.]
  {NormN}{N}{norm}{glossary}{Unitarily invariant norm on its domain of finiteness; in the sine theorem, normalized by $N(uv^*)=\|u\|\|v\|$.}{norm}
\dkDeclareSymbol[Sine-theta separation G_delta(B,C): finite interval/exterior separation in either spectral order, or ordered separation B <= a I and C >= (a+delta) I, or the reverse. The ordered alternatives allow both spectra to be unbounded.]
  {SineGap}{\mathcal G_\delta}{gap}{local}{Interval/exterior separation or either ordered separation, as in the sine-theta theorem.}{predicate}
\dkDeclareSymbol[Auxiliary unitarily invariant norm M witnessing the Q-norm property N(T) squared = M(T* T). For the Frobenius norm N, one can take M to be the trace norm.]
  {NormM}{M}{normaux}{local}{Auxiliary unitarily invariant norm used to define a Q-norm through the squared modulus.}{norm}
\dkDeclareSymbol[Q-norm property: a unitarily invariant norm N is a Q-norm when N(T) squared = M(T* T) for some unitarily invariant norm M. The Frobenius norm is an example, with M the trace norm; the repaired full-displacement theorem uses this norm class.]
  {NormQ}{Q}{normaux}{glossary}{Q-norm: its square is a unitarily invariant norm of the squared modulus.}{norm}
\dkDeclareSymbol[Singleton trace cost s_i = min(4 sin(theta_i/2), 2) for angle label i. It transports one real principal plane independently, choosing a rotation of cost 4 sin(theta_i/2) or a reflection of cost 2.]
  {SingletonCost}{s}{norm}{glossary}{Minimum trace displacement for one real principal plane: choose a rotation or a reflection.}{scalar-cost}
\dkDeclareSymbol[Paired trace cost e_ij = 4 sin((theta_i + theta_j)/2) for distinct angle labels i,j. It is attained by coupling their principal planes in a real four-dimensional block. The edge weight s_i + s_j - e_ij is the difference between the cost of two singletons and this paired cost.]
  {PairCost}{e}{norm}{glossary}{Trace displacement cost of coupling two principal planes in a four-dimensional block.}{scalar-cost}
\dkDeclareSymbol[Scalar phase cost f_p(t) = (2 sin(t/2)) to the power p, for 0 <= t <= pi. A real planar rotation through t contributes two copies, so its displacement power cost is 2 f_p(t).]
  {MatchingPhaseCost}{f_p}{norm}{local}{Displacement power cost for one scalar phase.}{scalar-cost}
\dkDeclareSymbol[Singleton power cost a_i(p) = min(2 f_p(theta_i), 2 to the power p). These are the rotation and reflection choices for one principal plane. At p=1 it equals the trace-cost weight s_i.]
  {SingletonPowerCost}{a}{norm}{local}{Minimum displacement power cost for one angle.}{scalar-cost}
\dkDeclareSymbol[Paired power cost b_ij(p) = min over d <= tau <= h of 2(f_p(h+tau) + f_p(h-tau)), where h=(theta_i+theta_j)/2 and d=abs(theta_i-theta_j)/2. This is the positive-determinant pair minimum; unrestricted block optimization also compares two singletons. At p=1 it equals e_ij.]
  {PairPowerCost}{b}{norm}{local}{Positive-determinant displacement power cost for a pair of angles.}{scalar-cost}
\dkDeclareSymbol[Matching M on the finite general graph of positive principal-angle labels. Each label belongs to at most one edge; unmatched labels use singleton transporters. Minimizing paired and singleton costs is equivalent to maximum-weight matching with edge weights g_ij = s_i + s_j - e_ij. The empty matching is allowed.]
  {Matching}{M}{normaux}{glossary}{Matching of positive principal-angle labels: each label is used in at most one pair.}{matching}
\dkDeclareSymbol[Displacement power sum Phi_p(W) = sum_j s_j(I - W)^p. The matching proof first minimizes it for 1 < p < 2, then lets p decrease to 1, where it becomes the trace norm ||I - W||_1.]
  {PowerCost}{\Phi}{norm}{local}{Sum of the p-th powers of displacement singular values.}{scalar-cost}
% Auxiliary matching-proof notation: neutral ink, with the same hover registry.
\dkDeclareSymbol[Off-diagonal gradient block G_01 mapping the source complement K perpendicular into K. Stationarity makes the diagonal blocks of G zero, so G = ((0,G_01),(-G_01 transpose,0)) in the source/complement decomposition.]
  {GradientOffDiagonal}{G_{01}}{fg}{local}{Complement-to-source block of the skew-symmetric gradient.}{linear-operator}
\dkDeclareSymbol[Orthogonal group O(n): real n by n matrices W satisfying W transpose times W = I. Its elements preserve Euclidean inner products and lengths. The notation O(2) refers to the same group in dimension two.]
  {OrthogonalGroup}{O(n)}{fg}{local}{Group of real orthogonal matrices.}{group}
\dkDeclareSymbol[Feasible transporters F = {W in O(n) : W(K) = L}. Equivalently W P_K W transpose = P_L. This nonempty compact set is the optimization domain throughout the matching proof.]
  {FeasibleTransport}{\mathcal F}{fg}{local}{Orthogonal maps carrying the source subspace onto the target.}{set}
\dkDeclareSymbol[Singular value s_j(T), the nonnegative square root of an eigenvalue of T transpose times T, counted with multiplicity. Summing s_j(I-W) gives the trace displacement; summing its p-th powers gives Phi_p(W). These singular values differ from the singleton weights s_i.]
  {TransportSingularValues}{\symglyph{SingularValueBase}_j}{fg}{local}{A displacement singular value in the matching proof.}{scalar}
\dkDeclareSymbol[Scalar singular-value magnitude s between zero and two for the displacement I-W of an orthogonal matrix. Uniform convergence of s to the power p to s on this interval permits the trace-norm limit. This variable differs from the rotation parameter s used earlier in the proof.]
  {DisplacementSize}{s}{fg}{local}{Scalar displacement singular-value magnitude in the uniform limit.}{scalar}
\dkDeclareSymbol[All-singleton cost C_0 = sum_i s_i, before any two angle labels are paired. A matching has cost C_0 minus the sum of its selected edge weights.]
  {SingletonBaseline}{C_0}{fg}{local}{Trace displacement when every angle uses its singleton transporter.}{scalar-cost}
\dkDeclareSymbol[Matching power cost C_p(M): sum of b_ij(p) over selected pairs plus a_i(p) over unmatched labels. Its minimum equals the minimum of Phi_p over feasible transporters. At p=1 it is C(M).]
  {MatchingPowerCost}{C_p}{fg}{local}{Total paired and singleton power cost of a matching.}{scalar-cost}
\dkDeclareSymbol[Invariant orthogonal block V_a preserved by the source projection, target projection, and minimizing transporter. The displacement cost adds over these blocks, and their principal-angle lists combine into the original list.]
  {MatchingBlock}{V_a}{fg}{local}{Simultaneous invariant block in the matching proof.}{linear-subspace}
\dkDeclareSymbol[Gradient-level space V_g: the eigenspace of -G squared for a positive eigenvalue g squared. It is preserved by G, W, and the source reflection J_src. Multiplication by i on it is G/g.]
  {GradientLevelSpace}{V_g}{fg}{local}{Positive gradient-level eigenspace.}{linear-subspace}
\dkDeclareSymbol[Eigenphase varphi of W on V_g viewed as a complex space with multiplication by i = G/g. It lies between zero and pi and satisfies f_p prime(varphi) = g. There are at most two such phases.]
  {GradientEigenphase}{\varphi}{fg}{local}{One of the at most two phases at a positive gradient level.}{angle}
\dkDeclareSymbol[Orthogonal projection E onto the first phase space of W on V_g. Here W = exp(i varphi_1) E + exp(i varphi_2) (I-E). In coordinates where the source reflection is conjugation, E = X + iY.]
  {GradientPhaseProjection}{E}{fg}{local}{Projection onto the first phase space in a gradient-level space.}{orthogonal-projection}
\dkDeclareSymbol[Real part X of the phase projection E = X + iY. It is a real symmetric matrix with spectrum in the closed interval from zero to one, satisfying X squared - Y squared = X and XY + YX = Y. This X differs from the source-coordinate and Sylvester matrices elsewhere in the paper.]
  {ProjectionRealPart}{X}{fg}{local}{Real symmetric part of the phase projection.}{linear-operator}
\dkDeclareSymbol[Imaginary part Y of the phase projection E = X + iY. It is a real skew-symmetric matrix. It maps the lambda eigenspace of X to the 1-lambda eigenspace and constructs the small invariant planes.]
  {ProjectionImaginaryPart}{Y}{fg}{local}{Real skew-symmetric imaginary part of the phase projection.}{linear-operator}
\dkDeclareSymbol[Eigenvalue lambda in the closed interval from zero to one of the real part X of the phase projection. It is local to the block construction and differs from the original operator's spectral parameter.]
  {ProjectionEigenvalue}{\lambda}{fg}{local}{Eigenvalue of the real part of the phase projection.}{scalar}
\dkDeclareSymbol[Plane coupling eta = sqrt(lambda (1-lambda)) for 0 < lambda < 1. If u is a unit lambda eigenvector of X, then v = Yu/eta completes its invariant real plane.]
  {ProjectionCoupling}{\eta}{fg}{local}{Normalization used to construct an invariant real plane.}{scalar}
\dkDeclareSymbol[Pair mean angle h = (theta_i + theta_j)/2. Together with d = (theta_i - theta_j)/2 after ordering theta_i >= theta_j, it parametrizes the two-angle block cost.]
  {PairMeanAngle}{h}{fg}{local}{Mean of the two ordered principal angles.}{angle}
\dkDeclareSymbol[Pair half-difference d = (theta_i - theta_j)/2 after ordering theta_i >= theta_j. It satisfies 0 <= d <= h and h+d <= pi/2.]
  {PairHalfDifference}{d}{fg}{local}{Half the difference of the two ordered principal angles.}{angle}
\dkDeclareSymbol[Scalar pair parameter tau in the closed interval from d to h. The paired power cost is the minimum of 2(f_p(h+tau) + f_p(h-tau)) over this interval. At p=1 a minimum is attained at tau=h.]
  {PairPhaseParameter}{\tau}{fg}{local}{Scalar parameter in the two-angle power-cost minimum.}{scalar}
\dkDeclareSymbol[Attaining rotation angle gamma = arccos(cos(tau)/cos(d)). Choosing U_0 = U_1 = R_gamma realizes phase parameters x=h and y=tau in the four-dimensional block.]
  {PairAttainingAngle}{\gamma}{fg}{local}{Rotation angle realizing a chosen scalar pair parameter.}{angle}
\dkDeclareSymbol[Four-dimensional direct rotation D = ((C,-S),(S,C)) in source and complementary principal coordinates, with C the diagonal principal cosines and S the diagonal principal sines. Every transporter on this block is D diag(U_0,U_1).]
  {PairDirectRotation}{D}{fg}{local}{Direct rotation on the two principal planes.}{unitary-operator}
\dkDeclareSymbol[Diagonal cosine block C = diag(cos(theta_i), cos(theta_j)) in the four-dimensional direct rotation D. This matrix differs from the scalar displacement cost C(M).]
  {PairCosineBlock}{C}{fg}{local}{Diagonal principal-cosine block.}{matrix}
\dkDeclareSymbol[Diagonal sine block S = diag(sin(theta_i), sin(theta_j)) in the four-dimensional direct rotation D. This matrix is local to the matching proof.]
  {PairSineBlock}{S}{fg}{local}{Diagonal principal-sine block.}{matrix}
\dkDeclareSymbol[Free planar orthogonal matrices U_0 and U_1 acting on the two-dimensional source and its two-dimensional complement, respectively. Multiplying D by diag(U_0,U_1) parametrizes every transporter on this four-dimensional block.]
  {PairFreeOrthogonal}{U}{fg}{local}{Free source or complementary planar orthogonal factor.}{unitary-operator}
\dkDeclareSymbol[Phase parameter x in the closed interval from h to pi-h for a positive-determinant four-dimensional transporter. Its eigenvalues are exp(plus or minus i(x+y)) and exp(plus or minus i(x-y)).]
  {PairPhaseX}{x}{fg}{local}{First phase parameter of a four-dimensional transporter.}{angle}
\dkDeclareSymbol[Phase parameter y in the closed interval from d to pi-d for a positive-determinant four-dimensional transporter. Together with x it gives rotational phases x+y and x-y, folded into the closed interval from zero to pi for the displacement cost.]
  {PairPhaseY}{y}{fg}{local}{Second phase parameter of a four-dimensional transporter.}{angle}
\dkDeclareSymbol[Scalar a = cos(x) in the four-dimensional phase calculation. Together with b = cos(y), it expresses the trace as 4ab and determines the reciprocal characteristic polynomial.]
  {PairPhaseCosineX}{a}{fg}{local}{Cosine of the first phase parameter in the trace calculation.}{scalar}
\dkDeclareSymbol[Scalar b = cos(y) in the four-dimensional phase calculation, paired with a = cos(x). These scalars are local to the characteristic-polynomial calculation.]
  {PairPhaseCosineY}{b}{fg}{local}{Cosine of the second phase parameter in the trace calculation.}{scalar}
\dkDeclareSymbol[Rotation parameter r: the two free orthogonal factors use planar rotation angles r+s and r-s. It appears in cos(y) = cos(d) cos(r) for two rotations.]
  {PairRotationR}{r}{fg}{local}{Half the sum of the two free planar rotation angles.}{angle}
\dkDeclareSymbol[Rotation parameter s: the two free orthogonal factors use planar rotation angles r+s and r-s. It appears in cos(x) = cos(h) cos(s) for two rotations. This scalar differs from the singular values and singleton costs.]
  {PairRotationS}{s}{fg}{local}{Half the difference of the two free planar rotation angles.}{angle}
\dkDeclareSymbol[Planar reflection J_2 = diag(1,-1). Every two-dimensional orthogonal matrix of determinant -1 can be written as a planar rotation times J_2.]
  {PairPlanarReflection}{J_2}{fg}{local}{Fixed planar reflection used to parametrize determinant -1 factors.}{unitary-operator}
\dkDeclareSymbol[Real planar rotation R_t = ((cos(t),-sin(t)),(sin(t),cos(t))). Its two displacement singular values are both 2 times abs(sin(t/2)). In the matching proof a rotation through t in the closed interval from zero to pi contributes 2 f_p(t).]
  {MatchingPlanarRotation}{\symglyph{PlanarRotationBase}_t}{fg}{local}{Planar rotation through the indicated angle.}{unitary-operator}
% Statistical application / operator provenance.
\dkDeclareSymbol{Sigma}{\Sigma}{old}{local}{Population covariance matrix.}{matrix}
\dkDeclareSymbol{SigmaHat}{\widehat\Sigma}{current}{local}{Estimated or perturbed covariance matrix.}{matrix}
\dkDeclareSymbol{YWSV}{V}{wanted}{local}{Orthonormal coordinates for the population eigenspace in the statistical application.}{isometric-embedding}
\dkDeclareSymbol{YWSVhat}{\widehat V}{trialframe}{local}{Estimated eigenspace frame in the statistical application.}{isometric-embedding}

% Figure-specific objects retain their own meanings when bare letters overlap.
\dkDeclareSymbol[Signed rotation angle phi of the displayed map on its invariant real two-plane. At theta = pi/4 the direct rotation turns its two principal planes by opposite signed angles of magnitude pi/4; the competitor turns its moving plane by pi/2 and fixes its other plane.]
  {PlaneRotationAngle}{\varphi}{sine}{local}{Signed rotation angle of the orthogonal map on the displayed invariant plane.}{angle}
\dkDeclareSymbol[Unit vector q in the displayed invariant plane. The map T is the direct rotation D or the competitor W in that panel; Tq is the image and (I - T)q = q - Tq is the displacement.]
  {FigureInput}{q}{ambient}{local}{Unit input to the orthogonal map shown in this invariant-plane panel.}{vector}
\dkDeclareSymbol[Image Tq of the panel unit input q under T = D or W. Both q and Tq lie on the unit circle; the green displacement arrow joins Tq to q, so it represents q - Tq.]
  {FigureImage}{Tq}{wantedvector}{local}{Image of the unit input under the orthogonal map shown in this panel.}{vector}
\dkDeclareSymbol[Displacement (I - T)q = q - Tq, drawn from the image Tq to the unit input q. On a plane rotated through phi its length is 2 |sin(phi/2)|; it vanishes on a fixed plane.]
  {DisplacementVector}{(I-T)q}{displacement}{local}{Difference between a unit input and its image under the displayed orthogonal map.}{vector}
\dkDeclareSymbol[First principal plane Pi_1 = span(e_0, e_3) for the source-target pair in the real four-dimensional counterexample. At theta = pi/4 the direct rotation acts on this plane by a quarter of pi, with displacement singular value sqrt(2 - sqrt(2)).]
  {CounterPrincipalOne}{\Pi_1}{fg}{glossary}{First principal plane in the real counterexample: the span of e0 and e3.}{linear-subspace}
\dkDeclareSymbol[Second principal plane Pi_2 = span(e_1, e_2) for the source-target pair in the real four-dimensional counterexample. The direct rotation acts by the same angle magnitude as on Pi_1, with the opposite sign in the displayed plane coordinates.]
  {CounterPrincipalTwo}{\Pi_2}{fg}{glossary}{Second principal plane in the real counterexample: the span of e1 and e2.}{linear-subspace}
\dkDeclareSymbol[Moving plane M = range(W - I) = span(e_0 + e_2, e_1 + e_3) for the competitor. W rotates it by 2 theta; at theta = pi/4 the two displacement singular values contributed by this plane are both sqrt(2).]
  {CounterMovingPlane}{M}{ambient}{glossary}{Moving plane of the competing transporter: the range of W minus the identity.}{linear-subspace}
\dkDeclareSymbol[Fixed plane Z = kernel(W - I) = span(e_0 - e_2, e_1 - e_3) for the competitor. W acts as the identity here, so this plane contributes two zero displacement singular values.]
  {CounterFixedPlane}{Z}{ambient}{glossary}{Fixed plane of the competing transporter: the kernel of W minus the identity.}{linear-subspace}
\dkDeclareSymbol[Direct-rotation displacement singular value a = 2 sin(pi/8) = sqrt(2 - sqrt(2)) at theta = pi/4. Each of its two principal planes contributes two copies, giving trace displacement 4a.]
  {CounterChordA}{a}{displacement}{glossary}{Direct-rotation displacement singular value in the counterexample: sqrt(2 - sqrt(2)).}{scalar}
\dkDeclareSymbol[Competitor displacement singular value b = 2 sin(pi/4) = sqrt(2) at theta = pi/4. Its moving plane contributes two copies; the fixed plane contributes two zeros, giving trace displacement 2b.]
  {CounterChordB}{b}{displacement}{glossary}{Competing-transporter displacement singular value in the counterexample: square root of 2.}{scalar}
\dkDeclareSymbol[Singular value s_j(I - T) of the full ambient displacement, ordered from largest to smallest. Each rotating real two-plane contributes two values 2 |sin(phi/2)|; fixed directions contribute zeros. Their sum is the trace cost.]
  {DisplacementSingular}{\symglyph{SingularValueBase}_j}{displacement}{local}{An individual singular value of the full displacement from the identity.}{scalar}
\dkDeclareSymbol[Edge weight g_ij = s_i + s_j - e_ij, the difference between the cost of two singletons and the paired cost. A matching has total cost C_0 minus the sum of its edge weights, with C_0 = sum_i s_i. Maximizing this sum gives the minimum transport cost; edges of negative weight may be omitted.]
  {MatchingEdgeWeight}{g_{ij}}{current}{glossary}{Edge weight: the cost of two singletons minus their paired cost.}{scalar-weight}
\dkDeclareSymbol[Total trace-norm displacement cost. For a transporter W, C(W) = ||I - W||_1, the sum of its displacement singular values. For a matching M, C(M) is the sum of its paired and unmatched singleton costs.]
  {TraceCost}{C}{norm}{glossary}{Trace-norm displacement cost of the indicated orthogonal transporter or matching.}{scalar-cost}

% Spectral values have their own key, distinct from an operator block.
\dkDeclareSymbol[Spectral parameter lambda of a self-adjoint operator. In a finite-dimensional eigenvalue list lambda_j(A) it denotes an individual eigenvalue; in A = integral lambda dE(lambda), it is the real integration variable.]{SpectralValue}{\lambda}{wantedspectrum}{local}{Real spectral parameter or individual eigenvalue of a self-adjoint operator.}{scalar}



% Conventional mathematical notation. Hover text applies even when ink is plain.
% These are local (non-glossary) unless promoted explicitly by the author.
% IdentityTwo and IdentityFour derive their I from Identity. Singular-value
% registrations derive s_j from SingularValueBase; SingletonCost has its own s.
\dkDeclareSymbol[Identity operator on the vector or Hilbert space in context. Its dimension is inferred from the expression and need not be finite.]{Identity}{I}{plain}{local}{Identity operator on the space in context.}{identity-operator}
\dkDeclareSymbol[Two-by-two identity matrix, acting on a two-dimensional coordinate space.]{IdentityTwo}{\symglyph{Identity}_2}{plain}{local}{Two-by-two identity matrix.}{identity-operator}
\dkDeclareSymbol[Four-by-four identity matrix, acting on a four-dimensional coordinate space.]{IdentityFour}{\symglyph{Identity}_4}{plain}{local}{Four-by-four identity matrix.}{identity-operator}
\dkDeclareSymbol[Two-by-two zero matrix. Unlike the bold zero vector, this is a linear operator/matrix with two rows and two columns.]{ZeroMatrixTwo}{0_2}{plain}{local}{Two-by-two zero matrix.}{matrix}
\dkDeclareSymbol[The constant pi, the ratio of the circumference of a circle to its diameter. Angles in this article use radians.]{PiConstant}{\pi}{plain}{local}{The constant pi (angles are in radians).}{constant}
\dkDeclareSymbol[The field of real numbers.]{RealField}{\mathbb R}{plain}{local}{The real numbers.}{field}
\dkDeclareSymbol[The field of complex numbers.]{ComplexField}{\mathbb C}{plain}{local}{The complex numbers.}{field}
\dkDeclareSymbol[Generic Hilbert space used in the two-projection discussion; distinct from the paper's ambient space symbol E in the sine-theta theorem.]{GeneralHilbertSpace}{\mathcal H}{plain}{local}{Generic Hilbert space in a local discussion.}{hilbert-space}
\dkDeclareSymbol[The spectrum sigma(T) of an operator T: the set of complex scalars for which T minus the scalar identity is not boundedly invertible. For self-adjoint T the spectrum is real.]{SpectrumSymbol}{\sigma}{plain}{local}{Spectrum of an operator.}{operator-notation}
\dkDeclareSymbol[Base letter for singular values, presently s. It is distinct from the base s for matching singleton costs.]{SingularValueBase}{s}{plain}{local}{Base letter for singular values.}{notation-base}
\dkDeclareSymbol[Ordered singular value s_j(T), equal to the j-th eigenvalue of |T| for a finite matrix; distinguish this from the matching singleton cost s_j.]{SingularValueS}{\symglyph{SingularValueBase}_j}{plain}{glossary}{The j-th singular value, under the currently printed s_j convention.}{scalar}
\dkDeclareSymbol[Alternate indexed singular-value notation sigma_j(T). The current manuscript also uses s_j(T); this registered alias is available for a later deliberate notation decision.]{SingularValueSigma}{\sigma_j}{plain}{local}{Indexed singular value under the sigma_j convention.}{scalar}
\dkDeclareSymbol[Operator norm ||T||_op, the supremum of ||Tx|| over unit vectors x.]{OperatorNormNotation}{\|\cdot\|_{\mathrm{op}}}{plain}{glossary}{Operator norm.}{norm}
\dkDeclareSymbol[Frobenius norm ||T||_F, the square root of the sum of squared matrix entries, equal to the finite-dimensional Hilbert--Schmidt norm.]{FrobeniusNormNotation}{\|\cdot\|_F}{plain}{glossary}{Frobenius norm.}{norm}
\dkDeclareSymbol[Trace norm ||T||_1, the sum of singular values when finite.]{TraceNormNotation}{\|\cdot\|_1}{plain}{glossary}{Trace (nuclear) norm.}{norm}
\dkDeclareSymbol[Hilbert--Schmidt norm ||T||_HS, the square root of the sum of squared singular values when finite.]{HilbertSchmidtNormNotation}{\|\cdot\|_{\mathrm{HS}}}{plain}{local}{Hilbert--Schmidt norm.}{norm}
\dkDeclareSymbol[Ky Fan j-norm, the sum of the largest j singular values.]{KyFanNormNotation}{\|\cdot\|_{(j)}}{plain}{glossary}{Ky Fan j-norm.}{norm}
\dkDeclareSymbol[Interval/exterior spectral separation predicate I_delta(B,C), specifying a central closed interval separated by delta from the other operator's spectrum.]{IntervalSeparation}{\mathcal I}{plain}{local}{Interval/exterior separation predicate.}{predicate}
\dkDeclareSymbol[Ordered spectral separation predicate O_delta(B,C), requiring one self-adjoint operator to lie below a spectral threshold and the other above it by at least delta.]{OrderedSeparation}{\mathcal O}{plain}{local}{Ordered spectral separation predicate.}{predicate}
\dkDeclareSymbol[Positive parameter epsilon controlling a small perturbation in the two-dimensional eigenvalue example.]{EigenvalueExampleEpsilon}{\varepsilon}{plain}{local}{Small positive parameter in the eigenvalue-example matrices.}{scalar}
\dkDeclareSymbol[Positive lower bound epsilon for the trial-coordinate map E_0, as in E_0* E_0 >= epsilon squared times the identity.]{FrameLowerBoundEpsilon}{\varepsilon}{plain}{local}{Lower-bound parameter for the trial-coordinate embedding.}{scalar}
\dkDeclareSymbol[Rotation angle vartheta of a planar orthogonal block in the matching-proof first-variation calculation.]{MatchingPlanarAngle}{\vartheta}{plain}{local}{Planar rotation angle in the matching-proof calculation.}{angle}


% Ordinary named mathematical operators already used throughout the manuscript.
% Keeping these here makes their spelling a single editorial choice.
\dkDeclareSymbol[Range of a linear operator: the set of its output vectors.]{RangeOperator}{\operatorname{ran}}{plain}{local}{Range (image) of an operator.}{operator-notation}
\dkDeclareSymbol[Kernel (null space) of a linear operator: inputs mapped to zero.]{KernelOperator}{\operatorname{ker}}{plain}{local}{Kernel (null space) of an operator.}{operator-notation}
\dkDeclareSymbol[Domain of an operator: the vectors on which the operator is defined, important for unbounded operators.]{DomainOperator}{\operatorname{dom}}{plain}{local}{Domain of an operator.}{operator-notation}
\DeclareRobustCommand{\DKRange}{\sym{RangeOperator}}
\DeclareRobustCommand{\DKKernel}{\sym{KernelOperator}}
\DeclareRobustCommand{\DKDom}{\sym{DomainOperator}}

% Norm-formatting macros: the operand changes by occurrence but the notation
% lives here. The outer hover describes the norm and does not add ink color.
\DeclareRobustCommand{\dkOperatorNorm}[1]{\like{OperatorNormNotation}{\|#1\|_{\mathrm{op}}}}
\DeclareRobustCommand{\dkTraceNorm}[1]{\like{TraceNormNotation}{\|#1\|_1}}
\DeclareRobustCommand{\dkFrobeniusNorm}[1]{\like{FrobeniusNormNotation}{\|#1\|_F}}
\DeclareRobustCommand{\dkHilbertSchmidtNorm}[1]{\like{HilbertSchmidtNormNotation}{\|#1\|_{\mathrm{HS}}}}
\DeclareRobustCommand{\dkKyFanNorm}[2]{\like{KyFanNormNotation}{\|#1\|_{(#2)}}}
\DeclareRobustCommand{\dkSchattenNorm}[2]{\like{SchattenNormNotation}{\|#1\|_{S_{#2}}}}

\dkDeclareSymbol[Schatten p-norm ||T||_{S_p}, obtained from the l^p-norm of the singular values and finite for operators in the Schatten p-class.]{SchattenNormNotation}{\|\cdot\|_{S_p}}{plain}{local}{Schatten p-norm.}{norm}


\dkDeclareSymbol[Positive modulus of an operator T: |T| = (T* T)^(1/2).]{OperatorModulusNotation}{|\cdot|}{plain}{glossary}{Positive modulus of an operator.}{operator-notation}
\dkDeclareSymbol[Linear span of the indicated vectors, i.e. the smallest linear subspace that contains them.]{SpanOperator}{\operatorname{span}}{plain}{local}{Linear span of vectors.}{operator-notation}
\dkDeclareSymbol[Diagonal or block-diagonal matrix/operator with the displayed entries or blocks.]{DiagonalOperator}{\operatorname{diag}}{plain}{local}{Diagonal or block-diagonal operator.}{linear-operator}
\dkDeclareSymbol[Trace of a trace-class operator, or ordinary matrix trace in finite dimensions.]{TraceOperator}{\operatorname{tr}}{plain}{local}{Trace of an operator.}{operator-notation}
\dkDeclareSymbol[Dimension of a vector space or subspace.]{DimensionOperator}{\dim}{plain}{local}{Dimension of a vector space.}{operator-notation}
\dkDeclareSymbol[Rank of a linear map or matrix: dimension of its range.]{RankOperator}{\operatorname{rank}}{plain}{local}{Rank of a linear map or matrix.}{operator-notation}
\dkDeclareSymbol[Distance between sets of spectral values: the infimum of pairwise distances.]{SpectralDistanceOperator}{\operatorname{dist}}{plain}{local}{Distance between two spectral sets.}{operator-notation}
\DeclareRobustCommand{\dkModulus}[1]{\like{OperatorModulusNotation}{|#1|}}

% Context-specific glyphs previously hidden in variants of other semantic keys.
% Keeping both the roles and the exact printed forms makes collisions explicit.
\dkDeclareSymbol[Nonzero vector x tested in the Rayleigh quotient, before imposing unit length.]{RayleighTestVector}{x}{trialvector}{local}{Test vector in the Rayleigh quotient.}{vector}
\dkDeclareSymbol[Positive cosine block C in the generic two-projection decomposition of Halmos; C and S commute and C squared plus S squared is the identity.]{HalmosCosineBlock}{C}{anglecos}{local}{Cosine block in the Halmos two-projection representation.}{linear-operator}
\dkDeclareSymbol[Positive sine block S in the generic two-projection decomposition of Halmos; C and S commute and C squared plus S squared is the identity.]{HalmosSineBlock}{S}{sineoperator}{local}{Sine block in the Halmos two-projection representation.}{linear-operator}
\dkDeclareSymbol[Riesz spectral projection P_Gamma defined by a contour Gamma that encloses selected spectral values.]{RieszProjection}{P_{\symglyph{RieszContour}}}{wantedprojection}{local}{Spectral projection associated with a contour Gamma.}{orthogonal-projection}
\dkDeclareSymbol[Historical Sylvester right-hand coefficient matrix A, in B X - X A = Y. Distinct from the current paper's reference operator A.]{HistoricalSylvesterRight}{A}{trialspectrum}{local}{Right-hand coefficient A in the historical Sylvester equation.}{self-adjoint-operator}
\dkDeclareSymbol[Scalar coefficient b in the Sylvester prototype (b-a)x=y, with separated scalar spectral values a and b.]{SylvesterScalarLeft}{b}{unwantedspectrum}{local}{Scalar left coefficient b in a Sylvester equation.}{scalar}
\dkDeclareSymbol[Scalar coefficient a in the Sylvester prototype (b-a)x=y.]{SylvesterScalarRight}{a}{trialspectrum}{local}{Scalar right coefficient a in a Sylvester equation.}{scalar}
\dkDeclareSymbol[Scalar unknown x in the Sylvester prototype (b-a)x=y.]{SylvesterScalarUnknown}{x}{sylvester}{local}{Scalar unknown x in a Sylvester equation.}{scalar}
\dkDeclareSymbol[Scalar right side y in the Sylvester prototype (b-a)x=y.]{SylvesterScalarRHS}{y}{resid}{local}{Scalar right side y in a Sylvester equation.}{scalar}
\dkDeclareSymbol[Source subspace K of a real Euclidean space in the matching proof; an orthogonal transporter sends K onto L.]{MatchingSourceSubspace}{K}{trial}{local}{Source subspace K in the matching proof.}{linear-subspace}
\dkDeclareSymbol[Target subspace L of a real Euclidean space in the matching proof; an orthogonal transporter sends K onto L.]{MatchingTargetSubspace}{L}{wanted}{local}{Target subspace L in the matching proof.}{linear-subspace}
\dkDeclareSymbol[Orthogonal projection P_L onto the matching proof's target subspace L.]{MatchingTargetProjection}{P_L}{wantedprojection}{local}{Orthogonal projection onto the matching target L.}{orthogonal-projection}
\dkDeclareSymbol[Real two-dimensional orthogonal group O(2): planar rotations and reflections.]{PlanarOrthogonalGroup}{O(2)}{fg}{local}{Orthogonal matrices of size two.}{group}
\dkDeclareSymbol[Base letter R used to denote planar rotation operators R_t in the matching proof.]{PlanarRotationBase}{R}{fg}{local}{Base letter for planar rotations.}{notation-base}
\dkDeclareSymbol[Planar rotation R_(r+s) through the sum of scalar phases r and s.]{MatchingPlanarRotationSum}{\symglyph{PlanarRotationBase}_{r+s}}{fg}{local}{Planar rotation through phase r+s.}{unitary-operator}
\dkDeclareSymbol[Planar rotation R_gamma through phase gamma in the block rotation calculation.]{MatchingPlanarRotationGamma}{\symglyph{PlanarRotationBase}_{\symglyph{MatchingPhaseGamma}}}{fg}{local}{Planar rotation through phase gamma.}{unitary-operator}

% Specific symbols in the historical and appendix examples. They retain their
% existing black-ink notation but receive locally accurate hover definitions.
\dkDeclareSymbol[Contour Gamma, a positively oriented closed curve in the resolvent set enclosing a selected isolated spectral cluster.]{RieszContour}{\Gamma}{plain}{local}{Contour used in the Riesz spectral-projection formula.}{contour}
\dkDeclareSymbol[Even integer nu in a tested Ky Fan comparison, with nu/2 principal-plane terms on the right side.]{KyFanEvenIndex}{\nu}{plain}{local}{Even integer in the principal-plane Ky Fan comparison.}{scalar}
\dkDeclareSymbol[Orthogonal projection Omega_k onto the k-th principal plane in the published displacement comparison.]{PrincipalPlaneProjectionBase}{\Omega}{plain}{local}{Projection onto the indicated principal plane.}{notation-base}
\dkDeclareSymbol[Scalar gap threshold gamma in the ordered Sylvester spectral separation B >= (gamma+delta)I and C <= gamma I.]{SylvesterGapThreshold}{\gamma}{plain}{local}{Spectral separation threshold in a Sylvester estimate.}{scalar}
\dkDeclareSymbol[Projection Omega_t from a net or sequence that converges strongly to the identity in the singular-value approximation lemma.]{StrongApproxProjection}{\Omega_t}{plain}{local}{Orthogonal projection converging strongly to the identity.}{orthogonal-projection}
\dkDeclareSymbol[First orthogonal projection Omega used in the two-block pinching inequality for a unitarily invariant norm.]{PinchingProjectionOne}{\Omega}{plain}{local}{First projection in the pinching inequality.}{orthogonal-projection}
\dkDeclareSymbol[Second orthogonal projection Upsilon used in the two-block pinching inequality for a unitarily invariant norm.]{PinchingProjectionTwo}{\Upsilon}{plain}{local}{Second projection in the pinching inequality.}{orthogonal-projection}
\dkDeclareSymbol[Finite-rank projection Gamma used to retain the image of a chosen finite-rank compression K Psi.]{CompressionProjectionImage}{\Gamma}{plain}{local}{Range projection in the finite-rank compression lemma.}{orthogonal-projection}
\dkDeclareSymbol[Finite-rank projection Psi used to select input vectors in a nearly maximizing singular-value compression.]{CompressionProjectionInput}{\Psi}{plain}{local}{Domain projection in the finite-rank compression lemma.}{orthogonal-projection}
\dkDeclareSymbol[Small positive tolerance eta in the finite-rank compression approximation estimate.]{CompressionTolerance}{\eta}{plain}{local}{Tolerance controlling the off-diagonal compression norm.}{scalar}

\dkDeclareSymbol[Rotation phase gamma in a planar block from the matching proof.]{MatchingPhaseGamma}{\gamma}{plain}{local}{Rotation phase gamma in the matching calculation.}{scalar}

% Two common indices with local mathematical meanings distinct from other deltas.
\dkDeclareSymbol[Kronecker delta: delta_ij is one when the indices i and j agree and zero otherwise. This is not the spectral-gap parameter delta.]{KroneckerDelta}{\delta}{plain}{local}{Kronecker delta: equals 1 for equal indices and 0 otherwise.}{scalar}
\dkDeclareSymbol[Sign variable epsilon taking the values -1 and +1 in the factorization over two phase choices.]{PhaseSignChoice}{\epsilon}{plain}{local}{Sign variable taking values minus one and plus one.}{scalar}

% Documented notation correspondence. The source column is a faithful
% transcription of each cited original context, NOT an asserted universal DK
% meaning. The original A_0 is deliberately listed for two distinct roles:
% its reference block and the comparison/trial operator. The source uses the
% same printed symbol for both; this paper is free to distinguish them.
\dkDeclareNotationSource{DK1970}{Davis and Kahan (1970), The Rotation of Eigenvectors by a Perturbation. III}{https://doi.org/10.1137/0707001}
\dkDeclareNotationSource{Seelmann2019}{Seelmann (2019), Semidefinite Perturbations in the Subspace Perturbation Problem}{https://arxiv.org/abs/1708.02463}
\dkDeclareNotationSource{Blechta2020}{Blechta (2020), Stability of Linear GMRES Convergence with Respect to Compact Perturbations}{https://arxiv.org/abs/2005.12960}
% Arguments: {key}{source ID}{source glyph}{section or equation}{role-specific qualification}
\dkDeclareNotationAlternative{A}{DK1970}{A}{\S1, (1.2)}{Reference self-adjoint operator.}
\dkDeclareNotationAlternative{H}{DK1970}{H}{\S1, (1.2)}{Additive perturbation.}
\dkDeclareNotationAlternative{Acurrent}{DK1970}{A+H}{\S1}{Perturbed operator.}
\dkDeclareNotationAlternative{E0}{DK1970}{E_0}{\S1, (1.1)}{Source isometric frame.}
\dkDeclareNotationAlternative{E1}{DK1970}{E_1}{\S1, (1.1)}{Reference complement frame.}
\dkDeclareNotationAlternative{F0}{DK1970}{F_0}{\S1}{Exact-target frame.}
\dkDeclareNotationAlternative{F1}{DK1970}{F_1}{\S1}{Exact-target complement frame.}
\dkDeclareNotationAlternative{ReferenceBlock0}{DK1970}{A_0}{\S1, (1.2)}{Reference source block.}
\dkDeclareNotationAlternative{ReferenceBlock1}{DK1970}{A_1}{\S1, (1.2)}{Reference complement block.}
\dkDeclareNotationAlternative{A0}{DK1970}{A_0}{\S1, residual discussion}{Residual comparison operator (the reference block in DK).}
% We deliberately leave Lambda0 and Lambda1 without DK source-glyph metadata:
% the available source transcription does not reliably preserve the source
% accents/decorations on the perturbed spectral-block symbols.
\dkDeclareNotationAlternative{Psource}{DK1970}{P}{\S1}{Reference/source projection.}
\dkDeclareNotationAlternative{Qtarget}{DK1970}{Q}{\S1}{Target projection.}
\dkDeclareNotationAlternative{Pexact}{DK1970}{Q}{\S1}{Exact-target projection.}
\dkDeclareNotationAlternative{B}{DK1970}{B}{\S1, (1.2)}{Off-diagonal coupling block.}
\dkDeclareNotationAlternative{Theta0}{DK1970}{\Theta_0}{\S1}{Source-side angle operator.}
\dkDeclareNotationAlternative{Theta1}{DK1970}{\Theta_1}{\S1}{Complementary angle operator.}
\dkDeclareNotationAlternative{Theta}{DK1970}{\Theta}{\S1}{Full angle operator.}
\dkDeclareNotationAlternative{R}{DK1970}{R}{\S1, residual discussion}{Trial residual.}
\dkDeclareNotationAlternative{Rrotation}{DK1970}{U}{\S1, (1.18)}{Canonical direct rotation.}
\dkDeclareNotationAlternative{Wcompetitor}{DK1970}{V}{\S4}{Competing unitary transporter.}
\dkDeclareNotationAlternative{delta}{DK1970}{\delta}{\S2}{Spectral separation parameter.}
% Contemporary alternatives are recorded separately and do not enter the DK
% correspondence table. Their citation locations must be individually checked.
\dkDeclareNotationAlternative{Eambient}{Seelmann2019}{\mathcal H}{p.~1}{Ambient Hilbert space.}
\dkDeclareNotationAlternative{SingularValueS}{Blechta2020}{\sigma_j}{Thm.~1.1}{Ordered singular values of a compact operator.}
 in paper.tex)
% or pass \def\DKMonochromeMath{1} before % One switch: false = journal style; true = colored symbols and author--year citations.
\newif\ifDKAuthorStyle

%\DKAuthorStylefalse
\DKAuthorStyletrue

\ifdefined\DKUseAuthorStyle\DKAuthorStyletrue\fi
\ifdefined\DKJournalStyle\DKAuthorStylefalse\fi
\ifDKAuthorStyle
  \PassOptionsToClass{sn-mathphys-ay}{sn-jnl}
\else
  \PassOptionsToClass{sn-mathphys-num}{sn-jnl}
\fi
\documentclass[pdflatex]{sn-jnl} % PDF engine option; compiled with LuaLaTeX.
% PDF version used for the manuscript and its symbol annotations.
\pdfvariable minorversion=7\relax

\usepackage{amsmath,amssymb,amsthm,mathtools}
\usepackage{booktabs,tabularx,array,longtable}
\usepackage[section]{placeins}
\usepackage{float}
\usepackage{microtype}
\usepackage{graphicx}
\usepackage{pgf} % Live vector figures with LaTeX symbol labels and tooltips
\usepackage{xurl}
\ifDKAuthorStyle
  \setcitestyle{authoryear,round,aysep={,}}
  \hypersetup{colorlinks=true,allcolors=blue}
\else
  \setcitestyle{numbers,square}
  \hypersetup{hidelinks}
\fi
\usepackage[nameinlink]{cleveref}
\usepackage[T1]{fontenc}
\usepackage{lmodern} % Scalable LaTeX lettering, shared with the Manim slides
\usepackage{xcolor}
% Draft controls. Set these independently for review, print, or arXiv versions.
\newif\ifDKAnonymous
\DKAnonymousfalse
\ifdefined\DKAnonymousReview\DKAnonymoustrue\fi
\newif\ifDKColoredSymbols
\ifDKAuthorStyle\DKColoredSymbolstrue\else\DKColoredSymbolsfalse\fi
\ifdefined\DKMonochromeMath\DKColoredSymbolsfalse\fi
% The colors retain the Manim glossary's light-theme mathematical roles.
\input{math_symbol_colors.tex}
\ifDKColoredSymbols\else\DKMathColorsfalse\fi
\ifDKAuthorStyle\else\DKSymbolTooltipsfalse\fi




\usepackage[most]{tcolorbox}
\usepackage{listings}
\usepackage{listingsutf8}
\usepackage{fontspec}
\IfFileExists{emoji.sty}{%
  \IfFontExistsTF{Noto Color Emoji}{\usepackage{emoji}}{}%
}{}
\providecommand{\emoji}[1]{\textnormal{[#1]}} % Review notes also work without an emoji font.

\setlength{\emergencystretch}{3em}
\raggedbottom

% ArXiv-friendly Lean code: no shell escape, minted, or external highlighter.
\definecolor{leankeyword}{HTML}{243B6B}
\definecolor{leantactic}{HTML}{70477A}
\definecolor{leantype}{HTML}{1F6473}
\definecolor{leancomment}{HTML}{39734D}
\definecolor{leanstring}{HTML}{8A4B2A}
\definecolor{leannumber}{HTML}{777777}
\definecolor{leanbackground}{HTML}{F7F7F3}
\lstdefinelanguage{lean}{%
  sensitive=true,
  morecomment=[l]{--},
  morecomment=[s]{/-}{-/},
  morestring=[b]",
  morekeywords=[1]{import,open,namespace,section,end,noncomputable,variable,variables,
    universe,universes,set_option,attribute,axiom,opaque,theorem,lemma,def,abbrev,
    example,instance,structure,class,inductive,notation,local,deriving,extends,
    by,where,let,have,show,fun,match,with,if,then,else,do,return,in,forall},
  morekeywords=[2]{simp,simpa,rw,erw,dsimp,unfold,ring,ring_nf,nlinarith,linarith,
    positivity,norm_num,field_simp,omega,exact,refine,apply,assumption,intro,intros,
    ext,funext,cases,rcases,obtain,constructor,left,right,classical,by_contra,
    contrapose,convert,using,suffices,calc},
  morekeywords=[3]{Prop,Type,Nat,Int,Real,Fin,Set,Submodule,Matrix,EuclideanSpace,
    LinearMap,ContinuousLinearMap,LinearPMap,RCLike,IsTrialResidual,
    IsExactSpectralDecomposition,FormBoundedSylvesterGap,
    NormalizedSymmetricOperatorIdealFamily},
}
\lstdefinestyle{leanpaper}{%
  language=lean,
  columns=fullflexible,
  backgroundcolor=\color{leanbackground},
  basicstyle=\ttfamily\scriptsize,
  commentstyle=\color{leancomment}\itshape,
  keywordstyle=[1]\color{leankeyword}\bfseries,
  keywordstyle=[2]\color{leantactic},
  keywordstyle=[3]\color{leantype},
  stringstyle=\color{leanstring},
  numberstyle=\tiny\color{leannumber},
  breakatwhitespace=false,
  breaklines=true,
  captionpos=b,
  keepspaces=true,
  numbers=left,
  numbersep=6pt,
  showspaces=false,
  showstringspaces=false,
  showtabs=false,
  tabsize=2,
  aboveskip=0.8em,
  belowskip=0.8em,
  literate=
    {\\KField}{{$\Bbbk$}}1
    {∃}{{$\exists$}}1
    {∧}{{$\land$}}1
    {∀}{{$\forall$}}1
    {≠}{{$\neq$}}1
    {⊤}{{$\top$}}1
}
\crefname{lstlisting}{formalization}{formalizations}
\Crefname{lstlisting}{Formalization}{Formalizations}
\renewcommand{\lstlistingname}{Formalization}

\theoremstyle{thmstyleone}
\newtheorem{theorem}{Theorem}
\newtheorem{proposition}[theorem]{Proposition}
\newtheorem{lemma}[theorem]{Lemma}
\newtheorem{remark}[theorem]{Remark}
\newtheorem{corollary}[theorem]{Corollary}

\newcommand{\Lean}{\textsc{Lean}}
\newcommand{\DK}{Davis--Kahan}
\newcommand{\TauCeti}{Tau Ceti}
% Use \ie{} and \eg{}; supply surrounding punctuation at the call site.
\newcommand{\ie}{i.e.\@}
\newcommand{\eg}{e.g.\@}
\newcommand{\DKResultCount}{29}
\newcommand{\DKProvedCount}{28}
\newcommand{\DKRefutedCount}{1}

% Usage: \margincomment{Check this formula!}
%\newcommand{\jonsays}[1]{%
%  \marginpar{%
%    \raggedright\color{red}\tiny% Change to \footnotesize if text is too small
%    \emoji{person-raising-hand} \textbf{Comment:}\\#1%
%  }%
%}

% Usage: \margincomment{Your comment that now wraps perfectly!}
\newcommand{\jonsays}[1]{%
  \marginpar{%
    \raggedright% Forces text to wrap beautifully without ugly hyphenation
    \color{red}\scriptsize% Using \scriptsize makes the text tighter
    \emoji{person-raising-hand}: \\#1%
  }%
}
%\newcommand{\jonsays}[1]{{\color{red}\small\textbf{[\emoji{person-raising-hand}: #1]}}}

% Usage: \margincomment{Your comment that now wraps perfectly!}
%\newcommand{\agentsays}[1]{%
%  \marginpar{%
%    \raggedright% Forces text to wrap beautifully without ugly hyphenation
%    \color{red}\scriptsize% Using \scriptsize makes the text tighter
%    \emoji{robot}: \\#1%
%  }%
%}
\newcommand{\agentsays}[1]{{\color{red}\small\textbf{[\emoji{robot}: #1]}}}

\newcommand{\jonrsp}[1]{{\color{red}\small\textbf{[\emoji{person-raising-hand}: #1]}}}

\title[Formalizing Davis--Kahan Perturbation Theory]{Formalizing Davis--Kahan Perturbation Theory in Lean 4}
% Additional authors are expected; keep the confirmed ordering below stable.
\ifDKAnonymous
  \author{Anonymous authors}
\else
  \author*[1]{\fnm{Jonathan} \sur{Crall}}\email{jon.crall@kitware.com}
  \author[2]{\fnm{Edward} \sur{Wang}}\email{ewang43@jhu.edu}
  \author[1]{\fnm{Brian} \sur{Hu}}\email{brian.hu@kitware.com}
  \author[2]{\fnm{Carey E.} \sur{Priebe}}\email{cep@jhu.edu}
  \affil[1]{\orgname{Kitware, Inc.},
    \orgaddress{\city{Clifton Park}, \state{NY}, \country{USA}}}
  \affil[2]{\orgdiv{Department of Applied Mathematics and Statistics},
    \orgname{Johns Hopkins University},
    \orgaddress{\city{Baltimore}, \state{MD}, \country{USA}}}
\fi

\abstract{We formalize the spectral-subspace perturbation theory of Davis and
Kahan's 1970 paper in Lean~4. Of \DKResultCount{} selected results, we prove
\DKProvedCount{} under their stated hypotheses and refute the remaining
printed statement. The development includes the sine and tangent theorems,
direct rotations, principal-angle operators, Sylvester separation, and
operator-ideal norms on real and complex Hilbert spaces, including unbounded
self-adjoint operators. We explain the mathematical lineage connecting these
constructions. The refuted statement concerns minimization of full ambient
displacement by direct rotation. We give real four-dimensional counterexamples
at every prescribed positive acute angle and identify the failed
singular-value comparison. A corrected optimality statement holds for Q-norms,
whose squares are unitarily invariant norms of the squared modulus, including
the Frobenius norm. The formalization builds on Mathlib and adapts Spectra's
projection-valued-measure and Borel-calculus constructions. Portions of the
operator theory developed here were subsequently adapted into Tau Ceti.
These foundations also support an exact reduction of trace-norm optimal real
subspace transport to weighted matching on principal angles.}

\keywords{Davis--Kahan theory, spectral perturbation, formal verification,
Lean 4, direct rotation, principal angles}
\pacs[MSC Classification (2020)]{68V20, 47A55, 15A42}

\begin{document}
\maketitle

\ifDKAnonymous\else
\enlargethispage{2\baselineskip}
\vspace{-18pt}
\noindent\begin{minipage}{\linewidth}\footnotesize
\noindent\textbf{Acknowledgments.}
This material is based upon work supported by the Defense Advanced Research
Projects Agency (DARPA) under Contract No. HR001125CE017. Any opinions, findings
and conclusions or recommendations expressed in this material are those of the
author(s) and do not necessarily reflect the views of the Defense Advanced
Research Projects Agency.
\end{minipage}
\fi
\clearpage

\section{Introduction}
\label{sec:introduction}

An approximate eigenvector gives a candidate direction for the action of an
operator. If a unit vector $\sym{v}$ is paired with an approximate eigenvalue $\sym{mu}$,
the residual $\sym{A}\sym{v}-\sym{mu} \sym{v}$ measures how closely it satisfies the eigenvector
equation. A small residual can identify an accurate approximation to an
invariant subspace even when repeated or nearby eigenvalues make individual
eigenvectors difficult to distinguish. Davis--Kahan theory bounds the angles
between approximate and exact spectral subspaces in terms of a residual or
perturbation and a separation between the relevant spectra
\citep{davis1970rotation}.

In applications, one often works with a perturbed operator
  $\sym{Acurrent}=\sym{A}+\sym{H}$ while the reference operator $\sym{A}$ is
  unknown.
Taking an eigenspace of $\sym{Acurrent}$ as a trial subspace for $\sym{A}$, the norm of its
  residual relative to $\sym{A}$ is bounded by the norm of $\sym{H}$, allowing
  Davis--Kahan theory to bound the angle between the two subspaces without
  requiring $\sym{A}$ to be known or its residuals to be evaluated directly.

Our project began with a variant of this bound needed for
\ifDKAnonymous a spectral-estimation application\else the data kernel
perspective space (DKPS) \citep{helm2026query}\fi.

We extended the formalization to the results established in the original 1970
paper, including real and complex Hilbert spaces, unbounded self-adjoint
operators, and unitarily invariant operator-ideal norms. In particular, the
spectral-separation hypotheses allow half-infinite intervals
(e.g.\ $(-\infty,a]$). These settings require operator domains, spectral
calculus, and norm estimates beyond finite matrices.

We formalized the four headline theorems \citep[Section~2]{davis1970rotation} and every
named theorem, proposition, lemma, and corollary established in the 1970 paper.

This gives \DKResultCount{} targets, enumerated in
Appendix~\ref{app:dk-inventory}. We prove \DKProvedCount{} of them and refute
the printed statement \citep[Proposition~4.4]{davis1970rotation}.
That proposition asserts a full-displacement
optimality property for real direct rotations. We give a family of real
four-dimensional examples exhibiting trace-norm failure at every prescribed
positive acute angle and identify
the invalid singular-value comparison in its proof. The phase construction
behind the counterexample occurs over the complex field in the authors'
example \citep[Example~4.2]{davis1970rotation}. To the best of our knowledge,
its contradiction of the printed real, small-angle claim has not previously
been made explicit.
For the constructed pair, whose two principal angles both equal
$\sym{theta}$, the competing transporter attains the global trace-norm minimum
$4\sin\sym{theta}$ for the constructed subspaces. Section~\ref{sec:prop44}
compares these statements and repairs the optimality claim by restricting
the norm class to Q-norms (N is a Q-norm when $N(T)^2 = M(T^* T)$ for some
unitarily invariant norm $M$). This repair follows from the squared-displacement
comparison \citep[Proposition~4.3]{davis1970rotation};
earlier Q-norm minimization results appear in \citet{bolotnikov2001minimal}.

% TODO (lineage): Consolidate the scattered lineage records through the
% maintained dev/tauceti/dkps-upstream-contribution-census.md and the evidence
% index in AUTHORING_NOTES.md. Audit each attribution against an immutable
% donor commit/path, the accepted upstream commit/PR, and the corresponding
% declaration statements and proofs. Preserve Mathlib prerequisites, Spectra
% Spectra-derived constructions where applicable, and independent work as distinct relationships.
% Courant--Fischer/Weyl checked 2026-10-28: developed here in June 2026;
% August donor df78246f965881705c9301f96feddc270c98ceb1 predates the
% accepted Tau Ceti port 7ebf17ba04725c1a575ecf3d099ce0630cd949a0 (#13599).
% See COURANT_FISCHER_LINEAGE_AUDIT.md for the proof comparison and sources.

The proofs required us to formalize a significant amount of operator
theory, namely: principal angles and direct rotations, approximation numbers
and symmetric operator ideals, Borel spectral calculus, and Sylvester
estimates for unbounded operators. These developments build on Mathlib and
adapt Spectra's projection-valued-measure and Borel-calculus constructions
\citep{mathlib2020,bornemann2026spectra}.
These operator-theory foundations include our
Courant--Fischer and Weyl proofs, rectangular unitarily invariant seminorms,
and projection-valued measures, which were subsequently adapted into Tau
Ceti \citep{tauceti2026}.
Section~\ref{sec:related} explains the mathematical history and formal
provenance; Section~\ref{sec:foundations} describes the foundations used in
our proofs.

The same principal-angle and operator-ideal results support a further
optimization theorem. In Section~\ref{sec:example-transport}, we prove that
minimizing trace-norm displacement over all real orthogonal maps carrying
one subspace onto another reduces to weighted matching on principal angles.

\noindent\textbf{Contributions.}
\begin{enumerate}
\item \textbf{Davis--Kahan formalization.} We prove 28 of the 29 source
results and refute the full-displacement assertion
\citep[Proposition~4.4]{davis1970rotation}. The sine-theta theorem covers real and
complex Hilbert spaces, unbounded self-adjoint operators, the source's
spectral-separation alternatives, and operator-ideal norms where finite
(\cref{sec:coverage,sec:sine-theta}).
\item \textbf{Formalized operator-theoretic foundations.} We formalized the
geometric and analytic results required by these theorems, adapted existing
formal mathematics, and contributed foundations later used in Tau Ceti
(\cref{sec:related-formal,sec:foundations}).
\item \textbf{Source correction.} We identify a false comparison
of Ky Fan sums in the printed claim \citep[Proposition~4.4]{davis1970rotation},
give a real counterexample at every positive acute angle,
and prove the Q-norm consequence of the valid squared-displacement theorem
(\cref{sec:prop44}).
\item \textbf{A reduction to weighted matching.} We reduce unrestricted
real trace-norm transport to a weighted-matching problem on principal angles
(\cref{sec:example-transport}).

\end{enumerate}

\section{Mathematical background}
\label{sec:related}

The Davis--Kahan angle estimates use three constructions. Principal angles
measure the error between subspaces; a residual measures how well a proposed
subspace represents the operator; and a Sylvester estimate bounds the first
quantity by the second under spectral separation. Operator ideals allow the
same argument to express several norms. We describe the development of these
constructions before stating the formalized sine-theta theorem.

Existing historical surveys cover these developments. Paige and Wei trace
the cosine--sine decomposition from Jordan's subspace geometry through the
principal ideas assembled by Davis and Kahan to Stewart's later completion
and explicit formulation \citep{paige1994history}.
Simon's historical survey of Kato's work in non-relativistic quantum mechanics
traces Sz.-Nagy's unitary construction for pairs of projections and Kato's
later algebraic formulation \citep[\S5]{simon2018kato}.
We use this account for the history of the projection transport described in
\cref{sec:related-projections}.
B\"ottcher and Spitkovsky survey pairs of projections and their operator algebras.
Their historical discussion relates Halmos's 1969 block representation to the
earlier formulations of Dixmier (1948) and Davis (1958), and to the related
projection constructions in Davis and Kahan's 1970 paper
\citep[Remark~1.3]{bottcher2010gentle}.
We use these accounts together with the primary
sources. Table~\ref{tab:historical-lineage} summarizes the mathematical
constructions used in our formalization.

\begin{table}[!htbp]
\centering
\small
\setlength{\tabcolsep}{5pt}
\renewcommand{\arraystretch}{1.18}
\begin{tabularx}{\linewidth}{@{}>{\raggedright\arraybackslash}p{0.26\linewidth}>{\raggedright\arraybackslash}X>{\raggedright\arraybackslash}p{0.28\linewidth}@{}}
\toprule
\textbf{Lineage} & \textbf{Mathematical role} & \textbf{Formal dependency}\\
\midrule
Jordan (1875); Davis (1958); Halmos (1969) & Principal angles, invariant subspace geometry,
  and pairs of projections & Angle operators; projection blocks\\
Rayleigh (1877); Ritz (1909); Weyl (1912) & Approximate eigenpairs,
  variational spectra, and eigenvalue stability & Trial coordinates;
  min--max and eigenvalue inequalities\\
Rellich (1937--1942); Sz.-Nagy (1947); Kato (1966) & Moving spectral projections and unitary
  transport between nearby subspaces & Spectral calculus; polar
  direct rotation\\
Sylvester (1884); Rosenblum (1956, 1969) & Spectral separation controls an
  off-diagonal operator & Domain-aware Sylvester estimates\\
Von Neumann (1937); Schatten (1950); Ky Fan (1951) & Comparison of spectral errors
  through singular-value gauges & Approximation numbers;
  symmetric operator ideals\\
Davis (1963, 1965); Davis--Kahan (1969, 1970) & Residual-to-angle bounds, tangent estimates,
  rotation extremality, and spectral branch selection & Source-level
  sine, tangent, and direct-rotation theorems\\
\bottomrule
\end{tabularx}
\caption{Mathematical sources used in the Davis--Kahan formalization.
Dates identify the publications discussed in the text.}
\label{tab:historical-lineage}
\end{table}
\FloatBarrier

\subsection{Principal angles}
\label{sec:related-classical}

Camille Jordan studied the geometry of linear subspaces in $n$ dimensions in
  1875 \citep{jordan1875essai}.

For two lines, a single acute angle describes their relative position.
For subspaces of higher dimension, one needs a sequence of \emph{principal
  angles}.
The successive principal vectors are selected under orthogonality constraints,
  so their angles are independent of the coordinates used to describe either
  subspace.
This separates the geometry of a subspace from the arbitrary choice of its
  orthonormal basis.

For finite-dimensional subspaces $\sym{Usource}$ and $\sym{Vtarget}$ of equal dimension, choose
isometries $\sym{Xframe}$ and $\sym{Yframe}$ with those ranges.
The singular values of the overlap matrix are the principal cosines:
\begin{equation}
  \sym{SingularValue}(\like{Xframe}{\symglyph{Xframe}^*}\sym{Yframe})=\cos\like{theta}{\symglyph{theta}_j},
  \qquad 0\leq\like{theta}{\symglyph{theta}_j}\leq\frac{\sym{PiConstant}}{2}.
  \label{eq:history-principal-cosines}
\end{equation}

The corresponding principal vectors $\like{SourceVector}{\symglyph{SourceVector}_j}\in \sym{Usource}$ and $\like{TargetVector}{\symglyph{TargetVector}_j}\in \sym{Vtarget}$ satisfy
$\langle \like{SourceVector}{\symglyph{SourceVector}_i},\like{TargetVector}{\symglyph{TargetVector}_j}\rangle=\like{KroneckerDelta}{\symglyph{KroneckerDelta}_{ij}}\cos\like{theta}{\symglyph{theta}_j}$ in suitably oriented
coordinates.  For nonzero acute angles, each pair determines a real
two-dimensional plane in which the relative position is an ordinary
rotation.  An entire pair of subspaces is described by these planes,
together with their shared and orthogonal components.

Principal angles also arise in statistics: Hotelling's 1936 canonical
  correlation analysis sought pairs of maximally correlated linear combinations
  of two sets of variables \citep{hotelling1936relations}.
After appropriate whitening, the canonical correlations are singular values of a
  cross-coordinate map.
In numerical linear algebra, Bj\"orck and Golub developed algorithms for
  computing principal angles and vectors in 1973 \citep{bjorck1973numerical}.
Their computational formulation makes the basis-invariant angles available from
  orthonormal matrix factorizations.

A further formulation uses the cosine--sine decomposition of a
partitioned orthogonal or unitary matrix.
The blocks of the decomposition contain compatible diagonal cosine and sine
  factors.
Davis and Kahan describe the geometry using orthogonal projections,
angle operators, and the direct rotation carrying one subspace onto the other
\citep[Section~3]{davis1970rotation}. Their work assembled the principal ideas
of the cosine--sine decomposition; Stewart completed and explicitly formulated
the matrix decomposition in 1977
\citep{stewart1977perturbation,paige1994history}.
Paige and Wei explain how this decomposition and direct rotation represent
the same principal-angle data
\citep{paige1994history}. Our formalization uses projections and functional
calculus for the Hilbert-space results and principal-plane coordinates for
the finite-dimensional counterexample.

\subsection{Eigenvalue approximation}
\label{sec:related-ritz-weyl}

Jordan's angles answered how to compare subspaces once they are given.
An independent question came from mechanics and numerical analysis:
how can one approximate an eigenvalue or eigenmode when the underlying
differential equation cannot be solved exactly?  In the nineteenth
century Rayleigh estimated vibration frequencies using trial shapes
and their energy ratios \citep{rayleigh1877sound}.  Ritz developed a systematic finite-dimensional
variational approximation in the work published in 1909
\citep{ritz1909variations,leissa2005historical}.  The method selects
a trial space and studies the action of the operator through that
space. A trial space is a chosen space of candidate approximate solutions.
For an eigenvalue problem, one searches for an approximate eigenvector among
its vectors. This reduces an infinite-dimensional problem to a smaller
problem whose residual can then be assessed.

For a Hermitian matrix $\sym{A}$, the Rayleigh quotient
$\langle \sym{A}\sym{RayleighTestVector},\sym{RayleighTestVector}\rangle/\langle \sym{RayleighTestVector},\sym{RayleighTestVector}\rangle$ expresses the scalar action
of $\sym{A}$ along a trial vector. A normalized trial vector $\sym{v}$ gives the
Rayleigh--Ritz approximate eigenvalue $\sym{mu}=\langle \sym{A}\sym{v},\sym{v}\rangle$ and the
residual $\sym{A}\sym{v}-\sym{mu} \sym{v}$. The Courant--Fischer min--max
principle organizes the eigenvalues through extremal quotients over
subspaces \citep{bhatia1997matrix}.  Weyl's eigenvalue comparison
inequalities then control the motion of ordered eigenvalues under a
Hermitian perturbation \citep{weyl1912eigenwerte,bhatia1997matrix}.  In modern notation,
if $\sym{A}$ and $\sym{Acurrent}=\sym{A}+\sym{H}$ are Hermitian $n$-by-$n$ matrices and their eigenvalues
are arranged in the same order, then
\begin{equation}
  \max_j |\like{SpectralValue}{\symglyph{SpectralValue}_j}(\sym{A})-\like{SpectralValue}{\symglyph{SpectralValue}_j}(\sym{Acurrent})|\leq\dkOperatorNorm{\sym{H}}.
  \label{eq:history-weyl}
\end{equation}
This is an estimate on \emph{eigenvalues}; it does not tell us which
individual eigenvector of $\sym{A}$ corresponds to which eigenvector of $\sym{Acurrent}$.

The limitation already appears in two dimensions.  For any $\sym{EigenvalueExampleEpsilon}>0$,
put
\begin{equation}
 \like{A}{\symglyph{A}_\varepsilon}=\begin{pmatrix}0&0\\0&\sym{EigenvalueExampleEpsilon}\end{pmatrix},
 \qquad
 \like{Acurrent}{\symglyph{Acurrent}_\varepsilon}=\begin{pmatrix}\sym{EigenvalueExampleEpsilon}&0\\0&0\end{pmatrix}.
 \label{eq:history-swap-example}
\end{equation}
The matrices have identical ordered eigenvalues and
$\dkOperatorNorm{\like{A}{\symglyph{A}_\varepsilon}-\like{Acurrent}{\symglyph{Acurrent}_\varepsilon}}=\sym{EigenvalueExampleEpsilon}$,
yet their lowest-eigenvalue lines are perpendicular. The gap between the
two eigenvalues of each matrix is $\sym{EigenvalueExampleEpsilon}$, equal to the perturbation
norm. An eigenspace estimate therefore needs to compare the perturbation
size with the spectral separation.  For a cluster with repeated eigenvalues, an
entire invariant subspace is likewise more natural than a selected
orthonormal eigenbasis.

Rellich developed a rigorous perturbation theory of self-adjoint operators
in papers published from 1937 to 1942. Kato's 1966 book gives a systematic
account \citep{simon2018kato,kato1966perturbation}. The Davis--Kahan residual
bounds address a related approximation question: given a trial space and its
residual, spectral separation bounds its angle from an exact spectral
subspace. The estimate can therefore assess the accuracy of a
Rayleigh--Ritz approximation.

\subsection{Two orthogonal projections}
\label{sec:related-two-projections}

For a closed subspace of a Hilbert space $\sym{GeneralHilbertSpace}$, the orthogonal
projection remains a bounded operator even when the subspace is
infinite dimensional.  Denote the projections onto two closed
subspaces by $\sym{Psource}$ and $\sym{Qtarget}$. Compressing an operator to $\DKRange(\sym{Psource})$ means
restricting its input to that subspace and projecting its output back with
$\sym{Psource}$. The positive compression of $\sym{Qtarget}$ is therefore $\sym{Psource}\sym{Qtarget}\sym{Psource}$ restricted to
$\DKRange(\sym{Psource})$. It is the operator version of the squared overlap matrix: in finite dimensions its eigenvalues are
$\cos^2\like{theta}{\symglyph{theta}_j}$.  In Hilbert space it need not have a discrete
spectrum.  Functional calculus nevertheless defines its positive
square root and the associated cosine and sine of an angle operator.

Write $\sym{Theta0}$ for this directed angle operator on
$\DKRange(\sym{Psource})$, and $\sym{Theta1}$ for the corresponding
operator for the complementary pair $\sym{Identity}-\sym{Psource},\sym{Identity}-\sym{Qtarget}$.
The full angle operator is
\[
 \sym{Theta}=\sym{DiagonalOperator}(\sym{Theta0},\sym{Theta1})
 \quad\text{on }\DKRange(\sym{Psource})\oplus\DKKernel(\sym{Psource}).
\]
For $c>0$, the comparisons $\sym{Theta}\le c$ and $\sym{Theta}<c$ mean that its
operator norm is at most, or strictly less than, $c$, respectively.
In finite dimensions this norm is the largest eigenvalue of either directed
angle block, comparing the source with the target and their orthogonal
complements.
This includes right angles from unmatched directions: if a source vector is
orthogonal to the target, its cosine compression is zero and its directed
angle is $\sym{PiConstant}/2$. The compression construction applies even when the subspace
dimensions differ. Target directions orthogonal to the source appear as
right angles in the complementary block. Later, isometric frames express these angle blocks on
coordinate Hilbert spaces.

The relative position of two closed subspaces was studied by Dixmier
in 1948 and by Davis in his 1958 paper \emph{Separation of Two Linear
Subspaces} \citep{dixmier1948position,davis1958separation}.  Davis's
closeness and separation operators place the cosine/sine data in an
operator-theoretic setting.  Halmos's 1969 \emph{Two Subspaces} gives a
block representation of the two projections
\citep{halmos1969two}.  One first splits off the four intersection
spaces corresponding to the combinations of membership in the ranges
and kernels of $\sym{Psource}$ and $\sym{Qtarget}$.  On their remaining, generic part, the pair
is unitarily equivalent to
\begin{equation}
 \like{Psource}{\symglyph{Psource}_{\mathrm{gen}}}=\begin{pmatrix}\sym{Identity}&0\\0&0\end{pmatrix},
 \qquad
 \like{Qtarget}{\symglyph{Qtarget}_{\mathrm{gen}}}=\begin{pmatrix}\like{HalmosCosineBlock}{\symglyph{HalmosCosineBlock}^2}&\sym{HalmosCosineBlock}\sym{HalmosSineBlock}\\\sym{HalmosCosineBlock}\sym{HalmosSineBlock}&\like{HalmosSineBlock}{\symglyph{HalmosSineBlock}^2}\end{pmatrix},
 \qquad \like{HalmosCosineBlock}{\symglyph{HalmosCosineBlock}^2}+\like{HalmosSineBlock}{\symglyph{HalmosSineBlock}^2}=\sym{Identity},
 \quad \sym{HalmosCosineBlock}\sym{HalmosSineBlock}=\sym{HalmosSineBlock}\sym{HalmosCosineBlock},
 \label{eq:history-halmos-block}
\end{equation}
where $\sym{HalmosCosineBlock}$ and $\sym{HalmosSineBlock}$ are commuting positive contractions on a common
Hilbert space.  Thus an operator-valued angle replaces a finite list
of principal angles.  B\"ottcher and Spitkovsky survey this
representation and its consequences for the algebra generated by two
projections \citep{bottcher2010gentle}.

Among these four intersections, the two right-angle intersection
subspaces (called defect spaces below) are
\[
 \sym{CrossedSourceDefect}=\DKRange(\sym{Psource})\cap\DKKernel(\sym{Qtarget}),
 \qquad
 \sym{CrossedTargetDefect}=\DKKernel(\sym{Psource})\cap\DKRange(\sym{Qtarget}).
\]
These contain source directions orthogonal to the target, and target directions
orthogonal to the source, respectively. The pair is called \emph{acute} when
both are $\{\sym{ZeroVector}\}$, the zero subspace consisting only of the zero
vector. For a general pair these spaces affect whether a specified
intertwining operator exists. An operator $\sym{IntertwinerT}$ intertwines $\sym{Psource}$ and $\sym{Qtarget}$ when
$\sym{IntertwinerT}\sym{Psource}=\sym{Qtarget}\sym{IntertwinerT}$: applying $\sym{IntertwinerT}$ after the source projection agrees with applying the
target projection after $\sym{IntertwinerT}$. In particular, $\sym{IntertwinerT}$ maps the source range into
the target range and its orthogonal complement into the target complement.
The existence criteria for direct rotations include conditions on the
dimensions of these defect spaces \citep[Section~3]{davis1970rotation}.

A finite-dimensional orthogonal matrix decomposes into invariant real
lines and two-dimensional rotation planes.  In a general Hilbert space an
angle operator can instead have continuous spectrum.  The spectral
calculus of the pair of projections takes the place of a finite list
of angles, and the intersection and defect spaces cannot be discarded.
Our Hilbert-space direct-rotation proofs therefore use projection identities
and functional calculus. In finite dimensions these identities recover the
principal-plane construction.

\subsection{Transport between nearby spectral subspaces}
\label{sec:related-projections}

A perturbation can change the invariant subspace associated with an isolated
part of a self-adjoint operator's spectrum. Let $\sym{Psource}$ and
$\sym{Qtarget}$ be the orthogonal projections onto the corresponding spectral
subspaces before and after perturbation. To compare these subspaces, we seek
a unitary carrying $\DKRange(\sym{Psource})$ onto
$\DKRange(\sym{Qtarget})$ that depends continuously, or analytically, on the
projections.
Sz.-Nagy developed the nearby-projection construction in work
published in 1947 \citep{sznagy1947perturbations,simon2018kato}.  Kato gave a streamlined
algebraic construction, described in his 1955 report and later in
\emph{Perturbation Theory for Linear Operators}
\citep[Chapter~I, \S4.6, equations~(4.33)--(4.42)]{kato1966perturbation};
see also \citet[Theorem~5.1]{simon2018kato}.

For orthogonal projections the useful intertwiner is
\begin{equation}
 \sym{IntertwinerT}=\sym{Qtarget}\sym{Psource}+(\sym{Identity}-\sym{Qtarget})(\sym{Identity}-\sym{Psource}), \qquad \sym{IntertwinerT}\sym{Psource}=\sym{Qtarget}\sym{IntertwinerT}.
 \label{eq:history-kato-T}
\end{equation}
The identity
\begin{equation}
 \like{IntertwinerT}{\symglyph{IntertwinerT}^*}\sym{IntertwinerT}=\sym{Identity}-(\sym{Psource}-\sym{Qtarget})^2
 \label{eq:history-kato-square}
\end{equation}
shows why $\|\sym{Psource}-\sym{Qtarget}\|<1$ is the natural threshold: it makes $\like{IntertwinerT}{\symglyph{IntertwinerT}^*}\sym{IntertwinerT}$
strictly positive and invertible.  Hence
\begin{equation}
 \sym{Rrotation}=\sym{IntertwinerT}\bigl(\sym{Identity}-(\sym{Psource}-\sym{Qtarget})^2\bigr)^{-1/2},
 \qquad \sym{Rrotation}\sym{Psource}\like{Rrotation}{\symglyph{Rrotation}^*}=\sym{Qtarget}
 \label{eq:history-kato-unitary}
\end{equation}
is unitary. This is the polar decomposition $\sym{IntertwinerT}=\sym{Rrotation}|\sym{IntertwinerT}|$, where
$|\sym{IntertwinerT}|=(\like{IntertwinerT}{\symglyph{IntertwinerT}^*}\sym{IntertwinerT})^{1/2}$ is positive. The factor $|\sym{IntertwinerT}|$ describes stretching;
the unitary factor $\sym{Rrotation}$ preserves lengths. Here $\sym{IntertwinerT}$ is invertible, so
$\sym{Rrotation}=\sym{IntertwinerT}|\sym{IntertwinerT}|^{-1}$. In principal-plane coordinates it is the direct rotation through the
principal angle, acting trivially on common directions.

The transport equation
$\sym{Wcompetitor}\sym{Psource}\sym{Wcompetitor}^*=\sym{Qtarget}$
has many unitary solutions.
The polar-factor formula selects one determined by the two projections
and continuously connected to the identity under the norm bound above. That
choice makes the construction valuable in perturbation theory: the
transport identifies moving spectral subspaces while avoiding an
arbitrary choice of orthonormal bases. In infinite dimensions,
$\|\sym{Psource}-\sym{Qtarget}\|<1$ is stronger than acuteness: it also
bounds the operator angles away from $\sym{PiConstant}/2$. For equal-dimensional
finite subspaces the conditions coincide. For an orthogonal or unitary transporter $\sym{Wcompetitor}$, the operator
$\sym{Identity}-\sym{Wcompetitor}$ records the displacement of every vector. Davis and Kahan ask which
norms of this displacement are minimized by the direct rotation
\citep[Section~4]{davis1970rotation}.

Kato used such transformations to identify nearby spectral
subspaces of analytic operator families, thereby reducing the study
of an isolated moving eigenvalue cluster to a fixed coordinate
space.  Simon's historical account explains the earlier Sz.-Nagy
construction, Kato's formula, and subsequent work on pairs of
projections \citep[\S5]{simon2018kato}. Our proofs of direct-rotation
existence, uniqueness, and symmetry use the square-root and
polar-decomposition identities in
\cref{eq:history-kato-square,eq:history-kato-unitary}. The condition
$\|\sym{Psource}-\sym{Qtarget}\|<1$ applies in arbitrary Hilbert dimension.

\subsection{Spectral projections and residuals}
\label{sec:related-spectral}

The spectral theorem represents a self-adjoint operator as
$\sym{A}=\int\sym{SpectralValue}\,d\sym{SpectralMeasure}(\sym{SpectralValue})$ for a projection-valued measure $\sym{SpectralMeasure}$.
For a Borel set $\sym{SpectralSet}$, the operator $\sym{SpectralMeasure}(\sym{SpectralSet})$ is the orthogonal projection onto
the corresponding spectral subspace. This gives a second way of obtaining
a subspace: select a region of the spectrum and take the range of its
spectral projection. The functional-calculus
framework developed from the work of Riesz, von Neumann, and others;
Kato made spectral projections central to the perturbation theory of
closed operators \citep{kato1966perturbation,simon2018kato}.
Unlike the finite-dimensional eigenspace picture, a spectral
subspace need not be spanned by a finite list of eigenvectors.
A self-adjoint operator may be unbounded, while each of its spectral
projections is bounded and everywhere defined. The Hellinger--Toeplitz
theorem states that an everywhere-defined symmetric operator on a Hilbert
space is bounded \citep{hellinger1910grundlagen}. Thus an unbounded
self-adjoint operator has a proper domain. Its action $\sym{A}x$ is defined only
for $x\in\DKDom(\sym{A})$, so residual and intertwining equations must include
the required domain inclusions.

In the complex case, a number $z$ belongs to the resolvent set when
$z\sym{Identity}-\sym{A}$ has a bounded, everywhere-defined inverse. That inverse,
$(z\sym{Identity}-\sym{A})^{-1}$, is the resolvent at $z$. For an isolated bounded spectral
cluster, choose a positively oriented contour $\sym{RieszContour}$ in the resolvent
set enclosing the cluster and no other spectrum. Its projection is
\begin{equation}
 \sym{RieszProjection}=\frac{1}{2\sym{PiConstant} i}\int_{\sym{RieszContour}}(z\sym{Identity}-\sym{A})^{-1}\,dz.
 \label{eq:history-riesz-projection}
\end{equation}
The resolvent is bounded at each point of $\sym{RieszContour}$ even when $\sym{A}$ is unbounded.
This makes the spectral projection a stable object for perturbation analysis,
  while preserving the separate domain requirements whenever $\sym{A}$ itself appears
  in an identity.
More general Borel spectral projections cover finite and half-infinite intervals,
  including spectral configurations for which a single isolated
  finite-eigenvalue cluster is not the right model.
Our formalization uses Borel functional calculus to construct projections for
  these spectral sets, including for unbounded self-adjoint operators.

The Courant--Fischer min--max principle and Weyl-type eigenvalue
perturbation bounds provide complementary information about spectral
\emph{values} \citep{bhatia1997matrix,kato1966perturbation}.  They do
not by themselves identify a stable basis inside a multiple-eigenvalue
cluster, or quantify the movement of its invariant subspace.  The Courant--Fischer and Weyl results that we formalized describe extremal
Rayleigh quotients and eigenvalue stability. The angle estimates add control
of the associated subspaces.

Davis studied the rotation of eigenvectors under perturbation in
1963 and 1965 \citep{davis1963rotation,davis1965rotation}.  The
1969 Davis--Kahan announcement preceded their extended 1970 treatment
\citep{daviskahan1969bounds,davis1970rotation}.  Their problem has
two closely related formulations.  In one, the original self-adjoint
operator is perturbed and the movement of an exact spectral subspace
is estimated.  In the other, one has an approximate trial subspace and
a computable residual: how far can it be from the desired exact
spectral subspace?  Unlike a comparison of individual eigenvectors,
these are meaningful even when an eigenvalue has multiplicity and its
orthonormal eigenbasis is nonunique.

Write $\sym{A}$ for a reference operator and $\sym{Acurrent}=\sym{A}+\sym{H}$ for the perturbed
ambient operator. We use this convention in the sine-theta statement below.
Let $\sym{E0}$ embed a trial-coordinate Hilbert space into the ambient space
isometrically, and let $\sym{TrialOperator}$ be a self-adjoint operator on those coordinates.
The residual is the bounded operator agreeing, where defined, with
\begin{equation}
 \sym{R}=\sym{Acurrent}\sym{E0}-\sym{E0}\sym{TrialOperator}.
 \label{eq:history-residual}
\end{equation}
If $\sym{R}=0$, applying $\sym{Acurrent}$ after the embedding agrees with applying
$\sym{TrialOperator}$ before it: the trial subspace is invariant and $\sym{TrialOperator}$ represents its
operator. For $\sym{R}\ne0$, the residual measures the failure of this identity.  A second isometric coordinate map for the
desired exact spectral space defines its orthogonal projection $\sym{Pexact}$.
The positive operator
\begin{equation}
 \sym{sinTheta0}=\bigl(\sym{Identity}-\like{E0}{\symglyph{E0}^*}\sym{Pexact}\sym{E0}\bigr)^{1/2}
 \label{eq:history-angle-sine}
\end{equation}
measures the component of a trial vector lying outside the exact
subspace.  The critical distinction for the formal proof is that this
positive sine acts on the trial-coordinate space, whereas
$(\sym{Identity}-\sym{Pexact})\sym{E0}$ maps into the ambient space.  For an operator $B$, its positive modulus $|B|=(B^*B)^{1/2}$ acts on its
input space and records its stretching factors. The modulus of $(\sym{Identity}-\sym{Pexact})\sym{E0}$
is $\sym{sinTheta0}$, so the two maps have equal unitarily invariant ideal
norms whenever those norms are finite.

\subsection{Sylvester estimates}
\label{sec:related-sylvester}

The matrix equation $\sym{SylvesterLeft}\sym{X}-\sym{X}\sym{HistoricalSylvesterRight}=\sym{SylvesterRHS}$ is associated with Sylvester's 1884
work \citep{sylvester1884equation}.  Rosenblum studied separated
operator spectra in 1956 and the equation with self-adjoint coefficients
in 1969 \citep{rosenblum1956equation,rosenblum1969selfadjoint}.

The scalar prototype is $(\sym{SylvesterScalarLeft}-\sym{SylvesterScalarRight})\sym{SylvesterScalarUnknown}=\sym{SylvesterScalarRHS}$: if $|\sym{SylvesterScalarLeft}-\sym{SylvesterScalarRight}|\geq\sym{delta}>0$,
then $\sym{delta}|\sym{SylvesterScalarUnknown}|\leq|\sym{SylvesterScalarRHS}|$.  In the operator problem the scalar
multipliers are replaced by the spectral actions of two operators.
The operator estimate has the form $\sym{delta} \sym{NormN}(\sym{X})\leq \sym{NormN}(\sym{SylvesterLeft}\sym{X}-\sym{X}\sym{HistoricalSylvesterRight})$.
Its hypotheses specify the spectral positions of $\sym{HistoricalSylvesterRight}$ and $\sym{SylvesterLeft}$, their domains
when unbounded, and the unitarily invariant norm $\sym{NormN}$.

Let $\sym{F1}$ give isometric coordinates for the exact complementary
spectral subspace, with self-adjoint block $\sym{Lambda1}$.  Writing
$\sym{X}=\like{F1}{\symglyph{F1}^*}\sym{E0}$ converts the residual identity into a Sylvester equation:
\begin{equation}
 \sym{Lambda1}\sym{X}-\sym{X}\sym{TrialOperator}=\like{F1}{\symglyph{F1}^*}\sym{R}
 \label{eq:history-sylvester}
\end{equation}
on the appropriate domain. Thus $\sym{X}$ represents the trial frame's
``leakage'' (i.e., its orthogonal projection onto the exact complementary
spectral subspace, expressed in $\sym{F1}$ coordinates), while the right
side is controlled by the computable residual.  Under the interval/exterior
or ordered spectral-gap hypotheses of the sine-theta theorem,
separated Sylvester estimates give
\begin{equation}
 \sym{delta}\,\sym{NormN}(\sym{X})\leq \sym{NormN}(\like{F1}{\symglyph{F1}^*}\sym{R})\leq \sym{NormN}(\sym{R}).
 \label{eq:history-separation-bound}
\end{equation}
This is the bridge from a spectral-gap assumption to the bound on the
sine of a subspace angle.

The coefficient of $\sym{NormN}(\sym{R})/\sym{delta}$ in this bound is one.
This coefficient uses the ordered or interval/exterior spectral
positions specified by Davis and Kahan. Arbitrarily interlaced spectra can
require a different constant. Our theorem includes both of the source's
spectral configurations, including ordered cases where both operators are
unbounded.  A formal implementation
therefore needs the domains of unbounded self-adjoint operators,
their spectral projections, strong spectral cutoffs, and the
ideal-norm inequalities used to pass through bounded approximations.

\subsection{Operator ideals and displacement norms}
\label{sec:related-norms}

The geometry of a pair of subspaces is fixed, but there are many ways
to assign a cost to a unitary transporting one onto the other.
Von Neumann's 1937 study of unitarily invariant matrix norms identified
the role of symmetric gauge functions of singular values
\citep{vonneumann1937matrix}. Schatten's 1950 monograph developed normed
operator ideals; Ky Fan's 1951 eigenvalue and singular-value inequalities
provided the partial-sum comparisons underlying majorization
arguments \citep{schatten1950crossspaces,fan1951maximum,bhatia1997matrix}.
The singular values $\sym{SingularValue}(Z)$ are the nonnegative square roots of the
eigenvalues of $Z^*Z$, ordered decreasingly. They record the stretching
factors of the matrix $Z$. For $1\leq p<\infty$, the Schatten norms include
\begin{equation}
 \|Z\|_{S_p}=\left(\sum_j \sym{SingularValue}(Z)^p\right)^{1/p},
 \qquad
 \|Z\|_{S_1}=\sum_j \sym{SingularValue}(Z),
 \quad
 \|Z\|_{S_2}=\left(\sum_j \sym{SingularValue}(Z)^2\right)^{1/2}.
 \label{eq:history-schatten}
\end{equation}
An operator ideal is a linear class of operators closed under multiplication
on either side by bounded operators. A normed ideal also controls the norm
of these products. For compact operators in infinite dimensions, the
Schatten $p$ ideal consists of those whose singular values have a finite
$p$-power sum. Thus a trace-norm inequality requires trace-class membership;
this is an instance of the source's qualification that its norms be defined.

The largest singular value gives the operator norm, the Euclidean
norm of the singular-value sequence gives the Hilbert--Schmidt norm
(the Frobenius norm $\dkFrobeniusNorm{Z}=\|Z\|_{S_2}$ for matrices),
and its sum gives the trace norm. These norms can rank transporters
differently. In Section~\ref{sec:prop44}, the two displacement singular-value
lists are $(2\sin(\sym{theta}/2),\ldots,2\sin(\sym{theta}/2))$ with four entries and
$(2\sin\sym{theta},2\sin\sym{theta},0,0)$. The second has a smaller sum and a larger
maximum. Ky Fan norms compare partial sums of the ordered singular values;
these comparisons determine which norm inequalities follow.

Davis established Hilbert--Schmidt and finite-dimensional Schatten
extremal properties for exponents at least two in his 1958 subspace
paper \citep[Theorems~7.3--7.4]{davis1958separation}.  Davis and
Kahan compare three objectives \citep[Section~4]{davis1970rotation} for a transporter $\sym{Wcompetitor}$ carrying
$\sym{Usource}$ onto $\sym{Vtarget}$: its restricted displacement $(\sym{Identity}-\sym{Wcompetitor})\like{Psource}{\symglyph{Psource}_U}$, its squared full
displacement $(\sym{Identity}-\sym{Wcompetitor})^*(\sym{Identity}-\sym{Wcompetitor})$, and its full displacement $\sym{Identity}-\sym{Wcompetitor}$. The last acts
on every ambient vector, including vectors outside $\sym{Usource}$. They gave a
complex-phase modification that decreases its trace norm \citep[Example~4.2]{davis1970rotation}.
The full-displacement claim \citep[Proposition~4.4]{davis1970rotation}
asserts that direct rotation minimizes
full displacement for real subspaces if their principal angles are no
larger than $\sym{PiConstant}/3$.  The four-dimensional example in
\cref{sec:prop44} gives a smaller trace norm at every positive acute
angle.  It exposes the failed step: adding singular-value sums of
restrictions to separate principal planes is not a valid global lower
bound for an arbitrary ambient transporter.

The subsequent minimal-distortion theory of Bolotnikov, Li, and Rodman
studied the same class of subspace-transport constraints, including
Frobenius and Q-norm optimality and explicit subquadratic
counterexamples \citep{bolotnikov2001minimal}.  Later work on
Grassmannian metrics and canonical-angle inequalities also separates
intrinsic paths between subspaces from minimizing the extrinsic
operator displacement \citep{qiu2006grassmann,zhang2010canonicalangles}.
These results give the norm-theoretic context for the source
correction and for the trace-norm matching problem in
\cref{sec:example-transport}.

\subsection{The Davis--Kahan angle estimates}
\label{sec:related-synthesis}

Davis's papers of 1963 and 1965 and the Davis--Kahan announcement of 1969
preceded the extended 1970 treatment
\citep{davis1963rotation,davis1965rotation,daviskahan1969bounds,davis1970rotation}.
For the sine-theta estimate, let $\sym{R}=\sym{Acurrent} \sym{E0}-\sym{E0}\sym{TrialOperator}$ be the residual
and $\sym{Pexact}$ the desired exact spectral projection. The operator $(\sym{Identity}-\sym{Pexact})\sym{E0}$
extracts the part of a trial vector outside that spectral subspace. In
isometric coordinates $\sym{F1}$ for the complementary space, this component is
represented by $\sym{X}=\like{F1}{\symglyph{F1}^*}\sym{E0}$. The residual identity gives
\cref{eq:history-sylvester}, and a separated Sylvester estimate bounds $\sym{NormN}(\sym{X})$.
Since $\sym{S}=\sym{F1}\sym{X}$ and $\sym{F1}$ is an isometry,
$\sym{NormN}(\sym{X})=\sym{NormN}(\sym{S})$.
The modulus identity $|\sym{S}|=\sym{sinTheta0}$ gives
$\sym{NormN}(\sym{S})=\sym{NormN}(\sym{sinTheta0})$.

The resulting chain of equalities and inequalities is
\begin{equation}
 \underbrace{\sym{delta} \sym{NormN}(\sym{sinTheta0})}_{\text{subspace error}}
 =\sym{delta} \sym{NormN}(\sym{X})
 \leq \sym{NormN}(\like{F1}{\symglyph{F1}^*}\sym{R})
 \leq\underbrace{\sym{NormN}(\sym{R})}_{\text{computable residual}}.
 \label{eq:history-unified-chain}
\end{equation}
The first equality follows from the positive-modulus identity for the
complementary projection. The first inequality applies the Sylvester
estimate to $\sym{X}$; the second uses $\sym{NormN}(BC)\leq\dkOperatorNorm{B}\sym{NormN}(C)$ and
$\dkOperatorNorm{\like{F1}{\symglyph{F1}^*}}\leq1$. For unbounded operators the Sylvester
identity holds on the trial operator's domain, with a bounded residual.
These are the geometric, spectral, and norm estimates formalized in our
proof of the theorem.

The tangent and double-angle estimates use related projection blocks.
For example, an estimate for $\sin(2\sym{Theta})$ allows angles near zero and
near $\sym{PiConstant}/2$. The spectral-selection results \citep[Section~8]{davis1970rotation}
identify which spectral subspaces give the small-angle alternative.
The direct-rotation analysis \citep[Section~4]{davis1970rotation}
studies the orthogonal or unitary map carrying one subspace onto another
and asks which displacement norms the direct rotation minimizes.

\subsection{Numerical and statistical applications}
\label{sec:related-numerical}

The 1970 theory influenced more than the analysis of exact
self-adjoint eigenvectors.  Wedin's 1972 singular-subspace
perturbation bounds address the corresponding problem for singular
vectors of general rectangular matrices
\citep{wedin1972perturbation}.  Bj\"orck and Golub's angle algorithms
made the subspace quantities computable
\citep{bjorck1973numerical}.  Stewart's work on pseudoinverses,
projections, and least-squares sensitivity helped connect these
geometric notions with matrix computations
\citep{stewart1977perturbation,stewart1990matrix}.
Paige and Wei's 1994 history shows how the cosine--sine decomposition
links the principal-angle geometry of Davis and Kahan with numerical
matrix factorizations \citep{paige1994history}.

The same subspace-versus-eigenvector distinction remains central in
statistical spectral estimation.  For example, Yu, Wang, and
Samworth give a statistical variant whose eigenspace error bound uses
separation in the population spectrum \citep{yu2015davis}, while Chen and
collaborators survey spectral methods across statistical learning
problems \citep{chen2021spectral}.  In that setting, an empirical covariance matrix approximates a population
operator and the perturbation norm controls the error of estimated
eigenspaces. Our general theorem supplies the deterministic angle estimate;
using it in a statistical result also requires the corresponding probabilistic
control of the perturbation.

\subsection{Foundations used in our formalization}
\label{sec:related-formal}

Each historical component above corresponds to a specific requirement
in the formalization.  Jordan's angle geometry and Halmos's
projection blocks lead to positive principal-angle operators rather
than chosen eigenbases.  The Sz.-Nagy--Kato intertwiner requires
polar decomposition, square roots, and real/complex transfer.
Spectral projections require Borel functional calculus, including for
unbounded self-adjoint operators.  The Sylvester proof requires a
representation carrying operator domains and a controlled passage
through spectral cutoffs.  Schatten and Ky Fan theory requires
approximation numbers and symmetric operator ideals to express
inequalities with the source's norm generality.

We formalized principal-angle geometry, symmetric operator ideals, and
unbounded Sylvester estimates using Lean~4 and Mathlib
\citep{demoura2021lean4,mathlib2020}. Mathlib supplies the algebraic,
topological, and Hilbert-space framework. Its continuous functional calculus
is described in the preprint by Anatole Dedecker and Jireh Loreaux
(\href{https://arxiv.org/abs/2501.15639}{arXiv:2501.15639}).
Our projection-valued-measure and Borel-calculus
work also adapts Adam Bornemann's Spectra formalization
\citep{bornemann2026spectra}.

Our Courant--Fischer and Weyl development, rectangular unitarily invariant
seminorms, projection-valued measures, and supporting polar-decomposition
results were subsequently adapted into Tau Ceti \citep{tauceti2026}.
% Editorial note (not manuscript prose): Do not add a sentence emphasizing
% the current Tau Ceti canonical polar-decomposition interface here.
% This paragraph records work developed in DKPS and subsequently
% contributed upstream; current library interfaces are discussed elsewhere.
% Evidence: dev/tauceti/dkps-upstream-contribution-census.md, PRs 13599,
% 14241, 17527 (adaptations), and 16127/16450 (independent polar proofs).
% AUTHORING_NOTES.md indexes accepted-source and campaign evidence.
% Mathematical origin and formal proof provenance are separate attributions.

This history also separates the two mathematical questions that meet
in the project.  The sine-theta theorem estimates \emph{how far} an
approximate or perturbed spectral subspace lies from an exact one.
The direct-rotation problem asks \emph{which} ambient transformation
carries one given subspace onto another at minimum displacement cost.
Our proofs of the optimization results use the principal-angle,
projection, and operator-ideal foundations developed for the formalization.
The matching theorem has its own proof; the shared foundations explain the
relationship between the two parts of the work. The resulting counterexample
to the full-displacement claim \citep[Proposition~4.4]{davis1970rotation}
and the follow-on matching theorem are examined in
\cref{sec:prop44,sec:example-transport}.

\section{Davis--Kahan formalization and the sine-theta theorem}
\label{sec:coverage}

Our targets are the four headline theorems \citep[Section~2]{davis1970rotation} and every
named theorem, proposition, lemma, and corollary established by Davis and
Kahan in the 1970 paper. These \DKResultCount{} results include the sine and
tangent inequalities, the geometry and extremal properties of direct
rotations, bounded and unbounded Sylvester estimates, the generalized
perturbation bounds \citep[Section~6]{davis1970rotation}, and the spectral-selection
theorems \citep[Section~8]{davis1970rotation}. We prove \DKProvedCount{} results and refute the printed
full-displacement claim \citep[Proposition~4.4]{davis1970rotation}.
Appendix~\ref{app:dk-inventory} lists these targets.

We develop the sine-theta theorem and its proof below; the remaining source
results are summarized in Appendix~\ref{app:dk-inventory}, with hypotheses
and conclusions in Section~\ref{app:dk-result-summaries}. The main text also
explains the spectral-selection results needed to interpret the angle bounds.
Appendix~\ref{app:notation} collects the notation used here and in the
direct-rotation argument.

\subsection{The sine-theta theorem for self-adjoint operators}
\label{sec:sine-theta}

\subsubsection{Trial and exact subspaces}

We use $\sym{A}$ for the reference operator and
$\sym{Acurrent}=\sym{A}+\sym{H}$ for the perturbed ambient operator.
The symbol $\sym{Acurrent}$ corresponds to $\sym{A}+\sym{H}$ in the source.

For a unit trial vector $\sym{v}$ and an approximate eigenvalue $\sym{mu}$, the residual
is $\sym{Acurrent}\sym{v}-\sym{mu} \sym{v}$. An isometric coordinate map $\sym{E0}$ generalizes
this construction to a trial subspace: its range consists of the candidate
approximate vectors, and a self-adjoint operator $\sym{TrialOperator}$ represents their
approximate dynamics. For a finite-dimensional Rayleigh--Ritz choice,
$\sym{TrialOperator}=\like{E0}{\symglyph{E0}^*}\sym{Acurrent}\sym{E0}$. The theorem allows other self-adjoint trial
operators as well. We write $\sym{TrialOperator}$ for the general trial operator,
reserving $\sym{ReferenceBlock0}$ for the reference spectral block that
Davis and Kahan also call $A_0$.

Let $\sym{AmbientHilbertSpace}$ be the ambient Hilbert space.  Let $\sym{TrialOperator}$ be a self-adjoint operator on a
trial-coordinate space $\sym{TrialCoordinateSpace}$, and let
\[
  \sym{E0}:\sym{TrialCoordinateSpace}\to \sym{AmbientHilbertSpace}
\]
be an isometric coordinate map for the trial subspace.  The residual compares
applying the ambient operator after embedding with applying the trial operator
before embedding:
\begin{equation}
  \sym{R} = \sym{Acurrent}\sym{E0}-\sym{E0}\sym{TrialOperator}.
  \label{eq:residual}
\end{equation}
When the operators are unbounded, \cref{eq:residual} is asserted on
$\DKDom(\sym{TrialOperator})$ and $\sym{R}$ is the bounded residual that agrees there.
The Lean predicate \texttt{IsTrialResidual} packages the isometry, the required
domain inclusion, and this residual identity.

Let $\sym{F0}:\sym{ExactCoordinateSpace}\to \sym{AmbientHilbertSpace}$ be isometric coordinates for the desired exact spectral
subspace and let $\sym{F1}:\sym{ComplementaryCoordinateSpace}\to \sym{AmbientHilbertSpace}$ be isometric coordinates for its orthogonal
complement.  A self-adjoint operator $\sym{Lambda1}$ on $\sym{ComplementaryCoordinateSpace}$ describes the
complementary spectral block through
\[
  \sym{Acurrent}\sym{F1} = \sym{F1}\sym{Lambda1}
\]
on its domain. This intertwining identity says that $\sym{Lambda1}$ gives
the action of $\sym{Acurrent}$ in complementary coordinates. The predicate
\texttt{IsExactSpectralDecomposition} also requires the domain inclusion
$\sym{F1}(\DKDom(\sym{Lambda1}))\subseteq\DKDom(\sym{Acurrent})$
and states orthogonality and completeness through
\[
 \like{F0}{\symglyph{F0}^*}\sym{F1}=0,\qquad
 \sym{F0}\like{F0}{\symglyph{F0}^*}+\sym{F1}\like{F1}{\symglyph{F1}^*}=\sym{Identity}.
\]

\subsubsection{The sine operator and its rectangular representative}

Here $\sym{Theta0}$ acts on $\sym{TrialCoordinateSpace}$; under the isometry
$\sym{E0}$ it represents the directed angle operator on the trial subspace
introduced in \cref{sec:related-two-projections}. If
$\sym{Pexact}=\sym{F0}\like{F0}{\symglyph{F0}^*}$ is the orthogonal projection onto the desired exact subspace, then
\[
  \cos\sym{Theta0} = (\like{E0}{\symglyph{E0}^*}\sym{Pexact}\sym{E0})^{1/2},
  \qquad
  \sym{sinTheta0} = (\sym{Identity}-\like{E0}{\symglyph{E0}^*}\sym{Pexact}\sym{E0})^{1/2}.
\]
The analytic proof uses the rectangular complementary block
\begin{equation}
  \sym{S}=(\sym{Identity}-\sym{Pexact})\sym{E0}.
  \label{eq:rectangular-sine}
\end{equation}
The theorem uses the positive sine on the trial-coordinate space.  The rectangular block has codomain $\sym{AmbientHilbertSpace}$, while the
positive sine has codomain $\sym{TrialCoordinateSpace}$.  Their relationship is
\[
  \sym{S}^*\sym{S}
  = \like{E0}{\symglyph{E0}^*}(\sym{Identity}-\sym{Pexact})\sym{E0}
  = \sym{Identity}-\like{E0}{\symglyph{E0}^*}\sym{Pexact}\sym{E0},
\]
so
\begin{equation}
  |\sym{S}|=(\sym{S}^*\sym{S})^{1/2}=\sym{sinTheta0}.
  \label{eq:modulus-sine}
\end{equation}
% It might be worth noting that our formalization actually defines \sym{Theta0}
% and proves eq:modulus-sine whereas most treatments of the subject seam to
% just start from |\sym{S}|. Maybe it's worth enough to put in the paper, but
% it at least deserves a latex comment!
The polar decomposition $\sym{S}=\sym{PolarU}|\sym{S}|$ separates the positive
  modulus from a partial isometry.
Unitarily invariant ideal norms assign the same value to an operator and its
  modulus, so \[ \sym{NormN}(\sym{S})=\sym{NormN}(\sym{sinTheta0}) \] whenever
  the norm is finite.
This equality transfers an estimate for the complementary block $\sym{S}$ to the
  positive sine operator.

\subsubsection{Gap geometry and the inequality}
\label{sec:gap-geometry}

Fix $\sym{delta}>0$. There are two geometric forms of separation.
\emph{Interval/exterior separation:} for some real
$\sym{betaEndpoint}\leq\sym{alphaEndpoint}$, one of the two spectra
$\sym{SpectrumSymbol}(\sym{TrialOperator})$ and $\sym{SpectrumSymbol}(\sym{Lambda1})$ lies in $[\sym{betaEndpoint},\sym{alphaEndpoint}]$ and the other
lies in
\[
 (-\infty,\sym{betaEndpoint}-\sym{delta}]\;\cup\;[\sym{alphaEndpoint}+\sym{delta},\infty).
\]
\emph{Ordered separation:} for some real $b$, the operator forms satisfy either
\[
 \sym{TrialOperator}\geq (b+\sym{delta})\sym{Identity},\quad \sym{Lambda1}\leq bI,
 \qquad\text{or}\qquad
 \sym{TrialOperator}\leq bI,\quad \sym{Lambda1}\geq (b+\sym{delta})\sym{Identity}.
\]
Each form inequality is asserted on the domain of the corresponding operator;
for example, $\sym{TrialOperator}\geq aI$ means
$\langle \sym{TrialOperator}x,x\rangle\geq a\|x\|^2$ for every
$x\in\DKDom(\sym{TrialOperator})$. For self-adjoint operators, the source's ordered
hypotheses on half-infinite spectral intervals, such as $(-\infty,b]$ and
$[b+\sym{delta},\infty)$, imply these form bounds. These cases are included
in the source's unbounded sine theorem
\citep[Section~6, discussion of unbounded operators]{davis1970rotation}.
Write $\sym{SineGap}(\sym{TrialOperator},\sym{Lambda1})$ for their union.
The Lean predicate \texttt{FormBoundedSylvesterGap} has three cases:
the finite-interval alternative, which includes either spectral order,
and the two ordered form alternatives.

Both blocks may be unbounded in an ordered configuration. For instance,
on $\ell^2(\mathbb N_{\geq1})$, take diagonal operators
$\sym{TrialOperator}e_n=(n+\sym{delta})e_n$ and $\sym{Lambda1}e_n=-ne_n$, on the domains where their
images are square summable. They satisfy the first ordered configuration
with $b=0$, while neither spectrum lies in a finite interval. This example
explains why retaining the half-infinite alternatives changes the theorem's
applicability.

For the unbounded trial configuration, the source assumes that
$\sym{Acurrent}\sym{E0}$ and $\sym{E0}\sym{TrialOperator}$ have a common dense domain.
Since $\sym{E0}$ is everywhere defined, this gives
\[
 \DKDom(\sym{Acurrent}\sym{E0})=\DKDom(\sym{TrialOperator}),
 \qquad
 \sym{E0}(\DKDom(\sym{TrialOperator}))\subseteq\DKDom(\sym{Acurrent}).
\]
The residual agrees with the difference on this domain and extends to a
bounded operator. The formal predicate \texttt{IsTrialResidual} includes the
domain inclusion and the residual identity.

The norm $\sym{NormN}$ is a normalized symmetric operator-ideal norm.
Normalization means that the rank-one map $uv^*:x\mapsto u\langle v,x\rangle$
has norm
\[
 \sym{NormN}(uv^*)=\|u\|\,\|v\|;
\]
in particular, its norm is one when $u,v$ are unit vectors.
Its hypotheses also include Ky Fan dominance whenever both ideal norms
are finite: domination of all finite singular-value sums implies domination
of the ideal norms. This comparison does not assert ideal membership.
The norm is invariant under isometries and positive modulus, and satisfies
$\sym{NormN}(BT)\leq\dkOperatorNorm{B}\sym{NormN}(T)$. Trace-class and
Hilbert--Schmidt operators with their respective norms are examples; the
operator norm on bounded operators is another. For an ideal norm, both the
residual and the positive sine must belong to its domain.

In two dimensions, take a one-dimensional trial frame $\sym{E0}(x)=x\sym{v}$ with
$\|\sym{v}\|=1$. A scalar trial operator $\sym{TrialOperator}=\sym{rho}$ gives
$\sym{r}=\sym{R}(1)=\sym{Acurrent}\sym{v}-\sym{rho}\sym{v}$.
Choosing $\sym{rho}=\langle\sym{v},\sym{Acurrent}\sym{v}\rangle$
makes $\sym{rho}\sym{v}$ the orthogonal projection of
$\sym{Acurrent}\sym{v}$ onto the trial line. The residual is therefore
perpendicular to that line. Figure~\ref{fig:sine-theta-residual} pairs this
residual triangle with the exact complementary projection that measures
the subspace error.

\begin{figure}[tb]
\centering
\resizebox{119mm}{!}{\input{figures/Fig1.pgf}}
\caption{Rayleigh residual geometry for a one-dimensional trial space in
$\sym{RealField}^2$ (the rank-one case). With $\sym{E0}(x)=x\sym{v}$, the residual
operator satisfies $\sym{R}(x)=x\sym{r}$, where
$\sym{r}=\sym{Acurrent}\sym{v}-\sym{rho}\sym{v}$.
For a higher-dimensional trial space,
$\sym{R}=\sym{Acurrent}\sym{E0}-\sym{E0}\sym{TrialOperator}$ assigns a residual vector
to each trial-coordinate vector. In a finite-dimensional eigenbasis of
$\sym{TrialOperator}$, its columns are the individual eigenpair residuals.
The exact frames in the drawing are
$\sym{F0}(h)=(h,0)$ and $\sym{F1}(g)=(0,g)$, with
$\sym{Pexact}=\sym{F0}\like{F0}{\symglyph{F0}^*}$ and $\sym{PexactPerp}=\sym{Identity}-\sym{Pexact}$.
(a) The pink leg $\sym{S}(1)=\sym{PexactPerp}\sym{v}$ has length $\sin\sym{theta}$.
(b) The green residual joins $\sym{rho}\sym{v}$ to $\sym{Acurrent}\sym{v}$;
for the Rayleigh choice it is perpendicular to the trial line. The dashed
legs are its exact wanted and unwanted components. Their lengths give
$\sym{delta}\sin\sym{theta}=\|\sym{PexactPerp}\sym{r}\|_2
\leq\dkOperatorNorm{\sym{r}\|_2=\|\sym{R}}$.
The drawing uses $\sym{Acurrent}=\sym{DiagonalOperator}(1.9,0.7)$,
$\sym{theta}=40^\circ$, $\sym{rho}=\langle\sym{v},\sym{Acurrent}\sym{v}\rangle$,
and $\sym{delta}=\sym{rho}-0.7$}
\label{fig:sine-theta-residual}
\end{figure}
% Place the residual figure before the theorem and its Lean signature.
\FloatBarrier

Appendix~\ref{app:lean-guide} explains the general Lean syntax in the
following theorem signature; the geometric and spectral hypotheses have been
defined above.

\begin{theorem}[Sine-theta estimate]
\label{thm:sine-theta}
Let the scalar field be $\sym{RealField}$ or $\sym{ComplexField}$, let $\sym{AmbientHilbertSpace}$ be separable,
and let $\sym{Acurrent}$, $\sym{TrialOperator}$, and $\sym{Lambda1}$ be self-adjoint operators on the
Hilbert spaces above. Suppose the isometries, domain inclusions, residual
identity, exact complementary decomposition, and the separation condition
$\sym{SineGap}(\sym{TrialOperator},\sym{Lambda1})$ of
Section~\ref{sec:gap-geometry} hold. Let $\sym{NormN}$ be a normalized symmetric
operator-ideal norm. Whenever $\sym{sinTheta0}$ and $\sym{R}$ belong to its domain,
\begin{equation}
 \sym{delta}\,\sym{NormN}(\sym{sinTheta0})\leq \sym{NormN}(\sym{R}).
 \label{eq:dk-sine}
\end{equation}
\end{theorem}

% Keep the Lean signature and its caption together across page breaks.
\noindent\begin{minipage}{\linewidth}
\begin{lstlisting}[style=leanpaper,
  caption={Lean signature for the sine-theta estimate. The surrounding
  scalar-field and Hilbert-space context is omitted.
  This is a readable ASCII transcription of the Lean declaration,
  rather than a verbatim copy of its Unicode binders and operator arrows.
  The paper-to-Lean names are registered centrally and displayed in
  Table~\ref{tab:paper-lean-notation}. In particular, the Lean binder
  displayed here as \texttt{A0} names our $\sym{TrialOperator}$, not the
  distinct reference block $\sym{ReferenceBlock0}$.
  The positive trial-coordinate sine is
  \texttt{sourceDirectedSinThetaOperator E0 F0};
  \texttt{FormBoundedSylvesterGap} encodes the interval/exterior and ordered
  half-infinite separation hypotheses described above.},
  label={lst:dk-sintheta-headline}]
theorem sinTheta_unbounded_formGap_whereDefinedUIN_rclike
    [TopologicalSpace.SeparableSpace E]
    (N : NormalizedSymmetricOperatorIdealFamily K)
    (A : LinearPMap K E E)
    (A0 : LinearPMap K F F)
    (Lambda1 : LinearPMap K G G)
    (E0 : ContinuousLinearMap K F E)
    (F0 : ContinuousLinearMap K H E)
    (F1 : ContinuousLinearMap K G E)
    (R : ContinuousLinearMap K F E)
    (hA : IsSelfAdjoint A)
    (hA0 : IsSelfAdjoint A0)
    (hLambda1 : IsSelfAdjoint Lambda1)
    (htrial : IsTrialResidual A A0 E0 R)
    (hexact : IsExactSpectralDecomposition A Lambda1 F0 F1)
    (delta : Real) (hdelta : 0 < delta)
    (hgap : FormBoundedSylvesterGap A0 Lambda1 delta) :
    N.Mem (sourceDirectedSinThetaOperator E0 F0) ->
    N.Mem R ->
      delta * N.gaugeReal (sourceDirectedSinThetaOperator E0 F0) <=
        N.gaugeReal R
\end{lstlisting}
\end{minipage}\par

The two ideal-membership hypotheses ensure that the norm values in the
conclusion are finite. The gap predicate includes the finite
interval/exterior case and both ordered cases. In the ordered cases,
$\sym{TrialOperator}$ and $\sym{Lambda1}$ may both be unbounded.

\subsubsection{Why the Sylvester estimate gives the angle bound}

Set $\sym{X}=\like{F1}{\symglyph{F1}^*}\sym{E0}$. The exact complementary decomposition gives
$\sym{S}=\sym{F1}\sym{X}$, and multiplication by the isometry $\sym{F1}$ preserves the
singular-value gauges. Self-adjointness, the domain hypotheses, and the
residual identity yield the Sylvester relation
\[
 \sym{Lambda1} \sym{X}-\sym{X}\sym{TrialOperator}=\like{F1}{\symglyph{F1}^*}\sym{R}
\]
on $\DKDom(\sym{TrialOperator})$, with the left side represented by the bounded
right side. The separated Sylvester estimate gives
\[
 \sym{delta} \sym{NormN}(\sym{X})\leq \sym{NormN}(\like{F1}{\symglyph{F1}^*}\sym{R})\leq \sym{NormN}(\sym{R}).
\]
Combining this with $\sym{NormN}(\sym{X})=\sym{NormN}(\sym{S})=\sym{NormN}(\sym{sinTheta0})$ proves the conclusion.
For unbounded operators, deriving the domain compatibility for $\sym{X}$ is part of
the argument. The displayed relation holds on $\DKDom(\sym{TrialOperator})$ and
has a bounded right side. The specified interval/exterior or ordered
spectral positions give the factor $1/\sym{delta}$, with no additional
multiplicative constant.

For a one-dimensional trial space, take
$\sym{Acurrent}=\sym{DiagonalOperator}(0,\sym{delta})$, the trial scalar
$\sym{TrialOperator}=0$, and the unit vector
$\like{v}{\symglyph{v}_{\symglyph{theta}}}=(\cos\sym{theta},\sin\sym{theta})$.
Choose the desired exact subspace $\sym{SpanOperator}\{(1,0)\}$;
its orthogonal complement is $\sym{SpanOperator}\{(0,1)\}$, on which
$\sym{Lambda1}=\sym{delta}$.
The complementary component is $(0,\sin\sym{theta})$ and the
residual is $(0,\sym{delta}\sin\sym{theta})$. Thus equality holds
in the sine-theta estimate for this finite-dimensional configuration.

\section{Formalized operator-theoretic foundations}
\label{sec:foundations}

We formalized the geometry of two subspaces, symmetric operator ideals and
majorization, and Sylvester estimates for unbounded self-adjoint operators.
These results supply the constructions and inequalities used in the
Davis--Kahan proofs.

\subsection{Direct rotations and two-subspace geometry}

Let $\sym{Psource}$ and $\sym{Qtarget}$ project onto closed subspaces
$\sym{Usource}$ and $\sym{Vtarget}$. For an acute pair,
$\|\sym{Psource}-\sym{Qtarget}\|<1$. The compression
$\sym{Psource}\sym{Qtarget}\sym{Psource}$ on $\sym{Usource}$ determines
the positive cosine operator. Its square root and the complementary
compression give the cosine and sine factors in the direct-rotation
construction. Polar decomposition gives the isometries between the
corresponding parts of the two subspaces.

The resulting orthogonal or unitary map $\sym{Rrotation}$ carries
$\sym{Usource}$ onto $\sym{Vtarget}$. In finite-dimensional principal-plane
coordinates it rotates each plane through its principal angle and fixes
common directions. The projection formula from
\cref{eq:history-kato-unitary} defines it in Hilbert space without requiring
a discrete list of principal vectors.

Our formalization proves acute existence and uniqueness, the
principal-square-root characterization, commutation with the angle
operators, and reversal symmetry: interchanging the two subspaces replaces
the rotation by its adjoint. For general pairs, the existence criteria
include the right-angle defect spaces described in
\cref{sec:related-two-projections}. When the angle operator is compact,
its spectral decomposition recovers a sequence of principal angles
\citep[Section~3]{davis1970rotation}.

The double-angle theorems for $\sin(2\sym{Theta})$ and $\tan(2\sym{Theta})$
  require selecting the branch on which $\sym{Theta}\le\sym{PiConstant}/4$.
Here $\sym{Theta}$ includes both the source and complementary angle blocks;
$\sym{Theta0}$ alone is the directed angle on source coordinates.
Because the function $\sin(2\sym{theta})$ is small both when $\sym{theta}$ is
  near $0$ and when $\sym{theta}$ is near $\sym{PiConstant}/2$, a bound on
  $\sin(2\sym{Theta})$ alone cannot determine whether an approximate subspace is
  close to or nearly orthogonal to the target.
The branch-selection theorem \citep[Theorem~8.1]{davis1970rotation}
uses an ordered reference gap. Let $\sym{Psource}$ reduce a self-adjoint
reference operator $\sym{A}$, with
blocks $\sym{ReferenceBlock0}$ on $\DKRange(\sym{Psource})$ and
$\sym{ReferenceBlock1}$ on $\DKKernel(\sym{Psource})$, satisfying
\[
 \sym{ReferenceBlock0}\le\sym{alphaEndpoint}\sym{Identity},
 \qquad
 \sym{ReferenceBlock1}\ge(\sym{alphaEndpoint}+\sym{delta})\sym{Identity},
 \qquad \sym{delta}>0,
\]
in the quadratic-form sense. Suppose $\sym{H}$ is bounded and self-adjoint,
with both diagonal blocks zero:
$\sym{Psource}\sym{H}\sym{Psource}
 =(\sym{Identity}-\sym{Psource})\sym{H}(\sym{Identity}-\sym{Psource})=0$.
These hypotheses imply that $\sym{Acurrent}=\sym{A}+\sym{H}$ has no spectrum
in $(\sym{alphaEndpoint},\sym{alphaEndpoint}+\sym{delta})$.

Fix a reducing subspace $M$ of $\sym{Acurrent}$, and write
$\sym{Lambda0},\sym{Lambda1}$ for the blocks of $\sym{Acurrent}$ on $M,M^\perp$.
The full angle $\sym{Theta}$ comparing $\DKRange(\sym{Psource})$ with $M$
satisfies $\sym{Theta}\le\sym{PiConstant}/4$ if and only if
$\sym{Lambda0}\le\sym{alphaEndpoint}\sym{Identity}$ and
$\sym{Lambda1}\ge(\sym{alphaEndpoint}+\sym{delta})\sym{Identity}$.
Such a subspace exists: take the spectral subspace of $\sym{Acurrent}$
for $(-\infty,\sym{alphaEndpoint}]$. Our formalization also proves that
it is the unique reducing subspace satisfying this angle bound.

To compare the reference and perturbed blocks on common coordinate spaces,
choose isometric frames for $\DKRange(\sym{Psource})$ and
$\DKKernel(\sym{Psource})$. Apply the direct rotation $\sym{Rrotation}$ to
these frames to obtain frames for $M$ and $M^\perp$. Express
$\sym{ReferenceBlock0},\sym{ReferenceBlock1}$ in the original frames and
$\sym{Lambda0},\sym{Lambda1}$ in the transported frames. In these coordinates
the branch satisfies the stronger quadratic-form inequality
\[
 \sym{ReferenceBlock1}-\sym{alphaEndpoint}\sym{Identity}
 \le\sym{ComplementCosineBlock}
       (\sym{Lambda1}-\sym{alphaEndpoint}\sym{Identity})\sym{ComplementCosineBlock},
 \qquad \sym{ComplementCosineBlock}=\cos\sym{Theta1},
\]
with the lower-block analogue. These inequalities quantify spectral
repulsion in terms of the angles and yield eigenvalue and symmetric-gauge
comparisons in finite dimensions.

The smallness criterion \citep[Theorem~8.2]{davis1970rotation} uses a central
spectral interval instead of an ordered gap. Let $\sym{Psource}$ reduce
the self-adjoint operator $\sym{A}$ and $\sym{Qtarget}$ reduce the
self-adjoint operator $\sym{Acurrent}$. The perturbed block on
$\DKRange(\sym{Qtarget})$ has spectrum in
$[\sym{betaEndpoint},\sym{alphaEndpoint}]$, while its complementary spectrum
lies in
\[
 (-\infty,\sym{betaEndpoint}-\sym{delta}]
 \ \cup\ [\sym{alphaEndpoint}+\sym{delta},\infty),
 \qquad \sym{betaEndpoint}\le\sym{alphaEndpoint},\quad\sym{delta}>0.
\]
The reference block $\sym{ReferenceBlock0}$, representing
$\sym{A}|_{\DKRange(\sym{Psource})}$ in source coordinates, must have spectrum
in $[\sym{betaEndpoint}-\sym{delta}/2,\sym{alphaEndpoint}+\sym{delta}/2]$.
For an isometric frame $\sym{E0}$ of the reference subspace, its residual
is $\sym{R}=\sym{Acurrent}\sym{E0}-\sym{E0}\sym{ReferenceBlock0}$
on the operator domain, as in \cref{eq:history-residual}.
Assume
$\sym{DimensionOperator}\sym{CrossedSourceDefect}=\sym{DimensionOperator}\sym{CrossedTargetDefect}$: the space of source
vectors orthogonal to the target and the space of target vectors orthogonal
to the source have equal Hilbert dimension.

If $\dkOperatorNorm{\sym{H}}<\sym{delta}/2$ for a
bounded self-adjoint perturbation, or $\dkOperatorNorm{\sym{R}}<\sym{delta}/2$
for the bounded residual on the reference subspace, then
$\sym{Theta}<\sym{PiConstant}/4$. For each ideal norm whose displayed arguments belong
to its domain, the corresponding double-angle estimate also holds:
\[
 \sym{delta}\sym{NormN}(\sin(2\sym{Theta}))\le2\sym{NormN}(\sym{H}),
 \qquad\text{or}\qquad
 \sym{delta}\sym{NormN}(\sin(2\sym{Theta0}))\le2\sym{NormN}(\sym{R}).
\]
% Note (not for print): the factor 2 in the residual form is correct and
% matches the 1970 Section 2 statement. The 1970 Section 7 proof, after (7.5),
% rewrites delta||sin 2Theta|| <= ||H - XHX|| with right side ||[[0,B*],[B,0]]||
% and concludes delta||sin 2Theta_0|| <= ||B|| <= ||R||; since H - XHX has
% off-diagonal block 2B, the right side should be 2||B|| <= 2||R||. The
% Section 7 example (A = diag(0,d), H = [[0,1],[1,-d]], delta = 2,
% theta = pi/4, ||R|| = 1) attains delta sin 2theta = 2 = 2||R||. Proof-line
% slip only; the stated theorem is unaffected.
The perturbation estimate uses the full angle, and the residual estimate
the directed angle on source coordinates. In the unbounded residual form,
$\sym{A}$ and $\sym{Acurrent}$ are self-adjoint on a common domain; only
the residual is required to be bounded.

We formalized both theorems in Hilbert space using spectral calculus.

\subsection{Approximation numbers and operator-ideal norms}

In finite dimensions, a unitarily invariant norm is a symmetric gauge
function of the singular values. Approximation numbers extend the
singular-value comparisons needed for bounded operators; for compact
operators they give the usual singular-value sequence. We formalized
symmetric norming functions and Ky Fan partial sums to pass from
majorization to norm inequalities.

For example, let $\sym{Rrotation}$ be the direct rotation and
$\sym{Wcompetitor}$ any unitary carrying $\sym{Usource}$ onto
$\sym{Vtarget}$. The restricted displacement of the direct rotation has
singular-value minimality, which implies minimality for unitarily invariant
norms. The squared-displacement result compares
$(\sym{Identity}-\sym{Rrotation})^*(\sym{Identity}-\sym{Rrotation})$ and
$(\sym{Identity}-\sym{Wcompetitor})^*(\sym{Identity}-\sym{Wcompetitor})$
\citep[Proposition~4.3]{davis1970rotation}. Their singular values are the
squares of those of the unsquared displacements. Section~\ref{sec:prop44}
shows why the claimed trace-norm comparison fails and derives the valid
Q-norm conclusion.
% Additional formalized quantitative results, omitted from this article:
% DavisKahan/FiniteDimensional/DirectRotation/SharpDisplacement.lean and
% SharpDisplacementLimit.lean. AUTHORING_NOTES.md indexes the evidence.

The same norm inequalities enter the sine-theta proof. The assertion
$\sym{delta}\sym{NormN}(\sym{S})\leq\sym{NormN}(\sym{R})$ applies when
$\sym{S}$ and $\sym{R}$ have finite ideal norm. Our ideal-membership
hypotheses state this condition explicitly in arbitrary Hilbert dimension.

\subsection{Sylvester estimates and unbounded operators}

For operators $\sym{SylvesterLeft}$ and $\sym{SylvesterRight}$ on the target and source spaces of $\sym{X}$, the
Sylvester map is $\sym{X}\mapsto \sym{SylvesterLeft}\sym{X}-\sym{X}\sym{SylvesterRight}$. Under the interval/exterior or ordered
spectral-separation hypotheses, the estimate has the form
\[
 \sym{delta} \sym{NormN}(\sym{X})\leq \sym{NormN}(\sym{SylvesterLeft}\sym{X}-\sym{X}\sym{SylvesterRight}).
\]
The sine-theta proof uses $\sym{SylvesterLeft}=\sym{Lambda1}$, $\sym{SylvesterRight}=\sym{TrialOperator}$, and
$\sym{X}=\like{F1}{\symglyph{F1}^*}\sym{E0}$. Other perturbation estimates choose related
projection blocks for $\sym{X}$.

For unbounded $\sym{SylvesterLeft}$ and $\sym{SylvesterRight}$, the identity is asserted for
$x\in\DKDom(\sym{SylvesterRight})$, with $\sym{X}x\in\DKDom(\sym{SylvesterLeft})$, and $\sym{SylvesterLeft}\sym{X}-\sym{X}\sym{SylvesterRight}$ agrees there with a
bounded residual. We formalized these domain conditions using linear
partial maps. Spectral cutoffs give bounded approximations; convergence
of approximation numbers passes the norm estimates to the unbounded
operators \citep[Section~5 and Appendix to Section~6]{davis1970rotation}.

The generalized sine and tangent theorems also compare block operators.
Reflection pinching replaces a block matrix $T$ by
$(T+\sym{Jreflect}T\sym{Jreflect})/2$, where $\sym{Jreflect}$ acts as $+\sym{Identity}$ on one coordinate subspace and $-\sym{Identity}$
on the other. This removes the off-diagonal blocks and contracts a unitarily
invariant norm. Direct-sum comparisons relate the norms of the component
blocks to the whole operator. Together with Sylvester separation these give
the generalized angle estimates
\citep[Lemmas~6.1--6.2 and Theorems~6.1--6.3]{davis1970rotation}.

\subsection{Real and complex Hilbert spaces}

The spectral constructions use complex functional calculus. We formalized
the scalar-transfer results needed to carry these arguments to real Hilbert
spaces through complexification. This permits the source theorems to be
stated over either field, including their domain and ideal-norm hypotheses.
The underlying algebra, topology, and continuous functional calculus come
from Mathlib; the adapted Spectra constructions and later Tau Ceti
contributions are described in Section~\ref{sec:related-formal}.

\input{counterexample.tex}

\input{formalization_enabled_extension.tex}

\section{Discussion}
\label{sec:discussion}

The formalization covers the original perturbation theory on real and
complex Hilbert spaces, including unbounded self-adjoint operators and
operator-ideal norms. The unbounded statements require explicit domain
conditions and bounded residuals. The operator-ideal statements require
finite norm values. These hypotheses determine which approximation problems
the estimates address.

The correction of the full-displacement claim
\citep[Proposition~4.4]{davis1970rotation} separates two extremal statements.
The valid comparison \citep[Proposition~4.3]{davis1970rotation}
gives squared-displacement minimality; its Q-norm consequence
is valid. Our real four-dimensional family exhibits full trace-norm failure
at every prescribed positive acute angle. In each example, the competitor has
two nonzero displacement singular values while the direct rotation has four.
The sums and squared sums of these values account for their different norm
comparisons. The earlier complex example supplied the construction; our
calculation makes its contradiction of the printed real claim explicit.

The matching theorem addresses the resulting optimization problem for any
equal-dimensional pair of finite-dimensional real subspaces. Its proof uses
principal-angle geometry, real orthogonal block decompositions, and
singular-value norms. These are also among the foundations required for the
Davis--Kahan formalization. The proof of the matching theorem uses additional
optimization arguments, including a passage to the trace-norm endpoint.

\section{Conclusion}
\label{sec:conclusion}

We formalized 28 of the 29 selected Davis--Kahan results and refuted the
remaining printed statement \citep[Proposition~4.4]{davis1970rotation}.
The sine-theta theorem includes
unbounded self-adjoint operators and symmetric operator-ideal norms with
their domain and finiteness assumptions. The direct-rotation calculation
identifies the failed singular-value comparison and gives an explicit real
counterexample and a valid Q-norm consequence.

We also formalized the operator theory used in these results: principal
angles and direct rotations, Borel spectral calculus, approximation numbers
and symmetric ideals, and Sylvester separation. Parts of this development
were subsequently adapted into Tau Ceti. Using these foundations, we proved
an exact weighted-matching formula for the real trace-norm transport
minimum.

\backmatter
\section*{Statements and Declarations}

\subsection*{Funding}
\ifDKAnonymous
Funding information is omitted for anonymous review.
\else
This work was supported by the Defense Advanced Research Projects Agency
under Contract No. HR001125CE017.
\fi

\subsection*{Competing interests}
The authors declare no competing interests.

\subsection*{Data and code availability}
\ifDKAnonymous
Artifact information is omitted for anonymous review.
\else
The Lean proofs, result-to-Lean mapping, provenance records, figure-generation
sources, and manuscript sources are available at
\url{https://github.com/AIQ-Kitware/aiq-dkps-formalization}.
\fi

\subsection*{AI assistance}
ChatGPT and Claude assisted with the formalization and manuscript preparation.
Formalizations were carried out by large language models with human oversight.
The matching result was derived using GPT-6.1 Sol and reviewed by the human
  authors.

\bibliography{references}

\clearpage
\appendix

\input{appendices.tex}

\end{document}
, without editing.
% EVERY \sym, \like, and \dkrole obeys this switch, including the glossary.
% No color is required to identify a mathematical object: the glyphs, indices,
% definitions and adjacent prose also distinguish them.
%
% Families (base hue): exact/target=blue, trial/source=amber, exact unwanted=
% cyan, angles and their operators=magenta, residual/displacement=green,
% gaps=violet. Norm gauges=slate; provenance retains brown/steel/vermilion.
% P_ex is the exact-target projection in sine-theta; P is the source
% projection in transport. R denotes a residual and D a direct rotation.
% Semantic keys retain their roles even when the printed notation changes.
\newif\ifDKMathColors
\DKMathColorstrue
\ifdefined\DKMonochromeMath
  \DKMathColorsfalse
\fi

% Family anchors (print-compatible versions of the Manim light palette).
\definecolor{dkfg}{HTML}{1B1F24}
\definecolor{dkmuted}{HTML}{6B7280}
\definecolor{dkwanted}{HTML}{1F6FD1}
\definecolor{dktrial}{HTML}{A66705}
\definecolor{dkunwanted}{HTML}{007C91}
\definecolor{dksine}{HTML}{D12A7B}
\definecolor{dkresid}{HTML}{087D43}
\definecolor{dkgap}{HTML}{6E4BD8}
\definecolor{dkold}{HTML}{705519}
\definecolor{dkcurrent}{HTML}{054669}
\definecolor{dkperturb}{HTML}{BB491B}
\definecolor{dkrefuted}{HTML}{C62828}

% Related yet distinguishable shades. These are NOT new unrelated categories:
% objects inherit their parent family; a different shade signals a different
% mathematical object (e.g., a frame versus its subspace or its projection).
\definecolor{dkwantedframe}{HTML}{1254A9}
\definecolor{dkwantedprojection}{HTML}{08618D}
\definecolor{dkwantedrotation}{HTML}{4248B0}
\definecolor{dkwantedspectrum}{HTML}{345F8F}
\definecolor{dkwantedvector}{HTML}{1468AA}
\definecolor{dktrialframe}{HTML}{905600}
\definecolor{dktrialprojection}{HTML}{79511B}
\definecolor{dktrialcompetitor}{HTML}{A04D10}
\definecolor{dktrialspectrum}{HTML}{79591C}
\definecolor{dktrialvector}{HTML}{9B522A}
\definecolor{dkunwantedframe}{HTML}{006477}
\definecolor{dkunwantedspectrum}{HTML}{00546B}
\definecolor{dkunwantedprojection}{HTML}{145B6D}
\definecolor{dkunwantedvector}{HTML}{256879}
\definecolor{dkangleoperator}{HTML}{923478}
\definecolor{dksineoperator}{HTML}{AA275A}
\definecolor{dksineblock}{HTML}{B43F70}
\definecolor{dksylvester}{HTML}{804074}
\definecolor{dkanglemax}{HTML}{B32185}
\definecolor{dkanglecos}{HTML}{824276}
\definecolor{dkresidvector}{HTML}{086038}
\definecolor{dkdisplacement}{HTML}{226C50}
\definecolor{dkgaporiginal}{HTML}{4E4D89}
\definecolor{dkgapstat}{HTML}{7039A7}
\definecolor{dkgapinterval}{HTML}{574A82}
\definecolor{dkambient}{HTML}{36516A}
\definecolor{dknorm}{HTML}{475875}
\definecolor{dknormaux}{HTML}{66526E}
\definecolor{dkintertwiner}{HTML}{4A6280}
\definecolor{dkturn}{HTML}{47647D}
\definecolor{dkmix}{HTML}{505B65}

% Coloring primitives. "plain" suppresses color without suppressing hovers.
% All descendants can retain a semantic role through \like{key}{expression}.
\expandafter\def\csname dknocolor@plain\endcsname{1}
\newcommand{\dkrole}[2]{%
  \ifDKMathColors
    \ifcsname dknocolor@#1\endcsname #2\else\textcolor{dk#1}{#2}\fi
  \else #2\fi
}
% An optional plain-text reminder supplies definitions and types on hover.
% Required fields: {semantic key}{printed glyph}{color role}{glossary|local}
%                  {short definition}{mathematical type}. The "glossary" flag records whether an
% entry appears in the curated printed notation table; "local" does not.
% This registry is the sole authority for symbol meaning and printed form.
% Use role "plain" for an uncolored symbol that still carries hover text.
% Changing the glyph is a separate editorial decision: this audit preserves
% the current E, F, F_ex, G, and other existing printed choices.
\newcommand{\dkDeclareSymbol}[7][]{%
  \expandafter\def\csname dkgloss@#2\endcsname{#6}%
  \expandafter\def\csname dkglossaryflag@#2\endcsname{#5}%
  \if\relax\detokenize{#1}\relax
    \expandafter\pdfstringdef\csname dkpdfgloss@#2\endcsname{#7: #6}%
  \else
    \expandafter\pdfstringdef\csname dkpdfgloss@#2\endcsname{#7: \detokenize{#1}}%
  \fi
  \expandafter\def\csname dktype@#2\endcsname{#7}%
  \expandafter\def\csname dkglyph@#2\endcsname{#3}%
  \expandafter\def\csname dkrole@#2\endcsname{#4}%
  \expandafter\def\csname dksym@#2\endcsname{\dkrole{#4}{#3}}%
}
% Unstyled printed glyph, for variants constructed through \like.
% Example: \like{E0}{\symglyph{E0}^*} derives E_0^* from E0's
% one canonical glyph; retyping E_0 here would defeat central control.
\DeclareRobustCommand{\symglyph}[1]{\csname dkglyph@#1\endcsname}
\DeclareRobustCommand{\sym}[1]{%
  \ifcsname dksym@#1\endcsname
    \dkSymbolTooltip{#1}{\csname dksym@#1\endcsname}%
  \else
    \PackageError{dk-math-colors}{Unknown symbol key #1}{See math_symbol_colors.tex}%
  \fi
}
\DeclareRobustCommand{\like}[2]{%
  \ifcsname dkrole@#1\endcsname
    \dkSymbolTooltip{#1}{\dkrole{\csname dkrole@#1\endcsname}{#2}}%
  \else
    \PackageError{dk-math-colors}{Unknown role key #1}{See math_symbol_colors.tex}%
  \fi
}

% A short description belongs to the semantic key, not the bare letter.
% The same descriptions are available to the notation appendix.
\DeclareRobustCommand{\dkSymbolDescription}[1]{\csname dkgloss@#1\endcsname}
\DeclareRobustCommand{\dkSymbolGlossaryStatus}[1]{\csname dkglossaryflag@#1\endcsname}
\DeclareRobustCommand{\dkSymbolType}[1]{\csname dktype@#1\endcsname}

% Citable alternative notation is metadata, not a request to rename glyphs.
% The source identifier, the source's exact printed glyph, and the specific
% source location are all mandatory. No alternative is inferred from a name.
% One key may have alternatives from several works or multiple passages.
\newcommand{\dkDeclareNotationSource}[3]{%
  \expandafter\def\csname dksourcetitle@#1\endcsname{#2}%
  \expandafter\def\csname dksourceurl@#1\endcsname{#3}%
}
\newcommand{\dkDavisKahanRows}{}
\makeatletter
\newcommand{\dkDeclareNotationAlternative}[5]{%
  \ifcsname dkglyph@#1\endcsname\else
    \PackageError{dk-math-colors}{Unknown alternative key #1}{Declare the symbol first}%
  \fi
  \ifcsname dksourcetitle@#2\endcsname\else
    \PackageError{dk-math-colors}{Unknown notation source #2}{Declare the source first}%
  \fi
  \expandafter\def\csname dkaltglyph@#2@#1\endcsname{#3}%
  \expandafter\def\csname dkaltlocator@#2@#1\endcsname{#4}%
  \expandafter\def\csname dkaltnote@#2@#1\endcsname{#5}%
  \def\dk@source{#2}%
  \def\dk@original{DK1970}%
  \ifx\dk@source\dk@original
    \g@addto@macro\dkDavisKahanRows{\dkDavisKahanRow{#1}}%
  \fi
}
\makeatother
% A correspondence row is produced from the SAME registry records used by
% \sym and by the source-census script. It is never spelled a second time.
\newcommand{\dkDavisKahanRow}[1]{%
  \csname dkaltnote@DK1970@#1\endcsname &
  $\csname dkaltglyph@DK1970@#1\endcsname$ &
  $\sym{#1}$ &
  \csname dkaltlocator@DK1970@#1\endcsname\\
}

\newif\ifDKSymbolTooltips
\DKSymbolTooltipstrue
\ifdefined\DKNoSymbolTooltips\DKSymbolTooltipsfalse\fi
\newsavebox{\DKTooltipBox}
\newcount\DKTooltipCount
% Register only shipped widgets through hyperref's existing form-field
% machinery. Poppler/Okular need the catalog's field list, not just page
% annotations. The deferred aux write also discards measuring-only copies.
\makeatletter
\newcommand{\dkRegisterTooltipField}{%
  \expandafter\HyField@@AddToFields\expandafter{\number\pdffeedback lastannot}%
}
\makeatother
\AtBeginDocument{\ifDKSymbolTooltips\Form[NeedAppearances=false]\fi}
\AtEndDocument{\ifDKSymbolTooltips\endForm\fi}
% Native LuaLaTeX PDF widgets: no navigation action or JavaScript. Avoid
% pdfcomment's soulpos dependency, which previously broke local Vim builds.
% Measure in the CURRENT math style, then emit the original expression as-is:
% boxing the visible math would change spacing, scripts and line breaking.
\newcommand{\dkSymbolTooltip}[2]{%
  \ifDKSymbolTooltips
    \begingroup
      \DKSymbolTooltipsfalse % suppress nested annotations while measuring
      \ifmmode
        \edef\DKTooltipStyle{\ifcase\mathstyle
          \displaystyle\or\displaystyle\or\textstyle\or\textstyle\or
          \scriptstyle\or\scriptstyle\or\scriptscriptstyle\or\scriptscriptstyle\fi}%
        \sbox{\DKTooltipBox}{$\DKTooltipStyle#2$}%
      \else
        \leavevmode\sbox{\DKTooltipBox}{#2}%
      \fi
      \global\advance\DKTooltipCount by 1
      \pdfextension annot width\wd\DKTooltipBox
        height\ht\DKTooltipBox depth\dp\DKTooltipBox
        {/Subtype /Widget /FT /Btn /Ff 65536 /F 512
         /T (dk-symbol-\the\DKTooltipCount) /TU (\csname dkpdfgloss@#1\endcsname)
         /H /N /BS << /W 0 >> /Border [0 0 0]}%
      \dkRegisterTooltipField
      #2%
    \endgroup
  \else
    #2%
  \fi
}
\pdfstringdefDisableCommands{%
  \def\sym#1{\csname dkglyph@#1\endcsname}%
  \def\like#1#2{#2}%
  \def\symglyph#1{\csname dkglyph@#1\endcsname}%
}

% Declarations: [detailed hover], key, glyph, role, glossary|local, short text, type.
% A local entry has full semantic/hover support, but no notation-table row.
% When two keys have the same glyph, it is an auditable collision, not an
% instruction to change the printed notation automatically.
% This registry includes local proof notation. The notation appendix selects
% the symbols useful for reference; registration does not add a printed row.
% Manim's original symbol keys, now split into semantic subshades.
% Ambient and perturbed operators (provenance is independent of geometry).
\dkDeclareSymbol[Reference self-adjoint operator A on the ambient Hilbert space; the perturbed operator is Ahat = A + H. In a single-operator eigenproblem, A itself is the operator being approximated. An unbounded A acts only on domain(A).]
  {A}{A}{old}{glossary}{Reference self-adjoint operator before perturbation.}{self-adjoint-operator}
\dkDeclareSymbol{A1}{A_1}{old}{local}{Complementary block of the reference operator.}{self-adjoint-operator}
\dkDeclareSymbol{a1}{a_1}{old}{local}{Reference eigenvalue or complementary spectral value.}{scalar}
\dkDeclareSymbol[Reference block A_0: the restriction of A to the range of the reducing source projection P, expressed in source coordinates. In Theorem 8.1 it lies at or below alpha; in Theorem 8.2 its spectrum lies in the enlarged central interval. The perturbed blocks are denoted Lambda_0 and Lambda_1.]
  {ReferenceBlock0}{A_0}{old}{glossary}{Reference operator restricted to the source subspace in Section 8.}{self-adjoint-operator}
\dkDeclareSymbol[Reference block A_1: the restriction of A to the kernel of the reducing source projection P, expressed in complementary coordinates. In Theorem 8.1 it lies at or above alpha + delta in quadratic-form order.]
  {ReferenceBlock1}{A_1}{old}{glossary}{Reference operator restricted to the source complement in Section 8.}{self-adjoint-operator}
\dkDeclareSymbol[Self-adjoint ambient operator Ahat being studied; in the perturbation setting Ahat = A + H. Its action on a trial vector is compared with the proposed eigenvalue times that vector.]
  {Ahat}{\widehat A}{current}{local}{Operator under study; in the perturbation setting, the reference operator plus its perturbation.}{self-adjoint-operator}
\dkDeclareSymbol[Self-adjoint ambient operator Ahat being studied. In the perturbation setting Ahat = A + H. The trial residual is R = Ahat E_0 - E_0 A_tr, and the exact complementary block satisfies Ahat F_1 = F_1 Lambda_1 on its domain.]
  {Acurrent}{\widehat A}{current}{glossary}{Operator under study; in the perturbation setting, the reference operator plus its perturbation.}{self-adjoint-operator}
\dkDeclareSymbol[Additive self-adjoint perturbation H = Ahat - A. A bound on its norm can control the residual of an eigenspace of one operator relative to the other.]
  {H}{H}{perturb}{glossary}{Additive self-adjoint perturbation of the reference operator.}{self-adjoint-operator}
\dkDeclareSymbol{H0}{H_0}{perturb}{local}{Wanted diagonal block of the perturbation.}{self-adjoint-operator}
\dkDeclareSymbol{H1}{H_1}{perturb}{local}{Unwanted diagonal block of the perturbation.}{self-adjoint-operator}
\dkDeclareSymbol{B}{B}{perturb}{local}{Off-diagonal perturbation or coupling block.}{linear-operator}
\dkDeclareSymbol{b}{b}{perturb}{local}{Scalar perturbation or coupling coefficient.}{scalar}
% Trial/source family: subspace, frame, compressed operator, projector, vector.
\dkDeclareSymbol{U}{U}{wanted}{local}{Desired exact subspace in the local geometric notation.}{linear-subspace}% Manim's U remains wanted in local slide notation.
\dkDeclareSymbol{V}{V}{trial}{local}{Trial subspace in the local geometric notation.}{linear-subspace}% Manim's V remains trial in local slide notation.
\dkDeclareSymbol[Source subspace U to be carried onto the target V by an orthogonal or unitary map W. In the transport sections P projects onto U and Q onto V. The matching theorem also names the source subspace K.]
  {Usource}{U}{trial}{glossary}{Source subspace to be transported onto the target subspace.}{linear-subspace}
\dkDeclareSymbol[Trial frame E_0: F -> E, a linear isometry from trial-coordinate space F into ambient space E. Its range is the trial subspace; E_0* returns trial coordinates and E_0* E_0 = I_F. In rank one, E_0(x) = x v.]
  {E0}{E_0}{trialframe}{glossary}{Isometric embedding of the trial subspace; its adjoint returns trial coordinates.}{isometric-embedding}
\dkDeclareSymbol[Self-adjoint trial operator A_tr on trial-coordinate space F. The residual compares Ahat E_0 with E_0 A_tr. Rayleigh--Ritz chooses A_tr = E_0* Ahat E_0 when this compression is defined; the theorem allows other trial operators.]
  {A0}{A_{\mathrm{tr}}}{trialspectrum}{glossary}{Self-adjoint operator on trial coordinates; its eigenvalues are the trial spectral values.}{self-adjoint-operator}
\dkDeclareSymbol[Trial vector: a proposed approximate eigenvector. In the rank-one geometry v has unit length and E_0(x) = x v. Its residual for trial scalar rho is r = Ahat v - rho v.]
  {v}{v}{trialvector}{local}{Trial vector or vector in the trial subspace.}{vector}
\dkDeclareSymbol[Scalar trial operator in rank one, A_tr(x) = rho x. Figure 1 uses the Rayleigh value rho = <v, Ahat v> for unit v; then rho v is the projection of Ahat v onto the trial line and r is perpendicular to v.]
  {rho}{\rho}{trialspectrum}{local}{Trial eigenvalue or scalar trial operator.}{scalar}
\dkDeclareSymbol[Approximate eigenvalue assigned to a trial vector v. For the local operator A, its residual is A v - mu v. For unit v, the Rayleigh--Ritz choice is mu = <A v, v>.]
  {mu}{\mu}{trialspectrum}{local}{Trial spectral value.}{scalar}
\dkDeclareSymbol{a0}{a_0}{trialspectrum}{local}{Wanted trial eigenvalue or trial spectral value.}{scalar}
\dkDeclareSymbol[Ambient orthogonal or unitary transporter W with W(U) = V, equivalently W P W* = Q. Its full displacement is I - W on the whole ambient space. The real matching theorem minimizes its trace norm over every such orthogonal W.]
  {Wcompetitor}{W}{trialcompetitor}{glossary}{Orthogonal or unitary competitor carrying the source subspace onto the target.}{unitary-operator}
\dkDeclareSymbol[Orthogonal projection P onto the source subspace U in the transport problem; P = P* = P squared. A feasible orthogonal or unitary W satisfies W P W* = Q. Restricted displacement (I - W)P tests only source vectors.]
  {Psource}{P}{trialprojection}{glossary}{Orthogonal projection onto the source subspace in the transport problem.}{orthogonal-projection}
\dkDeclareSymbol{ProjSource}{P_U}{trialprojection}{local}{Orthogonal projection onto the source subspace in the transport problem.}{orthogonal-projection}
\dkDeclareSymbol[Vector x in the source subspace U. For principal vectors, unit x_j in U and y_j in V are chosen with <x_i, y_j> = 0 for i different from j and <x_j, y_j> = cos(theta_j).]
  {SourceVector}{x}{trialvector}{local}{Vector in the source subspace.}{vector}
\dkDeclareSymbol[Orthonormal coordinate frame X for the source subspace, so X* X = I. Together with a target frame Y, the singular values of X* Y are the principal cosines. In the four-dimensional counterexample X is the vertical stack of I_2 and 0.]
  {Xframe}{X}{trialframe}{glossary}{Orthonormal coordinate matrix for the source subspace in the counterexample.}{isometric-embedding}
\dkDeclareSymbol[Hilbert space F of trial coordinates. E_0 embeds F isometrically into ambient space E; A_tr and the positive sine operator sin(Theta_0) act on F, while the residual R maps F into E.]
  {Ftrial}{F}{trialframe}{glossary}{Hilbert space of trial coordinates.}{hilbert-space}
\dkDeclareSymbol[Ambient Hilbert space E over the real or complex scalars. It contains the trial subspace range(E_0) and the orthogonal exact subspaces range(F_0) and range(F_1). The source theorem permits infinite-dimensional separable E.]
  {Eambient}{E}{ambient}{glossary}{Ambient Hilbert space containing the trial and exact subspaces.}{hilbert-space}
% Zero and all-ones vectors are distinguished from scalar 0 and 1.
\dkDeclareSymbol[Zero vector of the vector or Hilbert space in context: the additive identity. The zero subspace consists only of this vector.]
  {ZeroVector}{\mathbf{0}}{fg}{glossary}{Zero vector; the additive identity of the ambient vector space.}{vector}
\dkDeclareSymbol[All-ones vector in a specified finite-dimensional coordinate space: every coordinate equals 1. This vector depends on the chosen coordinates and dimension.]
  {OnesVector}{\mathbf{1}}{fg}{local}{All-ones vector in a chosen finite-dimensional coordinate space.}{vector}
% Exact/wanted family: the target subspace, its frame, projection and rotation.
\dkDeclareSymbol[Target subspace V in the transport problem: feasible maps satisfy W(U) = V, or equivalently W P W* = Q. The matching theorem also names this subspace L.]
  {Vtarget}{V}{wanted}{glossary}{Target subspace in the transport problem.}{linear-subspace}
\dkDeclareSymbol[Exact wanted frame F_0: F_ex -> E, a linear isometry onto the desired exact spectral subspace. F_0* returns its coordinates. P_ex = F_0 F_0* is its ambient orthogonal projection; in Figure 1, F_0(h) = (h, 0).]
  {F0}{F_0}{wantedframe}{glossary}{Isometric coordinates for the desired exact spectral subspace; star denotes the adjoint.}{isometric-embedding}
\dkDeclareSymbol[Self-adjoint operator block on coordinates for the desired exact spectral subspace. Its spectrum consists of the wanted exact spectral values; in the corresponding decomposition Ahat F_0 = F_0 Lambda_0 on the operator domain.]
  {Lambda0}{\Lambda_0}{wantedspectrum}{local}{Self-adjoint block on the desired exact spectral subspace.}{self-adjoint-operator}
\dkDeclareSymbol{u}{u}{wantedvector}{local}{Vector in the desired exact subspace.}{vector}
\dkDeclareSymbol{f0}{f_0}{wantedvector}{local}{Component or coordinate vector in the desired exact subspace.}{vector}
\dkDeclareSymbol{lambda0}{\lambda_0}{wantedspectrum}{local}{Wanted exact eigenvalue or wanted spectral value.}{scalar}
\dkDeclareSymbol{lambdau}{\lambda_u}{wantedspectrum}{local}{Eigenvalue associated with a desired exact vector.}{scalar}
\dkDeclareSymbol[Exact wanted orthogonal projection P_ex: E -> E, with P_ex = F_0 F_0*. It retains the desired exact spectral component. Its complementary projection is P_ex_perp = I - P_ex = F_1 F_1*; the projected trial frame is S = P_ex_perp E_0.]
  {Pexact}{P_{\mathrm{ex}}}{wantedprojection}{glossary}{Orthogonal projection onto the desired exact subspace in the sine-theta theorem.}{orthogonal-projection}
\dkDeclareSymbol[Exact complementary orthogonal projection P_ex_perp = I - P_ex = F_1 F_1*: E -> E. It retains the unwanted exact component. In Figure 1, P_ex_perp v = S(1), whereas P_ex_perp r is the unwanted component of the residual.]
  {PexactPerp}{P_{\mathrm{ex}}^\perp}{unwantedprojection}{local}{Orthogonal projection onto the exact complementary subspace.}{orthogonal-projection}
\dkDeclareSymbol{ReflectionSigma}{\Sigma}{wantedprojection}{local}{Reflection associated with the desired exact spectral projection.}{unitary-operator}
\dkDeclareSymbol[Orthogonal projection Q onto the target subspace V in the transport problem. A feasible orthogonal or unitary W satisfies W P W* = Q, where P projects onto the source U.]
  {Qtarget}{Q}{wantedprojection}{glossary}{Orthogonal projection onto the target subspace in the transport problem.}{orthogonal-projection}
\dkDeclareSymbol{ProjTarget}{P_V}{wantedprojection}{local}{Orthogonal projection onto the target subspace in the transport problem.}{orthogonal-projection}
\dkDeclareSymbol[Canonical direct rotation D from the source U to the target V for an acute pair. It is the unitary polar factor of T = QP + (I - Q)(I - P), and acts by the principal angle on each principal plane. R is reserved for the trial residual.]
  {Rrotation}{D}{wantedrotation}{glossary}{Canonical direct rotation carrying the source subspace onto the target.}{unitary-operator}
\dkDeclareSymbol[Direct rotation D between the ranges of orthogonal projections P and Q, including the nonacute case. Relative to range(P) and kernel(P), its diagonal blocks are positive and its off-diagonal blocks are opposite adjoints. It exists when the two crossed defect spaces have equal dimension, and is unique when both defects vanish.]
  {SourceDirectRotation}{D}{wantedrotation}{local}{Direct rotation, including the nonacute case in the source-result table.}{unitary-operator}
\dkDeclareSymbol[Orthonormal coordinate frame Y for the target subspace, so Y* Y = I. In the counterexample Y is the vertical stack of cos(theta) I_2 and sin(theta) J_0; hence X* Y = cos(theta) I_2 and both principal angles equal theta.]
  {Yframe}{Y}{wantedframe}{glossary}{Orthonormal coordinate matrix for the target subspace in the counterexample.}{isometric-embedding}
\dkDeclareSymbol[Vector y in the target subspace V. For principal vectors, unit x_j in U and y_j in V are paired so that <x_j, y_j> = cos(theta_j); distinct pairs have zero cross inner product.]
  {TargetVector}{y}{wantedvector}{local}{Vector in the target subspace.}{vector}
\dkDeclareSymbol[Hilbert space F_ex of exact wanted coordinates. The isometry F_0: F_ex -> E identifies it with the desired exact spectral subspace. This space differs from H, the additive perturbation.]
  {Hexact}{F_{\mathrm{ex}}}{wantedframe}{glossary}{Hilbert space of coordinates for the desired exact spectral subspace.}{hilbert-space}% Coordinate Hilbert space (not perturbation).
% Exact/unwanted family; the trial error operators are MAGENTA, not cyan.
\dkDeclareSymbol{Uperp}{U^{\perp}}{unwanted}{local}{Orthogonal complement of the desired exact subspace in the local notation.}{linear-subspace}
\dkDeclareSymbol[Exact complementary frame F_1: G -> E, a linear isometry onto the orthogonal complement of range(F_0). F_1* returns unwanted coordinates; F_1 F_1* = I - P_ex. In Figure 1, F_1(g) = (0, g).]
  {F1}{F_1}{unwantedframe}{glossary}{Isometric coordinates for the exact complementary subspace; star denotes the adjoint.}{isometric-embedding}
\dkDeclareSymbol[Self-adjoint unwanted block Lambda_1 on G, satisfying Ahat F_1 = F_1 Lambda_1 on its domain. Its spectrum is separated from the trial spectrum of A_tr by the specified positive delta. Figure 1 uses the scalar block 0.7.]
  {Lambda1}{\Lambda_1}{unwantedspectrum}{glossary}{Self-adjoint operator on exact complementary coordinates: the unwanted spectral block.}{self-adjoint-operator}
% Spectral measure and set (distinct from the ambient space E and sine block S).
\dkDeclareSymbol[Projection-valued spectral measure E associated with a self-adjoint operator A: A = integral lambda dE(lambda). For a Borel set S, E(S) is the orthogonal projection onto the corresponding spectral subspace.]
  {SpectralMeasure}{E}{wantedprojection}{glossary}{Projection-valued spectral measure of a self-adjoint operator.}{projection-valued-measure}
\dkDeclareSymbol[Borel subset S of the real spectral axis. The spectral projection E(S) selects the part of a self-adjoint operator with spectral values in S; its range is a reducing subspace.]
  {SpectralSet}{S}{wantedspectrum}{glossary}{Borel spectral set used to select a spectral subspace.}{set}
% Generic Sylvester coefficients and right side, by source/target role.
\dkDeclareSymbol[Left coefficient B of the Sylvester map X -> B X - X C; B acts on the output space of X. For the sine bound, B = Lambda_1 on the unwanted-coordinate space G.]
  {SylvesterLeft}{B}{unwantedspectrum}{glossary}{Left coefficient of the Sylvester equation; corresponds to the unwanted spectral block.}{self-adjoint-operator}
\dkDeclareSymbol[Right coefficient C of the Sylvester map X -> B X - X C; C acts on the input space of X. For the sine bound, C = A_tr on trial space F. The historical equation also denotes this coefficient by A.]
  {SylvesterRight}{C}{trialspectrum}{glossary}{Right coefficient of the Sylvester equation; corresponds to the trial operator.}{self-adjoint-operator}
\dkDeclareSymbol[Right-hand side Y in B X - X C = Y, with the same input and output spaces as X. For the sine bound, Y = F_1* R: F -> G, the residual expressed in unwanted exact coordinates.]
  {SylvesterRHS}{Y}{resid}{glossary}{Right-hand side of the Sylvester equation, obtained by projecting the residual.}{linear-operator}
\dkDeclareSymbol{w}{w}{unwantedvector}{local}{Vector in an unwanted or complementary spectral subspace.}{vector}
\dkDeclareSymbol{f1}{f_1}{unwantedvector}{local}{Component or coordinate vector in the unwanted exact subspace.}{vector}
\dkDeclareSymbol{lambda1}{\lambda_1}{unwantedspectrum}{local}{Unwanted exact eigenvalue or complementary spectral value.}{scalar}
\dkDeclareSymbol{lambdaw}{\lambda_w}{unwantedspectrum}{local}{Eigenvalue associated with an unwanted exact vector.}{scalar}
\dkDeclareSymbol[Hilbert space G of exact complementary coordinates. F_1: G -> E embeds it onto the unwanted exact subspace; Lambda_1 acts on G and X = F_1* E_0 maps trial coordinates into G.]
  {Gunwanted}{G}{unwantedframe}{glossary}{Hilbert space of exact complementary coordinates.}{hilbert-space}
\dkDeclareSymbol{E1}{E_1}{unwantedframe}{local}{Isometric coordinates for the complementary trial subspace.}{isometric-embedding}
\dkDeclareSymbol{Vperp}{V^{\perp}}{unwanted}{local}{Orthogonal complement of the trial subspace in the local notation.}{linear-subspace}
% Angle and error geometry, including distinct rectangular and positive maps.
\dkDeclareSymbol[Individual principal angle, between 0 and pi/2. For orthonormal subspace frames X and Y, the singular values of X* Y are the principal cosines. In Figure 1, theta is the angle between the unit trial vector v and the exact wanted line.]
  {theta}{\theta}{sine}{glossary}{Principal angle between the subspaces under comparison.}{angle}
\dkDeclareSymbol[Full positive angle operator Theta = diag(Theta_0, Theta_1) for the source/target decompositions. Theta_0 compares the source with the target; Theta_1 compares their orthogonal complements. Its operator norm is the largest spectral angle. For positive c, the conditions Theta <= c and Theta < c mean operator norm at most c and strictly below c, respectively.]
  {Theta}{\Theta}{angleoperator}{glossary}{Full angle operator, the diagonal sum of the source and complementary angle blocks.}{angle-operator}
\dkDeclareSymbol[Directed principal-angle operator Theta_0 on source or trial coordinates. For a trial frame E_0 and exact wanted projection P_ex, cos(Theta_0) = sqrt(E_0* P_ex E_0), and sin(Theta_0) = sqrt(I_F - E_0* P_ex E_0). In the two-projection geometry it comes from the positive cosine compression P Q P restricted to range(P).]
  {Theta0}{\Theta_0}{angleoperator}{glossary}{Directed principal-angle operator acting on source or trial coordinates.}{angle-operator}
\dkDeclareSymbol[Complementary directed angle operator Theta_1, obtained from the projections I-P and I-Q by the same cosine-compression construction as Theta_0. It acts on the source complement, or its isometric coordinates. The full angle operator is diag(Theta_0, Theta_1).]
  {Theta1}{\Theta_1}{angleoperator}{glossary}{Directed angle operator comparing the orthogonal complements.}{angle-operator}
\dkDeclareSymbol[Right-angle source defect D_0 = range(P) intersect kernel(Q). These are source vectors orthogonal to the target. The full-angle conclusion in the Section 8 smallness criterion assumes dim(D_0) = dim(D_1); in Lean their identification is CrossedDefectsEquivalent.]
  {CrossedSourceDefect}{D_0}{fg}{glossary}{Source directions perpendicular to the target subspace.}{linear-subspace}
\dkDeclareSymbol[Right-angle target defect D_1 = kernel(P) intersect range(Q). These are target vectors orthogonal to the source. Equality of its Hilbert dimension with that of D_0 is assumed for the full-angle conclusion of Theorem 8.2.]
  {CrossedTargetDefect}{D_1}{fg}{glossary}{Target directions perpendicular to the source subspace.}{linear-subspace}
\dkDeclareSymbol[Positive complementary cosine block C_1 = cos(Theta_1). The reference complement is identified with the perturbed complement by the direct rotation D, so A_1 and Lambda_1 act on common coordinates. Theorem 8.1 compares A_1 - alpha I with C_1 (Lambda_1 - alpha I) C_1 as quadratic forms.]
  {ComplementCosineBlock}{C_1}{anglecos}{local}{Cosine of the complementary angle block in the spectral-repulsion inequality.}{self-adjoint-operator}
\dkDeclareSymbol[Scalar sine of a principal angle. For a unit principal trial vector v and exact wanted projection P_ex, sin(theta) = ||(I - P_ex)v||. In Figure 1 it is the length of the pink leg S(1).]
  {sintheta}{\sin\theta}{sine}{local}{Sine of an individual principal angle: the length of the complementary component of a unit principal trial vector.}{scalar}
\dkDeclareSymbol[Positive sine operator sin(Theta_0) = sqrt(I_F - E_0* P_ex E_0): F -> F. It equals |S| for S = (I - P_ex)E_0: F -> E. Its eigenvalues in finite dimensions are the principal sines; the sine theorem bounds its ideal norm by N(R)/delta.]
  {sinTheta0}{\sin\Theta_0}{sineoperator}{glossary}{Positive sine operator on trial coordinates; the modulus of the rectangular complementary error.}{angle-operator}
\dkDeclareSymbol[Rectangular complementary error S = (I - P_ex) E_0: F -> E. It records how much each trial vector leaks outside the exact wanted subspace. Its modulus |S| = sqrt(S* S) is sin(Theta_0): F -> F. In rank one S(1) = P_ex_perp v.]
  {S}{S}{sineblock}{glossary}{Rectangular complementary error: the trial frame projected onto the unwanted exact subspace.}{linear-operator}
\dkDeclareSymbol[Complementary-coordinate map X = F_1* E_0: F -> G, with S = F_1 X. It satisfies Lambda_1 X - X A_tr = F_1* R on domain(A_tr). Isometric multiplication gives N(X) = N(S); the separate modulus identity abs(S) = sin(Theta_0) gives N(S) = N(sin(Theta_0)), whenever these norms are finite.]
  {X}{X}{sylvester}{glossary}{Solution of the Sylvester equation; for the sine bound, the trial frame in unwanted exact coordinates.}{linear-operator}
\dkDeclareSymbol{alphaAngle}{\alpha}{anglemax}{local}{Largest principal angle.}{angle}
\dkDeclareSymbol[Scalar c = cos(theta) in the counterexample. In the general two-projection decomposition, the same role denotes a positive cosine operator C; the cosine and sine blocks commute and satisfy C squared + S squared = I.]
  {CosAngle}{c}{anglecos}{local}{Cosine of a principal angle or the corresponding positive cosine block.}{scalar-or-operator}
\dkDeclareSymbol[Scalar s = sin(theta) in the counterexample. In the general two-projection decomposition, the same role denotes a positive sine operator S; it commutes with the positive cosine block C and satisfies C squared + S squared = I.]
  {SinAngle}{s}{sineoperator}{local}{Sine of a principal angle or the corresponding positive sine block.}{scalar-or-operator}
\dkDeclareSymbol[Positive small-angle parameter in the local subspace comparison. It controls the allowed angular deviation; it is an angle parameter rather than a perturbation-operator norm.]
  {EpsilonAngle}{\varepsilon}{anglemax}{local}{Angle parameter controlling the maximum subspace deviation.}{angle}
\dkDeclareSymbol{PhiAngle}{\varphi}{angleoperator}{local}{Auxiliary angle or angle operator in the two-subspace comparison.}{angle}
\dkDeclareSymbol{Bblock}{B}{sineblock}{local}{Off-diagonal block in the two-subspace decomposition.}{linear-operator}
\dkDeclareSymbol{Cblock}{C}{sylvester}{local}{Complementary or Sylvester block in the two-subspace decomposition.}{linear-operator}
\dkDeclareSymbol{Ablock}{A}{anglecos}{local}{Positive cosine-related block in the two-subspace decomposition.}{linear-operator}
% Residual/error-displacement family.
\dkDeclareSymbol[Residual vector r = Ahat v - rho v = R(1) for the rank-one frame E_0(x) = x v. Figure 1 draws it from rho v to Ahat v. For unit v and the Rayleigh value rho = <v, Ahat v>, <v, r> = 0; a general trial scalar need not give orthogonality.]
  {r}{r}{residvector}{local}{Residual vector of an approximate eigenpair.}{vector}
\dkDeclareSymbol[Bounded residual R: F -> E, agreeing with Ahat E_0 - E_0 A_tr on domain(A_tr). The domain condition requires E_0(domain(A_tr)) inside domain(Ahat). R = 0 means that E_0 intertwines A_tr with Ahat; its norm measures the defect.]
  {R}{R}{resid}{glossary}{Bounded residual of the trial eigenproblem, agreeing with the operator difference on the trial domain.}{linear-operator}
\dkDeclareSymbol[Full displacement Delta_W = I - W on the ambient space. In the failed Ky Fan comparison W = W_(pi/4); Delta_W has singular values sqrt(2), sqrt(2), 0, 0. Multiplying by a principal-plane projection restricts its input to that plane.]
  {Kdisplacement}{\Delta_W}{displacement}{glossary}{Full displacement of a transporter from the identity.}{operator-valued-function}
\dkDeclareSymbol[Squared full displacement Dcal(W) = (I-W)* (I-W), a positive operator. Proposition 4.3 compares its unitarily invariant norm for a direct rotation D with that for every feasible transporter W.]
  {SquaredDisplacement}{\mathcal D}{displacement}{local}{Squared full displacement as a function of the transporter.}{operator-valued-function}
% Matching proof: the gradient and its spectral level are not a space or gap.
\dkDeclareSymbol[Skew-symmetric spectral gradient G of Phi_p. It commutes with W and represents first variations by d/dt Phi_p(W exp(tB)) at t = 0 equal to -tr(G B) for skew-symmetric B. Feasible stationarity gives J_src G J_src = -G.]
  {SpectralGradient}{G}{displacement}{local}{Skew-symmetric spectral gradient of the displacement power sum in the matching proof.}{linear-operator}
\dkDeclareSymbol[Positive spectral level g of the skew-symmetric gradient G: on this invariant subspace, -G squared = g squared I. Consequently G/g squares to -I and supplies the complex structure used in the real matching proof.]
  {GradientLevel}{g}{displacement}{local}{Positive gradient level; the gradient squared is minus this level squared times the identity.}{scalar}
\dkDeclareSymbol[Complex structure I_g = G/g on the positive gradient-level space V_g. It is a real isometry with square -I. Multiplication by a+ib means a times identity plus b times I_g; source reflection J_src reverses I_g and acts as conjugation.]
  {GradientComplexStructure}{\mathcal I_g}{displacement}{local}{Complex multiplication constructed from a positive gradient level.}{linear-operator}
\dkDeclareSymbol[Real skew-symmetric matrix B, with B transpose = -B, giving the orthogonal variation W exp(tB). If B commutes with the source projection P, this curve preserves the constraint W(U) = V.]
  {SkewDirection}{B}{ambient}{local}{Skew-symmetric direction for a first variation of an orthogonal transporter.}{linear-operator}
% Spectral separation family: gap and original gap are NOT identical.
\dkDeclareSymbol[Positive spectral separation in the stated configuration. In the sine theorem it separates the trial operator from the exact complementary block; Theorem 8.1 uses it for the ordered reference gap, and Theorem 8.2 for the perturbed central-band separation. Figure 1 uses delta = abs(rho - 0.7).]
  {delta}{\delta}{gap}{glossary}{Positive spectral separation in the stated gap configuration.}{scalar}
\dkDeclareSymbol{g}{g}{gaporiginal}{local}{Original spectral gap of the reference operator; distinct from the certified sine-theta separation.}{scalar}
\dkDeclareSymbol{Delta}{\Delta}{gapstat}{local}{Population or statistical spectral gap.}{scalar}
\dkDeclareSymbol[Spectral endpoint alpha. In the interval/exterior configuration, one spectrum lies between beta and alpha, the other at or below beta - delta or at or above alpha + delta. Theorem 8.2 places the selected perturbed block in the central interval. In Theorem 8.1, alpha bounds the lower reference block and alpha + delta bounds the upper block from below.]
  {alphaEndpoint}{\alpha}{gapinterval}{glossary}{Upper endpoint of the central band, or lower-block cut in the ordered gap.}{scalar}
\dkDeclareSymbol[Lower endpoint beta of the finite interval from beta to alpha in the gap hypothesis. One spectrum lies in that interval; the other lies at or below beta - delta or at or above alpha + delta.]
  {betaEndpoint}{\beta}{gapinterval}{glossary}{Lower endpoint of the finite interval in the interval/exterior gap hypothesis.}{scalar}
% Spectral/projector and auxiliary matrix objects, context-scoped.
\dkDeclareSymbol{P}{P}{fg}{local}{Orthogonal projection; its source or exact-target role is specified locally.}{orthogonal-projection}% Neutral historical Manim key; prefer Pexact/Psource.
\dkDeclareSymbol{Q}{Q}{fg}{local}{Orthogonal projection; its role is specified locally.}{orthogonal-projection}% Neutral historical Manim key; prefer Qtarget.
\dkDeclareSymbol[Two-projection intertwiner T = QP + (I - Q)(I - P), satisfying TP = QT and T* T = I - (P - Q) squared. For ||P - Q|| < 1 it is invertible; its unitary polar factor T |T| inverse is the canonical direct rotation.]
  {IntertwinerT}{T}{intertwiner}{glossary}{Two-projection intertwiner; its unitary polar factor is the acute direct rotation.}{linear-operator}
\dkDeclareSymbol[Planar pi/2-rotation matrix J_0 = ((0, -1), (1, 0)), with J_0 squared = -I_2 and J_0 transpose = -J_0. It defines the rotation C = cos(2 theta) I_2 + sin(2 theta) J_0 used to construct the real competitor.]
  {Jturn}{J_0}{turn}{glossary}{Planar rotation through $\pi/2$, with square minus the identity.}{matrix}
\dkDeclareSymbol[Partial isometry V_pol in the polar decomposition S = V_pol |S|, where |S| = sqrt(S* S). It is isometric on the orthogonal complement of kernel(S) and zero on kernel(S). For an invertible operator its polar factor is unitary.]
  {PolarU}{V_{\mathrm{pol}}}{intertwiner}{glossary}{Partial isometry in the polar decomposition; unitary when the operator is invertible.}{partial-isometry}
\dkDeclareSymbol[Planar rotation C = cos(2 theta) I_2 + sin(2 theta) J_0 through twice the principal angle. The competitor W_theta is orthogonally conjugate to C direct-sum I_2, so it rotates one two-plane and fixes the other.]
  {Cturn}{C}{turn}{glossary}{Planar rotation through twice the principal angle in the counterexample.}{matrix}
\dkDeclareSymbol[Orthogonal mixing matrix H_mix = (1/sqrt(2)) times the block matrix ((I_2, I_2), (I_2, -I_2)). It constructs W_theta = H_mix (C direct-sum I_2) H_mix transpose.]
  {Hmix}{H_{\mathrm{mix}}}{mix}{glossary}{Orthogonal mixing matrix in the counterexample.}{matrix}
\dkDeclareSymbol[Source reflection J_src = 2P - I: it acts as +I on the source subspace and -I on its orthogonal complement, so its square is I. In the matching proof stationarity gives J_src G J_src = -G for the skew-symmetric spectral gradient G.]
  {Jreflect}{J_{\mathrm{src}}}{trialprojection}{glossary}{Reflection equal to twice the source projection minus the identity.}{linear-operator}
\dkDeclareSymbol[Skew angle generator J_ang for a direct rotation D = exp(J_ang Theta). On each principal plane it is the oriented rotation through pi/2; it vanishes on the common fixed directions. It commutes with the full angle operator Theta.]
  {AngleGenerator}{J_{\mathrm{ang}}}{turn}{glossary}{Skew angle generator in the direct-rotation result summaries.}{linear-operator}
\dkDeclareSymbol{Dblock}{D}{mix}{local}{Auxiliary diagonal or block matrix in the counterexample.}{matrix}
\dkDeclareSymbol[Orthogonal projection onto the first principal plane span(e_0, e_3) in the four-dimensional counterexample. For Delta_W = I - W_(pi/4), the Ky Fan 2-norm of Delta_W Omega_1 equals 2.]
  {OmegaOne}{\Omega_1}{turn}{glossary}{First principal-plane projection in the failed Ky Fan comparison.}{orthogonal-projection}
\dkDeclareSymbol[Orthogonal projection onto the second principal plane span(e_1, e_2) in the four-dimensional counterexample. For Delta_W = I - W_(pi/4), the Ky Fan 2-norm of Delta_W Omega_2 equals 2, although the trace norm of Delta_W is only 2 sqrt(2).]
  {OmegaTwo}{\Omega_2}{intertwiner}{glossary}{Second principal-plane projection in the failed Ky Fan comparison.}{orthogonal-projection}
\dkDeclareSymbol{Dplanar}{D}{wantedrotation}{local}{Planar direct-rotation block.}{unitary-operator}
\dkDeclareSymbol{Hreflection}{H}{trialcompetitor}{local}{Planar reflection carrying the source line to the target.}{linear-operator}
% Unitarily invariant gauge family; N, M are NOT a norm versus its input.
\dkDeclareSymbol[Unitarily invariant norm N, or a symmetric operator-ideal norm on its domain of finiteness. It satisfies N(U T V) = N(T) for unitaries and N(B T) <= ||B||op N(T). In the sine theorem it is normalized by N(u v*) = ||u|| ||v|| and has Ky Fan dominance where both norms are finite. Both R and sin(Theta_0) must belong to its ideal.]
  {NormN}{N}{norm}{glossary}{Unitarily invariant norm on its domain of finiteness; in the sine theorem, normalized by $N(uv^*)=\|u\|\|v\|$.}{norm}
\dkDeclareSymbol[Sine-theta separation G_delta(B,C): finite interval/exterior separation in either spectral order, or ordered separation B <= a I and C >= (a+delta) I, or the reverse. The ordered alternatives allow both spectra to be unbounded.]
  {SineGap}{\mathcal G_\delta}{gap}{local}{Interval/exterior separation or either ordered separation, as in the sine-theta theorem.}{predicate}
\dkDeclareSymbol[Auxiliary unitarily invariant norm M witnessing the Q-norm property N(T) squared = M(T* T). For the Frobenius norm N, one can take M to be the trace norm.]
  {NormM}{M}{normaux}{local}{Auxiliary unitarily invariant norm used to define a Q-norm through the squared modulus.}{norm}
\dkDeclareSymbol[Q-norm property: a unitarily invariant norm N is a Q-norm when N(T) squared = M(T* T) for some unitarily invariant norm M. The Frobenius norm is an example, with M the trace norm; the repaired full-displacement theorem uses this norm class.]
  {NormQ}{Q}{normaux}{glossary}{Q-norm: its square is a unitarily invariant norm of the squared modulus.}{norm}
\dkDeclareSymbol[Singleton trace cost s_i = min(4 sin(theta_i/2), 2) for angle label i. It transports one real principal plane independently, choosing a rotation of cost 4 sin(theta_i/2) or a reflection of cost 2.]
  {SingletonCost}{s}{norm}{glossary}{Minimum trace displacement for one real principal plane: choose a rotation or a reflection.}{scalar-cost}
\dkDeclareSymbol[Paired trace cost e_ij = 4 sin((theta_i + theta_j)/2) for distinct angle labels i,j. It is attained by coupling their principal planes in a real four-dimensional block. The edge weight s_i + s_j - e_ij is the difference between the cost of two singletons and this paired cost.]
  {PairCost}{e}{norm}{glossary}{Trace displacement cost of coupling two principal planes in a four-dimensional block.}{scalar-cost}
\dkDeclareSymbol[Scalar phase cost f_p(t) = (2 sin(t/2)) to the power p, for 0 <= t <= pi. A real planar rotation through t contributes two copies, so its displacement power cost is 2 f_p(t).]
  {MatchingPhaseCost}{f_p}{norm}{local}{Displacement power cost for one scalar phase.}{scalar-cost}
\dkDeclareSymbol[Singleton power cost a_i(p) = min(2 f_p(theta_i), 2 to the power p). These are the rotation and reflection choices for one principal plane. At p=1 it equals the trace-cost weight s_i.]
  {SingletonPowerCost}{a}{norm}{local}{Minimum displacement power cost for one angle.}{scalar-cost}
\dkDeclareSymbol[Paired power cost b_ij(p) = min over d <= tau <= h of 2(f_p(h+tau) + f_p(h-tau)), where h=(theta_i+theta_j)/2 and d=abs(theta_i-theta_j)/2. This is the positive-determinant pair minimum; unrestricted block optimization also compares two singletons. At p=1 it equals e_ij.]
  {PairPowerCost}{b}{norm}{local}{Positive-determinant displacement power cost for a pair of angles.}{scalar-cost}
\dkDeclareSymbol[Matching M on the finite general graph of positive principal-angle labels. Each label belongs to at most one edge; unmatched labels use singleton transporters. Minimizing paired and singleton costs is equivalent to maximum-weight matching with edge weights g_ij = s_i + s_j - e_ij. The empty matching is allowed.]
  {Matching}{M}{normaux}{glossary}{Matching of positive principal-angle labels: each label is used in at most one pair.}{matching}
\dkDeclareSymbol[Displacement power sum Phi_p(W) = sum_j s_j(I - W)^p. The matching proof first minimizes it for 1 < p < 2, then lets p decrease to 1, where it becomes the trace norm ||I - W||_1.]
  {PowerCost}{\Phi}{norm}{local}{Sum of the p-th powers of displacement singular values.}{scalar-cost}
% Auxiliary matching-proof notation: neutral ink, with the same hover registry.
\dkDeclareSymbol[Off-diagonal gradient block G_01 mapping the source complement K perpendicular into K. Stationarity makes the diagonal blocks of G zero, so G = ((0,G_01),(-G_01 transpose,0)) in the source/complement decomposition.]
  {GradientOffDiagonal}{G_{01}}{fg}{local}{Complement-to-source block of the skew-symmetric gradient.}{linear-operator}
\dkDeclareSymbol[Orthogonal group O(n): real n by n matrices W satisfying W transpose times W = I. Its elements preserve Euclidean inner products and lengths. The notation O(2) refers to the same group in dimension two.]
  {OrthogonalGroup}{O(n)}{fg}{local}{Group of real orthogonal matrices.}{group}
\dkDeclareSymbol[Feasible transporters F = {W in O(n) : W(K) = L}. Equivalently W P_K W transpose = P_L. This nonempty compact set is the optimization domain throughout the matching proof.]
  {FeasibleTransport}{\mathcal F}{fg}{local}{Orthogonal maps carrying the source subspace onto the target.}{set}
\dkDeclareSymbol[Singular value s_j(T), the nonnegative square root of an eigenvalue of T transpose times T, counted with multiplicity. Summing s_j(I-W) gives the trace displacement; summing its p-th powers gives Phi_p(W). These singular values differ from the singleton weights s_i.]
  {TransportSingularValues}{\symglyph{SingularValueBase}_j}{fg}{local}{A displacement singular value in the matching proof.}{scalar}
\dkDeclareSymbol[Scalar singular-value magnitude s between zero and two for the displacement I-W of an orthogonal matrix. Uniform convergence of s to the power p to s on this interval permits the trace-norm limit. This variable differs from the rotation parameter s used earlier in the proof.]
  {DisplacementSize}{s}{fg}{local}{Scalar displacement singular-value magnitude in the uniform limit.}{scalar}
\dkDeclareSymbol[All-singleton cost C_0 = sum_i s_i, before any two angle labels are paired. A matching has cost C_0 minus the sum of its selected edge weights.]
  {SingletonBaseline}{C_0}{fg}{local}{Trace displacement when every angle uses its singleton transporter.}{scalar-cost}
\dkDeclareSymbol[Matching power cost C_p(M): sum of b_ij(p) over selected pairs plus a_i(p) over unmatched labels. Its minimum equals the minimum of Phi_p over feasible transporters. At p=1 it is C(M).]
  {MatchingPowerCost}{C_p}{fg}{local}{Total paired and singleton power cost of a matching.}{scalar-cost}
\dkDeclareSymbol[Invariant orthogonal block V_a preserved by the source projection, target projection, and minimizing transporter. The displacement cost adds over these blocks, and their principal-angle lists combine into the original list.]
  {MatchingBlock}{V_a}{fg}{local}{Simultaneous invariant block in the matching proof.}{linear-subspace}
\dkDeclareSymbol[Gradient-level space V_g: the eigenspace of -G squared for a positive eigenvalue g squared. It is preserved by G, W, and the source reflection J_src. Multiplication by i on it is G/g.]
  {GradientLevelSpace}{V_g}{fg}{local}{Positive gradient-level eigenspace.}{linear-subspace}
\dkDeclareSymbol[Eigenphase varphi of W on V_g viewed as a complex space with multiplication by i = G/g. It lies between zero and pi and satisfies f_p prime(varphi) = g. There are at most two such phases.]
  {GradientEigenphase}{\varphi}{fg}{local}{One of the at most two phases at a positive gradient level.}{angle}
\dkDeclareSymbol[Orthogonal projection E onto the first phase space of W on V_g. Here W = exp(i varphi_1) E + exp(i varphi_2) (I-E). In coordinates where the source reflection is conjugation, E = X + iY.]
  {GradientPhaseProjection}{E}{fg}{local}{Projection onto the first phase space in a gradient-level space.}{orthogonal-projection}
\dkDeclareSymbol[Real part X of the phase projection E = X + iY. It is a real symmetric matrix with spectrum in the closed interval from zero to one, satisfying X squared - Y squared = X and XY + YX = Y. This X differs from the source-coordinate and Sylvester matrices elsewhere in the paper.]
  {ProjectionRealPart}{X}{fg}{local}{Real symmetric part of the phase projection.}{linear-operator}
\dkDeclareSymbol[Imaginary part Y of the phase projection E = X + iY. It is a real skew-symmetric matrix. It maps the lambda eigenspace of X to the 1-lambda eigenspace and constructs the small invariant planes.]
  {ProjectionImaginaryPart}{Y}{fg}{local}{Real skew-symmetric imaginary part of the phase projection.}{linear-operator}
\dkDeclareSymbol[Eigenvalue lambda in the closed interval from zero to one of the real part X of the phase projection. It is local to the block construction and differs from the original operator's spectral parameter.]
  {ProjectionEigenvalue}{\lambda}{fg}{local}{Eigenvalue of the real part of the phase projection.}{scalar}
\dkDeclareSymbol[Plane coupling eta = sqrt(lambda (1-lambda)) for 0 < lambda < 1. If u is a unit lambda eigenvector of X, then v = Yu/eta completes its invariant real plane.]
  {ProjectionCoupling}{\eta}{fg}{local}{Normalization used to construct an invariant real plane.}{scalar}
\dkDeclareSymbol[Pair mean angle h = (theta_i + theta_j)/2. Together with d = (theta_i - theta_j)/2 after ordering theta_i >= theta_j, it parametrizes the two-angle block cost.]
  {PairMeanAngle}{h}{fg}{local}{Mean of the two ordered principal angles.}{angle}
\dkDeclareSymbol[Pair half-difference d = (theta_i - theta_j)/2 after ordering theta_i >= theta_j. It satisfies 0 <= d <= h and h+d <= pi/2.]
  {PairHalfDifference}{d}{fg}{local}{Half the difference of the two ordered principal angles.}{angle}
\dkDeclareSymbol[Scalar pair parameter tau in the closed interval from d to h. The paired power cost is the minimum of 2(f_p(h+tau) + f_p(h-tau)) over this interval. At p=1 a minimum is attained at tau=h.]
  {PairPhaseParameter}{\tau}{fg}{local}{Scalar parameter in the two-angle power-cost minimum.}{scalar}
\dkDeclareSymbol[Attaining rotation angle gamma = arccos(cos(tau)/cos(d)). Choosing U_0 = U_1 = R_gamma realizes phase parameters x=h and y=tau in the four-dimensional block.]
  {PairAttainingAngle}{\gamma}{fg}{local}{Rotation angle realizing a chosen scalar pair parameter.}{angle}
\dkDeclareSymbol[Four-dimensional direct rotation D = ((C,-S),(S,C)) in source and complementary principal coordinates, with C the diagonal principal cosines and S the diagonal principal sines. Every transporter on this block is D diag(U_0,U_1).]
  {PairDirectRotation}{D}{fg}{local}{Direct rotation on the two principal planes.}{unitary-operator}
\dkDeclareSymbol[Diagonal cosine block C = diag(cos(theta_i), cos(theta_j)) in the four-dimensional direct rotation D. This matrix differs from the scalar displacement cost C(M).]
  {PairCosineBlock}{C}{fg}{local}{Diagonal principal-cosine block.}{matrix}
\dkDeclareSymbol[Diagonal sine block S = diag(sin(theta_i), sin(theta_j)) in the four-dimensional direct rotation D. This matrix is local to the matching proof.]
  {PairSineBlock}{S}{fg}{local}{Diagonal principal-sine block.}{matrix}
\dkDeclareSymbol[Free planar orthogonal matrices U_0 and U_1 acting on the two-dimensional source and its two-dimensional complement, respectively. Multiplying D by diag(U_0,U_1) parametrizes every transporter on this four-dimensional block.]
  {PairFreeOrthogonal}{U}{fg}{local}{Free source or complementary planar orthogonal factor.}{unitary-operator}
\dkDeclareSymbol[Phase parameter x in the closed interval from h to pi-h for a positive-determinant four-dimensional transporter. Its eigenvalues are exp(plus or minus i(x+y)) and exp(plus or minus i(x-y)).]
  {PairPhaseX}{x}{fg}{local}{First phase parameter of a four-dimensional transporter.}{angle}
\dkDeclareSymbol[Phase parameter y in the closed interval from d to pi-d for a positive-determinant four-dimensional transporter. Together with x it gives rotational phases x+y and x-y, folded into the closed interval from zero to pi for the displacement cost.]
  {PairPhaseY}{y}{fg}{local}{Second phase parameter of a four-dimensional transporter.}{angle}
\dkDeclareSymbol[Scalar a = cos(x) in the four-dimensional phase calculation. Together with b = cos(y), it expresses the trace as 4ab and determines the reciprocal characteristic polynomial.]
  {PairPhaseCosineX}{a}{fg}{local}{Cosine of the first phase parameter in the trace calculation.}{scalar}
\dkDeclareSymbol[Scalar b = cos(y) in the four-dimensional phase calculation, paired with a = cos(x). These scalars are local to the characteristic-polynomial calculation.]
  {PairPhaseCosineY}{b}{fg}{local}{Cosine of the second phase parameter in the trace calculation.}{scalar}
\dkDeclareSymbol[Rotation parameter r: the two free orthogonal factors use planar rotation angles r+s and r-s. It appears in cos(y) = cos(d) cos(r) for two rotations.]
  {PairRotationR}{r}{fg}{local}{Half the sum of the two free planar rotation angles.}{angle}
\dkDeclareSymbol[Rotation parameter s: the two free orthogonal factors use planar rotation angles r+s and r-s. It appears in cos(x) = cos(h) cos(s) for two rotations. This scalar differs from the singular values and singleton costs.]
  {PairRotationS}{s}{fg}{local}{Half the difference of the two free planar rotation angles.}{angle}
\dkDeclareSymbol[Planar reflection J_2 = diag(1,-1). Every two-dimensional orthogonal matrix of determinant -1 can be written as a planar rotation times J_2.]
  {PairPlanarReflection}{J_2}{fg}{local}{Fixed planar reflection used to parametrize determinant -1 factors.}{unitary-operator}
\dkDeclareSymbol[Real planar rotation R_t = ((cos(t),-sin(t)),(sin(t),cos(t))). Its two displacement singular values are both 2 times abs(sin(t/2)). In the matching proof a rotation through t in the closed interval from zero to pi contributes 2 f_p(t).]
  {MatchingPlanarRotation}{\symglyph{PlanarRotationBase}_t}{fg}{local}{Planar rotation through the indicated angle.}{unitary-operator}
% Statistical application / operator provenance.
\dkDeclareSymbol{Sigma}{\Sigma}{old}{local}{Population covariance matrix.}{matrix}
\dkDeclareSymbol{SigmaHat}{\widehat\Sigma}{current}{local}{Estimated or perturbed covariance matrix.}{matrix}
\dkDeclareSymbol{YWSV}{V}{wanted}{local}{Orthonormal coordinates for the population eigenspace in the statistical application.}{isometric-embedding}
\dkDeclareSymbol{YWSVhat}{\widehat V}{trialframe}{local}{Estimated eigenspace frame in the statistical application.}{isometric-embedding}

% Figure-specific objects retain their own meanings when bare letters overlap.
\dkDeclareSymbol[Signed rotation angle phi of the displayed map on its invariant real two-plane. At theta = pi/4 the direct rotation turns its two principal planes by opposite signed angles of magnitude pi/4; the competitor turns its moving plane by pi/2 and fixes its other plane.]
  {PlaneRotationAngle}{\varphi}{sine}{local}{Signed rotation angle of the orthogonal map on the displayed invariant plane.}{angle}
\dkDeclareSymbol[Unit vector q in the displayed invariant plane. The map T is the direct rotation D or the competitor W in that panel; Tq is the image and (I - T)q = q - Tq is the displacement.]
  {FigureInput}{q}{ambient}{local}{Unit input to the orthogonal map shown in this invariant-plane panel.}{vector}
\dkDeclareSymbol[Image Tq of the panel unit input q under T = D or W. Both q and Tq lie on the unit circle; the green displacement arrow joins Tq to q, so it represents q - Tq.]
  {FigureImage}{Tq}{wantedvector}{local}{Image of the unit input under the orthogonal map shown in this panel.}{vector}
\dkDeclareSymbol[Displacement (I - T)q = q - Tq, drawn from the image Tq to the unit input q. On a plane rotated through phi its length is 2 |sin(phi/2)|; it vanishes on a fixed plane.]
  {DisplacementVector}{(I-T)q}{displacement}{local}{Difference between a unit input and its image under the displayed orthogonal map.}{vector}
\dkDeclareSymbol[First principal plane Pi_1 = span(e_0, e_3) for the source-target pair in the real four-dimensional counterexample. At theta = pi/4 the direct rotation acts on this plane by a quarter of pi, with displacement singular value sqrt(2 - sqrt(2)).]
  {CounterPrincipalOne}{\Pi_1}{fg}{glossary}{First principal plane in the real counterexample: the span of e0 and e3.}{linear-subspace}
\dkDeclareSymbol[Second principal plane Pi_2 = span(e_1, e_2) for the source-target pair in the real four-dimensional counterexample. The direct rotation acts by the same angle magnitude as on Pi_1, with the opposite sign in the displayed plane coordinates.]
  {CounterPrincipalTwo}{\Pi_2}{fg}{glossary}{Second principal plane in the real counterexample: the span of e1 and e2.}{linear-subspace}
\dkDeclareSymbol[Moving plane M = range(W - I) = span(e_0 + e_2, e_1 + e_3) for the competitor. W rotates it by 2 theta; at theta = pi/4 the two displacement singular values contributed by this plane are both sqrt(2).]
  {CounterMovingPlane}{M}{ambient}{glossary}{Moving plane of the competing transporter: the range of W minus the identity.}{linear-subspace}
\dkDeclareSymbol[Fixed plane Z = kernel(W - I) = span(e_0 - e_2, e_1 - e_3) for the competitor. W acts as the identity here, so this plane contributes two zero displacement singular values.]
  {CounterFixedPlane}{Z}{ambient}{glossary}{Fixed plane of the competing transporter: the kernel of W minus the identity.}{linear-subspace}
\dkDeclareSymbol[Direct-rotation displacement singular value a = 2 sin(pi/8) = sqrt(2 - sqrt(2)) at theta = pi/4. Each of its two principal planes contributes two copies, giving trace displacement 4a.]
  {CounterChordA}{a}{displacement}{glossary}{Direct-rotation displacement singular value in the counterexample: sqrt(2 - sqrt(2)).}{scalar}
\dkDeclareSymbol[Competitor displacement singular value b = 2 sin(pi/4) = sqrt(2) at theta = pi/4. Its moving plane contributes two copies; the fixed plane contributes two zeros, giving trace displacement 2b.]
  {CounterChordB}{b}{displacement}{glossary}{Competing-transporter displacement singular value in the counterexample: square root of 2.}{scalar}
\dkDeclareSymbol[Singular value s_j(I - T) of the full ambient displacement, ordered from largest to smallest. Each rotating real two-plane contributes two values 2 |sin(phi/2)|; fixed directions contribute zeros. Their sum is the trace cost.]
  {DisplacementSingular}{\symglyph{SingularValueBase}_j}{displacement}{local}{An individual singular value of the full displacement from the identity.}{scalar}
\dkDeclareSymbol[Edge weight g_ij = s_i + s_j - e_ij, the difference between the cost of two singletons and the paired cost. A matching has total cost C_0 minus the sum of its edge weights, with C_0 = sum_i s_i. Maximizing this sum gives the minimum transport cost; edges of negative weight may be omitted.]
  {MatchingEdgeWeight}{g_{ij}}{current}{glossary}{Edge weight: the cost of two singletons minus their paired cost.}{scalar-weight}
\dkDeclareSymbol[Total trace-norm displacement cost. For a transporter W, C(W) = ||I - W||_1, the sum of its displacement singular values. For a matching M, C(M) is the sum of its paired and unmatched singleton costs.]
  {TraceCost}{C}{norm}{glossary}{Trace-norm displacement cost of the indicated orthogonal transporter or matching.}{scalar-cost}

% Spectral values have their own key, distinct from an operator block.
\dkDeclareSymbol[Spectral parameter lambda of a self-adjoint operator. In a finite-dimensional eigenvalue list lambda_j(A) it denotes an individual eigenvalue; in A = integral lambda dE(lambda), it is the real integration variable.]{SpectralValue}{\lambda}{wantedspectrum}{local}{Real spectral parameter or individual eigenvalue of a self-adjoint operator.}{scalar}



% Conventional mathematical notation. Hover text applies even when ink is plain.
% These are local (non-glossary) unless promoted explicitly by the author.
% IdentityTwo and IdentityFour derive their I from Identity. Singular-value
% registrations derive s_j from SingularValueBase; SingletonCost has its own s.
\dkDeclareSymbol[Identity operator on the vector or Hilbert space in context. Its dimension is inferred from the expression and need not be finite.]{Identity}{I}{plain}{local}{Identity operator on the space in context.}{identity-operator}
\dkDeclareSymbol[Two-by-two identity matrix, acting on a two-dimensional coordinate space.]{IdentityTwo}{\symglyph{Identity}_2}{plain}{local}{Two-by-two identity matrix.}{identity-operator}
\dkDeclareSymbol[Four-by-four identity matrix, acting on a four-dimensional coordinate space.]{IdentityFour}{\symglyph{Identity}_4}{plain}{local}{Four-by-four identity matrix.}{identity-operator}
\dkDeclareSymbol[Two-by-two zero matrix. Unlike the bold zero vector, this is a linear operator/matrix with two rows and two columns.]{ZeroMatrixTwo}{0_2}{plain}{local}{Two-by-two zero matrix.}{matrix}
\dkDeclareSymbol[The constant pi, the ratio of the circumference of a circle to its diameter. Angles in this article use radians.]{PiConstant}{\pi}{plain}{local}{The constant pi (angles are in radians).}{constant}
\dkDeclareSymbol[The field of real numbers.]{RealField}{\mathbb R}{plain}{local}{The real numbers.}{field}
\dkDeclareSymbol[The field of complex numbers.]{ComplexField}{\mathbb C}{plain}{local}{The complex numbers.}{field}
\dkDeclareSymbol[Generic Hilbert space used in the two-projection discussion; distinct from the paper's ambient space symbol E in the sine-theta theorem.]{GeneralHilbertSpace}{\mathcal H}{plain}{local}{Generic Hilbert space in a local discussion.}{hilbert-space}
\dkDeclareSymbol[The spectrum sigma(T) of an operator T: the set of complex scalars for which T minus the scalar identity is not boundedly invertible. For self-adjoint T the spectrum is real.]{SpectrumSymbol}{\sigma}{plain}{local}{Spectrum of an operator.}{operator-notation}
\dkDeclareSymbol[Base letter for singular values, presently s. It is distinct from the base s for matching singleton costs.]{SingularValueBase}{s}{plain}{local}{Base letter for singular values.}{notation-base}
\dkDeclareSymbol[Ordered singular value s_j(T), equal to the j-th eigenvalue of |T| for a finite matrix; distinguish this from the matching singleton cost s_j.]{SingularValueS}{\symglyph{SingularValueBase}_j}{plain}{glossary}{The j-th singular value, under the currently printed s_j convention.}{scalar}
\dkDeclareSymbol[Alternate indexed singular-value notation sigma_j(T). The current manuscript also uses s_j(T); this registered alias is available for a later deliberate notation decision.]{SingularValueSigma}{\sigma_j}{plain}{local}{Indexed singular value under the sigma_j convention.}{scalar}
\dkDeclareSymbol[Operator norm ||T||_op, the supremum of ||Tx|| over unit vectors x.]{OperatorNormNotation}{\|\cdot\|_{\mathrm{op}}}{plain}{glossary}{Operator norm.}{norm}
\dkDeclareSymbol[Frobenius norm ||T||_F, the square root of the sum of squared matrix entries, equal to the finite-dimensional Hilbert--Schmidt norm.]{FrobeniusNormNotation}{\|\cdot\|_F}{plain}{glossary}{Frobenius norm.}{norm}
\dkDeclareSymbol[Trace norm ||T||_1, the sum of singular values when finite.]{TraceNormNotation}{\|\cdot\|_1}{plain}{glossary}{Trace (nuclear) norm.}{norm}
\dkDeclareSymbol[Hilbert--Schmidt norm ||T||_HS, the square root of the sum of squared singular values when finite.]{HilbertSchmidtNormNotation}{\|\cdot\|_{\mathrm{HS}}}{plain}{local}{Hilbert--Schmidt norm.}{norm}
\dkDeclareSymbol[Ky Fan j-norm, the sum of the largest j singular values.]{KyFanNormNotation}{\|\cdot\|_{(j)}}{plain}{glossary}{Ky Fan j-norm.}{norm}
\dkDeclareSymbol[Interval/exterior spectral separation predicate I_delta(B,C), specifying a central closed interval separated by delta from the other operator's spectrum.]{IntervalSeparation}{\mathcal I}{plain}{local}{Interval/exterior separation predicate.}{predicate}
\dkDeclareSymbol[Ordered spectral separation predicate O_delta(B,C), requiring one self-adjoint operator to lie below a spectral threshold and the other above it by at least delta.]{OrderedSeparation}{\mathcal O}{plain}{local}{Ordered spectral separation predicate.}{predicate}
\dkDeclareSymbol[Positive parameter epsilon controlling a small perturbation in the two-dimensional eigenvalue example.]{EigenvalueExampleEpsilon}{\varepsilon}{plain}{local}{Small positive parameter in the eigenvalue-example matrices.}{scalar}
\dkDeclareSymbol[Positive lower bound epsilon for the trial-coordinate map E_0, as in E_0* E_0 >= epsilon squared times the identity.]{FrameLowerBoundEpsilon}{\varepsilon}{plain}{local}{Lower-bound parameter for the trial-coordinate embedding.}{scalar}
\dkDeclareSymbol[Rotation angle vartheta of a planar orthogonal block in the matching-proof first-variation calculation.]{MatchingPlanarAngle}{\vartheta}{plain}{local}{Planar rotation angle in the matching-proof calculation.}{angle}


% Ordinary named mathematical operators already used throughout the manuscript.
% Keeping these here makes their spelling a single editorial choice.
\dkDeclareSymbol[Range of a linear operator: the set of its output vectors.]{RangeOperator}{\operatorname{ran}}{plain}{local}{Range (image) of an operator.}{operator-notation}
\dkDeclareSymbol[Kernel (null space) of a linear operator: inputs mapped to zero.]{KernelOperator}{\operatorname{ker}}{plain}{local}{Kernel (null space) of an operator.}{operator-notation}
\dkDeclareSymbol[Domain of an operator: the vectors on which the operator is defined, important for unbounded operators.]{DomainOperator}{\operatorname{dom}}{plain}{local}{Domain of an operator.}{operator-notation}
\DeclareRobustCommand{\DKRange}{\sym{RangeOperator}}
\DeclareRobustCommand{\DKKernel}{\sym{KernelOperator}}
\DeclareRobustCommand{\DKDom}{\sym{DomainOperator}}

% Norm-formatting macros: the operand changes by occurrence but the notation
% lives here. The outer hover describes the norm and does not add ink color.
\DeclareRobustCommand{\dkOperatorNorm}[1]{\like{OperatorNormNotation}{\|#1\|_{\mathrm{op}}}}
\DeclareRobustCommand{\dkTraceNorm}[1]{\like{TraceNormNotation}{\|#1\|_1}}
\DeclareRobustCommand{\dkFrobeniusNorm}[1]{\like{FrobeniusNormNotation}{\|#1\|_F}}
\DeclareRobustCommand{\dkHilbertSchmidtNorm}[1]{\like{HilbertSchmidtNormNotation}{\|#1\|_{\mathrm{HS}}}}
\DeclareRobustCommand{\dkKyFanNorm}[2]{\like{KyFanNormNotation}{\|#1\|_{(#2)}}}
\DeclareRobustCommand{\dkSchattenNorm}[2]{\like{SchattenNormNotation}{\|#1\|_{S_{#2}}}}

\dkDeclareSymbol[Schatten p-norm ||T||_{S_p}, obtained from the l^p-norm of the singular values and finite for operators in the Schatten p-class.]{SchattenNormNotation}{\|\cdot\|_{S_p}}{plain}{local}{Schatten p-norm.}{norm}


\dkDeclareSymbol[Positive modulus of an operator T: |T| = (T* T)^(1/2).]{OperatorModulusNotation}{|\cdot|}{plain}{glossary}{Positive modulus of an operator.}{operator-notation}
\dkDeclareSymbol[Linear span of the indicated vectors, i.e. the smallest linear subspace that contains them.]{SpanOperator}{\operatorname{span}}{plain}{local}{Linear span of vectors.}{operator-notation}
\dkDeclareSymbol[Diagonal or block-diagonal matrix/operator with the displayed entries or blocks.]{DiagonalOperator}{\operatorname{diag}}{plain}{local}{Diagonal or block-diagonal operator.}{linear-operator}
\dkDeclareSymbol[Trace of a trace-class operator, or ordinary matrix trace in finite dimensions.]{TraceOperator}{\operatorname{tr}}{plain}{local}{Trace of an operator.}{operator-notation}
\dkDeclareSymbol[Dimension of a vector space or subspace.]{DimensionOperator}{\dim}{plain}{local}{Dimension of a vector space.}{operator-notation}
\dkDeclareSymbol[Rank of a linear map or matrix: dimension of its range.]{RankOperator}{\operatorname{rank}}{plain}{local}{Rank of a linear map or matrix.}{operator-notation}
\dkDeclareSymbol[Distance between sets of spectral values: the infimum of pairwise distances.]{SpectralDistanceOperator}{\operatorname{dist}}{plain}{local}{Distance between two spectral sets.}{operator-notation}
\DeclareRobustCommand{\dkModulus}[1]{\like{OperatorModulusNotation}{|#1|}}

% Context-specific glyphs previously hidden in variants of other semantic keys.
% Keeping both the roles and the exact printed forms makes collisions explicit.
\dkDeclareSymbol[Nonzero vector x tested in the Rayleigh quotient, before imposing unit length.]{RayleighTestVector}{x}{trialvector}{local}{Test vector in the Rayleigh quotient.}{vector}
\dkDeclareSymbol[Positive cosine block C in the generic two-projection decomposition of Halmos; C and S commute and C squared plus S squared is the identity.]{HalmosCosineBlock}{C}{anglecos}{local}{Cosine block in the Halmos two-projection representation.}{linear-operator}
\dkDeclareSymbol[Positive sine block S in the generic two-projection decomposition of Halmos; C and S commute and C squared plus S squared is the identity.]{HalmosSineBlock}{S}{sineoperator}{local}{Sine block in the Halmos two-projection representation.}{linear-operator}
\dkDeclareSymbol[Riesz spectral projection P_Gamma defined by a contour Gamma that encloses selected spectral values.]{RieszProjection}{P_{\symglyph{RieszContour}}}{wantedprojection}{local}{Spectral projection associated with a contour Gamma.}{orthogonal-projection}
\dkDeclareSymbol[Historical Sylvester right-hand coefficient matrix A, in B X - X A = Y. Distinct from the current paper's reference operator A.]{HistoricalSylvesterRight}{A}{trialspectrum}{local}{Right-hand coefficient A in the historical Sylvester equation.}{self-adjoint-operator}
\dkDeclareSymbol[Scalar coefficient b in the Sylvester prototype (b-a)x=y, with separated scalar spectral values a and b.]{SylvesterScalarLeft}{b}{unwantedspectrum}{local}{Scalar left coefficient b in a Sylvester equation.}{scalar}
\dkDeclareSymbol[Scalar coefficient a in the Sylvester prototype (b-a)x=y.]{SylvesterScalarRight}{a}{trialspectrum}{local}{Scalar right coefficient a in a Sylvester equation.}{scalar}
\dkDeclareSymbol[Scalar unknown x in the Sylvester prototype (b-a)x=y.]{SylvesterScalarUnknown}{x}{sylvester}{local}{Scalar unknown x in a Sylvester equation.}{scalar}
\dkDeclareSymbol[Scalar right side y in the Sylvester prototype (b-a)x=y.]{SylvesterScalarRHS}{y}{resid}{local}{Scalar right side y in a Sylvester equation.}{scalar}
\dkDeclareSymbol[Source subspace K of a real Euclidean space in the matching proof; an orthogonal transporter sends K onto L.]{MatchingSourceSubspace}{K}{trial}{local}{Source subspace K in the matching proof.}{linear-subspace}
\dkDeclareSymbol[Target subspace L of a real Euclidean space in the matching proof; an orthogonal transporter sends K onto L.]{MatchingTargetSubspace}{L}{wanted}{local}{Target subspace L in the matching proof.}{linear-subspace}
\dkDeclareSymbol[Orthogonal projection P_L onto the matching proof's target subspace L.]{MatchingTargetProjection}{P_L}{wantedprojection}{local}{Orthogonal projection onto the matching target L.}{orthogonal-projection}
\dkDeclareSymbol[Real two-dimensional orthogonal group O(2): planar rotations and reflections.]{PlanarOrthogonalGroup}{O(2)}{fg}{local}{Orthogonal matrices of size two.}{group}
\dkDeclareSymbol[Base letter R used to denote planar rotation operators R_t in the matching proof.]{PlanarRotationBase}{R}{fg}{local}{Base letter for planar rotations.}{notation-base}
\dkDeclareSymbol[Planar rotation R_(r+s) through the sum of scalar phases r and s.]{MatchingPlanarRotationSum}{\symglyph{PlanarRotationBase}_{r+s}}{fg}{local}{Planar rotation through phase r+s.}{unitary-operator}
\dkDeclareSymbol[Planar rotation R_gamma through phase gamma in the block rotation calculation.]{MatchingPlanarRotationGamma}{\symglyph{PlanarRotationBase}_{\symglyph{MatchingPhaseGamma}}}{fg}{local}{Planar rotation through phase gamma.}{unitary-operator}

% Specific symbols in the historical and appendix examples. They retain their
% existing black-ink notation but receive locally accurate hover definitions.
\dkDeclareSymbol[Contour Gamma, a positively oriented closed curve in the resolvent set enclosing a selected isolated spectral cluster.]{RieszContour}{\Gamma}{plain}{local}{Contour used in the Riesz spectral-projection formula.}{contour}
\dkDeclareSymbol[Even integer nu in a tested Ky Fan comparison, with nu/2 principal-plane terms on the right side.]{KyFanEvenIndex}{\nu}{plain}{local}{Even integer in the principal-plane Ky Fan comparison.}{scalar}
\dkDeclareSymbol[Orthogonal projection Omega_k onto the k-th principal plane in the published displacement comparison.]{PrincipalPlaneProjectionBase}{\Omega}{plain}{local}{Projection onto the indicated principal plane.}{notation-base}
\dkDeclareSymbol[Scalar gap threshold gamma in the ordered Sylvester spectral separation B >= (gamma+delta)I and C <= gamma I.]{SylvesterGapThreshold}{\gamma}{plain}{local}{Spectral separation threshold in a Sylvester estimate.}{scalar}
\dkDeclareSymbol[Projection Omega_t from a net or sequence that converges strongly to the identity in the singular-value approximation lemma.]{StrongApproxProjection}{\Omega_t}{plain}{local}{Orthogonal projection converging strongly to the identity.}{orthogonal-projection}
\dkDeclareSymbol[First orthogonal projection Omega used in the two-block pinching inequality for a unitarily invariant norm.]{PinchingProjectionOne}{\Omega}{plain}{local}{First projection in the pinching inequality.}{orthogonal-projection}
\dkDeclareSymbol[Second orthogonal projection Upsilon used in the two-block pinching inequality for a unitarily invariant norm.]{PinchingProjectionTwo}{\Upsilon}{plain}{local}{Second projection in the pinching inequality.}{orthogonal-projection}
\dkDeclareSymbol[Finite-rank projection Gamma used to retain the image of a chosen finite-rank compression K Psi.]{CompressionProjectionImage}{\Gamma}{plain}{local}{Range projection in the finite-rank compression lemma.}{orthogonal-projection}
\dkDeclareSymbol[Finite-rank projection Psi used to select input vectors in a nearly maximizing singular-value compression.]{CompressionProjectionInput}{\Psi}{plain}{local}{Domain projection in the finite-rank compression lemma.}{orthogonal-projection}
\dkDeclareSymbol[Small positive tolerance eta in the finite-rank compression approximation estimate.]{CompressionTolerance}{\eta}{plain}{local}{Tolerance controlling the off-diagonal compression norm.}{scalar}

\dkDeclareSymbol[Rotation phase gamma in a planar block from the matching proof.]{MatchingPhaseGamma}{\gamma}{plain}{local}{Rotation phase gamma in the matching calculation.}{scalar}

% Two common indices with local mathematical meanings distinct from other deltas.
\dkDeclareSymbol[Kronecker delta: delta_ij is one when the indices i and j agree and zero otherwise. This is not the spectral-gap parameter delta.]{KroneckerDelta}{\delta}{plain}{local}{Kronecker delta: equals 1 for equal indices and 0 otherwise.}{scalar}
\dkDeclareSymbol[Sign variable epsilon taking the values -1 and +1 in the factorization over two phase choices.]{PhaseSignChoice}{\epsilon}{plain}{local}{Sign variable taking values minus one and plus one.}{scalar}

% Documented notation correspondence. The source column is a faithful
% transcription of each cited original context, NOT an asserted universal DK
% meaning. The original A_0 is deliberately listed for two distinct roles:
% its reference block and the comparison/trial operator. The source uses the
% same printed symbol for both; this paper is free to distinguish them.
\dkDeclareNotationSource{DK1970}{Davis and Kahan (1970), The Rotation of Eigenvectors by a Perturbation. III}{https://doi.org/10.1137/0707001}
\dkDeclareNotationSource{Seelmann2019}{Seelmann (2019), Semidefinite Perturbations in the Subspace Perturbation Problem}{https://arxiv.org/abs/1708.02463}
\dkDeclareNotationSource{Blechta2020}{Blechta (2020), Stability of Linear GMRES Convergence with Respect to Compact Perturbations}{https://arxiv.org/abs/2005.12960}
% Arguments: {key}{source ID}{source glyph}{section or equation}{role-specific qualification}
\dkDeclareNotationAlternative{A}{DK1970}{A}{\S1, (1.2)}{Reference self-adjoint operator.}
\dkDeclareNotationAlternative{H}{DK1970}{H}{\S1, (1.2)}{Additive perturbation.}
\dkDeclareNotationAlternative{Acurrent}{DK1970}{A+H}{\S1}{Perturbed operator.}
\dkDeclareNotationAlternative{E0}{DK1970}{E_0}{\S1, (1.1)}{Source isometric frame.}
\dkDeclareNotationAlternative{E1}{DK1970}{E_1}{\S1, (1.1)}{Reference complement frame.}
\dkDeclareNotationAlternative{F0}{DK1970}{F_0}{\S1}{Exact-target frame.}
\dkDeclareNotationAlternative{F1}{DK1970}{F_1}{\S1}{Exact-target complement frame.}
\dkDeclareNotationAlternative{ReferenceBlock0}{DK1970}{A_0}{\S1, (1.2)}{Reference source block.}
\dkDeclareNotationAlternative{ReferenceBlock1}{DK1970}{A_1}{\S1, (1.2)}{Reference complement block.}
\dkDeclareNotationAlternative{A0}{DK1970}{A_0}{\S1, residual discussion}{Residual comparison operator (the reference block in DK).}
% We deliberately leave Lambda0 and Lambda1 without DK source-glyph metadata:
% the available source transcription does not reliably preserve the source
% accents/decorations on the perturbed spectral-block symbols.
\dkDeclareNotationAlternative{Psource}{DK1970}{P}{\S1}{Reference/source projection.}
\dkDeclareNotationAlternative{Qtarget}{DK1970}{Q}{\S1}{Target projection.}
\dkDeclareNotationAlternative{Pexact}{DK1970}{Q}{\S1}{Exact-target projection.}
\dkDeclareNotationAlternative{B}{DK1970}{B}{\S1, (1.2)}{Off-diagonal coupling block.}
\dkDeclareNotationAlternative{Theta0}{DK1970}{\Theta_0}{\S1}{Source-side angle operator.}
\dkDeclareNotationAlternative{Theta1}{DK1970}{\Theta_1}{\S1}{Complementary angle operator.}
\dkDeclareNotationAlternative{Theta}{DK1970}{\Theta}{\S1}{Full angle operator.}
\dkDeclareNotationAlternative{R}{DK1970}{R}{\S1, residual discussion}{Trial residual.}
\dkDeclareNotationAlternative{Rrotation}{DK1970}{U}{\S1, (1.18)}{Canonical direct rotation.}
\dkDeclareNotationAlternative{Wcompetitor}{DK1970}{V}{\S4}{Competing unitary transporter.}
\dkDeclareNotationAlternative{delta}{DK1970}{\delta}{\S2}{Spectral separation parameter.}
% Contemporary alternatives are recorded separately and do not enter the DK
% correspondence table. Their citation locations must be individually checked.
\dkDeclareNotationAlternative{Eambient}{Seelmann2019}{\mathcal H}{p.~1}{Ambient Hilbert space.}
\dkDeclareNotationAlternative{SingularValueS}{Blechta2020}{\sigma_j}{Thm.~1.1}{Ordered singular values of a compact operator.}

\ifDKColoredSymbols\else\DKMathColorsfalse\fi
\ifDKAuthorStyle\else\DKSymbolTooltipsfalse\fi




\usepackage[most]{tcolorbox}
\usepackage{listings}
\usepackage{listingsutf8}
\usepackage{fontspec}
\IfFileExists{emoji.sty}{%
  \IfFontExistsTF{Noto Color Emoji}{\usepackage{emoji}}{}%
}{}
\providecommand{\emoji}[1]{\textnormal{[#1]}} % Review notes also work without an emoji font.

\setlength{\emergencystretch}{3em}
\raggedbottom

% ArXiv-friendly Lean code: no shell escape, minted, or external highlighter.
\definecolor{leankeyword}{HTML}{243B6B}
\definecolor{leantactic}{HTML}{70477A}
\definecolor{leantype}{HTML}{1F6473}
\definecolor{leancomment}{HTML}{39734D}
\definecolor{leanstring}{HTML}{8A4B2A}
\definecolor{leannumber}{HTML}{777777}
\definecolor{leanbackground}{HTML}{F7F7F3}
\lstdefinelanguage{lean}{%
  sensitive=true,
  morecomment=[l]{--},
  morecomment=[s]{/-}{-/},
  morestring=[b]",
  morekeywords=[1]{import,open,namespace,section,end,noncomputable,variable,variables,
    universe,universes,set_option,attribute,axiom,opaque,theorem,lemma,def,abbrev,
    example,instance,structure,class,inductive,notation,local,deriving,extends,
    by,where,let,have,show,fun,match,with,if,then,else,do,return,in,forall},
  morekeywords=[2]{simp,simpa,rw,erw,dsimp,unfold,ring,ring_nf,nlinarith,linarith,
    positivity,norm_num,field_simp,omega,exact,refine,apply,assumption,intro,intros,
    ext,funext,cases,rcases,obtain,constructor,left,right,classical,by_contra,
    contrapose,convert,using,suffices,calc},
  morekeywords=[3]{Prop,Type,Nat,Int,Real,Fin,Set,Submodule,Matrix,EuclideanSpace,
    LinearMap,ContinuousLinearMap,LinearPMap,RCLike,IsTrialResidual,
    IsExactSpectralDecomposition,FormBoundedSylvesterGap,
    NormalizedSymmetricOperatorIdealFamily},
}
\lstdefinestyle{leanpaper}{%
  language=lean,
  columns=fullflexible,
  backgroundcolor=\color{leanbackground},
  basicstyle=\ttfamily\scriptsize,
  commentstyle=\color{leancomment}\itshape,
  keywordstyle=[1]\color{leankeyword}\bfseries,
  keywordstyle=[2]\color{leantactic},
  keywordstyle=[3]\color{leantype},
  stringstyle=\color{leanstring},
  numberstyle=\tiny\color{leannumber},
  breakatwhitespace=false,
  breaklines=true,
  captionpos=b,
  keepspaces=true,
  numbers=left,
  numbersep=6pt,
  showspaces=false,
  showstringspaces=false,
  showtabs=false,
  tabsize=2,
  aboveskip=0.8em,
  belowskip=0.8em,
  literate=
    {\\KField}{{$\Bbbk$}}1
    {∃}{{$\exists$}}1
    {∧}{{$\land$}}1
    {∀}{{$\forall$}}1
    {≠}{{$\neq$}}1
    {⊤}{{$\top$}}1
}
\crefname{lstlisting}{formalization}{formalizations}
\Crefname{lstlisting}{Formalization}{Formalizations}
\renewcommand{\lstlistingname}{Formalization}

\theoremstyle{thmstyleone}
\newtheorem{theorem}{Theorem}
\newtheorem{proposition}[theorem]{Proposition}
\newtheorem{lemma}[theorem]{Lemma}
\newtheorem{remark}[theorem]{Remark}
\newtheorem{corollary}[theorem]{Corollary}

\newcommand{\Lean}{\textsc{Lean}}
\newcommand{\DK}{Davis--Kahan}
\newcommand{\TauCeti}{Tau Ceti}
% Use \ie{} and \eg{}; supply surrounding punctuation at the call site.
\newcommand{\ie}{i.e.\@}
\newcommand{\eg}{e.g.\@}
\newcommand{\DKResultCount}{29}
\newcommand{\DKProvedCount}{28}
\newcommand{\DKRefutedCount}{1}

% Usage: \margincomment{Check this formula!}
%\newcommand{\jonsays}[1]{%
%  \marginpar{%
%    \raggedright\color{red}\tiny% Change to \footnotesize if text is too small
%    \emoji{person-raising-hand} \textbf{Comment:}\\#1%
%  }%
%}

% Usage: \margincomment{Your comment that now wraps perfectly!}
\newcommand{\jonsays}[1]{%
  \marginpar{%
    \raggedright% Forces text to wrap beautifully without ugly hyphenation
    \color{red}\scriptsize% Using \scriptsize makes the text tighter
    \emoji{person-raising-hand}: \\#1%
  }%
}
%\newcommand{\jonsays}[1]{{\color{red}\small\textbf{[\emoji{person-raising-hand}: #1]}}}

% Usage: \margincomment{Your comment that now wraps perfectly!}
%\newcommand{\agentsays}[1]{%
%  \marginpar{%
%    \raggedright% Forces text to wrap beautifully without ugly hyphenation
%    \color{red}\scriptsize% Using \scriptsize makes the text tighter
%    \emoji{robot}: \\#1%
%  }%
%}
\newcommand{\agentsays}[1]{{\color{red}\small\textbf{[\emoji{robot}: #1]}}}

\newcommand{\jonrsp}[1]{{\color{red}\small\textbf{[\emoji{person-raising-hand}: #1]}}}

\title[Formalizing Davis--Kahan Perturbation Theory]{Formalizing Davis--Kahan Perturbation Theory in Lean 4}
% Additional authors are expected; keep the confirmed ordering below stable.
\ifDKAnonymous
  \author{Anonymous authors}
\else
  \author*[1]{\fnm{Jonathan} \sur{Crall}}\email{jon.crall@kitware.com}
  \author[2]{\fnm{Edward} \sur{Wang}}\email{ewang43@jhu.edu}
  \author[1]{\fnm{Brian} \sur{Hu}}\email{brian.hu@kitware.com}
  \author[2]{\fnm{Carey E.} \sur{Priebe}}\email{cep@jhu.edu}
  \affil[1]{\orgname{Kitware, Inc.},
    \orgaddress{\city{Clifton Park}, \state{NY}, \country{USA}}}
  \affil[2]{\orgdiv{Department of Applied Mathematics and Statistics},
    \orgname{Johns Hopkins University},
    \orgaddress{\city{Baltimore}, \state{MD}, \country{USA}}}
\fi

\abstract{We formalize the spectral-subspace perturbation theory of Davis and
Kahan's 1970 paper in Lean~4. Of \DKResultCount{} selected results, we prove
\DKProvedCount{} under their stated hypotheses and refute the remaining
printed statement. The development includes the sine and tangent theorems,
direct rotations, principal-angle operators, Sylvester separation, and
operator-ideal norms on real and complex Hilbert spaces, including unbounded
self-adjoint operators. We explain the mathematical lineage connecting these
constructions. The refuted statement concerns minimization of full ambient
displacement by direct rotation. We give real four-dimensional counterexamples
at every prescribed positive acute angle and identify the failed
singular-value comparison. A corrected optimality statement holds for Q-norms,
whose squares are unitarily invariant norms of the squared modulus, including
the Frobenius norm. The formalization builds on Mathlib and adapts Spectra's
projection-valued-measure and Borel-calculus constructions. Portions of the
operator theory developed here were subsequently adapted into Tau Ceti.
These foundations also support an exact reduction of trace-norm optimal real
subspace transport to weighted matching on principal angles.}

\keywords{Davis--Kahan theory, spectral perturbation, formal verification,
Lean 4, direct rotation, principal angles}
\pacs[MSC Classification (2020)]{68V20, 47A55, 15A42}

\begin{document}
\maketitle

\ifDKAnonymous\else
\enlargethispage{2\baselineskip}
\vspace{-18pt}
\noindent\begin{minipage}{\linewidth}\footnotesize
\noindent\textbf{Acknowledgments.}
This material is based upon work supported by the Defense Advanced Research
Projects Agency (DARPA) under Contract No. HR001125CE017. Any opinions, findings
and conclusions or recommendations expressed in this material are those of the
author(s) and do not necessarily reflect the views of the Defense Advanced
Research Projects Agency.
\end{minipage}
\fi
\clearpage

\section{Introduction}
\label{sec:introduction}

An approximate eigenvector gives a candidate direction for the action of an
operator. If a unit vector $\sym{v}$ is paired with an approximate eigenvalue $\sym{mu}$,
the residual $\sym{A}\sym{v}-\sym{mu} \sym{v}$ measures how closely it satisfies the eigenvector
equation. A small residual can identify an accurate approximation to an
invariant subspace even when repeated or nearby eigenvalues make individual
eigenvectors difficult to distinguish. Davis--Kahan theory bounds the angles
between approximate and exact spectral subspaces in terms of a residual or
perturbation and a separation between the relevant spectra
\citep{davis1970rotation}.

In applications, one often works with a perturbed operator
  $\sym{Acurrent}=\sym{A}+\sym{H}$ while the reference operator $\sym{A}$ is
  unknown.
Taking an eigenspace of $\sym{Acurrent}$ as a trial subspace for $\sym{A}$, the norm of its
  residual relative to $\sym{A}$ is bounded by the norm of $\sym{H}$, allowing
  Davis--Kahan theory to bound the angle between the two subspaces without
  requiring $\sym{A}$ to be known or its residuals to be evaluated directly.

Our project began with a variant of this bound needed for
\ifDKAnonymous a spectral-estimation application\else the data kernel
perspective space (DKPS) \citep{helm2026query}\fi.

We extended the formalization to the results established in the original 1970
paper, including real and complex Hilbert spaces, unbounded self-adjoint
operators, and unitarily invariant operator-ideal norms. In particular, the
spectral-separation hypotheses allow half-infinite intervals
(e.g.\ $(-\infty,a]$). These settings require operator domains, spectral
calculus, and norm estimates beyond finite matrices.

We formalized the four headline theorems \citep[Section~2]{davis1970rotation} and every
named theorem, proposition, lemma, and corollary established in the 1970 paper.

This gives \DKResultCount{} targets, enumerated in
Appendix~\ref{app:dk-inventory}. We prove \DKProvedCount{} of them and refute
the printed statement \citep[Proposition~4.4]{davis1970rotation}.
That proposition asserts a full-displacement
optimality property for real direct rotations. We give a family of real
four-dimensional examples exhibiting trace-norm failure at every prescribed
positive acute angle and identify
the invalid singular-value comparison in its proof. The phase construction
behind the counterexample occurs over the complex field in the authors'
example \citep[Example~4.2]{davis1970rotation}. To the best of our knowledge,
its contradiction of the printed real, small-angle claim has not previously
been made explicit.
For the constructed pair, whose two principal angles both equal
$\sym{theta}$, the competing transporter attains the global trace-norm minimum
$4\sin\sym{theta}$ for the constructed subspaces. Section~\ref{sec:prop44}
compares these statements and repairs the optimality claim by restricting
the norm class to Q-norms (N is a Q-norm when $N(T)^2 = M(T^* T)$ for some
unitarily invariant norm $M$). This repair follows from the squared-displacement
comparison \citep[Proposition~4.3]{davis1970rotation};
earlier Q-norm minimization results appear in \citet{bolotnikov2001minimal}.

% TODO (lineage): Consolidate the scattered lineage records through the
% maintained dev/tauceti/dkps-upstream-contribution-census.md and the evidence
% index in AUTHORING_NOTES.md. Audit each attribution against an immutable
% donor commit/path, the accepted upstream commit/PR, and the corresponding
% declaration statements and proofs. Preserve Mathlib prerequisites, Spectra
% Spectra-derived constructions where applicable, and independent work as distinct relationships.
% Courant--Fischer/Weyl checked 2026-10-28: developed here in June 2026;
% August donor df78246f965881705c9301f96feddc270c98ceb1 predates the
% accepted Tau Ceti port 7ebf17ba04725c1a575ecf3d099ce0630cd949a0 (#13599).
% See COURANT_FISCHER_LINEAGE_AUDIT.md for the proof comparison and sources.

The proofs required us to formalize a significant amount of operator
theory, namely: principal angles and direct rotations, approximation numbers
and symmetric operator ideals, Borel spectral calculus, and Sylvester
estimates for unbounded operators. These developments build on Mathlib and
adapt Spectra's projection-valued-measure and Borel-calculus constructions
\citep{mathlib2020,bornemann2026spectra}.
These operator-theory foundations include our
Courant--Fischer and Weyl proofs, rectangular unitarily invariant seminorms,
and projection-valued measures, which were subsequently adapted into Tau
Ceti \citep{tauceti2026}.
Section~\ref{sec:related} explains the mathematical history and formal
provenance; Section~\ref{sec:foundations} describes the foundations used in
our proofs.

The same principal-angle and operator-ideal results support a further
optimization theorem. In Section~\ref{sec:example-transport}, we prove that
minimizing trace-norm displacement over all real orthogonal maps carrying
one subspace onto another reduces to weighted matching on principal angles.

\noindent\textbf{Contributions.}
\begin{enumerate}
\item \textbf{Davis--Kahan formalization.} We prove 28 of the 29 source
results and refute the full-displacement assertion
\citep[Proposition~4.4]{davis1970rotation}. The sine-theta theorem covers real and
complex Hilbert spaces, unbounded self-adjoint operators, the source's
spectral-separation alternatives, and operator-ideal norms where finite
(\cref{sec:coverage,sec:sine-theta}).
\item \textbf{Formalized operator-theoretic foundations.} We formalized the
geometric and analytic results required by these theorems, adapted existing
formal mathematics, and contributed foundations later used in Tau Ceti
(\cref{sec:related-formal,sec:foundations}).
\item \textbf{Source correction.} We identify a false comparison
of Ky Fan sums in the printed claim \citep[Proposition~4.4]{davis1970rotation},
give a real counterexample at every positive acute angle,
and prove the Q-norm consequence of the valid squared-displacement theorem
(\cref{sec:prop44}).
\item \textbf{A reduction to weighted matching.} We reduce unrestricted
real trace-norm transport to a weighted-matching problem on principal angles
(\cref{sec:example-transport}).

\end{enumerate}

\section{Mathematical background}
\label{sec:related}

The Davis--Kahan angle estimates use three constructions. Principal angles
measure the error between subspaces; a residual measures how well a proposed
subspace represents the operator; and a Sylvester estimate bounds the first
quantity by the second under spectral separation. Operator ideals allow the
same argument to express several norms. We describe the development of these
constructions before stating the formalized sine-theta theorem.

Existing historical surveys cover these developments. Paige and Wei trace
the cosine--sine decomposition from Jordan's subspace geometry through the
principal ideas assembled by Davis and Kahan to Stewart's later completion
and explicit formulation \citep{paige1994history}.
Simon's historical survey of Kato's work in non-relativistic quantum mechanics
traces Sz.-Nagy's unitary construction for pairs of projections and Kato's
later algebraic formulation \citep[\S5]{simon2018kato}.
We use this account for the history of the projection transport described in
\cref{sec:related-projections}.
B\"ottcher and Spitkovsky survey pairs of projections and their operator algebras.
Their historical discussion relates Halmos's 1969 block representation to the
earlier formulations of Dixmier (1948) and Davis (1958), and to the related
projection constructions in Davis and Kahan's 1970 paper
\citep[Remark~1.3]{bottcher2010gentle}.
We use these accounts together with the primary
sources. Table~\ref{tab:historical-lineage} summarizes the mathematical
constructions used in our formalization.

\begin{table}[!htbp]
\centering
\small
\setlength{\tabcolsep}{5pt}
\renewcommand{\arraystretch}{1.18}
\begin{tabularx}{\linewidth}{@{}>{\raggedright\arraybackslash}p{0.26\linewidth}>{\raggedright\arraybackslash}X>{\raggedright\arraybackslash}p{0.28\linewidth}@{}}
\toprule
\textbf{Lineage} & \textbf{Mathematical role} & \textbf{Formal dependency}\\
\midrule
Jordan (1875); Davis (1958); Halmos (1969) & Principal angles, invariant subspace geometry,
  and pairs of projections & Angle operators; projection blocks\\
Rayleigh (1877); Ritz (1909); Weyl (1912) & Approximate eigenpairs,
  variational spectra, and eigenvalue stability & Trial coordinates;
  min--max and eigenvalue inequalities\\
Rellich (1937--1942); Sz.-Nagy (1947); Kato (1966) & Moving spectral projections and unitary
  transport between nearby subspaces & Spectral calculus; polar
  direct rotation\\
Sylvester (1884); Rosenblum (1956, 1969) & Spectral separation controls an
  off-diagonal operator & Domain-aware Sylvester estimates\\
Von Neumann (1937); Schatten (1950); Ky Fan (1951) & Comparison of spectral errors
  through singular-value gauges & Approximation numbers;
  symmetric operator ideals\\
Davis (1963, 1965); Davis--Kahan (1969, 1970) & Residual-to-angle bounds, tangent estimates,
  rotation extremality, and spectral branch selection & Source-level
  sine, tangent, and direct-rotation theorems\\
\bottomrule
\end{tabularx}
\caption{Mathematical sources used in the Davis--Kahan formalization.
Dates identify the publications discussed in the text.}
\label{tab:historical-lineage}
\end{table}
\FloatBarrier

\subsection{Principal angles}
\label{sec:related-classical}

Camille Jordan studied the geometry of linear subspaces in $n$ dimensions in
  1875 \citep{jordan1875essai}.

For two lines, a single acute angle describes their relative position.
For subspaces of higher dimension, one needs a sequence of \emph{principal
  angles}.
The successive principal vectors are selected under orthogonality constraints,
  so their angles are independent of the coordinates used to describe either
  subspace.
This separates the geometry of a subspace from the arbitrary choice of its
  orthonormal basis.

For finite-dimensional subspaces $\sym{Usource}$ and $\sym{Vtarget}$ of equal dimension, choose
isometries $\sym{Xframe}$ and $\sym{Yframe}$ with those ranges.
The singular values of the overlap matrix are the principal cosines:
\begin{equation}
  \sym{SingularValue}(\like{Xframe}{\symglyph{Xframe}^*}\sym{Yframe})=\cos\like{theta}{\symglyph{theta}_j},
  \qquad 0\leq\like{theta}{\symglyph{theta}_j}\leq\frac{\sym{PiConstant}}{2}.
  \label{eq:history-principal-cosines}
\end{equation}

The corresponding principal vectors $\like{SourceVector}{\symglyph{SourceVector}_j}\in \sym{Usource}$ and $\like{TargetVector}{\symglyph{TargetVector}_j}\in \sym{Vtarget}$ satisfy
$\langle \like{SourceVector}{\symglyph{SourceVector}_i},\like{TargetVector}{\symglyph{TargetVector}_j}\rangle=\like{KroneckerDelta}{\symglyph{KroneckerDelta}_{ij}}\cos\like{theta}{\symglyph{theta}_j}$ in suitably oriented
coordinates.  For nonzero acute angles, each pair determines a real
two-dimensional plane in which the relative position is an ordinary
rotation.  An entire pair of subspaces is described by these planes,
together with their shared and orthogonal components.

Principal angles also arise in statistics: Hotelling's 1936 canonical
  correlation analysis sought pairs of maximally correlated linear combinations
  of two sets of variables \citep{hotelling1936relations}.
After appropriate whitening, the canonical correlations are singular values of a
  cross-coordinate map.
In numerical linear algebra, Bj\"orck and Golub developed algorithms for
  computing principal angles and vectors in 1973 \citep{bjorck1973numerical}.
Their computational formulation makes the basis-invariant angles available from
  orthonormal matrix factorizations.

A further formulation uses the cosine--sine decomposition of a
partitioned orthogonal or unitary matrix.
The blocks of the decomposition contain compatible diagonal cosine and sine
  factors.
Davis and Kahan describe the geometry using orthogonal projections,
angle operators, and the direct rotation carrying one subspace onto the other
\citep[Section~3]{davis1970rotation}. Their work assembled the principal ideas
of the cosine--sine decomposition; Stewart completed and explicitly formulated
the matrix decomposition in 1977
\citep{stewart1977perturbation,paige1994history}.
Paige and Wei explain how this decomposition and direct rotation represent
the same principal-angle data
\citep{paige1994history}. Our formalization uses projections and functional
calculus for the Hilbert-space results and principal-plane coordinates for
the finite-dimensional counterexample.

\subsection{Eigenvalue approximation}
\label{sec:related-ritz-weyl}

Jordan's angles answered how to compare subspaces once they are given.
An independent question came from mechanics and numerical analysis:
how can one approximate an eigenvalue or eigenmode when the underlying
differential equation cannot be solved exactly?  In the nineteenth
century Rayleigh estimated vibration frequencies using trial shapes
and their energy ratios \citep{rayleigh1877sound}.  Ritz developed a systematic finite-dimensional
variational approximation in the work published in 1909
\citep{ritz1909variations,leissa2005historical}.  The method selects
a trial space and studies the action of the operator through that
space. A trial space is a chosen space of candidate approximate solutions.
For an eigenvalue problem, one searches for an approximate eigenvector among
its vectors. This reduces an infinite-dimensional problem to a smaller
problem whose residual can then be assessed.

For a Hermitian matrix $\sym{A}$, the Rayleigh quotient
$\langle \sym{A}\sym{RayleighTestVector},\sym{RayleighTestVector}\rangle/\langle \sym{RayleighTestVector},\sym{RayleighTestVector}\rangle$ expresses the scalar action
of $\sym{A}$ along a trial vector. A normalized trial vector $\sym{v}$ gives the
Rayleigh--Ritz approximate eigenvalue $\sym{mu}=\langle \sym{A}\sym{v},\sym{v}\rangle$ and the
residual $\sym{A}\sym{v}-\sym{mu} \sym{v}$. The Courant--Fischer min--max
principle organizes the eigenvalues through extremal quotients over
subspaces \citep{bhatia1997matrix}.  Weyl's eigenvalue comparison
inequalities then control the motion of ordered eigenvalues under a
Hermitian perturbation \citep{weyl1912eigenwerte,bhatia1997matrix}.  In modern notation,
if $\sym{A}$ and $\sym{Acurrent}=\sym{A}+\sym{H}$ are Hermitian $n$-by-$n$ matrices and their eigenvalues
are arranged in the same order, then
\begin{equation}
  \max_j |\like{SpectralValue}{\symglyph{SpectralValue}_j}(\sym{A})-\like{SpectralValue}{\symglyph{SpectralValue}_j}(\sym{Acurrent})|\leq\dkOperatorNorm{\sym{H}}.
  \label{eq:history-weyl}
\end{equation}
This is an estimate on \emph{eigenvalues}; it does not tell us which
individual eigenvector of $\sym{A}$ corresponds to which eigenvector of $\sym{Acurrent}$.

The limitation already appears in two dimensions.  For any $\sym{EigenvalueExampleEpsilon}>0$,
put
\begin{equation}
 \like{A}{\symglyph{A}_\varepsilon}=\begin{pmatrix}0&0\\0&\sym{EigenvalueExampleEpsilon}\end{pmatrix},
 \qquad
 \like{Acurrent}{\symglyph{Acurrent}_\varepsilon}=\begin{pmatrix}\sym{EigenvalueExampleEpsilon}&0\\0&0\end{pmatrix}.
 \label{eq:history-swap-example}
\end{equation}
The matrices have identical ordered eigenvalues and
$\dkOperatorNorm{\like{A}{\symglyph{A}_\varepsilon}-\like{Acurrent}{\symglyph{Acurrent}_\varepsilon}}=\sym{EigenvalueExampleEpsilon}$,
yet their lowest-eigenvalue lines are perpendicular. The gap between the
two eigenvalues of each matrix is $\sym{EigenvalueExampleEpsilon}$, equal to the perturbation
norm. An eigenspace estimate therefore needs to compare the perturbation
size with the spectral separation.  For a cluster with repeated eigenvalues, an
entire invariant subspace is likewise more natural than a selected
orthonormal eigenbasis.

Rellich developed a rigorous perturbation theory of self-adjoint operators
in papers published from 1937 to 1942. Kato's 1966 book gives a systematic
account \citep{simon2018kato,kato1966perturbation}. The Davis--Kahan residual
bounds address a related approximation question: given a trial space and its
residual, spectral separation bounds its angle from an exact spectral
subspace. The estimate can therefore assess the accuracy of a
Rayleigh--Ritz approximation.

\subsection{Two orthogonal projections}
\label{sec:related-two-projections}

For a closed subspace of a Hilbert space $\sym{GeneralHilbertSpace}$, the orthogonal
projection remains a bounded operator even when the subspace is
infinite dimensional.  Denote the projections onto two closed
subspaces by $\sym{Psource}$ and $\sym{Qtarget}$. Compressing an operator to $\DKRange(\sym{Psource})$ means
restricting its input to that subspace and projecting its output back with
$\sym{Psource}$. The positive compression of $\sym{Qtarget}$ is therefore $\sym{Psource}\sym{Qtarget}\sym{Psource}$ restricted to
$\DKRange(\sym{Psource})$. It is the operator version of the squared overlap matrix: in finite dimensions its eigenvalues are
$\cos^2\like{theta}{\symglyph{theta}_j}$.  In Hilbert space it need not have a discrete
spectrum.  Functional calculus nevertheless defines its positive
square root and the associated cosine and sine of an angle operator.

Write $\sym{Theta0}$ for this directed angle operator on
$\DKRange(\sym{Psource})$, and $\sym{Theta1}$ for the corresponding
operator for the complementary pair $\sym{Identity}-\sym{Psource},\sym{Identity}-\sym{Qtarget}$.
The full angle operator is
\[
 \sym{Theta}=\sym{DiagonalOperator}(\sym{Theta0},\sym{Theta1})
 \quad\text{on }\DKRange(\sym{Psource})\oplus\DKKernel(\sym{Psource}).
\]
For $c>0$, the comparisons $\sym{Theta}\le c$ and $\sym{Theta}<c$ mean that its
operator norm is at most, or strictly less than, $c$, respectively.
In finite dimensions this norm is the largest eigenvalue of either directed
angle block, comparing the source with the target and their orthogonal
complements.
This includes right angles from unmatched directions: if a source vector is
orthogonal to the target, its cosine compression is zero and its directed
angle is $\sym{PiConstant}/2$. The compression construction applies even when the subspace
dimensions differ. Target directions orthogonal to the source appear as
right angles in the complementary block. Later, isometric frames express these angle blocks on
coordinate Hilbert spaces.

The relative position of two closed subspaces was studied by Dixmier
in 1948 and by Davis in his 1958 paper \emph{Separation of Two Linear
Subspaces} \citep{dixmier1948position,davis1958separation}.  Davis's
closeness and separation operators place the cosine/sine data in an
operator-theoretic setting.  Halmos's 1969 \emph{Two Subspaces} gives a
block representation of the two projections
\citep{halmos1969two}.  One first splits off the four intersection
spaces corresponding to the combinations of membership in the ranges
and kernels of $\sym{Psource}$ and $\sym{Qtarget}$.  On their remaining, generic part, the pair
is unitarily equivalent to
\begin{equation}
 \like{Psource}{\symglyph{Psource}_{\mathrm{gen}}}=\begin{pmatrix}\sym{Identity}&0\\0&0\end{pmatrix},
 \qquad
 \like{Qtarget}{\symglyph{Qtarget}_{\mathrm{gen}}}=\begin{pmatrix}\like{HalmosCosineBlock}{\symglyph{HalmosCosineBlock}^2}&\sym{HalmosCosineBlock}\sym{HalmosSineBlock}\\\sym{HalmosCosineBlock}\sym{HalmosSineBlock}&\like{HalmosSineBlock}{\symglyph{HalmosSineBlock}^2}\end{pmatrix},
 \qquad \like{HalmosCosineBlock}{\symglyph{HalmosCosineBlock}^2}+\like{HalmosSineBlock}{\symglyph{HalmosSineBlock}^2}=\sym{Identity},
 \quad \sym{HalmosCosineBlock}\sym{HalmosSineBlock}=\sym{HalmosSineBlock}\sym{HalmosCosineBlock},
 \label{eq:history-halmos-block}
\end{equation}
where $\sym{HalmosCosineBlock}$ and $\sym{HalmosSineBlock}$ are commuting positive contractions on a common
Hilbert space.  Thus an operator-valued angle replaces a finite list
of principal angles.  B\"ottcher and Spitkovsky survey this
representation and its consequences for the algebra generated by two
projections \citep{bottcher2010gentle}.

Among these four intersections, the two right-angle intersection
subspaces (called defect spaces below) are
\[
 \sym{CrossedSourceDefect}=\DKRange(\sym{Psource})\cap\DKKernel(\sym{Qtarget}),
 \qquad
 \sym{CrossedTargetDefect}=\DKKernel(\sym{Psource})\cap\DKRange(\sym{Qtarget}).
\]
These contain source directions orthogonal to the target, and target directions
orthogonal to the source, respectively. The pair is called \emph{acute} when
both are $\{\sym{ZeroVector}\}$, the zero subspace consisting only of the zero
vector. For a general pair these spaces affect whether a specified
intertwining operator exists. An operator $\sym{IntertwinerT}$ intertwines $\sym{Psource}$ and $\sym{Qtarget}$ when
$\sym{IntertwinerT}\sym{Psource}=\sym{Qtarget}\sym{IntertwinerT}$: applying $\sym{IntertwinerT}$ after the source projection agrees with applying the
target projection after $\sym{IntertwinerT}$. In particular, $\sym{IntertwinerT}$ maps the source range into
the target range and its orthogonal complement into the target complement.
The existence criteria for direct rotations include conditions on the
dimensions of these defect spaces \citep[Section~3]{davis1970rotation}.

A finite-dimensional orthogonal matrix decomposes into invariant real
lines and two-dimensional rotation planes.  In a general Hilbert space an
angle operator can instead have continuous spectrum.  The spectral
calculus of the pair of projections takes the place of a finite list
of angles, and the intersection and defect spaces cannot be discarded.
Our Hilbert-space direct-rotation proofs therefore use projection identities
and functional calculus. In finite dimensions these identities recover the
principal-plane construction.

\subsection{Transport between nearby spectral subspaces}
\label{sec:related-projections}

A perturbation can change the invariant subspace associated with an isolated
part of a self-adjoint operator's spectrum. Let $\sym{Psource}$ and
$\sym{Qtarget}$ be the orthogonal projections onto the corresponding spectral
subspaces before and after perturbation. To compare these subspaces, we seek
a unitary carrying $\DKRange(\sym{Psource})$ onto
$\DKRange(\sym{Qtarget})$ that depends continuously, or analytically, on the
projections.
Sz.-Nagy developed the nearby-projection construction in work
published in 1947 \citep{sznagy1947perturbations,simon2018kato}.  Kato gave a streamlined
algebraic construction, described in his 1955 report and later in
\emph{Perturbation Theory for Linear Operators}
\citep[Chapter~I, \S4.6, equations~(4.33)--(4.42)]{kato1966perturbation};
see also \citet[Theorem~5.1]{simon2018kato}.

For orthogonal projections the useful intertwiner is
\begin{equation}
 \sym{IntertwinerT}=\sym{Qtarget}\sym{Psource}+(\sym{Identity}-\sym{Qtarget})(\sym{Identity}-\sym{Psource}), \qquad \sym{IntertwinerT}\sym{Psource}=\sym{Qtarget}\sym{IntertwinerT}.
 \label{eq:history-kato-T}
\end{equation}
The identity
\begin{equation}
 \like{IntertwinerT}{\symglyph{IntertwinerT}^*}\sym{IntertwinerT}=\sym{Identity}-(\sym{Psource}-\sym{Qtarget})^2
 \label{eq:history-kato-square}
\end{equation}
shows why $\|\sym{Psource}-\sym{Qtarget}\|<1$ is the natural threshold: it makes $\like{IntertwinerT}{\symglyph{IntertwinerT}^*}\sym{IntertwinerT}$
strictly positive and invertible.  Hence
\begin{equation}
 \sym{Rrotation}=\sym{IntertwinerT}\bigl(\sym{Identity}-(\sym{Psource}-\sym{Qtarget})^2\bigr)^{-1/2},
 \qquad \sym{Rrotation}\sym{Psource}\like{Rrotation}{\symglyph{Rrotation}^*}=\sym{Qtarget}
 \label{eq:history-kato-unitary}
\end{equation}
is unitary. This is the polar decomposition $\sym{IntertwinerT}=\sym{Rrotation}|\sym{IntertwinerT}|$, where
$|\sym{IntertwinerT}|=(\like{IntertwinerT}{\symglyph{IntertwinerT}^*}\sym{IntertwinerT})^{1/2}$ is positive. The factor $|\sym{IntertwinerT}|$ describes stretching;
the unitary factor $\sym{Rrotation}$ preserves lengths. Here $\sym{IntertwinerT}$ is invertible, so
$\sym{Rrotation}=\sym{IntertwinerT}|\sym{IntertwinerT}|^{-1}$. In principal-plane coordinates it is the direct rotation through the
principal angle, acting trivially on common directions.

The transport equation
$\sym{Wcompetitor}\sym{Psource}\sym{Wcompetitor}^*=\sym{Qtarget}$
has many unitary solutions.
The polar-factor formula selects one determined by the two projections
and continuously connected to the identity under the norm bound above. That
choice makes the construction valuable in perturbation theory: the
transport identifies moving spectral subspaces while avoiding an
arbitrary choice of orthonormal bases. In infinite dimensions,
$\|\sym{Psource}-\sym{Qtarget}\|<1$ is stronger than acuteness: it also
bounds the operator angles away from $\sym{PiConstant}/2$. For equal-dimensional
finite subspaces the conditions coincide. For an orthogonal or unitary transporter $\sym{Wcompetitor}$, the operator
$\sym{Identity}-\sym{Wcompetitor}$ records the displacement of every vector. Davis and Kahan ask which
norms of this displacement are minimized by the direct rotation
\citep[Section~4]{davis1970rotation}.

Kato used such transformations to identify nearby spectral
subspaces of analytic operator families, thereby reducing the study
of an isolated moving eigenvalue cluster to a fixed coordinate
space.  Simon's historical account explains the earlier Sz.-Nagy
construction, Kato's formula, and subsequent work on pairs of
projections \citep[\S5]{simon2018kato}. Our proofs of direct-rotation
existence, uniqueness, and symmetry use the square-root and
polar-decomposition identities in
\cref{eq:history-kato-square,eq:history-kato-unitary}. The condition
$\|\sym{Psource}-\sym{Qtarget}\|<1$ applies in arbitrary Hilbert dimension.

\subsection{Spectral projections and residuals}
\label{sec:related-spectral}

The spectral theorem represents a self-adjoint operator as
$\sym{A}=\int\sym{SpectralValue}\,d\sym{SpectralMeasure}(\sym{SpectralValue})$ for a projection-valued measure $\sym{SpectralMeasure}$.
For a Borel set $\sym{SpectralSet}$, the operator $\sym{SpectralMeasure}(\sym{SpectralSet})$ is the orthogonal projection onto
the corresponding spectral subspace. This gives a second way of obtaining
a subspace: select a region of the spectrum and take the range of its
spectral projection. The functional-calculus
framework developed from the work of Riesz, von Neumann, and others;
Kato made spectral projections central to the perturbation theory of
closed operators \citep{kato1966perturbation,simon2018kato}.
Unlike the finite-dimensional eigenspace picture, a spectral
subspace need not be spanned by a finite list of eigenvectors.
A self-adjoint operator may be unbounded, while each of its spectral
projections is bounded and everywhere defined. The Hellinger--Toeplitz
theorem states that an everywhere-defined symmetric operator on a Hilbert
space is bounded \citep{hellinger1910grundlagen}. Thus an unbounded
self-adjoint operator has a proper domain. Its action $\sym{A}x$ is defined only
for $x\in\DKDom(\sym{A})$, so residual and intertwining equations must include
the required domain inclusions.

In the complex case, a number $z$ belongs to the resolvent set when
$z\sym{Identity}-\sym{A}$ has a bounded, everywhere-defined inverse. That inverse,
$(z\sym{Identity}-\sym{A})^{-1}$, is the resolvent at $z$. For an isolated bounded spectral
cluster, choose a positively oriented contour $\sym{RieszContour}$ in the resolvent
set enclosing the cluster and no other spectrum. Its projection is
\begin{equation}
 \sym{RieszProjection}=\frac{1}{2\sym{PiConstant} i}\int_{\sym{RieszContour}}(z\sym{Identity}-\sym{A})^{-1}\,dz.
 \label{eq:history-riesz-projection}
\end{equation}
The resolvent is bounded at each point of $\sym{RieszContour}$ even when $\sym{A}$ is unbounded.
This makes the spectral projection a stable object for perturbation analysis,
  while preserving the separate domain requirements whenever $\sym{A}$ itself appears
  in an identity.
More general Borel spectral projections cover finite and half-infinite intervals,
  including spectral configurations for which a single isolated
  finite-eigenvalue cluster is not the right model.
Our formalization uses Borel functional calculus to construct projections for
  these spectral sets, including for unbounded self-adjoint operators.

The Courant--Fischer min--max principle and Weyl-type eigenvalue
perturbation bounds provide complementary information about spectral
\emph{values} \citep{bhatia1997matrix,kato1966perturbation}.  They do
not by themselves identify a stable basis inside a multiple-eigenvalue
cluster, or quantify the movement of its invariant subspace.  The Courant--Fischer and Weyl results that we formalized describe extremal
Rayleigh quotients and eigenvalue stability. The angle estimates add control
of the associated subspaces.

Davis studied the rotation of eigenvectors under perturbation in
1963 and 1965 \citep{davis1963rotation,davis1965rotation}.  The
1969 Davis--Kahan announcement preceded their extended 1970 treatment
\citep{daviskahan1969bounds,davis1970rotation}.  Their problem has
two closely related formulations.  In one, the original self-adjoint
operator is perturbed and the movement of an exact spectral subspace
is estimated.  In the other, one has an approximate trial subspace and
a computable residual: how far can it be from the desired exact
spectral subspace?  Unlike a comparison of individual eigenvectors,
these are meaningful even when an eigenvalue has multiplicity and its
orthonormal eigenbasis is nonunique.

Write $\sym{A}$ for a reference operator and $\sym{Acurrent}=\sym{A}+\sym{H}$ for the perturbed
ambient operator. We use this convention in the sine-theta statement below.
Let $\sym{E0}$ embed a trial-coordinate Hilbert space into the ambient space
isometrically, and let $\sym{TrialOperator}$ be a self-adjoint operator on those coordinates.
The residual is the bounded operator agreeing, where defined, with
\begin{equation}
 \sym{R}=\sym{Acurrent}\sym{E0}-\sym{E0}\sym{TrialOperator}.
 \label{eq:history-residual}
\end{equation}
If $\sym{R}=0$, applying $\sym{Acurrent}$ after the embedding agrees with applying
$\sym{TrialOperator}$ before it: the trial subspace is invariant and $\sym{TrialOperator}$ represents its
operator. For $\sym{R}\ne0$, the residual measures the failure of this identity.  A second isometric coordinate map for the
desired exact spectral space defines its orthogonal projection $\sym{Pexact}$.
The positive operator
\begin{equation}
 \sym{sinTheta0}=\bigl(\sym{Identity}-\like{E0}{\symglyph{E0}^*}\sym{Pexact}\sym{E0}\bigr)^{1/2}
 \label{eq:history-angle-sine}
\end{equation}
measures the component of a trial vector lying outside the exact
subspace.  The critical distinction for the formal proof is that this
positive sine acts on the trial-coordinate space, whereas
$(\sym{Identity}-\sym{Pexact})\sym{E0}$ maps into the ambient space.  For an operator $B$, its positive modulus $|B|=(B^*B)^{1/2}$ acts on its
input space and records its stretching factors. The modulus of $(\sym{Identity}-\sym{Pexact})\sym{E0}$
is $\sym{sinTheta0}$, so the two maps have equal unitarily invariant ideal
norms whenever those norms are finite.

\subsection{Sylvester estimates}
\label{sec:related-sylvester}

The matrix equation $\sym{SylvesterLeft}\sym{X}-\sym{X}\sym{HistoricalSylvesterRight}=\sym{SylvesterRHS}$ is associated with Sylvester's 1884
work \citep{sylvester1884equation}.  Rosenblum studied separated
operator spectra in 1956 and the equation with self-adjoint coefficients
in 1969 \citep{rosenblum1956equation,rosenblum1969selfadjoint}.

The scalar prototype is $(\sym{SylvesterScalarLeft}-\sym{SylvesterScalarRight})\sym{SylvesterScalarUnknown}=\sym{SylvesterScalarRHS}$: if $|\sym{SylvesterScalarLeft}-\sym{SylvesterScalarRight}|\geq\sym{delta}>0$,
then $\sym{delta}|\sym{SylvesterScalarUnknown}|\leq|\sym{SylvesterScalarRHS}|$.  In the operator problem the scalar
multipliers are replaced by the spectral actions of two operators.
The operator estimate has the form $\sym{delta} \sym{NormN}(\sym{X})\leq \sym{NormN}(\sym{SylvesterLeft}\sym{X}-\sym{X}\sym{HistoricalSylvesterRight})$.
Its hypotheses specify the spectral positions of $\sym{HistoricalSylvesterRight}$ and $\sym{SylvesterLeft}$, their domains
when unbounded, and the unitarily invariant norm $\sym{NormN}$.

Let $\sym{F1}$ give isometric coordinates for the exact complementary
spectral subspace, with self-adjoint block $\sym{Lambda1}$.  Writing
$\sym{X}=\like{F1}{\symglyph{F1}^*}\sym{E0}$ converts the residual identity into a Sylvester equation:
\begin{equation}
 \sym{Lambda1}\sym{X}-\sym{X}\sym{TrialOperator}=\like{F1}{\symglyph{F1}^*}\sym{R}
 \label{eq:history-sylvester}
\end{equation}
on the appropriate domain. Thus $\sym{X}$ represents the trial frame's
``leakage'' (i.e., its orthogonal projection onto the exact complementary
spectral subspace, expressed in $\sym{F1}$ coordinates), while the right
side is controlled by the computable residual.  Under the interval/exterior
or ordered spectral-gap hypotheses of the sine-theta theorem,
separated Sylvester estimates give
\begin{equation}
 \sym{delta}\,\sym{NormN}(\sym{X})\leq \sym{NormN}(\like{F1}{\symglyph{F1}^*}\sym{R})\leq \sym{NormN}(\sym{R}).
 \label{eq:history-separation-bound}
\end{equation}
This is the bridge from a spectral-gap assumption to the bound on the
sine of a subspace angle.

The coefficient of $\sym{NormN}(\sym{R})/\sym{delta}$ in this bound is one.
This coefficient uses the ordered or interval/exterior spectral
positions specified by Davis and Kahan. Arbitrarily interlaced spectra can
require a different constant. Our theorem includes both of the source's
spectral configurations, including ordered cases where both operators are
unbounded.  A formal implementation
therefore needs the domains of unbounded self-adjoint operators,
their spectral projections, strong spectral cutoffs, and the
ideal-norm inequalities used to pass through bounded approximations.

\subsection{Operator ideals and displacement norms}
\label{sec:related-norms}

The geometry of a pair of subspaces is fixed, but there are many ways
to assign a cost to a unitary transporting one onto the other.
Von Neumann's 1937 study of unitarily invariant matrix norms identified
the role of symmetric gauge functions of singular values
\citep{vonneumann1937matrix}. Schatten's 1950 monograph developed normed
operator ideals; Ky Fan's 1951 eigenvalue and singular-value inequalities
provided the partial-sum comparisons underlying majorization
arguments \citep{schatten1950crossspaces,fan1951maximum,bhatia1997matrix}.
The singular values $\sym{SingularValue}(Z)$ are the nonnegative square roots of the
eigenvalues of $Z^*Z$, ordered decreasingly. They record the stretching
factors of the matrix $Z$. For $1\leq p<\infty$, the Schatten norms include
\begin{equation}
 \|Z\|_{S_p}=\left(\sum_j \sym{SingularValue}(Z)^p\right)^{1/p},
 \qquad
 \|Z\|_{S_1}=\sum_j \sym{SingularValue}(Z),
 \quad
 \|Z\|_{S_2}=\left(\sum_j \sym{SingularValue}(Z)^2\right)^{1/2}.
 \label{eq:history-schatten}
\end{equation}
An operator ideal is a linear class of operators closed under multiplication
on either side by bounded operators. A normed ideal also controls the norm
of these products. For compact operators in infinite dimensions, the
Schatten $p$ ideal consists of those whose singular values have a finite
$p$-power sum. Thus a trace-norm inequality requires trace-class membership;
this is an instance of the source's qualification that its norms be defined.

The largest singular value gives the operator norm, the Euclidean
norm of the singular-value sequence gives the Hilbert--Schmidt norm
(the Frobenius norm $\dkFrobeniusNorm{Z}=\|Z\|_{S_2}$ for matrices),
and its sum gives the trace norm. These norms can rank transporters
differently. In Section~\ref{sec:prop44}, the two displacement singular-value
lists are $(2\sin(\sym{theta}/2),\ldots,2\sin(\sym{theta}/2))$ with four entries and
$(2\sin\sym{theta},2\sin\sym{theta},0,0)$. The second has a smaller sum and a larger
maximum. Ky Fan norms compare partial sums of the ordered singular values;
these comparisons determine which norm inequalities follow.

Davis established Hilbert--Schmidt and finite-dimensional Schatten
extremal properties for exponents at least two in his 1958 subspace
paper \citep[Theorems~7.3--7.4]{davis1958separation}.  Davis and
Kahan compare three objectives \citep[Section~4]{davis1970rotation} for a transporter $\sym{Wcompetitor}$ carrying
$\sym{Usource}$ onto $\sym{Vtarget}$: its restricted displacement $(\sym{Identity}-\sym{Wcompetitor})\like{Psource}{\symglyph{Psource}_U}$, its squared full
displacement $(\sym{Identity}-\sym{Wcompetitor})^*(\sym{Identity}-\sym{Wcompetitor})$, and its full displacement $\sym{Identity}-\sym{Wcompetitor}$. The last acts
on every ambient vector, including vectors outside $\sym{Usource}$. They gave a
complex-phase modification that decreases its trace norm \citep[Example~4.2]{davis1970rotation}.
The full-displacement claim \citep[Proposition~4.4]{davis1970rotation}
asserts that direct rotation minimizes
full displacement for real subspaces if their principal angles are no
larger than $\sym{PiConstant}/3$.  The four-dimensional example in
\cref{sec:prop44} gives a smaller trace norm at every positive acute
angle.  It exposes the failed step: adding singular-value sums of
restrictions to separate principal planes is not a valid global lower
bound for an arbitrary ambient transporter.

The subsequent minimal-distortion theory of Bolotnikov, Li, and Rodman
studied the same class of subspace-transport constraints, including
Frobenius and Q-norm optimality and explicit subquadratic
counterexamples \citep{bolotnikov2001minimal}.  Later work on
Grassmannian metrics and canonical-angle inequalities also separates
intrinsic paths between subspaces from minimizing the extrinsic
operator displacement \citep{qiu2006grassmann,zhang2010canonicalangles}.
These results give the norm-theoretic context for the source
correction and for the trace-norm matching problem in
\cref{sec:example-transport}.

\subsection{The Davis--Kahan angle estimates}
\label{sec:related-synthesis}

Davis's papers of 1963 and 1965 and the Davis--Kahan announcement of 1969
preceded the extended 1970 treatment
\citep{davis1963rotation,davis1965rotation,daviskahan1969bounds,davis1970rotation}.
For the sine-theta estimate, let $\sym{R}=\sym{Acurrent} \sym{E0}-\sym{E0}\sym{TrialOperator}$ be the residual
and $\sym{Pexact}$ the desired exact spectral projection. The operator $(\sym{Identity}-\sym{Pexact})\sym{E0}$
extracts the part of a trial vector outside that spectral subspace. In
isometric coordinates $\sym{F1}$ for the complementary space, this component is
represented by $\sym{X}=\like{F1}{\symglyph{F1}^*}\sym{E0}$. The residual identity gives
\cref{eq:history-sylvester}, and a separated Sylvester estimate bounds $\sym{NormN}(\sym{X})$.
Since $\sym{S}=\sym{F1}\sym{X}$ and $\sym{F1}$ is an isometry,
$\sym{NormN}(\sym{X})=\sym{NormN}(\sym{S})$.
The modulus identity $|\sym{S}|=\sym{sinTheta0}$ gives
$\sym{NormN}(\sym{S})=\sym{NormN}(\sym{sinTheta0})$.

The resulting chain of equalities and inequalities is
\begin{equation}
 \underbrace{\sym{delta} \sym{NormN}(\sym{sinTheta0})}_{\text{subspace error}}
 =\sym{delta} \sym{NormN}(\sym{X})
 \leq \sym{NormN}(\like{F1}{\symglyph{F1}^*}\sym{R})
 \leq\underbrace{\sym{NormN}(\sym{R})}_{\text{computable residual}}.
 \label{eq:history-unified-chain}
\end{equation}
The first equality follows from the positive-modulus identity for the
complementary projection. The first inequality applies the Sylvester
estimate to $\sym{X}$; the second uses $\sym{NormN}(BC)\leq\dkOperatorNorm{B}\sym{NormN}(C)$ and
$\dkOperatorNorm{\like{F1}{\symglyph{F1}^*}}\leq1$. For unbounded operators the Sylvester
identity holds on the trial operator's domain, with a bounded residual.
These are the geometric, spectral, and norm estimates formalized in our
proof of the theorem.

The tangent and double-angle estimates use related projection blocks.
For example, an estimate for $\sin(2\sym{Theta})$ allows angles near zero and
near $\sym{PiConstant}/2$. The spectral-selection results \citep[Section~8]{davis1970rotation}
identify which spectral subspaces give the small-angle alternative.
The direct-rotation analysis \citep[Section~4]{davis1970rotation}
studies the orthogonal or unitary map carrying one subspace onto another
and asks which displacement norms the direct rotation minimizes.

\subsection{Numerical and statistical applications}
\label{sec:related-numerical}

The 1970 theory influenced more than the analysis of exact
self-adjoint eigenvectors.  Wedin's 1972 singular-subspace
perturbation bounds address the corresponding problem for singular
vectors of general rectangular matrices
\citep{wedin1972perturbation}.  Bj\"orck and Golub's angle algorithms
made the subspace quantities computable
\citep{bjorck1973numerical}.  Stewart's work on pseudoinverses,
projections, and least-squares sensitivity helped connect these
geometric notions with matrix computations
\citep{stewart1977perturbation,stewart1990matrix}.
Paige and Wei's 1994 history shows how the cosine--sine decomposition
links the principal-angle geometry of Davis and Kahan with numerical
matrix factorizations \citep{paige1994history}.

The same subspace-versus-eigenvector distinction remains central in
statistical spectral estimation.  For example, Yu, Wang, and
Samworth give a statistical variant whose eigenspace error bound uses
separation in the population spectrum \citep{yu2015davis}, while Chen and
collaborators survey spectral methods across statistical learning
problems \citep{chen2021spectral}.  In that setting, an empirical covariance matrix approximates a population
operator and the perturbation norm controls the error of estimated
eigenspaces. Our general theorem supplies the deterministic angle estimate;
using it in a statistical result also requires the corresponding probabilistic
control of the perturbation.

\subsection{Foundations used in our formalization}
\label{sec:related-formal}

Each historical component above corresponds to a specific requirement
in the formalization.  Jordan's angle geometry and Halmos's
projection blocks lead to positive principal-angle operators rather
than chosen eigenbases.  The Sz.-Nagy--Kato intertwiner requires
polar decomposition, square roots, and real/complex transfer.
Spectral projections require Borel functional calculus, including for
unbounded self-adjoint operators.  The Sylvester proof requires a
representation carrying operator domains and a controlled passage
through spectral cutoffs.  Schatten and Ky Fan theory requires
approximation numbers and symmetric operator ideals to express
inequalities with the source's norm generality.

We formalized principal-angle geometry, symmetric operator ideals, and
unbounded Sylvester estimates using Lean~4 and Mathlib
\citep{demoura2021lean4,mathlib2020}. Mathlib supplies the algebraic,
topological, and Hilbert-space framework. Its continuous functional calculus
is described in the preprint by Anatole Dedecker and Jireh Loreaux
(\href{https://arxiv.org/abs/2501.15639}{arXiv:2501.15639}).
Our projection-valued-measure and Borel-calculus
work also adapts Adam Bornemann's Spectra formalization
\citep{bornemann2026spectra}.

Our Courant--Fischer and Weyl development, rectangular unitarily invariant
seminorms, projection-valued measures, and supporting polar-decomposition
results were subsequently adapted into Tau Ceti \citep{tauceti2026}.
% Editorial note (not manuscript prose): Do not add a sentence emphasizing
% the current Tau Ceti canonical polar-decomposition interface here.
% This paragraph records work developed in DKPS and subsequently
% contributed upstream; current library interfaces are discussed elsewhere.
% Evidence: dev/tauceti/dkps-upstream-contribution-census.md, PRs 13599,
% 14241, 17527 (adaptations), and 16127/16450 (independent polar proofs).
% AUTHORING_NOTES.md indexes accepted-source and campaign evidence.
% Mathematical origin and formal proof provenance are separate attributions.

This history also separates the two mathematical questions that meet
in the project.  The sine-theta theorem estimates \emph{how far} an
approximate or perturbed spectral subspace lies from an exact one.
The direct-rotation problem asks \emph{which} ambient transformation
carries one given subspace onto another at minimum displacement cost.
Our proofs of the optimization results use the principal-angle,
projection, and operator-ideal foundations developed for the formalization.
The matching theorem has its own proof; the shared foundations explain the
relationship between the two parts of the work. The resulting counterexample
to the full-displacement claim \citep[Proposition~4.4]{davis1970rotation}
and the follow-on matching theorem are examined in
\cref{sec:prop44,sec:example-transport}.

\section{Davis--Kahan formalization and the sine-theta theorem}
\label{sec:coverage}

Our targets are the four headline theorems \citep[Section~2]{davis1970rotation} and every
named theorem, proposition, lemma, and corollary established by Davis and
Kahan in the 1970 paper. These \DKResultCount{} results include the sine and
tangent inequalities, the geometry and extremal properties of direct
rotations, bounded and unbounded Sylvester estimates, the generalized
perturbation bounds \citep[Section~6]{davis1970rotation}, and the spectral-selection
theorems \citep[Section~8]{davis1970rotation}. We prove \DKProvedCount{} results and refute the printed
full-displacement claim \citep[Proposition~4.4]{davis1970rotation}.
Appendix~\ref{app:dk-inventory} lists these targets.

We develop the sine-theta theorem and its proof below; the remaining source
results are summarized in Appendix~\ref{app:dk-inventory}, with hypotheses
and conclusions in Section~\ref{app:dk-result-summaries}. The main text also
explains the spectral-selection results needed to interpret the angle bounds.
Appendix~\ref{app:notation} collects the notation used here and in the
direct-rotation argument.

\subsection{The sine-theta theorem for self-adjoint operators}
\label{sec:sine-theta}

\subsubsection{Trial and exact subspaces}

We use $\sym{A}$ for the reference operator and
$\sym{Acurrent}=\sym{A}+\sym{H}$ for the perturbed ambient operator.
The symbol $\sym{Acurrent}$ corresponds to $\sym{A}+\sym{H}$ in the source.

For a unit trial vector $\sym{v}$ and an approximate eigenvalue $\sym{mu}$, the residual
is $\sym{Acurrent}\sym{v}-\sym{mu} \sym{v}$. An isometric coordinate map $\sym{E0}$ generalizes
this construction to a trial subspace: its range consists of the candidate
approximate vectors, and a self-adjoint operator $\sym{TrialOperator}$ represents their
approximate dynamics. For a finite-dimensional Rayleigh--Ritz choice,
$\sym{TrialOperator}=\like{E0}{\symglyph{E0}^*}\sym{Acurrent}\sym{E0}$. The theorem allows other self-adjoint trial
operators as well. We write $\sym{TrialOperator}$ for the general trial operator,
reserving $\sym{ReferenceBlock0}$ for the reference spectral block that
Davis and Kahan also call $A_0$.

Let $\sym{AmbientHilbertSpace}$ be the ambient Hilbert space.  Let $\sym{TrialOperator}$ be a self-adjoint operator on a
trial-coordinate space $\sym{TrialCoordinateSpace}$, and let
\[
  \sym{E0}:\sym{TrialCoordinateSpace}\to \sym{AmbientHilbertSpace}
\]
be an isometric coordinate map for the trial subspace.  The residual compares
applying the ambient operator after embedding with applying the trial operator
before embedding:
\begin{equation}
  \sym{R} = \sym{Acurrent}\sym{E0}-\sym{E0}\sym{TrialOperator}.
  \label{eq:residual}
\end{equation}
When the operators are unbounded, \cref{eq:residual} is asserted on
$\DKDom(\sym{TrialOperator})$ and $\sym{R}$ is the bounded residual that agrees there.
The Lean predicate \texttt{IsTrialResidual} packages the isometry, the required
domain inclusion, and this residual identity.

Let $\sym{F0}:\sym{ExactCoordinateSpace}\to \sym{AmbientHilbertSpace}$ be isometric coordinates for the desired exact spectral
subspace and let $\sym{F1}:\sym{ComplementaryCoordinateSpace}\to \sym{AmbientHilbertSpace}$ be isometric coordinates for its orthogonal
complement.  A self-adjoint operator $\sym{Lambda1}$ on $\sym{ComplementaryCoordinateSpace}$ describes the
complementary spectral block through
\[
  \sym{Acurrent}\sym{F1} = \sym{F1}\sym{Lambda1}
\]
on its domain. This intertwining identity says that $\sym{Lambda1}$ gives
the action of $\sym{Acurrent}$ in complementary coordinates. The predicate
\texttt{IsExactSpectralDecomposition} also requires the domain inclusion
$\sym{F1}(\DKDom(\sym{Lambda1}))\subseteq\DKDom(\sym{Acurrent})$
and states orthogonality and completeness through
\[
 \like{F0}{\symglyph{F0}^*}\sym{F1}=0,\qquad
 \sym{F0}\like{F0}{\symglyph{F0}^*}+\sym{F1}\like{F1}{\symglyph{F1}^*}=\sym{Identity}.
\]

\subsubsection{The sine operator and its rectangular representative}

Here $\sym{Theta0}$ acts on $\sym{TrialCoordinateSpace}$; under the isometry
$\sym{E0}$ it represents the directed angle operator on the trial subspace
introduced in \cref{sec:related-two-projections}. If
$\sym{Pexact}=\sym{F0}\like{F0}{\symglyph{F0}^*}$ is the orthogonal projection onto the desired exact subspace, then
\[
  \cos\sym{Theta0} = (\like{E0}{\symglyph{E0}^*}\sym{Pexact}\sym{E0})^{1/2},
  \qquad
  \sym{sinTheta0} = (\sym{Identity}-\like{E0}{\symglyph{E0}^*}\sym{Pexact}\sym{E0})^{1/2}.
\]
The analytic proof uses the rectangular complementary block
\begin{equation}
  \sym{S}=(\sym{Identity}-\sym{Pexact})\sym{E0}.
  \label{eq:rectangular-sine}
\end{equation}
The theorem uses the positive sine on the trial-coordinate space.  The rectangular block has codomain $\sym{AmbientHilbertSpace}$, while the
positive sine has codomain $\sym{TrialCoordinateSpace}$.  Their relationship is
\[
  \sym{S}^*\sym{S}
  = \like{E0}{\symglyph{E0}^*}(\sym{Identity}-\sym{Pexact})\sym{E0}
  = \sym{Identity}-\like{E0}{\symglyph{E0}^*}\sym{Pexact}\sym{E0},
\]
so
\begin{equation}
  |\sym{S}|=(\sym{S}^*\sym{S})^{1/2}=\sym{sinTheta0}.
  \label{eq:modulus-sine}
\end{equation}
% It might be worth noting that our formalization actually defines \sym{Theta0}
% and proves eq:modulus-sine whereas most treatments of the subject seam to
% just start from |\sym{S}|. Maybe it's worth enough to put in the paper, but
% it at least deserves a latex comment!
The polar decomposition $\sym{S}=\sym{PolarU}|\sym{S}|$ separates the positive
  modulus from a partial isometry.
Unitarily invariant ideal norms assign the same value to an operator and its
  modulus, so \[ \sym{NormN}(\sym{S})=\sym{NormN}(\sym{sinTheta0}) \] whenever
  the norm is finite.
This equality transfers an estimate for the complementary block $\sym{S}$ to the
  positive sine operator.

\subsubsection{Gap geometry and the inequality}
\label{sec:gap-geometry}

Fix $\sym{delta}>0$. There are two geometric forms of separation.
\emph{Interval/exterior separation:} for some real
$\sym{betaEndpoint}\leq\sym{alphaEndpoint}$, one of the two spectra
$\sym{SpectrumSymbol}(\sym{TrialOperator})$ and $\sym{SpectrumSymbol}(\sym{Lambda1})$ lies in $[\sym{betaEndpoint},\sym{alphaEndpoint}]$ and the other
lies in
\[
 (-\infty,\sym{betaEndpoint}-\sym{delta}]\;\cup\;[\sym{alphaEndpoint}+\sym{delta},\infty).
\]
\emph{Ordered separation:} for some real $b$, the operator forms satisfy either
\[
 \sym{TrialOperator}\geq (b+\sym{delta})\sym{Identity},\quad \sym{Lambda1}\leq bI,
 \qquad\text{or}\qquad
 \sym{TrialOperator}\leq bI,\quad \sym{Lambda1}\geq (b+\sym{delta})\sym{Identity}.
\]
Each form inequality is asserted on the domain of the corresponding operator;
for example, $\sym{TrialOperator}\geq aI$ means
$\langle \sym{TrialOperator}x,x\rangle\geq a\|x\|^2$ for every
$x\in\DKDom(\sym{TrialOperator})$. For self-adjoint operators, the source's ordered
hypotheses on half-infinite spectral intervals, such as $(-\infty,b]$ and
$[b+\sym{delta},\infty)$, imply these form bounds. These cases are included
in the source's unbounded sine theorem
\citep[Section~6, discussion of unbounded operators]{davis1970rotation}.
Write $\sym{SineGap}(\sym{TrialOperator},\sym{Lambda1})$ for their union.
The Lean predicate \texttt{FormBoundedSylvesterGap} has three cases:
the finite-interval alternative, which includes either spectral order,
and the two ordered form alternatives.

Both blocks may be unbounded in an ordered configuration. For instance,
on $\ell^2(\mathbb N_{\geq1})$, take diagonal operators
$\sym{TrialOperator}e_n=(n+\sym{delta})e_n$ and $\sym{Lambda1}e_n=-ne_n$, on the domains where their
images are square summable. They satisfy the first ordered configuration
with $b=0$, while neither spectrum lies in a finite interval. This example
explains why retaining the half-infinite alternatives changes the theorem's
applicability.

For the unbounded trial configuration, the source assumes that
$\sym{Acurrent}\sym{E0}$ and $\sym{E0}\sym{TrialOperator}$ have a common dense domain.
Since $\sym{E0}$ is everywhere defined, this gives
\[
 \DKDom(\sym{Acurrent}\sym{E0})=\DKDom(\sym{TrialOperator}),
 \qquad
 \sym{E0}(\DKDom(\sym{TrialOperator}))\subseteq\DKDom(\sym{Acurrent}).
\]
The residual agrees with the difference on this domain and extends to a
bounded operator. The formal predicate \texttt{IsTrialResidual} includes the
domain inclusion and the residual identity.

The norm $\sym{NormN}$ is a normalized symmetric operator-ideal norm.
Normalization means that the rank-one map $uv^*:x\mapsto u\langle v,x\rangle$
has norm
\[
 \sym{NormN}(uv^*)=\|u\|\,\|v\|;
\]
in particular, its norm is one when $u,v$ are unit vectors.
Its hypotheses also include Ky Fan dominance whenever both ideal norms
are finite: domination of all finite singular-value sums implies domination
of the ideal norms. This comparison does not assert ideal membership.
The norm is invariant under isometries and positive modulus, and satisfies
$\sym{NormN}(BT)\leq\dkOperatorNorm{B}\sym{NormN}(T)$. Trace-class and
Hilbert--Schmidt operators with their respective norms are examples; the
operator norm on bounded operators is another. For an ideal norm, both the
residual and the positive sine must belong to its domain.

In two dimensions, take a one-dimensional trial frame $\sym{E0}(x)=x\sym{v}$ with
$\|\sym{v}\|=1$. A scalar trial operator $\sym{TrialOperator}=\sym{rho}$ gives
$\sym{r}=\sym{R}(1)=\sym{Acurrent}\sym{v}-\sym{rho}\sym{v}$.
Choosing $\sym{rho}=\langle\sym{v},\sym{Acurrent}\sym{v}\rangle$
makes $\sym{rho}\sym{v}$ the orthogonal projection of
$\sym{Acurrent}\sym{v}$ onto the trial line. The residual is therefore
perpendicular to that line. Figure~\ref{fig:sine-theta-residual} pairs this
residual triangle with the exact complementary projection that measures
the subspace error.

\begin{figure}[tb]
\centering
\resizebox{119mm}{!}{\input{figures/Fig1.pgf}}
\caption{Rayleigh residual geometry for a one-dimensional trial space in
$\sym{RealField}^2$ (the rank-one case). With $\sym{E0}(x)=x\sym{v}$, the residual
operator satisfies $\sym{R}(x)=x\sym{r}$, where
$\sym{r}=\sym{Acurrent}\sym{v}-\sym{rho}\sym{v}$.
For a higher-dimensional trial space,
$\sym{R}=\sym{Acurrent}\sym{E0}-\sym{E0}\sym{TrialOperator}$ assigns a residual vector
to each trial-coordinate vector. In a finite-dimensional eigenbasis of
$\sym{TrialOperator}$, its columns are the individual eigenpair residuals.
The exact frames in the drawing are
$\sym{F0}(h)=(h,0)$ and $\sym{F1}(g)=(0,g)$, with
$\sym{Pexact}=\sym{F0}\like{F0}{\symglyph{F0}^*}$ and $\sym{PexactPerp}=\sym{Identity}-\sym{Pexact}$.
(a) The pink leg $\sym{S}(1)=\sym{PexactPerp}\sym{v}$ has length $\sin\sym{theta}$.
(b) The green residual joins $\sym{rho}\sym{v}$ to $\sym{Acurrent}\sym{v}$;
for the Rayleigh choice it is perpendicular to the trial line. The dashed
legs are its exact wanted and unwanted components. Their lengths give
$\sym{delta}\sin\sym{theta}=\|\sym{PexactPerp}\sym{r}\|_2
\leq\dkOperatorNorm{\sym{r}\|_2=\|\sym{R}}$.
The drawing uses $\sym{Acurrent}=\sym{DiagonalOperator}(1.9,0.7)$,
$\sym{theta}=40^\circ$, $\sym{rho}=\langle\sym{v},\sym{Acurrent}\sym{v}\rangle$,
and $\sym{delta}=\sym{rho}-0.7$}
\label{fig:sine-theta-residual}
\end{figure}
% Place the residual figure before the theorem and its Lean signature.
\FloatBarrier

Appendix~\ref{app:lean-guide} explains the general Lean syntax in the
following theorem signature; the geometric and spectral hypotheses have been
defined above.

\begin{theorem}[Sine-theta estimate]
\label{thm:sine-theta}
Let the scalar field be $\sym{RealField}$ or $\sym{ComplexField}$, let $\sym{AmbientHilbertSpace}$ be separable,
and let $\sym{Acurrent}$, $\sym{TrialOperator}$, and $\sym{Lambda1}$ be self-adjoint operators on the
Hilbert spaces above. Suppose the isometries, domain inclusions, residual
identity, exact complementary decomposition, and the separation condition
$\sym{SineGap}(\sym{TrialOperator},\sym{Lambda1})$ of
Section~\ref{sec:gap-geometry} hold. Let $\sym{NormN}$ be a normalized symmetric
operator-ideal norm. Whenever $\sym{sinTheta0}$ and $\sym{R}$ belong to its domain,
\begin{equation}
 \sym{delta}\,\sym{NormN}(\sym{sinTheta0})\leq \sym{NormN}(\sym{R}).
 \label{eq:dk-sine}
\end{equation}
\end{theorem}

% Keep the Lean signature and its caption together across page breaks.
\noindent\begin{minipage}{\linewidth}
\begin{lstlisting}[style=leanpaper,
  caption={Lean signature for the sine-theta estimate. The surrounding
  scalar-field and Hilbert-space context is omitted.
  This is a readable ASCII transcription of the Lean declaration,
  rather than a verbatim copy of its Unicode binders and operator arrows.
  The paper-to-Lean names are registered centrally and displayed in
  Table~\ref{tab:paper-lean-notation}. In particular, the Lean binder
  displayed here as \texttt{A0} names our $\sym{TrialOperator}$, not the
  distinct reference block $\sym{ReferenceBlock0}$.
  The positive trial-coordinate sine is
  \texttt{sourceDirectedSinThetaOperator E0 F0};
  \texttt{FormBoundedSylvesterGap} encodes the interval/exterior and ordered
  half-infinite separation hypotheses described above.},
  label={lst:dk-sintheta-headline}]
theorem sinTheta_unbounded_formGap_whereDefinedUIN_rclike
    [TopologicalSpace.SeparableSpace E]
    (N : NormalizedSymmetricOperatorIdealFamily K)
    (A : LinearPMap K E E)
    (A0 : LinearPMap K F F)
    (Lambda1 : LinearPMap K G G)
    (E0 : ContinuousLinearMap K F E)
    (F0 : ContinuousLinearMap K H E)
    (F1 : ContinuousLinearMap K G E)
    (R : ContinuousLinearMap K F E)
    (hA : IsSelfAdjoint A)
    (hA0 : IsSelfAdjoint A0)
    (hLambda1 : IsSelfAdjoint Lambda1)
    (htrial : IsTrialResidual A A0 E0 R)
    (hexact : IsExactSpectralDecomposition A Lambda1 F0 F1)
    (delta : Real) (hdelta : 0 < delta)
    (hgap : FormBoundedSylvesterGap A0 Lambda1 delta) :
    N.Mem (sourceDirectedSinThetaOperator E0 F0) ->
    N.Mem R ->
      delta * N.gaugeReal (sourceDirectedSinThetaOperator E0 F0) <=
        N.gaugeReal R
\end{lstlisting}
\end{minipage}\par

The two ideal-membership hypotheses ensure that the norm values in the
conclusion are finite. The gap predicate includes the finite
interval/exterior case and both ordered cases. In the ordered cases,
$\sym{TrialOperator}$ and $\sym{Lambda1}$ may both be unbounded.

\subsubsection{Why the Sylvester estimate gives the angle bound}

Set $\sym{X}=\like{F1}{\symglyph{F1}^*}\sym{E0}$. The exact complementary decomposition gives
$\sym{S}=\sym{F1}\sym{X}$, and multiplication by the isometry $\sym{F1}$ preserves the
singular-value gauges. Self-adjointness, the domain hypotheses, and the
residual identity yield the Sylvester relation
\[
 \sym{Lambda1} \sym{X}-\sym{X}\sym{TrialOperator}=\like{F1}{\symglyph{F1}^*}\sym{R}
\]
on $\DKDom(\sym{TrialOperator})$, with the left side represented by the bounded
right side. The separated Sylvester estimate gives
\[
 \sym{delta} \sym{NormN}(\sym{X})\leq \sym{NormN}(\like{F1}{\symglyph{F1}^*}\sym{R})\leq \sym{NormN}(\sym{R}).
\]
Combining this with $\sym{NormN}(\sym{X})=\sym{NormN}(\sym{S})=\sym{NormN}(\sym{sinTheta0})$ proves the conclusion.
For unbounded operators, deriving the domain compatibility for $\sym{X}$ is part of
the argument. The displayed relation holds on $\DKDom(\sym{TrialOperator})$ and
has a bounded right side. The specified interval/exterior or ordered
spectral positions give the factor $1/\sym{delta}$, with no additional
multiplicative constant.

For a one-dimensional trial space, take
$\sym{Acurrent}=\sym{DiagonalOperator}(0,\sym{delta})$, the trial scalar
$\sym{TrialOperator}=0$, and the unit vector
$\like{v}{\symglyph{v}_{\symglyph{theta}}}=(\cos\sym{theta},\sin\sym{theta})$.
Choose the desired exact subspace $\sym{SpanOperator}\{(1,0)\}$;
its orthogonal complement is $\sym{SpanOperator}\{(0,1)\}$, on which
$\sym{Lambda1}=\sym{delta}$.
The complementary component is $(0,\sin\sym{theta})$ and the
residual is $(0,\sym{delta}\sin\sym{theta})$. Thus equality holds
in the sine-theta estimate for this finite-dimensional configuration.

\section{Formalized operator-theoretic foundations}
\label{sec:foundations}

We formalized the geometry of two subspaces, symmetric operator ideals and
majorization, and Sylvester estimates for unbounded self-adjoint operators.
These results supply the constructions and inequalities used in the
Davis--Kahan proofs.

\subsection{Direct rotations and two-subspace geometry}

Let $\sym{Psource}$ and $\sym{Qtarget}$ project onto closed subspaces
$\sym{Usource}$ and $\sym{Vtarget}$. For an acute pair,
$\|\sym{Psource}-\sym{Qtarget}\|<1$. The compression
$\sym{Psource}\sym{Qtarget}\sym{Psource}$ on $\sym{Usource}$ determines
the positive cosine operator. Its square root and the complementary
compression give the cosine and sine factors in the direct-rotation
construction. Polar decomposition gives the isometries between the
corresponding parts of the two subspaces.

The resulting orthogonal or unitary map $\sym{Rrotation}$ carries
$\sym{Usource}$ onto $\sym{Vtarget}$. In finite-dimensional principal-plane
coordinates it rotates each plane through its principal angle and fixes
common directions. The projection formula from
\cref{eq:history-kato-unitary} defines it in Hilbert space without requiring
a discrete list of principal vectors.

Our formalization proves acute existence and uniqueness, the
principal-square-root characterization, commutation with the angle
operators, and reversal symmetry: interchanging the two subspaces replaces
the rotation by its adjoint. For general pairs, the existence criteria
include the right-angle defect spaces described in
\cref{sec:related-two-projections}. When the angle operator is compact,
its spectral decomposition recovers a sequence of principal angles
\citep[Section~3]{davis1970rotation}.

The double-angle theorems for $\sin(2\sym{Theta})$ and $\tan(2\sym{Theta})$
  require selecting the branch on which $\sym{Theta}\le\sym{PiConstant}/4$.
Here $\sym{Theta}$ includes both the source and complementary angle blocks;
$\sym{Theta0}$ alone is the directed angle on source coordinates.
Because the function $\sin(2\sym{theta})$ is small both when $\sym{theta}$ is
  near $0$ and when $\sym{theta}$ is near $\sym{PiConstant}/2$, a bound on
  $\sin(2\sym{Theta})$ alone cannot determine whether an approximate subspace is
  close to or nearly orthogonal to the target.
The branch-selection theorem \citep[Theorem~8.1]{davis1970rotation}
uses an ordered reference gap. Let $\sym{Psource}$ reduce a self-adjoint
reference operator $\sym{A}$, with
blocks $\sym{ReferenceBlock0}$ on $\DKRange(\sym{Psource})$ and
$\sym{ReferenceBlock1}$ on $\DKKernel(\sym{Psource})$, satisfying
\[
 \sym{ReferenceBlock0}\le\sym{alphaEndpoint}\sym{Identity},
 \qquad
 \sym{ReferenceBlock1}\ge(\sym{alphaEndpoint}+\sym{delta})\sym{Identity},
 \qquad \sym{delta}>0,
\]
in the quadratic-form sense. Suppose $\sym{H}$ is bounded and self-adjoint,
with both diagonal blocks zero:
$\sym{Psource}\sym{H}\sym{Psource}
 =(\sym{Identity}-\sym{Psource})\sym{H}(\sym{Identity}-\sym{Psource})=0$.
These hypotheses imply that $\sym{Acurrent}=\sym{A}+\sym{H}$ has no spectrum
in $(\sym{alphaEndpoint},\sym{alphaEndpoint}+\sym{delta})$.

Fix a reducing subspace $M$ of $\sym{Acurrent}$, and write
$\sym{Lambda0},\sym{Lambda1}$ for the blocks of $\sym{Acurrent}$ on $M,M^\perp$.
The full angle $\sym{Theta}$ comparing $\DKRange(\sym{Psource})$ with $M$
satisfies $\sym{Theta}\le\sym{PiConstant}/4$ if and only if
$\sym{Lambda0}\le\sym{alphaEndpoint}\sym{Identity}$ and
$\sym{Lambda1}\ge(\sym{alphaEndpoint}+\sym{delta})\sym{Identity}$.
Such a subspace exists: take the spectral subspace of $\sym{Acurrent}$
for $(-\infty,\sym{alphaEndpoint}]$. Our formalization also proves that
it is the unique reducing subspace satisfying this angle bound.

To compare the reference and perturbed blocks on common coordinate spaces,
choose isometric frames for $\DKRange(\sym{Psource})$ and
$\DKKernel(\sym{Psource})$. Apply the direct rotation $\sym{Rrotation}$ to
these frames to obtain frames for $M$ and $M^\perp$. Express
$\sym{ReferenceBlock0},\sym{ReferenceBlock1}$ in the original frames and
$\sym{Lambda0},\sym{Lambda1}$ in the transported frames. In these coordinates
the branch satisfies the stronger quadratic-form inequality
\[
 \sym{ReferenceBlock1}-\sym{alphaEndpoint}\sym{Identity}
 \le\sym{ComplementCosineBlock}
       (\sym{Lambda1}-\sym{alphaEndpoint}\sym{Identity})\sym{ComplementCosineBlock},
 \qquad \sym{ComplementCosineBlock}=\cos\sym{Theta1},
\]
with the lower-block analogue. These inequalities quantify spectral
repulsion in terms of the angles and yield eigenvalue and symmetric-gauge
comparisons in finite dimensions.

The smallness criterion \citep[Theorem~8.2]{davis1970rotation} uses a central
spectral interval instead of an ordered gap. Let $\sym{Psource}$ reduce
the self-adjoint operator $\sym{A}$ and $\sym{Qtarget}$ reduce the
self-adjoint operator $\sym{Acurrent}$. The perturbed block on
$\DKRange(\sym{Qtarget})$ has spectrum in
$[\sym{betaEndpoint},\sym{alphaEndpoint}]$, while its complementary spectrum
lies in
\[
 (-\infty,\sym{betaEndpoint}-\sym{delta}]
 \ \cup\ [\sym{alphaEndpoint}+\sym{delta},\infty),
 \qquad \sym{betaEndpoint}\le\sym{alphaEndpoint},\quad\sym{delta}>0.
\]
The reference block $\sym{ReferenceBlock0}$, representing
$\sym{A}|_{\DKRange(\sym{Psource})}$ in source coordinates, must have spectrum
in $[\sym{betaEndpoint}-\sym{delta}/2,\sym{alphaEndpoint}+\sym{delta}/2]$.
For an isometric frame $\sym{E0}$ of the reference subspace, its residual
is $\sym{R}=\sym{Acurrent}\sym{E0}-\sym{E0}\sym{ReferenceBlock0}$
on the operator domain, as in \cref{eq:history-residual}.
Assume
$\sym{DimensionOperator}\sym{CrossedSourceDefect}=\sym{DimensionOperator}\sym{CrossedTargetDefect}$: the space of source
vectors orthogonal to the target and the space of target vectors orthogonal
to the source have equal Hilbert dimension.

If $\dkOperatorNorm{\sym{H}}<\sym{delta}/2$ for a
bounded self-adjoint perturbation, or $\dkOperatorNorm{\sym{R}}<\sym{delta}/2$
for the bounded residual on the reference subspace, then
$\sym{Theta}<\sym{PiConstant}/4$. For each ideal norm whose displayed arguments belong
to its domain, the corresponding double-angle estimate also holds:
\[
 \sym{delta}\sym{NormN}(\sin(2\sym{Theta}))\le2\sym{NormN}(\sym{H}),
 \qquad\text{or}\qquad
 \sym{delta}\sym{NormN}(\sin(2\sym{Theta0}))\le2\sym{NormN}(\sym{R}).
\]
% Note (not for print): the factor 2 in the residual form is correct and
% matches the 1970 Section 2 statement. The 1970 Section 7 proof, after (7.5),
% rewrites delta||sin 2Theta|| <= ||H - XHX|| with right side ||[[0,B*],[B,0]]||
% and concludes delta||sin 2Theta_0|| <= ||B|| <= ||R||; since H - XHX has
% off-diagonal block 2B, the right side should be 2||B|| <= 2||R||. The
% Section 7 example (A = diag(0,d), H = [[0,1],[1,-d]], delta = 2,
% theta = pi/4, ||R|| = 1) attains delta sin 2theta = 2 = 2||R||. Proof-line
% slip only; the stated theorem is unaffected.
The perturbation estimate uses the full angle, and the residual estimate
the directed angle on source coordinates. In the unbounded residual form,
$\sym{A}$ and $\sym{Acurrent}$ are self-adjoint on a common domain; only
the residual is required to be bounded.

We formalized both theorems in Hilbert space using spectral calculus.

\subsection{Approximation numbers and operator-ideal norms}

In finite dimensions, a unitarily invariant norm is a symmetric gauge
function of the singular values. Approximation numbers extend the
singular-value comparisons needed for bounded operators; for compact
operators they give the usual singular-value sequence. We formalized
symmetric norming functions and Ky Fan partial sums to pass from
majorization to norm inequalities.

For example, let $\sym{Rrotation}$ be the direct rotation and
$\sym{Wcompetitor}$ any unitary carrying $\sym{Usource}$ onto
$\sym{Vtarget}$. The restricted displacement of the direct rotation has
singular-value minimality, which implies minimality for unitarily invariant
norms. The squared-displacement result compares
$(\sym{Identity}-\sym{Rrotation})^*(\sym{Identity}-\sym{Rrotation})$ and
$(\sym{Identity}-\sym{Wcompetitor})^*(\sym{Identity}-\sym{Wcompetitor})$
\citep[Proposition~4.3]{davis1970rotation}. Their singular values are the
squares of those of the unsquared displacements. Section~\ref{sec:prop44}
shows why the claimed trace-norm comparison fails and derives the valid
Q-norm conclusion.
% Additional formalized quantitative results, omitted from this article:
% DavisKahan/FiniteDimensional/DirectRotation/SharpDisplacement.lean and
% SharpDisplacementLimit.lean. AUTHORING_NOTES.md indexes the evidence.

The same norm inequalities enter the sine-theta proof. The assertion
$\sym{delta}\sym{NormN}(\sym{S})\leq\sym{NormN}(\sym{R})$ applies when
$\sym{S}$ and $\sym{R}$ have finite ideal norm. Our ideal-membership
hypotheses state this condition explicitly in arbitrary Hilbert dimension.

\subsection{Sylvester estimates and unbounded operators}

For operators $\sym{SylvesterLeft}$ and $\sym{SylvesterRight}$ on the target and source spaces of $\sym{X}$, the
Sylvester map is $\sym{X}\mapsto \sym{SylvesterLeft}\sym{X}-\sym{X}\sym{SylvesterRight}$. Under the interval/exterior or ordered
spectral-separation hypotheses, the estimate has the form
\[
 \sym{delta} \sym{NormN}(\sym{X})\leq \sym{NormN}(\sym{SylvesterLeft}\sym{X}-\sym{X}\sym{SylvesterRight}).
\]
The sine-theta proof uses $\sym{SylvesterLeft}=\sym{Lambda1}$, $\sym{SylvesterRight}=\sym{TrialOperator}$, and
$\sym{X}=\like{F1}{\symglyph{F1}^*}\sym{E0}$. Other perturbation estimates choose related
projection blocks for $\sym{X}$.

For unbounded $\sym{SylvesterLeft}$ and $\sym{SylvesterRight}$, the identity is asserted for
$x\in\DKDom(\sym{SylvesterRight})$, with $\sym{X}x\in\DKDom(\sym{SylvesterLeft})$, and $\sym{SylvesterLeft}\sym{X}-\sym{X}\sym{SylvesterRight}$ agrees there with a
bounded residual. We formalized these domain conditions using linear
partial maps. Spectral cutoffs give bounded approximations; convergence
of approximation numbers passes the norm estimates to the unbounded
operators \citep[Section~5 and Appendix to Section~6]{davis1970rotation}.

The generalized sine and tangent theorems also compare block operators.
Reflection pinching replaces a block matrix $T$ by
$(T+\sym{Jreflect}T\sym{Jreflect})/2$, where $\sym{Jreflect}$ acts as $+\sym{Identity}$ on one coordinate subspace and $-\sym{Identity}$
on the other. This removes the off-diagonal blocks and contracts a unitarily
invariant norm. Direct-sum comparisons relate the norms of the component
blocks to the whole operator. Together with Sylvester separation these give
the generalized angle estimates
\citep[Lemmas~6.1--6.2 and Theorems~6.1--6.3]{davis1970rotation}.

\subsection{Real and complex Hilbert spaces}

The spectral constructions use complex functional calculus. We formalized
the scalar-transfer results needed to carry these arguments to real Hilbert
spaces through complexification. This permits the source theorems to be
stated over either field, including their domain and ideal-norm hypotheses.
The underlying algebra, topology, and continuous functional calculus come
from Mathlib; the adapted Spectra constructions and later Tau Ceti
contributions are described in Section~\ref{sec:related-formal}.

\section{A counterexample to full-displacement optimality}
\label{sec:prop44}

This section gives an explicit real counterexample to the small-angle
optimality assertion \citep[Proposition~4.4]{davis1970rotation} and
identifies the invalid inequality in its proof. To the best of our knowledge,
this contradiction of the printed real-space claim has not previously been
published. Its complex-phase construction is already in the authors'
example \citep[Example~4.2]{davis1970rotation}; the distinction between trace-norm failure and Q-norm optimality
also has earlier treatments \citep{bolotnikov2001minimal}. We give the
counterexample, the failed comparison, and a Q-norm repair below.

Write $\sym{Psource}$ and $\sym{Qtarget}$ for the orthogonal projections
onto subspaces $\sym{Usource}$ and $\sym{Vtarget}$. For an orthogonal or
unitary map $\sym{Wcompetitor}$ with
$\sym{Wcompetitor}(\sym{Usource})=\sym{Vtarget}$, the vector
$(\sym{Identity}-\sym{Wcompetitor})x$ is the displacement of $x$ up to sign. Full
displacement measures this operator on the whole ambient space; restricted
displacement measures $(\sym{Identity}-\sym{Wcompetitor})\sym{Psource}$, which tests only
vectors in the source subspace. We use the trace norm
$\dkTraceNorm{Z}=\sum_j \sym{SingularValue}(Z)$.

For a finite-dimensional acute pair, the direct rotation $\sym{Rrotation}$
is the orthogonal or unitary polar factor of
\[
 \sym{IntertwinerT}=\sym{Qtarget}\sym{Psource}
 +(\sym{Identity}-\sym{Qtarget})(\sym{Identity}-\sym{Psource}).
\]
The printed claim \citep[Proposition~4.4]{davis1970rotation} asserts that,
for real subspaces with largest principal angle
at most $\sym{PiConstant}/3$, it minimizes $\sym{NormN}(\sym{Identity}-\sym{Wcompetitor})$ for every
unitarily invariant norm $\sym{NormN}$.

\subsection{The real four-dimensional counterexample}

The competitor below rotates one real two-plane and fixes its orthogonal
complement. The direct rotation moves both principal planes. The competitor
has a smaller trace norm but a larger operator norm of its displacement.

\begin{theorem}[Direct rotation can fail to minimize the trace norm at any acute angle]
\label{thm:small-angle}
For every $0<\sym{theta}<\sym{PiConstant}/2$, there are two-dimensional subspaces
$\sym{Usource},\like{Vtarget}{\symglyph{Vtarget}_{\symglyph{theta}}}\subset\sym{RealField}^4$ and an orthogonal map $\like{Wcompetitor}{\symglyph{Wcompetitor}_{\symglyph{theta}}}$ with
$\like{Wcompetitor}{\symglyph{Wcompetitor}_{\symglyph{theta}}} \sym{Usource}=\like{Vtarget}{\symglyph{Vtarget}_{\symglyph{theta}}}$, such that both principal angles equal $\sym{theta}$ and
the canonical direct rotation $\like{Rrotation}{\symglyph{Rrotation}_{\symglyph{theta}}}$ satisfies
\begin{align}
 \|\sym{Identity}-\like{Wcompetitor}{\symglyph{Wcompetitor}_{\symglyph{theta}}}\|_1&=4\sin\sym{theta},\label{eq:competitor-trace}\\
 \|\sym{Identity}-\like{Rrotation}{\symglyph{Rrotation}_{\symglyph{theta}}}\|_1&=8\sin(\sym{theta}/2)>4\sin\sym{theta}.
 \label{eq:direct-trace}
\end{align}
\end{theorem}

\begin{proof}
Put $\sym{CosAngle}=\cos\sym{theta}$, $\sym{SinAngle}=\sin\sym{theta}$, and let
\[
 \sym{Jturn}=\begin{pmatrix}0&-1\\1&0\end{pmatrix},\quad
 \sym{Cturn}=\cos(2\sym{theta})\sym{IdentityTwo}+\sin(2\sym{theta})\sym{Jturn},\quad
 \sym{Hmix}=\frac{1}{\sqrt2}\begin{pmatrix}\sym{IdentityTwo}&\sym{IdentityTwo}\\I_2&-\sym{IdentityTwo}\end{pmatrix}.
\]
Define
\begin{equation}
 \like{Wcompetitor}{\symglyph{Wcompetitor}_{\symglyph{theta}}}=\sym{Hmix}\begin{pmatrix}\sym{Cturn}&0\\0&\sym{IdentityTwo}\end{pmatrix}\sym{Hmix}^T
 =\begin{pmatrix}\like{CosAngle}{\symglyph{CosAngle}^2}\sym{IdentityTwo}+\sym{SinAngle}\sym{CosAngle}\sym{Jturn}&-\like{SinAngle}{\symglyph{SinAngle}^2}\sym{IdentityTwo}+\sym{SinAngle}\sym{CosAngle}\sym{Jturn}\\
                  -\like{SinAngle}{\symglyph{SinAngle}^2}\sym{IdentityTwo}+\sym{SinAngle}\sym{CosAngle}\sym{Jturn}&\like{CosAngle}{\symglyph{CosAngle}^2}\sym{IdentityTwo}+\sym{SinAngle}\sym{CosAngle}\sym{Jturn}\end{pmatrix}.
 \label{eq:wtheta}
\end{equation}
Both $\sym{Hmix}$ and $\sym{Cturn}$ are orthogonal, so $\like{Wcompetitor}{\symglyph{Wcompetitor}_{\symglyph{theta}}}$ is orthogonal. Let
$\sym{Usource}=\sym{RealField}^2\oplus\{\sym{ZeroVector}\}$ and $\like{Vtarget}{\symglyph{Vtarget}_{\symglyph{theta}}}=\like{Wcompetitor}{\symglyph{Wcompetitor}_{\symglyph{theta}}} \sym{Usource}$.
The first two columns of $\like{Wcompetitor}{\symglyph{Wcompetitor}_{\symglyph{theta}}}$ factor as
\[
 \begin{pmatrix}\like{CosAngle}{\symglyph{CosAngle}^2}\sym{IdentityTwo}+\sym{SinAngle}\sym{CosAngle}\sym{Jturn}\\-\like{SinAngle}{\symglyph{SinAngle}^2}\sym{IdentityTwo}+\sym{SinAngle}\sym{CosAngle}\sym{Jturn}\end{pmatrix}
 =\begin{pmatrix}\sym{CosAngle}\sym{IdentityTwo}\\\sym{SinAngle}\sym{Jturn}\end{pmatrix}(\sym{CosAngle}\sym{IdentityTwo}+\sym{SinAngle}\sym{Jturn}).
\]
Since $\sym{CosAngle}\sym{IdentityTwo}+\sym{SinAngle}\sym{Jturn}$ is orthogonal, the columns of
$\sym{Yframe}=\binom{\sym{CosAngle}\sym{IdentityTwo}}{\sym{SinAngle}\sym{Jturn}}$ form an orthonormal basis of $\like{Vtarget}{\symglyph{Vtarget}_{\symglyph{theta}}}$.
With $\sym{Xframe}=\binom{\sym{IdentityTwo}}{0}$, one has $\sym{Xframe}^T\sym{Yframe}=\sym{CosAngle}\sym{IdentityTwo}$.
Thus both principal cosines are $\sym{CosAngle}>0$, and both principal angles are
$\sym{theta}$.

The corresponding projections and canonical intertwiner are
\[
 \sym{Psource}=\begin{pmatrix}\sym{IdentityTwo}&0\\0&0\end{pmatrix},\qquad
 \sym{Qtarget}=\sym{Yframe}\sym{Yframe}^T=\begin{pmatrix}\like{CosAngle}{\symglyph{CosAngle}^2}\sym{IdentityTwo}&-\sym{SinAngle}\sym{CosAngle}\sym{Jturn}\\\sym{SinAngle}\sym{CosAngle}\sym{Jturn}&\like{SinAngle}{\symglyph{SinAngle}^2}\sym{IdentityTwo}\end{pmatrix},
\]
\begin{equation}
 \sym{IntertwinerT}=\sym{Qtarget}\sym{Psource}+(\sym{Identity}-\sym{Qtarget})(\sym{Identity}-\sym{Psource})
   =\sym{CosAngle}\begin{pmatrix}\sym{CosAngle}\sym{IdentityTwo}&\sym{SinAngle}\sym{Jturn}\\\sym{SinAngle}\sym{Jturn}&\sym{CosAngle}\sym{IdentityTwo}\end{pmatrix}.
 \label{eq:intertwiner-theta}
\end{equation}
The last block matrix is orthogonal and $\sym{CosAngle}>0$, so $|\sym{IntertwinerT}|=\sym{CosAngle}\sym{IdentityFour}$ and
its polar factor is exactly
\[
 \like{Rrotation}{\symglyph{Rrotation}_{\symglyph{theta}}}=\begin{pmatrix}\sym{CosAngle}\sym{IdentityTwo}&\sym{SinAngle}\sym{Jturn}\\\sym{SinAngle}\sym{Jturn}&\sym{CosAngle}\sym{IdentityTwo}\end{pmatrix}.
\]
By the polar-factor characterization in \cref{eq:history-kato-unitary},
this is the canonical direct rotation for the constructed pair.
Since $\like{Rrotation}{\symglyph{Rrotation}_{\symglyph{theta}}}+\like{Rrotation}{\symglyph{Rrotation}_{\symglyph{theta}}}^T=2\sym{CosAngle}\sym{IdentityFour}$,
\[
 (\sym{Identity}-\like{Rrotation}{\symglyph{Rrotation}_{\symglyph{theta}}})^T(\sym{Identity}-\like{Rrotation}{\symglyph{Rrotation}_{\symglyph{theta}}})=(2-2\sym{CosAngle})\sym{IdentityFour}.
\]
Consequently all four displacement singular values equal
$\sqrt{2-2\sym{CosAngle}}=2\sin(\sym{theta}/2)$.

On the other hand, \cref{eq:wtheta} shows that the singular values of
$\sym{Identity}-\like{Wcompetitor}{\symglyph{Wcompetitor}_{\symglyph{theta}}}$ are those of $(\sym{IdentityTwo}-\sym{Cturn})\oplus0_2$.
The identity
$(\sym{IdentityTwo}-\sym{Cturn})^T(\sym{IdentityTwo}-\sym{Cturn})=4\like{SinAngle}{\symglyph{SinAngle}^2}\sym{IdentityTwo}$ gives the list
$(2\sym{SinAngle},2\sym{SinAngle},0,0)$. The two trace formulas follow. Finally,
\[
 4\sin\sym{theta}
 =8\sin(\sym{theta}/2)\cos(\sym{theta}/2)
 <8\sin(\sym{theta}/2),
\]
because $0<\sym{theta}/2<\sym{PiConstant}/4$.
\end{proof}

\begin{corollary}[No positive small-angle threshold]
\label{cor:no-threshold}
For every $\sym{EpsilonAngle}>0$, there are real two-dimensional subspaces of
$\sym{RealField}^4$ whose largest principal angle is strictly between zero and
$\sym{EpsilonAngle}$ and is at most $\sym{PiConstant}/3$, together with an orthogonal competitor
having strictly smaller trace-norm displacement than the direct rotation.
\end{corollary}

\begin{proof}
Apply \cref{thm:small-angle} with
$\sym{theta}=\min\{\sym{EpsilonAngle}/2,\sym{PiConstant}/4\}$.
\end{proof}

\noindent\begin{minipage}{\linewidth}
Take both principal angles to be $\sym{theta}=30^\circ$. The orthogonal map
in \cref{eq:wtheta} has two nonzero displacement singular values, both
equal to $1$. The direct rotation has four displacement singular values,
each equal to $2\sin 15^\circ$. Hence
\begin{equation}\label{eq:prop44-thirty-degrees}
 \|\sym{Identity}-\like{Wcompetitor}{\symglyph{Wcompetitor}_{\symglyph{theta}}}\|_1
 =4\sin 30^\circ=2
 <8\sin 15^\circ
 =\|\sym{Identity}-\like{Rrotation}{\symglyph{Rrotation}_{\symglyph{theta}}}\|_1 .
\end{equation}
The direct rotation costs approximately $2.071$, and the largest angle
$30^\circ$ is below the printed bound of $60^\circ$.
\Cref{lst:prop44-refutation} verifies the transport equation, the principal
angles, and the strict trace-norm comparison for every positive acute angle.

% Source: DavisKahan/FiniteDimensional/DirectRotation/SmallAngleCounterexample.lean,
% TauCeti.DavisKahan.FiniteDimensional.SmallAngleCounterexample.explicit_traceNorm_counterexample.
% U4 is defined in ShortRotationCounterexample.lean.
\begin{lstlisting}[style=leanpaper,
  caption={The explicit counterexample to Proposition~4.4.
  In $\sym{RealField}^4$, \texttt{U4} is the first coordinate plane,
  \texttt{Wequiv theta} is the orthogonal map in~\eqref{eq:wtheta},
  and \texttt{V4 theta} is its image of \texttt{U4}.
  Principal angles are indexed from zero in decreasing order.
  \texttt{traceDisplacement W} is $\|\sym{Identity}-W\|_1$.},
  label={lst:prop44-refutation}]
theorem explicit_traceNorm_counterexample
    (theta : Real) (hpos : 0 < theta)
    (hangle : theta < Real.pi / 2) :
    let V := V4 theta
    let W := Wequiv theta
    let R := directRotation U4 V
      (acute theta (cos_pos theta hpos hangle))
    U4.map W.toLinearMap = V (*@$\land$@*)
      U4.principalAngles V 0 = theta (*@$\land$@*)
      U4.principalAngles V 1 = theta (*@$\land$@*)
      traceDisplacement W < traceDisplacement R := by (*@\laaettr@*)
\end{lstlisting}
\end{minipage}\par

The competitor rotates a diagonal real two-plane through $2\sym{theta}$ and
fixes its orthogonal complement; its displacement has rank two, compared
with rank four for $\like{Rrotation}{\symglyph{Rrotation}_{\symglyph{theta}}}$. The operator norm favors the canonical
rotation instead:
$\dkOperatorNorm{\sym{Identity}-\like{Rrotation}{\symglyph{Rrotation}_{\symglyph{theta}}}}=2\sin(\sym{theta}/2)
 <2\sin\sym{theta}=\dkOperatorNorm{\sym{Identity}-\like{Wcompetitor}{\symglyph{Wcompetitor}_{\symglyph{theta}}}}$.
The competitor attains the global trace-norm minimum $4\sin\sym{theta}$.
Indeed, the projection difference has trace norm
$\|\sym{Psource}-\sym{Qtarget}\|_1=4\sin\sym{theta}$, and every orthogonal
transporter $\sym{Wcompetitor}$ satisfies
$\|\sym{Psource}-\sym{Qtarget}\|_1\leq\|\sym{Identity}-\sym{Wcompetitor}\|_1$.
Equation~\eqref{eq:competitor-trace} attains this lower bound.
% Formal proof: DavisKahan/FiniteDimensional/DirectRotation/TraceOptimality.lean,
% SmallAngleCounterexample.traceDisplacement_lower_bound and Wequiv_traceMinimizer.
% Dimension minimality is now proved in DirectRotation/DimensionMinimality.lean:
% nuclear_directRotation_le_of_finrank_le_three compares the actual direct
% rotation with every ambient real orthogonal transporter in dimension <= 3,
% under the acute-pair and largest-principal-angle <= pi/3 hypotheses.
% isLeast_real_shortAngle_nuclear_counterexample_dimension proves the least
% dimension is four; the same file retains arbitrarily small positive witnesses.
% Proposed corollary for the next prose review: Four is the least real ambient
% dimension for trace-norm failure under the printed small-angle hypotheses.

At $\sym{theta}=\sym{PiConstant}/4$, \cref{fig:prop44-displacement-planes}
shows the different displacement patterns of the canonical direct rotation
and the competing transporter. Both carry the same source subspace onto the
same target, with both principal angles equal to $45^\circ$. The competitor
fixes half of the ambient space; the direct rotation moves both principal planes.

\begin{figure}[tbp]
\centering
\resizebox{119mm}{!}{\input{figures/Fig2.pgf}}
\caption{A counterexample to the full-displacement claim
\citep[Proposition~4.4]{davis1970rotation} at
$\sym{theta}=\sym{PiConstant}/4$. Both orthogonal maps carry $\sym{Usource}$ onto
$\sym{Vtarget}$. Each circle shows the action on one invariant real
two-plane: (a,b) the principal planes $\Pi_1,\Pi_2$ for $\sym{Rrotation}$,
and (c,d) the moving and fixed planes $M,Z$ for $\sym{Wcompetitor}$.
The unit inputs are $u_1=e_0$, $u_2=e_1$,
$x=(e_0+e_2)/\sqrt2$, and $z=(e_0-e_2)/\sqrt2$.
For $T=\sym{Rrotation}$ or $\sym{Wcompetitor}$, the green arrow from
the image $Tq$ to the input $q$ is the displacement $(\sym{Identity}-T)q$.
The bars list the four singular values of $\sym{Identity}-T$, grouped by contributing
plane, on the same vertical scale. Their sums give
$\|\sym{Identity}-\sym{Wcompetitor}\|_1=2\sqrt2<4\sqrt{2-\sqrt2}=\|\sym{Identity}-\sym{Rrotation}\|_1$.
The two maps have different invariant-plane decompositions}
\label{fig:prop44-displacement-planes}
\end{figure}

\subsection{Failure of the published proof}

At $\sym{theta}=\sym{PiConstant}/4$, the family gives
\[
 \like{Wcompetitor}{\symglyph{Wcompetitor}_{\symglyph{PiConstant}/4}}=\frac12
 \begin{pmatrix}
 1&-1&-1&-1\\1&1&1&-1\\-1&-1&1&-1\\1&-1&1&1
 \end{pmatrix}.
\]
The two trace norms are
\[
 \|\sym{Identity}-\like{Wcompetitor}{\symglyph{Wcompetitor}_{\symglyph{PiConstant}/4}}\|_1=2\sqrt2
 <4\sqrt{2-\sqrt2}=\|\sym{Identity}-\like{Rrotation}{\symglyph{Rrotation}_{\symglyph{PiConstant}/4}}\|_1.
\]

For $\sym{Kdisplacement}=\sym{Identity}-\like{Wcompetitor}{\symglyph{Wcompetitor}_{\symglyph{PiConstant}/4}}$, the even-index part
of the published comparison \citep[equation~(4.3)]{davis1970rotation}
reads, in our notation,
\[
 \|\sym{Kdisplacement}\|_{(\sym{KyFanEvenIndex})}\geq
 \sum_{k=1}^{\sym{KyFanEvenIndex}/2}\|\sym{Kdisplacement}\like{PrincipalPlaneProjectionBase}{\symglyph{PrincipalPlaneProjectionBase}_k}\|_{(2)},
 \qquad \sym{KyFanEvenIndex}\ \text{even},
\]
where $\sym{KyFanNormNotation}$ is the Ky Fan $j$-norm and $\like{PrincipalPlaneProjectionBase}{\symglyph{PrincipalPlaneProjectionBase}_k}$ projects onto
the $k$th principal plane. Choose $\sym{KyFanEvenIndex}=4$ and let
$\sym{OmegaOne}$ and $\sym{OmegaTwo}$ project onto
$\sym{SpanOperator}\{e_0,e_3\}$ and
$\sym{SpanOperator}\{e_1,e_2\}$. These are principal planes for the pair.
The comparison then requires
\begin{equation}
 \|\sym{Kdisplacement}\|_{(4)}\geq
 \|\sym{Kdisplacement}\sym{OmegaOne}\|_{(2)}+\|\sym{Kdisplacement}\sym{OmegaTwo}\|_{(2)},
 \label{eq:dk43-false}
\end{equation}
Direct calculation gives
\[
 \|\sym{Kdisplacement}\|_{(4)}=2\sqrt2,\qquad
 \|\sym{Kdisplacement}\sym{OmegaOne}\|_{(2)}=\|\sym{Kdisplacement}\sym{OmegaTwo}\|_{(2)}=2.
\]
Thus \cref{eq:dk43-false} would imply $2\sqrt2\geq4$.
The displayed matrix also gives
\[
 \DKRange(\sym{Kdisplacement}\sym{OmegaOne})
 =\DKRange(\sym{Kdisplacement}\sym{OmegaTwo})
 =\sym{CounterMovingPlane}
 =\sym{SpanOperator}\{e_0+e_2,e_1+e_3\}.
\]
Both restrictions map onto the same two-dimensional output plane, even
though their input planes are orthogonal. Each has singular values $1,1$,
whereas the full displacement has just two nonzero singular values,
$\sqrt2,\sqrt2$. The comparison in \cref{eq:dk43-false} would follow if the
singular-value list of the full displacement were the concatenation of
the two lists for the restrictions. This argument would be valid for an
orthogonal direct sum, but it does not apply here because the output
planes coincide.

\subsection{Repair of Davis and Kahan's Proposition~4.4: direct-rotation optimality for Q-norms}

Call a unitarily invariant norm $\sym{NormN}$ a $\sym{NormQ}$-norm when there is a unitarily
invariant norm $\sym{NormM}$ satisfying
\begin{equation}
 \sym{NormN}(Z)^2=\sym{NormM}(Z^*Z)\qquad\text{for every }Z.
 \label{eq:qnorm-def}
\end{equation}

For example, the Frobenius norm is a $\sym{NormQ}$-norm because
$\dkFrobeniusNorm{Z}^2=\dkTraceNorm{Z^*Z}$: its squared norm is the sum of the eigenvalues of
the positive matrix $Z^*Z$. This identity uses the trace norm for
$\sym{NormM}$. The squared-displacement comparison
\citep[Proposition~4.3]{davis1970rotation} gives the following Q-norm consequence.
Earlier Q-norm minimization results appear in
\citet[Theorem~3.1]{bolotnikov2001minimal}.

\begin{proposition}[$\sym{NormQ}$-norm consequence]
\label{prop:qnorm-repair}
Let $\sym{Usource},\sym{Vtarget}$ be an acute pair of subspaces of a finite-dimensional real or
complex Hilbert space, and let $\sym{Rrotation}$ be their canonical direct rotation.
For every $\sym{NormQ}$-norm $\sym{NormN}$ and every orthogonal or unitary map $\sym{Wcompetitor}$ with $\sym{Wcompetitor}\sym{Usource}=\sym{Vtarget}$,
\begin{equation}
 \sym{NormN}(\sym{Identity}-\sym{Rrotation})\leq \sym{NormN}(\sym{Identity}-\sym{Wcompetitor}).
 \label{eq:qnorm-repair}
\end{equation}
No $\sym{PiConstant}/3$ restriction is required.
\end{proposition}

\begin{proof}
The valid squared-displacement theorem gives
\[
 \sym{NormM}((\sym{Identity}-\sym{Rrotation})^*(\sym{Identity}-\sym{Rrotation}))\leq \sym{NormM}((\sym{Identity}-\sym{Wcompetitor})^*(\sym{Identity}-\sym{Wcompetitor})).
\]
By \cref{eq:qnorm-def}, these quantities are $\sym{NormN}(\sym{Identity}-\sym{Rrotation})^2$ and $\sym{NormN}(\sym{Identity}-\sym{Wcompetitor})^2$.
Both norm values are nonnegative, so taking square roots proves the claim.
\end{proof}

The trace norm on the counterexample space does not satisfy the
$\sym{NormQ}$-norm condition. Its smaller competing displacement in
\cref{thm:small-angle} contradicts the minimality that this condition
would imply.

\subsection{Historical context}

Davis and Kahan already exhibited a complex two-dimensional phase
modification that decreases the trace norm of full displacement
\citep[Example~4.2]{davis1970rotation}. Our real family is its
realification. Indeed, under the identification $\sym{ComplexField}^2\simeq
\sym{RealField}^2\oplus\sym{RealField}^2$ in which multiplication by $i$ is
represented by $\sym{Jturn}$ from \cref{eq:wtheta}, the matrices
\[
 \like{Rrotation}{\symglyph{Rrotation}_{\symglyph{theta}}}^{\sym{ComplexField}}=
 \begin{pmatrix}\sym{CosAngle}&i\sym{SinAngle}\\i\sym{SinAngle}&\sym{CosAngle}\end{pmatrix},\qquad
 \like{Wcompetitor}{\symglyph{Wcompetitor}_{\symglyph{theta}}}^{\sym{ComplexField}}=e^{i\sym{theta}}\like{Rrotation}{\symglyph{Rrotation}_{\symglyph{theta}}}^{\sym{ComplexField}}
\]
become precisely $\like{Rrotation}{\symglyph{Rrotation}_{\symglyph{theta}}}$ and $\like{Wcompetitor}{\symglyph{Wcompetitor}_{\symglyph{theta}}}$ above. Realification doubles
the complex singular-value multiplicities. Thus the real small-angle
restriction in the full-displacement claim
\citep[Proposition~4.4]{davis1970rotation} does not eliminate the mechanism of the
authors' preceding example.

\citet{bolotnikov2001minimal} studied related unitary completion and
subspace-mapping problems. Their results give a $\sym{NormQ}$-norm
minimization theorem and a counterexample for Schatten $p$-norms with
$1\leq p<2$ \citep[Theorem~3.1 and Lemma~3.3]{bolotnikov2001minimal}.
Those results establish the earlier norm-class
distinction. \citet[Section~5]{qiu2006grassmann} subsequently repeated
the $\sym{PiConstant}/3$ assertion, whereas
\citet[Section~4, p.~1615]{zhang2010canonicalangles} explicitly noted
general trace-norm nonminimality without addressing that particular
real-space restriction. The present refutation identifies the invalid
inequality in the printed full-displacement claim
\citep[Proposition~4.4]{davis1970rotation}. The Q-norm conclusion above follows
from the earlier squared-displacement extremality statement.


\section{A follow-on theorem enabled by the formal development}
\label{sec:example-transport}

We prove an exact formula for the smallest trace-norm displacement of a real
orthogonal map carrying one given subspace onto another in finite dimension.
The formula reduces this optimization to maximum-weight matching on a finite
general graph whose vertices label the positive principal angles.
Its proof uses the principal-angle geometry, real and complex operator
  constructions, and singular-value norms developed for the formalization,
  together with additional optimization arguments.

Let $\sym{MatchingSourceSubspace},\sym{MatchingTargetSubspace}\subseteq\sym{RealField}^n$ have the same dimension and let
$\like{theta}{\symglyph{theta}_1},\ldots,\like{theta}{\symglyph{theta}_m}\in(0,\sym{PiConstant}/2]$ be their positive principal angles,
counted with multiplicity. A matching $\sym{Matching}$ chooses unordered
pairs of labels from $\{1,\ldots,m\}$, using each label at most once.
Labels that are not used in a pair are unmatched.
For three labels, $\sym{Matching}=\{\{1,2\}\}$ pairs labels $1,2$ and
leaves $3$ unmatched. The choice $\{\{1,2\},\{2,3\}\}$ is disallowed:
label $2$ would occur in two pairs.

Each positive principal angle describes a source line and a target line
in a two-dimensional principal plane. Transporting that plane independently
allows a rotation through $\like{theta}{\symglyph{theta}_i}$, with trace displacement
$4\sin(\like{theta}{\symglyph{theta}_i}/2)$, or a reflection carrying the source line
to the target, with trace displacement $2$. Pairing two labels $\{i,j\}$
allows a single orthogonal map to act on both principal planes together,
on their four-dimensional sum.
For each possible pair of distinct labels $i,j$, define its paired cost
$\like{PairCost}{\symglyph{PairCost}_{ij}}$; for each unmatched label $i$, define its singleton
cost $\like{SingletonCost}{\symglyph{SingletonCost}_i}$:
\begin{equation}\label{eq:nuclear-matching-costs}
 \like{PairCost}{\symglyph{PairCost}_{ij}}=4\sin\bigl((\like{theta}{\symglyph{theta}_i}+\like{theta}{\symglyph{theta}_j})/2\bigr),\qquad
 \like{SingletonCost}{\symglyph{SingletonCost}_i}=\min\{4\sin(\like{theta}{\symglyph{theta}_i}/2),2\}.
\end{equation}
The paired cost is attained by an orthogonal transporter on the two-plane sum.
It fixes two real directions and rotates the other two through
$\like{theta}{\symglyph{theta}_i}+\like{theta}{\symglyph{theta}_j}$, as the competitor does for
equal angles in \cref{fig:prop44-displacement-planes}. Step~6 below derives
this construction and its cost from the block minimum.
Write $\sym{OrthogonalGroup}$ for the orthogonal group of real
$n\times n$ matrices, defined by $\sym{Wcompetitor}^T\sym{Wcompetitor}=\sym{Identity}$.

% Note:
% this corresponds to exists_minimizing_real_nuclear_transport_matching
% in ../../DavisKahan/FiniteDimensional/DirectRotation/RealNuclearMatching.lean
%  It states that an orthogonal transporter attains the matching cost and
%  minimizes trace displacement over every orthogonal competitor.
%  The matching costs are defined in RealMatchingExponent.lean realNuclearMatchingCost
%  (DavisKahan/FiniteDimensional/DirectRotation/RealMatchingExponent.lean:136).
%  The endpoint proof uses the 1<p<2 theorem in RealMatchingAttainment.lean:296
%  (DavisKahan/FiniteDimensional/DirectRotation/RealMatchingAttainment.lean:296)
%  exists_minimizing_real_schatten_transport_matching
\begin{theorem}[Real trace-norm matching reduction]%
\label{thm:real-matching-example}
For equal-dimensional real subspaces $\sym{MatchingSourceSubspace},\sym{MatchingTargetSubspace}$ as above,
\begin{equation}\label{eq:real-nuclear-match}
\min_{\substack{\sym{Wcompetitor}\in \sym{OrthogonalGroup}\\\sym{Wcompetitor}(\sym{MatchingSourceSubspace})=\sym{MatchingTargetSubspace}}}\|\sym{Identity}-\sym{Wcompetitor}\|_1
=
\min_{\sym{Matching}\text{ a matching}}
\left(
\sum_{\{i,j\}\in \sym{Matching}}\like{PairCost}{\symglyph{PairCost}_{ij}}
+
\sum_{i\text{ unmatched}}\like{SingletonCost}{\symglyph{SingletonCost}_i}
\right).
\end{equation}
Both minima are attained.  The equality includes repeated angles,
right angles, and common ambient directions; the left side has no
principal-plane or determinant constraint.
\end{theorem}

We formalized this theorem in Lean~4, together with the corresponding
matching reduction for the Schatten power cost
$\|\sym{Identity}-\sym{Wcompetitor}\|_p^p$ for $1\le p<2$.
The exponent-dependent matching costs are developed in the proof below.
\Cref{lst:real-transport-matching} states both results.

% Source declarations (namespace TauCeti.DavisKahan.FiniteDimensional.RealTransport):
% DavisKahan/FiniteDimensional/DirectRotation/RealNuclearMatching.lean:
%   positivePrincipalAngles, realTransportCosts,
%   real_nuclear_transport_matching, real_schatten_transport_matching.
% The endpoint identity is RealMatchingExponent.lean:
%   realPositiveSchattenMatchingMinimum_one.
\noindent\begin{minipage}{\linewidth}
\begin{lstlisting}[style=leanpaper,
  caption={Lean presentation theorems for the trace-norm and Schatten
  matching reductions. The omitted context supplies a finite-dimensional
  real inner-product space \texttt{E}. The set \texttt{realTransportCosts p K L}
  contains the power costs $\|\sym{Identity}-T\|_p^p$ of orthogonal transporters
  $T$ carrying $K$ onto $L$. The function \texttt{positivePrincipalAngles}
  retains the labels of the strictly positive principal angles.},
  label={lst:real-transport-matching}]
theorem real_nuclear_transport_matching
    (K L : Submodule Real E)
    (hKL : Module.finrank Real K = Module.finrank Real L) :
    IsLeast (realTransportCosts 1 K L)
      (realNuclearMatchingMinimum
        (positivePrincipalAngles K L)) := by
  (*@\laaettr@*)

theorem real_schatten_transport_matching
    (p : Real) (hp1 : 1 <= p) (hp2 : p < 2)
    (K L : Submodule Real E)
    (hKL : Module.finrank Real K = Module.finrank Real L) :
    IsLeast (realTransportCosts p K L)
      (realPositiveSchattenMatchingMinimum p
        (positivePrincipalAngles K L)) := by
  (*@\laaettr@*)
\end{lstlisting}
\end{minipage}\par

In Lean, \texttt{IsLeast S x} states that $x$ belongs to $S$ and is less
than or equal to every element of $S$. Thus both signatures assert
attainment and global optimality. At $p=1$ the two matching minima coincide.

To express the formula as maximum-weight matching, start from the
all-singleton cost $\sym{SingletonBaseline}=\sum_i \like{SingletonCost}{\symglyph{SingletonCost}_i}$ and assign each edge $\{i,j\}$ the weight
\[
 \sym{MatchingEdgeWeight}=\like{SingletonCost}{\symglyph{SingletonCost}_i}+\like{SingletonCost}{\symglyph{SingletonCost}_j}-\like{PairCost}{\symglyph{PairCost}_{ij}}.
\]
This weight is the difference between the cost of two singletons and the
paired cost. Thus a matching $\sym{Matching}$ has cost
\[
 \like{TraceCost}{\symglyph{TraceCost}(M)}=\sym{SingletonBaseline}-
       \sum_{\{i,j\}\in\sym{Matching}}\sym{MatchingEdgeWeight}.
\]
Minimizing this cost is equivalent to maximizing the sum of edge weights.
The empty matching is allowed, and edges of negative weight can be omitted.
Figure~\ref{fig:real-nuclear-matching} shows a three-angle example.

\begin{figure}[tb]
\centering
\resizebox{119mm}{!}{\input{figures/Fig3.pgf}}
\caption{
Graph of principal-angle pairing choices for two three-dimensional subspaces
of $\sym{RealField}^6$ whose intersection is $\{\sym{ZeroVector}\}$ and whose principal angles are
$30^\circ$, $45^\circ$, and $75^\circ$.
(a) Vertex positions are arbitrary; each vertex represents an angle label and
each edge a possible pairing. Vertices carry singleton costs $\like{SingletonCost}{\symglyph{SingletonCost}_i}$;
edges carry weights $\sym{MatchingEdgeWeight}=\like{SingletonCost}{\symglyph{SingletonCost}_i}+\like{SingletonCost}{\symglyph{SingletonCost}_j}-\like{PairCost}{\symglyph{PairCost}_{ij}}$. The solid edge
pairs labels $1,2$ and leaves label $3$ unmatched. (b) The table lists all four
matchings; costs and edge weights are rounded to four decimals. Geometrically,
the selected pair is realized by a four-dimensional transporter and the
unmatched angle by a two-dimensional reflection. Their orthogonal sum
attains the minimum on $\sym{RealField}^6$, with trace displacement
$4\sin(37.5^\circ)+2\approx4.4350$}
\label{fig:real-nuclear-matching}
\end{figure}

\begin{proof}[Proof and explanation]
The trace norm, corresponding to $p=1$, need not be differentiable.
We first replace it by a differentiable displacement cost with $1<p<2$.
At a global minimizer, derivatives along constraint-preserving variations
vanish. We use these equations to decompose the minimizer into orthogonal
blocks carrying one or two principal angles. The blocks determine a matching;
conversely, every matching constructs a transporter with the same cost.
This proves equality of the two minimum values for $1<p<2$.
We then let $p$ approach $1$ from above to recover the trace norm.

\emph{Step 1: choose a differentiable cost and a global minimizer.}
For $1<p<2$, set
\[
 \sym{FeasibleTransport}=\{\sym{Wcompetitor}\in \sym{OrthogonalGroup}:
                 \sym{Wcompetitor}(\sym{MatchingSourceSubspace})=\sym{MatchingTargetSubspace}\},
 \qquad
 \like{PowerCost}{\symglyph{PowerCost}_p}(\sym{Wcompetitor})
       =\sum_{j=1}^n \sym{TransportSingularValues}(\sym{Identity}-\sym{Wcompetitor})^p.
\]
Here $\sym{TransportSingularValues}(T)$ denotes a singular value of $T$. For $p>1$ this cost is
differentiable, including at zero displacement singular values.

\emph{(a) The cost attains a global minimum on $\sym{FeasibleTransport}$.}
Write $\like{Psource}{\symglyph{Psource}_K}$ and $\sym{MatchingTargetProjection}$ for the orthogonal
projections onto the source and target.
\emph{Justification.} An orthogonal map carrying an
orthonormal basis of $\sym{MatchingSourceSubspace}$ to one of $\sym{MatchingTargetSubspace}$,
extended to their orthogonal complements, belongs to $\sym{FeasibleTransport}$.
Thus $\sym{FeasibleTransport}$ is nonempty. The group $\sym{OrthogonalGroup}$ is closed and bounded
in $\sym{RealField}^{n\times n}$, hence compact by Heine--Borel.
The equivalent transport equation
$\sym{Wcompetitor}\like{Psource}{\symglyph{Psource}_K}\sym{Wcompetitor}^T
 =\sym{MatchingTargetProjection}$ is a closed constraint.
Since $\like{PowerCost}{\symglyph{PowerCost}_p}$ is continuous, the extreme value theorem
gives a global minimizer $\sym{Wcompetitor}\in\sym{FeasibleTransport}$.

\emph{(b) A simultaneous invariant decomposition splits both the cost
and the principal-angle list.}
The decomposition we seek is an orthogonal sum
$\sym{RealField}^n=\bigoplus_a\sym{MatchingBlock}$ in which each $\sym{MatchingBlock}$ is preserved by
$\like{Psource}{\symglyph{Psource}_K}$, $\sym{MatchingTargetProjection}$, and $\sym{Wcompetitor}$.
\emph{Justification.} In a basis adapted to this sum all three matrices are block diagonal.
Consequently, the source and target split as the sums of their
intersections with the $\sym{MatchingBlock}$, and
\begin{equation}\label{eq:matching-block-additivity}
 \like{PowerCost}{\symglyph{PowerCost}_p}(\sym{Wcompetitor})
 =\sum_a\like{PowerCost}{\symglyph{PowerCost}_p}(\sym{Wcompetitor}|_{\sym{MatchingBlock}}).
\end{equation}
Indeed, the displacement is block diagonal too, and its singular-value
list is the union of the block singular-value lists. The principal-angle
lists also combine in this way: the compression
$\like{Psource}{\symglyph{Psource}_K}\sym{MatchingTargetProjection}\like{Psource}{\symglyph{Psource}_K}$
is block diagonal, and its restriction to $\sym{MatchingSourceSubspace}$ has the
squared principal cosines as eigenvalues.
These two additivity facts will turn the block decomposition into a
matching of the original angle labels.

\emph{Step 2: use derivatives along maps that still transport the source.}
Consider smooth curves $\sym{Wcompetitor}(t)\in\sym{FeasibleTransport}$, for real $t$
near zero, with $\sym{Wcompetitor}(0)=\sym{Wcompetitor}$.
The cost derivative at zero vanishes along every such curve because
$\sym{Wcompetitor}$ is a minimizer. This is the stationarity condition.
We use it as a necessary condition on the global minimizer already
obtained in Step~1.

\emph{(a) Source-preserving rotations give feasible curves.}
Let $\sym{SkewDirection}$ be skew-symmetric (\ie{},
$\sym{SkewDirection}^T=-\sym{SkewDirection}$).
\emph{Justification.} For real $t$, its exponential is orthogonal, since
\[
 \exp(t\sym{SkewDirection})^T
 =\exp(-t\sym{SkewDirection})
 =\exp(t\sym{SkewDirection})^{-1}.
\]
If $\sym{SkewDirection}$ commutes with $\like{Psource}{\symglyph{Psource}_K}$, so does its
exponential (every power of $\sym{SkewDirection}$ commutes with
$\like{Psource}{\symglyph{Psource}_K}$, hence so does the exponential power series).
Its inverse $\exp(-t\sym{SkewDirection})$ also commutes with
$\like{Psource}{\symglyph{Psource}_K}$, so the exponential maps the source onto itself.
Consequently,
$\sym{Wcompetitor}\exp(t\sym{SkewDirection})$ remains in $\sym{FeasibleTransport}$.
This gives explicit curves on which to test stationarity.

\emph{(b) The first variation is a trace pairing with a spectral gradient.}
Define the scalar phase cost
\[
 \sym{MatchingPhaseCost}(t)=(2\sin(t/2))^p,\qquad 0\le t\le\sym{PiConstant}.
\]
A planar rotation $\sym{MatchingPlanarRotation}$ through $t$ satisfies
\[
 (\sym{Identity}-\sym{MatchingPlanarRotation})^T(\sym{Identity}-\sym{MatchingPlanarRotation})
   =2\sym{Identity}-\sym{MatchingPlanarRotation}-\sym{MatchingPlanarRotation}^T
   =4\sin^2(t/2)\sym{Identity}.
\]
Its two displacement singular values are therefore both $2\sin(t/2)$,
and its contribution is $2\sym{MatchingPhaseCost}(t)$.
On a plane of rotation angle $\sym{MatchingPlanarAngle}$, use the planar rotation matrix through $\sym{PiConstant}/2$
$\sym{Jturn}=\bigl(\begin{smallmatrix}0&-1\\1&0\end{smallmatrix}\bigr)$
and define the skew-symmetric spectral gradient by
\[
 \sym{Wcompetitor}|_{\mathrm{plane}}
   =\cos\sym{MatchingPlanarAngle}\,\sym{IdentityTwo}+\sin\sym{MatchingPlanarAngle}\,\sym{Jturn},
 \qquad
 \sym{SpectralGradient}|_{\mathrm{plane}}
   =\sym{MatchingPhaseCost}'(\sym{MatchingPlanarAngle})\sym{Jturn}.
\]
On the $+1$ and $-1$ eigenspaces of the transporter, set the spectral
gradient $\sym{SpectralGradient}$ to zero.
Every real orthogonal matrix splits into these scalar blocks and planar
rotations: diagonalize over $\sym{ComplexField}$ and pair conjugate eigenvectors
to obtain the real planes. This defines $\sym{SpectralGradient}$ on the
whole space. On each plane the displayed blocks are combinations of
$\sym{IdentityTwo}$ and the same $\sym{Jturn}$, so they commute; on the scalar blocks
the gradient is zero. Thus $\sym{SpectralGradient}\sym{Wcompetitor}
=\sym{Wcompetitor}\sym{SpectralGradient}$ on the whole space.
These rotation planes may mix source and complementary directions;
stationarity will relate them to $\like{Psource}{\symglyph{Psource}_K}$.
Then for every skew-symmetric $\sym{SkewDirection}$,
\begin{equation}\label{eq:matching-first-variation}
 \left.\frac{d}{dt}
   \like{PowerCost}{\symglyph{PowerCost}_p}(\sym{Wcompetitor}\exp(t\sym{SkewDirection}))
 \right|_{t=0}
 =-\sym{TraceOperator}(\sym{SpectralGradient}\sym{SkewDirection}).
\end{equation}
\emph{Justification.} For $p>1$ the cost is differentiable, including at
zero displacement singular values, and its derivative is linear in
$\sym{SkewDirection}$. Change the sign on any whole rotation plane or
scalar line. This orthogonal sign matrix commutes with
$\sym{Wcompetitor}$, so conjugating $\sym{SkewDirection}$ by it conjugates
the whole curve and leaves its cost unchanged. Average over these sign
choices. Linearity of the derivative leaves its value unchanged, while
every off-diagonal block of $\sym{SkewDirection}$ cancels. This averaging
computes the derivative for arbitrary skew directions; feasibility is
imposed afterward, when we apply stationarity.

On a rotation plane of angle $\sym{MatchingPlanarAngle}$, the diagonal skew block is
$b\sym{Jturn}$. Since
$\exp(tb\sym{Jturn})$ is a rotation through $tb$, the phase becomes
$\sym{MatchingPlanarAngle}+tb$ and the cost derivative is
$2b\sym{MatchingPhaseCost}'(\sym{MatchingPlanarAngle})$. The trace pairing gives exactly
the same value:
\[
 -\sym{TraceOperator}\bigl(
   (\sym{MatchingPhaseCost}'(\sym{MatchingPlanarAngle})\sym{Jturn})(b\sym{Jturn})
  \bigr)
 =-b\sym{MatchingPhaseCost}'(\sym{MatchingPlanarAngle})\sym{TraceOperator}(-\sym{IdentityTwo})
 =2b\sym{MatchingPhaseCost}'(\sym{MatchingPlanarAngle}).
\]
On a scalar line the diagonal skew-symmetric block is zero.
Summing proves \cref{eq:matching-first-variation}, including directions
that mix the rotation planes.

\emph{(c) Stationarity forces the gradient to couple only source and
complementary directions.}
\emph{Justification.} Take the particular direction
$\sym{SkewDirection}=\like{Psource}{\symglyph{Psource}_K}\sym{SpectralGradient}\like{Psource}{\symglyph{Psource}_K}$.
It is skew-symmetric and commutes with $\like{Psource}{\symglyph{Psource}_K}$, so stationarity
and cyclicity of trace give
\[
 0=-\sym{TraceOperator}\bigl(
      (\like{Psource}{\symglyph{Psource}_K}\sym{SpectralGradient}\like{Psource}{\symglyph{Psource}_K})^2
     \bigr)
   =\|\like{Psource}{\symglyph{Psource}_K}\sym{SpectralGradient}\like{Psource}{\symglyph{Psource}_K}\|_F^2.
\]
The source diagonal block of the gradient is therefore zero.
Using $\sym{Identity}-\like{Psource}{\symglyph{Psource}_K}$ instead gives the same conclusion for the
complementary diagonal block. Relative to
$\sym{MatchingSourceSubspace}\oplus\sym{MatchingSourceSubspace}^{\perp}$, the gradient and source
reflection thus have the forms
\[
 \sym{SpectralGradient}
  =\begin{pmatrix}0&\sym{GradientOffDiagonal}\\-\like{GradientOffDiagonal}{\symglyph{GradientOffDiagonal}^T}&0\end{pmatrix},
 \qquad
 \sym{Jreflect}=2\like{Psource}{\symglyph{Psource}_K}-\sym{Identity}
  =\begin{pmatrix}\sym{Identity}&0\\0&-\sym{Identity}\end{pmatrix}.
\]
Multiplication gives
$\sym{Jreflect}\sym{SpectralGradient}\sym{Jreflect}
 =-\sym{SpectralGradient}$.
This identity is what makes the source projection compatible with the
decomposition constructed next.

\emph{Step 3: obtain blocks with one or two source directions.}
The rotation planes used to define the gradient need not preserve the
source. We now use the relation between the gradient and the source
reflection to construct blocks that preserve the source as well as the
minimizer. Their size will explain why the final optimization needs only
single labels and pairs.
\emph{(a) Each positive gradient level (an eigenspace of
$-\sym{SpectralGradient}^2$ with positive eigenvalue) admits a complex structure.}
\emph{Justification.}
The matrix $-\like{SpectralGradient}{\symglyph{SpectralGradient}^2}$ is symmetric and nonnegative,
since
$\langle-\like{SpectralGradient}{\symglyph{SpectralGradient}^2}x,x\rangle
  =\|\sym{SpectralGradient}x\|^2$.
Decompose the space into its eigenspaces. The gradient, transporter,
and source reflection all preserve them. On an eigenspace $\sym{GradientLevelSpace}$
with positive eigenvalue $\like{GradientLevel}{\symglyph{GradientLevel}^2}$, put
$\sym{GradientComplexStructure}=\sym{SpectralGradient}/\sym{GradientLevel}$.
It satisfies $\sym{GradientComplexStructure}^2=-\sym{Identity}$ and preserves lengths.
We can therefore treat $\sym{GradientLevelSpace}$ as a complex space by defining
\[
 (a+ib)x=ax+b\sym{GradientComplexStructure}x.
\]
A real plane spanned by $x,\sym{GradientComplexStructure}x$ becomes one
complex line. Commutation with the gradient makes the transporter
complex-linear and unitary in these coordinates.
The source reflection reverses the complex structure:
\[
 \sym{Jreflect}\sym{GradientComplexStructure}
   =-\sym{GradientComplexStructure}\sym{Jreflect}.
\]
Consequently
$\sym{Jreflect}((a+ib)x)=(a-ib)\sym{Jreflect}x$.
Together with $\sym{Jreflect}^2=\sym{Identity}$, this says that it acts as conjugation.

\emph{(b) The transporter has at most two phases on each positive level.}
\emph{Justification.} The definition of the gradient says that an eigenphase $t$ of the
transporter on this complex space satisfies
$\sym{MatchingPhaseCost}'(t)=\sym{GradientLevel}$, with $0<t<\sym{PiConstant}$.
There are at most two such phases. To check this directly,
\[
 \sym{MatchingPhaseCost}'(t)
   =p(2\sin(t/2))^{p-1}\cos(t/2),
 \qquad
 \sym{MatchingPhaseCost}''(t)
   =p(2\sin(t/2))^{p-2}
       \bigl((p-1)\cos^2(t/2)-\sin^2(t/2)\bigr).
\]
The second derivative changes sign once, where
$\tan^2(t/2)=p-1$. The first derivative therefore increases and then
decreases, so a positive horizontal level has at most two intersections.
Writing the two phases as $\like{GradientEigenphase}{\symglyph{GradientEigenphase}_1},\like{GradientEigenphase}{\symglyph{GradientEigenphase}_2}$, with equality allowed,
we obtain
\[
 \sym{Wcompetitor}|_{\sym{GradientLevelSpace}}
   =e^{i\like{GradientEigenphase}{\symglyph{GradientEigenphase}_1}}\sym{GradientPhaseProjection}+e^{i\like{GradientEigenphase}{\symglyph{GradientEigenphase}_2}}(\sym{Identity}-\sym{GradientPhaseProjection}),
\]
where $\sym{GradientPhaseProjection}$ is the orthogonal projection onto the first phase space.
A subspace preserved by $\sym{GradientPhaseProjection}$ is preserved by this transporter.
To preserve the source projection as well, we require invariance under
the conjugation $\sym{Jreflect}$.

\emph{(c) The phase projection and conjugation preserve common blocks of
real dimension at most four.}
\emph{Justification.} Choose complex coordinates in which $\sym{Jreflect}$ is coordinate
conjugation and write $\sym{GradientPhaseProjection}=\sym{ProjectionRealPart}+i\sym{ProjectionImaginaryPart}$ with real matrices $\sym{ProjectionRealPart},\sym{ProjectionImaginaryPart}$.
The equations $\sym{GradientPhaseProjection}^*=\sym{GradientPhaseProjection}$ and $\sym{GradientPhaseProjection}^2=\sym{GradientPhaseProjection}$ give
\[
 \sym{ProjectionRealPart}^T=\sym{ProjectionRealPart},\qquad \sym{ProjectionImaginaryPart}^T=-\sym{ProjectionImaginaryPart},\qquad
 \sym{ProjectionRealPart}^2-\sym{ProjectionImaginaryPart}^2=\sym{ProjectionRealPart},\qquad \sym{ProjectionRealPart}\sym{ProjectionImaginaryPart}+\sym{ProjectionImaginaryPart}\sym{ProjectionRealPart}=\sym{ProjectionImaginaryPart}.
\]
For real coordinate vectors $x$,
$\langle \sym{ProjectionRealPart}x,x\rangle=\operatorname{Re}\langle \sym{GradientPhaseProjection}x,x\rangle$
lies between zero and $\|x\|^2$, so the eigenvalues of $\sym{ProjectionRealPart}$ lie in $[0,1]$.
If $\sym{ProjectionRealPart}u=\sym{ProjectionEigenvalue} u$, the last two identities imply
\[
 \sym{ProjectionRealPart}(\sym{ProjectionImaginaryPart}u)=(1-\sym{ProjectionEigenvalue})\sym{ProjectionImaginaryPart}u,\qquad
 \sym{ProjectionImaginaryPart}^2u=-\sym{ProjectionEigenvalue}(1-\sym{ProjectionEigenvalue})u.
\]
For a unit $u$ with $0<\sym{ProjectionEigenvalue}<1$, set
$\sym{ProjectionCoupling}=\sqrt{\sym{ProjectionEigenvalue}(1-\sym{ProjectionEigenvalue})}$ and $v=\sym{ProjectionImaginaryPart}u/\sym{ProjectionCoupling}$.
Skew-symmetry gives
\[
 \langle u,\sym{ProjectionImaginaryPart}u\rangle=0,\qquad
 \|\sym{ProjectionImaginaryPart}u\|^2
   =-\langle u,\sym{ProjectionImaginaryPart}^2u\rangle
   =\sym{ProjectionEigenvalue}(1-\sym{ProjectionEigenvalue})
   =\sym{ProjectionCoupling}^2.
\]
Thus $u,v$ are orthonormal. The preceding identities now become
\[
 \sym{ProjectionRealPart}u=\sym{ProjectionEigenvalue} u,\quad \sym{ProjectionRealPart}v=(1-\sym{ProjectionEigenvalue})v,\quad
 \sym{ProjectionImaginaryPart}u=\sym{ProjectionCoupling} v,\quad \sym{ProjectionImaginaryPart}v=-\sym{ProjectionCoupling} u.
\]
Thus their real span is preserved by both $\sym{ProjectionRealPart}$ and $\sym{ProjectionImaginaryPart}$.
This includes $\sym{ProjectionEigenvalue}=1/2$. At $\sym{ProjectionEigenvalue}=0$ or $1$, $\sym{ProjectionImaginaryPart}u=0$ and a
single eigenvector gives an invariant real line.
By symmetry of $\sym{ProjectionRealPart}$ and skew-symmetry of $\sym{ProjectionImaginaryPart}$, the orthogonal complement
of each chosen line or plane is again preserved by $\sym{ProjectionRealPart},\sym{ProjectionImaginaryPart}$.
The construction can be repeated until the real
coordinate space is exhausted.

The complex spans of these lines and planes are invariant under
$\sym{GradientPhaseProjection},\sym{Jreflect}$, and $\sym{Wcompetitor}$. Returning to the original
real space gives blocks of dimension two or four. Each is also
preserved by
$\like{Psource}{\symglyph{Psource}_K}=(\sym{Jreflect}+\sym{Identity})/2$, hence by
$\sym{MatchingTargetProjection}
 =\sym{Wcompetitor}\like{Psource}{\symglyph{Psource}_K}\sym{Wcompetitor}^T$.
For example, the complex span of $u,v$ is the real span of
$u,v,\sym{GradientComplexStructure}u,\sym{GradientComplexStructure}v$.
The first two vectors are fixed by conjugation and hence belong to the
source; conjugation negates the last two, which belong to its complement.
This gives a four-dimensional block with two source directions.
An invariant real coordinate line similarly gives a two-dimensional block
with one source direction. Their principal-angle lists therefore have
two and one entries, respectively.

\emph{(d) The zero-gradient space contributes only fixed directions and
one-angle planes.}
\emph{Justification.} On $\ker\sym{SpectralGradient}$ the transporter has eigenvalues $+1,-1$,
because these are exactly the phases where the scalar derivative vanishes.
It is therefore a reflection there. Apply the two-projection decomposition
to $\like{Psource}{\symglyph{Psource}_K}$ and the projection onto its $-1$ eigenspace:
this gives real lines and planes preserved by both projections and the
transporter. On an invariant line the source and target are either both
zero or both the whole line. Replacing a one-dimensional $-1$ block by
$+1$ preserves them and decreases the cost by $2^p$. Such a block cannot
occur at a global minimum. The remaining blocks are fixed directions
or planes with one source direction.
This completes the decomposition needed for
\cref{eq:matching-block-additivity}.

\emph{Step 4: calculate the costs on a plane and on two planes.}
\emph{(a) The optimal map on a one-angle block is a rotation or a reflection.}
\emph{Justification.} For a principal angle $\like{theta}{\symglyph{theta}_i}$, the least rotation cost
is $2\sym{MatchingPhaseCost}(\like{theta}{\symglyph{theta}_i})$.
The alternative rotation sends a source unit vector to the opposite
target unit vector and turns through $\sym{PiConstant}-\like{theta}{\symglyph{theta}_i}$,
which is at least $\like{theta}{\symglyph{theta}_i}$.
A planar reflection has eigenvalues $+1,-1$, hence displacement singular values $0,2$,
so its cost is $2^p$. The one-angle minimum is
\[
 \like{SingletonPowerCost}{\symglyph{SingletonPowerCost}_i(p)}
   =\min\{2\sym{MatchingPhaseCost}(\like{theta}{\symglyph{theta}_i}),\,2^p\}.
\]

\emph{(b) A positive-determinant two-angle block reduces to a scalar minimum.}
\emph{Justification.} For two angles, relabel within the pair so that
$\like{theta}{\symglyph{theta}_i}\ge\like{theta}{\symglyph{theta}_j}$, and set
\[
 \sym{PairMeanAngle}=(\like{theta}{\symglyph{theta}_i}+\like{theta}{\symglyph{theta}_j})/2,\qquad
 \sym{PairHalfDifference}=(\like{theta}{\symglyph{theta}_i}-\like{theta}{\symglyph{theta}_j})/2.
\]
Then $0\le \sym{PairHalfDifference}\le \sym{PairMeanAngle}$ and $\sym{PairMeanAngle}+\sym{PairHalfDifference}\le\sym{PiConstant}/2$.
In principal coordinates the direct rotation is
\[
 \sym{PairDirectRotation}=\begin{pmatrix}\sym{PairCosineBlock}&-\sym{PairSineBlock}\\\sym{PairSineBlock}&\sym{PairCosineBlock}\end{pmatrix},\qquad
 \sym{PairCosineBlock}=\sym{DiagonalOperator}(\cos\like{theta}{\symglyph{theta}_i},
                      \cos\like{theta}{\symglyph{theta}_j}),\quad
 \sym{PairSineBlock}=\sym{DiagonalOperator}(\sin\like{theta}{\symglyph{theta}_i},
                      \sin\like{theta}{\symglyph{theta}_j}).
\]
Every orthogonal transporter on this four-dimensional space is
$\sym{PairDirectRotation}\sym{DiagonalOperator}(\like{PairFreeOrthogonal}{\symglyph{PairFreeOrthogonal}_0},\like{PairFreeOrthogonal}{\symglyph{PairFreeOrthogonal}_1})$ with $\like{PairFreeOrthogonal}{\symglyph{PairFreeOrthogonal}_0},\like{PairFreeOrthogonal}{\symglyph{PairFreeOrthogonal}_1}\in \sym{PlanarOrthogonalGroup}$:
multiplying by $\sym{PairDirectRotation}^T$ gives an orthogonal map preserving the source
and its complement. This parametrization allows mixing between the
principal planes.

For determinant $+1$, the free factors are either both rotations or both
reflections. Write their rotation angles as
$\sym{PairRotationR}+\sym{PairRotationS}$ and
$\sym{PairRotationR}-\sym{PairRotationS}$: the parameters
$\sym{PairRotationR},\sym{PairRotationS}$ are the half-sum and half-difference
of those angles. Define $\sym{PairPhaseX},\sym{PairPhaseY}\in[0,\sym{PiConstant}]$ by
\[
 \begin{array}{c|cc}
  &\cos\sym{PairPhaseX}&\cos\sym{PairPhaseY}\\ \hline
  \text{two rotations}
    &\cos\sym{PairMeanAngle}\cos\sym{PairRotationS}
    &\cos\sym{PairHalfDifference}\cos\sym{PairRotationR}\\
  \text{two reflections}
    &-\sin\sym{PairHalfDifference}\cos\sym{PairRotationR}
    &\sin\sym{PairMeanAngle}\cos\sym{PairRotationS}
 \end{array}
\]
Here a reflection factor is a planar rotation followed by
$\sym{PairPlanarReflection}=\sym{DiagonalOperator}(1,-1)$; for example,
$\like{PairFreeOrthogonal}{\symglyph{PairFreeOrthogonal}_0}
 =\sym{MatchingPlanarRotationSum}\sym{PairPlanarReflection}$.
The cosine bounds imply
\[
 \sym{PairMeanAngle}\le \sym{PairPhaseX}\le\sym{PiConstant}-\sym{PairMeanAngle},\qquad
 \sym{PairHalfDifference}\le \sym{PairPhaseY}\le\sym{PiConstant}-\sym{PairHalfDifference}.
\]
For reflections, the initial bounds are
$\sym{PiConstant}/2-\sym{PairHalfDifference}\le\sym{PairPhaseX}\le\sym{PiConstant}/2+\sym{PairHalfDifference}$
and $\sym{PiConstant}/2-\sym{PairMeanAngle}\le\sym{PairPhaseY}\le\sym{PiConstant}/2+\sym{PairMeanAngle}$;
these lie in the displayed ranges because
$\sym{PairMeanAngle}+\sym{PairHalfDifference}\le\sym{PiConstant}/2$.

To check that these parameters give the actual eigenphases, put
$\sym{PairPhaseCosineX}=\cos\sym{PairPhaseX}$ and
$\sym{PairPhaseCosineY}=\cos\sym{PairPhaseY}$. Block multiplication gives
\begin{align*}
 \sym{TraceOperator}\sym{Wcompetitor}
   &=\sym{TraceOperator}(\sym{PairCosineBlock}\like{PairFreeOrthogonal}{\symglyph{PairFreeOrthogonal}_0})
     +\sym{TraceOperator}(\sym{PairCosineBlock}\like{PairFreeOrthogonal}{\symglyph{PairFreeOrthogonal}_1})
     =4\sym{PairPhaseCosineX}\sym{PairPhaseCosineY},\\
 \sym{TraceOperator}(\sym{Wcompetitor}^2)
   &=\sym{TraceOperator}((\sym{PairCosineBlock}\like{PairFreeOrthogonal}{\symglyph{PairFreeOrthogonal}_0})^2)
     +\sym{TraceOperator}((\sym{PairCosineBlock}\like{PairFreeOrthogonal}{\symglyph{PairFreeOrthogonal}_1})^2)
     -2\sym{TraceOperator}(\sym{PairSineBlock}\like{PairFreeOrthogonal}{\symglyph{PairFreeOrthogonal}_1}
                        \sym{PairSineBlock}\like{PairFreeOrthogonal}{\symglyph{PairFreeOrthogonal}_0})\\
   &=16\sym{PairPhaseCosineX}^2\sym{PairPhaseCosineY}^2
     -8(\sym{PairPhaseCosineX}^2+\sym{PairPhaseCosineY}^2)+4.
\end{align*}
The coefficient of $z^2$ in the characteristic polynomial is
$((\sym{TraceOperator}\sym{Wcompetitor})^2
   -\sym{TraceOperator}(\sym{Wcompetitor}^2))/2$.
Orthogonality and determinant one make this polynomial reciprocal, so
\begin{align*}
 \det(zI-\sym{Wcompetitor})
   &=z^4-4\sym{PairPhaseCosineX}\sym{PairPhaseCosineY}z^3
     +\bigl(4(\sym{PairPhaseCosineX}^2+\sym{PairPhaseCosineY}^2)-2\bigr)z^2
     -4\sym{PairPhaseCosineX}\sym{PairPhaseCosineY}z+1\\
   &=\prod_{\sym{PhaseSignChoice}=\pm1}
       \bigl(z^2-2\cos(\sym{PairPhaseX}+\sym{PhaseSignChoice}\sym{PairPhaseY})z+1\bigr).
\end{align*}
Its roots are therefore
$e^{\pm i(\sym{PairPhaseX}+\sym{PairPhaseY})},
 e^{\pm i(\sym{PairPhaseX}-\sym{PairPhaseY})}$.
Folding their angles into $[0,\sym{PiConstant}]$ gives the displacement cost
\[
 2\bigl(\sym{MatchingPhaseCost}(\min\{\sym{PairPhaseX}+\sym{PairPhaseY},
                            2\sym{PiConstant}-\sym{PairPhaseX}-\sym{PairPhaseY}\})
       +\sym{MatchingPhaseCost}(|\sym{PairPhaseX}-\sym{PairPhaseY}|)\bigr).
\]

We can replace $(\sym{PairPhaseX},\sym{PairPhaseY})$ by $(\sym{PiConstant}-\sym{PairPhaseX},\sym{PiConstant}-\sym{PairPhaseY})$ when $\sym{PairPhaseX}+\sym{PairPhaseY}>\sym{PiConstant}$ without
changing this cost. Suppose therefore that $\sym{PairPhaseX}+\sym{PairPhaseY}\le\sym{PiConstant}$.
To replace the first parameter by $\sym{PairMeanAngle}$ without increasing the smaller
phase, we need $\sym{PairMeanAngle}-\sym{PairPhaseParameter}\le|\sym{PairPhaseX}-\sym{PairPhaseY}|$, as well as $\sym{PairPhaseParameter}\ge \sym{PairHalfDifference}$.
Choose the smallest such parameter,
$\sym{PairPhaseParameter}=\max\{\sym{PairHalfDifference},\sym{PairMeanAngle}-|\sym{PairPhaseX}-\sym{PairPhaseY}|\}$.
The identities $\sym{PairPhaseX}+\sym{PairPhaseY}+|\sym{PairPhaseX}-\sym{PairPhaseY}|=2\max\{\sym{PairPhaseX},\sym{PairPhaseY}\}\ge2\sym{PairMeanAngle}$ and $\sym{PairPhaseX}+\sym{PairPhaseY}\ge \sym{PairMeanAngle}+\sym{PairHalfDifference}$
also give $\sym{PairMeanAngle}+\sym{PairPhaseParameter}\le \sym{PairPhaseX}+\sym{PairPhaseY}$. Thus
\[
 \sym{PairHalfDifference}\le\sym{PairPhaseParameter}\le \sym{PairMeanAngle},\qquad
 \sym{PairMeanAngle}+\sym{PairPhaseParameter}\le \sym{PairPhaseX}+\sym{PairPhaseY},\qquad \sym{PairMeanAngle}-\sym{PairPhaseParameter}\le|\sym{PairPhaseX}-\sym{PairPhaseY}|.
\]
Since $\sym{MatchingPhaseCost}$ is increasing on $[0,\sym{PiConstant}]$, this
transporter costs at least
$2\bigl(\sym{MatchingPhaseCost}(\sym{PairMeanAngle}+\sym{PairPhaseParameter})+
        \sym{MatchingPhaseCost}(\sym{PairMeanAngle}-\sym{PairPhaseParameter})\bigr)$.
Every $\sym{PairPhaseParameter}\in[\sym{PairHalfDifference},\sym{PairMeanAngle}]$ is attained: take
$\like{PairFreeOrthogonal}{\symglyph{PairFreeOrthogonal}_0}=\like{PairFreeOrthogonal}{\symglyph{PairFreeOrthogonal}_1}=\sym{MatchingPlanarRotationGamma}$, where
$\sym{PairAttainingAngle}=\arccos(\cos\sym{PairPhaseParameter}/\cos \sym{PairHalfDifference})$.
The denominator is positive since $\sym{PairHalfDifference}\le\sym{PiConstant}/4$.
The preceding phase equations give $\sym{PairPhaseX}=\sym{PairMeanAngle},\sym{PairPhaseY}=\sym{PairPhaseParameter}$.
Thus the positive-determinant pair minimum is
\begin{equation}\label{eq:matching-pair-power-cost}
 \like{PairPowerCost}{\symglyph{PairPowerCost}_{ij}(p)}
 =\min_{\sym{PairPhaseParameter}\in[\sym{PairHalfDifference},\sym{PairMeanAngle}]}
     2\bigl(\sym{MatchingPhaseCost}(\sym{PairMeanAngle}+\sym{PairPhaseParameter})+
             \sym{MatchingPhaseCost}(\sym{PairMeanAngle}-\sym{PairPhaseParameter})\bigr).
\end{equation}

\emph{(c) The unrestricted two-angle minimum also compares two singletons.}
\emph{Justification.} For determinant $-1$, the eigenvalues are $+1,-1$ and a conjugate
rotational pair $e^{\pm it}$, with $0\le t\le\sym{PiConstant}$.
One free planar factor is a rotation and the other a reflection, so
the absolute trace is at most
$|\cos\like{theta}{\symglyph{theta}_i}+\cos\like{theta}{\symglyph{theta}_j}|
 +|\cos\like{theta}{\symglyph{theta}_i}-\cos\like{theta}{\symglyph{theta}_j}|
 =2\cos\like{theta}{\symglyph{theta}_j}$; thus
$|\sym{TraceOperator}\sym{Wcompetitor}|\le2\cos\like{theta}{\symglyph{theta}_j}$.
Since this trace is $2\cos t$, the rotational phase satisfies
$t\ge\like{theta}{\symglyph{theta}_j}$.
Its least cost is therefore
$2^p+2\sym{MatchingPhaseCost}(\like{theta}{\symglyph{theta}_j})$.
It is attained by reflecting on the larger-angle plane and rotating
on the smaller-angle plane.
Taking $\sym{PairPhaseParameter}=\sym{PairHalfDifference}$ bounds the positive-determinant minimum by the cost
of two rotations. Taking $\sym{PairPhaseParameter}=\sym{PairMeanAngle}$ bounds it by
$2\sym{MatchingPhaseCost}(\like{theta}{\symglyph{theta}_i}+\like{theta}{\symglyph{theta}_j})
 \le2\cdot2^p$, the cost of two reflections.
The less expensive mixed rotation/reflection choice is the
negative-determinant minimum just computed.
Comparing these four independent choices shows that the unrestricted
four-dimensional minimum is
\[
 \min\{\like{PairPowerCost}{\symglyph{PairPowerCost}_{ij}(p)},\,
         \like{SingletonPowerCost}{\symglyph{SingletonPowerCost}_i(p)}
          +\like{SingletonPowerCost}{\symglyph{SingletonPowerCost}_j(p)}\}.
\]
A four-dimensional block therefore offers either a paired choice or
two independent singleton choices.

\emph{Step 5: prove equality of the minimum costs and recover a transporter.}
We construct a transporter from every matching and extract a matching
from a minimizing transporter, preserving the cost in both constructions.
Thus solving the matching problem gives the minimum cost and an attaining
transporter, although individual transporters and matchings need not
correspond one-to-one.
For a matching $\sym{Matching}$ define
\[
 \sym{MatchingPowerCost}(\sym{Matching})
   =\sum_{\{i,j\}\in\sym{Matching}}\like{PairPowerCost}{\symglyph{PairPowerCost}_{ij}(p)}
     +\sum_{i\text{ unmatched}}\like{SingletonPowerCost}{\symglyph{SingletonPowerCost}_i(p)}.
\]
\emph{(a) Every matching gives a feasible transporter of the same cost.}
\emph{Justification.} The principal-angle decomposition gives mutually orthogonal planes
for the positive angle labels. Choose an attaining two-angle transporter for each selected pair,
an attaining planar rotation or reflection for each unmatched label,
and identity on the remaining common directions.
Their orthogonal sum is an element of $\sym{FeasibleTransport}$.
By \cref{eq:matching-block-additivity}, its cost is
$\sym{MatchingPowerCost}(\sym{Matching})$. Hence
$\min_{\sym{FeasibleTransport}}\like{PowerCost}{\symglyph{PowerCost}_p}
 \le\min_{\sym{Matching}}\sym{MatchingPowerCost}(\sym{Matching})$.

\emph{(b) A minimizing transporter gives a matching of the same cost.}
\emph{Justification.} Use the global minimizer from Step~1 and its
decomposition from Step~3. Its restriction to each block must minimize
the corresponding block problem: otherwise replacing that restriction
by a cheaper transporter would lower the global cost while retaining
the same source and target. By Step~4, a plane supplies a singleton,
and a four-dimensional block supplies a pair or two singletons.
Each occurrence of an angle belongs to one block, so no label can be used twice.
This produces a matching with the same total cost as the minimizer.

Zero angles do not supply cheaper extra partners. Their singleton cost
is zero; pairing an angle $\sym{theta}$ with zero forces $\sym{PairHalfDifference}=\sym{PairMeanAngle}=\sym{theta}/2$
in \cref{eq:matching-pair-power-cost}, giving $2\sym{MatchingPhaseCost}(\sym{theta})$,
which is at least its singleton cost.
Thus we can leave zero-angle directions fixed and use only the original
positive labels. Removing such pairs cannot strictly lower the minimizing
cost, by the construction in (a), so it preserves that cost. We have proved
\begin{equation}\label{eq:matching-power-equality}
 \min_{\sym{FeasibleTransport}}\like{PowerCost}{\symglyph{PowerCost}_p}
   =\min_{\sym{Matching}}\sym{MatchingPowerCost}(\sym{Matching}),
 \qquad 1<p<2.
\end{equation}
For the matching in \cref{fig:real-nuclear-matching}, this construction
uses one four-dimensional block for labels $1,2$ and one plane for label~$3$.

\emph{Step 6: pass to the trace norm and evaluate the weights.}
\emph{(a) The minimum values converge as $p$ approaches $1$ from above.}
\emph{Justification.} Every displacement singular value lies in $[0,2]$.
Since $\sym{DisplacementSize}^p$ converges uniformly to $\sym{DisplacementSize}$ on this interval as $p$ approaches $1$ from above,
\[
 \sup_{\sym{Wcompetitor}\in\sym{FeasibleTransport}}
   \left|\like{PowerCost}{\symglyph{PowerCost}_p}(\sym{Wcompetitor})
                  -\|\sym{Identity}-\sym{Wcompetitor}\|_1\right|
 \le n\sup_{0\le\sym{DisplacementSize}\le2}|\sym{DisplacementSize}^p-\sym{DisplacementSize}|\ \longrightarrow\ 0.
\]
The minimum values on the left of \cref{eq:matching-power-equality}
therefore converge to the minimum trace norm.
The scalar block minima also converge: the one-angle minimum is a
minimum of two continuous quantities, and the pair minimum is over
the fixed compact interval $[\sym{PairHalfDifference},\sym{PairMeanAngle}]$ with a uniformly convergent cost.
There are finitely many matchings, so their minimum values converge too.

\emph{(b) The limiting block weights are the stated trace-norm costs.}
\emph{Justification.} At $p=1$ the singleton formula gives
$\like{SingletonPowerCost}{\symglyph{SingletonPowerCost}_i(1)}
 =\min\{4\sin(\like{theta}{\symglyph{theta}_i}/2),2\}
 =\like{SingletonCost}{\symglyph{SingletonCost}_i}$.
For the paired cost, the scalar expression becomes
\[
 2\bigl(\sym{MatchingPhaseCost}(\sym{PairMeanAngle}+\sym{PairPhaseParameter})
          +\sym{MatchingPhaseCost}(\sym{PairMeanAngle}-\sym{PairPhaseParameter})\bigr)\big|_{p=1}
 =8\sin(\sym{PairMeanAngle}/2)\cos(\sym{PairPhaseParameter}/2).
\]
It is nonincreasing as $\sym{PairPhaseParameter}$ runs from $\sym{PairHalfDifference}$ to $\sym{PairMeanAngle}$, so its minimum is at $\sym{PairPhaseParameter}=\sym{PairMeanAngle}$,
with value $4\sin \sym{PairMeanAngle}=\like{PairCost}{\symglyph{PairCost}_{ij}}$.
At this endpoint one rotational phase is $2\sym{PairMeanAngle}$ and the other is zero:
the paired transporter fixes two real directions and rotates the other
two through $\like{theta}{\symglyph{theta}_i}+\like{theta}{\symglyph{theta}_j}$.
Taking the limit in \cref{eq:matching-power-equality} now gives
\cref{eq:real-nuclear-match}. Compactness supplies an attaining transporter
for the trace norm as well.
\end{proof}
% Formal proof locations: RealNuclearMatching.lean (endpoint continuation),
% SchattenRealDecomposition.lean / SchattenRealRestrictions.lean (block
% reduction), SchattenRealMatchingWeights.lean / RealMatchingAttainment.lean
% (matching and attainment), SchattenRealAngleCosts.lean (angle costs),
% RealFourCoordinates.lean / RealFourPositive.lean / RealFourOptimization.lean
% (all two-angle competitors and the scalar edge reduction).
% All paths are under DavisKahan/FiniteDimensional/DirectRotation/.
% See AUTHORING_NOTES.md for the mathematical and provenance evidence.

For $\like{theta}{\symglyph{theta}_1}=\like{theta}{\symglyph{theta}_2}=\sym{PiConstant}/6$, the direct rotation costs
$8\sin(\sym{PiConstant}/12)>2$, whereas the paired competitor costs
$4\sin(\sym{PiConstant}/6)=2$. The theorem includes comparison with every ambient
orthogonal transporter, so this paired cost is globally minimal in the
example.

Davis and Kahan's complex-phase example already shows how a phase change
can decrease full trace displacement \citep[Example~4.2]{davis1970rotation}.
Bolotnikov, Li, and Rodman formulate the ambient unitary minimal-distortion
problem and prove Q-norm and Frobenius optimality results
\citep[Problem~1.1 and Theorems~3.1, 4.2]{bolotnikov2001minimal}. Our formula gives the exact minimum trace-norm displacement over all real
orthogonal maps carrying one subspace onto the other, in terms of singleton
and paired angle costs. We did not locate this exact matching formula in
the sources compared here. The real field permits pairing between distinct
principal planes, which is why the answer is an optimization over matchings.


\section{Discussion}
\label{sec:discussion}

The formalization covers the original perturbation theory on real and
complex Hilbert spaces, including unbounded self-adjoint operators and
operator-ideal norms. The unbounded statements require explicit domain
conditions and bounded residuals. The operator-ideal statements require
finite norm values. These hypotheses determine which approximation problems
the estimates address.

The correction of the full-displacement claim
\citep[Proposition~4.4]{davis1970rotation} separates two extremal statements.
The valid comparison \citep[Proposition~4.3]{davis1970rotation}
gives squared-displacement minimality; its Q-norm consequence
is valid. Our real four-dimensional family exhibits full trace-norm failure
at every prescribed positive acute angle. In each example, the competitor has
two nonzero displacement singular values while the direct rotation has four.
The sums and squared sums of these values account for their different norm
comparisons. The earlier complex example supplied the construction; our
calculation makes its contradiction of the printed real claim explicit.

The matching theorem addresses the resulting optimization problem for any
equal-dimensional pair of finite-dimensional real subspaces. Its proof uses
principal-angle geometry, real orthogonal block decompositions, and
singular-value norms. These are also among the foundations required for the
Davis--Kahan formalization. The proof of the matching theorem uses additional
optimization arguments, including a passage to the trace-norm endpoint.

\section{Conclusion}
\label{sec:conclusion}

We formalized 28 of the 29 selected Davis--Kahan results and refuted the
remaining printed statement \citep[Proposition~4.4]{davis1970rotation}.
The sine-theta theorem includes
unbounded self-adjoint operators and symmetric operator-ideal norms with
their domain and finiteness assumptions. The direct-rotation calculation
identifies the failed singular-value comparison and gives an explicit real
counterexample and a valid Q-norm consequence.

We also formalized the operator theory used in these results: principal
angles and direct rotations, Borel spectral calculus, approximation numbers
and symmetric ideals, and Sylvester separation. Parts of this development
were subsequently adapted into Tau Ceti. Using these foundations, we proved
an exact weighted-matching formula for the real trace-norm transport
minimum.

\backmatter
\section*{Statements and Declarations}

\subsection*{Funding}
\ifDKAnonymous
Funding information is omitted for anonymous review.
\else
This work was supported by the Defense Advanced Research Projects Agency
under Contract No. HR001125CE017.
\fi

\subsection*{Competing interests}
The authors declare no competing interests.

\subsection*{Data and code availability}
\ifDKAnonymous
Artifact information is omitted for anonymous review.
\else
The Lean proofs, result-to-Lean mapping, provenance records, figure-generation
sources, and manuscript sources are available at
\url{https://github.com/AIQ-Kitware/aiq-dkps-formalization}.
\fi

\subsection*{AI assistance}
ChatGPT and Claude assisted with the formalization and manuscript preparation.
Formalizations were carried out by large language models with human oversight.
The matching result was derived using GPT-6.1 Sol and reviewed by the human
  authors.

\bibliography{references}

\clearpage
\appendix

\section{Symbols, conventions, and terminology}
\label{app:notation}

We write $\sym{A}$ for a reference operator and
$\sym{Acurrent}=\sym{A}+\sym{H}$ for the perturbed operator. The
sine-theta theorem uses coordinate maps for trial and exact subspaces.
The transport results name the subspaces $\sym{Usource},\sym{Vtarget}$
and their projections $\sym{Psource},\sym{Qtarget}$ directly.
The sine-theta projection is $\sym{Pexact}$, the residual is $\sym{R}$,
and direct rotations are denoted by $\sym{Rrotation}$ throughout.

In the colored version, blue identifies exact or target data, amber trial
or source data, cyan the exact complementary block, magenta angles, green
residuals, and violet spectral separation. The residual $\sym{R}$ and
direct rotation $\sym{Rrotation}$ have distinct roles and colors. The
notation and definitions also apply in monochrome.

\begingroup
\small
\setlength{\tabcolsep}{4pt}
\renewcommand{\arraystretch}{1.15}
\begin{longtable}{@{}>{\raggedright\arraybackslash}p{0.29\textwidth}>{\raggedright\arraybackslash}p{0.65\textwidth}@{}}
\toprule
\textbf{Symbol} & \textbf{Definition or interpretation}\\
\midrule
\endfirsthead
\multicolumn{2}{@{}l}{\textit{Symbols and conventions (continued)}}\\[2pt]
\toprule
\textbf{Symbol} & \textbf{Definition or interpretation}\\
\midrule
\endhead
\bottomrule
\endfoot
$\sym{A},\sym{H},\sym{Acurrent}$ & Reference operator, additive perturbation, and perturbed operator $\sym{Acurrent}=\sym{A}+\sym{H}$.\\
$\sym{AmbientHilbertSpace},\ \sym{TrialCoordinateSpace},\ \sym{ExactCoordinateSpace},\ \sym{ComplementaryCoordinateSpace}$ & Ambient, trial-coordinate, exact-coordinate, and complementary-coordinate Hilbert spaces.\\
$\sym{E0}:\sym{TrialCoordinateSpace}\to\sym{AmbientHilbertSpace}$ & \dkSymbolDescription{E0}\\
$\sym{TrialOperator}$ & Self-adjoint trial operator. A Rayleigh--Ritz choice is $\sym{TrialOperator}=\like{E0}{\symglyph{E0}^*}\sym{Acurrent}\sym{E0}$.\\
$\sym{ReferenceBlock0},\sym{ReferenceBlock1}$ & Reference blocks on $\DKRange(\sym{Psource})$ and $\DKKernel(\sym{Psource})$ in Section~\ref{sec:foundations}.\\
$\sym{F0}:\sym{ExactCoordinateSpace}\to\sym{AmbientHilbertSpace}$ & \dkSymbolDescription{F0}\\
$\sym{F1}:\sym{ComplementaryCoordinateSpace}\to\sym{AmbientHilbertSpace}$ & \dkSymbolDescription{F1}\\
$\sym{Lambda1}$ & Complementary self-adjoint block, with $\sym{Acurrent}\sym{F1}=\sym{F1}\sym{Lambda1}$ on its domain.\\
$\sym{SpectralMeasure},\sym{SpectralSet}$ & Projection-valued measure and Borel spectral set; $\sym{SpectralMeasure}(\sym{SpectralSet})$ is the corresponding spectral projection. These symbols differ from the ambient space $\sym{AmbientHilbertSpace}$ and complementary block $\sym{S}$.\\
$\sym{Pexact}=\sym{F0}\like{F0}{\symglyph{F0}^*}$ & Projection onto the desired exact subspace in the sine-theta theorem.\\
$\sym{Theta0},\sym{sinTheta0}$ & Trial-coordinate angle operator and its positive sine, with $\sym{sinTheta0}=(\sym{Identity}-\like{E0}{\symglyph{E0}^*}\sym{Pexact}\sym{E0})^{1/2}$.\\
$\sym{Theta}=\sym{DiagonalOperator}(\sym{Theta0},\sym{Theta1})$ & Full angle operator; $\sym{Theta1}$ compares the orthogonal complements. Its operator norm is the largest spectral angle.\\
$\sym{CrossedSourceDefect},\sym{CrossedTargetDefect}$ & Source directions perpendicular to the target, and target directions perpendicular to the source. Section~\ref{sec:foundations} uses the condition $\sym{DimensionOperator}\sym{CrossedSourceDefect}=\sym{DimensionOperator}\sym{CrossedTargetDefect}$.\\
$\{\sym{ZeroVector}\}$ & Zero subspace: the set containing only the zero vector.\\
$\sym{S}=(\sym{Identity}-\sym{Pexact})\sym{E0}$ & Complementary block, with modulus $\dkModulus{\sym{S}}=\sym{sinTheta0}$.\\
$\sym{X}=\like{F1}{\symglyph{F1}^*}\sym{E0}$ & Coordinates of the trial component in the unwanted space; $\sym{S}=\sym{F1}\sym{X}$.\\
$\sym{SylvesterLeft},\sym{SylvesterRight},\sym{SylvesterRHS}$ & Generic left and right Sylvester coefficients and right side: $\sym{SylvesterLeft}\sym{X}-\sym{X}\sym{SylvesterRight}=\sym{SylvesterRHS}$. The historical equation writes the right coefficient as $\sym{HistoricalSylvesterRight}$.\\
$\sym{R}=\sym{Acurrent}\sym{E0}-\sym{E0}\sym{TrialOperator}$ & Bounded residual agreeing with this difference on $\DKDom(\sym{TrialOperator})$.\\
$\sym{delta}>0$ & \dkSymbolDescription{delta}\\
$[\sym{betaEndpoint},\sym{alphaEndpoint}]$ & Interval in the interval/exterior gap hypothesis.\\
$\DKRange(T),\DKKernel(T),\DKDom(T)$ & Range, kernel, and domain of an operator.\\
$\sym{Usource},\sym{Vtarget}$ & Source and target subspaces in the displacement problem.\\
$\sym{Psource},\sym{Qtarget}$ & Orthogonal projections onto $\sym{Usource},\sym{Vtarget}$ in that problem.\\
$\sym{Rrotation},\sym{Wcompetitor}$ & Direct rotation and an arbitrary orthogonal or unitary transporter carrying $\sym{Usource}$ onto $\sym{Vtarget}$.\\
$\sym{PolarU}$ & Partial isometry in a polar decomposition, as in $\sym{S}=\sym{PolarU}|\sym{S}|$.\\
$\sym{IntertwinerT}$ & The intertwiner $\sym{Qtarget}\sym{Psource}+(\sym{Identity}-\sym{Qtarget})(\sym{Identity}-\sym{Psource})$, whose unitary polar factor is the acute direct rotation.\\
$\sym{theta},\like{theta}{\symglyph{theta}_j}$ & Individual principal angles, in $[0,\sym{PiConstant}/2]$.\\
$\sym{Xframe},\sym{Yframe}$ & Orthonormal coordinate matrices for the source and target in the real counterexample.\\
$\sym{Jturn},\sym{Cturn},\sym{Hmix}$ & Planar rotation through $\sym{PiConstant}/2$, rotation through $2\sym{theta}$, and orthogonal mixing matrix used in that counterexample.\\
$\sym{Jreflect};\ \sym{AngleGenerator}$ & Source reflection $2\sym{Psource}-\sym{Identity}$, with square $\sym{Identity}$; skew angle generator of a direct rotation. The $\sym{PiConstant}/2$-rotation matrix $\sym{Jturn}$ has square $-\sym{IdentityTwo}$.\\
$\sym{CounterPrincipalOne},\sym{CounterPrincipalTwo};\ \sym{CounterMovingPlane},\sym{CounterFixedPlane}$ & Principal planes of the direct rotation and moving/fixed planes of the competitor in \cref{fig:prop44-displacement-planes}.\\
$\sym{CounterChordA},\sym{CounterChordB}$ & Displacement singular values $\sqrt{2-\sqrt2}$ and $\sqrt2$ in that counterexample.\\
$\sym{Kdisplacement}=\sym{Identity}-\sym{Wcompetitor}$ & Full displacement operator in the failed Ky Fan comparison.\\
$\sym{OmegaOne},\sym{OmegaTwo}$ & Principal-plane projections used to test that comparison.\\
$\dkModulus{T}=(T^*T)^{1/2}$ & Positive modulus of an operator.\\
$\sym{SingularValue}(T)$ & Ordered singular values; approximation numbers provide the corresponding comparisons for bounded operators in Hilbert space.\\
$\sym{NormN}(T)$ & \dkSymbolDescription{NormN}\\
$\dkOperatorNorm{T}$ & Operator norm; in finite dimension, the largest singular value.\\
$\dkFrobeniusNorm{T}$ & Frobenius norm, $\bigl(\sum_j \sym{SingularValue}(T)^2\bigr)^{1/2}$.\\
$\dkTraceNorm{T}$ & Trace norm, $\sum_j \sym{SingularValue}(T)$.\\
$\dkKyFanNorm{T}{j}$ & Ky Fan norm, the sum of the $j$ largest singular values.\\
$\sym{NormQ}\text{-norm}$ & A norm satisfying $\sym{NormN}(T)^2=\sym{NormM}(T^*T)$ for a unitarily invariant norm $\sym{NormM}$.\\
$\sym{Matching},\like{SingletonCost}{\symglyph{SingletonCost}_i},\like{PairCost}{\symglyph{PairCost}_{ij}}$ & Matching of positive principal-angle labels, with singleton and paired costs in \cref{eq:nuclear-matching-costs}.\\
$\sym{MatchingEdgeWeight};\ \like{TraceCost}{\symglyph{TraceCost}(M)}$ & Edge weight $\like{SingletonCost}{\symglyph{SingletonCost}_i}+\like{SingletonCost}{\symglyph{SingletonCost}_j}-\like{PairCost}{\symglyph{PairCost}_{ij}}$; total trace-norm displacement cost of the matching.\\
\end{longtable}
\endgroup

\subsection{Terminology}
\label{app:terminology}

The following terms describe the roles of the mathematical objects used in
\cref{sec:sine-theta,sec:example-transport}.

\begin{description}
\item[Ambient space and operator.] The whole Hilbert space
$\sym{AmbientHilbertSpace}$ in which the subspaces lie, and the operator acting
on that space. In the sine-theta statement the ambient operator is
$\sym{Acurrent}$.

\item[Reference and perturbed operators.] The reference operator is
$\sym{A}$; an additive perturbation $\sym{H}$ gives the perturbed operator
$\sym{Acurrent}=\sym{A}+\sym{H}$. The choice of which operator supplies the
desired spectral subspace depends on the approximation problem.

\item[Exact subspace.] The desired subspace specified by the spectral
decomposition of the operator of interest. In the sine-theta statement,
it is $\DKRange(\sym{F0})$ for $\sym{Acurrent}$, with orthogonal complement
$\DKRange(\sym{F1})$. In the perturbation application in the introduction,
the desired subspace instead belongs to the reference operator $\sym{A}$.
An exact formula, as in the transport theorem, gives an equality for the
quantity being determined.

\item[Trial subspace and trial operator.] The trial subspace is a candidate
approximation to the desired subspace, represented by $\DKRange(\sym{E0})$.
The self-adjoint trial operator $\sym{TrialOperator}$ represents an
approximation to the action of the ambient operator in trial coordinates.

\item[Residual.] The difference
$\sym{Acurrent}\sym{E0}-\sym{E0}\sym{TrialOperator}$ between the ambient
action and the embedded trial action. For unbounded operators this identity
holds on the appropriate domain, with a bounded residual $\sym{R}$.

\item[Spectral separation.] A lower bound $\sym{delta}>0$ on the distance
between the relevant spectra. In the sine-theta theorem these are the
spectra of the trial operator and the complementary block $\sym{Lambda1}$;
the interval and ordered separation hypotheses are given in
\cref{sec:gap-geometry}.

\item[Transporter and competitor.] A transporter is an ambient orthogonal
or unitary map carrying the source subspace onto the target subspace.
A competitor is another map satisfying the same constraint.

\item[Full displacement.] The operator $\sym{Identity}-\sym{Wcompetitor}$
on the whole ambient space for a transporter $\sym{Wcompetitor}$. Its
restriction to source directions is represented by
$(\sym{Identity}-\sym{Wcompetitor})\sym{Psource}$.
\end{description}

\section{Reading the Lean signatures}
\label{app:lean-guide}
\label{app:reading-lean}

A Lean theorem declaration has the form
\[
  \texttt{theorem name (parameters) (hypotheses) : conclusion := proof}.
\]
Parameters introduce the mathematical objects under discussion.  A hypothesis
such as \texttt{(hA : IsSelfAdjoint A)} supplies the corresponding fact;
bundled predicates such as \texttt{IsTrialResidual} and
\texttt{IsExactSpectralDecomposition} collect related assumptions.  Braces mark
arguments normally inferred by Lean, while square brackets request typeclass
instances supplying ambient algebraic, analytic, or topological structure.  The proposition after the final colon is the theorem's conclusion
\citep{demoura2021lean4}.

In \Cref{lst:dk-sintheta-headline}, \texttt{RCLike} is the Mathlib abstraction
covering $\sym{RealField}$ and $\sym{ComplexField}$.  A type
\texttt{ContinuousLinearMap K F E} is a continuous linear map.  A type
\texttt{LinearPMap K E E} is a linear partial map whose
domain is a submodule; this is the representation used for possibly unbounded
operators.  \texttt{N.Mem T} states that the ideal gauge associated with $\sym{NormN}$ is
finite on $T$, and \texttt{N.gaugeReal T} is the corresponding real value.  The
suffix \texttt{Real} refers to the codomain of the gauge and not to the scalar
field.

\subsection{Paper notation and Lean binder names}
\label{app:paper-lean-notation}
The following correspondence is for the particular sine-theta theorem
in \Cref{lst:dk-sintheta-headline}, not a global naming rule for the Lean
library. The Lean source is
\texttt{DavisKahan/Sources/DavisKahan1970/SineTheta/Presentation.lean},
the declaration \texttt{sinTheta\_unbounded\_formGap\_whereDefinedUIN\_rclike}.
Lean's exact source uses Unicode subscripts and Greek letters (for example,
its trial-operator binder is typeset here as $\mathtt{A_0}$). The listing
in the main text uses ASCII transliterations and readable type names.
In particular, the Lean binder $\mathtt{A_0}$ denotes the general
$\sym{TrialOperator}$, while $\sym{ReferenceBlock0}$ is the separate
reference spectral block and does not occur as a parameter of this theorem.
The variable $\mathtt{H}$ is a coordinate-space carrier, not the
perturbation $\sym{H}$. Lean's scalar-field type variable is rendered as \texttt{K} in the
schematic listing (the source uses a blackboard-bold k).

\begingroup
\small
\setlength{\tabcolsep}{4pt}
\renewcommand{\arraystretch}{1.11}
\begin{longtable}{@{}>{\centering\arraybackslash}p{0.22\linewidth}>{\centering\arraybackslash}p{0.16\linewidth}>{\raggedright\arraybackslash}p{0.55\linewidth}@{}}
\caption{Paper symbols and the corresponding Lean binders in the sine-theta declaration.}
\label{tab:paper-lean-notation}\\
\toprule
This paper & Lean binder & Mathematical role \\
\midrule
\endfirsthead
\toprule
This paper & Lean binder & Mathematical role \\
\midrule
\endhead
\bottomrule
\endfoot
\dkLeanSineThetaRows
\end{longtable}
\endgroup

\section{Formalized Davis--Kahan results}
\label{app:dk-inventory}

Table~\ref{tab:dk-full-inventory} lists the \DKResultCount{} source targets:
the four headline theorems \citep[Section~2]{davis1970rotation} and every named theorem,
proposition, lemma, and corollary established in the paper. The table keeps
compact IDs for reference. The artifact maps each ID to the corresponding
Lean declarations.
\ifDKAnonymous\else
The \href{https://github.com/AIQ-Kitware/aiq-dkps-formalization/blob/main/dev/davis-kahan-1970-formalization-result-inventory.json}{result-to-Lean mapping}
contains that correspondence.
\fi

\begin{longtable}{>{\raggedright\arraybackslash}p{0.24\linewidth}>{\raggedright\arraybackslash}p{0.68\linewidth}}
\caption{The \DKResultCount{} Davis--Kahan results selected for formalization.}
\label{tab:dk-full-inventory}\\
\toprule
ID & Result \\
\midrule
\endfirsthead
\toprule
ID & Result \\
\midrule
\endhead
\path{S2-sin-theta} & Single-angle sine theorem \\
\path{S2-tan-theta} & Single-angle tangent theorem \\
\path{S2-sin-two-theta} & Double-angle sine theorem \\
\path{S2-tan-two-theta} & Double-angle tangent theorem \\
\path{DK-3.1-prop} & Acute direct rotation existence and uniqueness \\
\path{DK-3.2-prop} & Nonacute existence criterion \\
\path{DK-3.3-prop} & Principal square-root characterization \\
\path{DK-3.4-prop} & Square as a direct rotation \\
\path{DK-3.1-thm} & Classification of pairs of subspaces \\
\path{DK-3.1-cor} & Compact classification by angle eigenvalues \\
\path{DK-3.5-prop} & Angle commutation and eigenspace geometry \\
\path{DK-3.2-cor} & Reversal symmetry \\
\path{DK-4.1-prop} & Pointwise and singular-value extremality of the direct rotation \\
\path{DK-4.1-cor} & Unitarily invariant norm minimality of restricted displacement \\
\path{DK-4.2-prop} & Basis-angle square-sum extremality \\
\path{DK-4.3-prop} & Squared-displacement unitarily invariant norm minimality \\
\path{DK-4.4-prop} & Full-displacement optimality claim \\
\path{DK-5.1-thm} & Banach-space Sylvester lower bound \\
\path{DK-5.2-thm} & Semibounded self-adjoint Sylvester theorem \\
\path{DK-5.1-lem} & Strong-cutoff convergence of singular values \\
\path{DK-6.1-lem} & Direct-sum norm comparison and converse \\
\path{DK-6.2-lem} & Reflection-pinch contraction \\
\path{DK-6.1-prop} & Sine proof, ambient limitation, and symmetric sine theorem \\
\path{DK-6.1-thm} & Generalized sine theorem \\
\path{DK-6.2-thm} & Pairwise-gap square-norm sine theorem \\
\path{DK-6.3-thm} & Generalized tangent theorem and supporting machinery \\
\path{DK-6.3-lem} & Finite-rank near-maximizer off-diagonal block estimate \\
\path{DK-8.1-thm} & Branch selection and spectral repulsion \\
\path{DK-8.2-thm} & Smallness selects the acute branch \\
\bottomrule
\end{longtable}

A single source result may use several Lean declarations. The IDs identify
source statements independently of the organization of the formal library.
The printed claim \citep[Proposition~4.4]{davis1970rotation}
is refuted by the counterexample in
Section~\ref{sec:prop44}; its Q-norm consequence is stated separately there.

% Detailed table following the compact summary in app:dk-inventory.
% Summaries checked against the registered Lean statements and their supporting
% declarations; see TABLE4_LEAN_AUDIT.md for the row-by-row comparison.
\subsection{Hypotheses and conclusions}
\label{app:dk-result-summaries}

Table~\ref{tab:dk-result-summaries} summarizes the hypotheses and main
conclusions of the \DKResultCount{} results listed in
Table~\ref{tab:dk-full-inventory}, using the same IDs
\citep[Sections~2--8]{davis1970rotation}.
Hilbert spaces are separable and real or complex unless specified otherwise.
$\sym{NormN}$ denotes an arbitrary unitarily invariant norm on its domain of
finiteness; each displayed norm is evaluated on an operator in that domain.
$\dkOperatorNorm{\cdot}$ and $\dkHilbertSchmidtNorm{\cdot}$ denote the operator
and Hilbert--Schmidt norms. Self-adjoint operators may be unbounded:
operator identities hold on their domains, and residuals or right-hand sides
in norm estimates have bounded extensions.
For the ambient tangent bound, the tangent is defined as a bounded operator,
equivalently $\dkOperatorNorm{\sym{Psource}-\sym{Qtarget}}<1$.

For the geometric results, $\sym{Psource},\sym{Qtarget}$ are orthogonal
projections, $\sym{SourceDirectRotation}$ is a direct rotation, and
$\sym{Wcompetitor}$ is any orthogonal or unitary transporter satisfying
$\sym{Wcompetitor}\sym{Psource}\sym{Wcompetitor}^*=\sym{Qtarget}$.
Write
\[
 \sym{CrossedSourceDefect}=\DKRange(\sym{Psource})\cap\DKKernel(\sym{Qtarget}),\qquad
 \sym{CrossedTargetDefect}=\DKKernel(\sym{Psource})\cap\DKRange(\sym{Qtarget}).
\]
The pair is acute when $\sym{CrossedSourceDefect}=\sym{CrossedTargetDefect}=\{\sym{ZeroVector}\}$.
For geometric and displacement rows from \path{DK-3.3-prop} onward,
the ranges and kernels of
$\sym{Psource},\sym{Qtarget}$ have matching dimensions and
$\sym{DimensionOperator} \sym{CrossedSourceDefect}=\sym{DimensionOperator} \sym{CrossedTargetDefect}$; the displacement results also assume
the directed sine block
$(\sym{Identity}-\sym{Qtarget})|_{\DKRange(\sym{Psource})}$ compact.
$\sym{Theta0},\sym{Theta1}$ are the source and complementary angle operators,
$\sym{Theta}$ their ambient direct sum, and $\like{theta}{\symglyph{theta}_j}$ the ordered principal
angles. In the source decomposition the direct rotation has blocks
$\bigl(\begin{smallmatrix}C_0&-S_0^*\\S_0&C_1\end{smallmatrix}\bigr)$
with $C_j\ge0$; $\sym{AngleGenerator}$ is its skew angle generator, so
$\sym{SourceDirectRotation}=e^{\sym{AngleGenerator}\sym{Theta}}$.
The symbol $\sym{SourceDirectRotation}$ includes nonacute direct rotations;
for an acute pair it is the unique rotation denoted by the same letter
in the main text.
For real unitaries, spectral half-plane conditions refer to their
complexifications.

For perturbation results, $\sym{TrialOperator}$ is a self-adjoint trial operator.
In a reducing decomposition of the reference operator $\sym{A}$, the
source and complementary reference blocks are $\sym{ReferenceBlock0}$
and $\sym{ReferenceBlock1}$. The trial operator may be chosen to equal the
source reference block when the trial coordinates identify that subspace.
The operators $\sym{Lambda0},\sym{Lambda1}$ are the exact blocks of
$\sym{Acurrent}=\sym{A}+\sym{H}$.
The maps $\sym{E0},\sym{F0}$ give source and target coordinates,
$\sym{H0},\sym{H1}$ are the diagonal perturbation blocks, and
$\sym{R}=\sym{Acurrent}\sym{E0}-\sym{E0}\sym{TrialOperator}$.
Coordinate maps are isometries except in the two lower-frame rows.
The exact frames have orthogonal ranges whose sum is the ambient space;
reducing projections preserve operator domains.
We abbreviate the two spectral hypotheses, for $\sym{delta}>0$, by
\begin{align*}
 \sym{IntervalSeparation}_{\sym{delta}}(B,C):\quad&
 \sym{SpectrumSymbol}(B)\subset[\sym{betaEndpoint},\sym{alphaEndpoint}],\quad
 \sym{SpectrumSymbol}(C)\cap(\sym{betaEndpoint}-\sym{delta},\sym{alphaEndpoint}+\sym{delta})=\varnothing,
 \quad\text{or the roles interchanged};\\
 \sym{OrderedSeparation}_{\sym{delta}}(B,C):\quad&
 B\le\sym{alphaEndpoint} \sym{Identity},\quad C\ge(\sym{alphaEndpoint}+\sym{delta})\sym{Identity}.
\end{align*}
The endpoints are real, satisfy $\sym{betaEndpoint}\le\sym{alphaEndpoint}$, and are chosen for each row.
The ordered condition $\sym{OrderedSeparation}_{\sym{delta}}(B,C)$ places the spectra in
the half-infinite intervals $(-\infty,\sym{alphaEndpoint}]$ and
$[\sym{alphaEndpoint}+\sym{delta},\infty)$; it allows both operators to be unbounded.
To include both orientations, set
\[
 \sym{SineGap}(B,C)\quad\Longleftrightarrow\quad
 \sym{IntervalSeparation}_{\sym{delta}}(B,C)\quad\text{or}\quad
 \sym{OrderedSeparation}_{\sym{delta}}(B,C)\quad\text{or}\quad
 \sym{OrderedSeparation}_{\sym{delta}}(C,B).
\]
This is the separation condition of Section~\ref{sec:gap-geometry}.

\begingroup
\small
\setlength{\tabcolsep}{4pt}
\renewcommand{\arraystretch}{1.15}
\begin{longtable}{@{}>{\raggedright\arraybackslash}p{0.20\linewidth}>{\raggedright\arraybackslash}p{0.34\linewidth}>{\raggedright\arraybackslash}p{\dimexpr0.46\linewidth-4\tabcolsep\relax}@{}}
\caption{Hypotheses and main conclusions of the \DKResultCount{} Davis--Kahan results.}
\label{tab:dk-result-summaries}\\
\toprule
\textbf{ID} & \textbf{Hypotheses} & \textbf{Conclusion}\\
\midrule
\endfirsthead
\multicolumn{3}{@{}l}{\textit{Hypotheses and conclusions (continued)}}\\[2pt]
\toprule
\textbf{ID} & \textbf{Hypotheses} & \textbf{Conclusion}\\
\midrule
\endhead
\midrule
\multicolumn{3}{r@{}}{\textit{Continued on the next page}}\\
\endfoot
\bottomrule
\endlastfoot
\multicolumn{3}{@{}l}{\textbf{Headline perturbation bounds}}\\*\addlinespace[3pt]
\path{S2-sin-theta} &
 $\sym{SineGap}(\sym{TrialOperator},\sym{Lambda1})$. &
 $\sym{delta}\sym{NormN}(\sin\sym{Theta0})\le\sym{NormN}(\sym{R})$.\\\addlinespace
\path{S2-tan-theta} &
 $\sym{OrderedSeparation}_{\sym{delta}}(\sym{TrialOperator},\sym{Lambda1})$;
 Rayleigh--Ritz condition $\sym{H0}=0$. &
 $\sym{delta}\sym{NormN}(\tan\sym{Theta0})\le\sym{NormN}(\sym{R})$;
 $\sym{delta}\sym{NormN}(\tan\sym{Theta})\le\sym{NormN}(\sym{H})$.\\\addlinespace
\path{S2-sin-two-theta} &
 $\sym{A},\sym{Acurrent}$ have a common domain;
 $\sym{SpectrumSymbol}(\sym{Lambda0})\subset[\sym{betaEndpoint},\sym{alphaEndpoint}]$;
 $\sym{SpectrumSymbol}(\sym{Lambda1})$ outside
 $(\sym{betaEndpoint}-\sym{delta},\sym{alphaEndpoint}+\sym{delta})$. &
 $\sym{delta}\sym{NormN}(\sin(2\sym{Theta0}))\le2\sym{NormN}(\sym{R})$;
 $\sym{delta}\sym{NormN}(\sin(2\sym{Theta}))\le2\sym{NormN}(\sym{H})$.\\\addlinespace
\path{S2-tan-two-theta} &
 $\sym{OrderedSeparation}_{\sym{delta}}(\sym{TrialOperator},\sym{A1})$;
 bounded $\sym{H}$, $\sym{H0}=\sym{H1}=0$. &
 $\sym{delta}\sym{NormN}(\tan(2\sym{Theta0}))\le2\sym{NormN}(\sym{R})$;
 $\sym{delta}\sym{NormN}(\tan(2\sym{Theta}))\le2\sym{NormN}(\sym{H})$.\\\addlinespace[5pt]
\multicolumn{3}{@{}l}{\textbf{Two subspaces and direct rotations}}\\*\addlinespace[3pt]
\path{DK-3.1-prop} &
 $\sym{CrossedSourceDefect}=\sym{CrossedTargetDefect}=\{\sym{ZeroVector}\}$. &
 A unique direct rotation exists. Positive diagonal blocks characterize it
 among transporters.\\\addlinespace
\path{DK-3.2-prop} &
 General pair of projections. &
 A direct rotation exists iff $\sym{DimensionOperator} \sym{CrossedSourceDefect}=\sym{DimensionOperator} \sym{CrossedTargetDefect}$.
 On $\sym{CrossedSourceDefect}\oplus \sym{CrossedTargetDefect}$, $\sym{SourceDirectRotation}^2=-\sym{Identity}$;
 uniqueness can fail.\\\addlinespace
\path{DK-3.3-prop} &
 A direct rotation $\sym{SourceDirectRotation}$; conversely a unitary
 square root with spectrum in the closed right half-plane that maps
 $\sym{CrossedSourceDefect}$ onto $\sym{CrossedTargetDefect}$. &
 $\sym{SourceDirectRotation}^2=(2\sym{Qtarget}-\sym{Identity})(2\sym{Psource}-\sym{Identity})$,
 and its spectrum lies in the closed right half-plane.
 Conversely, such a square root is a direct rotation.\\\addlinespace
\path{DK-3.4-prop} &
 $C_0^2\ge\tfrac12 \sym{Identity}$;
 $Q_-=(2\sym{Psource}-\sym{Identity})\sym{Qtarget}(2\sym{Psource}-\sym{Identity})$. &
 $\sym{SourceDirectRotation}^2$ is a direct rotation from
 $\DKRange(Q_-)$ to $\DKRange(\sym{Qtarget})$.\\\addlinespace
\path{DK-3.1-thm} &
 Matching dimensions as above; angle operators
 $0\le\sym{Theta0},\sym{Theta1}\le\sym{PiConstant}/2$. &
 The four intersection dimensions and the generic angle's spectral
 multiplicity completely classify pairs up to isometric equivalence.
 Any angle data with matching nonzero multiplicities
 and the prescribed total dimension are realizable.\\\addlinespace
\path{DK-3.1-cor} &
 Classification hypotheses;
 $\sym{Psource}(\sym{Identity}-\sym{Qtarget})\sym{Psource}$ compact. &
 Classification reduces to angle eigenvalues with multiplicity:
 $\like{theta}{\symglyph{theta}_j}\downarrow0$. The two angle operators may have different
 zero multiplicities.\\\addlinespace
\path{DK-3.5-prop} &
 A direct rotation and its angle operator $\sym{Theta}$; the maximal
 eigenspace characterization assumes the acute case. &
 $\sym{Theta}$ commutes with $\sym{Psource},\sym{Qtarget},\sym{AngleGenerator},
 \sym{SourceDirectRotation}$. If $x\ne0$ and
 $\sym{Theta}x=\sym{theta} x$, then
 $\angle(x,\sym{SourceDirectRotation}x)=\sym{theta}$.
 In the acute case these eigenspaces are the unique maximal reducing
 subspaces of the prescribed angle.\\\addlinespace
\path{DK-3.2-cor} &
 Interchange $\sym{Psource}$ and $\sym{Qtarget}$. &
 $\sym{Theta}$ is unchanged, $\sym{AngleGenerator}$ changes to $-\sym{AngleGenerator}$, and the corresponding
 direct rotation becomes $\sym{SourceDirectRotation}^*$.\\\addlinespace[5pt]
\multicolumn{3}{@{}l}{\textbf{Displacement extremality}}\\*\addlinespace[3pt]
\path{DK-4.1-prop} &
 Compact angle data; arbitrary transporter $\sym{Wcompetitor}$. &
 For the positive principal angles, there are orthonormal
 $v_j\in\DKRange(\sym{Psource})$ with
 $\angle(v_j,\sym{Wcompetitor}v_j)\ge\like{theta}{\symglyph{theta}_j}$.
 Each singular value of restricted displacement is minimized by
 $\sym{SourceDirectRotation}$, with value $2\sin(\like{theta}{\symglyph{theta}_j}/2)$.\\\addlinespace
\path{DK-4.1-cor} &
 Same displacement setting. &
 $\sym{NormN}((\sym{Identity}-\sym{SourceDirectRotation})\sym{Psource})
 \le\sym{NormN}((\sym{Identity}-\sym{Wcompetitor})\sym{Psource})$.\\\addlinespace
\path{DK-4.2-prop} &
 Any orthonormal basis $(v_j)$ of $\DKRange(\sym{Psource})$;
 arbitrary transporter $\sym{Wcompetitor}$. &
 $\sum_j\sin^2\angle(v_j,\sym{Wcompetitor}v_j)
 \ge\sum_j\sin^2\like{theta}{\symglyph{theta}_j}$, including an infinite right-hand side.\\\addlinespace
\path{DK-4.3-prop} &
 Same displacement setting; put
 $\sym{SquaredDisplacement}(W)=(\sym{Identity}-W)^*(\sym{Identity}-W)$. &
 $\sym{NormN}(\sym{SquaredDisplacement}(\sym{SourceDirectRotation}))
 \le\sym{NormN}(\sym{SquaredDisplacement}(\sym{Wcompetitor}))$:
 the direct rotation minimizes squared full displacement.\\\addlinespace
\path{DK-4.4-prop} &
 Real space; largest principal angle at most $\sym{PiConstant}/3$. &
 The printed assertion that $\sym{SourceDirectRotation}$ minimizes
 every $\sym{NormN}(\sym{Identity}-\sym{Wcompetitor})$ is false.
 A real four-dimensional trace-norm counterexample is given in
 \cref{sec:prop44}; the Q-norm repair is stated there separately.\\\addlinespace[5pt]
\multicolumn{3}{@{}l}{\textbf{Sylvester estimates and spectral cutoffs}}\\*\addlinespace[3pt]
\path{DK-5.1-thm} &
 Banach spaces; bounded $C,X,Y$; $BX-XC=Y$;
 $\dkOperatorNorm{C}\le a$,
 $\dkOperatorNorm{B^{-1}}\le(a+\sym{delta})^{-1}$, $a\ge0$.
 If $B$ is unbounded, it is closed and densely defined with bounded inverse.
 The map norm is compatible with these operator norms. &
 $\sym{delta}\|X\|\le\|Y\|$.
 The coefficient hypotheses may be interchanged.\\\addlinespace
\path{DK-5.2-thm} &
 Self-adjoint $B\ge(\sym{SylvesterGapThreshold}+\sym{delta})\sym{Identity}$,
 $C\le\sym{SylvesterGapThreshold} \sym{Identity}$; bounded $X,Y$;
 $X\DKDom(C)\subset\DKDom(B)$ and $BX-XC=Y$ there. &
 $\sym{delta}\sym{NormN}(X)\le\sym{NormN}(Y)$,
 including unbounded $B,C$.\\\addlinespace
\path{DK-5.1-lem} &
 Orthogonal projections $\sym{StrongApproxProjection}\to \sym{Identity}$ strongly; bounded $K$. &
 $\sym{SingularValue}(K\sym{StrongApproxProjection})\to \sym{SingularValue}(K)$ for each singular-value index $j$.\\\addlinespace[5pt]
\multicolumn{3}{@{}l}{\textbf{Generalized single-angle bounds}}\\*\addlinespace[3pt]
\path{DK-6.1-lem} &
 Each diagonal block of $K$ is dominated by the corresponding block of
 $L$ in every $\sym{NormN}$. &
 The same domination holds for their block-diagonal sums.
 The converse holds if the two $K$ blocks have identical singular values
 (counting multiplicities), as do the two $L$ blocks.\\\addlinespace
\path{DK-6.2-lem} &
 Orthogonal projections $\sym{PinchingProjectionOne},\sym{PinchingProjectionTwo}$; bounded $K$. &
 $\sym{NormN}(\sym{PinchingProjectionOne} K\sym{PinchingProjectionTwo}+
 (\sym{Identity}-\sym{PinchingProjectionOne})K(\sym{Identity}-\sym{PinchingProjectionTwo}))\le\sym{NormN}(K)$.\\\addlinespace
\path{DK-6.1-prop} &
 $\sym{A},\sym{Acurrent}$ have a common domain and bounded difference
 $\sym{H}$;
 both $\sym{SineGap}(\sym{TrialOperator},\sym{Lambda1})$ and
 $\sym{SineGap}(\sym{A1},\sym{Lambda0})$. &
 $\sym{delta}\sym{NormN}(\sin\sym{Theta})\le\sym{NormN}(\sym{H})$.\\\addlinespace
\path{DK-6.1-thm} &
 $\sym{E0}^*\sym{E0}\ge\sym{FrameLowerBoundEpsilon}^2\sym{Identity}$, $\sym{FrameLowerBoundEpsilon}>0$;
 $\sym{SineGap}(\sym{TrialOperator},\sym{Lambda1})$.
 Source and target dimensions may differ. &
 $\sym{delta}\sym{FrameLowerBoundEpsilon}\sym{NormN}(\sin\sym{Theta0})
 \le\sym{NormN}(\sym{R})$.
 The directed sine has the singular values of
 $\sym{Psource}(\sym{Identity}-\sym{Qtarget})$.\\\addlinespace
\path{DK-6.2-thm} &
 Same lower-frame condition; only
 $\sym{SpectralDistanceOperator}(\sym{SpectrumSymbol}(\sym{TrialOperator}),\sym{SpectrumSymbol}(\sym{Lambda1}))
 \ge\sym{delta}$. &
 $\sym{delta}\sym{FrameLowerBoundEpsilon}\|\sin\sym{Theta0}\|_{\mathrm{HS}}
 \le\|\sym{R}\|_{\mathrm{HS}}$.
 If $\sym{R}$ has finite rank, also
 $\sym{delta}\sym{FrameLowerBoundEpsilon}\dkOperatorNorm{\sin\sym{Theta0}}
 \le\dkOperatorNorm{\sym{R}}\sqrt{\sym{RankOperator}\sym{R}}$.\\\addlinespace
\path{DK-6.3-thm} &
 Isometric Rayleigh--Ritz trial data;
 $\sym{OrderedSeparation}_{\sym{delta}}(\sym{TrialOperator},\sym{Lambda1})$.
 The source dimension may be smaller than the target dimension. &
 $\sym{delta}\sym{NormN}(\tan\sym{Theta0})
 \le\sym{NormN}(\sym{R})$.\\\addlinespace
\path{DK-6.3-lem} &
 Bounded $K$; projections $\sym{CompressionProjectionImage},\sym{CompressionProjectionInput}$ of rank at most $n\ge1$;
 $K\sym{CompressionProjectionInput}=\sym{CompressionProjectionImage} K\sym{CompressionProjectionInput}$;
 $\sum_{j\le n}\sym{SingularValue}(K\sym{CompressionProjectionInput})^2$\newline
 ${}>\sum_{j\le n}\sym{SingularValue}(K)^2-\sym{CompressionTolerance}^2$, $\sym{CompressionTolerance}>0$. &
 $\dkOperatorNorm{\sym{CompressionProjectionImage} K(\sym{Identity}-\sym{CompressionProjectionInput})}<\sym{CompressionTolerance}$:
 a nearly maximizing finite-rank compression has small ``leakage''
 (i.e., the norm of this off-diagonal block is below $\sym{CompressionTolerance}$).\\\addlinespace[5pt]
\multicolumn{3}{@{}l}{\textbf{Spectral selection and repulsion}}\\*\addlinespace[3pt]
\path{DK-8.1-thm} &
 Ordered reference gap
 $\sym{OrderedSeparation}_{\sym{delta}}(\sym{ReferenceBlock0},\sym{ReferenceBlock1})$ with endpoint $\sym{alphaEndpoint}$;
 bounded $\sym{H}$, $\sym{H0}=\sym{H1}=0$.
 For the equivalence, $\sym{DimensionOperator}\sym{CrossedSourceDefect}=\sym{DimensionOperator}\sym{CrossedTargetDefect}$. &
 For a chosen reducing target,
 $\sym{Theta}\le\sym{PiConstant}/4$ iff
 $\sym{Lambda0}\le\sym{alphaEndpoint} \sym{Identity}$ and
 $\sym{Lambda1}\ge(\sym{alphaEndpoint}+\sym{delta})\sym{Identity}$; such a reducing target exists.
 With exact frames obtained by applying $\sym{SourceDirectRotation}$ to the reference frames,
 $\sym{ReferenceBlock1}-\sym{alphaEndpoint} \sym{Identity}\le\sym{ComplementCosineBlock}(\sym{Lambda1}-\sym{alphaEndpoint} \sym{Identity})\sym{ComplementCosineBlock}$,
 as quadratic forms, with the lower-block analogue.
 In finite dimensions there are eigenvalue and symmetric-gauge
 repulsion inequalities.\\\addlinespace
\path{DK-8.2-thm} &
 Double-angle sine hypotheses; $\sym{DimensionOperator}\sym{CrossedSourceDefect}=\sym{DimensionOperator}\sym{CrossedTargetDefect}$;
 either bounded $\sym{H}$ with
 $\dkOperatorNorm{\sym{H}}<\sym{delta}/2$, or bounded
 $\sym{R}$ with $\dkOperatorNorm{\sym{R}}<\sym{delta}/2$;
 $\sym{SpectrumSymbol}(\sym{ReferenceBlock0})\subset
 [\sym{betaEndpoint}-\sym{delta}/2,\sym{alphaEndpoint}+\sym{delta}/2]$. &
 The corresponding ambient or directed double-angle sine bound holds
 on the branch with angles strictly below $\sym{PiConstant}/4$.\\
\end{longtable}
\endgroup

\clearpage
\section{Notation correspondence with Davis--Kahan (1970)}
\label{app:dk-notation-correspondence}
This table compares our notation with the corresponding symbols of
Davis and Kahan~\citep[especially \S1]{davis1970rotation}. The correspondences
are by mathematical role, not just by letter. In the original, the reference
block $A_0$ also serves as the residual's comparison operator. Here
$\sym{TrialOperator}$ is a general trial operator, which need not equal that block.
Our additional matching constructions have no direct counterpart in the source.

\begingroup
\small
\setlength{\tabcolsep}{3pt}
\renewcommand{\arraystretch}{1.12}
\begin{longtable}{@{}>{\raggedright\arraybackslash}p{0.43\linewidth}>{\centering\arraybackslash}p{0.12\linewidth}>{\centering\arraybackslash}p{0.16\linewidth}>{\raggedright\arraybackslash}p{0.23\linewidth}@{}}
\caption{Documented correspondence with the source symbols of Davis--Kahan (1970).}
\label{tab:dk-source-notation}\\
\toprule
Mathematical object & DK (1970) & This paper & DK location \\
\midrule
\endfirsthead
\toprule
Mathematical object & DK (1970) & This paper & DK location \\
\midrule
\endhead
\bottomrule
\endfoot
\dkDavisKahanRows
\end{longtable}
\endgroup


\end{document}
, without editing.
% EVERY \sym, \like, and \dkrole obeys this switch, including the glossary.
% No color is required to identify a mathematical object: the glyphs, indices,
% definitions and adjacent prose also distinguish them.
%
% Families (base hue): exact/target=blue, trial/source=amber, exact unwanted=
% cyan, angles and their operators=magenta, residual/displacement=green,
% gaps=violet. Norm gauges=slate; provenance retains brown/steel/vermilion.
% P_ex is the exact-target projection in sine-theta; P is the source
% projection in transport. R denotes a residual and D a direct rotation.
% Semantic keys retain their roles even when the printed notation changes.
\newif\ifDKMathColors
\DKMathColorstrue
\ifdefined\DKMonochromeMath
  \DKMathColorsfalse
\fi

% Family anchors (print-compatible versions of the Manim light palette).
\definecolor{dkfg}{HTML}{1B1F24}
\definecolor{dkmuted}{HTML}{6B7280}
\definecolor{dkwanted}{HTML}{1F6FD1}
\definecolor{dktrial}{HTML}{A66705}
\definecolor{dkunwanted}{HTML}{007C91}
\definecolor{dksine}{HTML}{D12A7B}
\definecolor{dkresid}{HTML}{087D43}
\definecolor{dkgap}{HTML}{6E4BD8}
\definecolor{dkold}{HTML}{705519}
\definecolor{dkcurrent}{HTML}{054669}
\definecolor{dkperturb}{HTML}{BB491B}
\definecolor{dkrefuted}{HTML}{C62828}

% Related yet distinguishable shades. These are NOT new unrelated categories:
% objects inherit their parent family; a different shade signals a different
% mathematical object (e.g., a frame versus its subspace or its projection).
\definecolor{dkwantedframe}{HTML}{1254A9}
\definecolor{dkwantedprojection}{HTML}{08618D}
\definecolor{dkwantedrotation}{HTML}{4248B0}
\definecolor{dkwantedspectrum}{HTML}{345F8F}
\definecolor{dkwantedvector}{HTML}{1468AA}
\definecolor{dktrialframe}{HTML}{905600}
\definecolor{dktrialprojection}{HTML}{79511B}
\definecolor{dktrialcompetitor}{HTML}{A04D10}
\definecolor{dktrialspectrum}{HTML}{79591C}
\definecolor{dktrialvector}{HTML}{9B522A}
\definecolor{dkunwantedframe}{HTML}{006477}
\definecolor{dkunwantedspectrum}{HTML}{00546B}
\definecolor{dkunwantedprojection}{HTML}{145B6D}
\definecolor{dkunwantedvector}{HTML}{256879}
\definecolor{dkangleoperator}{HTML}{923478}
\definecolor{dksineoperator}{HTML}{AA275A}
\definecolor{dksineblock}{HTML}{B43F70}
\definecolor{dksylvester}{HTML}{804074}
\definecolor{dkanglemax}{HTML}{B32185}
\definecolor{dkanglecos}{HTML}{824276}
\definecolor{dkresidvector}{HTML}{086038}
\definecolor{dkdisplacement}{HTML}{226C50}
\definecolor{dkgaporiginal}{HTML}{4E4D89}
\definecolor{dkgapstat}{HTML}{7039A7}
\definecolor{dkgapinterval}{HTML}{574A82}
\definecolor{dkambient}{HTML}{36516A}
\definecolor{dknorm}{HTML}{475875}
\definecolor{dknormaux}{HTML}{66526E}
\definecolor{dkintertwiner}{HTML}{4A6280}
\definecolor{dkturn}{HTML}{47647D}
\definecolor{dkmix}{HTML}{505B65}

% Coloring primitives. "plain" suppresses color without suppressing hovers.
% All descendants can retain a semantic role through \like{key}{expression}.
\expandafter\def\csname dknocolor@plain\endcsname{1}
\newcommand{\dkrole}[2]{%
  \ifDKMathColors
    \ifcsname dknocolor@#1\endcsname #2\else\textcolor{dk#1}{#2}\fi
  \else #2\fi
}
% An optional plain-text reminder supplies definitions and types on hover.
% Required fields: {semantic key}{printed glyph}{color role}{glossary|local}
%                  {short definition}{mathematical type}. The "glossary" flag records whether an
% entry appears in the curated printed notation table; "local" does not.
% This registry is the sole authority for symbol meaning and printed form.
% Use role "plain" for an uncolored symbol that still carries hover text.
% Changing the glyph is a separate editorial decision: this audit preserves
% the current E, F, F_ex, G, and other existing printed choices.
\newcommand{\dkDeclareSymbol}[7][]{%
  \expandafter\def\csname dkgloss@#2\endcsname{#6}%
  \expandafter\def\csname dkglossaryflag@#2\endcsname{#5}%
  \if\relax\detokenize{#1}\relax
    \expandafter\pdfstringdef\csname dkpdfgloss@#2\endcsname{#7: #6}%
  \else
    \expandafter\pdfstringdef\csname dkpdfgloss@#2\endcsname{#7: \detokenize{#1}}%
  \fi
  \expandafter\def\csname dktype@#2\endcsname{#7}%
  \expandafter\def\csname dkglyph@#2\endcsname{#3}%
  \expandafter\def\csname dkrole@#2\endcsname{#4}%
  \expandafter\def\csname dksym@#2\endcsname{\dkrole{#4}{#3}}%
}
% Unstyled printed glyph, for variants constructed through \like.
% Example: \like{E0}{\symglyph{E0}^*} derives E_0^* from E0's
% one canonical glyph; retyping E_0 here would defeat central control.
\DeclareRobustCommand{\symglyph}[1]{\csname dkglyph@#1\endcsname}
\DeclareRobustCommand{\sym}[1]{%
  \ifcsname dksym@#1\endcsname
    \dkSymbolTooltip{#1}{\csname dksym@#1\endcsname}%
  \else
    \PackageError{dk-math-colors}{Unknown symbol key #1}{See math_symbol_colors.tex}%
  \fi
}
\DeclareRobustCommand{\like}[2]{%
  \ifcsname dkrole@#1\endcsname
    \dkSymbolTooltip{#1}{\dkrole{\csname dkrole@#1\endcsname}{#2}}%
  \else
    \PackageError{dk-math-colors}{Unknown role key #1}{See math_symbol_colors.tex}%
  \fi
}

% A short description belongs to the semantic key, not the bare letter.
% The same descriptions are available to the notation appendix.
\DeclareRobustCommand{\dkSymbolDescription}[1]{\csname dkgloss@#1\endcsname}
\DeclareRobustCommand{\dkSymbolGlossaryStatus}[1]{\csname dkglossaryflag@#1\endcsname}
\DeclareRobustCommand{\dkSymbolType}[1]{\csname dktype@#1\endcsname}

% Citable alternative notation is metadata, not a request to rename glyphs.
% The source identifier, the source's exact printed glyph, and the specific
% source location are all mandatory. No alternative is inferred from a name.
% One key may have alternatives from several works or multiple passages.
\newcommand{\dkDeclareNotationSource}[3]{%
  \expandafter\def\csname dksourcetitle@#1\endcsname{#2}%
  \expandafter\def\csname dksourceurl@#1\endcsname{#3}%
}
\newcommand{\dkDavisKahanRows}{}
\makeatletter
\newcommand{\dkDeclareNotationAlternative}[5]{%
  \ifcsname dkglyph@#1\endcsname\else
    \PackageError{dk-math-colors}{Unknown alternative key #1}{Declare the symbol first}%
  \fi
  \ifcsname dksourcetitle@#2\endcsname\else
    \PackageError{dk-math-colors}{Unknown notation source #2}{Declare the source first}%
  \fi
  \expandafter\def\csname dkaltglyph@#2@#1\endcsname{#3}%
  \expandafter\def\csname dkaltlocator@#2@#1\endcsname{#4}%
  \expandafter\def\csname dkaltnote@#2@#1\endcsname{#5}%
  \def\dk@source{#2}%
  \def\dk@original{DK1970}%
  \ifx\dk@source\dk@original
    \g@addto@macro\dkDavisKahanRows{\dkDavisKahanRow{#1}}%
  \fi
}
\makeatother
% A correspondence row is produced from the SAME registry records used by
% \sym and by the source-census script. It is never spelled a second time.
\newcommand{\dkDavisKahanRow}[1]{%
  \csname dkaltnote@DK1970@#1\endcsname &
  $\csname dkaltglyph@DK1970@#1\endcsname$ &
  $\sym{#1}$ &
  \csname dkaltlocator@DK1970@#1\endcsname\\
}

\newif\ifDKSymbolTooltips
\DKSymbolTooltipstrue
\ifdefined\DKNoSymbolTooltips\DKSymbolTooltipsfalse\fi
\newsavebox{\DKTooltipBox}
\newcount\DKTooltipCount
% Register only shipped widgets through hyperref's existing form-field
% machinery. Poppler/Okular need the catalog's field list, not just page
% annotations. The deferred aux write also discards measuring-only copies.
\makeatletter
\newcommand{\dkRegisterTooltipField}{%
  \expandafter\HyField@@AddToFields\expandafter{\number\pdffeedback lastannot}%
}
\makeatother
\AtBeginDocument{\ifDKSymbolTooltips\Form[NeedAppearances=false]\fi}
\AtEndDocument{\ifDKSymbolTooltips\endForm\fi}
% Native LuaLaTeX PDF widgets: no navigation action or JavaScript. Avoid
% pdfcomment's soulpos dependency, which previously broke local Vim builds.
% Measure in the CURRENT math style, then emit the original expression as-is:
% boxing the visible math would change spacing, scripts and line breaking.
\newcommand{\dkSymbolTooltip}[2]{%
  \ifDKSymbolTooltips
    \begingroup
      \DKSymbolTooltipsfalse % suppress nested annotations while measuring
      \ifmmode
        \edef\DKTooltipStyle{\ifcase\mathstyle
          \displaystyle\or\displaystyle\or\textstyle\or\textstyle\or
          \scriptstyle\or\scriptstyle\or\scriptscriptstyle\or\scriptscriptstyle\fi}%
        \sbox{\DKTooltipBox}{$\DKTooltipStyle#2$}%
      \else
        \leavevmode\sbox{\DKTooltipBox}{#2}%
      \fi
      \global\advance\DKTooltipCount by 1
      \pdfextension annot width\wd\DKTooltipBox
        height\ht\DKTooltipBox depth\dp\DKTooltipBox
        {/Subtype /Widget /FT /Btn /Ff 65536 /F 512
         /T (dk-symbol-\the\DKTooltipCount) /TU (\csname dkpdfgloss@#1\endcsname)
         /H /N /BS << /W 0 >> /Border [0 0 0]}%
      \dkRegisterTooltipField
      #2%
    \endgroup
  \else
    #2%
  \fi
}
\pdfstringdefDisableCommands{%
  \def\sym#1{\csname dkglyph@#1\endcsname}%
  \def\like#1#2{#2}%
  \def\symglyph#1{\csname dkglyph@#1\endcsname}%
}

% Declarations: [detailed hover], key, glyph, role, glossary|local, short text, type.
% A local entry has full semantic/hover support, but no notation-table row.
% When two keys have the same glyph, it is an auditable collision, not an
% instruction to change the printed notation automatically.
% This registry includes local proof notation. The notation appendix selects
% the symbols useful for reference; registration does not add a printed row.
% Manim's original symbol keys, now split into semantic subshades.
% Ambient and perturbed operators (provenance is independent of geometry).
\dkDeclareSymbol[Reference self-adjoint operator A on the ambient Hilbert space; the perturbed operator is Ahat = A + H. In a single-operator eigenproblem, A itself is the operator being approximated. An unbounded A acts only on domain(A).]
  {A}{A}{old}{glossary}{Reference self-adjoint operator before perturbation.}{self-adjoint-operator}
\dkDeclareSymbol{A1}{A_1}{old}{local}{Complementary block of the reference operator.}{self-adjoint-operator}
\dkDeclareSymbol{a1}{a_1}{old}{local}{Reference eigenvalue or complementary spectral value.}{scalar}
\dkDeclareSymbol[Reference block A_0: the restriction of A to the range of the reducing source projection P, expressed in source coordinates. In Theorem 8.1 it lies at or below alpha; in Theorem 8.2 its spectrum lies in the enlarged central interval. The perturbed blocks are denoted Lambda_0 and Lambda_1.]
  {ReferenceBlock0}{A_0}{old}{glossary}{Reference operator restricted to the source subspace in Section 8.}{self-adjoint-operator}
\dkDeclareSymbol[Reference block A_1: the restriction of A to the kernel of the reducing source projection P, expressed in complementary coordinates. In Theorem 8.1 it lies at or above alpha + delta in quadratic-form order.]
  {ReferenceBlock1}{A_1}{old}{glossary}{Reference operator restricted to the source complement in Section 8.}{self-adjoint-operator}
\dkDeclareSymbol[Self-adjoint ambient operator Ahat being studied; in the perturbation setting Ahat = A + H. Its action on a trial vector is compared with the proposed eigenvalue times that vector.]
  {Ahat}{\widehat A}{current}{local}{Operator under study; in the perturbation setting, the reference operator plus its perturbation.}{self-adjoint-operator}
\dkDeclareSymbol[Self-adjoint ambient operator Ahat being studied. In the perturbation setting Ahat = A + H. The trial residual is R = Ahat E_0 - E_0 A_tr, and the exact complementary block satisfies Ahat F_1 = F_1 Lambda_1 on its domain.]
  {Acurrent}{\widehat A}{current}{glossary}{Operator under study; in the perturbation setting, the reference operator plus its perturbation.}{self-adjoint-operator}
\dkDeclareSymbol[Additive self-adjoint perturbation H = Ahat - A. A bound on its norm can control the residual of an eigenspace of one operator relative to the other.]
  {H}{H}{perturb}{glossary}{Additive self-adjoint perturbation of the reference operator.}{self-adjoint-operator}
\dkDeclareSymbol{H0}{H_0}{perturb}{local}{Wanted diagonal block of the perturbation.}{self-adjoint-operator}
\dkDeclareSymbol{H1}{H_1}{perturb}{local}{Unwanted diagonal block of the perturbation.}{self-adjoint-operator}
\dkDeclareSymbol{B}{B}{perturb}{local}{Off-diagonal perturbation or coupling block.}{linear-operator}
\dkDeclareSymbol{b}{b}{perturb}{local}{Scalar perturbation or coupling coefficient.}{scalar}
% Trial/source family: subspace, frame, compressed operator, projector, vector.
\dkDeclareSymbol{U}{U}{wanted}{local}{Desired exact subspace in the local geometric notation.}{linear-subspace}% Manim's U remains wanted in local slide notation.
\dkDeclareSymbol{V}{V}{trial}{local}{Trial subspace in the local geometric notation.}{linear-subspace}% Manim's V remains trial in local slide notation.
\dkDeclareSymbol[Source subspace U to be carried onto the target V by an orthogonal or unitary map W. In the transport sections P projects onto U and Q onto V. The matching theorem also names the source subspace K.]
  {Usource}{U}{trial}{glossary}{Source subspace to be transported onto the target subspace.}{linear-subspace}
\dkDeclareSymbol[Trial frame E_0: F -> E, a linear isometry from trial-coordinate space F into ambient space E. Its range is the trial subspace; E_0* returns trial coordinates and E_0* E_0 = I_F. In rank one, E_0(x) = x v.]
  {E0}{E_0}{trialframe}{glossary}{Isometric embedding of the trial subspace; its adjoint returns trial coordinates.}{isometric-embedding}
\dkDeclareSymbol[Self-adjoint trial operator A_tr on trial-coordinate space F. The residual compares Ahat E_0 with E_0 A_tr. Rayleigh--Ritz chooses A_tr = E_0* Ahat E_0 when this compression is defined; the theorem allows other trial operators.]
  {A0}{A_{\mathrm{tr}}}{trialspectrum}{glossary}{Self-adjoint operator on trial coordinates; its eigenvalues are the trial spectral values.}{self-adjoint-operator}
\dkDeclareSymbol[Trial vector: a proposed approximate eigenvector. In the rank-one geometry v has unit length and E_0(x) = x v. Its residual for trial scalar rho is r = Ahat v - rho v.]
  {v}{v}{trialvector}{local}{Trial vector or vector in the trial subspace.}{vector}
\dkDeclareSymbol[Scalar trial operator in rank one, A_tr(x) = rho x. Figure 1 uses the Rayleigh value rho = <v, Ahat v> for unit v; then rho v is the projection of Ahat v onto the trial line and r is perpendicular to v.]
  {rho}{\rho}{trialspectrum}{local}{Trial eigenvalue or scalar trial operator.}{scalar}
\dkDeclareSymbol[Approximate eigenvalue assigned to a trial vector v. For the local operator A, its residual is A v - mu v. For unit v, the Rayleigh--Ritz choice is mu = <A v, v>.]
  {mu}{\mu}{trialspectrum}{local}{Trial spectral value.}{scalar}
\dkDeclareSymbol{a0}{a_0}{trialspectrum}{local}{Wanted trial eigenvalue or trial spectral value.}{scalar}
\dkDeclareSymbol[Ambient orthogonal or unitary transporter W with W(U) = V, equivalently W P W* = Q. Its full displacement is I - W on the whole ambient space. The real matching theorem minimizes its trace norm over every such orthogonal W.]
  {Wcompetitor}{W}{trialcompetitor}{glossary}{Orthogonal or unitary competitor carrying the source subspace onto the target.}{unitary-operator}
\dkDeclareSymbol[Orthogonal projection P onto the source subspace U in the transport problem; P = P* = P squared. A feasible orthogonal or unitary W satisfies W P W* = Q. Restricted displacement (I - W)P tests only source vectors.]
  {Psource}{P}{trialprojection}{glossary}{Orthogonal projection onto the source subspace in the transport problem.}{orthogonal-projection}
\dkDeclareSymbol{ProjSource}{P_U}{trialprojection}{local}{Orthogonal projection onto the source subspace in the transport problem.}{orthogonal-projection}
\dkDeclareSymbol[Vector x in the source subspace U. For principal vectors, unit x_j in U and y_j in V are chosen with <x_i, y_j> = 0 for i different from j and <x_j, y_j> = cos(theta_j).]
  {SourceVector}{x}{trialvector}{local}{Vector in the source subspace.}{vector}
\dkDeclareSymbol[Orthonormal coordinate frame X for the source subspace, so X* X = I. Together with a target frame Y, the singular values of X* Y are the principal cosines. In the four-dimensional counterexample X is the vertical stack of I_2 and 0.]
  {Xframe}{X}{trialframe}{glossary}{Orthonormal coordinate matrix for the source subspace in the counterexample.}{isometric-embedding}
\dkDeclareSymbol[Hilbert space F of trial coordinates. E_0 embeds F isometrically into ambient space E; A_tr and the positive sine operator sin(Theta_0) act on F, while the residual R maps F into E.]
  {Ftrial}{F}{trialframe}{glossary}{Hilbert space of trial coordinates.}{hilbert-space}
\dkDeclareSymbol[Ambient Hilbert space E over the real or complex scalars. It contains the trial subspace range(E_0) and the orthogonal exact subspaces range(F_0) and range(F_1). The source theorem permits infinite-dimensional separable E.]
  {Eambient}{E}{ambient}{glossary}{Ambient Hilbert space containing the trial and exact subspaces.}{hilbert-space}
% Zero and all-ones vectors are distinguished from scalar 0 and 1.
\dkDeclareSymbol[Zero vector of the vector or Hilbert space in context: the additive identity. The zero subspace consists only of this vector.]
  {ZeroVector}{\mathbf{0}}{fg}{glossary}{Zero vector; the additive identity of the ambient vector space.}{vector}
\dkDeclareSymbol[All-ones vector in a specified finite-dimensional coordinate space: every coordinate equals 1. This vector depends on the chosen coordinates and dimension.]
  {OnesVector}{\mathbf{1}}{fg}{local}{All-ones vector in a chosen finite-dimensional coordinate space.}{vector}
% Exact/wanted family: the target subspace, its frame, projection and rotation.
\dkDeclareSymbol[Target subspace V in the transport problem: feasible maps satisfy W(U) = V, or equivalently W P W* = Q. The matching theorem also names this subspace L.]
  {Vtarget}{V}{wanted}{glossary}{Target subspace in the transport problem.}{linear-subspace}
\dkDeclareSymbol[Exact wanted frame F_0: F_ex -> E, a linear isometry onto the desired exact spectral subspace. F_0* returns its coordinates. P_ex = F_0 F_0* is its ambient orthogonal projection; in Figure 1, F_0(h) = (h, 0).]
  {F0}{F_0}{wantedframe}{glossary}{Isometric coordinates for the desired exact spectral subspace; star denotes the adjoint.}{isometric-embedding}
\dkDeclareSymbol[Self-adjoint operator block on coordinates for the desired exact spectral subspace. Its spectrum consists of the wanted exact spectral values; in the corresponding decomposition Ahat F_0 = F_0 Lambda_0 on the operator domain.]
  {Lambda0}{\Lambda_0}{wantedspectrum}{local}{Self-adjoint block on the desired exact spectral subspace.}{self-adjoint-operator}
\dkDeclareSymbol{u}{u}{wantedvector}{local}{Vector in the desired exact subspace.}{vector}
\dkDeclareSymbol{f0}{f_0}{wantedvector}{local}{Component or coordinate vector in the desired exact subspace.}{vector}
\dkDeclareSymbol{lambda0}{\lambda_0}{wantedspectrum}{local}{Wanted exact eigenvalue or wanted spectral value.}{scalar}
\dkDeclareSymbol{lambdau}{\lambda_u}{wantedspectrum}{local}{Eigenvalue associated with a desired exact vector.}{scalar}
\dkDeclareSymbol[Exact wanted orthogonal projection P_ex: E -> E, with P_ex = F_0 F_0*. It retains the desired exact spectral component. Its complementary projection is P_ex_perp = I - P_ex = F_1 F_1*; the projected trial frame is S = P_ex_perp E_0.]
  {Pexact}{P_{\mathrm{ex}}}{wantedprojection}{glossary}{Orthogonal projection onto the desired exact subspace in the sine-theta theorem.}{orthogonal-projection}
\dkDeclareSymbol[Exact complementary orthogonal projection P_ex_perp = I - P_ex = F_1 F_1*: E -> E. It retains the unwanted exact component. In Figure 1, P_ex_perp v = S(1), whereas P_ex_perp r is the unwanted component of the residual.]
  {PexactPerp}{P_{\mathrm{ex}}^\perp}{unwantedprojection}{local}{Orthogonal projection onto the exact complementary subspace.}{orthogonal-projection}
\dkDeclareSymbol{ReflectionSigma}{\Sigma}{wantedprojection}{local}{Reflection associated with the desired exact spectral projection.}{unitary-operator}
\dkDeclareSymbol[Orthogonal projection Q onto the target subspace V in the transport problem. A feasible orthogonal or unitary W satisfies W P W* = Q, where P projects onto the source U.]
  {Qtarget}{Q}{wantedprojection}{glossary}{Orthogonal projection onto the target subspace in the transport problem.}{orthogonal-projection}
\dkDeclareSymbol{ProjTarget}{P_V}{wantedprojection}{local}{Orthogonal projection onto the target subspace in the transport problem.}{orthogonal-projection}
\dkDeclareSymbol[Canonical direct rotation D from the source U to the target V for an acute pair. It is the unitary polar factor of T = QP + (I - Q)(I - P), and acts by the principal angle on each principal plane. R is reserved for the trial residual.]
  {Rrotation}{D}{wantedrotation}{glossary}{Canonical direct rotation carrying the source subspace onto the target.}{unitary-operator}
\dkDeclareSymbol[Direct rotation D between the ranges of orthogonal projections P and Q, including the nonacute case. Relative to range(P) and kernel(P), its diagonal blocks are positive and its off-diagonal blocks are opposite adjoints. It exists when the two crossed defect spaces have equal dimension, and is unique when both defects vanish.]
  {SourceDirectRotation}{D}{wantedrotation}{local}{Direct rotation, including the nonacute case in the source-result table.}{unitary-operator}
\dkDeclareSymbol[Orthonormal coordinate frame Y for the target subspace, so Y* Y = I. In the counterexample Y is the vertical stack of cos(theta) I_2 and sin(theta) J_0; hence X* Y = cos(theta) I_2 and both principal angles equal theta.]
  {Yframe}{Y}{wantedframe}{glossary}{Orthonormal coordinate matrix for the target subspace in the counterexample.}{isometric-embedding}
\dkDeclareSymbol[Vector y in the target subspace V. For principal vectors, unit x_j in U and y_j in V are paired so that <x_j, y_j> = cos(theta_j); distinct pairs have zero cross inner product.]
  {TargetVector}{y}{wantedvector}{local}{Vector in the target subspace.}{vector}
\dkDeclareSymbol[Hilbert space F_ex of exact wanted coordinates. The isometry F_0: F_ex -> E identifies it with the desired exact spectral subspace. This space differs from H, the additive perturbation.]
  {Hexact}{F_{\mathrm{ex}}}{wantedframe}{glossary}{Hilbert space of coordinates for the desired exact spectral subspace.}{hilbert-space}% Coordinate Hilbert space (not perturbation).
% Exact/unwanted family; the trial error operators are MAGENTA, not cyan.
\dkDeclareSymbol{Uperp}{U^{\perp}}{unwanted}{local}{Orthogonal complement of the desired exact subspace in the local notation.}{linear-subspace}
\dkDeclareSymbol[Exact complementary frame F_1: G -> E, a linear isometry onto the orthogonal complement of range(F_0). F_1* returns unwanted coordinates; F_1 F_1* = I - P_ex. In Figure 1, F_1(g) = (0, g).]
  {F1}{F_1}{unwantedframe}{glossary}{Isometric coordinates for the exact complementary subspace; star denotes the adjoint.}{isometric-embedding}
\dkDeclareSymbol[Self-adjoint unwanted block Lambda_1 on G, satisfying Ahat F_1 = F_1 Lambda_1 on its domain. Its spectrum is separated from the trial spectrum of A_tr by the specified positive delta. Figure 1 uses the scalar block 0.7.]
  {Lambda1}{\Lambda_1}{unwantedspectrum}{glossary}{Self-adjoint operator on exact complementary coordinates: the unwanted spectral block.}{self-adjoint-operator}
% Spectral measure and set (distinct from the ambient space E and sine block S).
\dkDeclareSymbol[Projection-valued spectral measure E associated with a self-adjoint operator A: A = integral lambda dE(lambda). For a Borel set S, E(S) is the orthogonal projection onto the corresponding spectral subspace.]
  {SpectralMeasure}{E}{wantedprojection}{glossary}{Projection-valued spectral measure of a self-adjoint operator.}{projection-valued-measure}
\dkDeclareSymbol[Borel subset S of the real spectral axis. The spectral projection E(S) selects the part of a self-adjoint operator with spectral values in S; its range is a reducing subspace.]
  {SpectralSet}{S}{wantedspectrum}{glossary}{Borel spectral set used to select a spectral subspace.}{set}
% Generic Sylvester coefficients and right side, by source/target role.
\dkDeclareSymbol[Left coefficient B of the Sylvester map X -> B X - X C; B acts on the output space of X. For the sine bound, B = Lambda_1 on the unwanted-coordinate space G.]
  {SylvesterLeft}{B}{unwantedspectrum}{glossary}{Left coefficient of the Sylvester equation; corresponds to the unwanted spectral block.}{self-adjoint-operator}
\dkDeclareSymbol[Right coefficient C of the Sylvester map X -> B X - X C; C acts on the input space of X. For the sine bound, C = A_tr on trial space F. The historical equation also denotes this coefficient by A.]
  {SylvesterRight}{C}{trialspectrum}{glossary}{Right coefficient of the Sylvester equation; corresponds to the trial operator.}{self-adjoint-operator}
\dkDeclareSymbol[Right-hand side Y in B X - X C = Y, with the same input and output spaces as X. For the sine bound, Y = F_1* R: F -> G, the residual expressed in unwanted exact coordinates.]
  {SylvesterRHS}{Y}{resid}{glossary}{Right-hand side of the Sylvester equation, obtained by projecting the residual.}{linear-operator}
\dkDeclareSymbol{w}{w}{unwantedvector}{local}{Vector in an unwanted or complementary spectral subspace.}{vector}
\dkDeclareSymbol{f1}{f_1}{unwantedvector}{local}{Component or coordinate vector in the unwanted exact subspace.}{vector}
\dkDeclareSymbol{lambda1}{\lambda_1}{unwantedspectrum}{local}{Unwanted exact eigenvalue or complementary spectral value.}{scalar}
\dkDeclareSymbol{lambdaw}{\lambda_w}{unwantedspectrum}{local}{Eigenvalue associated with an unwanted exact vector.}{scalar}
\dkDeclareSymbol[Hilbert space G of exact complementary coordinates. F_1: G -> E embeds it onto the unwanted exact subspace; Lambda_1 acts on G and X = F_1* E_0 maps trial coordinates into G.]
  {Gunwanted}{G}{unwantedframe}{glossary}{Hilbert space of exact complementary coordinates.}{hilbert-space}
\dkDeclareSymbol{E1}{E_1}{unwantedframe}{local}{Isometric coordinates for the complementary trial subspace.}{isometric-embedding}
\dkDeclareSymbol{Vperp}{V^{\perp}}{unwanted}{local}{Orthogonal complement of the trial subspace in the local notation.}{linear-subspace}
% Angle and error geometry, including distinct rectangular and positive maps.
\dkDeclareSymbol[Individual principal angle, between 0 and pi/2. For orthonormal subspace frames X and Y, the singular values of X* Y are the principal cosines. In Figure 1, theta is the angle between the unit trial vector v and the exact wanted line.]
  {theta}{\theta}{sine}{glossary}{Principal angle between the subspaces under comparison.}{angle}
\dkDeclareSymbol[Full positive angle operator Theta = diag(Theta_0, Theta_1) for the source/target decompositions. Theta_0 compares the source with the target; Theta_1 compares their orthogonal complements. Its operator norm is the largest spectral angle. For positive c, the conditions Theta <= c and Theta < c mean operator norm at most c and strictly below c, respectively.]
  {Theta}{\Theta}{angleoperator}{glossary}{Full angle operator, the diagonal sum of the source and complementary angle blocks.}{angle-operator}
\dkDeclareSymbol[Directed principal-angle operator Theta_0 on source or trial coordinates. For a trial frame E_0 and exact wanted projection P_ex, cos(Theta_0) = sqrt(E_0* P_ex E_0), and sin(Theta_0) = sqrt(I_F - E_0* P_ex E_0). In the two-projection geometry it comes from the positive cosine compression P Q P restricted to range(P).]
  {Theta0}{\Theta_0}{angleoperator}{glossary}{Directed principal-angle operator acting on source or trial coordinates.}{angle-operator}
\dkDeclareSymbol[Complementary directed angle operator Theta_1, obtained from the projections I-P and I-Q by the same cosine-compression construction as Theta_0. It acts on the source complement, or its isometric coordinates. The full angle operator is diag(Theta_0, Theta_1).]
  {Theta1}{\Theta_1}{angleoperator}{glossary}{Directed angle operator comparing the orthogonal complements.}{angle-operator}
\dkDeclareSymbol[Right-angle source defect D_0 = range(P) intersect kernel(Q). These are source vectors orthogonal to the target. The full-angle conclusion in the Section 8 smallness criterion assumes dim(D_0) = dim(D_1); in Lean their identification is CrossedDefectsEquivalent.]
  {CrossedSourceDefect}{D_0}{fg}{glossary}{Source directions perpendicular to the target subspace.}{linear-subspace}
\dkDeclareSymbol[Right-angle target defect D_1 = kernel(P) intersect range(Q). These are target vectors orthogonal to the source. Equality of its Hilbert dimension with that of D_0 is assumed for the full-angle conclusion of Theorem 8.2.]
  {CrossedTargetDefect}{D_1}{fg}{glossary}{Target directions perpendicular to the source subspace.}{linear-subspace}
\dkDeclareSymbol[Positive complementary cosine block C_1 = cos(Theta_1). The reference complement is identified with the perturbed complement by the direct rotation D, so A_1 and Lambda_1 act on common coordinates. Theorem 8.1 compares A_1 - alpha I with C_1 (Lambda_1 - alpha I) C_1 as quadratic forms.]
  {ComplementCosineBlock}{C_1}{anglecos}{local}{Cosine of the complementary angle block in the spectral-repulsion inequality.}{self-adjoint-operator}
\dkDeclareSymbol[Scalar sine of a principal angle. For a unit principal trial vector v and exact wanted projection P_ex, sin(theta) = ||(I - P_ex)v||. In Figure 1 it is the length of the pink leg S(1).]
  {sintheta}{\sin\theta}{sine}{local}{Sine of an individual principal angle: the length of the complementary component of a unit principal trial vector.}{scalar}
\dkDeclareSymbol[Positive sine operator sin(Theta_0) = sqrt(I_F - E_0* P_ex E_0): F -> F. It equals |S| for S = (I - P_ex)E_0: F -> E. Its eigenvalues in finite dimensions are the principal sines; the sine theorem bounds its ideal norm by N(R)/delta.]
  {sinTheta0}{\sin\Theta_0}{sineoperator}{glossary}{Positive sine operator on trial coordinates; the modulus of the rectangular complementary error.}{angle-operator}
\dkDeclareSymbol[Rectangular complementary error S = (I - P_ex) E_0: F -> E. It records how much each trial vector leaks outside the exact wanted subspace. Its modulus |S| = sqrt(S* S) is sin(Theta_0): F -> F. In rank one S(1) = P_ex_perp v.]
  {S}{S}{sineblock}{glossary}{Rectangular complementary error: the trial frame projected onto the unwanted exact subspace.}{linear-operator}
\dkDeclareSymbol[Complementary-coordinate map X = F_1* E_0: F -> G, with S = F_1 X. It satisfies Lambda_1 X - X A_tr = F_1* R on domain(A_tr). Isometric multiplication gives N(X) = N(S); the separate modulus identity abs(S) = sin(Theta_0) gives N(S) = N(sin(Theta_0)), whenever these norms are finite.]
  {X}{X}{sylvester}{glossary}{Solution of the Sylvester equation; for the sine bound, the trial frame in unwanted exact coordinates.}{linear-operator}
\dkDeclareSymbol{alphaAngle}{\alpha}{anglemax}{local}{Largest principal angle.}{angle}
\dkDeclareSymbol[Scalar c = cos(theta) in the counterexample. In the general two-projection decomposition, the same role denotes a positive cosine operator C; the cosine and sine blocks commute and satisfy C squared + S squared = I.]
  {CosAngle}{c}{anglecos}{local}{Cosine of a principal angle or the corresponding positive cosine block.}{scalar-or-operator}
\dkDeclareSymbol[Scalar s = sin(theta) in the counterexample. In the general two-projection decomposition, the same role denotes a positive sine operator S; it commutes with the positive cosine block C and satisfies C squared + S squared = I.]
  {SinAngle}{s}{sineoperator}{local}{Sine of a principal angle or the corresponding positive sine block.}{scalar-or-operator}
\dkDeclareSymbol[Positive small-angle parameter in the local subspace comparison. It controls the allowed angular deviation; it is an angle parameter rather than a perturbation-operator norm.]
  {EpsilonAngle}{\varepsilon}{anglemax}{local}{Angle parameter controlling the maximum subspace deviation.}{angle}
\dkDeclareSymbol{PhiAngle}{\varphi}{angleoperator}{local}{Auxiliary angle or angle operator in the two-subspace comparison.}{angle}
\dkDeclareSymbol{Bblock}{B}{sineblock}{local}{Off-diagonal block in the two-subspace decomposition.}{linear-operator}
\dkDeclareSymbol{Cblock}{C}{sylvester}{local}{Complementary or Sylvester block in the two-subspace decomposition.}{linear-operator}
\dkDeclareSymbol{Ablock}{A}{anglecos}{local}{Positive cosine-related block in the two-subspace decomposition.}{linear-operator}
% Residual/error-displacement family.
\dkDeclareSymbol[Residual vector r = Ahat v - rho v = R(1) for the rank-one frame E_0(x) = x v. Figure 1 draws it from rho v to Ahat v. For unit v and the Rayleigh value rho = <v, Ahat v>, <v, r> = 0; a general trial scalar need not give orthogonality.]
  {r}{r}{residvector}{local}{Residual vector of an approximate eigenpair.}{vector}
\dkDeclareSymbol[Bounded residual R: F -> E, agreeing with Ahat E_0 - E_0 A_tr on domain(A_tr). The domain condition requires E_0(domain(A_tr)) inside domain(Ahat). R = 0 means that E_0 intertwines A_tr with Ahat; its norm measures the defect.]
  {R}{R}{resid}{glossary}{Bounded residual of the trial eigenproblem, agreeing with the operator difference on the trial domain.}{linear-operator}
\dkDeclareSymbol[Full displacement Delta_W = I - W on the ambient space. In the failed Ky Fan comparison W = W_(pi/4); Delta_W has singular values sqrt(2), sqrt(2), 0, 0. Multiplying by a principal-plane projection restricts its input to that plane.]
  {Kdisplacement}{\Delta_W}{displacement}{glossary}{Full displacement of a transporter from the identity.}{operator-valued-function}
\dkDeclareSymbol[Squared full displacement Dcal(W) = (I-W)* (I-W), a positive operator. Proposition 4.3 compares its unitarily invariant norm for a direct rotation D with that for every feasible transporter W.]
  {SquaredDisplacement}{\mathcal D}{displacement}{local}{Squared full displacement as a function of the transporter.}{operator-valued-function}
% Matching proof: the gradient and its spectral level are not a space or gap.
\dkDeclareSymbol[Skew-symmetric spectral gradient G of Phi_p. It commutes with W and represents first variations by d/dt Phi_p(W exp(tB)) at t = 0 equal to -tr(G B) for skew-symmetric B. Feasible stationarity gives J_src G J_src = -G.]
  {SpectralGradient}{G}{displacement}{local}{Skew-symmetric spectral gradient of the displacement power sum in the matching proof.}{linear-operator}
\dkDeclareSymbol[Positive spectral level g of the skew-symmetric gradient G: on this invariant subspace, -G squared = g squared I. Consequently G/g squares to -I and supplies the complex structure used in the real matching proof.]
  {GradientLevel}{g}{displacement}{local}{Positive gradient level; the gradient squared is minus this level squared times the identity.}{scalar}
\dkDeclareSymbol[Complex structure I_g = G/g on the positive gradient-level space V_g. It is a real isometry with square -I. Multiplication by a+ib means a times identity plus b times I_g; source reflection J_src reverses I_g and acts as conjugation.]
  {GradientComplexStructure}{\mathcal I_g}{displacement}{local}{Complex multiplication constructed from a positive gradient level.}{linear-operator}
\dkDeclareSymbol[Real skew-symmetric matrix B, with B transpose = -B, giving the orthogonal variation W exp(tB). If B commutes with the source projection P, this curve preserves the constraint W(U) = V.]
  {SkewDirection}{B}{ambient}{local}{Skew-symmetric direction for a first variation of an orthogonal transporter.}{linear-operator}
% Spectral separation family: gap and original gap are NOT identical.
\dkDeclareSymbol[Positive spectral separation in the stated configuration. In the sine theorem it separates the trial operator from the exact complementary block; Theorem 8.1 uses it for the ordered reference gap, and Theorem 8.2 for the perturbed central-band separation. Figure 1 uses delta = abs(rho - 0.7).]
  {delta}{\delta}{gap}{glossary}{Positive spectral separation in the stated gap configuration.}{scalar}
\dkDeclareSymbol{g}{g}{gaporiginal}{local}{Original spectral gap of the reference operator; distinct from the certified sine-theta separation.}{scalar}
\dkDeclareSymbol{Delta}{\Delta}{gapstat}{local}{Population or statistical spectral gap.}{scalar}
\dkDeclareSymbol[Spectral endpoint alpha. In the interval/exterior configuration, one spectrum lies between beta and alpha, the other at or below beta - delta or at or above alpha + delta. Theorem 8.2 places the selected perturbed block in the central interval. In Theorem 8.1, alpha bounds the lower reference block and alpha + delta bounds the upper block from below.]
  {alphaEndpoint}{\alpha}{gapinterval}{glossary}{Upper endpoint of the central band, or lower-block cut in the ordered gap.}{scalar}
\dkDeclareSymbol[Lower endpoint beta of the finite interval from beta to alpha in the gap hypothesis. One spectrum lies in that interval; the other lies at or below beta - delta or at or above alpha + delta.]
  {betaEndpoint}{\beta}{gapinterval}{glossary}{Lower endpoint of the finite interval in the interval/exterior gap hypothesis.}{scalar}
% Spectral/projector and auxiliary matrix objects, context-scoped.
\dkDeclareSymbol{P}{P}{fg}{local}{Orthogonal projection; its source or exact-target role is specified locally.}{orthogonal-projection}% Neutral historical Manim key; prefer Pexact/Psource.
\dkDeclareSymbol{Q}{Q}{fg}{local}{Orthogonal projection; its role is specified locally.}{orthogonal-projection}% Neutral historical Manim key; prefer Qtarget.
\dkDeclareSymbol[Two-projection intertwiner T = QP + (I - Q)(I - P), satisfying TP = QT and T* T = I - (P - Q) squared. For ||P - Q|| < 1 it is invertible; its unitary polar factor T |T| inverse is the canonical direct rotation.]
  {IntertwinerT}{T}{intertwiner}{glossary}{Two-projection intertwiner; its unitary polar factor is the acute direct rotation.}{linear-operator}
\dkDeclareSymbol[Planar pi/2-rotation matrix J_0 = ((0, -1), (1, 0)), with J_0 squared = -I_2 and J_0 transpose = -J_0. It defines the rotation C = cos(2 theta) I_2 + sin(2 theta) J_0 used to construct the real competitor.]
  {Jturn}{J_0}{turn}{glossary}{Planar rotation through $\pi/2$, with square minus the identity.}{matrix}
\dkDeclareSymbol[Partial isometry V_pol in the polar decomposition S = V_pol |S|, where |S| = sqrt(S* S). It is isometric on the orthogonal complement of kernel(S) and zero on kernel(S). For an invertible operator its polar factor is unitary.]
  {PolarU}{V_{\mathrm{pol}}}{intertwiner}{glossary}{Partial isometry in the polar decomposition; unitary when the operator is invertible.}{partial-isometry}
\dkDeclareSymbol[Planar rotation C = cos(2 theta) I_2 + sin(2 theta) J_0 through twice the principal angle. The competitor W_theta is orthogonally conjugate to C direct-sum I_2, so it rotates one two-plane and fixes the other.]
  {Cturn}{C}{turn}{glossary}{Planar rotation through twice the principal angle in the counterexample.}{matrix}
\dkDeclareSymbol[Orthogonal mixing matrix H_mix = (1/sqrt(2)) times the block matrix ((I_2, I_2), (I_2, -I_2)). It constructs W_theta = H_mix (C direct-sum I_2) H_mix transpose.]
  {Hmix}{H_{\mathrm{mix}}}{mix}{glossary}{Orthogonal mixing matrix in the counterexample.}{matrix}
\dkDeclareSymbol[Source reflection J_src = 2P - I: it acts as +I on the source subspace and -I on its orthogonal complement, so its square is I. In the matching proof stationarity gives J_src G J_src = -G for the skew-symmetric spectral gradient G.]
  {Jreflect}{J_{\mathrm{src}}}{trialprojection}{glossary}{Reflection equal to twice the source projection minus the identity.}{linear-operator}
\dkDeclareSymbol[Skew angle generator J_ang for a direct rotation D = exp(J_ang Theta). On each principal plane it is the oriented rotation through pi/2; it vanishes on the common fixed directions. It commutes with the full angle operator Theta.]
  {AngleGenerator}{J_{\mathrm{ang}}}{turn}{glossary}{Skew angle generator in the direct-rotation result summaries.}{linear-operator}
\dkDeclareSymbol{Dblock}{D}{mix}{local}{Auxiliary diagonal or block matrix in the counterexample.}{matrix}
\dkDeclareSymbol[Orthogonal projection onto the first principal plane span(e_0, e_3) in the four-dimensional counterexample. For Delta_W = I - W_(pi/4), the Ky Fan 2-norm of Delta_W Omega_1 equals 2.]
  {OmegaOne}{\Omega_1}{turn}{glossary}{First principal-plane projection in the failed Ky Fan comparison.}{orthogonal-projection}
\dkDeclareSymbol[Orthogonal projection onto the second principal plane span(e_1, e_2) in the four-dimensional counterexample. For Delta_W = I - W_(pi/4), the Ky Fan 2-norm of Delta_W Omega_2 equals 2, although the trace norm of Delta_W is only 2 sqrt(2).]
  {OmegaTwo}{\Omega_2}{intertwiner}{glossary}{Second principal-plane projection in the failed Ky Fan comparison.}{orthogonal-projection}
\dkDeclareSymbol{Dplanar}{D}{wantedrotation}{local}{Planar direct-rotation block.}{unitary-operator}
\dkDeclareSymbol{Hreflection}{H}{trialcompetitor}{local}{Planar reflection carrying the source line to the target.}{linear-operator}
% Unitarily invariant gauge family; N, M are NOT a norm versus its input.
\dkDeclareSymbol[Unitarily invariant norm N, or a symmetric operator-ideal norm on its domain of finiteness. It satisfies N(U T V) = N(T) for unitaries and N(B T) <= ||B||op N(T). In the sine theorem it is normalized by N(u v*) = ||u|| ||v|| and has Ky Fan dominance where both norms are finite. Both R and sin(Theta_0) must belong to its ideal.]
  {NormN}{N}{norm}{glossary}{Unitarily invariant norm on its domain of finiteness; in the sine theorem, normalized by $N(uv^*)=\|u\|\|v\|$.}{norm}
\dkDeclareSymbol[Sine-theta separation G_delta(B,C): finite interval/exterior separation in either spectral order, or ordered separation B <= a I and C >= (a+delta) I, or the reverse. The ordered alternatives allow both spectra to be unbounded.]
  {SineGap}{\mathcal G_\delta}{gap}{local}{Interval/exterior separation or either ordered separation, as in the sine-theta theorem.}{predicate}
\dkDeclareSymbol[Auxiliary unitarily invariant norm M witnessing the Q-norm property N(T) squared = M(T* T). For the Frobenius norm N, one can take M to be the trace norm.]
  {NormM}{M}{normaux}{local}{Auxiliary unitarily invariant norm used to define a Q-norm through the squared modulus.}{norm}
\dkDeclareSymbol[Q-norm property: a unitarily invariant norm N is a Q-norm when N(T) squared = M(T* T) for some unitarily invariant norm M. The Frobenius norm is an example, with M the trace norm; the repaired full-displacement theorem uses this norm class.]
  {NormQ}{Q}{normaux}{glossary}{Q-norm: its square is a unitarily invariant norm of the squared modulus.}{norm}
\dkDeclareSymbol[Singleton trace cost s_i = min(4 sin(theta_i/2), 2) for angle label i. It transports one real principal plane independently, choosing a rotation of cost 4 sin(theta_i/2) or a reflection of cost 2.]
  {SingletonCost}{s}{norm}{glossary}{Minimum trace displacement for one real principal plane: choose a rotation or a reflection.}{scalar-cost}
\dkDeclareSymbol[Paired trace cost e_ij = 4 sin((theta_i + theta_j)/2) for distinct angle labels i,j. It is attained by coupling their principal planes in a real four-dimensional block. The edge weight s_i + s_j - e_ij is the difference between the cost of two singletons and this paired cost.]
  {PairCost}{e}{norm}{glossary}{Trace displacement cost of coupling two principal planes in a four-dimensional block.}{scalar-cost}
\dkDeclareSymbol[Scalar phase cost f_p(t) = (2 sin(t/2)) to the power p, for 0 <= t <= pi. A real planar rotation through t contributes two copies, so its displacement power cost is 2 f_p(t).]
  {MatchingPhaseCost}{f_p}{norm}{local}{Displacement power cost for one scalar phase.}{scalar-cost}
\dkDeclareSymbol[Singleton power cost a_i(p) = min(2 f_p(theta_i), 2 to the power p). These are the rotation and reflection choices for one principal plane. At p=1 it equals the trace-cost weight s_i.]
  {SingletonPowerCost}{a}{norm}{local}{Minimum displacement power cost for one angle.}{scalar-cost}
\dkDeclareSymbol[Paired power cost b_ij(p) = min over d <= tau <= h of 2(f_p(h+tau) + f_p(h-tau)), where h=(theta_i+theta_j)/2 and d=abs(theta_i-theta_j)/2. This is the positive-determinant pair minimum; unrestricted block optimization also compares two singletons. At p=1 it equals e_ij.]
  {PairPowerCost}{b}{norm}{local}{Positive-determinant displacement power cost for a pair of angles.}{scalar-cost}
\dkDeclareSymbol[Matching M on the finite general graph of positive principal-angle labels. Each label belongs to at most one edge; unmatched labels use singleton transporters. Minimizing paired and singleton costs is equivalent to maximum-weight matching with edge weights g_ij = s_i + s_j - e_ij. The empty matching is allowed.]
  {Matching}{M}{normaux}{glossary}{Matching of positive principal-angle labels: each label is used in at most one pair.}{matching}
\dkDeclareSymbol[Displacement power sum Phi_p(W) = sum_j s_j(I - W)^p. The matching proof first minimizes it for 1 < p < 2, then lets p decrease to 1, where it becomes the trace norm ||I - W||_1.]
  {PowerCost}{\Phi}{norm}{local}{Sum of the p-th powers of displacement singular values.}{scalar-cost}
% Auxiliary matching-proof notation: neutral ink, with the same hover registry.
\dkDeclareSymbol[Off-diagonal gradient block G_01 mapping the source complement K perpendicular into K. Stationarity makes the diagonal blocks of G zero, so G = ((0,G_01),(-G_01 transpose,0)) in the source/complement decomposition.]
  {GradientOffDiagonal}{G_{01}}{fg}{local}{Complement-to-source block of the skew-symmetric gradient.}{linear-operator}
\dkDeclareSymbol[Orthogonal group O(n): real n by n matrices W satisfying W transpose times W = I. Its elements preserve Euclidean inner products and lengths. The notation O(2) refers to the same group in dimension two.]
  {OrthogonalGroup}{O(n)}{fg}{local}{Group of real orthogonal matrices.}{group}
\dkDeclareSymbol[Feasible transporters F = {W in O(n) : W(K) = L}. Equivalently W P_K W transpose = P_L. This nonempty compact set is the optimization domain throughout the matching proof.]
  {FeasibleTransport}{\mathcal F}{fg}{local}{Orthogonal maps carrying the source subspace onto the target.}{set}
\dkDeclareSymbol[Singular value s_j(T), the nonnegative square root of an eigenvalue of T transpose times T, counted with multiplicity. Summing s_j(I-W) gives the trace displacement; summing its p-th powers gives Phi_p(W). These singular values differ from the singleton weights s_i.]
  {TransportSingularValues}{\symglyph{SingularValueBase}_j}{fg}{local}{A displacement singular value in the matching proof.}{scalar}
\dkDeclareSymbol[Scalar singular-value magnitude s between zero and two for the displacement I-W of an orthogonal matrix. Uniform convergence of s to the power p to s on this interval permits the trace-norm limit. This variable differs from the rotation parameter s used earlier in the proof.]
  {DisplacementSize}{s}{fg}{local}{Scalar displacement singular-value magnitude in the uniform limit.}{scalar}
\dkDeclareSymbol[All-singleton cost C_0 = sum_i s_i, before any two angle labels are paired. A matching has cost C_0 minus the sum of its selected edge weights.]
  {SingletonBaseline}{C_0}{fg}{local}{Trace displacement when every angle uses its singleton transporter.}{scalar-cost}
\dkDeclareSymbol[Matching power cost C_p(M): sum of b_ij(p) over selected pairs plus a_i(p) over unmatched labels. Its minimum equals the minimum of Phi_p over feasible transporters. At p=1 it is C(M).]
  {MatchingPowerCost}{C_p}{fg}{local}{Total paired and singleton power cost of a matching.}{scalar-cost}
\dkDeclareSymbol[Invariant orthogonal block V_a preserved by the source projection, target projection, and minimizing transporter. The displacement cost adds over these blocks, and their principal-angle lists combine into the original list.]
  {MatchingBlock}{V_a}{fg}{local}{Simultaneous invariant block in the matching proof.}{linear-subspace}
\dkDeclareSymbol[Gradient-level space V_g: the eigenspace of -G squared for a positive eigenvalue g squared. It is preserved by G, W, and the source reflection J_src. Multiplication by i on it is G/g.]
  {GradientLevelSpace}{V_g}{fg}{local}{Positive gradient-level eigenspace.}{linear-subspace}
\dkDeclareSymbol[Eigenphase varphi of W on V_g viewed as a complex space with multiplication by i = G/g. It lies between zero and pi and satisfies f_p prime(varphi) = g. There are at most two such phases.]
  {GradientEigenphase}{\varphi}{fg}{local}{One of the at most two phases at a positive gradient level.}{angle}
\dkDeclareSymbol[Orthogonal projection E onto the first phase space of W on V_g. Here W = exp(i varphi_1) E + exp(i varphi_2) (I-E). In coordinates where the source reflection is conjugation, E = X + iY.]
  {GradientPhaseProjection}{E}{fg}{local}{Projection onto the first phase space in a gradient-level space.}{orthogonal-projection}
\dkDeclareSymbol[Real part X of the phase projection E = X + iY. It is a real symmetric matrix with spectrum in the closed interval from zero to one, satisfying X squared - Y squared = X and XY + YX = Y. This X differs from the source-coordinate and Sylvester matrices elsewhere in the paper.]
  {ProjectionRealPart}{X}{fg}{local}{Real symmetric part of the phase projection.}{linear-operator}
\dkDeclareSymbol[Imaginary part Y of the phase projection E = X + iY. It is a real skew-symmetric matrix. It maps the lambda eigenspace of X to the 1-lambda eigenspace and constructs the small invariant planes.]
  {ProjectionImaginaryPart}{Y}{fg}{local}{Real skew-symmetric imaginary part of the phase projection.}{linear-operator}
\dkDeclareSymbol[Eigenvalue lambda in the closed interval from zero to one of the real part X of the phase projection. It is local to the block construction and differs from the original operator's spectral parameter.]
  {ProjectionEigenvalue}{\lambda}{fg}{local}{Eigenvalue of the real part of the phase projection.}{scalar}
\dkDeclareSymbol[Plane coupling eta = sqrt(lambda (1-lambda)) for 0 < lambda < 1. If u is a unit lambda eigenvector of X, then v = Yu/eta completes its invariant real plane.]
  {ProjectionCoupling}{\eta}{fg}{local}{Normalization used to construct an invariant real plane.}{scalar}
\dkDeclareSymbol[Pair mean angle h = (theta_i + theta_j)/2. Together with d = (theta_i - theta_j)/2 after ordering theta_i >= theta_j, it parametrizes the two-angle block cost.]
  {PairMeanAngle}{h}{fg}{local}{Mean of the two ordered principal angles.}{angle}
\dkDeclareSymbol[Pair half-difference d = (theta_i - theta_j)/2 after ordering theta_i >= theta_j. It satisfies 0 <= d <= h and h+d <= pi/2.]
  {PairHalfDifference}{d}{fg}{local}{Half the difference of the two ordered principal angles.}{angle}
\dkDeclareSymbol[Scalar pair parameter tau in the closed interval from d to h. The paired power cost is the minimum of 2(f_p(h+tau) + f_p(h-tau)) over this interval. At p=1 a minimum is attained at tau=h.]
  {PairPhaseParameter}{\tau}{fg}{local}{Scalar parameter in the two-angle power-cost minimum.}{scalar}
\dkDeclareSymbol[Attaining rotation angle gamma = arccos(cos(tau)/cos(d)). Choosing U_0 = U_1 = R_gamma realizes phase parameters x=h and y=tau in the four-dimensional block.]
  {PairAttainingAngle}{\gamma}{fg}{local}{Rotation angle realizing a chosen scalar pair parameter.}{angle}
\dkDeclareSymbol[Four-dimensional direct rotation D = ((C,-S),(S,C)) in source and complementary principal coordinates, with C the diagonal principal cosines and S the diagonal principal sines. Every transporter on this block is D diag(U_0,U_1).]
  {PairDirectRotation}{D}{fg}{local}{Direct rotation on the two principal planes.}{unitary-operator}
\dkDeclareSymbol[Diagonal cosine block C = diag(cos(theta_i), cos(theta_j)) in the four-dimensional direct rotation D. This matrix differs from the scalar displacement cost C(M).]
  {PairCosineBlock}{C}{fg}{local}{Diagonal principal-cosine block.}{matrix}
\dkDeclareSymbol[Diagonal sine block S = diag(sin(theta_i), sin(theta_j)) in the four-dimensional direct rotation D. This matrix is local to the matching proof.]
  {PairSineBlock}{S}{fg}{local}{Diagonal principal-sine block.}{matrix}
\dkDeclareSymbol[Free planar orthogonal matrices U_0 and U_1 acting on the two-dimensional source and its two-dimensional complement, respectively. Multiplying D by diag(U_0,U_1) parametrizes every transporter on this four-dimensional block.]
  {PairFreeOrthogonal}{U}{fg}{local}{Free source or complementary planar orthogonal factor.}{unitary-operator}
\dkDeclareSymbol[Phase parameter x in the closed interval from h to pi-h for a positive-determinant four-dimensional transporter. Its eigenvalues are exp(plus or minus i(x+y)) and exp(plus or minus i(x-y)).]
  {PairPhaseX}{x}{fg}{local}{First phase parameter of a four-dimensional transporter.}{angle}
\dkDeclareSymbol[Phase parameter y in the closed interval from d to pi-d for a positive-determinant four-dimensional transporter. Together with x it gives rotational phases x+y and x-y, folded into the closed interval from zero to pi for the displacement cost.]
  {PairPhaseY}{y}{fg}{local}{Second phase parameter of a four-dimensional transporter.}{angle}
\dkDeclareSymbol[Scalar a = cos(x) in the four-dimensional phase calculation. Together with b = cos(y), it expresses the trace as 4ab and determines the reciprocal characteristic polynomial.]
  {PairPhaseCosineX}{a}{fg}{local}{Cosine of the first phase parameter in the trace calculation.}{scalar}
\dkDeclareSymbol[Scalar b = cos(y) in the four-dimensional phase calculation, paired with a = cos(x). These scalars are local to the characteristic-polynomial calculation.]
  {PairPhaseCosineY}{b}{fg}{local}{Cosine of the second phase parameter in the trace calculation.}{scalar}
\dkDeclareSymbol[Rotation parameter r: the two free orthogonal factors use planar rotation angles r+s and r-s. It appears in cos(y) = cos(d) cos(r) for two rotations.]
  {PairRotationR}{r}{fg}{local}{Half the sum of the two free planar rotation angles.}{angle}
\dkDeclareSymbol[Rotation parameter s: the two free orthogonal factors use planar rotation angles r+s and r-s. It appears in cos(x) = cos(h) cos(s) for two rotations. This scalar differs from the singular values and singleton costs.]
  {PairRotationS}{s}{fg}{local}{Half the difference of the two free planar rotation angles.}{angle}
\dkDeclareSymbol[Planar reflection J_2 = diag(1,-1). Every two-dimensional orthogonal matrix of determinant -1 can be written as a planar rotation times J_2.]
  {PairPlanarReflection}{J_2}{fg}{local}{Fixed planar reflection used to parametrize determinant -1 factors.}{unitary-operator}
\dkDeclareSymbol[Real planar rotation R_t = ((cos(t),-sin(t)),(sin(t),cos(t))). Its two displacement singular values are both 2 times abs(sin(t/2)). In the matching proof a rotation through t in the closed interval from zero to pi contributes 2 f_p(t).]
  {MatchingPlanarRotation}{\symglyph{PlanarRotationBase}_t}{fg}{local}{Planar rotation through the indicated angle.}{unitary-operator}
% Statistical application / operator provenance.
\dkDeclareSymbol{Sigma}{\Sigma}{old}{local}{Population covariance matrix.}{matrix}
\dkDeclareSymbol{SigmaHat}{\widehat\Sigma}{current}{local}{Estimated or perturbed covariance matrix.}{matrix}
\dkDeclareSymbol{YWSV}{V}{wanted}{local}{Orthonormal coordinates for the population eigenspace in the statistical application.}{isometric-embedding}
\dkDeclareSymbol{YWSVhat}{\widehat V}{trialframe}{local}{Estimated eigenspace frame in the statistical application.}{isometric-embedding}

% Figure-specific objects retain their own meanings when bare letters overlap.
\dkDeclareSymbol[Signed rotation angle phi of the displayed map on its invariant real two-plane. At theta = pi/4 the direct rotation turns its two principal planes by opposite signed angles of magnitude pi/4; the competitor turns its moving plane by pi/2 and fixes its other plane.]
  {PlaneRotationAngle}{\varphi}{sine}{local}{Signed rotation angle of the orthogonal map on the displayed invariant plane.}{angle}
\dkDeclareSymbol[Unit vector q in the displayed invariant plane. The map T is the direct rotation D or the competitor W in that panel; Tq is the image and (I - T)q = q - Tq is the displacement.]
  {FigureInput}{q}{ambient}{local}{Unit input to the orthogonal map shown in this invariant-plane panel.}{vector}
\dkDeclareSymbol[Image Tq of the panel unit input q under T = D or W. Both q and Tq lie on the unit circle; the green displacement arrow joins Tq to q, so it represents q - Tq.]
  {FigureImage}{Tq}{wantedvector}{local}{Image of the unit input under the orthogonal map shown in this panel.}{vector}
\dkDeclareSymbol[Displacement (I - T)q = q - Tq, drawn from the image Tq to the unit input q. On a plane rotated through phi its length is 2 |sin(phi/2)|; it vanishes on a fixed plane.]
  {DisplacementVector}{(I-T)q}{displacement}{local}{Difference between a unit input and its image under the displayed orthogonal map.}{vector}
\dkDeclareSymbol[First principal plane Pi_1 = span(e_0, e_3) for the source-target pair in the real four-dimensional counterexample. At theta = pi/4 the direct rotation acts on this plane by a quarter of pi, with displacement singular value sqrt(2 - sqrt(2)).]
  {CounterPrincipalOne}{\Pi_1}{fg}{glossary}{First principal plane in the real counterexample: the span of e0 and e3.}{linear-subspace}
\dkDeclareSymbol[Second principal plane Pi_2 = span(e_1, e_2) for the source-target pair in the real four-dimensional counterexample. The direct rotation acts by the same angle magnitude as on Pi_1, with the opposite sign in the displayed plane coordinates.]
  {CounterPrincipalTwo}{\Pi_2}{fg}{glossary}{Second principal plane in the real counterexample: the span of e1 and e2.}{linear-subspace}
\dkDeclareSymbol[Moving plane M = range(W - I) = span(e_0 + e_2, e_1 + e_3) for the competitor. W rotates it by 2 theta; at theta = pi/4 the two displacement singular values contributed by this plane are both sqrt(2).]
  {CounterMovingPlane}{M}{ambient}{glossary}{Moving plane of the competing transporter: the range of W minus the identity.}{linear-subspace}
\dkDeclareSymbol[Fixed plane Z = kernel(W - I) = span(e_0 - e_2, e_1 - e_3) for the competitor. W acts as the identity here, so this plane contributes two zero displacement singular values.]
  {CounterFixedPlane}{Z}{ambient}{glossary}{Fixed plane of the competing transporter: the kernel of W minus the identity.}{linear-subspace}
\dkDeclareSymbol[Direct-rotation displacement singular value a = 2 sin(pi/8) = sqrt(2 - sqrt(2)) at theta = pi/4. Each of its two principal planes contributes two copies, giving trace displacement 4a.]
  {CounterChordA}{a}{displacement}{glossary}{Direct-rotation displacement singular value in the counterexample: sqrt(2 - sqrt(2)).}{scalar}
\dkDeclareSymbol[Competitor displacement singular value b = 2 sin(pi/4) = sqrt(2) at theta = pi/4. Its moving plane contributes two copies; the fixed plane contributes two zeros, giving trace displacement 2b.]
  {CounterChordB}{b}{displacement}{glossary}{Competing-transporter displacement singular value in the counterexample: square root of 2.}{scalar}
\dkDeclareSymbol[Singular value s_j(I - T) of the full ambient displacement, ordered from largest to smallest. Each rotating real two-plane contributes two values 2 |sin(phi/2)|; fixed directions contribute zeros. Their sum is the trace cost.]
  {DisplacementSingular}{\symglyph{SingularValueBase}_j}{displacement}{local}{An individual singular value of the full displacement from the identity.}{scalar}
\dkDeclareSymbol[Edge weight g_ij = s_i + s_j - e_ij, the difference between the cost of two singletons and the paired cost. A matching has total cost C_0 minus the sum of its edge weights, with C_0 = sum_i s_i. Maximizing this sum gives the minimum transport cost; edges of negative weight may be omitted.]
  {MatchingEdgeWeight}{g_{ij}}{current}{glossary}{Edge weight: the cost of two singletons minus their paired cost.}{scalar-weight}
\dkDeclareSymbol[Total trace-norm displacement cost. For a transporter W, C(W) = ||I - W||_1, the sum of its displacement singular values. For a matching M, C(M) is the sum of its paired and unmatched singleton costs.]
  {TraceCost}{C}{norm}{glossary}{Trace-norm displacement cost of the indicated orthogonal transporter or matching.}{scalar-cost}

% Spectral values have their own key, distinct from an operator block.
\dkDeclareSymbol[Spectral parameter lambda of a self-adjoint operator. In a finite-dimensional eigenvalue list lambda_j(A) it denotes an individual eigenvalue; in A = integral lambda dE(lambda), it is the real integration variable.]{SpectralValue}{\lambda}{wantedspectrum}{local}{Real spectral parameter or individual eigenvalue of a self-adjoint operator.}{scalar}



% Conventional mathematical notation. Hover text applies even when ink is plain.
% These are local (non-glossary) unless promoted explicitly by the author.
% IdentityTwo and IdentityFour derive their I from Identity. Singular-value
% registrations derive s_j from SingularValueBase; SingletonCost has its own s.
\dkDeclareSymbol[Identity operator on the vector or Hilbert space in context. Its dimension is inferred from the expression and need not be finite.]{Identity}{I}{plain}{local}{Identity operator on the space in context.}{identity-operator}
\dkDeclareSymbol[Two-by-two identity matrix, acting on a two-dimensional coordinate space.]{IdentityTwo}{\symglyph{Identity}_2}{plain}{local}{Two-by-two identity matrix.}{identity-operator}
\dkDeclareSymbol[Four-by-four identity matrix, acting on a four-dimensional coordinate space.]{IdentityFour}{\symglyph{Identity}_4}{plain}{local}{Four-by-four identity matrix.}{identity-operator}
\dkDeclareSymbol[Two-by-two zero matrix. Unlike the bold zero vector, this is a linear operator/matrix with two rows and two columns.]{ZeroMatrixTwo}{0_2}{plain}{local}{Two-by-two zero matrix.}{matrix}
\dkDeclareSymbol[The constant pi, the ratio of the circumference of a circle to its diameter. Angles in this article use radians.]{PiConstant}{\pi}{plain}{local}{The constant pi (angles are in radians).}{constant}
\dkDeclareSymbol[The field of real numbers.]{RealField}{\mathbb R}{plain}{local}{The real numbers.}{field}
\dkDeclareSymbol[The field of complex numbers.]{ComplexField}{\mathbb C}{plain}{local}{The complex numbers.}{field}
\dkDeclareSymbol[Generic Hilbert space used in the two-projection discussion; distinct from the paper's ambient space symbol E in the sine-theta theorem.]{GeneralHilbertSpace}{\mathcal H}{plain}{local}{Generic Hilbert space in a local discussion.}{hilbert-space}
\dkDeclareSymbol[The spectrum sigma(T) of an operator T: the set of complex scalars for which T minus the scalar identity is not boundedly invertible. For self-adjoint T the spectrum is real.]{SpectrumSymbol}{\sigma}{plain}{local}{Spectrum of an operator.}{operator-notation}
\dkDeclareSymbol[Base letter for singular values, presently s. It is distinct from the base s for matching singleton costs.]{SingularValueBase}{s}{plain}{local}{Base letter for singular values.}{notation-base}
\dkDeclareSymbol[Ordered singular value s_j(T), equal to the j-th eigenvalue of |T| for a finite matrix; distinguish this from the matching singleton cost s_j.]{SingularValueS}{\symglyph{SingularValueBase}_j}{plain}{glossary}{The j-th singular value, under the currently printed s_j convention.}{scalar}
\dkDeclareSymbol[Alternate indexed singular-value notation sigma_j(T). The current manuscript also uses s_j(T); this registered alias is available for a later deliberate notation decision.]{SingularValueSigma}{\sigma_j}{plain}{local}{Indexed singular value under the sigma_j convention.}{scalar}
\dkDeclareSymbol[Operator norm ||T||_op, the supremum of ||Tx|| over unit vectors x.]{OperatorNormNotation}{\|\cdot\|_{\mathrm{op}}}{plain}{glossary}{Operator norm.}{norm}
\dkDeclareSymbol[Frobenius norm ||T||_F, the square root of the sum of squared matrix entries, equal to the finite-dimensional Hilbert--Schmidt norm.]{FrobeniusNormNotation}{\|\cdot\|_F}{plain}{glossary}{Frobenius norm.}{norm}
\dkDeclareSymbol[Trace norm ||T||_1, the sum of singular values when finite.]{TraceNormNotation}{\|\cdot\|_1}{plain}{glossary}{Trace (nuclear) norm.}{norm}
\dkDeclareSymbol[Hilbert--Schmidt norm ||T||_HS, the square root of the sum of squared singular values when finite.]{HilbertSchmidtNormNotation}{\|\cdot\|_{\mathrm{HS}}}{plain}{local}{Hilbert--Schmidt norm.}{norm}
\dkDeclareSymbol[Ky Fan j-norm, the sum of the largest j singular values.]{KyFanNormNotation}{\|\cdot\|_{(j)}}{plain}{glossary}{Ky Fan j-norm.}{norm}
\dkDeclareSymbol[Interval/exterior spectral separation predicate I_delta(B,C), specifying a central closed interval separated by delta from the other operator's spectrum.]{IntervalSeparation}{\mathcal I}{plain}{local}{Interval/exterior separation predicate.}{predicate}
\dkDeclareSymbol[Ordered spectral separation predicate O_delta(B,C), requiring one self-adjoint operator to lie below a spectral threshold and the other above it by at least delta.]{OrderedSeparation}{\mathcal O}{plain}{local}{Ordered spectral separation predicate.}{predicate}
\dkDeclareSymbol[Positive parameter epsilon controlling a small perturbation in the two-dimensional eigenvalue example.]{EigenvalueExampleEpsilon}{\varepsilon}{plain}{local}{Small positive parameter in the eigenvalue-example matrices.}{scalar}
\dkDeclareSymbol[Positive lower bound epsilon for the trial-coordinate map E_0, as in E_0* E_0 >= epsilon squared times the identity.]{FrameLowerBoundEpsilon}{\varepsilon}{plain}{local}{Lower-bound parameter for the trial-coordinate embedding.}{scalar}
\dkDeclareSymbol[Rotation angle vartheta of a planar orthogonal block in the matching-proof first-variation calculation.]{MatchingPlanarAngle}{\vartheta}{plain}{local}{Planar rotation angle in the matching-proof calculation.}{angle}


% Ordinary named mathematical operators already used throughout the manuscript.
% Keeping these here makes their spelling a single editorial choice.
\dkDeclareSymbol[Range of a linear operator: the set of its output vectors.]{RangeOperator}{\operatorname{ran}}{plain}{local}{Range (image) of an operator.}{operator-notation}
\dkDeclareSymbol[Kernel (null space) of a linear operator: inputs mapped to zero.]{KernelOperator}{\operatorname{ker}}{plain}{local}{Kernel (null space) of an operator.}{operator-notation}
\dkDeclareSymbol[Domain of an operator: the vectors on which the operator is defined, important for unbounded operators.]{DomainOperator}{\operatorname{dom}}{plain}{local}{Domain of an operator.}{operator-notation}
\DeclareRobustCommand{\DKRange}{\sym{RangeOperator}}
\DeclareRobustCommand{\DKKernel}{\sym{KernelOperator}}
\DeclareRobustCommand{\DKDom}{\sym{DomainOperator}}

% Norm-formatting macros: the operand changes by occurrence but the notation
% lives here. The outer hover describes the norm and does not add ink color.
\DeclareRobustCommand{\dkOperatorNorm}[1]{\like{OperatorNormNotation}{\|#1\|_{\mathrm{op}}}}
\DeclareRobustCommand{\dkTraceNorm}[1]{\like{TraceNormNotation}{\|#1\|_1}}
\DeclareRobustCommand{\dkFrobeniusNorm}[1]{\like{FrobeniusNormNotation}{\|#1\|_F}}
\DeclareRobustCommand{\dkHilbertSchmidtNorm}[1]{\like{HilbertSchmidtNormNotation}{\|#1\|_{\mathrm{HS}}}}
\DeclareRobustCommand{\dkKyFanNorm}[2]{\like{KyFanNormNotation}{\|#1\|_{(#2)}}}
\DeclareRobustCommand{\dkSchattenNorm}[2]{\like{SchattenNormNotation}{\|#1\|_{S_{#2}}}}

\dkDeclareSymbol[Schatten p-norm ||T||_{S_p}, obtained from the l^p-norm of the singular values and finite for operators in the Schatten p-class.]{SchattenNormNotation}{\|\cdot\|_{S_p}}{plain}{local}{Schatten p-norm.}{norm}


\dkDeclareSymbol[Positive modulus of an operator T: |T| = (T* T)^(1/2).]{OperatorModulusNotation}{|\cdot|}{plain}{glossary}{Positive modulus of an operator.}{operator-notation}
\dkDeclareSymbol[Linear span of the indicated vectors, i.e. the smallest linear subspace that contains them.]{SpanOperator}{\operatorname{span}}{plain}{local}{Linear span of vectors.}{operator-notation}
\dkDeclareSymbol[Diagonal or block-diagonal matrix/operator with the displayed entries or blocks.]{DiagonalOperator}{\operatorname{diag}}{plain}{local}{Diagonal or block-diagonal operator.}{linear-operator}
\dkDeclareSymbol[Trace of a trace-class operator, or ordinary matrix trace in finite dimensions.]{TraceOperator}{\operatorname{tr}}{plain}{local}{Trace of an operator.}{operator-notation}
\dkDeclareSymbol[Dimension of a vector space or subspace.]{DimensionOperator}{\dim}{plain}{local}{Dimension of a vector space.}{operator-notation}
\dkDeclareSymbol[Rank of a linear map or matrix: dimension of its range.]{RankOperator}{\operatorname{rank}}{plain}{local}{Rank of a linear map or matrix.}{operator-notation}
\dkDeclareSymbol[Distance between sets of spectral values: the infimum of pairwise distances.]{SpectralDistanceOperator}{\operatorname{dist}}{plain}{local}{Distance between two spectral sets.}{operator-notation}
\DeclareRobustCommand{\dkModulus}[1]{\like{OperatorModulusNotation}{|#1|}}

% Context-specific glyphs previously hidden in variants of other semantic keys.
% Keeping both the roles and the exact printed forms makes collisions explicit.
\dkDeclareSymbol[Nonzero vector x tested in the Rayleigh quotient, before imposing unit length.]{RayleighTestVector}{x}{trialvector}{local}{Test vector in the Rayleigh quotient.}{vector}
\dkDeclareSymbol[Positive cosine block C in the generic two-projection decomposition of Halmos; C and S commute and C squared plus S squared is the identity.]{HalmosCosineBlock}{C}{anglecos}{local}{Cosine block in the Halmos two-projection representation.}{linear-operator}
\dkDeclareSymbol[Positive sine block S in the generic two-projection decomposition of Halmos; C and S commute and C squared plus S squared is the identity.]{HalmosSineBlock}{S}{sineoperator}{local}{Sine block in the Halmos two-projection representation.}{linear-operator}
\dkDeclareSymbol[Riesz spectral projection P_Gamma defined by a contour Gamma that encloses selected spectral values.]{RieszProjection}{P_{\symglyph{RieszContour}}}{wantedprojection}{local}{Spectral projection associated with a contour Gamma.}{orthogonal-projection}
\dkDeclareSymbol[Historical Sylvester right-hand coefficient matrix A, in B X - X A = Y. Distinct from the current paper's reference operator A.]{HistoricalSylvesterRight}{A}{trialspectrum}{local}{Right-hand coefficient A in the historical Sylvester equation.}{self-adjoint-operator}
\dkDeclareSymbol[Scalar coefficient b in the Sylvester prototype (b-a)x=y, with separated scalar spectral values a and b.]{SylvesterScalarLeft}{b}{unwantedspectrum}{local}{Scalar left coefficient b in a Sylvester equation.}{scalar}
\dkDeclareSymbol[Scalar coefficient a in the Sylvester prototype (b-a)x=y.]{SylvesterScalarRight}{a}{trialspectrum}{local}{Scalar right coefficient a in a Sylvester equation.}{scalar}
\dkDeclareSymbol[Scalar unknown x in the Sylvester prototype (b-a)x=y.]{SylvesterScalarUnknown}{x}{sylvester}{local}{Scalar unknown x in a Sylvester equation.}{scalar}
\dkDeclareSymbol[Scalar right side y in the Sylvester prototype (b-a)x=y.]{SylvesterScalarRHS}{y}{resid}{local}{Scalar right side y in a Sylvester equation.}{scalar}
\dkDeclareSymbol[Source subspace K of a real Euclidean space in the matching proof; an orthogonal transporter sends K onto L.]{MatchingSourceSubspace}{K}{trial}{local}{Source subspace K in the matching proof.}{linear-subspace}
\dkDeclareSymbol[Target subspace L of a real Euclidean space in the matching proof; an orthogonal transporter sends K onto L.]{MatchingTargetSubspace}{L}{wanted}{local}{Target subspace L in the matching proof.}{linear-subspace}
\dkDeclareSymbol[Orthogonal projection P_L onto the matching proof's target subspace L.]{MatchingTargetProjection}{P_L}{wantedprojection}{local}{Orthogonal projection onto the matching target L.}{orthogonal-projection}
\dkDeclareSymbol[Real two-dimensional orthogonal group O(2): planar rotations and reflections.]{PlanarOrthogonalGroup}{O(2)}{fg}{local}{Orthogonal matrices of size two.}{group}
\dkDeclareSymbol[Base letter R used to denote planar rotation operators R_t in the matching proof.]{PlanarRotationBase}{R}{fg}{local}{Base letter for planar rotations.}{notation-base}
\dkDeclareSymbol[Planar rotation R_(r+s) through the sum of scalar phases r and s.]{MatchingPlanarRotationSum}{\symglyph{PlanarRotationBase}_{r+s}}{fg}{local}{Planar rotation through phase r+s.}{unitary-operator}
\dkDeclareSymbol[Planar rotation R_gamma through phase gamma in the block rotation calculation.]{MatchingPlanarRotationGamma}{\symglyph{PlanarRotationBase}_{\symglyph{MatchingPhaseGamma}}}{fg}{local}{Planar rotation through phase gamma.}{unitary-operator}

% Specific symbols in the historical and appendix examples. They retain their
% existing black-ink notation but receive locally accurate hover definitions.
\dkDeclareSymbol[Contour Gamma, a positively oriented closed curve in the resolvent set enclosing a selected isolated spectral cluster.]{RieszContour}{\Gamma}{plain}{local}{Contour used in the Riesz spectral-projection formula.}{contour}
\dkDeclareSymbol[Even integer nu in a tested Ky Fan comparison, with nu/2 principal-plane terms on the right side.]{KyFanEvenIndex}{\nu}{plain}{local}{Even integer in the principal-plane Ky Fan comparison.}{scalar}
\dkDeclareSymbol[Orthogonal projection Omega_k onto the k-th principal plane in the published displacement comparison.]{PrincipalPlaneProjectionBase}{\Omega}{plain}{local}{Projection onto the indicated principal plane.}{notation-base}
\dkDeclareSymbol[Scalar gap threshold gamma in the ordered Sylvester spectral separation B >= (gamma+delta)I and C <= gamma I.]{SylvesterGapThreshold}{\gamma}{plain}{local}{Spectral separation threshold in a Sylvester estimate.}{scalar}
\dkDeclareSymbol[Projection Omega_t from a net or sequence that converges strongly to the identity in the singular-value approximation lemma.]{StrongApproxProjection}{\Omega_t}{plain}{local}{Orthogonal projection converging strongly to the identity.}{orthogonal-projection}
\dkDeclareSymbol[First orthogonal projection Omega used in the two-block pinching inequality for a unitarily invariant norm.]{PinchingProjectionOne}{\Omega}{plain}{local}{First projection in the pinching inequality.}{orthogonal-projection}
\dkDeclareSymbol[Second orthogonal projection Upsilon used in the two-block pinching inequality for a unitarily invariant norm.]{PinchingProjectionTwo}{\Upsilon}{plain}{local}{Second projection in the pinching inequality.}{orthogonal-projection}
\dkDeclareSymbol[Finite-rank projection Gamma used to retain the image of a chosen finite-rank compression K Psi.]{CompressionProjectionImage}{\Gamma}{plain}{local}{Range projection in the finite-rank compression lemma.}{orthogonal-projection}
\dkDeclareSymbol[Finite-rank projection Psi used to select input vectors in a nearly maximizing singular-value compression.]{CompressionProjectionInput}{\Psi}{plain}{local}{Domain projection in the finite-rank compression lemma.}{orthogonal-projection}
\dkDeclareSymbol[Small positive tolerance eta in the finite-rank compression approximation estimate.]{CompressionTolerance}{\eta}{plain}{local}{Tolerance controlling the off-diagonal compression norm.}{scalar}

\dkDeclareSymbol[Rotation phase gamma in a planar block from the matching proof.]{MatchingPhaseGamma}{\gamma}{plain}{local}{Rotation phase gamma in the matching calculation.}{scalar}

% Two common indices with local mathematical meanings distinct from other deltas.
\dkDeclareSymbol[Kronecker delta: delta_ij is one when the indices i and j agree and zero otherwise. This is not the spectral-gap parameter delta.]{KroneckerDelta}{\delta}{plain}{local}{Kronecker delta: equals 1 for equal indices and 0 otherwise.}{scalar}
\dkDeclareSymbol[Sign variable epsilon taking the values -1 and +1 in the factorization over two phase choices.]{PhaseSignChoice}{\epsilon}{plain}{local}{Sign variable taking values minus one and plus one.}{scalar}

% Documented notation correspondence. The source column is a faithful
% transcription of each cited original context, NOT an asserted universal DK
% meaning. The original A_0 is deliberately listed for two distinct roles:
% its reference block and the comparison/trial operator. The source uses the
% same printed symbol for both; this paper is free to distinguish them.
\dkDeclareNotationSource{DK1970}{Davis and Kahan (1970), The Rotation of Eigenvectors by a Perturbation. III}{https://doi.org/10.1137/0707001}
\dkDeclareNotationSource{Seelmann2019}{Seelmann (2019), Semidefinite Perturbations in the Subspace Perturbation Problem}{https://arxiv.org/abs/1708.02463}
\dkDeclareNotationSource{Blechta2020}{Blechta (2020), Stability of Linear GMRES Convergence with Respect to Compact Perturbations}{https://arxiv.org/abs/2005.12960}
% Arguments: {key}{source ID}{source glyph}{section or equation}{role-specific qualification}
\dkDeclareNotationAlternative{A}{DK1970}{A}{\S1, (1.2)}{Reference self-adjoint operator.}
\dkDeclareNotationAlternative{H}{DK1970}{H}{\S1, (1.2)}{Additive perturbation.}
\dkDeclareNotationAlternative{Acurrent}{DK1970}{A+H}{\S1}{Perturbed operator.}
\dkDeclareNotationAlternative{E0}{DK1970}{E_0}{\S1, (1.1)}{Source isometric frame.}
\dkDeclareNotationAlternative{E1}{DK1970}{E_1}{\S1, (1.1)}{Reference complement frame.}
\dkDeclareNotationAlternative{F0}{DK1970}{F_0}{\S1}{Exact-target frame.}
\dkDeclareNotationAlternative{F1}{DK1970}{F_1}{\S1}{Exact-target complement frame.}
\dkDeclareNotationAlternative{ReferenceBlock0}{DK1970}{A_0}{\S1, (1.2)}{Reference source block.}
\dkDeclareNotationAlternative{ReferenceBlock1}{DK1970}{A_1}{\S1, (1.2)}{Reference complement block.}
\dkDeclareNotationAlternative{A0}{DK1970}{A_0}{\S1, residual discussion}{Residual comparison operator (the reference block in DK).}
% We deliberately leave Lambda0 and Lambda1 without DK source-glyph metadata:
% the available source transcription does not reliably preserve the source
% accents/decorations on the perturbed spectral-block symbols.
\dkDeclareNotationAlternative{Psource}{DK1970}{P}{\S1}{Reference/source projection.}
\dkDeclareNotationAlternative{Qtarget}{DK1970}{Q}{\S1}{Target projection.}
\dkDeclareNotationAlternative{Pexact}{DK1970}{Q}{\S1}{Exact-target projection.}
\dkDeclareNotationAlternative{B}{DK1970}{B}{\S1, (1.2)}{Off-diagonal coupling block.}
\dkDeclareNotationAlternative{Theta0}{DK1970}{\Theta_0}{\S1}{Source-side angle operator.}
\dkDeclareNotationAlternative{Theta1}{DK1970}{\Theta_1}{\S1}{Complementary angle operator.}
\dkDeclareNotationAlternative{Theta}{DK1970}{\Theta}{\S1}{Full angle operator.}
\dkDeclareNotationAlternative{R}{DK1970}{R}{\S1, residual discussion}{Trial residual.}
\dkDeclareNotationAlternative{Rrotation}{DK1970}{U}{\S1, (1.18)}{Canonical direct rotation.}
\dkDeclareNotationAlternative{Wcompetitor}{DK1970}{V}{\S4}{Competing unitary transporter.}
\dkDeclareNotationAlternative{delta}{DK1970}{\delta}{\S2}{Spectral separation parameter.}
% Contemporary alternatives are recorded separately and do not enter the DK
% correspondence table. Their citation locations must be individually checked.
\dkDeclareNotationAlternative{Eambient}{Seelmann2019}{\mathcal H}{p.~1}{Ambient Hilbert space.}
\dkDeclareNotationAlternative{SingularValueS}{Blechta2020}{\sigma_j}{Thm.~1.1}{Ordered singular values of a compact operator.}
